%% ---- begin preamble.tex
\documentclass[11pt]{amsart}
\usepackage[margin=1.15in]{geometry}
\usepackage{amssymb,amsmath,amsthm,mathtools}
\usepackage{needspace,microtype,etoolbox}
\usepackage{tikz-cd}
\usepackage{listings,xcolor,textcomp}
\definecolor{LeanKeyword}{HTML}{254E91}
\definecolor{LeanType}{HTML}{006B65}
\definecolor{LeanComment}{HTML}{526B46}
\definecolor{LeanString}{HTML}{8A4727}
\lstdefinelanguage{Lean}{
  sensitive=true,
  morekeywords={abbrev,by,class,def,end,exact,fun,import,in,instance,let,
    namespace,noncomputable,open,scoped,section,theorem,variable,where},
  morekeywords=[2]{Prop,Type,Set,NontriviallyNormedField,NormedAddCommGroup,
    NormedSpace,TopologicalSpace,ChartedSpace,ModelWithCorners,IsOpen,
    ContMDiffOn,ContMDiffVectorBundle},
  morecomment=[l]{--},morecomment=[s]{/-}{-/},morestring=[b]"
}
% Render Lean's Unicode notation with both the editor and PDF export engines.
\lstnewenvironment{leancode}{\lstset{
  language=Lean,
  basicstyle=\small\ttfamily,columns=fullflexible,keepspaces=true,
  keywordstyle=\color{LeanKeyword}\bfseries,
  keywordstyle=[2]\color{LeanType},
  commentstyle=\color{LeanComment}\itshape,
  stringstyle=\color{LeanString},
  showstringspaces=false,upquote=true,
  literate={∀}{{\ensuremath{\forall}}}1
           {γ}{{\ensuremath{\gamma}}}1
           {→}{{\ensuremath{\to}}}1
           {↦}{{\ensuremath{\mapsto}}}1
           {ℕ}{{\ensuremath{\mathbb N}}}1
           {∞}{{\ensuremath{\infty}}}1
           {ω}{{\ensuremath{\omega}}}1
           {×}{{\ensuremath{\times}}}1
           {⋀}{{\ensuremath{\bigwedge}}}1
           {𝓘}{{\ensuremath{\mathcal I}}}1
           {₁}{{\ensuremath{{}_1}}}1
           {₂}{{\ensuremath{{}_2}}}1
}}{}
\usepackage[colorlinks=true,linkcolor=blue,urlcolor=blue,citecolor=blue]{hyperref}
\hypersetup{pdftitle={Alternating bundles and the splitting of exterior powers},
 pdfauthor={Jack McCarthy},pdfsubject={Alternating bundles over nontrivially normed fields},
 pdfkeywords={alternating maps, vector bundles, analytic descent, splitting, exterior powers, nonarchimedean analysis, Lean}}
\newtheorem{theorem}{Theorem}[section]
\newtheorem{lemma}[theorem]{Lemma}
\newtheorem{proposition}[theorem]{Proposition}
\newtheorem{corollary}[theorem]{Corollary}
\theoremstyle{definition}
\newtheorem{definition}[theorem]{Definition}
\theoremstyle{remark}
\newtheorem{remark}[theorem]{Remark}
\numberwithin{equation}{section}
\newcommand{\N}{\mathbb N}
\newcommand{\Z}{\mathbb Z}
\newcommand{\F}{\mathbb F}
\newcommand{\cL}{\mathcal L}
\newcommand{\Alt}{\operatorname{Alt}}
\newcommand{\Mult}{\operatorname{Mult}}
\newcommand{\End}{\operatorname{End}}
\newcommand{\linf}{\ell^\infty}
\newcommand{\sgn}{\operatorname{sgn}}
\newcommand{\spn}{\operatorname{span}}
\newcommand{\id}{\mathrm{id}}
\newcommand{\abs}[1]{\lvert #1\rvert}
\newcommand{\norm}[1]{\lVert #1\rVert}
\newcommand{\cC}{\mathcal{C}}
\newcommand{\cD}{\mathcal{D}}
\newcommand{\coeff}{\operatorname{coeff}}
\newcommand{\ord}{\operatorname{ord}}
\newcommand{\rank}{\operatorname{rank}}
\newcommand{\type}{\operatorname{type}}
\newcommand{\sdim}{\operatorname{sdim}}
\newcommand{\incl}{\mathrm{incl}}
\newcommand{\Cop}{C^{\mathrm{op}}}
\newcommand{\Cvec}{C^{\mathrm{vec}}}
\newcommand{\Kh}{\widehat K}
\theoremstyle{plain}
\newtheorem{introtheorem}{Theorem}
\renewcommand{\theintrotheorem}{\Alph{introtheorem}}
\theoremstyle{definition}
\newtheorem{example}[theorem]{Example}
\newtheorem{conjecture}[theorem]{Conjecture}
\newtheorem{question}[theorem]{Question}
\theoremstyle{plain}
\newcommand{\alt}{\operatorname{alt}}
\newcommand{\sh}{\operatorname{sh}}
\newcommand{\ev}{\operatorname{ev}}
\newcommand{\ob}{\operatorname{ob}}
\newcommand{\Ext}{\operatorname{Ext}}
\newcommand{\Poly}{\operatorname{Poly}}
\newcommand{\hot}{\mathbin{\widehat\otimes}}
\newcommand{\hLambda}{\widehat\Lambda}
\newcommand{\hGamma}{\widehat\Gamma}
\newcommand{\cN}{\mathcal N}
\newcommand{\cS}{\mathcal S}
\newcommand{\cR}{\mathcal R}
\newcommand{\cG}{\mathcal G}
\newcommand{\cM}{\mathcal M}
\newcommand{\cP}{\mathcal P}
\newcommand{\fD}{\mathfrak D}
\newcommand{\fX}{\mathfrak X}
\newcommand{\Lsym}{\Lambda_{\mathrm{sym}}}
%% ---- end preamble.tex
%% ---- begin front.tex
\setlength{\emergencystretch}{1em}
\title{Alternating bundles and the splitting of exterior powers}
\author{Jack McCarthy}
\thanks{The original counterexample and its initial write-up were developed with Anthropic's
Claude. Subsequent research and exposition used Claude and OpenAI's Codex, and the organization of the paper was prepared
with Claude and Codex. AI assistants contributed to mathematical arguments and independent proof audits. The original Lean
formalization used Codex and Autoform. The remaining claims and the per-claim ledger were formalized with Anthropic's
Claude Code in October 2026. The splitting theory and the classification of good pairs, that is, the results of
Sections~\ref{sec:duality}--\ref{sec:questions} and Appendices~\ref{app:ordering}--\ref{app:families} not listed in
Section~\ref{sec:verification}, and Proposition~\ref{prop:tangent-values} were developed and written with Claude and are
not formalized. Section~\ref{sec:verification} records the verification scope.}
\makeatletter
% Keep the complete front matter on one page without changing the body layout.
\patchcmd{\@maketitle}{\global\topskip42\p@}{\global\topskip30\p@}{}%
  {\PackageError{charp}{Unable to adjust title spacing}{Check the amsart title definition.}}
\patchcmd{\@maketitle}{\dimen@34\p@}{\dimen@10\p@}{}%
  {\PackageError{charp}{Unable to adjust abstract spacing}{Check the amsart title definition.}}
\patchcmd{\@setauthors}{\@topsep30\p@}{\@topsep18\p@}{}%
  {\PackageError{charp}{Unable to adjust author spacing}{Check the amsart author definition.}}
\g@addto@macro\@setauthors{\vskip6pt{\centering\footnotesize October 2026\par}}
\makeatother
\begin{document}
\begin{abstract}
The alternating-map bifunctor over a nontrivially normed field acts analytically on analytic bundles when
precomposition is analytic. For normed $E$ and $D$, precomposition is analytic for all Banach values exactly
when one class vanishes: the class of the exact sequence of completed projective tensor and exterior powers of $E$,
pushed along an evaluation map. When $k!=0$, no splitting of this sequence for $E=K^k$ with the maximum norm commutes with
the permutation matrix of a single $p$-cycle, where $p=\operatorname{char}K$. Source, target
and parameter splittings, chiefly ordering splittings, make the class vanish; Ramsey forcing over finite residue fields
obstructs it. Over subfields of $\kappa((X))$, $\kappa$ finite, this decides several classes of pairs, and rigid Banach
spaces give vanishing classes that no splitting decomposition explains. Seventy-six statements, listed in
Section~\ref{sec:verification}, have Lean proofs of corresponding statements.
\end{abstract}
\maketitle
\setcounter{tocdepth}{1}
\begingroup
\small
\linespacing=8pt
\tableofcontents\removelastskip
\endgroup
\clearpage
%% ---- end front.tex
%% ---- begin sec-01-intro.tex
\section{Introduction}\label{sec:intro}

A functor on normed spaces extends fiberwise to $C^n$ vector bundles, with its natural transformations, when $C^n$
families of linear maps induce $C^n$ families (Theorem~\ref{thm:bundle-lifting}). We ask when the bifunctor
$(E,F)\mapsto\Alt^k(E;F)$ of alternating $k$-linear maps over a nontrivially normed field $K$ carries analytic bundles
to analytic bundles. For normed spaces $E$, $D$ and $F$ this concerns
$A^k_{E,D;F}\colon\cL(E,D)\to\cL(\Alt^k(D;F),\Alt^k(E;F))$, $A^k_{E,D;F}(u)(m)=m\circ(u,\dots,u)$. Let
$j_{E,F}\colon\Alt^k(E;F)\to\Mult^k(E;F)$ be the isometric inclusion. Composed with the induced isometric inclusion
$j_{E,F}\circ{-}\colon\cL(\Alt^k(D;F),\Alt^k(E;F))\to\cL(\Alt^k(D;F),\Mult^k(E;F))$, the map $A^k_{E,D;F}$ is the
diagonal of a bounded $k$-linear map (Proposition~\ref{prop:ambient-polynomial} with $g=\id_F$), and
Theorem~\ref{thm:coefficient-descent} decides when such a diagonal descends analytically through $j_{E,F}\circ{-}$. By Proposition~\ref{prop:homogeneous-lift}, $A^k_{E,D;F}$ is analytic at some
point if and only if it has a \emph{lift}: a bounded $k$-linear map $Q\colon\cL(E,D)^k\to\cL(\Alt^k(D;F),\Alt^k(E;F))$
with $Q(u,\dots,u)=A^k_{E,D;F}(u)$.

\subsection*{The splitting question}
Let $\cN^k_E$ be the kernel of the quotient map $\widehat\wedge^k_E$ from the completed projective tensor power of $E$
onto its completed projective exterior power (Section~\ref{sec:duality}):
\begin{equation}\label{intro:eq:source}
 0\longrightarrow\cN^k_E\longrightarrow\widehat\otimes^k_\pi E\xrightarrow{\ \widehat\wedge^k_E\ }
 \widehat\Lambda^k_\pi E\longrightarrow0.
\end{equation}
If $k!\ne0$ in $K$, antisymmetrization divided by $k!$ splits~\eqref{intro:eq:source} naturally in $E$ and solves every
lift problem (Propositions~\ref{can:prop:natural} and~\ref{prop:split-lift}(1)). If $k!=0$ in $K$, one $p$-cycle of
the coordinates of $K^k$, $p=\operatorname{char}K$, already forbids natural splittings (Theorem~\ref{can:thm:monomial},
Corollary~\ref{can:cor:p-cycle}). The question becomes when
the class $\ob_k(E,D)$ obtained from~\eqref{intro:eq:source} vanishes. It vanishes exactly when $(E,D)$ is \emph{good
for Banach values in degree $k$}: $A^k_{E,D;F}$ is analytic for every Banach space $F$; the pair is \emph{good in
degree $k$} if this holds for every normed space $F$ (Definition~\ref{fib:def:good}, Theorem~\ref{intro:thm:A}).

The class vanishes under a \emph{source splitting} $(\mathrm S^k_E)$, a bounded section of $\widehat\wedge^k_E$; a
\emph{target splitting} $(\mathrm R^k_D)$, a bounded retraction of the antisymmetrization $\alt^k_D$; or a
\emph{parameter splitting}, a bounded projection of $\widehat\otimes^k_\pi\cL(E,D)$ onto the closed span of the powers
$u^{\otimes k}$, which makes $\cL(E,D)$ \emph{polarizable in degree $k$}. Relative splittings, decompositions and
block assemblies combine them. The class can be nonzero: \emph{forced asymmetry} makes lifts asymmetric
(Section~\ref{sec:canonical}), and over a finite residue field \emph{Ramsey forcing} turns this asymmetry into a
triangular pattern of coefficients that norm estimates forbid (Section~\ref{sec:sharpness}).

\subsection*{Main results}


\begin{introtheorem}[descent and one obstruction class]\label{intro:thm:A}
Let $E$ and $D$ be normed spaces and $k\ge1$.
\begin{enumerate}
\item For every normed space $F$, $A^k_{E,D;F}$ is $C^\infty$ on $\cL(E,D)$; it is analytic at one point if and only if
it has a lift, and then it is a continuous polynomial.
\item The following are equivalent: \textup{(a)} $(E,D)$ is good for Banach values in degree $k$; \textup{(b)}
$u\mapsto\Lambda^k_\pi u$ has a bounded $k$-linear lift $\cL(E,D)^k\to\cL(\Lambda^k_\pi E,\widehat\Lambda^k_\pi D)$;
\textup{(c)} $\ob_k(E,D)=0$, where $\ob_k(E,D)$ is the pushout of the class of~\eqref{intro:eq:source} along the
restriction $\ev^k_{E,D}|_{\cN^k_E}\colon\cN^k_E\to\Mult^k_0(\cL(E,D);\widehat\Lambda^k_\pi D)$ of the evaluation map
$\ev^k_{E,D}(x_1\otimes\dots\otimes x_k)(u_1,\dots,u_k)=u_1x_1\wedge\dots\wedge u_kx_k$, with $\Mult^k_0$ the bounded
$k$-linear maps with zero diagonal; \textup{(d)} some bounded $k$-linear map
$\Psi\colon E^k\to\Mult^k_0(\cL(E,D);\widehat\Lambda^k_\pi D)$ satisfies
$\Psi(x_1,\dots,x_k)=\ev^k_{E,D}(x_1\otimes\dots\otimes x_k)$ whenever two of the $x_a$ coincide.
\item Let $(e_t)_{t\in\Sigma}$ be a linearly independent family in $E$ with dense span $V$, indexed by a totally ordered
set $\Sigma$; every normed space has one. The \emph{ordering lift} $Q_<(u)$, for $u\in\cL(E,D)^k$, is the linear map
$\Lambda^k_KV\to\widehat\Lambda^k_\pi D$ sending $e_{t_1}\wedge\dots\wedge e_{t_k}$, $t_1<\dots<t_k$, to
$u_1e_{t_1}\wedge\dots\wedge u_ke_{t_k}$. Composed with the canonical map $\Lambda^k_KV\to\Lambda^k_\pi E$, every lift
$Q$ as in \textup{(b)} is $Q_<+\Phi$ for a correction $\Phi$, $k$-linear in $u$ with $\Phi(v,\dots,v)=0$ for
$v\in\cL(E,D)$ (Theorem~\ref{ob:thm:ordering}). The pair $(E,D)$ is good for Banach values in degree $k$ if and only
if, for some such $\Phi$ and some $C$, $\norm{Q_<(u)y+\Phi(u)y}\le C\norm{u_1}\cdots\norm{u_k}\norm y_\pi$ for all $u$
and $y\in\Lambda^k_KV$, where $\norm y_\pi$ is the projective seminorm.
\end{enumerate}
\end{introtheorem}

Theorem~\ref{intro:thm:A} collects Lemma~\ref{lem:closed-smooth}, Propositions~\ref{prop:ambient-polynomial}
and~\ref{prop:homogeneous-lift}, Lemma~\ref{fib:lem:test}, Theorem~\ref{ob:thm:obstruction},
Corollary~\ref{ob:cor:adjacent} and Theorem~\ref{ob:thm:ordering}.

\begin{introtheorem}[canonical splittings exist exactly when $k!\ne0$]\label{intro:thm:B}
Let $k\ge2$.
\begin{enumerate}
\item If $k!\ne0$ in $K$, then $\alt^k_E/k!$, $\wedge^k_D/k!$ and the symmetrizer $(1/k!)\sum_{\sigma\in S_k}\sigma$ on
$\widehat\otimes^k_\pi P$ split the source, target and parameter sequences naturally in $E$, $D$ and $P$. Every pair is
good in degree $k$, and $\Alt^k$ is analytic on every hom space.
\item Suppose $k!=0$ in $K$. For every norm on $K^k$ with bounded coordinate functionals, no section of
$\wedge^k_{K^k}$, retraction of $\alt^k_{K^k}$ or projection onto the symmetric tensors commutes with the permutation
matrix of one $p$-cycle of the coordinates. If $E$ has $k$ linearly independent bounded functionals and
$\widehat\Lambda^k_\pi D\ne0$, no lift of $u\mapsto\Lambda^k_\pi u$ is symmetric, no section of $\widehat\wedge^k_E$
takes sign-isotypic values and no retraction of $\alt^k_D$ is sign-equivariant.
\item If $k!=0$ in $K$, then $A^k_{E,E;F}$ is analytic at no point for some Banach spaces $E$ and $F$, so $\Alt^k$ is
analytic on all hom spaces, or on those of Banach spaces, if and only if $k!\ne0$. If $K$ is moreover complete and not
spherically complete, $A^k_{E,D;K}$ is analytic at no point for some nonarchimedean Banach spaces $E$ and $D$.
\item If $K$ is complete and nonarchimedean, $\Alt^k(-;K)$ is analytic on all hom spaces of normed spaces if and only
if $k!\ne0$ in $K$ or $K$ is spherically complete.
\item Let a group act on a set $I$, and on $K^{(I)}$ by permuting the unit vectors. On the algebraic powers of
$K^{(I)}$, $\wedge^k$ has an equivariant section and $\alt^k$ an equivariant retraction if and only if the stabilizer of
each $k$-subset of $I$ induces on it a group of order nonzero in $K$; an equivariant projection onto the span of the
pure powers exists if and only if the stabilizer of each $k$-multiset on $I$ induces on its support a group of order
nonzero in $K$. The obstruction to an equivariant section is a family of Tate classes. One element inducing a
permutation of order $p$ on some $k$-subset forbids all three equivariant splittings, also on completed powers of normed
spaces with a compatible biorthogonal system of vectors and bounded functionals.
\end{enumerate}
\end{introtheorem}

Theorem~\ref{intro:thm:B} collects Propositions~\ref{can:prop:natural} and~\ref{prop:split-lift},
Theorem~\ref{can:thm:monomial}, Proposition~\ref{can:prop:asymmetry}, Theorem~\ref{thm:operator-obstructions},
Corollary~\ref{cor:classification}, Theorems~\ref{thm:scalar-classification} and~\ref{can:thm:permutation},
Corollary~\ref{can:cor:p-cycle} and Proposition~\ref{can:prop:tate}. Corollary~\ref{cor:bundle-cases} lists the bundle
settings in which fiberwise alternating maps form an analytic bifunctor; Proposition~\ref{prop:shear},
Corollary~\ref{cor:shear-bundle} and Proposition~\ref{prop:tangent-obstruction} give analytic atlases with nonanalytic
induced alternating or cotangent atlases, the latter over complete fields of characteristic $p>0$ that are not
spherically complete, in degrees $k\ge p$; over complete fields, ambient analytic forms keep pullback, wedge products
and exterior derivatives (Theorem~\ref{thm:ambient-forms}). Over nonarchimedean fields, with bounded polynomial laws
in place of bounded multilinear maps as Taylor coefficients, divided-power analytic maps compose
(Theorem~\ref{dp:thm:calculus}(2)), and in every characteristic the alternating and projective exterior power bundles of
divided-power analytic vector bundles are divided-power analytic (Corollary~\ref{dp:cor:bundles}). If $K$ has characteristic $p>0$, then in every degree at least $p$ there are bounded polynomial laws that are not
diagonals of bounded multilinear maps, even after an isometric embedding of their target
(Propositions~\ref{dp:prop:nonambient} and~\ref{dp:prop:degree-p}).

\begin{introtheorem}[how the class vanishes]\label{intro:thm:C}
Let $E$ and $D$ be normed spaces and $k\ge1$.
\begin{enumerate}
\item The pair $(E,D)$ is good for Banach values in degree $k$ if $(\mathrm S^k_E)$ or $(\mathrm R^k_D)$ holds, if
$\cL(E,D)$ is polarizable in degree $k$, if $(E,D)$ has a splitting decomposition of degree $k$, or if $(E,D)$ is an
$\ell^1$ block assembly of pairs with uniformly bounded lifts in all degrees $j\le k$. For decompositions into
$\mathrm S^{\le k}$-split and $\mathrm R^{\le k}$-split pieces, two pieces suffice.
\item If $E$ has a linearly independent family, indexed by a totally ordered set, with dense span $V$ such that the
strict triangle projection $\cP_<$ is bounded on $V\otimes_\pi V$, for instance one with a lower $r$-estimate for some
$r<2$, then $(\mathrm S^n_E)$ and $(\mathrm R^n_E)$ hold and $E$ is polarizable in degree $n$, for every $n\ge1$.
\item If $K$ is nonarchimedean, $E$ has a bounded linearly independent family indexed by a totally ordered set, with
dense span on which $E$ is $\ell^2$-dominated, and $D$ has an equivalent $r$-norm, $r\in(2,\infty)$, or an
equivalent nonarchimedean norm, then $(E,D)$ is good for Banach values in every degree.
\item Suppose $\operatorname{char}K=p>0$. The least degree, if any, in which $(E,D)$ is not good, or not good for
Banach values, is a power of $p$. If $E^*$ and $D^*$ are infinite-dimensional, the set of degrees in which $(E,D)$ is
good, or good for Banach values, is $\{1,\dots,p^h-1\}$ for some $h\ge1$ or $\{1,2,\dots\}$. For every normed space
$G$, each of $(\mathrm S^p_G)$, $(\mathrm R^p_G)$ and polarizability in degree $p$ implies the same condition in every
degree.
\item If $K$ is nonarchimedean and $E$ has a bounded linearly independent family $(e_t)_{t\in\Sigma}$, indexed by a
totally ordered set, with dense span $V$ and coordinate functionals $x\mapsto x(t)$ uniformly bounded on $V$, then the
class of~\eqref{intro:eq:source} pushes out to zero along every bounded linear map from $\cN^k_E$ to a Banach space with
an equivalent nonarchimedean norm. In particular $\ob_k(E,D)=0$ if $\widehat\Lambda^k_\pi D$ has such a norm.
\end{enumerate}
\end{introtheorem}

Theorem~\ref{intro:thm:C} collects Proposition~\ref{fib:prop:split}, Corollary~\ref{ob:cor:mechanisms},
Theorem~\ref{cmb:thm:decomposition}, Propositions~\ref{cmb:prop:normal-form} and~\ref{cmb:prop:blocks},
Theorem~\ref{fib:thm:lower}, Proposition~\ref{spl:prop:ordering}, Theorem~\ref{spl:thm:relative-ordering},
Corollaries~\ref{spl:cor:relative-ordering-na} and~\ref{fib:cor:degrees}, and Theorems~\ref{cmb:thm:initial-segments},
\ref{cmb:thm:column}, \ref{cmb:thm:polarizable-height} and~\ref{ob:thm:ultrametric}. A bounded retraction of $j_{E,F}$
gives analyticity for that value space $F$ only (Proposition~\ref{prop:split-lift}, Theorems~\ref{thm:spherical}
and~\ref{fib:thm:ultra}).

\begin{introtheorem}[forcing]\label{intro:thm:D}
Assume condition~\textup{(L)}: $\kappa$ is a finite field of characteristic $p$, $K$ is a subfield of $\kappa((X))$ on
which the absolute value $\abs a=\varrho^{\ord a}$, $0<\varrho<1$, is nontrivial, and $k\ge p$.
\begin{enumerate}
\item For $2\le r\le\infty$, the pair $(\ell^r(\N,K),\ell^r(\N,K))$, with $c_0(\N,K)$ for $r=\infty$, is not good for
Banach values in degree $k$.
\item If symmetric sequence spaces $E$ and $D$ form a pair good for Banach values in degree $k$, then for some $C$, the
\emph{product triangle bound} \textup{(PTB)}
\[
 \norm{\cP_<\bigl((ax)\otimes(by)\bigr)}_{D\otimes_\pi D}
 \le C\norm{\mu_a}_{E\to D}\norm x_E\norm{\mu_b}_{E\to D}\norm y_E
\]
holds for all finitely supported $a,x,b,y$, where $\mu_a$ is multiplication by $a$. If moreover
$\norm\xi_D\le C'\norm\xi_r$ for all finitely supported $\xi$ and some $r\in(2,\infty]$, then $E$ is $\ell^2$-dominated.
\item If $E$ has an equivalent nonarchimedean norm with $\dim_{\Kh}\widehat E=\infty$, then the pairs
$(E,c_0(\N,\Kh))$ and $(c_0(\N,K),E)$ are not good for Banach values in degree $k$.
\end{enumerate}
\end{introtheorem}

Theorem~\ref{intro:thm:D} collects Theorems~\ref{fib:thm:ell2} and~\ref{obs:thm:ptb},
Proposition~\ref{obs:prop:l2-domination}, and Theorem~\ref{fib:thm:local} with its proof, all proved in
Appendix~\ref{app:forcing}; Lemma~\ref{obs:lem:normalization} reduces condition~\textup{(L)} to fields dense in a
Laurent series field. Over every nontrivially normed field $K$ of characteristic~$2$, polynomial identities replace
Ramsey's theorem: for $E=c_0(\N,K)$, $\linf(\N,K)$ and $\ell^r(\N,K)$, $2<r<\infty$, the pair $(E,E)$ is not good for
Banach values in any degree $k\ge2$, and neither $(\mathrm S^k_E)$ nor $(\mathrm R^k_E)$ holds
(Theorem~\ref{obs:thm:char-two}). Over Laurent series fields over infinite fields and over densely valued fields, the
case $E=c_0(\N,K)$ for odd $p$ and the case $E=\ell^2(\N,K)$ for $p=2$ are open (Question~\ref{q:quest:residue}). In
characteristic~$2$, Proposition~\ref{fam2:prop:char-two} removes condition~\textup{(L)} from
Theorems~\ref{fam2:thm:staircase} and~\ref{fam2:thm:groups}(4), from Theorem~\ref{fam2:thm:tangent}(3) except its last
sentence, and from Corollary~\ref{fam2:cor:table}(2) for $E=c_0(\N,K)$ and for $E=\ell^{r_2}(\N,K)$, $r_2>2$.

\begin{introtheorem}[good pairs in classes]\label{intro:thm:E}
Let $\kappa$, $\varrho$ and $K$ be as in condition~\textup{(L)}, and let good mean good for Banach values.
\begin{enumerate}
\item Let $E$ be a symmetric sequence space in which $c_{00}$ is dense. The following are equivalent: $(E,E)$ is good
in some degree $k\ge p$; $(c_0,E)$ is good in some degree $k\ge p$; $\cP_<$ is bounded on $E\otimes_\pi E$;
$(\mathrm S^k_E)$ holds for some $k\ge p$; $(\mathrm R^k_E)$ holds for some $k\ge p$; $E$ is polarizable in some degree $k\ge p$; $(E,D)$ and $(D,E)$ are good in every
degree for every normed space $D$. If they hold, $E$ is $\ell^2$-dominated and $\varphi_E(m)^2\ge c\,m\ln m$ for some
$c>0$ and all $m\ge1$, with $\varphi_E(m)$ the norm of the indicator of $\{1,\dots,m\}$.
\item Let $E$ be as in \textup{(1)} and $D$ a symmetric sequence space with an equivalent $r$-norm, $r\in(2,\infty)$,
or an equivalent nonarchimedean norm. Then $(E,D)$ is good in some degree $k\ge p$ if and only if \textup{(PTB)}
holds, if and only if $E$ is $\ell^2$-dominated, if and only if $(E,D)$ is good in every degree.
\item For $1\le r_1,r_2\le\infty$, with $c_0$ in place of $\ell^\infty$, the pair $(\ell^{r_1},\ell^{r_2})$ is good in
some degree $k\ge p$ if and only if it is good in every degree, if and only if $\min(r_1,r_2)<2$ or $r_1=2<r_2$.
\item If $E$ has an equivalent nonarchimedean norm and $k\ge p$, then $(E,D)$ is good in degree $k$ for every normed
space $D$ if and only if $(D,E)$ is good in degree $k$ for every normed space $D$, if and only if $\dim_{\Kh}\widehat E<\infty$. If $K$ is complete and $k\ge p$, two
spaces with equivalent nonarchimedean norms form a pair good in degree $k$ if and only if one of them is
finite-dimensional.
\end{enumerate}
For the pairs in \textup{(1)}--\textup{(3)} and in the last sentence of \textup{(4)}, the good degrees form
$\{1,\dots,p-1\}$ or $\{1,2,\dots\}$, and every good pair has a splitting decomposition by absolute or relative
ordering splittings or by finite-coordinate splittings.
\end{introtheorem}

Theorem~\ref{intro:thm:E} collects Theorems~\ref{cls:thm:symmetric}, \ref{cls:thm:rnorm} and~\ref{fib:thm:local},
Corollaries~\ref{cls:cor:lr-table} and~\ref{cls:cor:dichotomy}, and Remark~\ref{cls:rem:mechanisms}. Items~(1)
and~(2), the last sentence for the pairs in them, and Theorem~\ref{intro:thm:D}(2) hold for spreading sequence spaces,
whose norms need only be invariant under increasing relabelings of the unit vectors (Theorem~\ref{cls:thm:spreading});
this class contains spaces with no equivalent symmetric norm (Example~\ref{cls:ex:garling}).

\begin{introtheorem}[goodness without splitting decompositions]\label{intro:thm:F}
\begin{enumerate}
\item Over every separable complete nonarchimedean field $K$ there are nonzero infinite-dimensional Banach spaces
$\fX_1,\fX_2,\dots$ with $\cL(\fX_\alpha,\fX_\alpha)=K\,\id$ and $\cL(\fX_\alpha,\fX_\beta)=0$ for $\alpha\ne\beta$.
\item Over $\kappa((X))$ with $\kappa$ finite, suitable choices in \textup{(1)} give $\fX=\fX_1$ with
$(\fX,\fX)$ good in every degree, while for every $k\ge p$ neither $(\mathrm S^k_{\fX})$ nor $(\mathrm R^k_{\fX})$
holds and $(\fX,\fX)$ has no splitting decomposition of degree $k$.
\item Over $\kappa((X))$ with $\kappa$ finite and these choices, the completed $\ell^1$-sum $E$ of the spaces $\fX_\alpha$
gives an $\ell^1$ block assembly $(E,E)$ that is good for Banach values in every degree and has no splitting
decomposition of any degree $k\ge p$, even with parameter pieces.
\end{enumerate}
\end{introtheorem}

Theorem~\ref{intro:thm:F} collects Theorems~\ref{bey:thm:rigid}, \ref{bey:thm:rigid-good} and~\ref{bey:thm:rigid-sum},
proved in Appendix~\ref{app:rigid}.

Bundles need goodness only for triples $(P;E,D)$, along analytic operator families on open subsets of the base model
$P$ (Section~\ref{sec:families}). Goodness of a triple reduces to homogeneous polynomial families and to familywise classes
(Theorems~\ref{fam2:thm:reduction} and~\ref{fam2:thm:class}). Under~\textup{(L)}, $(\ell^2(\N,K);E,E)$ with
$E=\ell^r(\N,K)$, $2<r<\infty$, is good in every degree, although neither is $\ell^2(\N,K)$ polarizable nor $(E,E)$
good for Banach values in degree $k$ (Theorems~\ref{fam2:thm:polarizable} and~\ref{fam2:thm:gamma-split},
Proposition~\ref{fam2:prop:fibers}, Corollary~\ref{cls:cor:symmetric-examples}, Theorem~\ref{fib:thm:ell2}). For a
structure group with an open subgroup that lies in a linear space of operators and contains a ball of it around the
identity, goodness is that of the linear space, not that of the closed linear span of the group
(Theorem~\ref{fam2:thm:groups}, Proposition~\ref{fam2:prop:span}).

\begin{introtheorem}[what bundles need]\label{intro:thm:G}
Let $P$ be a normed space and $k\ge1$.
\begin{enumerate}
\item Let $\cC$ be a class of Banach spaces closed under finite direct sums. Fiberwise alternating maps with Banach value
fibers preserve the analyticity of bundles with typical fibers in $\cC$ over every analytic manifold modeled on $P$ if
and only if $x\mapsto\Lambda^k_\pi u(x)$ is analytic for every analytic family $u$ of operators between spaces of $\cC$
on an open subset of $P$. For a single Banach fiber, families of invertible operators give the same condition.
\item The condition holds for all normed fibers if $P$ is polarizable in every degree $n\le k$, and over every $P$ for
fibers $E,D$ with $(E,\widehat D)$ good for Banach values in degree $k$. In characteristic $p$, the condition in all
degrees $j\le k$ is decided by the lifting constants of homogeneous polynomial families in the degrees $p^i\le k$.
Under~\textup{(L)}, it holds for all normed fibers over $P=\ell^r(\N,K)$, $1\le r<\infty$, if and only if $r<2$, and it
fails for some normed fibers over every $P$ with an equivalent nonarchimedean norm and $\dim_{\Kh}\widehat P=\infty$,
such as $c_0(\N,K)$.
\item Let $K$ be complete and $P$ a Banach space with $P\cong P^3$. The bundles $\widehat\Lambda^k_\pi TM$ are analytic
for every analytic manifold $M$ modeled on $P$ if and only if the condition holds for $E=D=P$. Under~\textup{(L)},
$c_0(\N,K)$ and $\ell^r(\N,K)$, $4<r<\infty$, model manifolds $M$ with a nonanalytic bundle $\widehat\Lambda^k_\pi TM$.
\end{enumerate}
\end{introtheorem}

Theorem~\ref{intro:thm:G} collects Proposition~\ref{fam2:prop:reductions}, Corollary~\ref{fam2:cor:bundles},
Theorems~\ref{fam2:thm:reduction}, \ref{fam2:thm:polarizable} and~\ref{fam2:thm:bases},
Proposition~\ref{fam2:prop:fibers} and Theorem~\ref{fam2:thm:tangent}.

Section~\ref{sec:questions} collects questions. Every theorem, proposition, lemma and corollary of the paper is proved
in the text, except Theorem~\ref{thm:ingleton} and Lemma~\ref{rig:lem:ramsey}, which are cited;
Section~\ref{sec:verification} lists the statements that also have Lean proofs of corresponding statements. Proofs
that would interrupt the text are collected in Appendix~\ref{app:deferred}.

\subsection*{Related work}
The bundle construction follows Lang's functors of class $C^n$~\cite[Ch.~III, \S4]{La} and Bourbaki's vector
functors~\cite[n\textsuperscript{o}~7.6]{Bou}. Bourbaki defines the alternating bifunctor when $K$ has characteristic
$0$ or the bundles have finite rank, and states that it is not known whether this restriction can be
removed~\cite[n\textsuperscript{o}~7.8]{Bou}; the paragraph \emph{Bourbaki's question} in Section~\ref{sec:geometry}
shows that it cannot over complete fields of characteristic $p>0$ that are not spherically complete. Gou\"ezel's
construction~\cite{G1,G2} covers characteristic zero; the finite-dimensional case in positive characteristic, noted
in~\cite{G1}, is Proposition~\ref{prop:split-lift}(2) over complete fields. Closed-range reflection of smoothness is
formalized in~\cite{PR}. The nonarchimedean extension mechanism is a norm-controlled form of spherical
completeness~\cite{In,Schi,Sch,Yuan}. Completed nonarchimedean exterior powers go back to Serre~\cite{Se} and
Gruson~\cite{Gr}, who do not study sections of the quotient map; we know no source on these splittings in positive
characteristic. The rigid spaces of Section~\ref{sec:beyond} transfer a construction of Kalton and Roberts~\cite{KR}.
Corollary~\ref{ob:cor:z-linear} detects the class by the distance of a $z$-linear map from the linear maps, as for
twisted sums~\cite{Ka,KP,CSC}.

\subsection*{Conventions}
$\cL(E,F)$ denotes bounded linear maps with the operator norm, $\Mult^k(E;F)$ bounded $k$-linear maps, and
$\Alt^k(E;F)$ those vanishing whenever two arguments coincide; this condition, not antisymmetry, defines alternation
in characteristic two. The degree-zero space is $F$. Products carry the maximum norm. $E$ denotes a source fiber, $D$
a target fiber, $F$ a value space (a \emph{target fiber} in Corollary~\ref{cor:bundle-cases} and
Appendix~\ref{app:laurent}), $G$ an intermediate space, and $P$ a base model or parameter space. \emph{Finite
coordinates} means a bounded linear isomorphism with $K^d$; over a complete field this is finite dimension. Hats
denote completions (Section~\ref{sec:duality}). Exponents are written $r$; $\ell^r$, lower $r$-estimates, $r$-norms
and $\ell^2$-domination of a family are defined in Section~\ref{sec:positive}, symmetric sequence spaces in
Definition~\ref{obs:def:sequence-spaces}. A finite field $\kappa$ has $\abs\kappa$ elements. Condition~\textup{(L)} is
fixed in Section~\ref{sec:sharpness}. A norm is \emph{nonarchimedean} if $\norm{x+y}\le\max(\norm x,\norm y)$. Every
normed field of characteristic $p>0$ is nonarchimedean: Frobenius gives $\abs{x+y}^{p^m}\le2\max(\abs x,\abs y)^{p^m}$; take roots
and let $m\to\infty$.

$C^n$ ($n\le\infty$) means iterated Fr\'echet differentiability; analyticity means local expansion in bounded
multilinear coefficients with a positive norm-controlled radius. On manifolds, $C^\omega$ also requires analytic
iterated derivatives, as in Mathlib~\cite{docs}. On open sets the two agree, even for incomplete ranges: the
derivative series converges in the completion to the Fr\'echet derivative, which exists in the operator space since
nearby points are points of analyticity.
%% ---- end sec-01-intro.tex
%% ---- begin sec-02-descent.tex
\section{Alternating bundles and analytic descent}\label{sec:bundles}

Throughout, $\mathsf{Vec}_K$ is the category of normed $K$-spaces, not necessarily complete, and bounded linear maps, with hom spaces $\cL(E,D)$. For $n\in\N\cup\{\infty,\omega\}$, $C^n$ is as in the conventions of Section~\ref{sec:intro}, and $C^0$ means continuous.

\subsection{The bundle construction}

Fix a $C^n$ manifold $M$ modeled on a normed space $P$. Write $\mathsf{VB}^n_K(M)$ for the category of $C^n$ normed vector bundles over $M$, with morphisms over $\id_M$ whose local \emph{operator-valued} expressions, not merely the maps of total spaces, are $C^n$. Thus a morphism from a bundle with fiber $E$ to one with fiber $F$ is locally a $C^n$ map $x\mapsto T(x)\in\cL(E,F)$. This condition is independent of the trivializations, since composition is bounded bilinear.

Let $\epsilon=(\epsilon_1,\ldots,\epsilon_r)$ specify a variance in each variable, and put $\mathbf C^\epsilon=\mathsf C_1\times\cdots\times\mathsf C_r$ with $\mathsf C_a=\mathsf{Vec}_K$ or $\mathsf{Vec}_K^{\mathrm{op}}$. A functor $\mathcal F\colon\mathbf C^\epsilon\to\mathsf{Vec}_K$ is $C^n$ if all its maps on products of hom spaces are jointly $C^n$. Let $\mathbf B^\epsilon_M$ denote the corresponding product of bundle categories, with the same variances. Natural transformations between such functors have bounded linear components, since their codomain is $\mathsf{Vec}_K$.

\begin{theorem}[lifting linear constructions to bundles]\label{thm:bundle-lifting}
For a fixed $C^n$ manifold $M$, fiberwise application defines, up to canonical natural isomorphism, a functor
\[
\begin{aligned}
 \mathcal B_M^n\colon
 \operatorname{Fun}^{C^n}(\mathbf C^\epsilon,\mathsf{Vec}_K)
 &\longrightarrow
 \operatorname{Fun}(\mathbf B^\epsilon_M,\mathsf{VB}^n_K(M)),\\
 \mathcal F&\longmapsto\mathcal F_M,\qquad
 \eta\longmapsto\eta_M.
\end{aligned}
\]
Here the category on the left is the full subcategory of the ordinary functor category on the functors regular on hom spaces. The resulting bundles, bundle morphisms, and natural transformations are given fiberwise by $\mathcal F$, its action on maps, and the components of $\eta$, respectively.

The same conclusion holds if regularity on whole hom spaces is replaced by preservation of $C^n$ families parametrized by open subsets of the model space $P$.
\end{theorem}

\begin{proof}
\emph{Step 1 (transitions).} Over a common trivializing cover of the input bundles, let the $C^n$ operator-valued maps $g^a_{ij}$ convert coordinates of the $a$th bundle from chart $j$ to chart $i$, so that $g^a_{ij}g^a_{j\ell}=g^a_{i\ell}$ and $g^a_{ii}=\id$. Let $\mathbf g_{ij}(x)$ be the tuple of arrows in $\mathbf C^\epsilon$ with entry $g^a_{ij}(x)$ in a covariant variable and $g^a_{ji}(x)$ in a contravariant one. By functoriality, $\theta_{ij}(x)=\mathcal F(\mathbf g_{ij}(x))$ satisfies $\theta_{ij}\theta_{j\ell}=\theta_{i\ell}$. The maps $\theta_{ij}$ are $C^n$ by composition, so they glue the products with typical fiber $\mathcal F(E_1,\ldots,E_r)$ into a $C^n$ vector bundle; its coordinate changes are $C^n$ because operator evaluation is bounded bilinear.

\emph{Step 2 (morphisms).} Applied to the local operator families of a tuple of bundle morphisms, reversed in the contravariant variables, $\mathcal F$ gives families compatible with the $\theta_{ij}$ by functoriality and $C^n$ by composition; this defines $\mathcal F_M$. The same argument for changes of trivialization identifies different choices canonically and preserves identities and composition.

\emph{Step 3 (natural transformations).} For $\eta\colon\mathcal F\Rightarrow\mathcal G$ and $\mathbf E=(E_1,\ldots,E_r)$, let $\eta_M$ be locally the constant operator $\eta_{\mathbf E}$. Naturality gives $\eta_{\mathbf E}\,\mathcal F(\mathbf g_{ij}(x))=\mathcal G(\mathbf g_{ij}(x))\,\eta_{\mathbf E}$, so $\eta_M$ is a bundle morphism, $C^n$ because its local expression is constant. The same identity for arbitrary input morphisms proves naturality of $\eta_M$, and identities and compositions of transformations are preserved fiberwise.

\emph{Step 4 (familywise version).} In a base chart, every transition and local morphism family is a $C^n$ family on an open subset of $P$, so preservation of such families gives the regularity used in Steps~1 and~2; Step~3 uses none.
\end{proof}

The familywise hypothesis, that $\mathcal F\circ\gamma$ be $C^n$ for each $C^n$ map $\gamma$ on an open subset of $P$ into $\prod_a\operatorname{Hom}_{\mathsf C_a}(E_a,D_a)$, is weaker than regularity on whole hom spaces. The theorem also holds for restricted categories of fibers containing the bundle morphisms; compare~\cite{G2,La} and~\cite[n\textsuperscript{os}~7.6.1--7.6.4]{Bou}.

\subsection{Alternating maps inside multilinear maps}

For normed spaces $E,F$, let $\Mult^k(E;F)$ be the normed space of bounded $k$-linear maps $E^k\to F$, and $\Alt^k(E;F)$ its subspace of alternating maps, those vanishing whenever two arguments agree. The inclusion
\[
 j_{E,F}\colon\Alt^k(E;F)\longrightarrow\Mult^k(E;F)
\]
is a linear isometry with closed range, since alternation is an intersection of closed conditions under evaluation at fixed tuples. Both constructions are contravariant in $E$ and covariant in $F$, and the inclusions form a natural transformation. If $F$ is complete, so are both spaces, since operator-norm Cauchy families have bounded pointwise limits satisfying the same identities; so the alternating bifunctor restricts to Banach spaces.

For $f\in\cL(E,D)$ and $g\in\cL(F,F')$, their joint action is
\[
 B(f,g)(m)(x_1,\ldots,x_k)
 =g\bigl(m(fx_1,\ldots,fx_k)\bigr).
\]
On alternating maps, write
\[
 A=A^k_{E,D;F}\colon\cL(E,D)\longrightarrow
 \cL(\Alt^k(D;F),\Alt^k(E;F)),\qquad A(f)=B(f,\id_F),
\]
for precomposition. Since $\norm{B(f,g)(m)}\le\norm g\,\norm m\,\norm f^k$, these actions are bounded operators.

\begin{proposition}[ambient polynomial action]\label{prop:ambient-polynomial}
The joint action on multilinear maps
\[
 B\colon\cL(E,D)\times\cL(F,F')
 \longrightarrow\cL\bigl(\Mult^k(D;F),\Mult^k(E;F')\bigr)
\]
is the diagonal of a bounded $(k+1)$-linear map, and hence is $C^\omega$. Its restriction to alternating inputs is analytic as a map into
\[
 \cL\bigl(\Alt^k(D;F),\Mult^k(E;F')\bigr),
\]
and all its values belong to the closed subspace of operators taking values in $\Alt^k(E;F')$.
\end{proposition}

\begin{proof}
For $v_i=(f_i,g_i)$ in the product of hom spaces, define
\[
 \Phi(v_0,\ldots,v_k)(m)(x_1,\ldots,x_k)
 =g_0\bigl(m(f_1x_1,\ldots,f_kx_k)\bigr).
\]
This map is multilinear in the $v_i$, has norm at most $1$, and has diagonal $B$. A diagonal of a bounded multilinear map has finite polynomial derivatives of every order, so it is $C^\omega$ even with incomplete range. Restricting $m$ to alternating maps is bounded linear and preserves analyticity. On the diagonal, two equal $x_i$ give two equal $fx_i$, so the output is alternating; the operator subspace is closed because $j_{E,F'}$ has closed range, applied pointwise. For $k=0$, the formula is postcomposition by $g_0$.
\end{proof}

For bundles with analytic transition families $g_{ij}$ and $h_{ij}$, the induced alternating transitions $m\mapsto h_{ij}(x)\circ m\circ(g_{ji}(x),\ldots,g_{ji}(x))$ are therefore analytic as maps into the multilinear operator space. Whether their power-series coefficients can be chosen in the alternating operator space is the descent problem of Section~\ref{sec:descent}.

\subsection{Analytic descent through closed subspaces}\label{sec:descent}

Throughout this subsection, $P,W,Z$ are normed $K$-spaces, not necessarily complete, and $j\colon W\to Z$ is a linear isometry with closed range; we identify $W$ with its image when convenient. If $U\subseteq P$ is open and $f\colon U\to W$, an ambient expansion at $x_0\in U$ means
\begin{equation}\label{eq:ambient-series}
 jf(x_0+h)=\sum_{n\ge0}b_n(h,\ldots,h),\qquad
 \sup_n\norm{b_n}R^n<\infty
\end{equation}
for some $R>0$, with bounded $n$-linear $b_n\colon P^n\to Z$ ($b_0$ a constant) and convergence for $\norm h<R$. For $j=\id_W$ this is, after shrinking $R$, the usual local definition of analyticity. For nonarchimedean spaces over a complete nonarchimedean field, this use of multilinear representatives agrees with the polynomial lifting norms of~\cite[7.14--7.19]{BGN}.

\begin{theorem}[the coefficient descent criterion]\label{thm:coefficient-descent}
Suppose \eqref{eq:ambient-series} holds. Then every diagonal $b_n(h,\ldots,h)$ belongs to $j(W)$. The map $f$ is analytic at $x_0$ if and only if there are bounded $n$-linear maps $q_n\colon P^n\to W$ and $r>0$ such that
\[
 jq_n(h,\ldots,h)=b_n(h,\ldots,h)\quad(n\ge0,\ h\in P),
 \qquad \sup_n\norm{q_n}r^n<\infty.
\]
Only the diagonals must agree; the original mixed values of $b_n$ need not lie in $j(W)$.
\end{theorem}

\begin{proof}
\emph{Step 1 (diagonals in $j(W)$).} Fix $h\in P$. On the scalar line $x_0+th$, the coefficients $b_n(h,\ldots,h)$ of \eqref{eq:ambient-series} satisfy a geometric bound, and $b_0=jf(x_0)$. If the diagonals of degree below $n$ lie in $j(W)$, then for small nonzero $t$
\[
 t^{-n}\Bigl(jf(x_0+th)-\sum_{m<n}t^m b_m(h,\ldots,h)\Bigr)\in j(W),
\]
and this vector tends to $b_n(h,\ldots,h)$ as $t\to0$, by the tail estimate in the proof of Lemma~\ref{lem:unique}; since $j(W)$ is closed, $b_n(h,\ldots,h)\in j(W)$.

\emph{Step 2 (the criterion).} If $f$ is analytic, apply $j$ to one of its expansions; on every scalar line, Lemma~\ref{lem:unique}, applied to the difference of the two series, gives the diagonal identities. Conversely, shrink $r$ below $R$ and let $\norm h<r$. The partial sums of $\sum_nq_n(h,\ldots,h)$ in $W$ have images converging to $jf(x_0+h)$, so they converge to $f(x_0+h)$ because $j$ is an isometry, and the bound on $\norm{q_n}r^n$ gives a positive radius.
\end{proof}

\begin{lemma}[smooth descent]\label{lem:closed-smooth}
For $n\in\mathbb N\cup\{\infty\}$, a map $f\colon U\to W$ is $C^n$ if and only if $jf$ is $C^n$.
\end{lemma}
\begin{proof}
Continuity reflects through an isometry. If $jf$ has Fr\'echet derivative $T$ at $x$, then $Th=\lim_{t\to0,\,t\ne0}t^{-1}(jf(x+th)-jf(x))\in j(W)$, so $T=jS$ for a unique bounded linear $S\colon P\to W$ of the same norm; the isometry transfers the remainder estimate, so $Df(x)=S$. Postcomposition by $j$ is a closed linear isometry $\cL(P,W)\to\cL(P,Z)$, since its image consists of the operators with values in $j(W)$, a closed condition under each evaluation. The same argument applied to the derivative maps proves every finite order by induction, hence $C^\infty$.
\end{proof}

By Proposition~\ref{prop:ambient-polynomial} and Lemma~\ref{lem:closed-smooth} with postcomposition by $j_{E,F'}$ as $j$, the joint action on alternating maps is $C^\infty$, and so is its case $g=\id_F$, the precomposition action $A^k_{E,D;F}$; this is the argument formalized in~\cite{PR}. With Theorem~\ref{thm:bundle-lifting}, this constructs alternating bundles and their morphisms at every finite order and at $C^\infty$; analyticity is the additional coefficient problem.

We call a finite sum of diagonals of bounded multilinear maps a \emph{continuous polynomial}.
\begin{proposition}[one lift, every base point]\label{prop:homogeneous-lift}
Let $k\ge1$, and suppose $a\colon H\to W$ satisfies
$ja(h)=\beta(h,\ldots,h)$ for a bounded $k$-linear map $\beta\colon H^k\to Z$. The following are equivalent:
\begin{enumerate}
\item $a$ is analytic at some point of $H$;
\item there is a bounded $k$-linear $Q\colon H^k\to W$ with $Q(h,\ldots,h)=a(h)$;
\item $a$ has a finite power-series expansion at every point, of infinite radius.
\end{enumerate}
Consequently each precomposition action $A^k_{E,D;F}$ is either a continuous polynomial on its whole domain or analytic nowhere.
\end{proposition}
\begin{proof}
An expansion at $h_0$, compared along $h_0+th$ with the polynomial $\beta(h_0+th,\ldots,h_0+th)$, has degree-$k$ diagonal $\beta(h,\ldots,h)$ by Lemma~\ref{lem:unique}; its degree-$k$ coefficient is the required $Q$. Conversely, expanding $Q(h_0+h,\ldots,h_0+h)$ by multilinearity gives a finite expansion whose degree-$i$ coefficient, the sum over the $i$ occupied slots, has norm at most $\binom ki\norm Q\norm{h_0}^{k-i}$. The last assertion follows from Proposition~\ref{prop:ambient-polynomial}.
\end{proof}

\paragraph{\textit{Retractions are splittings.}}
A bounded linear retraction $\rho$ of $j_{E,F}$ turns the representative $\beta(u_1,\ldots,u_k)(m)=m\circ(u_1,\ldots,u_k)$ of $j_{E,F}\circ A^k_{E,D;F}$ into the lift $\rho\circ\beta$ (Proposition~\ref{prop:split-lift}). Section~\ref{sec:duality} treats all Banach $F$ at once by one exact sequence and one class (Theorem~\ref{ob:thm:obstruction}).
%% ---- end sec-02-descent.tex
%% ---- begin sec-03-obstruction.tex
\section{Exterior powers and the obstruction class}\label{sec:duality}

\subsection{Projective tensor and exterior powers}\label{ob:ssec:powers}

For $k\ge1$ give the algebraic exterior power $\Lambda^k_KE$ the seminorm
\[
 \norm z_\pi=\inf_{z=\sum_jx_{j1}\wedge\cdots\wedge x_{jk}}
                 \sum_j\prod_{a=1}^k\norm{x_{ja}}.
\]
Its quotient by the zero-seminorm subspace is the normed space $\Lambda^k_\pi E$, with universal alternating map
$\omega_E(x_1,\ldots,x_k)=[x_1\wedge\cdots\wedge x_k]$. Put $\Lambda^0_\pi E=K$ and $\omega_E()=1$.

\begin{proposition}[representation]\label{prop:exterior}
There is a natural linear isometry
\[
 \Alt^k(E;F)\cong\cL(\Lambda^k_\pi E,F)
\]
for every normed $F$. Under it, precomposition by $u\colon E\to D$ is precomposition by
$\Lambda^k_\pi u\colon\Lambda^k_\pi E\to\Lambda^k_\pi D$.
\end{proposition}
\begin{proof}
The algebraic linearization $\widetilde m$ of a bounded alternating map $m$ satisfies
$\norm{\widetilde m(z)}\le\norm m\norm z_\pi$, by the definition of $\norm z_\pi$; so it kills the null subspace and
descends contractively. Conversely, $\norm{\omega_E(x_1,\ldots,x_k)}\le\prod_a\norm{x_a}$, so composing a bounded linear
map with $\omega_E$ is contractive. The two operations are inverse, hence isometric. Applying $u$ to each
wedge factor gives the intertwining identity and $\norm{\Lambda^k_\pi u}\le\norm u^k$ for $k\ge1$; in degree zero the
action is the identity on $K$.
\end{proof}

If $F$ is complete, the same correspondence holds with the completion $\widehat{\Lambda^k_\pi E}$.

The projective tensor product $X_1\otimes_\pi\dots\otimes_\pi X_n$ of normed spaces, and $\otimes^k_\pi E$ with $k$
factors $E$, are defined in the same way with $\otimes$ in place of $\wedge$. As in Proposition~\ref{prop:exterior},
bounded $n$-linear maps $m$ on $X_1\times\dots\times X_n$ correspond isometrically to bounded linear maps $m^\otimes$ on
$X_1\otimes_\pi\dots\otimes_\pi X_n$. For bounded operators $u_i$, $u_1\otimes\dots\otimes u_n$ has norm at most
$\prod_i\norm{u_i}$ and is multilinear in $(u_1,\dots,u_n)$; $u^{\otimes k}=u\otimes\dots\otimes u$. Hats denote
completions, maps extend to completions by continuity, and every normed space is a subspace of its completion. Let
\[
 \wedge^k_E\colon\otimes^k_\pi E\to\Lambda^k_\pi E,\qquad
 \operatorname{alt}^k_E\colon\Lambda^k_\pi E\to\otimes^k_\pi E
\]
be the maps with $x_1\otimes\dots\otimes x_k\mapsto x_1\wedge\dots\wedge x_k$ and
$x_1\wedge\dots\wedge x_k\mapsto\sum_{\sigma\in S_k}\sgn(\sigma)\,x_{\sigma(1)}\otimes\dots\otimes x_{\sigma(k)}$. The first is
$\omega_E^\otimes$, a contraction. The second is well defined in every characteristic, since for $x_a=x_b$ the
permutations $\sigma$ and $(a\ b)\circ\sigma$ give equal tensors with opposite signs; its norm is at most the real number
$k!$, and $\wedge^k_E\circ\operatorname{alt}^k_E=k!$. Both maps are natural:
$\wedge^k_D\circ u^{\otimes k}=\Lambda^k_\pi u\circ\wedge^k_E$ and
$u^{\otimes k}\circ\operatorname{alt}^k_E=\operatorname{alt}^k_D\circ\Lambda^k_\pi u$ for $u\in\cL(E,D)$.

\subsection{Universal and scalar tests}\label{ob:ssec:tests}

Write $Y^*=\cL(Y,K)$ and $J_Y(y)(\lambda)=\lambda(y)$ for the canonical contraction into the continuous bidual; no
point-separation or Hahn--Banach assumption is implicit.

\begin{theorem}[universal and scalar tests]\label{thm:duality}
Let $U$ be an open parameter domain and $W\colon U\to\cL(X,Y)$ a family. Define
\[
 R_F(t)(T)=T\circ W(t),\qquad
 R_F\colon U\to\cL(\cL(Y,F),\cL(X,F)).
\]
For power-series analyticity, or for any fixed class $C^n$, the following statements hold.
\begin{enumerate}
\item $W$ is regular if and only if $R_F$ is regular for every normed $F$; the single target $Y$ suffices.
\item $R_K$ is regular if and only if $J_YW$ is regular in $\cL(X,Y^{**})$, and if and only if
$R_{G^*}$ is regular for every normed $G$. Any one nonzero continuous dual $G^*$ suffices for the last test.
\item If $J_Y$ has a bounded linear left inverse, regularity of $R_K$ implies regularity of $W$.
If $J_F$ has such a left inverse, regularity of $R_K$ implies regularity of $R_F$.
\end{enumerate}
\end{theorem}
\begin{proof}
All identities below use bounded linear operations, which preserve each regularity class. Precomposition depends
bounded-linearly on $W$, and $W(t)=R_Y(t)(\id_Y)$; this proves~(1). Exchanging the arguments defines a linear isometry
\[
 \Theta\colon\cL(Y^*,X^*)\longrightarrow\cL(X,Y^{**}),
 \qquad (\Theta(B)x)(\lambda)=(B\lambda)(x),
\]
since both operator norms are the least $C$ bounding the displayed scalar by $C\norm x\norm\lambda$; its inverse exchanges
the arguments again, and $\Theta(R_K(t))=J_YW(t)$. The isometry
$\cL(X,G^*)\cong\cL(G,X^*)$, $T\mapsto[g\mapsto(x\mapsto T(x)(g))]$, identifies $R_{G^*}(t)$ with $S\mapsto R_K(t)\circ S$,
which gives regularity for every dual target. Conversely, choose $v\in G^*$ and $g\in G$ with $v(g)=1$; the injection
$\alpha\mapsto v\alpha$ and evaluation at $g$ compress $R_{G^*}$ to $R_K$. This proves~(2). For~(3), if $P_YJ_Y=\id_Y$,
then $W=P_Y\Theta(R_K)$; if $P_FJ_F=\id_F$, then $R_F(t)(T)=P_F\circ R_{F^{**}}(t)(J_F\circ T)$.
\end{proof}

\subsection{Good pairs and one test target}\label{ob:ssec:good}

A \emph{Banach space} is a complete normed $K$-space.

\begin{definition}\label{fib:def:good}
A pair $(E,D)$ of normed spaces is \emph{good in degree $k$} if $A^k_{E,D;F}$ is analytic for every normed $F$, and
\emph{good for Banach values in degree $k$} if this holds for every Banach space $F$.
\end{definition}

\begin{lemma}[one test target]\label{fib:lem:test}
The pair $(E,D)$ is good in degree $k$ if and only if $u\mapsto\Lambda^k_\pi u$ is analytic at some point of $\cL(E,D)$,
and if and only if there is a bounded $k$\nobreakdash-linear $Q\colon\cL(E,D)^k\to\cL(\Lambda^k_\pi E,\Lambda^k_\pi D)$
with $Q(u,\dots,u)=\Lambda^k_\pi u$. The same holds for Banach values, with $\cL(\Lambda^k_\pi E,\widehat\Lambda^k_\pi D)$
in place of $\cL(\Lambda^k_\pi E,\Lambda^k_\pi D)$, also as the space in which $u\mapsto\Lambda^k_\pi u$ is analytic.
\end{lemma}

\begin{proof}
Apply Theorem~\ref{thm:duality}(1) with $X=\Lambda^k_\pi E$, $W(u)=\Lambda^k_\pi u$, and $Y=\Lambda^k_\pi D$,
respectively $Y=\widehat\Lambda^k_\pi D$. By Proposition~\ref{prop:exterior} and its completed form, $R_F$ is
$A^k_{E,D;F}$ for every normed $F$, respectively every Banach $F$, and the target $Y$ suffices. Proposition~\ref{prop:homogeneous-lift}
applies to $W$, with the closed isometry $T\mapsto T\circ\omega_E$ of $\cL(\Lambda^k_\pi E,Y)$ onto
$\Alt^k(E;Y)\subseteq\Mult^k(E;Y)$ and the ambient map $(u_1,\dots,u_k)\mapsto\omega_D\circ(u_1,\dots,u_k)$.
\end{proof}

We use two splitting conditions on normed spaces:
\begin{enumerate}
\item[$(\mathrm S^k_E)$] there is a bounded linear $s\colon\Lambda^k_\pi E\to\widehat\otimes^k_\pi E$ with
$\widehat\wedge^k_E\circ s=\id$;
\item[$(\mathrm R^k_D)$] there is a bounded linear $\rho\colon\otimes^k_\pi D\to\widehat\Lambda^k_\pi D$ with
$\rho\circ\operatorname{alt}^k_D=\id$.
\end{enumerate}
They hold \emph{without completion} if $s$ takes values in $\otimes^k_\pi E$, respectively $\rho$ in
$\Lambda^k_\pi D$. Both hold without completion if $k!\ne0$ in $K$, with $s=\operatorname{alt}^k_E/k!$ and
$\rho=\wedge^k_D/k!$, and if the space has finite continuous coordinates $(e_i,e^*_i)$, with
\[
 \begin{gathered}
 s(x_1\wedge\dots\wedge x_k)=\sum_{i_1<\dots<i_k}\det\bigl(e^*_{i_a}(x_b)\bigr)_{a,b}\,
 e_{i_1}\otimes\dots\otimes e_{i_k},\\
 \rho(y_1\otimes\dots\otimes y_k)=\sum_{i_1<\dots<i_k}\prod_a e^*_{i_a}(y_a)\,e_{i_1}\wedge\dots\wedge e_{i_k}.
 \end{gathered}
\]

\begin{proposition}[splittings]\label{fib:prop:split}
Let $E$ and $D$ be normed spaces.
\begin{enumerate}
\item If $(\mathrm S^k_E)$ holds, then $(E,D)$ is good for Banach values in degree $k$ for every normed $D$:
$Q(u_1,\dots,u_k)=\widehat\wedge^k_D\circ(u_1\otimes\dots\otimes u_k)\circ s$ is a lift as in
Lemma~\ref{fib:lem:test}, of norm at most $\norm s$.
\item If $(\mathrm R^k_D)$ holds, then $(E,D)$ is good for Banach values in degree $k$ for every normed $E$, with the
lift $Q(u_1,\dots,u_k)=\rho\circ(u_1\otimes\dots\otimes u_k)\circ\operatorname{alt}^k_E$ of norm at most $k!\norm\rho$.
\item If the condition holds without completion, the pairs in~(1), respectively~(2), are good in degree $k$.
\item $(\mathrm S^k_E)$ holds if and only if the inclusions $j_{E,F}\colon\Alt^k(E;F)\to\Mult^k(E;F)$, for all Banach
spaces $F$, have bounded linear retractions that commute with postcomposition by bounded linear maps $F\to F'$. It
holds without completion if and only if such retractions exist for all normed $F$.
\end{enumerate}
\end{proposition}

The proof is in Appendix~\ref{app:deferred}.

\subsection{Three exact sequences}\label{ob:ssec:sequences}

A sequence $\mathcal Y\colon0\to Y_1\xrightarrow{i}Y_2\xrightarrow{q}Y_3\to0$ of Banach spaces and bounded linear maps is
\emph{exact} if $i$ is injective, $i(Y_1)=\ker q$, and for a constant $C$, $\norm y\le C\norm{iy}$ on $Y_1$ and every
$z\in Y_3$ is $qy$ for some $y$ with $\norm y\le C\norm z$. It \emph{splits} if $q$ has a bounded linear right inverse.
Exact sequences $\mathcal Y,\mathcal Y''$ with the same end terms are \emph{equivalent} if a bounded linear
$T\colon Y_2\to Y''_2$ satisfies $Ti=i''$ and $q''T=q$; such a $T$ is bijective with bounded inverse, since $q$ is open and
$i''$ is bounded below, so this is an equivalence relation, and it preserves splitting. $\Ext(Y_3,Y_1)$ is the set of
classes $[\mathcal Y]$, and $0$ is the class of $0\to Y_1\to Y_1\oplus Y_3\to Y_3\to0$, with the inclusion, the projection
and the norm $\norm{y_1}+\norm{y_3}$. Then $[\mathcal Y]=0$ if and only if $\mathcal Y$ splits: a left inverse $\lambda$ of
$i$ (Lemma~\ref{ob:lem:split-criteria}(2)) gives the equivalence $y\mapsto(\lambda y,qy)$. For a bounded linear
$\alpha\colon Y_1\to Y'$ into a Banach space, the \emph{pushout} $\alpha_*\mathcal Y$ is
$0\to Y'\xrightarrow{i_\alpha}Y_\alpha\xrightarrow{q_\alpha}Y_3\to0$, where $Y_\alpha$ is the quotient of $Y'\oplus Y_2$,
normed by $\norm{y'}+\norm y$, by $\{(\alpha y,-iy):y\in Y_1\}$, which is closed because $y\mapsto(\alpha y,-iy)$ is
bounded below on the complete space $Y_1$; $i_\alpha y'=[(y',0)]$ and $q_\alpha[(y',y)]=qy$. Put
$\alpha_*[\mathcal Y]=[\alpha_*\mathcal Y]$; an equivalence $T$ induces the equivalence $[(y',y)]\mapsto[(y',Ty)]$, so
this depends only on $[\mathcal Y]$.

\begin{lemma}[splitting criteria]\label{ob:lem:split-criteria}
Let $\mathcal Y$ and $\mathcal Y'\colon0\to Y'_1\xrightarrow{i'}Y'_2\xrightarrow{q'}Y'_3\to0$ be exact.
\begin{enumerate}
\item Let $\alpha_a\colon Y_a\to Y'_a$, $a=1,2,3$, be bounded linear maps with $\alpha_2i=i'\alpha_1$ and
$q'\alpha_2=\alpha_3q$. Then $\alpha_1=\tilde\alpha_1i$ for a bounded linear $\tilde\alpha_1\colon Y_2\to Y'_1$ if and only
if $\alpha_3=q'\tilde\alpha_3$ for a bounded linear $\tilde\alpha_3\colon Y_3\to Y'_2$, and the two can be chosen with
$\alpha_2=i'\tilde\alpha_1+\tilde\alpha_3q$.
\item $\mathcal Y$ splits if and only if $i$ has a bounded linear left inverse.
\item For a bounded linear $\alpha\colon Y_1\to Y'$ into a Banach space, $\alpha_*\mathcal Y$ is exact, and
$\alpha_*[\mathcal Y]=0$ if and only if $\alpha$ extends to $Y_2$, that is, $\alpha=\tilde\alpha i$ for a bounded linear
$\tilde\alpha\colon Y_2\to Y'$.
\end{enumerate}
\end{lemma}

The proof is in Appendix~\ref{app:deferred}.

\begin{lemma}\label{ob:lem:metric-surjection}
For a normed space $E$ and $k\ge1$, $\wedge^k_E$ is a metric surjection: $\norm y=\inf\{\norm z:\wedge^k_Ez=y\}$ for
$y\in\Lambda^k_\pi E$. Its extension $\widehat\wedge^k_E\colon\widehat\otimes^k_\pi E\to\widehat\Lambda^k_\pi E$ is
surjective and again a metric surjection.
\end{lemma}

The proof is in Appendix~\ref{app:deferred}.

\begin{definition}\label{ob:def:sequences}
Let $E$, $D$, $P$ be normed spaces and $k,n\ge1$.
\begin{enumerate}
\item The \emph{source sequence} of $E$ in degree $k$ is the exact sequence (Lemma~\ref{ob:lem:metric-surjection})
\begin{equation}\label{ob:eq:source}
 \cS^k_E\colon\quad0\to\cN^k_E\to\widehat\otimes^k_\pi E\xrightarrow{\;\widehat\wedge^k_E\;}\widehat\Lambda^k_\pi E\to0,
 \qquad \cN^k_E=\ker\widehat\wedge^k_E,
\end{equation}
whose first map is the inclusion.
\item The \emph{target sequence} of $D$ in degree $k$ is
$\cR^k_D\colon0\to\widehat\Lambda^k_\pi D\xrightarrow{\alt^k_D}\widehat\otimes^k_\pi D\to
\widehat\otimes^k_\pi D/\overline{\alt^k_D(\widehat\Lambda^k_\pi D)}\to0$, with the quotient map last. It is exact if and
only if $\alt^k_D$ is bounded below on $\widehat\Lambda^k_\pi D$.
\item Let $\widehat\Gamma^n_\pi P$ be the closed linear span in $\widehat\otimes^n_\pi P$ of the pure powers
$h^{\otimes n}=h\otimes\dots\otimes h$, $h\in P$. The \emph{parameter sequence} of $P$ in degree $n$ is the exact sequence
$\cG^n_P\colon0\to\widehat\Gamma^n_\pi P\to\widehat\otimes^n_\pi P\to\widehat\otimes^n_\pi P/\widehat\Gamma^n_\pi P\to0$.
The space $P$ is \emph{polarizable in degree $n$}, written $(\Gamma^n_P)$, if $\cG^n_P$ splits, that is
(Lemma~\ref{ob:lem:split-criteria}(2)), if there is a bounded linear projection $R_n$ of $\widehat\otimes^n_\pi P$ onto
$\widehat\Gamma^n_\pi P$. Let $g_n(P)$ be the infimum of the norms of such projections.
\item For a Banach space $Y$, $\Mult^n_0(P;Y)$ is the closed subspace of $\Mult^n(P;Y)$ of the maps $\Phi$ with zero
diagonal, $\Phi(h,\dots,h)=0$ for $h\in P$, and $\Poly^n(P;Y)$ is the quotient of $\Mult^n(P;Y)$ by $\Mult^n_0(P;Y)$,
with the quotient norm; through $\Phi\mapsto(h\mapsto\Phi(h,\dots,h))$ its elements correspond to the diagonals of bounded
$n$-linear maps $P^n\to Y$, the $n$-homogeneous continuous polynomials of Section~\ref{sec:descent}. The resulting exact sequence is
\[
 \cM^n(P;Y)\colon\quad0\to\Mult^n_0(P;Y)\to\Mult^n(P;Y)\to\Poly^n(P;Y)\to0.
\]
\end{enumerate}
\end{definition}

\begin{lemma}\label{ob:lem:source}
Condition $(\mathrm S^k_E)$ holds if and only if $\cS^k_E$ in~\eqref{ob:eq:source} splits, and if and only if $\cN^k_E$ is complemented in
$\widehat\otimes^k_\pi E$. A section $s$ gives the projection $\id-s\,\widehat\wedge^k_E$ onto $\cN^k_E$, of norm at most
$1+\norm s$; a bounded projection $\Pi$ onto $\cN^k_E$ gives the section $s(\widehat\wedge^k_Ez)=(\id-\Pi)z$, of norm at
most $\norm{\id-\Pi}$.
\end{lemma}

\begin{proof}
A section on $\Lambda^k_\pi E$ extends by continuity to a right inverse of $\widehat\wedge^k_E$, and a right inverse
restricts to a section. The map $\id-s\,\widehat\wedge^k_E$ is the identity on $\cN^k_E$ and is killed by
$\widehat\wedge^k_E$. Conversely, $\id-\Pi$ vanishes on $\cN^k_E$, so the formula defines $s$, and
Lemma~\ref{ob:lem:metric-surjection} bounds its norm.
\end{proof}

\begin{lemma}\label{ob:lem:target}
For a normed space $D$ and $k\ge1$ the following are equivalent: \textup{(1)} $(\mathrm R^k_D)$; \textup{(2)} $\cR^k_D$ is
exact and splits; \textup{(3)} $\omega_D=\mathcal A_D\eta$ for some $\eta\in\Mult^k(D;\widehat\Lambda^k_\pi D)$, where
$(\mathcal A_D\eta)(y_1,\dots,y_k)=\sum_{\sigma\in S_k}\sgn(\sigma)\,\eta(y_{\sigma(1)},\dots,y_{\sigma(k)})$ is
unnormalized antisymmetrization.
\end{lemma}

The proof is in Appendix~\ref{app:deferred}.

\begin{lemma}\label{ob:lem:parameter}
For a normed space $P$ and $n\ge1$ the following are equivalent.
\begin{enumerate}
\item $P$ is polarizable in degree $n$.
\item The power map $P\to\widehat\Gamma^n_\pi P$, $h\mapsto h^{\otimes n}$, is the diagonal of a bounded $n$-linear map
$P^n\to\widehat\Gamma^n_\pi P$.
\item For every closed subspace $W$ of a Banach space $Z$, every bounded $n$-linear $\Phi\colon P^n\to Z$ with
$\Phi(h,\dots,h)\in W$ for all $h\in P$, and every $\delta>0$, there is a bounded $n$-linear $Q\colon P^n\to W$ with
$Q(h,\dots,h)=\Phi(h,\dots,h)$ for all $h$ and $\norm Q\le(g_n(P)+\delta)\norm\Phi$.
\end{enumerate}
Under these conditions $\cM^n(P;Y)$ splits for every Banach space $Y$, by right inverses that commute with postcomposition
by bounded linear maps $Y\to Y'$.
\end{lemma}

The proof is in Appendix~\ref{app:deferred}.

\begin{definition}\label{ob:def:split-operators}
Let $w\colon E\to G$ be a bounded linear map of normed spaces and $k\ge1$. Then $w$ is \emph{$\mathrm S^k$-split} if
$\widehat\wedge^k_G\circ s=\Lambda^k_\pi w$ for a bounded linear $s\colon\Lambda^k_\pi E\to\widehat\otimes^k_\pi G$, and
\emph{$\mathrm R^k$-split} if $\rho\circ\alt^k_E=\Lambda^k_\pi w$ for a bounded linear
$\rho\colon\otimes^k_\pi E\to\widehat\Lambda^k_\pi G$; it is split \emph{without completion} if $s$ takes values in
$\otimes^k_\pi G$, respectively $\rho$ in $\Lambda^k_\pi G$. It is \emph{$\mathrm S^{\le k}$-split}
(\emph{$\mathrm R^{\le k}$-split}) if it is $\mathrm S^j$-split ($\mathrm R^j$-split) for every $j\le k$.
\end{definition}

Thus $(\mathrm S^k_E)$ says that $\id_E$ is $\mathrm S^k$-split, and $(\mathrm R^k_D)$ that $\id_D$ is
$\mathrm R^k$-split. Every $w$ is $\mathrm S^1$- and $\mathrm R^1$-split, and $\mathrm S^j$- and $\mathrm R^j$-split
whenever $j!\ne0$ in $K$, with $s=w^{\otimes j}\alt^j_E/j!$ and $\rho=\wedge^j_Gw^{\otimes j}/j!$, since
$\wedge^j_Gw^{\otimes j}\alt^j_E=\Lambda^j_\pi w\,\wedge^j_E\alt^j_E=j!\,\Lambda^j_\pi w$.

\begin{proposition}[split operators]\label{ob:prop:split-operators}
Let $w\colon E\to G$ be a bounded linear map of normed spaces and $k\ge1$.
\begin{enumerate}
\item The map $w^{\otimes k}$ sends $\cN^k_E$ into $\cN^k_G$, and the following are equivalent: (a) $w$ is
$\mathrm S^k$-split; (b) $\widehat\Lambda^k_\pi w=\widehat\wedge^k_G\circ\tilde s$ for a bounded linear
$\tilde s\colon\widehat\Lambda^k_\pi E\to\widehat\otimes^k_\pi G$; (c) $w^{\otimes k}|_{\cN^k_E}\colon\cN^k_E\to\cN^k_G$
extends to a bounded linear map $\widehat\otimes^k_\pi E\to\cN^k_G$; (d) the \emph{splitting class}
$\operatorname{sp}_k(w)=(w^{\otimes k}|_{\cN^k_E})_*[\cS^k_E]\in\Ext(\widehat\Lambda^k_\pi E,\cN^k_G)$ vanishes. In
particular $(\mathrm S^k_E)$ holds if and only if $\operatorname{sp}_k(\id_E)=0$, if and only if $[\cS^k_E]=0$
(Lemma~\ref{ob:lem:source}).
\item If $w$ is $\mathrm S^k$-split with $s$, and $v\colon G\to G_1$, $v'\colon G_0\to E$ are bounded linear, then $vwv'$
is $\mathrm S^k$-split with $v^{\otimes k}s\,\Lambda^k_\pi v'$. If $w$ is $\mathrm R^k$-split with $\rho$, then $vwv'$ is
$\mathrm R^k$-split with $\widehat\Lambda^k_\pi v\,\rho\,v'^{\otimes k}$. The same holds without completion.
\end{enumerate}
\end{proposition}

The proof is in Appendix~\ref{app:deferred}.

\subsection{The obstruction class}\label{ob:ssec:class}

Throughout this subsection $E$ and $D$ are normed spaces and $k\ge1$; $\Mult^k(\cL(E,D);\widehat\Lambda^k_\pi D)$ is a
Banach space, and $\Mult^k_0(\cL(E,D);\widehat\Lambda^k_\pi D)$ is its closed subspace of maps with zero diagonal.

\begin{lemma}[evaluation]\label{ob:lem:evaluation}
There is a unique bounded linear map
$\ev=\ev^k_{E,D}\colon\widehat\otimes^k_\pi E\to\Mult^k(\cL(E,D);\widehat\Lambda^k_\pi D)$ with
\begin{equation}\label{ob:eq:evaluation}
 \ev(x_1\otimes\dots\otimes x_k)(u_1,\dots,u_k)=u_1x_1\wedge\dots\wedge u_kx_k,
\end{equation}
and $\norm\ev\le1$. It satisfies $\ev(z)(u,\dots,u)=\widehat\Lambda^k_\pi u(\widehat\wedge^k_Ez)$ for
$z\in\widehat\otimes^k_\pi E$ and $u\in\cL(E,D)$; in particular $\ev$ maps $\cN^k_E$ into
$\Mult^k_0(\cL(E,D);\widehat\Lambda^k_\pi D)$.
\end{lemma}

\begin{proof}
For $x\in E^k$ the right side of \eqref{ob:eq:evaluation} is $k$-linear in $(u_1,\dots,u_k)$, of norm at most
$\prod_a\norm{x_a}$, and $k$-linear in $x$. Its linearization in $x$ is a contraction from $\otimes^k_\pi E$ into a
Banach space and extends to $\widehat\otimes^k_\pi E$; pure tensors span a dense subspace, which gives uniqueness. For
fixed $u$, both sides of the identity are bounded linear in $z$ and agree on pure tensors.
\end{proof}

Let $q$ be the quotient map of $\cM^k(\cL(E,D);\widehat\Lambda^k_\pi D)$. Since $q\circ\ev$ vanishes on $\cN^k_E$,
Lemma~\ref{ob:lem:metric-surjection} gives a unique bounded linear
$\overline\ev\colon\hLambda^k_\pi E\to\Poly^k(\cL(E,D);\widehat\Lambda^k_\pi D)$ with
$q\circ\ev=\overline\ev\circ\widehat\wedge^k_E$. By Lemma~\ref{ob:lem:evaluation}, $\overline\ev(y)$ is the class of every
map in $\Mult^k(\cL(E,D);\widehat\Lambda^k_\pi D)$ with diagonal $u\mapsto\widehat\Lambda^k_\pi u(y)$.

The \emph{obstruction class} of the pair $(E,D)$ in degree $k$ is
\begin{equation}\label{ob:eq:class}
 \begin{split}
 \ob_k(E,D)&=(\ev|_{\cN^k_E})_*[\cS^k_E]\\
 &\in\Ext\bigl(\widehat\Lambda^k_\pi E,\Mult^k_0(\cL(E,D);\widehat\Lambda^k_\pi D)\bigr).
 \end{split}
\end{equation}

\begin{theorem}[obstruction theorem]\label{ob:thm:obstruction}
Let $E$ and $D$ be normed spaces and $k\ge1$. The following are equivalent.
\begin{enumerate}
\item The pair $(E,D)$ is good for Banach values in degree $k$.
\item The restriction $\ev|_{\cN^k_E}$ extends to a bounded linear map
$\widehat\otimes^k_\pi E\to\Mult^k_0(\cL(E,D);\widehat\Lambda^k_\pi D)$.
\item $\ob_k(E,D)=0$ in~\eqref{ob:eq:class}.
\end{enumerate}
More precisely, a lift $Q\colon\cL(E,D)^k\to\cL(\Lambda^k_\pi E,\widehat\Lambda^k_\pi D)$ as in
Lemma~\ref{fib:lem:test} gives an extension in~(2) of norm at most $1+\norm Q$, and an extension $\xi$ in~(2) gives such a
lift of norm at most $1+\norm\xi$.
\end{theorem}

\begin{proof}
By Lemma~\ref{fib:lem:test}, (1) holds if and only if a lift $Q$ exists.

\emph{Step 1 (lifts give extensions).} Bounded $k$-linear maps
$Q\colon\cL(E,D)^k\to\cL(\Lambda^k_\pi E,\widehat\Lambda^k_\pi D)$ correspond to bounded linear maps
$Q^\sharp\colon\widehat\Lambda^k_\pi E\to\Mult^k(\cL(E,D);\widehat\Lambda^k_\pi D)$, by
$Q^\sharp(y)(u_1,\dots,u_k)=Q(u_1,\dots,u_k)y$ for $y\in\Lambda^k_\pi E$, with $\norm{Q^\sharp}=\norm Q$. Given a lift $Q$,
put $\xi=\ev-Q^\sharp\widehat\wedge^k_E$; it equals $\ev$ on $\cN^k_E=\ker\widehat\wedge^k_E$. By
Lemma~\ref{ob:lem:evaluation}, $\xi(z)(u,\dots,u)=\Lambda^k_\pi u(\wedge^k_Ez)-Q(u,\dots,u)(\wedge^k_Ez)=0$ for
$z\in\otimes^k_\pi E$, hence for all $z$ by continuity. So $\xi$ is an extension as in~(2). Since $\norm\ev\le1$ by
Lemma~\ref{ob:lem:evaluation} and $\widehat\wedge^k_E$ is contractive, $\norm\xi\le\norm\ev+\norm{Q^\sharp}\le1+\norm Q$.

\emph{Step 2 (extensions give lifts).} Given an extension $\xi$ as in~(2), the map $\ev-\xi$ vanishes on $\cN^k_E$. Since $\widehat\wedge^k_E$ is a
metric surjection with kernel $\cN^k_E$ (Lemma~\ref{ob:lem:metric-surjection}), $\ev-\xi=Q^\sharp\widehat\wedge^k_E$ for a
bounded linear $Q^\sharp$ with $\norm{Q^\sharp}=\norm{\ev-\xi}\le1+\norm\xi$, using $\norm\ev\le1$. The corresponding $Q$
satisfies $Q(u,\dots,u)(\wedge^k_Ez)=\ev(z)(u,\dots,u)-\xi(z)(u,\dots,u)=\Lambda^k_\pi u(\wedge^k_Ez)$ for
$z\in\otimes^k_\pi E$, by Lemma~\ref{ob:lem:evaluation} and since $\xi(z)$ has zero diagonal; so
$Q(u,\dots,u)=\Lambda^k_\pi u$, since $\wedge^k_E$ is surjective.

\emph{Step 3 (the class).} Conditions (2) and (3) are equivalent by Lemma~\ref{ob:lem:split-criteria}(3).
\end{proof}

If $D$ is dense in a normed $K$-space $D'$, the canonical maps $\otimes^k_\pi D\to\otimes^k_\pi D'$ and
$\Lambda^k_\pi D\to\Lambda^k_\pi D'$ are isometric with dense range; so, with $\widehat D$ regarded as a normed $K$-space,
$\widehat\otimes^k_\pi\widehat D=\widehat\otimes^k_\pi D$ and $\widehat\Lambda^k_\pi\widehat D=\widehat\Lambda^k_\pi D$.
The maps are contractive. By continuity $\omega_D$ extends to a contractive map $\omega'\colon D'^k\to\widehat\Lambda^k_\pi D$,
alternating because tuples with two equal entries in $D'^k$ are limits of such tuples in $D^k$; its linearization
(Proposition~\ref{prop:exterior}) is a contraction $\Lambda^k_\pi D'\to\widehat\Lambda^k_\pi D$ restricting to the
inclusion on $\Lambda^k_\pi D$, so the maps are isometric. Density follows from
$x_1\wedge\dots\wedge x_k-y_1\wedge\dots\wedge y_k=\sum_ay_1\wedge\dots\wedge y_{a-1}\wedge(x_a-y_a)\wedge x_{a+1}\wedge\dots\wedge x_k$.
The tensor case is the same.

\begin{theorem}[ways for the class to vanish]\label{ob:thm:mechanisms}
Let $E$ and $D$ be normed spaces and $k\ge1$.
\begin{enumerate}
\item \textup{(Source.)} If $(\mathrm S^k_E)$ holds, then $\ob_k(E,D)=0$. More generally, $\ob_k(E,D)=0$ if
$\ev|_{\cN^k_E}=\tau\circ w^{\otimes k}|_{\cN^k_E}$ for an $\mathrm S^k$-split $w\colon E\to G$ and a bounded linear
$\tau\colon\cN^k_G\to\Mult^k_0(\cL(E,D);\widehat\Lambda^k_\pi D)$. This applies when $w\colon E\to G$ is
$\mathrm S^k$-split and $\varepsilon\colon\cL(E,D)\to\cL(G,\widehat D)$ is bounded linear with $\varepsilon(u)\circ w=u$
for every $u$, with
$\tau(\zeta)(u_1,\dots,u_k)=\widehat\wedge^k_{\widehat D}\bigl(\varepsilon(u_1)\otimes\dots\otimes\varepsilon(u_k)\bigr)\zeta$.
\item \textup{(Parameter.)} If the sequence $\cM^k(\cL(E,D);\widehat\Lambda^k_\pi D)$ splits, and in particular if
$\cL(E,D)$ is polarizable in degree $k$, then $\ob_k(E,D)=0$. More generally, $\ob_k(E,D)=0$ if there are a normed space $P$ with $(\Gamma^k_P)$, a
bounded linear map $\psi\colon\cL(E,D)\to P$ and a bounded $k$-linear map
$\Phi\colon P^k\to\cL(\widehat\otimes^k_\pi E,\widehat\Lambda^k_\pi D)$ such that $\Phi(h,\dots,h)$ vanishes on $\cN^k_E$
for $h\in P$ and $\Phi(\psi u,\dots,\psi u)=\widehat\Lambda^k_\pi u\circ\widehat\wedge^k_E$ for $u\in\cL(E,D)$.
\item \textup{(Target.)} Let
$\alt_*\colon\Mult^k_0(\cL(E,D);\widehat\Lambda^k_\pi D)\to\Mult^k_0(\cL(E,D);\widehat\otimes^k_\pi D)$ be
postcomposition with $\alt^k_D$. Then
$(\alt_*\circ\ev|_{\cN^k_E})_*[\cS^k_E]=0$. If $(\mathrm R^k_D)$ holds, then $\ob_k(E,D)=0$.
\item[(3$'$)] \textup{(Target, relative.)} If $w'\colon G\to\widehat D$ is $\mathrm R^k$-split and
$\varepsilon\colon\cL(E,D)\to\cL(E,G)$ is bounded linear with $w'\circ\varepsilon(u)=u$ for every $u$, then
$\ob_k(E,D)=0$.
\end{enumerate}
\end{theorem}
The proof is in Appendix~\ref{app:deferred}.

\begin{corollary}\label{ob:cor:mechanisms}
Each of $(\mathrm S^k_E)$, $(\mathrm R^k_D)$ and $(\Gamma^k_{\cL(E,D)})$ makes $(E,D)$ good for Banach values in degree $k$;
so does each of the data $(G,w,\varepsilon)$ of Theorem~\ref{ob:thm:mechanisms}(1) and $(G,w',\varepsilon)$ of
Theorem~\ref{ob:thm:mechanisms}(3$'$).
\end{corollary}

\begin{proof}
Theorems~\ref{ob:thm:mechanisms} and~\ref{ob:thm:obstruction}.
\end{proof}

\subsection{The ordering lift}\label{ob:ssec:ordering}

Throughout this subsection $E$ and $D$ are normed spaces, $k\ge1$, and $(e_t)_{t\in\Sigma}$ is a linearly independent
family in $E$, indexed by a totally ordered set $\Sigma$, whose span $V$ is dense in $E$; a Hamel basis of $E$ in any total
order is one. Write $x=\sum_tx(t)e_t$ for $x\in V$, $\Mult^k=\Mult^k(\cL(E,D);\hLambda^k_\pi D)$ and
$\Mult^k_0=\Mult^k_0(\cL(E,D);\hLambda^k_\pi D)$, and let $q$ and $\overline\ev$ be the maps defined after
Lemma~\ref{ob:lem:evaluation}. By the description of $\overline\ev$ given there, in the notation of the proof of
Theorem~\ref{ob:thm:obstruction}, a bounded linear $T\colon\hLambda^k_\pi E\to\Mult^k$ is $Q^\sharp$ for a lift $Q$ as in
Lemma~\ref{fib:lem:test} if and only if $q\circ T=\overline\ev$. Let $c_k(E,D)\in[0,\infty]$ be the infimum of $\norm Q$
over these lifts; by Lemma~\ref{fib:lem:test} it is finite if and only if $(E,D)$ is good for Banach values in degree $k$.

Let $\Sigma^k_<\subseteq\Sigma^k$ be the set of strictly increasing $t$, and put $e^\wedge_t=e_{t_1}\wedge\dots\wedge e_{t_k}$
and $e^\otimes_t=e_{t_1}\otimes\dots\otimes e_{t_k}$ for $t\in\Sigma^k$; the $e^\wedge_t$ with $t\in\Sigma^k_<$ form a basis
of $\Lambda^k_KV$. Let $\wedge_V\colon\otimes^k_KV\to\Lambda^k_KV$ be the quotient map and
$j_\Lambda\colon\Lambda^k_KV\to\hLambda^k_\pi E$, $j_\otimes\colon\otimes^k_KV\to\widehat\otimes^k_\pi E$ the canonical maps,
so that $\widehat\wedge^k_E\circ j_\otimes=j_\Lambda\circ\wedge_V$, and put $\norm y_\pi=\norm{j_\Lambda y}$. By the paragraph
before Theorem~\ref{ob:thm:mechanisms}, $\norm y_\pi$ is the projective seminorm of $y$ computed from representations with
entries in $V$, and $j_\Lambda(\Lambda^k_KV)$ is dense in $\hLambda^k_\pi E$; the same holds for $j_\otimes$. Norms in
$[0,\infty]$ of linear maps on $\Lambda^k_KV$, for $\norm\cdot_\pi$, and of $k$-linear maps on $V^k$ are best constants,
as are the norms of lifts $Q$ and of the maps $Q^\sharp$.

The \emph{ordering section} is the linear map $s_<\colon\Lambda^k_KV\to\otimes^k_KV$ with $s_<(e^\wedge_t)=e^\otimes_t$ for
$t\in\Sigma^k_<$. Then $\wedge_V\circ s_<=\id$, and $s_<=\cP^{(k)}\circ\alt^k$ for the map $\cP^{(k)}$ of
Proposition~\ref{spl:prop:ordering} and the antisymmetrization $\alt^k\colon\Lambda^k_KV\to\otimes^k_KV$. The
\emph{ordering lift} assigns to $u_1,\dots,u_k\in\cL(E,D)$ the linear map
$Q_<(u_1,\dots,u_k)\colon\Lambda^k_KV\to\Lambda^k_\pi D$ with
\begin{equation}\label{ob:eq:ordering-lift}
 Q_<(u_1,\dots,u_k)(e^\wedge_t)=u_1e_{t_1}\wedge\dots\wedge u_ke_{t_k}\qquad(t\in\Sigma^k_<);
\end{equation}
on $x_1\wedge\dots\wedge x_k$ with $x_1,\dots,x_k\in V$ it is the right side of~\eqref{ord:eq:canonical-lift}. Let
$Q^\sharp_<\colon\Lambda^k_KV\to\Mult^k$ be the linear map with $Q^\sharp_<(y)(u_1,\dots,u_k)=Q_<(u_1,\dots,u_k)y$; it takes
values in $\Mult^k$ by~\eqref{ob:eq:ordering-lift}. A \emph{zero-diagonal correction} is a linear map
$L\colon\Lambda^k_KV\to\Mult^k_0$. The proofs of Theorems~\ref{ob:thm:ordering} and~\ref{ob:thm:ultrametric},
Lemma~\ref{ob:lem:repeated} and Corollaries~\ref{ob:cor:adjacent} and~\ref{ob:cor:z-linear} are in
Appendix~\ref{app:deferred}.

\begin{theorem}[ordering representation]\label{ob:thm:ordering}
\begin{enumerate}
\item $\ev\circ j_\otimes\circ s_<=Q^\sharp_<$ and $q\circ Q^\sharp_<=\overline\ev\circ j_\Lambda$.
\item For every lift $Q\colon\cL(E,D)^k\to\cL(\Lambda^k_\pi E,\hLambda^k_\pi D)$ as in Lemma~\ref{fib:lem:test},
$L=Q^\sharp\circ j_\Lambda-Q^\sharp_<$ is a zero-diagonal correction; equivalently,
$Q(u_1,\dots,u_k)\circ j_\Lambda=Q_<(u_1,\dots,u_k)+\Phi(u_1,\dots,u_k)$, where $\Phi(u_1,\dots,u_k)y=L(y)(u_1,\dots,u_k)$ is
$k$-linear in $(u_1,\dots,u_k)$ and $\Phi(u,\dots,u)=0$ for $u\in\cL(E,D)$.
\item For every zero-diagonal correction $L$ with $\norm{Q^\sharp_<+L}<\infty$ there is a unique lift $Q$ with
$Q^\sharp\circ j_\Lambda=Q^\sharp_<+L$, and $\norm Q=\norm{Q^\sharp_<+L}$. Hence $c_k(E,D)=\inf_L\norm{Q^\sharp_<+L}$ over the
zero-diagonal corrections $L$, and $(E,D)$ is good for Banach values in degree $k$ if and only if $\norm{Q^\sharp_<+L}<\infty$
for some zero-diagonal correction $L$.
\item \textup{(Defect form.)} The formula
$\Omega_<(x)=\ev(x_1\otimes\dots\otimes x_k)-Q^\sharp_<(x_1\wedge\dots\wedge x_k)$, for $x\in V^k$, defines a $k$-linear map
$\Omega_<\colon V^k\to\Mult^k_0$. For every zero-diagonal correction $L$, with $(L\circ\wedge)(x)=L(x_1\wedge\dots\wedge x_k)$,
\[
 \norm{Q^\sharp_<+L}\le1+\norm{\Omega_<-L\circ\wedge},\qquad\norm{\Omega_<-L\circ\wedge}\le1+\norm{Q^\sharp_<+L}.
\]
So $(E,D)$ is good for Banach values in degree $k$ if and only if $\Omega_<$ is the sum of a bounded $k$-linear map
$V^k\to\Mult^k_0$ and an alternating $k$-linear map $V^k\to\Mult^k_0$. The infima $c_k(E,D)$ and
$\inf_L\norm{\Omega_<-L\circ\wedge}$ are finite together, and then they differ by at most~$1$.
\end{enumerate}
\end{theorem}

If $\norm{\cP_<(x\otimes y)}\le C\norm x\norm y$ for $x,y\in V$, as in Proposition~\ref{spl:prop:ordering}, no correction is
needed: the lift $Q$ of
Proposition~\ref{fib:prop:split}(1) built from the section $s=\cP^{(k)}\circ\alt^k_E$ of
Proposition~\ref{spl:prop:ordering}(2) satisfies $Q(u_1,\dots,u_k)\circ j_\Lambda=Q_<(u_1,\dots,u_k)$, since
$s\circ j_\Lambda=j_\otimes\circ s_<$.

\begin{lemma}\label{ob:lem:repeated}
The space $\cN^k_E$ is the closure in $\widehat\otimes^k_\pi E$ of the span of the pure tensors $x_1\otimes\dots\otimes x_k$
with $x_a=x_{a+1}$ for some $a<k$, and also of the span of the pure tensors with two equal entries. The pure tensors of the
first kind with all entries in $V$ span a dense subspace of $\cN^k_E$.
\end{lemma}

\begin{corollary}[adjacent diagonals]\label{ob:cor:adjacent}
The pair $(E,D)$ is good for Banach values in degree $k$ if and only if there is a bounded $k$-linear
$\Psi\colon E^k\to\Mult^k_0$ with
\[
 \Psi(x_1,\dots,x_k)(u_1,\dots,u_k)=u_1x_1\wedge\dots\wedge u_kx_k
\]
whenever two of the $x_a$ coincide, equivalently whenever $x_a=x_{a+1}$ for some $a<k$. A lift $Q$ gives such a $\Psi$ with
$\norm\Psi\le1+\norm Q$, and such a $\Psi$ gives a lift $Q$ with $\norm Q\le1+\norm\Psi$.
\end{corollary}

\begin{corollary}[z-linear form]\label{ob:cor:z-linear}
Let $\delta>0$, and let $\nu\colon\hLambda^k_\pi E\to\widehat\otimes^k_\pi E$ be a homogeneous map, $\nu(cy)=c\,\nu(y)$ for $c\in K$,
with $\widehat\wedge^k_E\circ\nu=\id$ and $\norm{\nu(y)}\le(1+\delta)\norm y$; it exists by Lemma~\ref{ob:lem:metric-surjection},
choosing $\nu(y)$ for one point $y$ of each orbit $K^\times y$, $y\ne0$, and extending homogeneously. Put
$\Omega^\nu_<=\ev\circ\nu\circ j_\Lambda-Q^\sharp_<\colon\Lambda^k_KV\to\Mult^k$. Then $\Omega^\nu_<$ takes values in $\Mult^k_0$, is
homogeneous, and is \emph{z-linear}:
\[
 \Bigl\lVert\Omega^\nu_<\Bigl(\sum_iy_i\Bigr)-\sum_i\Omega^\nu_<(y_i)\Bigr\rVert\le2(1+\delta)\sum_i\norm{y_i}_\pi
\]
for every finite family $(y_i)$ in $\Lambda^k_KV$. Let $\operatorname{dist}(\Omega^\nu_<)$ be the infimum, over the linear maps
$L\colon\Lambda^k_KV\to\Mult^k_0$, of the infimum of the $C$ with $\norm{\Omega^\nu_<(y)-L(y)}\le C\norm y_\pi$ for all $y$. Then
$\operatorname{dist}(\Omega^\nu_<)$ and $c_k(E,D)$ are finite together, and then they differ by at most $1+\delta$. In
particular $\ob_k(E,D)=0$ if and only if $\Omega^\nu_<$ lies at finite distance from a linear map $\Lambda^k_KV\to\Mult^k_0$.
\end{corollary}

\begin{theorem}[nonarchimedean values]\label{ob:thm:ultrametric}
Let $K$ be nonarchimedean, and suppose that $C_e=\sup_t\norm{e_t}<\infty$ and that $\abs{x(t)}\le C_c\norm x$ for $x\in V$,
$t\in\Sigma$ and a constant $C_c$. Let $Y$ be a Banach space with a nonarchimedean norm $\lvert\!\lvert\!\lvert\cdot\rvert\!\rvert\!\rvert_Y$ and constants $0<c_-\le c_+$
with $c_-\lvert\!\lvert\!\lvert\cdot\rvert\!\rvert\!\rvert_Y\le\norm\cdot\le c_+\lvert\!\lvert\!\lvert\cdot\rvert\!\rvert\!\rvert_Y$, and let $\alpha\colon\cN^k_E\to Y$ be bounded linear. Then there is a bounded linear
$\tilde\alpha\colon\widehat\otimes^k_\pi E\to Y$ of norm at most $2(c_+/c_-)\norm\alpha(C_cC_e)^k$ with
$\tilde\alpha(j_\otimes z)=\alpha\bigl(j_\otimes(z-s_<(\wedge_Vz))\bigr)$ for $z\in\otimes^k_KV$ and $\tilde\alpha=\alpha$ on $\cN^k_E$. Hence
$\alpha_*[\cS^k_E]=0$. In particular $\ob_k(E,D)=0$ whenever $\hLambda^k_\pi D$ has an equivalent nonarchimedean norm.
\end{theorem}
%% ---- end sec-03-obstruction.tex
%% ---- begin sec-04-canonical.tex
\section{Canonical splittings and forced asymmetry}\label{sec:canonical}

Throughout this section $K$ is a nontrivially normed field and $[n]:=\{1,\dots,n\}$. For $\sigma\in S_n$ let
$\tau_\sigma$ be the isometry of the $n$\nobreakdash-th projective tensor power of a normed space, and of its completion,
with $\tau_\sigma(x_1\otimes\dots\otimes x_n)=x_{\sigma^{-1}(1)}\otimes\dots\otimes x_{\sigma^{-1}(n)}$; it commutes with
$u^{\otimes n}$ for every bounded operator $u$, and $\tau_\sigma\tau_{\sigma'}=\tau_{\sigma\sigma'}$. Evaluation on pure
tensors gives
\begin{equation}\label{can:eq:signs}
 \operatorname{alt}^k_E\circ\wedge^k_E=\sum_{\sigma\in S_k}\sgn(\sigma)\,\tau_\sigma,\qquad
 \tau_\sigma\circ\operatorname{alt}^k_E=\sgn(\sigma)\operatorname{alt}^k_E,\qquad
 \wedge^k_E\circ\tau_\sigma=\sgn(\sigma)\,\wedge^k_E ,
\end{equation}
and the first identity extends to $\widehat\otimes^k_\pi E$ by continuity. For vectors $x_1,\dots,x_n$ and
$S\subseteq[n]$ put $v_S:=\sum_{i\in S}x_i$. Expanding $v_S^{\otimes n}$ over the maps $f\colon[n]\to S$, and using that
$\sum_{S\supseteq f([n])}(-1)^{n-\abs S}$ is $1$ if $f$ is onto and $0$ otherwise, gives in every characteristic
\begin{equation}\label{can:eq:polarization}
 \sum_{S\subseteq[n]}(-1)^{n-\abs S}\,v_S^{\otimes n}=\sum_{\sigma\in S_n}x_{\sigma(1)}\otimes\dots\otimes x_{\sigma(n)} .
\end{equation}

\begin{proposition}[natural splittings]\label{can:prop:natural}
Let $E$, $D$ and $P$ be normed spaces.
\begin{enumerate}
\item If $k!\ne0$ in $K$, then $\operatorname{alt}^k_E/k!$ is a section of $\wedge^k_E$ of norm at most $k!/\abs{k!}$ and
$\wedge^k_D/k!$ is a retraction of $\operatorname{alt}^k_D$ of norm at most $1/\abs{k!}$; both are natural with respect
to every $u\in\cL(E,D)$ and satisfy the sign rules of \eqref{can:eq:signs}.
\item If $n!\ne0$ in $K$, then $R_n:=\frac1{n!}\sum_{\sigma\in S_n}\tau_\sigma$ is a projection of
$\widehat\otimes^n_\pi P$ onto $\widehat\Gamma^n_\pi P$ of norm at most $n!/\abs{n!}$; it commutes with $u^{\otimes n}$ for
bounded operators $u$ and satisfies $\tau_\sigma R_n=R_n\tau_\sigma=R_n$.
\end{enumerate}
\end{proposition}

\begin{proof}
(1) holds by the properties of $\wedge^k_E$ and $\operatorname{alt}^k_E$ in Section~\ref{sec:duality}.
(2) The norm bound and the commutation rules hold since each $\tau_\sigma$ is an isometry commuting with $u^{\otimes n}$
and $\tau_\sigma\tau_{\sigma'}=\tau_{\sigma\sigma'}$. $R_n$ fixes every pure power, hence $\widehat\Gamma^n_\pi P$
pointwise, and by~\eqref{can:eq:polarization} and continuity it maps $\widehat\otimes^n_\pi P$ into
$\widehat\Gamma^n_\pi P$.
\end{proof}

So every pair is good in degree $k$ when $k!\ne0$ (Propositions~\ref{fib:prop:split}(3) and~\ref{prop:split-lift}(1)).
A matrix in $GL_k(K)$ is \emph{monomial} if each row and each column has exactly one nonzero entry.

\begin{theorem}[no equivariant splitting]\label{can:thm:monomial}
Let $k\ge2$ with $k!=0$ in $K$, and let $E_0$ be $K^k$ with a norm for which the coordinate functionals
$e^*_1,\dots,e^*_k$ of the standard basis $e_1,\dots,e_k$ are bounded; by Lemma~\ref{lem:findim} every norm
qualifies if $K$ is complete. For a monomial $g$, extend $g^{\otimes k}$ and $\Lambda^k_\pi g$ to the completions.
\begin{enumerate}
\item No linear $s\colon\Lambda^k_\pi E_0\to\widehat\otimes^k_\pi E_0$ with $\widehat\wedge^k_{E_0}\circ s=\id$ satisfies
$s\circ\Lambda^k_\pi g=g^{\otimes k}\circ s$ for all monomial $g$.
\item No linear $\rho\colon\otimes^k_\pi E_0\to\widehat\Lambda^k_\pi E_0$ with $\rho\circ\operatorname{alt}^k_{E_0}=\id$
satisfies $\rho\circ g^{\otimes k}=\Lambda^k_\pi g\circ\rho$ for all monomial $g$.
\item No linear projection $R$ of $\widehat\otimes^k_\pi E_0$ onto $\widehat\Gamma^k_\pi E_0$ satisfies
$R\circ g^{\otimes k}=g^{\otimes k}\circ R$ for all monomial $g$.
\end{enumerate}
Statements (1)--(3) also hold for the algebraic powers $\otimes^k_KK^k$ and $\Lambda^k_KK^k$.
\end{theorem}

The proof is in Appendix~\ref{app:deferred}.

Let a group $\mathsf H$ act on a set $I$, and let $V=K^{(I)}$ be the space of finitely supported functions $I\to K$, with
unit vectors $e_i$. Each $h\in\mathsf H$ acts on $V$ by $he_i=e_{h(i)}$, on $\otimes^k_KV$ by $h^{\otimes k}$ and on
$\Lambda^k_KV$ by $\Lambda^k_Kh$; a linear map between two of these spaces is \emph{$\mathsf H$\nobreakdash-equivariant}
if it intertwines the actions of every $h\in\mathsf H$. The maps $\wedge^k\colon\otimes^k_KV\to\Lambda^k_KV$,
$\operatorname{alt}^k\colon\Lambda^k_KV\to\otimes^k_KV$ and $\tau_\sigma$ are given by the formulas of
Section~\ref{ob:ssec:powers} and of this section. A \emph{$k$\nobreakdash-subset} is a subset with $k$ elements; for a
finite $T\subseteq I$, $\mathsf H_T$ is its setwise stabilizer and $\bar{\mathsf H}_T\subseteq\operatorname{Sym}(T)$ is
the group of permutations of $T$ induced by $\mathsf H_T$. A \emph{$k$\nobreakdash-multiset} is a finitely supported
$\mathbf m\colon I\to\N$ with $\sum_i\mathbf m(i)=k$; $\mathsf H$ acts on these by
$(h\mathbf m)(i)=\mathbf m(h^{-1}(i))$, $\mathsf H_{\mathbf m}$ is the stabilizer of $\mathbf m$, and
$\bar{\mathsf H}_{\mathbf m}\subseteq\operatorname{Sym}(\operatorname{supp}\mathbf m)$ is the group induced by
$\mathsf H_{\mathbf m}$.
Let $\Gamma^k_KV\subseteq\otimes^k_KV$ be the span of the pure powers $x^{\otimes k}$, $x\in V$; since $K$ is infinite,
it is the space of tensors fixed by every $\tau_\sigma$, by the paragraph at the beginning of
Section~\ref{dgr:ssec:polarizable}.

\begin{theorem}[permutation groups]\label{can:thm:permutation}
Let $\mathsf H$ act on $I$, and let $k\ge1$.
\begin{enumerate}
\item There is an $\mathsf H$\nobreakdash-equivariant linear $s\colon\Lambda^k_KV\to\otimes^k_KV$ with
$\wedge^k\circ s=\id$ if and only if $\abs{\bar{\mathsf H}_T}\ne0$ in $K$ for every $k$\nobreakdash-subset $T$ of $I$.
\item There is an $\mathsf H$\nobreakdash-equivariant linear $\rho\colon\otimes^k_KV\to\Lambda^k_KV$ with
$\rho\circ\operatorname{alt}^k=\id$ if and only if the condition in~(1) holds.
\item There is an $\mathsf H$\nobreakdash-equivariant linear projection of $\otimes^k_KV$ onto $\Gamma^k_KV$ if and only
if $\abs{\bar{\mathsf H}_{\mathbf m}}\ne0$ in $K$ for every $k$\nobreakdash-multiset $\mathbf m$ on $I$.
\end{enumerate}
When the condition in (1), (2) or (3) holds, the corresponding map can be chosen to commute also with the operator
$e_i\mapsto\lambda_ie_i$ of $V$ for every $\lambda\in(K^\times)^I$.
\begin{enumerate}\setcounter{enumi}{3}
\item Let $E$ be a normed space with vectors $e_i\in E$ and bounded linear functionals $e^*_i$ on $E$, $i\in I$, with
$e^*_i(e_j)=\delta_{ij}$, and let $\mathsf H$ act on $I$ and on $E$ by bounded linear maps with
\begin{equation}\label{can:eq:compatible}
 he_i=e_{h(i)},\qquad e^*_i\circ h=e^*_{h^{-1}(i)}\qquad(h\in\mathsf H,\ i\in I);
\end{equation}
extend $h^{\otimes k}$ and $\Lambda^k_\pi h$ to the completions. If $\abs{\bar{\mathsf H}_T}=0$ in $K$ for some
$k$\nobreakdash-subset $T$, no linear $s\colon\Lambda^k_\pi E\to\widehat\otimes^k_\pi E$ with
$\widehat\wedge^k_E\circ s=\id$ satisfies $s\circ\Lambda^k_\pi h=h^{\otimes k}\circ s$ for all $h\in\mathsf H_T$, and no
linear $\rho\colon\otimes^k_\pi E\to\widehat\Lambda^k_\pi E$ with $\rho\circ\operatorname{alt}^k_E=\id$ satisfies
$\rho\circ h^{\otimes k}=\Lambda^k_\pi h\circ\rho$ for all $h\in\mathsf H_T$. If $\abs{\bar{\mathsf H}_{\mathbf m}}=0$ in
$K$ for some $k$\nobreakdash-multiset $\mathbf m$, no linear projection $R$ of $\widehat\otimes^k_\pi E$ onto
$\widehat\Gamma^k_\pi E$ satisfies $R\circ h^{\otimes k}=h^{\otimes k}\circ R$ for all $h\in\mathsf H_{\mathbf m}$.
\end{enumerate}
\end{theorem}

The proof is in Appendix~\ref{app:obstruction-class}.

\begin{corollary}[one element]\label{can:cor:p-cycle}
\begin{enumerate}
\item In each of (1)--(3) of Theorem~\ref{can:thm:permutation}, the following are equivalent: the
$\mathsf H$\nobreakdash-equivariant map exists; for every $h\in\mathsf H$, the map equivariant for the cyclic group
generated by $h$ exists; no $h\in\mathsf H$ maps a $k$\nobreakdash-subset $T$ onto itself, in (1) and (2), or fixes a
$k$\nobreakdash-multiset $\mathbf m$, in (3), and induces on $T$, respectively on $\operatorname{supp}\mathbf m$, a
permutation of order $\operatorname{char}K$. If $\operatorname{char}K=0$, the three statements hold.
\item Let $k!=0$ in $K$, so that $p=\operatorname{char}K$ satisfies $2\le p\le k$. Let $E_0$ be as in
Theorem~\ref{can:thm:monomial}, let $\nu\in S_k$ be a cycle of length $p$, and let $g$ be the permutation matrix with
$ge_i=e_{\nu(i)}$. No map $s$, $\rho$ or $R$ as in (1), (2) or (3) of Theorem~\ref{can:thm:monomial} satisfies the
corresponding equivariance identity for this single $g$.
\end{enumerate}
\end{corollary}

The proof is in Appendix~\ref{app:deferred}.

For a finite group $\mathsf G$ acting trivially on $K$, the zeroth Tate cohomology group is
$\widehat H^0(\mathsf G;K)=K/\abs{\mathsf G}K$~\cite[Ch.~VI]{Br}; it is $K$ if $\abs{\mathsf G}=0$ in $K$ and $0$
otherwise.

\begin{proposition}[Tate classes]\label{can:prop:tate}
Let $\mathsf H$ act on $I$, let $k\ge1$, and let $T\subseteq I$ be a $k$\nobreakdash-subset. For a bijection
$f\colon[k]\to T$ put $e_f:=e_{f(1)}\otimes\dots\otimes e_{f(k)}$. Fix a bijection $f_T\colon[k]\to T$, put
$\eta_T:=e_{f_T(1)}\wedge\dots\wedge e_{f_T(k)}$, let $U_T\subseteq\otimes^k_KV$ be the span of the tensors $e_f$ over the
bijections $f\colon[k]\to T$, put $\cN_T:=U_T\cap\ker\wedge^k$, and let $\chi_T\colon\mathsf H_T\to\{1,-1\}$ be the
homomorphism with $\Lambda^k_Kh(\eta_T)=\chi_T(h)\eta_T$. Let $\mathsf H_T$ act on $U_T$ by
$h\cdot z:=\chi_T(h)h^{\otimes k}z$. Write $H^1(\mathsf H_T;\cN_T)$ for the space of maps $c\colon\mathsf H_T\to\cN_T$
with $c(hh')=c(h)+h\cdot c(h')$ for $h,h'\in\mathsf H_T$, modulo the maps $h\mapsto h\cdot n-n$ with $n\in\cN_T$, and for
$a\in K$ let $\delta_T(a)\in H^1(\mathsf H_T;\cN_T)$ be the class of $h\mapsto a(h\cdot e_{f_T}-e_{f_T})$.
\begin{enumerate}
\item $\delta_T(1)=0$ if and only if some $z\in U_T$ satisfies $\wedge^kz=\eta_T$ and $h^{\otimes k}z=\chi_T(h)z$ for
every $h\in\mathsf H_T$.
\item $\ker\delta_T=\abs{\bar{\mathsf H}_T}K$. Hence $\delta_T$ induces an injective linear map from
$\widehat H^0(\bar{\mathsf H}_T;K)=K/\abs{\bar{\mathsf H}_T}K$ into $H^1(\mathsf H_T;\cN_T)$ that sends the class of $1$
to $\delta_T(1)$.
\item An $\mathsf H$\nobreakdash-equivariant linear section of $\wedge^k\colon\otimes^k_KV\to\Lambda^k_KV$ exists if and
only if $\delta_T(1)=0$ for every $k$\nobreakdash-subset $T$ of $I$.
\end{enumerate}
\end{proposition}

The proof is in Appendix~\ref{app:obstruction-class}. By~(2) and~(3), an $\mathsf H$\nobreakdash-equivariant section
exists if and only if the Tate classes $[1]\in\widehat H^0(\bar{\mathsf H}_T;K)$ vanish for all $k$\nobreakdash-subsets $T$.

A \emph{lift} is, as in Lemma~\ref{fib:lem:test}, a bounded $k$\nobreakdash-linear map
$Q\colon\cL(E,D)^k\to\cL(\Lambda^k_\pi E,\widehat\Lambda^k_\pi D)$ with $Q(u,\dots,u)=\Lambda^k_\pi u$. For
$u_1,\dots,u_k\in\cL(E,D)$ define $\Lsym(u_1,\dots,u_k)\in\cL(\Lambda^k_\pi E,\Lambda^k_\pi D)$ by
\begin{equation}\label{can:eq:sym}
 \Lsym(u_1,\dots,u_k)(x_1\wedge\dots\wedge x_k):=\sum_{\sigma\in S_k}u_{\sigma(1)}x_1\wedge\dots\wedge u_{\sigma(k)}x_k .
\end{equation}
The right side is alternating in $(x_b)$, since the terms for $\sigma$ and $\sigma\circ(a\ b)$ cancel when $x_a=x_b$; by
Proposition~\ref{prop:exterior}, $\Lsym$ is a bounded symmetric $k$\nobreakdash-linear map with
$\Lsym(u,\dots,u)=k!\,\Lambda^k_\pi u$. If $k!\ne0$, $\Lsym/k!$ is a symmetric lift, the normalized alternatization of
Remark~\ref{rem:FFunit}.

\begin{proposition}[forced asymmetry]\label{can:prop:asymmetry}
Let $E$ and $D$ be normed spaces.
\begin{enumerate}
\item Every lift $Q$ satisfies $\sum_{\sigma\in S_k}Q(u_{\sigma(1)},\dots,u_{\sigma(k)})=\Lsym(u_1,\dots,u_k)$.
\end{enumerate}
Suppose now that $k!=0$ in $K$.
\begin{enumerate}\setcounter{enumi}{1}
\item If $E$ has $k$ linearly independent bounded functionals and $\widehat\Lambda^k_\pi D\ne0$, no lift is symmetric.
\item No linear retraction $\rho\colon\otimes^k_\pi D\to\widehat\Lambda^k_\pi D$ of $\operatorname{alt}^k_D$ satisfies
$\rho\circ\tau_\sigma=\sgn(\sigma)\rho$ for all $\sigma\in S_k$, unless $\Lambda^k_\pi D=0$.
\item If $E$ has $k$ linearly independent bounded functionals, no linear section
$s\colon\Lambda^k_\pi E\to\widehat\otimes^k_\pi E$ of $\widehat\wedge^k_E$ takes values in the sign-isotypic tensors
$\{y:\tau_\sigma y=\sgn(\sigma)y\text{ for all }\sigma\}$.
\end{enumerate}
\end{proposition}

The proof is in Appendix~\ref{app:deferred}.

The functional hypothesis in~(2) is needed: for $k=p$, $K=\F_p(t)$ and the space
$E\subseteq A=\F_p((t))[\epsilon]/(\epsilon^p)$ of Lemma~\ref{dom:rigid}, every bounded functional on $E$ is zero and
$\cL(E,E)=K\cdot\id$, while $\Lambda^p_\pi E\ne0$ by Proposition~\ref{prop:exterior}, since the determinant in the basis
$(\epsilon^i)_{0\le i<p}$ of $A$ is a bounded alternating map on $E^p$ with value $1$ at that basis. So
$Q(c_1\id,\dots,c_p\id)=c_1\cdots c_p\,\id$ is a symmetric lift.
%% ---- end sec-04-canonical.tex
%% ---- begin sec-05-splittings.tex
\section{Splittings that exist}\label{sec:positive}

\subsection{Splittings for one value space}\label{spl:ssec:one-value}

\begin{proposition}[splitting the alternating inclusion]\label{prop:split-lift}
If the inclusion $j_{E,F}\colon\Alt^k(E;F)\hookrightarrow\Mult^k(E;F)$ has a bounded linear retraction $\rho$, then $A^k_{E,D;F}$ has a bounded $k$-linear lift of norm at most $\norm\rho$, for every normed $D$. Such a lift also exists in each of the following cases:
\begin{enumerate}
\item $k!\ne0$ in $K$;
\item either $E$ or $D$ has a finite algebraic basis with continuous coordinate functionals.
\end{enumerate}
No completeness hypothesis is required.
\end{proposition}
\begin{proof}
For a retraction, $Q(u_1,\ldots,u_k)(m)=\rho\bigl(m\circ(u_1,\ldots,u_k)\bigr)$ is a lift: the mixed composition has
norm at most $\norm m\prod_i\norm{u_i}$, and its diagonal is already alternating.

If $k!\ne0$, the lift
\[
 Q(u_1,\ldots,u_k)(m)(x)
 =\frac1{k!}\sum_{\sigma\in S_k}
 m(u_{\sigma(1)}x_1,\ldots,u_{\sigma(k)}x_k)
\]
vanishes when two $x$-arguments agree, by pairing the permutations that exchange their positions; on the diagonal all
$k!$ summands coincide; and its norm is at most $k!/\abs{k!}$, the numerator counting summands.

For continuous coordinates $(e_i,e^*_i)_{1\le i\le d}$ on $E$, sum over increasing words $i_1<\cdots<i_k$:
\[
 Q(u_1,\ldots,u_k)(m)(x)
 =\sum_{i_1<\cdots<i_k}\det\bigl(e^*_{i_a}(x_b)\bigr)_{a,b}
       m(u_1e_{i_1},\ldots,u_ke_{i_k}).
\]
If the coordinates are on $D$, replace the summand by $\det(e^*_{i_a}(u_ax_b))_{a,b}\,m(e_{i_1},\ldots,e_{i_k})$. Equal
columns give alternation, and the exterior-basis expansion gives the diagonal identity. Each is a finite sum of norm at
most $k!\sum_{i_1<\cdots<i_k}\prod_a\norm{e^*_{i_a}}\norm{e_{i_a}}$, empty if $k>d$.
\end{proof}

Over a complete field, finite-dimensional spaces have continuous coordinates (Lemma~\ref{lem:findim}), so cases~(1)
and~(2) contain Bourbaki's hypotheses, characteristic zero and finite rank~\cite[n\textsuperscript{o}~7.8]{Bou}; over
incomplete $K$ finite rank does not suffice, since the $a$-coordinate on $K+Ka\subseteq\widehat K$,
$a\in\widehat K\setminus K$, is discontinuous.

A nonarchimedean normed space is \emph{spherically complete} if every nonempty family of pairwise-intersecting
closed balls has a common point. In a nonarchimedean metric, intersecting balls are nested.

\begin{theorem}[spherically complete targets]\label{thm:spherical}
Suppose $K$ is nonarchimedean and $F$ is nonarchimedean and spherically complete. Then, for every normed $E$, the space $\Alt^k(E;F)$ is nonarchimedean and spherically complete, and its inclusion into $\Mult^k(E;F)$ has a contractive linear retraction. Thus $A^k_{E,D;F}$ is a continuous polynomial for every normed $E,D$ and every $k$. Neither the field nor the source spaces need be complete.
\end{theorem}
The proof is in Appendix~\ref{app:deferred}. An equivalent nonarchimedean spherically complete norm on $F$ suffices,
since equivalent norms preserve analyticity.

\begin{corollary}[discretely valued fields]\label{cor:discrete-targets}
Suppose $K$ is nonarchimedean with $\abs{K^\times}=r^\Z$ for some $0<r<1$.
Every Banach target $F$ admitting an equivalent nonarchimedean norm makes $A^k_{E,D;F}$
analytic in every degree and for all normed $E,D$. The field itself need not be complete.
\end{corollary}
\begin{proof}
Round the equivalent nonarchimedean norm upwards to $r^\Z$, with value zero at zero. The new norm lies between the old
one and $r^{-1}$ times it, and satisfies the norm axioms, since rounding is monotone and commutes with multiplication by
$r^\Z$. It is complete with nonzero distances in $r^\Z$, hence spherically complete by Lemma~\ref{lem:discdist}. Apply
Theorem~\ref{thm:spherical} and transport the resulting lift back through the equivalent norms.
\end{proof}

\paragraph{\textit{Changing the target category.}}
Over nonarchimedean $K$, the maximum in place of the sum in the projective seminorm gives a completed exterior power
$\widehat\Lambda^k_{\pi,\mathrm{na}}E$ representing $\Alt^k(E;F)$ for complete nonarchimedean $F$ (the universal-property
proof of Proposition~\ref{prop:exterior};~\cite[\S4]{Se}). Over a discretely valued field,
$u\mapsto\widehat\Lambda^k_{\pi,\mathrm{na}}u$ is analytic: it is $A^k_{E,D;F}$ for $F=\widehat\Lambda^k_{\pi,\mathrm{na}}D$,
evaluated at the universal alternating map, so Corollary~\ref{cor:discrete-targets} applies. The obstructions of
Section~\ref{sec:sharpness} concern the sum construction, whose target has no equivalent nonarchimedean norm.

In the next theorem and its proof, $K$ is complete and nonarchimedean, and $q_Y$, $Y^u$ and $j_Y\colon Y\to Y^u$ are the
seminorm, the ultrametric envelope and its canonical map of a normed space $Y$, constructed in the proof of
Proposition~\ref{prop:scalar-countable}.

\begin{theorem}[ultrametric values]\label{fib:thm:ultra}
Let $F$ be a Banach space with a nonarchimedean norm, and suppose that $E^u$ or $D^u$ is a bounded linear retract of
$c_0(I,K)$ for some set $I$: there are bounded linear maps $\iota$ into $c_0(I,K)$ and $\iota'$ out of it with
$\iota'\iota=\id$. Then $A^k_{E,D;F}$ has a bounded lift of norm at most $(\norm\iota\norm{\iota'})^k$, and is therefore a continuous
polynomial. If $E$ or $D$ is of countable type (Appendix~\ref{app:coordinate-methods}), there are such lifts of norm at most $t^{-k}$ for every $0<t<1$.
\end{theorem}

The proof is in Appendix~\ref{app:deferred}. If $\operatorname{char}K=p>0$ and $k\ge p$, both hypotheses matter:
Theorem~\ref{thm:operator-obstructions}(2) gives a failure with $F=K$ if $K$ is not spherically complete; for
$K=\kappa((X))$ and $E=D=\linf(\N,K)\cong c_0(I,K)$~\cite[\S1, Proposition~1]{Se}, Theorem~\ref{thm:laurent} gives one
with the target $F_1$, which has no equivalent nonarchimedean norm, so this pair is not good for Banach values
(Corollary~\ref{obs:cor:linf}).

\subsection{Source and target splittings}\label{spl:ssec:source-target}

By Proposition~\ref{fib:prop:split}, a source splitting $(\mathrm S^k_E)$ or a target splitting $(\mathrm R^k_D)$ makes
$(E,D)$ good for Banach values in degree $k$. A Banach space $Y$ is \emph{projective} if, for every bounded linear
surjection $\psi\colon Y_1\to Y_2$ of Banach spaces, every bounded linear $T\colon Y\to Y_2$ equals $\psi\circ T'$ for a bounded
linear $T'\colon Y\to Y_1$. For a totally ordered set $J$, let $J^k_<$ be the set of $i\in J^k$ with $i_1<\dots<i_k$, and
let $(e_i)_{i\in J}$ and $(\delta_i)_{i\in J^k_<}$ be the unit vectors of $\ell^1(J,K)$ and of $\ell^1(J^k_<,\widehat K)$.

\begin{proposition}[projective sources]\label{spl:prop:projective}
Let $E$ be a normed space and $k\ge1$.
\begin{enumerate}
\item For every totally ordered set $J$, the map $\delta_i\mapsto e_{i_1}\wedge\dots\wedge e_{i_k}$ extends to an isometric
isomorphism of $\ell^1(J^k_<,\widehat K)$ onto $\widehat\Lambda^k_\pi\ell^1(J,K)$.
\item If $\widehat\Lambda^k_\pi E$ is projective, then $(\mathrm S^k_E)$ holds.
\item If $E$ is a bounded linear retract of $\ell^1(J,K)$ for some set $J$, then $(\mathrm S^n_E)$ holds for every $n\ge1$.
\end{enumerate}
\end{proposition}

The proof is in Appendix~\ref{app:deferred}.

Let $\Sigma$ be a totally ordered set. An \emph{interval} of $\Sigma$ is a subset $I$ containing every $t$ with
$t_1<t<t_2$ for some $t_1,t_2\in I$. For a linearly independent family $(e_t)_{t\in\Sigma}$ in $E$ with span $V$, write
$x(t)$ for the coefficients of $x\in V$ and $\cP_Ix=\sum_{t\in I}x(t)e_t$.

\begin{theorem}[lower estimates]\label{fib:thm:lower}
Let $K$ be any nontrivially normed field and $E$ a normed space with a linearly independent family $(e_t)_{t\in\Sigma}$
whose span $V$ is dense. Suppose that for some $0<r<2$ and $C\ge1$,
\begin{equation}\label{fib:eq:lower}
 \Bigl(\sum_{i=1}^N\norm{\cP_{I_i}x}^r\Bigr)^{1/r}\le C\norm x
 \qquad(x\in V,\ I_1,\dots,I_N\text{ pairwise disjoint intervals}).
\end{equation}
Put $C_\sharp=2C^2/(1-2^{1-2/r})$.
\begin{enumerate}
\item The linear map $\cP_<$ on $V\otimes V$ with $\cP_<(e_t\otimes e_{t'})=e_t\otimes e_{t'}$ for $t<t'$ and
$\cP_<(e_t\otimes e_{t'})=0$ otherwise satisfies $\norm{\cP_<(x\otimes y)}\le C_\sharp\norm x\norm y$ for $x,y\in V$. It
therefore extends to a bounded linear map $E\otimes_\pi E\to E\widehat\otimes_\pi E$ of norm at most $C_\sharp$.
\item For every $k\ge1$, $(\mathrm S^k_E)$ and $(\mathrm R^k_E)$ hold, with $\norm s\le k!\,C_\sharp^{k-1}$ and
$\norm\rho\le C_\sharp^{k-1}$. If $E=V$, they hold without completion.
\item For every normed $D$ and every $k$, the pairs $(E,D)$ and $(D,E)$ are good for Banach values, and they are good
if $E=V$.
\end{enumerate}
\end{theorem}
The proof is in Appendix~\ref{app:ordering}.

A normed space $E$ has a \emph{lower $r$\nobreakdash-estimate} if it has a family $(e_t)_{t\in\Sigma}$ as in
Theorem~\ref{fib:thm:lower} satisfying~\eqref{fib:eq:lower} for some $C\ge1$. Examples with $C=1$ are $\ell^r(\Sigma,K)$,
$1\le r<2$, with its unit vectors, its subspace of finitely supported vectors, and, for $r_1,r_2\ge1$ with
$r=\max(r_1,r_2)<2$, the space $\ell^{r_1}(\N;\ell^{r_2}(\N,K))$ ordered lexicographically, by Minkowski's inequality if
$r_1\le r_2$ and by $\ell^{r_2}\subseteq\ell^{r_1}$ otherwise. Conclusions~(2), with norms multiplied by
$(\norm\iota\norm{\iota'})^k$, and~(3) pass to every $Y$ with bounded $\iota\colon Y\to E$ and $\iota'\colon E\to Y$,
$\iota'\iota=\id_Y$: the maps $\iota'^{\otimes k}\circ s\circ\Lambda^k_\pi\iota$ and
$\widehat\Lambda^k_\pi\iota'\circ\rho\circ\iota^{\otimes k}$ satisfy $(\mathrm S^k_Y)$ and $(\mathrm R^k_Y)$, by naturality.

\begin{proposition}[ordering splittings]\label{spl:prop:ordering}
Let $E$ be a normed space with a linearly independent family $(e_t)_{t\in\Sigma}$ whose span $V$ is dense, and suppose
that the map $\cP_<$ of Theorem~\ref{fib:thm:lower}(1) satisfies $\norm{\cP_<(x\otimes y)}\le C\norm x\norm y$ for
$x,y\in V$. Let $n\ge1$, and let $\cP^{(n)}$ and $R_n$ be the linear maps on $\otimes^nV$ that send
$e_{t_1}\otimes\dots\otimes e_{t_n}$ to itself if $t_1<\dots<t_n$, respectively to the sum $O_t$ of the tensors
$e_{\tilde t_1}\otimes\dots\otimes e_{\tilde t_n}$ over the distinct rearrangements $\tilde t$ of $t$ if $t_1\le\dots\le t_n$, and to $0$
otherwise.
\begin{enumerate}
\item $\cP^{(n)}$ extends to a bounded linear map $\otimes^n_\pi E\to\widehat\otimes^n_\pi E$ of norm at most $C^{n-1}$.
\item $s=\cP^{(n)}\circ\operatorname{alt}^n_E$ and $\rho=\widehat\wedge^n_E\circ\cP^{(n)}$ satisfy $(\mathrm S^n_E)$ and
$(\mathrm R^n_E)$, with $\norm s\le n!\,C^{n-1}$ and $\norm\rho\le C^{n-1}$. If $E=V$, both hold without completion.
\item $R_n$ extends to a projection of $\widehat\otimes^n_\pi E$ onto $\widehat\Gamma^n_\pi E$, so $(\Gamma^n_E)$ holds.
Its norm is at most the sum of $n!\,(c_1!\cdots c_m!)^{-1}C^{m-1}(1+2C)^{n-m}$ over the compositions
$c=(c_1,\dots,c_m)$ of $n$ into positive parts.
\end{enumerate}
The proof is in Appendix~\ref{app:ordering}.
\end{proposition}

By Theorem~\ref{fib:thm:lower}(1), the hypothesis holds with $C=C_\sharp$ for every family satisfying~\eqref{fib:eq:lower}.

In the rest of this subsection $K$ is nonarchimedean. A linearly independent family $(e_t)_{t\in\Sigma}$ in $E$ whose span
$V$ is dense is \emph{$\ell^2$-dominated} if $(\sum_t\abs{x(t)}^2)^{1/2}\le C_2\norm x$ for $x\in V$ and some constant
$C_2$. For $1\le r<\infty$, an \emph{$r$-norm} is a norm $\nu$ with $\nu(y+z)^r\le\nu(y)^r+\nu(z)^r$; a normed space $D$ has an \emph{equivalent
$r$-norm} if it carries an $r$-norm $\nu$ with $a\nu(y)\le\norm y\le b\nu(y)$ for constants $0<a\le b$.

\begin{proposition}[$r$-norms]\label{spl:prop:r-norm}
Let $1\le r<\infty$. A normed space $D$ has an equivalent $r$-norm if and only if there is a constant $C$ with
\begin{equation}\label{spl:eq:r-norm}
 \Bigl\lVert\sum_t\xi_ty_t\Bigr\rVert\le C\Bigl(\sum_t\abs{\xi_t}^r\Bigr)^{1/r}\max_t\norm{y_t}
\end{equation}
for all finite families $(\xi_t)$ in $K$ and $(y_t)$ in $D$ with a common index set. The proof is in
Appendix~\ref{app:ordering}.
\end{proposition}

\begin{theorem}[relative ordering splittings]\label{spl:thm:relative-ordering}
Let $E$ be a normed space with an $\ell^2$-dominated linearly independent family $(e_t)_{t\in\Sigma}$ whose span $V$ is
dense and with $\sup_t\norm{e_t}<\infty$.
\begin{enumerate}
\item For every $r\in(2,\infty)$, the coordinate map $V\to\ell^r(\Sigma,\widehat K)$, $x\mapsto(x(t))_{t\in\Sigma}$,
extends to a bounded linear map $w\colon E\to\ell^r(\Sigma,\widehat K)$, which is $\mathrm S^n$-split
(Definition~\ref{ob:def:split-operators}) for every $n\ge1$.
\item If $D$ has an equivalent $r$-norm for some $r\in(2,\infty)$, then $(E,D)$ is good for Banach values in every degree,
and good in every degree if $E=V$.
\end{enumerate}
The proof is in Appendix~\ref{app:ordering}.
\end{theorem}

\begin{corollary}\label{spl:cor:relative-ordering-na}
Let $E$ be as in Theorem~\ref{spl:thm:relative-ordering}, and let $D$ have an equivalent nonarchimedean norm. Then
$(E,D)$ is good for Banach values in every degree, and good in every degree if $E=V$. The extension
$E\to c_0(\Sigma,\widehat K)$ of the coordinate map is $\mathrm S^n$-split for every $n\ge1$.
\end{corollary}

The proof is in Appendix~\ref{app:deferred}. For example, the sources $\ell^r(\Sigma,K)$, $1\le r\le2$, are covered, with the targets $c_0$, $\linf$, spaces with an
equivalent nonarchimedean norm, and $\ell^r$, $r>2$, which is $r$-normed by its own norm, since
$\abs{\alpha+\beta}^r\le\max(\abs\alpha,\abs\beta)^r\le\abs\alpha^r+\abs\beta^r$.

\subsection{Parameter splittings}\label{spl:ssec:parameter}

A \emph{parameter splitting} in degree $n$, a bounded projection of $\widehat\otimes^n_\pi P$ onto $\widehat\Gamma^n_\pi P$,
exists if and only if $P$ is polarizable in degree $n$ (Definition~\ref{ob:def:sequences}).

\begin{theorem}[finite-coordinate parameters]\label{thm:finite-parameters}
Let $j\colon W\to Z$ be a closed linear isometry. If $P$ is boundedly linearly isomorphic to $K^d$ and $U\subseteq P$ is open, then $f\colon U\to W$ is analytic if and only if $jf$ is analytic. For the maximum norm on $K^d$, ambient coefficients $b_n$ can be replaced by $W$-valued coefficients $q_n$ with
\[
 \norm{q_n}\le d^n\norm{b_n}\quad(n\ge1).
\]
If $Z$ is nonarchimedean, the factor $d^n$ can be replaced by $1$.
\end{theorem}
\begin{proof}
Theorem~\ref{thm:coefficient-descent} puts every homogeneous diagonal in $W$. Expand it in coordinates as
$b_n(h,\ldots,h)=\sum_{|\alpha|=n}c_\alpha h^\alpha$. Every $c_\alpha$ lies in $W$, by polynomial uniqueness over the
infinite field $K$ in the algebraic quotient $Z/W$. For each $\alpha$, let $i_1(\alpha)\le\dots\le i_n(\alpha)$ be the coordinate word with
these multiplicities, and set
\[
 q_n(h_1,\ldots,h_n)
 =\sum_{|\alpha|=n}c_\alpha\prod_{a=1}^n(h_a)_{i_a(\alpha)}.
\]
Before grouping there are exactly $d^n$ ordered terms, so the triangle inequality gives the bound, and two maximum
estimates give the nonarchimedean refinement. The diagonals agree, and the exponential bounds preserve a positive radius
in Theorem~\ref{thm:coefficient-descent}. The case $d=0$ is immediate.
\end{proof}

Over an incomplete field, finite dimension of the base does not suffice: for $K=\F_p(t)$ and $k\ge p$, the end of
Appendix~\ref{app:no-largest} gives a base of dimension $p+1$ without continuous coordinates and a nonanalytic alternating atlas.

\begin{theorem}[summable parameter spaces]\label{thm:summable-parameters}
Let $j\colon W\to Z$ be a closed linear isometry.
\begin{enumerate}
\item Suppose $K$ and $Z$ are nonarchimedean and $W$ is complete. For $U\subseteq c_0(I,K)$ open, analyticity of $jf$ implies analyticity of $f\colon U\to W$. Coefficients can be lifted without increasing their norms.
\item Suppose $Z$ is complete. Let $a\colon H\to W$ have an ambient bounded homogeneous polynomial representation of fixed degree $d$, that is, $ja(h)=B(h,\ldots,h)$ for a bounded $d$-linear $B\colon H^d\to Z$ on a normed space $H$. For every analytic $\gamma\colon U\to H$, with $U\subseteq\ell^1(I,K)$ open, the composite $a\gamma$ is analytic.
\end{enumerate}
Here $c_0$ carries the supremum norm and $\ell^1$ the ordinary sum norm. The field need not be complete. Each analyticity assertion remains valid for a parameter space $P$ that is a bounded linear retract of the indicated space $V$, with bounded linear $i\colon P\to V$, $r\colon V\to P$ and $ri=\id_P$; in~(1) the lifted coefficients then satisfy $\norm{q_n}\le(\norm r\norm i)^n\norm{b_n}$.
\end{theorem}
The proof is in Appendix~\ref{app:deferred}. Theorems~\ref{thm:finite-parameters} and~\ref{thm:summable-parameters} give
analyticity at every point of an open set, hence the $C^\omega$ conclusion (see Conventions).

\begin{proposition}[parameter splittings of $\ell^1$]\label{spl:prop:gamma-l1}
Let $n\ge1$.
\begin{enumerate}
\item For every set $J$, $\ell^1(J,K)$ is polarizable in degree $n$, by a projection of norm at most $n!$.
\item Every normed space with finite continuous coordinates is polarizable in degree $n$.
\end{enumerate}
\end{proposition}

The proof is in Appendix~\ref{app:deferred}.

\begin{proposition}[degree two]\label{spl:prop:degree-two}
For every normed space $P$, $\cN^2_P=\widehat\Gamma^2_\pi P$. Hence $\cS^2_P$ and $\cG^2_P$ coincide up to the isometry
$\widehat\Lambda^2_\pi P\cong\widehat\otimes^2_\pi P/\cN^2_P$, and $(\mathrm S^2_P)$ holds if and only if $P$ is
polarizable in degree $2$. A section $s$ gives a projection onto $\widehat\Gamma^2_\pi P$ of norm at most $1+\norm s$, and
a projection $R_2$ onto $\widehat\Gamma^2_\pi P$ gives a section of norm at most $1+\norm{R_2}$.
\end{proposition}

The proof is in Appendix~\ref{app:deferred}.

\subsection{The bundle theorem}\label{spl:ssec:bundle}

\begin{corollary}\label{cor:bundle-cases}
Let $M$ be an analytic manifold over $K$, with normed model space $P$. Fiberwise alternating maps give an analytic bifunctor on analytic normed vector bundles and their operator-valued analytic morphisms in each of the following settings:
\begin{center}
\begin{tabular}{p{0.26\linewidth}p{0.63\linewidth}}
\hline
Base model & Hypotheses on the typical fibers\\
\hline
Arbitrary $P$ & $k!\ne0$ in $K$; or $K$ is nonarchimedean and the target fibers admit equivalent nonarchimedean spherically complete norms; or the source fibers have finite continuous coordinates.\\[3pt]
Finite continuous coordinates & Arbitrary normed fibers.\\[3pt]
$c_0(I,K)$, or a bounded linear retract & $K$ nonarchimedean; complete nonarchimedean target fibers.\\[3pt]
$\ell^1(I,K)$, or a bounded linear retract & Complete target fibers.\\
\hline
\end{tabular}
\end{center}
The fiber hypotheses apply to both pairs of bundles at the endpoints of a morphism. Each alternative in the first row specifies a separate setting. At every finite smoothness order and at $C^\infty$, no additional hypothesis on $K,P$ or the fibers is needed.
\end{corollary}
\begin{proof}
The fiber results give a bounded $k$-linear lift $Q$ of $A$, and
$(u_1,\ldots,u_k,g)\mapsto\bigl(m\mapsto g\circ Q(u_1,\ldots,u_k)(m)\bigr)$ is a bounded $(k+1)$-linear lift of the joint
action $B$, which is therefore analytic. The parameter results apply directly to the ambient polynomial along transition
families and local morphism families; the ambient multilinear and operator spaces are complete when the final
value-space fiber is complete, and nonarchimedean when it is nonarchimedean. Theorem~\ref{thm:bundle-lifting} supplies the
bundle construction in every case, and Lemma~\ref{lem:closed-smooth} the smooth assertion.
\end{proof}

Fields of characteristic zero, spherically complete or not, belong to the first row in every finite degree.

\paragraph{\textit{Good pairs of fibers.}}
If $(E,D)$ is good for Banach values in degree $k$ for all $E,D$ in a class $\cC$, then Theorem~\ref{thm:bundle-lifting}
applies over every analytic manifold to bundles with source fibers in $\cC$ and Banach value fibers: the joint action is
polynomial in the precomposition variable (Lemma~\ref{fib:lem:test}, Proposition~\ref{prop:homogeneous-lift}) and linear
in the other. Such classes are the spaces with a lower $r$-estimate, $r<2$ (Theorem~\ref{fib:thm:lower}), and the
symmetric spaces with $c_{00}$ dense and $\cP_<$ bounded (Propositions~\ref{spl:prop:ordering}(2), \ref{fib:prop:split}(1);
under \textup{(L)}, Theorem~\ref{cls:thm:symmetric}). Over complete nonarchimedean $K$, Theorem~\ref{fib:thm:ultra} gives
the same for source fibers of countable type and nonarchimedean Banach value fibers.
%% ---- end sec-05-splittings.tex
%% ---- begin sec-06-combining.tex
\section{Degrees, decompositions and permanence}\label{sec:combining}

\subsection{The degrees in which goodness can fail}\label{cmb:ssec:degrees}
For $a,b\ge1$ and $n=a+b$, let $\wedge^{a,b}_E\colon\Lambda^a_\pi E\otimes_\pi\Lambda^b_\pi E\to\Lambda^n_\pi E$ be the
wedge product, $z\otimes w\mapsto z\wedge w$, and define the shuffle coproduct by
\begin{equation}\label{fib:eq:shuffle}
 \operatorname{sh}^{a,b}_E(x_1\wedge\dots\wedge x_n)=\sum_{\abs S=a}\epsilon(S)\,x_S\otimes x_{S^c},
\end{equation}
where $S$ runs over the $a$\nobreakdash-element subsets of $\{1,\dots,n\}$, $x_S$ is the wedge of the $x_i$, $i\in S$, in
increasing order, and $\epsilon(S)$ is the sign of the permutation of $\{1,\dots,n\}$ that lists $S$ in increasing order
followed by $S^c$ in increasing order. We write $\binom na$ also for its image in $K$.

\begin{lemma}[the Hopf identity]\label{fib:lem:hopf}
The map $\wedge^{a,b}_E$ is contractive, and $\operatorname{sh}^{a,b}_E$ is a well-defined linear map of norm at most
$\binom na$, in every characteristic. Both are natural in bounded operators, and
\begin{equation}\label{fib:eq:hopf}
 \wedge^{a,b}_E\circ\operatorname{sh}^{a,b}_E=\tbinom na\,\id.
\end{equation}
Moreover $\wedge^n_E=\wedge^{a,b}_E\circ(\wedge^a_E\otimes\wedge^b_E)$ and
$\operatorname{alt}^n_E=(\operatorname{alt}^a_E\otimes\operatorname{alt}^b_E)\circ\operatorname{sh}^{a,b}_E$, where
$\otimes^n_\pi E$ is identified with $\otimes^a_\pi E\otimes_\pi\otimes^b_\pi E$.
\end{lemma}

The proof is in Appendix~\ref{app:deferred}.

\begin{theorem}[Lucas step]\label{fib:thm:lucas}
Let $n=a+b$ with $a,b\ge1$ and $\binom na\ne0$ in $K$, let $F$ be normed, and let $G=\Alt^b(D;F)$. If
$A^a_{E,D;G}$ and $A^b_{E,D;F}$ are analytic, so is $A^n_{E,D;F}$. In particular, if $n\ne0$ in $K$ and
$A^{n-1}_{E,D;F}$ is analytic, so is $A^n_{E,D;F}$.
\end{theorem}

The proof is in Appendix~\ref{app:deferred}.

\begin{corollary}[the good degrees]\label{fib:cor:degrees}
Let $\operatorname{char}K=p>0$, and let $\fD$ be one of the following sets of degrees $k\ge1$: the set $\fD(E,D)$ of
those in which $(E,D)$ is good; the set $\fD_B(E,D)$ of those in which it is good for Banach values; or the set
$\fD_K(E,D)$ of those for which $A^k_{E,D;K}$ is analytic.
\begin{enumerate}
\item $\{1,\dots,p-1\}\subseteq\fD$, and $a+b\in\fD$ whenever $a,b\in\fD$ and the addition of $a$
and $b$ in base $p$ has no carry.
\item If $\fD\ne\{1,2,\dots\}$, its least missing element is a power of $p$. Thus
$\fD=\{1,2,\dots\}$ if and only if $\fD$ contains every power of $p$.
\item Let $F$ be normed, and suppose that $A^n_{E,D;F}$ is not analytic for some $n$; let $n_F$ be the least such $n$.
Then $n_F$ is divisible by the least degree in which $(E,D)$ is not good.
\end{enumerate}
\end{corollary}

The proof is in Appendix~\ref{app:deferred}.

\begin{theorem}[hyperplane step]\label{fib:thm:hyperplane}
Give $E\oplus K$ and $D\oplus K$ the maximum norms, let $k\ge2$, and let $F$ be normed. If
$A^k_{E\oplus K,D\oplus K;F}$ is analytic, then so is $A^{k-1}_{E,D;F}$.
\end{theorem}

The proof is in Appendix~\ref{app:deferred}.

\begin{theorem}[no gaps for large duals]\label{cmb:thm:initial-segments}
Let $\operatorname{char}K=p>0$, and let $E$ and $D$ be normed spaces whose continuous duals $E^*$ and $D^*$ are
infinite-dimensional. Then each of $\fD(E,D)$ and $\fD_B(E,D)$ equals $\{1,\dots,p^h-1\}$ for a unique $h\ge1$, or
equals $\{1,2,\dots\}$.
\end{theorem}

If $E\cong E\oplus K$ and $D\cong D\oplus K$ by bounded linear isomorphisms, the same holds for $\fD_K(E,D)$: composing with
these isomorphisms on both sides shows that $A^k_{E,D;K}$ and $A^k_{E\oplus K,D\oplus K;K}$ are analytic together, so
Theorem~\ref{fib:thm:hyperplane} makes $\fD_K(E,D)$ closed under $k\mapsto k-1$ for $k\ge2$, and
Corollary~\ref{fib:cor:degrees}(2) identifies its first missing element.

\begin{theorem}[degree $p$ for one space]\label{cmb:thm:column}
Let $\operatorname{char}K=p>0$, and let $Y$ be a normed space. If $(\mathrm S^p_Y)$ holds, then $(\mathrm S^n_Y)$ holds for
every $n\ge1$. If $(\mathrm R^p_Y)$ holds, then $(\mathrm R^n_Y)$ holds for every $n\ge1$. Both statements also hold
without completion.
\end{theorem}

The proofs of Theorems~\ref{cmb:thm:initial-segments} and~\ref{cmb:thm:column} are in Appendix~\ref{app:degrees}.

\begin{corollary}\label{cmb:cor:column}
Let $\operatorname{char}K=p>0$, and let $E$ and $D$ be normed spaces. If $(\mathrm S^p_E)$ holds, or if
$(\mathrm R^p_D)$ holds, then $(E,D)$ is good for Banach values in every degree. If that condition holds without
completion, then $(E,D)$ is good in every degree.
\end{corollary}

\begin{proof}
By Theorem~\ref{cmb:thm:column} the condition holds in every degree, and Proposition~\ref{fib:prop:split} turns it into
goodness in every degree.
\end{proof}

\paragraph{Split operators.}
For bounded linear $w\colon E\to G$ and $n\ge1$, a \emph{section for $w$ in degree $n$} is a map $s$ as in
Definition~\ref{ob:def:split-operators} with $\widehat\wedge^n_G\circ s=\Lambda^n_\pi w$, and a \emph{retraction for $w$
in degree $n$} is a map $\rho$ with $\rho\circ\alt^n_E=\Lambda^n_\pi w$. In degree $0$ both are $\id_K$,
$\otimes^0_\pi E=K$, and $\wedge^{0,b}$, $\wedge^{a,0}$, $\operatorname{sh}^{0,b}$, $\operatorname{sh}^{a,0}$ are the
canonical identifications. For bounded linear $w_1,w_2\colon E\to G$ and $x_1,\dots,x_n\in E$, the term of
$(w_1+w_2)x_1\wedge\dots\wedge(w_1+w_2)x_n$ with $w_1$ in the positions of $S\subseteq\{1,\dots,n\}$ is
$\epsilon(S)\,(w_1x)_S\wedge(w_2x)_{S^c}$, with $(w_1x)_S$ formed as $x_S$ from the vectors $w_1x_i$; so
\eqref{fib:eq:shuffle} gives
\begin{equation}\label{cmb:eq:binomial}
 \Lambda^n_\pi(w_1+w_2)=\sum_{a=0}^n\wedge^{a,n-a}_G\circ\bigl(\Lambda^a_\pi w_1\otimes\Lambda^{n-a}_\pi w_2\bigr)\circ
 \operatorname{sh}^{a,n-a}_E .
\end{equation}
By~\eqref{cmb:eq:binomial} and Lemma~\ref{fib:lem:hopf}, sections $s^{(l)}_a$ for $w_l$ in all degrees $a\le n$ give the
section $\sum_a(s^{(1)}_a\otimes s^{(2)}_{n-a})\circ\operatorname{sh}^{a,n-a}_E$ for $w_1+w_2$ in degree $n$, and
retractions $\rho^{(l)}_a$ give the retraction $\sum_a\wedge^{a,n-a}_G\circ(\rho^{(1)}_a\otimes\rho^{(2)}_{n-a})$. By
Lemma~\ref{fib:lem:hopf}, if $a,b\ge1$ and $\binom{a+b}a\ne0$ in $K$, sections $s_a$, $s_b$ for $w$ in degrees $a$, $b$
give the section $\binom{a+b}a^{-1}(s_a\otimes s_b)\circ\operatorname{sh}^{a,b}_E$ in degree $a+b$, and retractions
$\rho_a$, $\rho_b$ give the retraction $\binom{a+b}a^{-1}\wedge^{a,b}_G\circ(\rho_a\otimes\rho_b)$; this is the Lucas step
for split operators. Every operator has sections and retractions in the degrees $n$ with $n!\ne0$ in $K$ (the paragraph
after Definition~\ref{ob:def:split-operators}); so if $\operatorname{char}K=p>0$, the degrees in which an operator is
$\mathrm S$-split, or $\mathrm R$-split, contain $1,\dots,p-1$ and are closed under carry-free sums. The
$\mathrm S^{\le k}$-split maps $E\to G$ form a linear space, as do the $\mathrm R^{\le k}$-split ones, and both are closed
under composition with bounded linear maps on either side (Proposition~\ref{ob:prop:split-operators}(2)).

\begin{theorem}[composition]\label{cmb:thm:composition}
Let $q,m\ge1$ and $n=qm$, and suppose that the number $N_{q,m}=(qm)!/\bigl((q!)^m\,m!\bigr)$ of partitions of
$\{1,\dots,n\}$ into $m$ blocks of size $q$ is nonzero in $K$; this holds, for instance, when $q$ and $m$ are powers of
$p=\operatorname{char}K>0$. Let $\abs{N_{q,m}}$ be the absolute value of its image in $K$; elsewhere $N_{q,m}$ and $m!$
are real numbers.
\begin{enumerate}
\item If $w_1\colon E\to G_1$ has a section $s_1$ in degree $q$ and $w_2\colon G_1\to G_2$ has a section $s_2$ in degree
$m$, then $w_2w_1$ has a section $s$ in degree $n$ with $\norm s\le N_{q,m}\norm{s_1}^m\norm{s_2}^q/\abs{N_{q,m}}$.
\item If $w_1\colon E\to G_1$ has a retraction $\rho_1$ in degree $m$ and $w_2\colon G_1\to G_2$ has a retraction
$\rho_2$ in degree $q$, then $w_2w_1$ has a retraction $\rho$ in degree $n$ with
$\norm\rho\le(m!)^{q-1}\norm{\rho_2}^m\norm{\rho_1}^q/\abs{N_{q,m}}$.
\end{enumerate}
If the given sections, respectively retractions, are without completion, so is the new one.
\end{theorem}

\begin{corollary}\label{cmb:cor:composition}
Let $\operatorname{char}K=p>0$ and $h\ge1$.
\begin{enumerate}
\item A composite of $h$ operators that are $\mathrm S^p$-split is $\mathrm S^n$-split for every $n<p^{h+1}$, and a
composite of $h$ operators that are $\mathrm R^p$-split is $\mathrm R^n$-split for every $n<p^{h+1}$.
\item An operator that factors through a normed space $Y$ with $(\mathrm S^p_Y)$ is $\mathrm S^n$-split for every $n$,
and an operator that factors through a normed space $Y$ with $(\mathrm R^p_Y)$ is $\mathrm R^n$-split for every $n$.
\end{enumerate}
\end{corollary}

\begin{theorem}[degree $p$ for polarizability]\label{cmb:thm:polarizable-height}
Let $\operatorname{char}K=p>0$. If a normed space $P$ is polarizable in degree $p$, then $P$ is polarizable in every degree
$n\ge1$.
\end{theorem}

The proofs of Theorems~\ref{cmb:thm:composition} and~\ref{cmb:thm:polarizable-height} and of
Corollary~\ref{cmb:cor:composition} are in Appendix~\ref{app:degrees}.

\subsection{Relative splittings and decompositions}\label{cmb:ssec:decompositions}

Throughout this subsection $E$ and $D$ are normed spaces and $k\ge1$. Since
$\widehat\Lambda^n_\pi\widehat D=\widehat\Lambda^n_\pi D$ (the paragraph before Theorem~\ref{ob:thm:mechanisms}),
$\Lambda^n_\pi v$ maps $\Lambda^n_\pi E$ into $\widehat\Lambda^n_\pi D$ for every $v\in\cL(E,\widehat D)$.

\begin{proposition}[relative splittings]\label{cmb:prop:relative}
Let $G$ be a normed space.
\begin{enumerate}
\item Let $w\colon E\to G$ have a section $s$ in degree $k$, and let $\varepsilon\colon\cL(E,D)\to\cL(G,\widehat D)$ be
bounded linear with $\varepsilon(u)\circ w=u$ for every $u\in\cL(E,D)$. Then $(E,D)$ is good for Banach values in degree
$k$, with the lift
\[
 Q(u_1,\dots,u_k)=\widehat\wedge^k_{\widehat D}\circ\bigl(\varepsilon(u_1)\otimes\dots\otimes\varepsilon(u_k)\bigr)\circ s
\]
of norm at most $\norm\varepsilon^k\norm s$.
\item Let $w'\colon G\to\widehat D$ have a retraction $\rho$ in degree $k$, and let
$\varepsilon\colon\cL(E,D)\to\cL(E,G)$ be bounded linear with $w'\circ\varepsilon(u)=u$ for every $u\in\cL(E,D)$. Then
$(E,D)$ is good for Banach values in degree $k$, with the lift
\[
 Q(u_1,\dots,u_k)=\rho\circ\bigl(\varepsilon(u_1)\otimes\dots\otimes\varepsilon(u_k)\bigr)\circ\alt^k_E
\]
of norm at most $k!\,\norm\varepsilon^k\norm\rho$.
\item If in~(1) the section $s$ takes values in $\otimes^k_\pi G$ and $\varepsilon$ takes values in $\cL(G,D)$, or if
in~(2) the retraction $\rho$ takes values in $\Lambda^k_\pi D$, then $(E,D)$ is good in degree $k$.
\end{enumerate}
\end{proposition}

\begin{proof}
Both maps $Q$ are bounded and $k$-linear, with the stated norms. In~(1), naturality of $\wedge^k$, extended to
completions, gives $Q(u,\dots,u)=\widehat\Lambda^k_\pi\varepsilon(u)\circ\widehat\wedge^k_G\circ s
=\Lambda^k_\pi\varepsilon(u)\circ\Lambda^k_\pi w=\Lambda^k_\pi u$. In~(2), naturality of $\alt^k$ gives
$Q(u,\dots,u)=\rho\circ\alt^k_G\circ\Lambda^k_\pi\varepsilon(u)=\Lambda^k_\pi w'\circ\Lambda^k_\pi\varepsilon(u)
=\Lambda^k_\pi u$. Under the hypotheses of~(3) both lifts take values in $\cL(\Lambda^k_\pi E,\Lambda^k_\pi D)$.
Lemma~\ref{fib:lem:test} concludes.
\end{proof}

\begin{definition}\label{cmb:def:decomposition}
\begin{enumerate}
\item A \emph{source-side relative splitting} of degree $k$ for $(E,D)$ is a triple $(G,w,\varepsilon)$ as in
Proposition~\ref{cmb:prop:relative}(1): $w\colon E\to G$ is $\mathrm S^k$-split and $\varepsilon(u)\circ w=u$. A
\emph{target-side relative splitting} of degree $k$ is a triple $(G,w',\varepsilon)$ as in
Proposition~\ref{cmb:prop:relative}(2): $w'\colon G\to\widehat D$ is $\mathrm R^k$-split and $w'\circ\varepsilon(u)=u$.
\item A \emph{piece} for $(E,D)$ is a tuple $(G,G',w,w',\varepsilon)$ of normed spaces $G$ and $G'$, bounded linear maps
$w\colon E\to G$ and $w'\colon G'\to\widehat D$, and a bounded linear map $\varepsilon\colon\cL(E,D)\to\cL(G,G')$. It is
of \emph{source type} in degree $k$ if $w$ is $\mathrm S^{\le k}$-split, and of \emph{target type} in degree $k$ if
$w'$ is $\mathrm R^{\le k}$-split.
\item A \emph{splitting decomposition} of degree $k$ is a finite family of pieces $(G_i,G'_i,w_i,w'_i,\varepsilon_i)$,
$1\le i\le m$, each of source type or of target type in degree $k$, such that
\begin{equation}\label{cmb:eq:decomposition}
 u=\sum_{i=1}^mw'_i\,\varepsilon_i(u)\,w_i\qquad\text{for every }u\in\cL(E,D).
\end{equation}
A source-side relative splitting $(G,w,\varepsilon)$, respectively a target-side relative splitting
$(G,w',\varepsilon)$, of degree $k$ also counts as a splitting decomposition of degree $k$ with the single piece
$(G,\widehat D,w,\id_{\widehat D},\varepsilon)$, respectively $(E,G,\id_E,w',\varepsilon)$.
\item These structures are \emph{without completion} if their sections take values in the uncompleted tensor powers,
their retractions take values in the uncompleted exterior powers $\Lambda^n_\pi D$, the maps $\varepsilon(u)$ of a
source-side relative splitting take values in $D$, and the maps $w'_i$ of a decomposition take values in $D$.
\end{enumerate}
\end{definition}

A \emph{linear family} is a bounded linear map $\theta\colon P\to\cL(E,\widehat D)$ on a normed space $P$. It is
\emph{liftable in degree $n\ge1$} if there is a bounded $n$-linear map
$Q\colon P^n\to\cL(\Lambda^n_\pi E,\widehat\Lambda^n_\pi D)$ with $Q(h,\dots,h)=\Lambda^n_\pi\theta(h)$ for every $h\in P$.
For $P=\cL(E,D)$, Lemma~\ref{fib:lem:test} shows that $(E,D)$ is good for Banach values in degree $n$ if and only if the
inclusion $\theta_{E,D}\colon\cL(E,D)\to\cL(E,\widehat D)$ is liftable in degree $n$. The proofs of
Theorem~\ref{cmb:thm:decomposition}, Propositions~\ref{cmb:prop:normal-form} and~\ref{cmb:prop:parameter-pieces} and
Lemma~\ref{cmb:lem:subpair} are in Appendix~\ref{app:deferred}.

\begin{theorem}[decompositions]\label{cmb:thm:decomposition}
\begin{enumerate}
\item For a piece $(G,G',w,w',\varepsilon)$ of source type or of target type in degree $k$, the map
$u\mapsto w'\varepsilon(u)w$ is liftable in every degree $n\le k$.
\item Suppose that $u=\sum_{i=1}^m\theta_i(u)$ for every $u\in\cL(E,D)$, where each
$\theta_i\colon\cL(E,D)\to\cL(E,\widehat D)$ is bounded linear and liftable in every degree $n\le k$. Then $(E,D)$ is good
for Banach values in degree $k$. If every $\theta_i(u)$ takes values in $D$ and every lift can be chosen with values in
$\cL(\Lambda^n_\pi E,\Lambda^n_\pi D)$, then $(E,D)$ is good in degree $k$.
\item A splitting decomposition of degree $k$ makes $(E,D)$ good for Banach values in degree $k$, and good in degree $k$
if it is without completion.
\end{enumerate}
\end{theorem}

\begin{proposition}[transfer]\label{cmb:prop:transfer}
Let $w\colon E\to G$ and $w'\colon G'\to\widehat D$ be bounded linear maps of normed spaces, and let
$\varepsilon\colon\cL(E,D)\to\cL(G,G')$ be bounded linear with $u=w'\varepsilon(u)w$ for every $u$. If $(G,G')$ is good for
Banach values in degree $k$, with a lift $Q'$, then so is $(E,D)$, with the lift
\[
 Q(u_1,\dots,u_k)=\widehat\Lambda^k_\pi w'\circ Q'\bigl(\varepsilon(u_1),\dots,\varepsilon(u_k)\bigr)\circ\Lambda^k_\pi w .
\]
\end{proposition}

\begin{proof}
The map $Q$ is bounded and $k$-linear, with values in $\cL(\Lambda^k_\pi E,\widehat\Lambda^k_\pi\widehat D)$ and
$\widehat\Lambda^k_\pi\widehat D=\widehat\Lambda^k_\pi D$; its diagonal is
$\widehat\Lambda^k_\pi w'\circ\Lambda^k_\pi\varepsilon(u)\circ\Lambda^k_\pi w=\Lambda^k_\pi u$. Lemma~\ref{fib:lem:test}
concludes.
\end{proof}

The same formula makes $u\mapsto w'\varepsilon(u)w$ liftable in degree $n$ for every piece $(G,G',w,w',\varepsilon)$ with
$(G,G')$ good for Banach values in degree $n$.

\begin{proposition}[two pieces suffice]\label{cmb:prop:normal-form}
For a bounded bilinear map $b\colon\cL(E,D)\times E\to\widehat D$ let $w_b\colon E\to\cL(\cL(E,D),\widehat D)$ be the
map $w_b(x)(u)=b(u,x)$, and let $w'_b\colon\cL(E,D)\otimes_\pi E\to\widehat D$ be the linear map with
$w'_b(u\otimes x)=b(u,x)$. The pair $(E,D)$ has a splitting decomposition of degree $k$ into pieces of source and target
type if and only if there are bounded bilinear maps $b_S,b_R\colon\cL(E,D)\times E\to\widehat D$ with
$b_S(u,x)+b_R(u,x)=ux$, $w_{b_S}$ $\mathrm S^{\le k}$-split and $w'_{b_R}$ $\mathrm R^{\le k}$-split. In that case $(E,D)$
has such a decomposition with at most two pieces.
\end{proposition}

\begin{lemma}[complemented sub-pairs]\label{cmb:lem:subpair}
Let $E_0$ and $D_0$ be normed spaces, $J\colon E_0\to E$ and $\Pi\colon\widehat D\to\widehat D_0$ bounded linear maps,
and $\Phi\colon\cL(E_0,D_0)\to\cL(E,D)$ a bounded linear map with $\Pi\,\Phi(v)\,J=v$ for every $v\in\cL(E_0,D_0)$. This
holds, for instance, if $\Pi'J=\id_{E_0}$ and $\Pi J'=\id_{D_0}$ for bounded linear maps $\Pi'\colon E\to E_0$ and
$J'\colon D_0\to D$, and $\Phi(v)=J'v\Pi'$.
\begin{enumerate}
\item If $(E,D)$ has a splitting decomposition of degree $k$, so does $(E_0,D_0)$.
\item If $(E,D)$ is good for Banach values in degree $k$, so is $(E_0,D_0)$.
\end{enumerate}
\end{lemma}

\begin{definition}\label{cmb:def:parameter-piece}
A normed space $P$ is \emph{polarizable in degrees $\le k$} if it is polarizable in degree $n$ for every $n\le k$. A
\emph{parameter piece} of degree $k$ for $(E,D)$ is a piece $(G,G',w,w',\varepsilon)$ with $\varepsilon=\varepsilon'\circ\psi$
for bounded linear maps $\psi\colon\cL(E,D)\to P$ and $\varepsilon'\colon P\to\cL(G,G')$, where $P$ is polarizable in
degrees $\le k$; the maps $w$ and $w'$ are arbitrary. A \emph{generalized decomposition} of degree $k$ is a finite family
of pieces satisfying~\eqref{cmb:eq:decomposition}, each of source type or of target type in degree $k$, or a parameter
piece of degree $k$. A source-side or target-side relative splitting of degree $k$ also counts as a generalized
decomposition of degree $k$, as in Definition~\ref{cmb:def:decomposition}(3).
\end{definition}

\begin{proposition}[parameter pieces]\label{cmb:prop:parameter-pieces}
\begin{enumerate}
\item For every set $I$, the space $\ell^1(I,K)$ is polarizable in every degree, and so is every normed space with finite
continuous coordinates, for instance every finite-dimensional normed space over a complete field
(Lemma~\ref{lem:findim}).
\item If $P_1$ and $P_2$ are polarizable in degrees $\le k$, so is $P_1\oplus P_2$ with the maximum norm.
\item If $P$ is polarizable in degree $n$ and $P_0$ is a retract of $P$, then $P_0$ is polarizable in degree $n$.
\item For a parameter piece $(G,G',w,w',\varepsilon)$ of degree $k$, the map $u\mapsto w'\varepsilon(u)w$ is liftable in
every degree $n\le k$. Hence a generalized decomposition of degree $k$ makes $(E,D)$ good for Banach values in degree
$k$.
\end{enumerate}
\end{proposition}

\begin{corollary}\label{cmb:cor:parameter-degrees}
Let $\operatorname{char}K=p>0$.
\begin{enumerate}
\item A parameter piece of degree $p$ is a parameter piece of every degree, and a generalized decomposition of degree $p$
into parameter pieces is one of every degree.
\item If $\cL(E,D)$, with the operator norm, is polarizable in degree $p$, then $(E,D)$ is good for Banach values in every
degree.
\end{enumerate}
\end{corollary}

\begin{theorem}[decompositions of degree $p$]\label{cmb:thm:mixed-degrees}
Let $\operatorname{char}K=p>0$. A splitting decomposition of degree $p$ is a splitting decomposition of degree $p^2-1$, with
the same pieces. Hence $(E,D)$ is then good for Banach values in every degree $n<p^2$, and good in every degree $n<p^2$ if
the decomposition is without completion.
\end{theorem}

\begin{proposition}[target pieces]\label{cmb:prop:target-pieces}
Suppose that $E^*\ne0$ and that $(E,D)$ has a splitting decomposition of degree $k$ all of whose pieces are of target type
in degree $k$. Then $(\mathrm R^k_D)$ holds, without completion if the decomposition is without completion. If moreover
$\operatorname{char}K=p>0$ and $k=p$, then $(E,D)$ is good for Banach values in every degree, and good in every degree if
the decomposition is without completion.
\end{proposition}

The proofs of Corollary~\ref{cmb:cor:parameter-degrees}, Theorem~\ref{cmb:thm:mixed-degrees} and
Proposition~\ref{cmb:prop:target-pieces} are in Appendix~\ref{app:degrees}.

\subsection{The class of an operator family}\label{cmb:ssec:families}

Throughout this subsection $E$, $D$, $F$, $P$ and $P'$ are normed spaces and $k\ge1$; linear families and liftability
are as defined before Theorem~\ref{cmb:thm:decomposition}. As in Lemma~\ref{ob:lem:evaluation}, with the same proof, a
linear family $\theta\colon P\to\cL(E,\widehat D)$ gives a unique bounded linear map
$\ev_\theta\colon\widehat\otimes^k_\pi E\to\Mult^k(P;\widehat\Lambda^k_\pi D)$ with
\[
 \ev_\theta(x_1\otimes\dots\otimes x_k)(h_1,\dots,h_k)=\theta(h_1)x_1\wedge\dots\wedge\theta(h_k)x_k
 \in\widehat\Lambda^k_\pi\widehat D=\widehat\Lambda^k_\pi D;
\]
it has norm at most $\norm\theta^k$, and
$\ev_\theta(z)(h,\dots,h)=\widehat\Lambda^k_\pi\theta(h)(\widehat\wedge^k_Ez)$ for $z\in\widehat\otimes^k_\pi E$ and $h\in P$,
so it maps $\cN^k_E$ into $\Mult^k_0(P;\widehat\Lambda^k_\pi D)$. Let $q$ be the quotient map of the \emph{coefficient
sequence} $\cM^k(P;\widehat\Lambda^k_\pi D)$ of $P$ (Definition~\ref{ob:def:sequences}(4)). Since $q\circ\ev_\theta$
vanishes on $\cN^k_E$, Lemma~\ref{ob:lem:metric-surjection} gives a bounded linear map
$\overline{\ev}_\theta\colon\widehat\Lambda^k_\pi E\to\Poly^k(P;\widehat\Lambda^k_\pi D)$ with
$\overline{\ev}_\theta\circ\widehat\wedge^k_E=q\circ\ev_\theta$. The \emph{class} of $\theta$ in degree $k$ is
\begin{equation}\label{cmb:eq:family-class}
 \ob_k(\theta)=(\ev_\theta|_{\cN^k_E})_*[\cS^k_E]\in\Ext\bigl(\widehat\Lambda^k_\pi E,\Mult^k_0(P;\widehat\Lambda^k_\pi D)\bigr),
\end{equation}
and $\ob_k(\theta_{E,D})=\ob_k(E,D)$, since $\ev_{\theta_{E,D}}=\ev^k_{E,D}$.

For exact sequences $\mathcal Y\colon0\to Y_1\xrightarrow iY_2\xrightarrow qY_3\to0$ and
$\mathcal Y'\colon0\to Y'_1\xrightarrow{i'}Y'_2\xrightarrow{q'}Y'_3\to0$, a \emph{morphism} $\mathcal Y\to\mathcal Y'$ is a
triple of bounded linear maps $\alpha_l\colon Y_l\to Y'_l$, $l=1,2,3$, with $\alpha_2i=i'\alpha_1$ and
$q'\alpha_2=\alpha_3q$. For a bounded linear $\gamma\colon Y'\to Y_3$ from a Banach space $Y'$, the \emph{pullback}
$\gamma^*\mathcal Y\colon0\to Y_1\to Y_\gamma\to Y'\to0$ has $Y_\gamma=\{(y,y')\in Y_2\oplus Y':qy=\gamma y'\}$, normed by
$\norm y+\norm{y'}$, with the maps $y_1\mapsto(iy_1,0)$ and $(y,y')\mapsto y'$. It is exact, its class
$\gamma^*[\mathcal Y]$ depends only on $[\mathcal Y]$, and $\gamma^*[\mathcal Y]=0$ if and only if $\gamma=q\circ L$
for a bounded linear $L\colon Y'\to Y_2$ (Lemma~\ref{ocl:lem:pullback}). Pushouts and pullbacks compose and commute, and a morphism
$(\alpha_1,\alpha_2,\alpha_3)\colon\mathcal Y\to\mathcal Y'$ gives $\alpha_{1*}[\mathcal Y]=\alpha_3^*[\mathcal Y']$
(Lemma~\ref{ocl:lem:morphism}, Corollary~\ref{ocl:cor:functorial}).

\begin{theorem}[the class of a family]\label{cmb:thm:family}
Let $\theta\colon P\to\cL(E,\widehat D)$ be a linear family.
\begin{enumerate}
\item The following are equivalent: $\theta$ is liftable in degree $k$; the restriction $\ev_\theta|_{\cN^k_E}$ extends to a
bounded linear map $\widehat\otimes^k_\pi E\to\Mult^k_0(P;\widehat\Lambda^k_\pi D)$; $\ob_k(\theta)=0$
in~\eqref{cmb:eq:family-class}. A lift $Q$ gives such
an extension of norm at most $\norm\theta^k+\norm Q$, and such an extension $\xi$ gives a lift of norm at most
$\norm\theta^k+\norm\xi$.
\item The triple $(\ev_\theta|_{\cN^k_E},\ev_\theta,\overline{\ev}_\theta)$ is a morphism
$\cS^k_E\to\cM^k(P;\widehat\Lambda^k_\pi D)$, and $\ob_k(\theta)=\overline{\ev}_\theta^{\,*}[\cM^k(P;\widehat\Lambda^k_\pi D)]$.
\end{enumerate}
\end{theorem}

\begin{proof}
(1) Steps~1--3 of the proof of Theorem~\ref{ob:thm:obstruction}, with its correspondence $Q\leftrightarrow Q^\sharp$,
apply with $P$, $\theta(h)$, $\ev_\theta$ and $\Mult^k(P;\widehat\Lambda^k_\pi D)$ in place of $\cL(E,D)$, $u$, $\ev$ and
$\Mult^k(\cL(E,D);\widehat\Lambda^k_\pi D)$, using $\ev_\theta(z)(h,\dots,h)=\widehat\Lambda^k_\pi\theta(h)(\widehat\wedge^k_Ez)$;
the bound $\norm{\ev_\theta}\le\norm\theta^k$ replaces $\norm\ev\le1$ in the constants. The proof of~(2) is in
Appendix~\ref{app:obstruction-class}.
\end{proof}

For $\psi\in\cL(P',P)$ and a Banach space $Y$ let $\psi^\sharp\colon\Mult^k(P;Y)\to\Mult^k(P';Y)$ be the map
$\Phi\mapsto\Phi\circ(\psi\times\dots\times\psi)$, and for a bounded linear map $\chi\colon Y\to Y'$ of Banach spaces let
$\chi_\flat\colon\Mult^k(P;Y)\to\Mult^k(P;Y')$ be the map $\Phi\mapsto\chi\circ\Phi$. Both preserve zero diagonals, and their
pushouts are taken along these restrictions.

\begin{proposition}[naturality]\label{cmb:prop:naturality}
Let $\theta\colon P\to\cL(E,\widehat D)$ be a linear family.
\begin{enumerate}
\item[(a)] For $\psi\in\cL(P',P)$, $\ob_k(\theta\circ\psi)=(\psi^\sharp)_*\ob_k(\theta)$.
\item[(b)] Let $E'$ be a normed space, $w\in\cL(E',E)$, and $\theta_w(h)=\theta(h)\circ w$ for $h\in P$. Then
$\ob_k(\theta_w)=(\widehat\Lambda^k_\pi w)^*\ob_k(\theta)=(\ev_\theta|_{\cN^k_E})_*\operatorname{sp}_k(w)$, with
$\operatorname{sp}_k(w)$ as in Proposition~\ref{ob:prop:split-operators}(1).
\item[(c)] Let $D'$ be a normed space, $w'\colon\widehat D\to\widehat D'$ bounded linear, and $(w'\theta)(h)=w'\circ\theta(h)$
for $h\in P$. Then $\ob_k(w'\theta)=((\widehat\Lambda^k_\pi w')_\flat)_*\ob_k(\theta)$.
\item[(d)] $((\alt^k_D)_\flat)_*\ob_k(\theta)=0$, with $\alt^k_D$ extended to $\widehat\Lambda^k_\pi D$. If $G$ is a normed
space, $w'\colon G\to\widehat D$ is $\mathrm R^k$-split and $\theta_0\colon P\to\cL(E,G)$ is bounded linear, then the family
$h\mapsto w'\circ\theta_0(h)$ is liftable in degree $k$.
\end{enumerate}
\end{proposition}

The proof is in Appendix~\ref{app:obstruction-class}.

\begin{corollary}[transport]\label{cmb:cor:transport}
\begin{enumerate}
\item Let $G$ and $G'$ be normed spaces, and let $w\colon E\to G$, $w'\colon G'\to\widehat D$ and
$\varepsilon\colon\cL(E,D)\to\cL(G,G')$ be bounded linear maps with $u=w'\varepsilon(u)w$ for every $u\in\cL(E,D)$. Then
\[
 \ob_k(E,D)=((\widehat\Lambda^k_\pi w')_\flat)_*(\varepsilon^\sharp)_*(\widehat\Lambda^k_\pi w)^*\ob_k(G,G'),
\]
with $w'$ extended to $\widehat G'$.
\item Let $G$ be a normed space, and let $w\colon E\to G$ and $\varepsilon\colon\cL(E,D)\to\cL(G,\widehat D)$ be bounded linear
maps with $u=\varepsilon(u)\circ w$ for every $u\in\cL(E,D)$. Then $\ob_k(E,D)=(\ev_\varepsilon|_{\cN^k_G})_*\operatorname{sp}_k(w)$,
which vanishes if $w$ is $\mathrm S^k$-split.
\end{enumerate}
\end{corollary}

The proof is in Appendix~\ref{app:deferred}.

\begin{theorem}[sum formula]\label{cmb:thm:cartan}
Let $P_1$ and $P_2$ be normed spaces, $\theta_1\colon P_1\to\cL(E,\widehat D)$ and $\theta_2\colon P_2\to\cL(E,\widehat D)$
linear families, $P=P_1\oplus P_2$ with the maximum norm, $\theta_1\boxplus\theta_2\colon P\to\cL(E,\widehat D)$ the family
$(h_1,h_2)\mapsto\theta_1(h_1)+\theta_2(h_2)$, and $\iota_1\colon P_1\to P$, $\iota_2\colon P_2\to P$ the inclusions. Put
$Y=\widehat\Lambda^k_\pi D$. For $S\subseteq\{1,\dots,k\}$ and $\Phi\in\Mult^k(P;Y)$ let $\Phi_S(h^1,\dots,h^k)$, where
$h^j=(h^j_1,h^j_2)\in P$, be the value of $\Phi$ at the tuple whose $j$-th entry is $(h^j_1,0)$ for $j\in S$ and $(0,h^j_2)$
for $j\notin S$, and let $\Pi^{(a)}\Phi=\sum_{\abs S=a}\Phi_S$. The maps $\Pi^{(a)}$, $0\le a\le k$, are bounded commuting
idempotents on $\Mult^k(P;Y)$ with sum $\id$ that map $\Mult^k_0(P;Y)$ into itself (Lemma~\ref{ocl:lem:idempotents}). Put
$\Mult^{(a)}_0=\Pi^{(a)}\Mult^k_0(P;Y)$ and
$\ob^{(a)}=(\Pi^{(a)}|_{\Mult^k_0(P;Y)})_*\ob_k(\theta_1\boxplus\theta_2)\in\Ext(\widehat\Lambda^k_\pi E,\Mult^{(a)}_0)$.
\begin{enumerate}
\item $\ob_k(\theta_1\boxplus\theta_2)=0$ if and only if $\ob^{(a)}=0$ for every $a$. The map $\iota_1^\sharp$ restricts to an
isomorphism of $\Mult^{(k)}_0$ onto $\Mult^k_0(P_1;Y)$ with $(\iota_1^\sharp)_*\ob^{(k)}=\ob_k(\theta_1)$, and $\iota_2^\sharp$
restricts to an isomorphism of $\Mult^{(0)}_0$ onto $\Mult^k_0(P_2;Y)$ with $(\iota_2^\sharp)_*\ob^{(0)}=\ob_k(\theta_2)$.
\item Let $0<a<k$. Then $\ob^{(a)}=0$ if and only if there is a bounded $k$-linear map
$Q^{(a)}\colon P_1^a\times P_2^{k-a}\to\cL(\widehat\Lambda^k_\pi E,\widehat\Lambda^k_\pi D)$ with
\[
 Q^{(a)}(h_1,\dots,h_1;h_2,\dots,h_2)=\wedge^{a,k-a}_{\widehat D}\circ\bigl(\widehat\Lambda^a_\pi\theta_1(h_1)\otimes
 \widehat\Lambda^{k-a}_\pi\theta_2(h_2)\bigr)\circ\operatorname{sh}^{a,k-a}_E
\]
for $h_1\in P_1$ and $h_2\in P_2$, the maps on the right extended to completions; the right side is the term of index $a$
in~\eqref{cmb:eq:binomial} for $w_1=\theta_1(h_1)$ and $w_2=\theta_2(h_2)$.
\item Let $0<a<k$. If $\theta_1$ is liftable in degree $a$ and $\theta_2$ in degree $k-a$, then $\ob^{(a)}=0$.
\item Suppose that $P_1=P_2$, and let $\Delta_{P_1}\colon P_1\to P$ be the map $h\mapsto(h,h)$. If $\ob^{(a)}=0$ for every
$a$, then $\theta_1+\theta_2=(\theta_1\boxplus\theta_2)\circ\Delta_{P_1}$ is liftable in degree $k$. In particular,
$\theta_1+\theta_2$ is liftable in degree $k$ if $\theta_1$ and $\theta_2$ are liftable in every degree at most $k$.
\item Suppose that $P_1=P_2$ and $\theta_1=\theta_2=\theta$, let $\Delta_{P_1}$ be as in~(4), and let $0<a<k$. Then
$\Delta_{P_1}^\sharp\circ\Pi^{(a)}\circ\ev_{\theta\boxplus\theta}=\binom ka\ev_\theta$. Hence, if $\binom ka\ne0$ in $K$, then
$\ob_k(\theta)=0$ if and only if $(\Delta_{P_1}^\sharp|_{\Mult^{(a)}_0})_*\ob^{(a)}=0$, and $\theta$ is liftable in degree $k$
if it is liftable in degrees $a$ and $k-a$.
\end{enumerate}
\end{theorem}

The proof is in Appendix~\ref{app:obstruction-class}.

\begin{proposition}[composite families]\label{cmb:prop:composition}
Let $\theta\colon P\to\cL(E,\widehat D)$ and $\theta'\colon P'\to\cL(D,\widehat F)$ be linear families. Extend each $\theta'(h')$ to
$\widehat D$, and regard $\theta'$ also as a family $P'\to\cL(\widehat D,\widehat F)$; its class is $\ob_k(\theta')$, since
$\cS^k_{\widehat D}=\cS^k_D$ (the paragraph before Theorem~\ref{ob:thm:mechanisms}). Let
$\theta'\odot\theta\colon P'\otimes_\pi P\to\cL(E,\widehat F)$ be the bounded linear map on the uncompleted projective tensor
product with $(\theta'\odot\theta)(h'\otimes h)=\theta'(h')\circ\theta(h)$.
\begin{enumerate}
\item For $h\in P$ let $\psi_h\colon P'\to P'\otimes_\pi P$ be the map $h'\mapsto h'\otimes h$. Then
$(\psi_h^\sharp)_*\ob_k(\theta'\odot\theta)=(\widehat\Lambda^k_\pi\theta(h))^*\ob_k(\theta')$. For $h'\in P'$ let
$\psi'_{h'}\colon P\to P'\otimes_\pi P$ be the map $h\mapsto h'\otimes h$. Then
$((\psi'_{h'})^\sharp)_*\ob_k(\theta'\odot\theta)=((\widehat\Lambda^k_\pi\theta'(h'))_\flat)_*\ob_k(\theta)$.
\item Let $\operatorname{comp}\colon\cL(D,F)\otimes_\pi\cL(E,D)\to\cL(E,F)$ be the composition map, $v\otimes u\mapsto vu$.
Then $\ob_k(\theta_{D,F}\odot\theta_{E,D})=(\operatorname{comp}^\sharp)_*\ob_k(E,F)$. In particular, if $(E,F)$ is good for
Banach values in degree $k$, then $\theta_{D,F}\odot\theta_{E,D}$ is liftable in degree $k$.
\end{enumerate}
\end{proposition}

The proof is in Appendix~\ref{app:obstruction-class}.

\subsection{Block assemblies and permanence}\label{cmb:ssec:blocks}\label{cmb:ssec:permanence}

The proofs of Propositions~\ref{cmb:prop:blocks} and~\ref{cmb:prop:permanence} are in Appendix~\ref{app:deferred}.

\begin{proposition}[block assemblies]\label{cmb:prop:blocks}
Let $I$ be a totally ordered set, and let $E$ be the $\ell^1$-sum, algebraic or completed, of normed spaces $E_i$,
$i\in I$, with the norm $\norm x=\sum_i\norm{\Pi_ix}$, inclusions $J_i\colon E_i\to E$ and block projections
$\Pi_i\colon E\to E_i$. Let $D_i$ be normed spaces and $J'_i\colon D_i\to D$ bounded linear maps with $\norm{J'_i}\le1$.
Let $\mathcal U\subseteq\cL(E,D)$ be a linear subspace with linear maps $\mathcal U\to\cL(E_i,D_i)$, $u\mapsto u_i$, such
that $\norm{u_i}\le c\norm u$ for a constant $c$, and $ux=\sum_iJ'_iu_i\Pi_ix$ for every $x\in E$ with finitely many
nonzero blocks. Suppose that there is a constant $C\ge1$ such that for every $i\in I$ and every $n\le k$ there is a
bounded $n$-linear map $Q^i_n$ on $\mathcal U_i=\{u_i:u\in\mathcal U\}$ with values in
$\cL(\Lambda^n_\pi E_i,\widehat\Lambda^n_\pi D_i)$, with $Q^i_n(v,\dots,v)=\Lambda^n_\pi v$ and $\norm{Q^i_n}\le C$. Then
there is a bounded $k$-linear map $Q\colon\mathcal U^k\to\cL(\Lambda^k_\pi E,\widehat\Lambda^k_\pi D)$ with
$Q(u,\dots,u)=\Lambda^k_\pi u$ and $\norm Q\le C^kc^k$. In particular, if $\mathcal U=\cL(E,D)$, then $(E,D)$ is good for
Banach values in degree $k$.
\end{proposition}

\begin{proposition}[permanence]\label{cmb:prop:permanence}
Let $k\ge1$, and give finite direct sums the maximum norm.
\begin{enumerate}
\item \textup{(Retracts.)} Let $Y_0$ be a retract of a normed space $Y$. If $Y$ satisfies $(\mathrm S^k_Y)$, or
$(\mathrm R^k_Y)$, or is polarizable in degree $k$, then $Y_0$ satisfies the same condition. If $(E,D)$ is good for Banach
values in degree $k$ and $E_0$, $D_0$ are retracts of $E$, $D$, then $(E_0,D_0)$ is good for Banach values in degree $k$.
\item \textup{(Finite sums.)} For normed spaces $Y_1$ and $Y_2$, $(\mathrm S^n_{Y_1\oplus Y_2})$ holds for every $n\le k$ if
and only if $(\mathrm S^n_{Y_1})$ and $(\mathrm S^n_{Y_2})$ hold for every $n\le k$; the same holds for $(\mathrm R^n)$ and
for polarizability in degrees $\le k$. If $(E_a,D_b)$ is good for Banach values in every degree $n\le k$ for all
$a,b\in\{1,2\}$, so is $(E_1\oplus E_2,D_1\oplus D_2)$.
\item \textup{(Completion of a space.)} Regard the completion $\widehat Y$ of a normed space $Y$ as a normed space over
$K$. Then $Y$ satisfies $(\mathrm S^k_Y)$, $(\mathrm R^k_Y)$, or polarizability in degree $k$, if and only if $\widehat Y$
satisfies the same condition.
\item \textup{(Completion of a pair.)} If $D$ is complete, then $(E,D)$ is good for Banach values in degree $k$ if and only
if $(\widehat E,D)$ is. If $(E,\widehat D)$ is good for Banach values in degree $k$, so is $(E,D)$.
\end{enumerate}
\end{proposition}

Every normed space embeds isometrically in a space with all three splittings, and every Banach space is isomorphic to a
quotient of such a space.

\begin{theorem}[a space with all splittings]\label{cmb:thm:cantor}
Let $G$ be a normed space and $\Omega=\{0,1\}^{\N}$. For $m\ge0$ and $\xi\in\{0,1\}^m$ let $C_\xi\subseteq\Omega$ be the
set of sequences that begin with $\xi$, a dyadic cell of level $m$, of Haar measure $2^{-m}$. A step function of level $m$
is a function $f=\sum_{\xi\in\{0,1\}^m}1_{C_\xi}g_\xi$ on $\Omega$ with $g_\xi\in G$; its norm
$\norm f_1=\sum_\xi2^{-m}\norm{g_\xi}$ does not depend on the level used to write it. Let $L^1(\Omega;G)$ be the
completion of the normed space of step functions. Then for every $k\ge1$ the space $Y=L^1(\Omega;G)$ satisfies
$(\mathrm S^k_Y)$ with a section of norm at most $1$ and $(\mathrm R^k_Y)$ with a retraction of norm at most $1$, and it is
polarizable in degree $k$ with a projection of norm at most $k!$.
\end{theorem}

\begin{corollary}[no hereditary permanence]\label{cmb:cor:no-hereditary}
Let $E$ and $D$ be Banach spaces, let $\operatorname{Ball}(E)$ be the closed unit ball of $E$, and let
$\ell^1(\operatorname{Ball}(E),\Kh)\to E$ be the bounded linear map that sends each unit vector $e_x$ to $x$.
\begin{enumerate}
\item The map $g\mapsto1_\Omega g$ is an isometry of $E$ onto a closed subspace of $L^1(\Omega;E)$, and
$(L^1(\Omega;E),D)$ is good for Banach values in every degree. Likewise $D$ is isometric to a closed subspace of
$L^1(\Omega;D)$, and $(E,L^1(\Omega;D))$ is good for Banach values in every degree.
\item The map $\ell^1(\operatorname{Ball}(E),\Kh)\to E$ is surjective and open, so $E$ is isomorphic to a quotient of $\ell^1(\operatorname{Ball}(E),\Kh)$, and
$(\ell^1(\operatorname{Ball}(E),\Kh),D)$ is good for Banach values in every degree. Likewise $D$ is isomorphic to a quotient of
$\ell^1(\operatorname{Ball}(D),\Kh)$, and $(E,\ell^1(\operatorname{Ball}(D),\Kh))$ is good for Banach values in every degree.
\item Assume condition~\textup{(L)}. Then $(c_0(\N,\Kh),c_0(\N,\Kh))$ is not good for Banach values in degree $k$, and
$c_0(\N,\Kh)$ satisfies neither $(\mathrm S^k)$ nor $(\mathrm R^k)$ and is not polarizable in degree $k$
(Theorem~\ref{fib:thm:ell2}, Corollary~\ref{cls:cor:c0-not-polarizable} and Proposition~\ref{cmb:prop:permanence}). Hence
goodness for Banach values in degree $k$ passes neither to closed subspaces nor to quotients, in either slot of a pair, and
none of $(\mathrm S^k)$, $(\mathrm R^k)$ and polarizability in degree $k$ passes to closed subspaces or to quotients.
\end{enumerate}
\end{corollary}

The proofs of Theorem~\ref{cmb:thm:cantor} and Corollary~\ref{cmb:cor:no-hereditary} are in Appendix~\ref{app:degrees}.

\begin{theorem}[no three-space property in the source slot]\label{cmb:thm:three-space}
Assume condition~\textup{(L)}, and let $D$ be a Banach space with an equivalent nonarchimedean norm and $\dim_{\Kh}D=\infty$,
for example $D=c_0(\N,\Kh)$. Let $Y_1=\ell^1(\operatorname{Ball}(D),\Kh)$, let $\varpi_D\colon Y_1\to D$ be the map
$e_x\mapsto x$, surjective and open by Corollary~\ref{cmb:cor:no-hereditary}(2), and let
$\eta\colon Y_1\to L^1(\Omega;Y_1)$ be the isometry $g\mapsto1_\Omega g$ of Corollary~\ref{cmb:cor:no-hereditary}(1). Put
$Y_2=D\oplus L^1(\Omega;Y_1)$ with the maximum norm and $Y_3=L^1(\Omega;Y_1)/\eta(\ker\varpi_D)$. The maps
$y\mapsto(-\varpi_Dy,\eta y)$ and $(d,v)\mapsto[v+\eta(y_d)]$, where $y_d\in Y_1$
is any element with $\varpi_Dy_d=d$, form an exact sequence $0\to Y_1\to Y_2\to Y_3\to0$ such that:
\begin{enumerate}
\item $(Y_1,D')$ is good for Banach values in every degree for every normed space $D'$;
\item $\cL(Y_3,D)=0$, so $(Y_3,D)$ is good in every degree;
\item $(Y_2,D)$ is not good for Banach values in degree $k$.
\end{enumerate}
\end{theorem}

The proof is in Appendix~\ref{app:obstruction-class}.
%% ---- end sec-06-combining.tex
%% ---- begin sec-07-obstructions.tex
\section{Obstructions}\label{sec:sharpness}

\subsection{The field}\label{obs:ssec:field}
We say that \emph{condition~\textup{(L)}} holds if $\kappa$ is a finite field of characteristic $p$, $0<\varrho<1$, the
Laurent series field $\kappa((X))$ carries the absolute value $\abs a=\varrho^{\ord a}$, $K$ is a subfield of
$\kappa((X))$ on which this absolute value is nontrivial, and $k\ge p$; for example $K=\kappa((X))$ or $K=\kappa(X)$.
Then $K$ is nonarchimedean, and $\Kh$ is the closure of $K$ in the complete field $\kappa((X))$. For $a\in\kappa((X))$
and $n\in\Z$, $\coeff_n(a)\in\kappa$ is the $X^n$\nobreakdash-coefficient of $a$, applied entrywise to vectors and
matrices, and $M_\varrho=\max_{l\ge0}(l+1)\varrho^l$, which exists since $(l+1)\varrho^l\to0$. The group
$\abs{K^\times}$ is a nontrivial subgroup of $\varrho^\Z$, hence $\abs{K^\times}=\varrho^{n_0\Z}$ for an integer
$n_0\ge1$; under~\textup{(L)} we fix $\vartheta\in K$ with $\abs\vartheta=\varrho^{n_0}$, so $\abs{K^\times}=\abs\vartheta^\Z$.

\begin{lemma}[normalization of the field]\label{obs:lem:normalization}
Assume~\textup{(L)}. There are a finite field $\kappa'$ of characteristic $p$ and an isometric field isomorphism
$\sigma$ from $\kappa'((X))$, with the absolute value $\abs g=\abs\vartheta^{\ord g}$, onto $\Kh$ with $\sigma(X)=\vartheta$.
Hence $K$ is isometrically isomorphic to a dense subfield of $\kappa'((X))$. Consequently, a statement whose hypotheses
and conclusion depend only on $K$ as a normed field holds for every $K$ satisfying~\textup{(L)} once it holds for every
subfield $K$ dense in $\kappa((X))$, for every finite field $\kappa$ of characteristic $p$ and every $\varrho\in(0,1)$.
\end{lemma}

The proof is in Appendix~\ref{app:deferred}. When $K$ is dense in $\kappa((X))$, we have $\Kh=\kappa((X))$ and
$\abs{K^\times}=\varrho^\Z$, so $\abs\vartheta=\varrho$.

\subsection{Forcing and the two counterexamples}\label{obs:ssec:counterexamples}
The factorial hypothesis of Proposition~\ref{prop:split-lift}(1) cannot be removed.

\begin{theorem}[operator obstructions]\label{thm:operator-obstructions}
Let $K$ be a nontrivially normed field of characteristic $p>0$, and let $k\ge p$.
\begin{enumerate}
\item There are Banach $K$\nobreakdash-spaces $E,F$ for which
\[
 A^k_{E,E;F}\colon\cL(E,E)\longrightarrow
 \cL\bigl(\Alt^k(E;F),\Alt^k(E;F)\bigr)
\]
is analytic at no point. The target $F$ admits no equivalent nonarchimedean norm.
\item If $K$ is complete and not spherically complete, there are nonarchimedean Banach spaces $E,D$ for which $A^k_{E,D;K}$ is analytic at no point. They can be chosen as closed subspaces of
\[
 \linf(\Lambda,K^{k+1})\quad\text{and}\quad\linf(\Lambda,K^k),
\]
respectively, containing the corresponding $c_0$ spaces, with $\Lambda$ countable.
\end{enumerate}
\end{theorem}

Part~(1) is assembled in Appendix~\ref{app:assembly}. For $\kappa$ finite of characteristic $p$, Theorem~\ref{thm:laurent}
excludes a bounded lift over $\kappa((X))$ for $E_1=\linf(\N,\kappa((X)))$ and $F_1=\widehat\Lambda^k_\pi E_1$
(Lemma~\ref{lem:proj}): through Proposition~\ref{prop:supp} a lift gives the map $\Psi$ of Lemma~\ref{lem:Psi}, which
Theorem~\ref{thm:FF} excludes by Ramsey's theorem, clusters and a staircase matrix. Proposition~\ref{prop:ascent}, which
rests on Ingleton's Theorem~\ref{thm:ingleton}, and Lemma~\ref{lem:descent} pass to every $K$; by
Proposition~\ref{prop:nonultra} the target $F$ has no equivalent nonarchimedean norm. Part~(2) is proved in
Appendix~\ref{app:scalar} from Lemmas~\ref{scalar:fibre} and~\ref{scalar:tail} and van der Put's description of the
continuous dual of a bounded sequence space~\cite{VdP}. In both parts, Proposition~\ref{prop:homogeneous-lift} turns the
absence of a lift into nonanalyticity at every point.

\begin{corollary}[regularity on whole hom spaces]\label{cor:classification}
For every nontrivially normed field $K$, the bifunctor
$\Alt^k\colon\mathsf{Vec}_K^{\mathrm{op}}\times\mathsf{Vec}_K\to\mathsf{Vec}_K$
is $C^\infty$, and is analytic if and only if $k!\ne0$ in $K$.
Both assertions remain true on the full subcategory of Banach spaces.
\end{corollary}
\begin{proof}
Smooth descent applies to the ambient polynomial action. An invertible factorial gives a finite
multilinear lift, and bounded bilinear postcomposition supplies the joint bifunctor action.
If the factorial vanishes, Theorem~\ref{thm:operator-obstructions}(1) excludes analyticity even on
Banach spaces, after fixing the output morphism to be the identity.
\end{proof}

\begin{theorem}[universal scalar classification]\label{thm:scalar-classification}
Let $K$ be a complete nonarchimedean, nontrivially normed field. For any $k\ge0$, the contravariant functor
\[
 \Alt^k(-;K)\colon\mathsf{Vec}_K^{\mathrm{op}}\longrightarrow\mathsf{Vec}_K
\]
is analytic on every hom space if and only if $k!\ne0$ in $K$ or $K$ is spherically complete. The same equivalence holds when the source category is restricted to Banach spaces, or to nonarchimedean Banach spaces.
\end{theorem}

\begin{proof}
An invertible factorial gives a multilinear lift by normalized antisymmetrization. If $K$ is spherically complete, Theorem~\ref{thm:spherical} applies with target $K$. Both arguments permit arbitrary normed sources. If neither condition holds, then $K$ has characteristic $p>0$ with $k\ge p$, and Theorem~\ref{thm:operator-obstructions}(2) supplies nonarchimedean Banach witnesses against all three assertions.
\end{proof}

By Theorem~\ref{thm:duality}(2), the same criterion holds for every nonzero continuous-dual target, and hence for every
nonzero finite-dimensional target, which over complete $K$ is isomorphic to a continuous dual (Lemma~\ref{lem:findim}).

Call a normed space $F$ a \emph{universal target} if $\Alt^k(-;F)$ is analytic on every hom space, that is, if
$A^k_{E,D;F}$ is analytic on $\cL(E,D)$ for all normed spaces $E$ and $D$. Over a nonarchimedean field, every
nonarchimedean spherically complete $F$ is a universal target, by Theorem~\ref{thm:spherical}.

\begin{proposition}[universal targets]\label{prop:universal-zero-dual}
If $\operatorname{char}K=p>0$, $k\ge p$, and $K$ is complete and not spherically complete, then every universal target
has zero continuous dual.
\end{proposition}

\begin{proof}
Let $F$ be a normed space with a nonzero $\lambda\in F^*$, and choose $v\in F$ with $\lambda(v)=1$. Let $E$ and $D$ be the
spaces of Theorem~\ref{thm:operator-obstructions}(2), and for $a\in\Alt^k(D;K)$ let $va\in\Alt^k(D;F)$ be $y\mapsto a(y)v$.
Then $A^k_{E,D;K}(u)(a)=\lambda\circ A^k_{E,D;F}(u)(va)$ for $u\in\cL(E,D)$. Composing with the bounded linear maps
$a\mapsto va$ and $m\mapsto\lambda\circ m$ is a bounded linear map
$\cL(\Alt^k(D;F),\Alt^k(E;F))\to\cL(\Alt^k(D;K),\Alt^k(E;K))$. So if $A^k_{E,D;F}$ were analytic at a point, so would be
$A^k_{E,D;K}$, which Theorem~\ref{thm:operator-obstructions}(2) excludes. Hence $F$ is not a universal target.
\end{proof}

\begin{corollary}[bounded sequences]\label{obs:cor:linf}
Let $\kappa$ be a finite field of characteristic $p$, $K_1=\kappa((X))$ with $\abs X=\varrho$, $k\ge p$, and
$E_1=\linf(\N,K_1)$. Then $(E_1,E_1)$ is not good for Banach values in degree $k$, and neither $(\mathrm S^k_{E_1})$ nor
$(\mathrm R^k_{E_1})$ holds, although $E_1$ is isometrically isomorphic to $c_0(I,K_1)$ for a set $I$, and
$A^k_{E_1,E_1;F}$ is a continuous polynomial for every Banach space $F$ with an equivalent nonarchimedean norm.
\end{corollary}

\begin{proof}
The space $F_1=\widehat\Lambda^k_\pi E_1$ is a Banach space (Lemma~\ref{lem:proj}), and $A^k_{E_1,E_1;F_1}$ is analytic
at no point by Theorem~\ref{thm:laurent}; so $(E_1,E_1)$ is not good for Banach values, and by
Proposition~\ref{fib:prop:split}(1),(2) neither $(\mathrm S^k_{E_1})$ nor $(\mathrm R^k_{E_1})$ holds. Since
$\abs{K_1^\times}=\varrho^\Z$, Corollary~\ref{cor:discrete-targets} and Proposition~\ref{prop:homogeneous-lift} give the
assertion about nonarchimedean $F$. The norm of $E_1$ takes values in $\abs{K_1}$, so $E_1$ has an orthonormal basis and
is isometrically isomorphic to $c_0(I,K_1)$~\cite[\S1, Prop.~1]{Se}.
\end{proof}

\subsection{The \texorpdfstring{$\ell^r$}{l\textasciicircum r} threshold}\label{fib:ssec:local}
Throughout this subsection $K$ satisfies~\textup{(L)}. For a sequence $a$, $\mu_a\colon x\mapsto ax$ denotes
multiplication by $a$. The proof of the next theorem follows the proof of Theorem~\ref{thm:FF}, with norm bounds on the
staircase matrix in place of the support bound.

\begin{theorem}[$\ell^r$ with $r\ge2$]\label{fib:thm:ell2}
Let $2\le r\le\infty$, and let $E$ be a normed space with $c_{00}(\N,K)\subseteq E\subseteq\ell^r(\N,K)$ and the
$\ell^r$ norm, the supremum norm if $r=\infty$. Suppose $Q\colon c_{00}(\N,K)^k\to\cL(\Lambda^k_\pi E,\widehat\Lambda^k_\pi E)$
is $k$\nobreakdash-linear, bounded for the supremum norm on $c_{00}(\N,K)$, and satisfies $Q(a,\dots,a)=\Lambda^k_\pi\mu_a$.
Then for every $m\ge1$, with $1/r=0$ when $r=\infty$,
\begin{equation}\label{fib:eq:quant}
 \begin{gathered}
 \norm Q\ge\frac{(1-\varrho^2)\ln m}{k(k-1)2^k(1+\ln(2\abs\kappa))}\quad\text{if }r=2,\\
 m-1\le M_\varrho k(k-1)\norm Q(4m)^{2/r}4^{(k-2)/r}\quad\text{if }r>2.
 \end{gathered}
\end{equation}
Hence no such $Q$ exists. Consequently $(E,E)$ is not good for Banach values in degree $k$; $A^k_{E,E;F}$ is analytic at
no point for $F=\Lambda^k_\pi E$ and for $F=\widehat\Lambda^k_\pi E$; $(\mathrm S^k_E)$ and $(\mathrm R^k_E)$ fail; and
$\widehat\Lambda^k_\pi E$ has no equivalent nonarchimedean norm.
\end{theorem}

\begin{corollary}[the $\ell^r$ threshold]\label{fib:cor:threshold}
Let $1\le r\le\infty$, and let $E=\ell^r(\N,K)$, or $E=c_0(\N,K)$ if $r=\infty$. Then $(E,E)$ is good for Banach
values in degree $k$ if and only if $r<2$; in that case $(E,D)$ and $(D,E)$ are good for Banach values for every normed
$D$. The same holds, with good in place of good for Banach values, for the finitely supported sequences with the
$\ell^r$ norm.
\end{corollary}

\begin{corollary}[pairs of $\ell^r$ spaces]\label{fib:cor:pairs}
Let $1\le r_1,r_2\le\infty$, and let $E=\ell^{r_1}(\N,K)$ and $D=\ell^{r_2}(\N,K)$, with $c_0(\N,K)$ in place of
$\ell^\infty$. The pair $(E,D)$ is good for Banach values in degree $k$ if $r_1<2$ or $r_2<2$, and it is not if
$r_2=2\le r_1$ or if $r_1,r_2>2$.
\end{corollary}

The proofs are in Appendix~\ref{app:forcing}. The pairs with $r_1=2<r_2$ are good for Banach values in every degree
(Corollary~\ref{cls:cor:lr-table}).

\subsection{The product triangle bound}\label{obs:ssec:ptb}
Throughout this subsection $K$ satisfies~\textup{(L)}. For $A\subseteq\N$ let $1_A$ be its indicator function, for
$n\ge1$ let $[n]$ be the set of the first $n$ elements of $\N$, let $e_t$ ($t\in\N$) be the unit vectors, and for a
sequence $x$ let $\abs x$ be the sequence of absolute values $\abs{x(t)}$.

\begin{definition}\label{obs:def:sequence-spaces}
A \emph{sequence space} is a normed space $D$ with $c_{00}(\N,K)\subseteq D\subseteq\linf(\N,K)$. It is
\emph{symmetric} if there is a norm on $c_{00}(\N,\mathbb R)$, again written $\norm\cdot_D$, with
$\norm f_D=\norm{\abs f}_D$, with $\norm f_D\le\norm g_D$ whenever $\abs f\le\abs g$ pointwise, invariant under every
permutation of $\N$ and equal to $1$ at the unit vectors, such that $\norm x_D=\norm{\abs x}_D$ for $x\in c_{00}(\N,K)$ and
$\norm x_D=\sup_n\norm{x1_{[n]}}_D$ for $x\in D$. A sequence space $E$ is \emph{$\ell^2$\nobreakdash-dominated} if there is
$C_2$ with $\norm x_2\le C_2\norm x_E$ for all $x\in E$; for symmetric $E$ in which $c_{00}(\N,K)$ is dense, this is the
$\ell^2$\nobreakdash-domination of the unit vectors in the sense of Section~\ref{sec:positive}, since $\norm x_E=\sup_n\norm{x1_{[n]}}_E$. The
\emph{fundamental function} of $E$ is $\varphi_E(m)=\norm{1_{[m]}}_E$. For a symmetric sequence space $E$, a sequence
space $D$ and $a\in c_{00}(\N,K)$, $\mu_a$ is bounded from $E$ to $D$ (Lemma~\ref{obs:lem:multiplier}(2)), and
$\norm{\mu_a}_{E\to D}$ denotes its norm.
\end{definition}

For example, $\ell^r(\N,K)$ with $1\le r<\infty$ and $c_0(\N,K)$ are symmetric sequence spaces.

\begin{lemma}[multiplier norms]\label{obs:lem:multiplier}
Let $E$ be a symmetric sequence space.
\begin{enumerate}
\item $\norm x_\infty\le\norm x_E$ for $x\in E$, and $\norm{bx}_E\le\norm x_E$ for $x\in E$ and every finitely
supported $b$ with $\abs b\le1$.
\item For every sequence space $D$ and $a\in c_{00}(\N,K)$, $\mu_a\in\cL(E,D)$ and
$\norm{\mu_a}_{E\to D}=\sup\{\norm{ax}_D:x\in c_{00}(\N,K),\ \operatorname{supp}x\subseteq\operatorname{supp}a,\
\norm x_E\le1\}$.
\item If $D$ is symmetric, then $a\mapsto\norm{\mu_a}_{E\to D}$ is a norm on $c_{00}(\N,K)$ with
$\norm{\mu_a}_{E\to D}\ge\norm a_\infty$ and $\norm{\mu_{e_t}}_{E\to D}=1$. It does not change when the coordinates
of $a$ are relabeled by an injective map $\N\to\N$ or when $a$ is multiplied by a sequence of absolute value $1$, and
$\norm{\mu_a}_{E\to D}\le\norm{\mu_b}_{E\to D}$ when $\abs a\le\abs b$.
\item If $\norm y_D\le C\norm y_E$ for $y\in c_{00}(\N,K)$, then $\norm{\mu_a}_{E\to D}\le C\norm a_\infty$. If $D$ is
symmetric, then $\norm{\mu_a}_{c_0\to D}=\norm a_D$ and $\norm{\mu_a}_{E\to c_0}=\norm a_\infty$, where $c_0=c_0(\N,K)$.
\end{enumerate}
\end{lemma}

The proof is in Appendix~\ref{app:deferred}.

\begin{definition}\label{obs:def:ptb}
Let $E$ be a symmetric sequence space, $D$ a sequence space and $C\ge0$, and let $\cP_<$ be the strict triangle
projection of the unit vectors: $\cP_<(e_t\otimes e_{t'})=e_t\otimes e_{t'}$ if $t<t'$ and $0$ otherwise. The pair $(E,D)$
satisfies the \emph{product triangle bound} \textup{(PTB)} with constant $C$ if
\begin{equation}\label{obs:eq:ptb}
 \norm{\cP_<((ax)\otimes(by))}_{D\otimes_\pi D}\le C\norm{\mu_a}_{E\to D}\norm x_E\norm{\mu_b}_{E\to D}\norm y_E
 \qquad(a,x,b,y\in c_{00}(\N,K)).
\end{equation}
We say that $\cP_<$ is bounded on $E\otimes_\pi E$ if $\norm{\cP_<(x\otimes y)}_{E\otimes_\pi E}\le C\norm x_E\norm y_E$
for some $C$ and all $x,y\in c_{00}(\N,K)$.
\end{definition}

\begin{theorem}[product triangle bound]\label{obs:thm:ptb}
Let $E$ and $D$ be symmetric sequence spaces, and suppose that $(E,D)$ is good for Banach values in degree $k$, with a
lift $Q\colon\cL(E,D)^k\to\cL(\Lambda^k_\pi E,\widehat\Lambda^k_\pi D)$ as in Lemma~\ref{fib:lem:test}. Then $(E,D)$
satisfies \textup{(PTB)} with a constant that depends only on $k$, $K$ and $\norm Q$. No inclusion between $E$ and $D$ is
assumed, and $c_{00}(\N,K)$ need not be dense in either.
\end{theorem}

\begin{lemma}[the triangle of one vector]\label{obs:lem:single-vector}
Suppose that $K$ is dense in $\kappa((X))$, and let $D$ be a sequence space with $\norm y_\infty\le\norm y_D$ for
$y\in D$. Then for every $x\in c_{00}(\N,K)$,
\[
 \norm x_2^2\le\frac{\norm{\cP_<(x\otimes x)}_{D\otimes_\pi D}}{1-\varrho}+\frac{\norm x_\infty^2}{1-\varrho^2}.
\]
\end{lemma}

\begin{proposition}[$\ell^2$-domination]\label{obs:prop:l2-domination}
Let $E$ be a symmetric sequence space and $D$ a sequence space with $\norm y_\infty\le\norm y_D$ for $y\in D$. Suppose
that $\norm\xi_D\le C'\norm\xi_r$ for all $\xi\in c_{00}(\N,K)$, for some $r\in(2,\infty]$ and $C'\ge0$, and that $(E,D)$
satisfies \textup{(PTB)} with constant $C$. Put $\mathfrak a=CC'^2/(1-\abs\vartheta)$ and
$\mathfrak b=1/(1-\abs\vartheta^2)$. Then $\norm x_2\le C_2\norm x_E$ for all $x\in E$, where $C_2^2=\mathfrak a+\mathfrak b$
if $r=\infty$ and $C_2^2=\max\bigl((2\mathfrak a)^{r/(r-2)},2\mathfrak b\bigr)$ if $r<\infty$. In particular $E$ is
$\ell^2$\nobreakdash-dominated.
\end{proposition}

\begin{proposition}[polarizable symmetric spaces]\label{obs:prop:gamma-ptb}
Let $E$ be a symmetric sequence space in which $c_{00}(\N,K)$ is dense. If $(\Gamma^k_E)$ holds, then $\cP_<$ is bounded
on $E\otimes_\pi E$.
\end{proposition}

The proofs of Theorem~\ref{obs:thm:ptb}, Lemma~\ref{obs:lem:single-vector} and Propositions~\ref{obs:prop:l2-domination}
and~\ref{obs:prop:gamma-ptb} are in Appendix~\ref{app:forcing}.

\subsection{Nonarchimedean fibers}\label{obs:ssec:nonarch}
Throughout this subsection $K$ satisfies~\textup{(L)}. The proofs, in Appendix~\ref{app:forcing}, take $K$ dense in
$\kappa((X))$; Lemma~\ref{obs:lem:normalization} extends them to every $K$ satisfying~\textup{(L)}.

\begin{lemma}[marking]\label{fib:lem:marking}
Let $E$ be a normed space with an equivalent nonarchimedean norm and $\dim_{\Kh}\widehat E=\infty$. There are
$e_n\in E$, $n\in\N$, an equivalent nonarchimedean norm $\norm\cdot'$ on $\widehat E$ with
$\norm{\sum_n\xi_ne_n}'=\max_n\abs{\xi_n}$ for finitely supported $(\xi_n)$ in $\Kh$, and a $\Kh$\nobreakdash-linear map
$\theta_E\colon\widehat E\to c_0(\N,\Kh)$, contractive for $\norm\cdot'$, that sends $e_n$ to the $n$th unit vector. If
$E$ is complete, $\widehat\psi_E(\xi)=\sum_n\xi_ne_n$ defines an isometry of $c_0(\N,\Kh)$ into $(E,\norm\cdot')$ with
$\theta_E\widehat\psi_E=\id$. If $K$ is complete, $\theta_E$ takes values in $c_0(\N,K)$.
\end{lemma}

\begin{theorem}[marking obstruction]\label{fib:thm:marking}
Let $E,E'$ be normed spaces, let $\psi\colon c_{00}(\N,K)\to E$ and $\theta\colon E'\to\linf(\N,\Kh)$ be bounded
linear maps, with the supremum norm on $c_{00}(\N,K)$, and let $a\mapsto u_a$ be a bounded linear map
$c_{00}(\N,K)\to\cL(E,E')$ with
\begin{equation}\label{fib:eq:marking}
 \theta\bigl(u_a(\psi x)\bigr)=ax\qquad(a,x\in c_{00}(\N,K)).
\end{equation}
Let $\linf_{\Kh}=\linf(\N,\Kh)$, regarded over $K$. Let $F$ be normed, $\nu\in\Alt^k(E';F)$, and
$\theta_F\colon F\to\widehat\Lambda^k_\pi\linf_{\Kh}$ bounded linear with
$\theta_F\circ\nu=\omega_{\linf_{\Kh}}\circ(\theta,\dots,\theta)$.
Then no bounded $k$\nobreakdash-linear $Q\colon c_{00}(\N,K)^k\to\Alt^k(E;F)$ satisfies $Q(a,\dots,a)=\nu\circ(u_a,\dots,u_a)$,
and $A^k_{E,E';F}$ is analytic at no point. For $F=\Lambda^k_\pi E'$ and $F=\widehat\Lambda^k_\pi E'$, with
$\nu=\omega_{E'}$ and $\theta_F$ induced by $\theta$, the pair $(E,E')$ is therefore not good, respectively not good for
Banach values, in degree $k$.
\end{theorem}

\begin{theorem}[nonarchimedean fibers over a local field]\label{fib:thm:local}
Let $E$ be a normed space with an equivalent nonarchimedean norm.
\begin{enumerate}
\item The following are equivalent: $\dim_{\Kh}\widehat E<\infty$; $(E,E')$ is good for Banach values in degree $k$ for every normed
$E'$; $(E',E)$ is good for Banach values in degree $k$ for every normed $E'$.
\item If $E$ is complete, these conditions are also equivalent to each of: $(E,E)$ is good in degree $k$; $(E,E)$ is
good for Banach values in degree $k$.
\item If $K$ is complete and $E'$ also has an equivalent nonarchimedean norm, then $(E,E')$ is good in degree $k$ if and
only if it is good for Banach values, if and only if $\dim E<\infty$ or $\dim E'<\infty$.
\end{enumerate}
Each failure is shown by an explicit pair whose action with values in the completed exterior power of its second space
is analytic at no point; these values admit no equivalent nonarchimedean norm, while complete nonarchimedean values
never fail, by Corollary~\ref{cor:discrete-targets}.
\end{theorem}

By Theorem~\ref{fib:thm:lower}, $(c_0(\N,K),\ell^1(\N,K))$ is good for Banach values; so the hypothesis on $E'$ in~(3)
cannot be dropped.

\subsection{Characteristic two}\label{obs:ssec:char-two}
Throughout this subsection condition~\textup{(L)} is not assumed. For a field $\Bbbk$, $e_t$ are the unit vectors of $\Bbbk^n$,
$[n]=\{1,\dots,n\}$, $1_S=\sum_{t\in S}e_t$ for $S\subseteq[n]$, also read in $c_{00}(\N,K)$, matrices are
written in sans-serif, $\cP_<\mathsf Z$ keeps the entries of $\mathsf Z$ above the diagonal and replaces the
others by $0$, and $\Omega_w$ is the alternating matrix of $w\in\Lambda^2\Bbbk^n$. In characteristic $2$, $\Omega_w$ is
symmetric with zero diagonal and $\Omega_{x\wedge y}=xy^{\mathsf T}+yx^{\mathsf T}$. For $t\ne t'$, $\pi_{tt'}$ is the
coordinate projection of $\Lambda^2\Bbbk^n$ onto $\Bbbk\,e_t\wedge e_{t'}$, and for $a\in\Bbbk^n$, $\mu_a$ is
multiplication by $a$, so that
\begin{equation}\label{obs:eq:lambda-mu}
\Lambda^2\mu_a=\sum_{t<t'}a_ta_{t'}\,\pi_{tt'} .
\end{equation}
For $L\subseteq[n]$ put $L^c=[n]\setminus L$, and for a vector $a$ of indeterminates $a_L=\prod_{t\in L}a_t$; for a
polynomial $g$, $[a_Lb_{L^c}x_Ly_{L^c}]g$ is the coefficient of the monomial $a_Lb_{L^c}x_Ly_{L^c}$ in $g$. The
Pfaffian of an alternating matrix of even size over a commutative ring of characteristic $2$ is the sum, over the
partitions of the index set into pairs, of the products of the entries at the pairs; its square is the determinant, as
for the generic alternating matrix over $\Z$. Sequence spaces are as in Definition~\ref{obs:def:sequence-spaces}, which
does not use condition~\textup{(L)}, and in this subsection a sequence space $E$ is also assumed to satisfy
$\norm x_\infty\le\norm x_E$ for $x\in E$; for finitely supported $a$, multiplication $\mu_a\colon x\mapsto ax$ then lies
in $\cL(E)$, as $\norm{ax}_E\le\sum_t\abs{a_t}\abs{x_t}\norm{e_t}_E\le\bigl(\sum_t\abs{a_t}\norm{e_t}_E\bigr)\norm x_E$.
Put $\bar\varphi_E(n)=\max_{S\subseteq[n]}\norm{1_S}_E$, which is nondecreasing in $n$. For $n\ge1$,
$c_0(n)=c_0([n],K)$.

Let $\Bbbk$ be a field of characteristic $2$ and $n\ge2$. A \emph{multiplier lift} on $\Bbbk^n$ is a $\Bbbk$-bilinear
map $\Psi\colon\Bbbk^n\times\Bbbk^n\to\operatorname{End}_\Bbbk(\Lambda^2\Bbbk^n)$,
$\Psi(a,b)=\sum_{t,t'}a_tb_{t'}\Psi_{tt'}$, with
\begin{equation}\label{obs:eq:multiplier-lift}
\Psi_{tt'}+\Psi_{t't}=\pi_{tt'}\qquad(t\ne t').
\end{equation}
Its \emph{output matrix} is $\mathsf M_\Psi(a,b,x,y)=\Omega_{\Psi(a,b)(x\wedge y)}$; regarding the entries of
$a,b,x,y$ as indeterminates, it is a matrix over $\Bbbk[a,b,x,y]$. The \emph{ordering lift} $\Psi^<$ has
$\Psi^<_{tt'}=\pi_{tt'}$ for $t<t'$ and $\Psi^<_{tt'}=0$ for $t\ge t'$; thus $\Psi^<(a,b)$ is given by
formula~\eqref{ob:eq:ordering-lift} with $k=2$, $u_1=\mu_a$, $u_2=\mu_b$ and the natural order of $[n]$.

If $\Psi(a,a)=\Lambda^2\mu_a$ for all $a\in\Bbbk^n$, then $\Psi$ is a multiplier lift: by \eqref{obs:eq:lambda-mu},
$\Psi_{tt}=\Lambda^2\mu_{e_t}=0$ and
$\Psi_{tt'}+\Psi_{t't}=\Lambda^2\mu_{e_t+e_{t'}}-\Lambda^2\mu_{e_t}-\Lambda^2\mu_{e_{t'}}=\pi_{tt'}$.

\begin{theorem}[characteristic two]\label{obs:thm:char-two}
\leavevmode
\begin{enumerate}
\item \textup{(Identities.)} Let $\Bbbk$ be a field of characteristic $2$ and $n\ge2$. If $n$ is even, every multiplier lift
$\Psi$ on $\Bbbk^n$ satisfies
\begin{equation}\label{obs:eq:pfaffian-identity}
\sum_{L\subseteq[n],\ \abs L=n/2}\bigl[a_L\,b_{L^c}\,x_L\,y_{L^c}\bigr]\operatorname{Pf}\mathsf M_\Psi(a,b,x,y)=1,
\end{equation}
and some $\mathsf M_\Psi(a,b,x,y)$ with $a,b,x,y\in\{0,1\}^n$ is invertible. For every $\Bbbk$-linear
$\sigma\colon\Lambda^2\Bbbk^n\to\Bbbk^{n\times n}$ with $\wedge\circ\sigma=\id$, where
$\wedge\mathsf Z=\sum_{t<t'}(\mathsf Z_{tt'}+\mathsf Z_{t't})\,e_t\wedge e_{t'}$, there are $S,T\subseteq[n]$ with
$\operatorname{rank}\sigma(1_S\wedge1_T)\ge n-1$, while $\sigma^<(x\wedge y)=\cP_<(xy^{\mathsf T}+yx^{\mathsf T})$,
$x,y\in\Bbbk^n$, defines such a map with strictly upper triangular values, of rank at most $n-1$. For every
$\Bbbk$-linear $\tau\colon\Bbbk^{n\times n}\to\Lambda^2\Bbbk^n$ with $\tau(\Omega_w)=w$ for every $w$, and even $n$,
there is $S\subseteq[n]$ with $\Omega_{\tau(1_S1_S^{\mathsf T})}$ invertible.
\item \textup{(Multiplier maps.)} Let $K$ be a nontrivially normed field of characteristic $2$, $k\ge2$, and $E$ a
sequence space, with $\norm{e_t}_E=1$ for all $t$ if $k\ge3$. Every $k$-linear
$Q\colon c_{00}(\N,K)^k\to\cL(\Lambda^k_\pi E,\widehat\Lambda^k_\pi E)$ with
$\norm{Q(a_1,\dots,a_k)}\le\norm Q\prod_i\norm{a_i}_\infty$ and $Q(a,\dots,a)=\Lambda^k_\pi\mu_a$ satisfies, for every
even $n\ge2$,
\begin{equation}\label{obs:eq:multiplier-bound}
n\le k(k-1)\,\norm Q\,\bar\varphi_E(n)^2 .
\end{equation}
\item \textup{(Splittings.)} Let $K$ be as in \textup{(2)}, $E$ a sequence space and $n\ge2$. Every section $s$ of
$\widehat\wedge{}^2_E$ has $\norm s\ge(n-1)/\bar\varphi_E(n)^2$, every retraction $\rho$ of $\alt^2_E$ has
$\norm\rho\ge\lfloor n/2\rfloor/\bar\varphi_E(n)^2$, and every bounded linear projection $R_2$ of
$\widehat\otimes{}^2_\pi E$ onto $\widehat\Gamma{}^2_\pi E$ has $\norm{R_2}\ge n/\bar\varphi_E(n)^2$. For the
corresponding maps of $c_0(n)$ the bounds are $n-1$, $\lfloor n/2\rfloor$ and $n$.
\item \textup{(Consequences.)} Let $K$ be as in \textup{(2)}, and let $E$ be $c_0(\N,K)$, $\ell^\infty(\N,K)$,
$\ell^r(\N,K)$ with $2<r<\infty$, or a symmetric sequence space with $\sup_nn/\varphi_E(n)^2=\infty$. For every
$k\ge2$, $(E,E)$ is not good for Banach values in degree $k$, $A^k_{E,E;F}$ is analytic at no point for
$F=\Lambda^k_\pi E$ and for $F=\widehat\Lambda^k_\pi E$, neither $(\mathrm S^k_E)$ nor $(\mathrm R^k_E)$ holds, and
$\widehat\Lambda^k_\pi E$ has no equivalent nonarchimedean norm; moreover $(\Gamma^2_E)$ fails. For
$2<r_1,r_2\le\infty$, with $c_0(\N,K)$ in place of
$\ell^\infty(\N,K)$, the pair $(\ell^{r_1}(\N,K),\ell^{r_2}(\N,K))$ is not good for Banach values in
any degree $k\ge2$.
\end{enumerate}
\end{theorem}

The proof is in Appendix~\ref{app:obstruction-class}. In \eqref{obs:eq:multiplier-bound}, $\bar\varphi_E(n)=1$ for
$c_0(\N,K)$ and $\ell^\infty(\N,K)$, and $\bar\varphi_E(n)=n^{1/r}$ for $\ell^r(\N,K)$; so $\norm Q\ge n/(k(k-1))$,
respectively $\norm Q\ge n^{1-2/r}/(k(k-1))$, and for $E=\ell^2(\N,K)$ only $\norm Q\ge1/(k(k-1))$. The proof uses
$Q$ only on vectors supported in $[n+k-2]$, and for $k=2$ the same argument applies to $E=c_0(n)$: for even $n$, every
lift for $(c_0(n),c_0(n))$ in degree $2$, for Banach values, has norm at least $n/2$. Over Laurent series fields over
infinite fields and over densely valued fields, the case $E=c_0(\N,K)$ for odd $p$ is Question~\ref{q:quest:residue}(1),
and the case $E=\ell^2(\N,K)$ for $p=2$ is Question~\ref{q:quest:residue}(2).

\begin{remark}[exactness without splitting]\label{obs:rem:c0-exact}
Let $K$ be nonarchimedean, of any characteristic, and let $D=c_0(\N,K)$. Then
\begin{equation}\label{obs:eq:alt-isometry}
 \norm{\alt^2_Dw}_\pi=2\norm w_\pi\qquad(w\in\hLambda^2_\pi D),
\end{equation}
so the target sequence $\cR^2_D$ of Definition~\ref{ob:def:sequences} is exact. If $K$ has characteristic $2$, then
$\cR^2_D$ does not split: by Lemma~\ref{ob:lem:target} a splitting would give $(\mathrm R^2_D)$, which
Theorem~\ref{obs:thm:char-two}(4) excludes. So in characteristic $2$ the failure of $(\mathrm R^2_D)$ is a failure of
$\cR^2_D$ to split, not of its exactness.
\end{remark}

\begin{proof}
Let $\mathcal O$ be the closed unit ball of $K$, and let $\norm\cdot_\infty$ be the supremum norm.

\emph{Step 1 (reduction).} Since $\alt^2_D(x\wedge y)=x\otimes y-y\otimes x$, the left side
of~\eqref{obs:eq:alt-isometry} is at most the right side. Both sides are continuous in $w$, and the union of the spaces
$\Lambda^2_KK^m$, where $K^m\subseteq D$ is the span of $e_1,\dots,e_m$, is dense in $\hLambda^2_\pi D$. The coordinate
projection of $D$ onto $K^m$ is contractive and fixes $K^m$, so for $w\in\Lambda^2_KK^m$ both norms may be computed in
$c_0(m)$. Put $T=\alt^2_Dw$, an alternating matrix in $K^{m\times m}$.

\emph{Step 2 (elimination).} Let $a=T_{ij}$ be an entry of $T$ of maximal absolute value; then $i\ne j$. Let $B$ be the
block of $T$ on the rows and columns $i,j$, and $C$ the part of the columns $i,j$ outside these rows. As
$\abs{\det B}=\abs a^2$, the matrix $CB^{-1}$ has entries in $\mathcal O$. Let $g$ and $h$ be the unit vectors $e_i$
and $e_j$ plus the first and second columns of $CB^{-1}$, placed in the rows other than $i,j$; so
$\norm g_\infty=\norm h_\infty=1$. The Schur complement gives $T=a(gh^{\mathsf T}-hg^{\mathsf T})+T'$, where
$T'=\alt^2_D(w-a\,g\wedge h)$ vanishes on the rows and columns $i,j$. Repeating with $T'$ gives
$w=\sum_ka_k\,g_k\wedge h_k$ with $\norm{g_k}_\infty=\norm{h_k}_\infty=1$, so $\norm w_\pi\le\sum_k\abs{a_k}$, and
$T=U\operatorname{diag}(a_1J,a_2J,\dots)U^{\mathsf T}$ with $U\in GL_m(\mathcal O)$, whose columns are the
$g_k$, the $h_k$ and unit vectors, and $J=\begin{psmallmatrix}0&1\\-1&0\end{psmallmatrix}$.

\emph{Step 3 (minors).} Let $\mu_1\ge\mu_2\ge\dots\ge\mu_{2r}$ be the numbers $\abs{a_k}$, each taken twice, and
$\mu_{2r+1}=0$. By the Cauchy--Binet formula the minors of $T$ and of $\operatorname{diag}(a_kJ)$ are
$\mathcal O$-linear combinations of each other, so the largest absolute value of an $l\times l$ minor of $T$ is
$\mu_1\cdots\mu_l$. Let $T=\sum_qx_qy_q^{\mathsf T}$ be a finite representation, with costs
$c_q=\norm{x_q}_\infty\norm{y_q}_\infty$ listed in decreasing order. The Cauchy--Binet formula and the ultrametric
inequality bound every $l\times l$ minor of $T$ by $c_1\cdots c_l$. So $\mu_1\cdots\mu_l\le c_1\cdots c_l$ for
$l\le2r$, and these $c_l$ are positive. Since $c-\mu\ge\mu\log(c/\mu)$ for $c,\mu>0$, Abel summation gives
\[
 \sum_{l\le2r}(c_l-\mu_l)\ge\sum_{l\le2r}\mu_l\log\frac{c_l}{\mu_l}
 =\sum_{l\le2r}(\mu_l-\mu_{l+1})\log\frac{c_1\cdots c_l}{\mu_1\cdots\mu_l}\ge0 .
\]
Hence $2\norm w_\pi\le2\sum_k\abs{a_k}=\sum_l\mu_l\le\sum_qc_q$, and taking the infimum over representations gives
$2\norm w_\pi\le\norm{\alt^2_Dw}_\pi$.
\end{proof}
%% ---- end sec-07-obstructions.tex
%% ---- begin sec-08-classification.tex
\section{Which pairs are good}\label{sec:classification}

Throughout this section $\kappa$, $\varrho$ and $K$ are as in condition~\textup{(L)}, and each result states its degrees. Symmetric sequence spaces,
their fundamental functions $\varphi_E$, $\ell^2$\nobreakdash-domination and the product triangle bound \textup{(PTB)}
are those of Definitions~\ref{obs:def:sequence-spaces} and~\ref{obs:def:ptb}. For a sequence space, $\cP_<$ is the
strict triangle projection of the unit vectors $(e_t)_{t\in\N}$, and $c_{00}$, $c_0$ and $\ell^r$ stand for
$c_{00}(\N,K)$, $c_0(\N,K)$ and $\ell^r(\N,K)$.

\subsection{Symmetric spaces}

\begin{theorem}[symmetric spaces]\label{cls:thm:symmetric}
Let $E$ be a symmetric sequence space in which $c_{00}$ is dense. The following are equivalent:
\begin{enumerate}
\item $(E,E)$ is good for Banach values in some degree $k\ge p$;
\item $(c_0,E)$ is good for Banach values in some degree $k\ge p$;
\item $\cP_<$ is bounded on $E\otimes_\pi E$: there is $C$ with
$\norm{\cP_<(x\otimes y)}_{E\otimes_\pi E}\le C\norm x_E\norm y_E$ for all $x,y\in c_{00}$;
\item $(\mathrm S^k_E)$ holds for some $k\ge p$;
\item $(\mathrm R^k_E)$ holds for some $k\ge p$;
\item $E$ is polarizable in some degree $k\ge p$;
\item for every normed $D$, the pairs $(E,D)$ and $(D,E)$ are good for Banach values in every degree.
\end{enumerate}
If these conditions hold, then $(\mathrm S^n_E)$, $(\mathrm R^n_E)$ and $(\Gamma^n_E)$ hold for every $n\ge1$, $E$ is
$\ell^2$\nobreakdash-dominated, and there is $c>0$ with $\varphi_E(m)^2\ge c\,m\ln m$ for every $m\ge1$.
\end{theorem}

The proof is in Appendix~\ref{app:forcing}.

\begin{corollary}[small fundamental functions]\label{cls:cor:symmetric-examples}
Let $E$ be $\ell^r$ with $2\le r<\infty$, or $c_0$, or a symmetric sequence space in which $c_{00}$ is dense and
$\varphi_E(m)^2\le Cm$ for some $C$ and every $m\ge1$. Then for every $k\ge p$ neither $(E,E)$ nor $(c_0,E)$ is good for
Banach values in degree $k$, and $(\mathrm S^k_E)$, $(\mathrm R^k_E)$ and $(\Gamma^k_E)$ fail.
\end{corollary}

\begin{proof}
For $E=\ell^r$ we have $\varphi_E(m)=m^{1/r}$, so $\varphi_E(m)^2\le m$ when $r\ge2$, and for $E=c_0$ we have
$\varphi_E(m)=1$; in both cases $c_{00}$ is dense. If one of the pairs were good in a degree $k\ge p$, or one of the
splittings held, Theorem~\ref{cls:thm:symmetric} would give $c\,m\ln m\le\varphi_E(m)^2\le Cm$ for every $m$, which
fails for large $m$.
\end{proof}

\begin{corollary}\label{cls:cor:c0-not-polarizable}
The space $c_0$ is not polarizable in any degree $k\ge p$.
\end{corollary}

\begin{proof}
This is Corollary~\ref{cls:cor:symmetric-examples} for $E=c_0$.
\end{proof}

\subsection{Symmetric sources against \texorpdfstring{$r$}{r}-normed targets}

For the targets of the next theorem, the necessary condition \textup{(PTB)} of Theorem~\ref{obs:thm:ptb} is also sufficient.

\begin{theorem}[$r$-normed targets]\label{cls:thm:rnorm}
Let $E$ be a symmetric sequence space in which $c_{00}$ is dense, and let $D$ be a symmetric sequence space with an
equivalent $r$\nobreakdash-norm for some $r\in(2,\infty)$ or with an equivalent nonarchimedean norm. The following are
equivalent:
\begin{enumerate}
\item $(E,D)$ is good for Banach values in some degree $k\ge p$;
\item $(E,D)$ satisfies \textup{(PTB)} with some constant;
\item $E$ is $\ell^2$\nobreakdash-dominated;
\item $(E,D)$ is good for Banach values in every degree.
\end{enumerate}
\end{theorem}

The proof is in Appendix~\ref{app:forcing}. Under~(3), the coordinate map of $E$ into $\ell^r(\N,\Kh)$, with $r=3$ when
$D$ has an equivalent nonarchimedean norm, is a relative ordering splitting: it is $\mathrm S^n$\nobreakdash-split for
every $n$, and every $u\in\cL(E,D)$ factors through it (Theorem~\ref{spl:thm:relative-ordering}(1) and Step~1 of its
proof, Corollary~\ref{spl:cor:relative-ordering-na}).

\subsection{The \texorpdfstring{$\ell^r$}{l\textasciicircum r} scale}

\begin{corollary}[the $\ell^r$ scale]\label{cls:cor:lr-table}
Let $1\le r_1,r_2\le\infty$. For the pair $(\ell^{r_1},\ell^{r_2})$, with $c_0$ in place of $\ell^\infty$, the following
are equivalent: it is good for Banach values in some degree $k\ge p$; it is good for Banach values in every degree;
$\min(r_1,r_2)<2$, or $r_1=2<r_2$. Hence $\fD_B(\ell^{r_1},\ell^{r_2})$ is given by the following table, in which $r_i=\infty$ stands
for $c_0$:
\begin{center}
\small
\begin{tabular}{c|cccc}
$r_1\backslash r_2$ & $[1,2)$ & $2$ & $(2,\infty)$ & $\infty$\\ \hline
$[1,2)$ & $\{1,2,\dots\}$ & $\{1,2,\dots\}$ & $\{1,2,\dots\}$ & $\{1,2,\dots\}$\\
$2$ & $\{1,2,\dots\}$ & $\{1,\dots,p-1\}$ & $\{1,2,\dots\}$ & $\{1,2,\dots\}$\\
$(2,\infty)$ & $\{1,2,\dots\}$ & $\{1,\dots,p-1\}$ & $\{1,\dots,p-1\}$ & $\{1,\dots,p-1\}$\\
$\infty$ & $\{1,2,\dots\}$ & $\{1,\dots,p-1\}$ & $\{1,\dots,p-1\}$ & $\{1,\dots,p-1\}$
\end{tabular}
\end{center}
\end{corollary}

The proof is in Appendix~\ref{app:forcing}.

\subsection{Degrees and mechanisms}

\begin{corollary}[one critical degree gives all]\label{cls:cor:dichotomy}
For each of the following pairs $(E,D)$, the set $\fD_B(E,D)$ is $\{1,\dots,p-1\}$ or $\{1,2,\dots\}$:
\begin{enumerate}
\item $(E,E)$ and $(c_0,E)$, for $E$ as in Theorem~\ref{cls:thm:symmetric};
\item $(E,D)$ as in Theorem~\ref{cls:thm:rnorm};
\item $(\ell^{r_1},\ell^{r_2})$ for $1\le r_1,r_2\le\infty$, with $c_0$ in place of $\ell^\infty$;
\item $(E,D)$ for complete $K$ and spaces $E$, $D$ with equivalent nonarchimedean norms.
\end{enumerate}
\end{corollary}

\begin{proof}
By Corollary~\ref{fib:cor:degrees}(1), $\{1,\dots,p-1\}\subseteq\fD_B(E,D)$. It remains to show that each listed pair
that is good for Banach values in one degree $k\ge p$ is good for Banach values in every degree. For~(1)--(3) this is the
equivalence of conditions~(1), (2) and~(7) of Theorem~\ref{cls:thm:symmetric}, of~(1) and~(4) of
Theorem~\ref{cls:thm:rnorm}, and Corollary~\ref{cls:cor:lr-table}. For~(4), Theorem~\ref{fib:thm:local}(3) shows that
goodness in a degree $k\ge p$ forces $\dim E<\infty$ or $\dim D<\infty$; that space has continuous coordinates by
Lemma~\ref{lem:findim}, and Proposition~\ref{prop:split-lift}(2) gives a lift in every degree.
\end{proof}

\begin{remark}[mechanisms]\label{cls:rem:mechanisms}
In every degree, each pair good for Banach values in Theorems~\ref{cls:thm:symmetric}, \ref{cls:thm:rnorm}
and~\ref{fib:thm:local} and Corollary~\ref{cls:cor:lr-table} has a splitting decomposition
(Definition~\ref{cmb:def:decomposition}) with one piece of one of three kinds: an ordering splitting of the source or the
target (Theorem~\ref{fib:thm:lower}, Proposition~\ref{spl:prop:ordering}); the relative ordering splitting of the source
described after Theorem~\ref{cls:thm:rnorm}; or a splitting of the source or the target by finite coordinates on its
completion (proof of Theorem~\ref{fib:thm:local}). In each of these results, goodness implies the stated criterion by
forcing (Theorems~\ref{fib:thm:ell2}, \ref{obs:thm:ptb} and~\ref{fib:thm:marking}).
\end{remark}

\subsection{Spreading spaces}\label{cls:ssec:spreading}

\begin{definition}\label{cls:def:spreading}
A sequence space $E$ is \emph{spreading} if it satisfies Definition~\ref{obs:def:sequence-spaces} for symmetric
sequence spaces with invariance under permutations of $\N$ replaced by the following condition: for $d\ge1$ and
$f\in\mathbb R^d$, the norm $\norm{\sum_{i=1}^df_ie_{t_i}}_E$ is the same for all $t_1<\dots<t_d$ in $\N$.
\end{definition}

Symmetric sequence spaces are spreading. In Banach space terms, the unit vectors of a spreading sequence space are
$1$\nobreakdash-unconditional and $1$\nobreakdash-spreading, that is, $1$\nobreakdash-subsymmetric. A map from a subset of
$\N$ to $\N$ is \emph{increasing} if it preserves the strict order.

\begin{theorem}[spreading spaces]\label{cls:thm:spreading}
Let $E$ and $D$ be spreading sequence spaces.
\begin{enumerate}
\item Lemma~\ref{obs:lem:multiplier} holds for $E$ and $D$, with relabeling by an injective map $\N\to\N$ replaced by
relabeling along a map increasing on the support of $a$; so \textup{(PTB)} is defined for $(E,D)$ by
Definition~\ref{obs:def:ptb}.
\item Theorem~\ref{obs:thm:ptb}, Propositions~\ref{obs:prop:l2-domination} and~\ref{obs:prop:gamma-ptb},
Theorems~\ref{cls:thm:symmetric} and~\ref{cls:thm:rnorm}, and Corollaries~\ref{cls:cor:symmetric-examples}
and~\ref{cls:cor:dichotomy}(1),(2) hold with spreading sequence spaces in place of symmetric ones, and so does
Remark~\ref{cls:rem:mechanisms} for the pairs of Theorems~\ref{cls:thm:symmetric} and~\ref{cls:thm:rnorm}.
\item Over every nontrivially normed field $K$ of characteristic~$2$, without condition~\textup{(L)},
Theorem~\ref{obs:thm:char-two}(4) holds for every spreading sequence space $E$ over $K$ with
$\sup_nn/\varphi_E(n)^2=\infty$.
\end{enumerate}
\end{theorem}

The proof is in Appendix~\ref{app:forcing}. The spreading spaces form a larger class than the symmetric ones, also up
to equivalent norms.

\begin{example}[Garling spaces]\label{cls:ex:garling}
In this example we identify $\N$ with $\{1,2,3,\dots\}$ through its order. Let $0<\alpha<1$, $1\le q<\infty$ and
$\beta=(1-\alpha)/q$. For $f\in c_{00}(\N,\mathbb R)$ put
\[
 \norm f_{\alpha,q}=\sup\Bigl\{\Bigl(\sum_{j=1}^Lj^{-\alpha}\abs{f(t_j)}^q\Bigr)^{1/q}:L\ge1,\ t_1<\dots<t_L\Bigr\},
\]
a norm of the kind introduced by Garling~\cite{Ga} and studied in~\cite{AAW}, and let $E_{\alpha,q}$ be the closure of $c_{00}$ in the space
of all $x\in\linf(\N,K)$ with $\norm x=\sup_n\norm{\abs{x1_{[n]}}}_{\alpha,q}<\infty$.
\begin{enumerate}
\item $E_{\alpha,q}$ is a spreading sequence space in which $c_{00}$ is dense, and
$m^{2\beta}\le\varphi_{E_{\alpha,q}}(m)^2\le(1-\alpha)^{-2/q}m^{2\beta}$ for $m\ge1$.
\item No norm on $c_{00}$ invariant under permutations of $\N$ is equivalent to that of $E_{\alpha,q}$, and the unit
vectors of $E_{\alpha,q}$, in the order of $\N$, satisfy~\eqref{fib:eq:lower} for no $r\in(0,2)$ and $C\ge1$.
\item If $2\beta\le1$, then for every $k\ge p$ neither $(E_{\alpha,q},E_{\alpha,q})$ nor $(c_0,E_{\alpha,q})$ is good for
Banach values in degree $k$, and $(\mathrm S^k_{E_{\alpha,q}})$, $(\mathrm R^k_{E_{\alpha,q}})$ and
$(\Gamma^k_{E_{\alpha,q}})$ fail. If $2\beta<1$, then over every nontrivially normed field $K$ of characteristic~$2$,
with $E_{\alpha,q}$ formed over $K$ in the same way, the pair $(E_{\alpha,q},E_{\alpha,q})$ is good for Banach values in no
degree $k\ge2$.
\item If $2\beta>1$, then $\norm x_2^2\le\bigl(\sum_{m\ge1}m^{-2\beta}\bigr)\norm x^2$ for $x\in E_{\alpha,q}$.
\end{enumerate}
\end{example}

The proof is in Appendix~\ref{app:forcing}. For $2\beta>1$ the necessary conditions of Theorem~\ref{cls:thm:symmetric}
hold: $E_{\alpha,q}$ is $\ell^2$\nobreakdash-dominated by~(4), and $\varphi_{E_{\alpha,q}}(m)^2\ge m^{2\beta}$ grows faster than
$m\ln m$; by~(2), Theorem~\ref{fib:thm:lower} does not apply to the unit vectors in their natural order
(Question~\ref{q:quest:garling}).
%% ---- end sec-08-classification.tex
%% ---- begin sec-09-beyond.tex
\section{Goodness without splitting decompositions}\label{sec:beyond}
The good pairs of Remark~\ref{cls:rem:mechanisms} all have splitting decompositions; this section shows that goodness
does not require one.

\subsection{One-dimensional operator spaces}

\begin{proposition}\label{bey:prop:one-dim}
Let $E$ and $D$ be normed spaces with $\cL(E,D)=K\,u_0$ for some nonzero $u_0\in\cL(E,D)$.
\begin{enumerate}
\item For every $k\ge1$ the pair $(E,D)$ is good in degree $k$, and the map
\[
Q(\lambda_1u_0,\dots,\lambda_ku_0)=\lambda_1\cdots\lambda_k\,\Lambda^k_\pi u_0
\]
is a lift as in Lemma~\ref{fib:lem:test}.
\item If $D$ is complete, then for every $k\ge1$ the pair $(E,D)$ has a splitting decomposition of degree $k$ if and
only if $u_0$ is $\mathrm S^k$-split or $\mathrm R^k$-split.
\item If $\dim D\ge2$, then every bounded linear functional on $E$ is zero.
\end{enumerate}
\end{proposition}

The proof is in Appendix~\ref{app:deferred}.

\subsection{Rigid Banach spaces}

The next theorem transfers the rigid subspace of $L_0$ of Kalton and Roberts~\cite{KR} to separable complete
nonarchimedean fields.

\begin{theorem}\label{bey:thm:rigid}
Let $K$ be a complete nonarchimedean field with nontrivial absolute value that is separable as a metric space. There
are infinite-dimensional Banach spaces $\fX_1,\fX_2,\dots$ over $K$ such that, for all $\alpha,\beta\ge1$,
\[
\cL(\fX_\alpha,\fX_\beta)=
\begin{cases}
K\,\id_{\fX_\alpha},&\alpha=\beta,\\
0,&\alpha\ne\beta.
\end{cases}
\]
\end{theorem}

The proof is in Appendix~\ref{app:rigid}. By Proposition~\ref{bey:prop:one-dim}(3) with $E=D=\fX_\alpha$ and
$u_0=\id_{\fX_\alpha}$, every bounded linear functional on $\fX_\alpha$ is zero.

\begin{theorem}\label{bey:thm:rigid-good}
Let $\kappa$ be a finite field of characteristic $p$ and $K=\kappa((X))$, with the absolute value of condition~(L) of
Section~\ref{sec:sharpness}. For suitable choices in the construction of Theorem~\ref{bey:thm:rigid}, the Banach space
$\fX=\fX_1$ has the following properties.
\begin{enumerate}
\item $\cL(\fX,\fX)=K\,\id_{\fX}$.
\item The pair $(\fX,\fX)$ is good in every degree, hence good for Banach values in every degree.
\item For every $k\ge p$, neither $(\mathrm S^k_{\fX})$ nor $(\mathrm R^k_{\fX})$ holds.
\item For every $k\ge p$, the pair $(\fX,\fX)$ has no splitting decomposition of degree $k$.
\end{enumerate}
\end{theorem}

The proof is in Appendix~\ref{app:rigid}. Since $\fX$ has no nonzero bounded functional, item~(3) is proved by
forcing through $L^r$-valued readouts (Corollary~\ref{rig:cor:readout}) instead of the coordinates of
Section~\ref{sec:sharpness}.

\begin{theorem}\label{bey:thm:rigid-sum}
Let $\kappa$ be a finite field of characteristic $p$ and $K=\kappa((X))$, with the absolute value of condition~(L).
Let $\fX_1,\fX_2,\dots$ be the spaces of Theorem~\ref{bey:thm:rigid}, built with the choices of
Theorem~\ref{bey:thm:rigid-good}, and let $E$ be the completed $\ell^1$-sum of the $\fX_\alpha$, with the norm
$\norm x=\sum_\alpha\norm{x_\alpha}$ for $x=(x_\alpha)_{\alpha\ge1}$.
\begin{enumerate}
\item For a bounded sequence $\lambda=(\lambda_\alpha)_{\alpha\ge1}$ in $K$ put $u_\lambda x=(\lambda_\alpha x_\alpha)_\alpha$.
Every bounded operator on $E$ is some $u_\lambda$, and $\norm{u_\lambda}=\sup_\alpha\abs{\lambda_\alpha}$.
\item The pair $(E,E)$ is good for Banach values in every degree.
\item For every $k\ge p$, the pair $(E,E)$ has no generalized decomposition of degree $k$, that is, no decomposition
into finitely many source, target and parameter pieces (Definition~\ref{cmb:def:parameter-piece}).
\end{enumerate}
\end{theorem}

The proof is in Appendix~\ref{app:rigid}. In Theorem~\ref{bey:thm:rigid-good} the lift $Q$ of
Proposition~\ref{bey:prop:one-dim}(1) is that of one parameter piece through $K$; in Theorem~\ref{bey:thm:rigid-sum}
the lift is a block assembly (Proposition~\ref{cmb:prop:blocks}) of infinitely many such lifts, which no finite family
of pieces replaces. Question~\ref{q:quest:generation} asks whether these mechanisms and the operations of
Section~\ref{sec:combining} yield every pair good for Banach values.
%% ---- end sec-09-beyond.tex
%% ---- begin sec-10-geometry.tex
\section{Bundles, tangent bundles and differential forms}\label{sec:geometry}

\subsection{From a bad operator family to changes of frame}

Transition functions take values in invertible operators; for a single Banach fiber, invertible families impose the same condition as all operator families (Proposition~\ref{fam2:prop:reductions}(1)).

\begin{proposition}[shear realization]\label{prop:shear}
Suppose $A=A^k_{E,D;F}$ is not analytic at $u_0\in\cL(E,D)$. Set $H=D\oplus E$, with the maximum norm, and
\[
 g(u)(d,e)=(d+ue,e).
\]
Both $g$ and its inverse family are affine analytic. Nevertheless, the family of pullback operators $g(u)^*$ on $\Alt^k(H;F)$ is not analytic at $u_0$.
\end{proposition}

\begin{proof}
The inverse family is $g(-u)$. Let $\pi_D\colon H\to D$ be the projection and $i_E(e)=(0,e)$. Define bounded linear maps $J\colon\Alt^k(D;F)\to\Alt^k(H;F)$ by $Jm=m\circ\pi_D^{\times k}$ and $R\colon\Alt^k(H;F)\to\Alt^k(E;F)$ by restriction along $i_E$. Since $\pi_Dg(u)i_E=u$,
\begin{equation}\label{eq:shear-compression}
 R\circ g(u)^*\circ J=A(u).
\end{equation}
Bounded linear compression preserves analyticity, so $u\mapsto g(u)^*$ is not analytic at $u_0$.
\end{proof}

\begin{corollary}[shear bundle]\label{cor:shear-bundle}
Suppose $A=A^k_{E,D;F}$ is not analytic at $u_0\in\cL(E,D)$, and let $H$ and $g$ be as in Proposition~\ref{prop:shear}. Take the base $M=\cL(E,D)$ and the trivial bundle $M\times H$. Its product trivialization and the trivialization $(u,z)\mapsto(u,g(u)z)$ form an analytic bundle atlas, and the induced alternating atlas on the trivial bundle with fiber $\Alt^k(H;F)$ is not analytic. Equivalently, the analytic bundle automorphism $(u,z)\mapsto(u,g(u)z)$ does not induce an analytic morphism on the alternating product bundle.
\end{corollary}

\begin{proof}
The transitions are $g(u)$ and $g(-u)$, affine analytic by Proposition~\ref{prop:shear}. The induced alternating transitions are pullback by $g(u)$ and $g(-u)$, and the induced morphism is locally pullback by $g(u)$; by Proposition~\ref{prop:shear}, $u\mapsto g(u)^*$ is not analytic at $u_0$.
\end{proof}

Thus every failure of the hom-space criterion has a bundle realization, and Theorem~\ref{thm:operator-obstructions} gives fibers for which it occurs at every point of the base. Applied to analytic families $x\mapsto u(x)$, the shear also makes the familywise hypothesis of Theorem~\ref{thm:bundle-lifting} necessary for all fibers and frames; Section~\ref{sec:families} studies it (Corollary~\ref{fam2:cor:bundles}, Theorems~\ref{fam2:thm:polarizable} and~\ref{fam2:thm:bases}).

\subsection{Tangent bundles and individual forms}

\begin{proposition}[tangent and pullback obstructions]\label{prop:tangent-obstruction}
Suppose $K$ is complete, has characteristic $p>0$, is not spherically complete, and $k\ge p$.
\begin{enumerate}
\item There is an analytic Banach manifold with two global analytic charts whose induced transition on scalar alternating $k$\nobreakdash-forms is analytic at no point.
\item There are Banach spaces $N,Y$, a polynomial analytic map $h\colon N\to Y$, and an analytic map $\omega_0\colon Y\to\Alt^k(Y;K)$ such that $h^*\omega_0$, expressed in the product coordinates on $N$, is analytic at no point as an $\Alt^k(N;K)$\nobreakdash-valued map. After inclusion in $\Mult^k(N;K)$ it is analytic everywhere.
\end{enumerate}
\end{proposition}

The proof is in Appendix~\ref{app:deferred}.

Assertion~(2) separates intrinsic analyticity of individual forms from analyticity of bundle transitions; Section~\ref{sec:forms} builds a calculus on the ambient notion.

The chart $\psi$ of the proof of Proposition~\ref{prop:tangent-obstruction} realizes any nowhere analytic action $A^k_{E,D;F}$; with Theorem~\ref{thm:operator-obstructions}(1) it works over every field of characteristic $p$.
\begin{proposition}[tangent obstructions with other values]\label{prop:tangent-values}
Let $K$ be a nontrivially normed field of characteristic $p>0$, complete or not, and let $k\ge p$.
\begin{enumerate}
\item There are a Banach space $F$ and an analytic Banach manifold with two global analytic charts whose induced transition on $F$\nobreakdash-valued alternating $k$\nobreakdash-forms is analytic at no point.
\item There is an analytic Banach manifold with two global analytic charts whose induced transition on tangent-valued alternating $k$\nobreakdash-forms is analytic at no point.
\item If $K$ is complete and not spherically complete, a manifold as in~(2) exists with a nonarchimedean Banach model.
\end{enumerate}
The target $F$ constructed for~(1) admits no equivalent nonarchimedean norm; for $K=\F_q((X))$ as in Corollary~\ref{lau:cor:dichotomy}, the model space can be chosen nonarchimedean. The model space constructed for~(2) admits no equivalent nonarchimedean norm.
\end{proposition}
The proof is in Appendix~\ref{app:deferred}.

\paragraph{\textit{Bourbaki's question.}}
For $k\ge1$, $\Alt^k$ on Banach spaces is the functor $\alpha_k$ of~\cite[n\textsuperscript{o}~7.8.1]{Bou}, and $A^k_{E,D;F}(u)=\alpha_k(\id_F,u)$. There $K$ is $\mathbb R$, $\mathbb C$ or complete nonarchimedean, norms over a nonarchimedean $K$ are nonarchimedean, and class $C^r$ means analytic if $K\ne\mathbb R$, in a sense that agrees with ours between Banach spaces. In that setting $\alpha_k$ is a vector functor of class $C^r$ if $k!\ne0$ in $K$ (Corollary~\ref{cor:classification}), and for every $k$ if $K$ is discretely valued (Corollary~\ref{cor:discrete-targets} and the proof of Corollary~\ref{cor:bundle-cases}). It is not one if $k!=0$ in $K$ and $K$ is not spherically complete, since Theorem~\ref{thm:operator-obstructions}(2) gives nonarchimedean Banach spaces $E,D$ with $A^k_{E,D;K}$ analytic at no point; so the definition in~\cite[n\textsuperscript{o}~7.8.1]{Bou} needs a restriction. By Corollary~\ref{cor:shear-bundle} and Proposition~\ref{prop:tangent-obstruction}, it is needed for the bundles themselves. Source spaces of countable type suffice in every degree, by the following corollary of Theorem~\ref{fib:thm:ultra}.

\begin{corollary}[countable type and Bourbaki's $\alpha_k$]\label{fib:cor:bourbaki}
Let $K$ be complete and nonarchimedean, let $E,D$ be normed spaces, at least one of them of countable type, and let
$F$ be a nonarchimedean Banach space. For every normed $F'$, the joint action $B$ of
Proposition~\ref{prop:ambient-polynomial} on $\cL(E,D)\times\cL(F,F')$ has a bounded $(k+1)$\nobreakdash-linear lift.
In the setting of the paragraph \emph{Bourbaki's question}, $\alpha_k$ is therefore a
vector functor of the class considered there on source spaces of countable type, in every degree.
\end{corollary}

\begin{proof}
If $Q$ lifts $A^k_{E,D;F}$ (Theorem~\ref{fib:thm:ultra}), then $(u_1,\dots,u_k,g)\mapsto(m\mapsto g\circ Q(u_1,\dots,u_k)(m))$
lifts $B$, as in the proof of Corollary~\ref{cor:bundle-cases}.
\end{proof}

Our results do not decide the case of a densely valued, spherically complete field with $k!=0$ (Section~\ref{sec:questions}).

\paragraph{\textit{A derivative adapted to positive characteristic.}}
Ya\"el Dillies asked whether $A=A^k_{E,D;F}$ is analytic in the calculus of~\cite[Definition~4.7 and Theorem~5.4]{BGN}, where a $C^k$ map $f$ has Taylor coefficients $a_n(x,h)$ with $n!\,a_n(x,h)=d^nf(x)(h,\ldots,h)$. There $A$ is smooth~\cite[Remark~4.2 and Lemma~10.1]{BGN} and $a_k(u_0,u)=A(u)$, while $d^kA(u_0)(u,\ldots,u)=k!\,A(u)$ vanishes if $k!=0$. Analyticity is not restored: it asks for a bounded $k$-linear lift (Proposition~\ref{prop:homogeneous-lift}), and Theorem~\ref{thm:operator-obstructions} gives spaces with none when $k!=0$ in $K$. Section~\ref{sec:dp} changes the coefficients instead: over every nonarchimedean field, with bounded polynomial laws as Taylor coefficients, $A$ is divided-power analytic on $\cL(E,D)$ in every characteristic.

\subsection{Analytic differential forms}\label{sec:forms}

Throughout this subsection $K$ is complete and nontrivially normed, manifold models are normed spaces, and charts and maps are $C^\omega$, so that their derivatives are analytic; on the Banach coefficient spaces this is power-series analyticity.

\begin{definition}
An \emph{ambient analytic scalar $k$\nobreakdash-form} on $M$ is a family $\eta_x\in\Alt^k(T_xM;K)$ whose representative in each chart is analytic after the inclusion
\[
 j\colon\Alt^k(P;K)\hookrightarrow\Mult^k(P;K).
\]
Write $\Omega^k_{\mathrm{amb}}(M)$ for the space of such forms. In degree zero this is the space of analytic scalar functions.
\end{definition}

Under a coordinate change $\psi=\phi'\circ\phi^{-1}$ the representatives satisfy $j\eta_\phi(y)=(j\eta_{\phi'}(\psi(y)))\circ(D\psi(y),\ldots,D\psi(y))$, which is analytic by bounded multilinear composition; so any analytic atlas may be used.

For an ambient $k$\nobreakdash-form $\eta$, an ambient $l$\nobreakdash-form $\zeta$ and $S\subseteq\{1,\ldots,k+l\}$, let $v_S$ list the arguments $v_i$, $i\in S$, in increasing order of $i$, and let $\epsilon(S)$ be the shuffle sign of Section~\ref{cmb:ssec:degrees}. The product is the shuffle formula
\begin{equation}\label{eq:form-shuffle}
 (\eta\wedge\zeta)(v_1,\ldots,v_{k+l})
 =\sum_{\abs{S}=k}\epsilon(S)\eta(v_S)\zeta(v_{S^c}),
\end{equation}
a bounded bilinear operation on multilinear coefficient spaces, of norm at most $\binom{k+l}{k}$, which is alternating on alternating inputs. On an open subset of a model space $P$, the exterior derivative of an ambient coefficient map $\eta$ of degree $k$ is
\begin{equation}\label{eq:exterior-derivative}
 (d\eta)_y(v_0,\ldots,v_k)
 =\sum_{i=0}^k(-1)^iD\eta(y)(v_i)
       (v_0,\ldots,\widehat v_i,\ldots,v_k).
\end{equation}

\begin{theorem}[ambient differential calculus]\label{thm:ambient-forms}
Ambient analytic forms define a sheaf of graded algebras on every analytic manifold. Analytic maps induce pullbacks, and there is an exterior derivative
\[
 d\colon\Omega^k_{\mathrm{amb}}(M)\longrightarrow\Omega^{k+1}_{\mathrm{amb}}(M).
\]
The product is associative and graded commutative, $d^2=0$, pullback commutes with both product and exterior derivative, and
\[
 d(\eta\wedge\zeta)=d\eta\wedge\zeta+(-1)^k\eta\wedge d\zeta
 \qquad(\eta\in\Omega^k_{\mathrm{amb}}(M)).
\]
Thus $\Omega^\bullet_{\mathrm{amb}}$ is a contravariant differential graded algebra construction in every characteristic.
\end{theorem}
The proof is in Appendix~\ref{app:deferred}.

\begin{proposition}[comparison with alternating-bundle sections]\label{prop:ambient-agreement}
Every intrinsically analytic alternating-valued coefficient map is ambient analytic. If, for the model space $P$, the inclusion $j$ has a bounded linear retraction, then the natural alternating cotangent atlas is analytic and its analytic sections are precisely the ambient analytic forms. The same conclusion holds when $P$ has finite continuous coordinates.
\end{proposition}

\begin{proof}
The first statement follows by composition with $j$. For a bounded retraction $R$, the identity $\eta=R(j\eta)$ gives the converse. Moreover, $(u_1,\ldots,u_k)\mapsto\bigl[m\mapsto R\bigl(m\circ(u_1,\ldots,u_k)\bigr)\bigr]$ is a bounded multilinear lift of the precomposition action, so the induced alternating cotangent transitions are analytic. For finite continuous coordinates, analytic descent through closed subspaces applies to the coefficient map itself, and Corollary~\ref{cor:bundle-cases} supplies the analytic bundle atlas.
\end{proof}

The retraction exists if $k!$ is invertible, by normalized antisymmetrization, and if $K$ is spherically complete and nonarchimedean, by Theorem~\ref{thm:spherical}.

\paragraph{\textit{Exact constant forms in degree $p$.}}
Here $K$ is any nontrivially normed field of characteristic $p>0$, $Y$ and $F$ are normed, and $\mathcal A_Y\colon\Mult^p(Y;F)\to\Alt^p(Y;F)$ is unnormalized antisymmetrization (Lemma~\ref{ob:lem:target}). Define $(d\eta)_y\in\Alt^p(Y;F)$ by~\eqref{eq:exterior-derivative} for every map $\eta$ from an open subset of $Y$ to $\Alt^{p-1}(Y;F)$ differentiable at $y$. Constant forms are closed, and $(\mathrm R^p_Y)$ is the exactness of one of them.
\begin{proposition}[exact constant forms]\label{fib:prop:derham}
For $\omega\in\Alt^p(Y;F)$ the following are equivalent:
\begin{enumerate}
\item there are a map $\eta$ from an open subset of $Y$ to $\Alt^{p-1}(Y;F)$ and a point $y$ at which $\eta$ is
differentiable with $(d\eta)_y=\omega$;
\item $\omega=d\eta$ on $Y$ for a bounded linear map $\eta\colon Y\to\Alt^{p-1}(Y;F)$;
\item $\omega=\mathcal A_Y\beta$ for some $\beta\in\Mult^p(Y;F)$.
\end{enumerate}
For $F=\Lambda^p_\pi Y$ and $\omega=\omega_Y$ they hold if and only if $\operatorname{alt}^p_Y$ has a bounded linear left
inverse, and then $(E,Y)$ is good in degree $p$ for every normed $E$. For $F=\widehat\Lambda^p_\pi Y$ and
$\omega=\omega_Y$ they hold if and only if $(\mathrm R^p_Y)$ holds, and then $(E,Y)$ is good for Banach values in degree
$p$ for every normed $E$.
\end{proposition}

The proof is in Appendix~\ref{app:deferred}.

Under condition~\textup{(L)} of Section~\ref{obs:ssec:field}, for $Y=c_0(\N,K)$ the closed constant form $\omega_Y$ with values in $\widehat\Lambda^p_\pi Y$ has no primitive near any point, by Theorem~\ref{fib:thm:ell2}.

\subsection{Divided-power analyticity}\label{sec:dp}

Throughout this subsection $K$ is nonarchimedean, $\varpi\in K$ satisfies $0<\abs\varpi<1$, and $P$, $W$ and $Y$ are
normed spaces, not necessarily complete or nonarchimedean.

\begin{definition}\label{dp:def:laws}
Let $n\ge0$.
\begin{enumerate}
\item A map $a\colon P\to W$ is a \emph{bounded polynomial law of degree $n$} if, for all $h_1,\dots,h_r\in P$, the map
$s\mapsto a(s_1h_1+\dots+s_rh_r)$ on $K^r$ is a homogeneous polynomial of degree $n$ with coefficients in $W$, and
$\norm a=\sup_{h\ne0}\norm{a(h)}/\norm h^n$ is finite. For $\alpha=(\alpha_1,\dots,\alpha_r)$ with integers
$\alpha_i\ge0$ and $\abs\alpha=\alpha_1+\dots+\alpha_r=n$, $a^{[\alpha]}(h_1,\dots,h_r)\in W$ is the coefficient of
$s_1^{\alpha_1}\cdots s_r^{\alpha_r}$, unique since $K$ is infinite.
\item A map $a\colon P\to W$ is \emph{ambient of degree $n$} if $ja(h)=\beta(h,\dots,h)$ for all $h\in P$, for a linear
isometry $j$ of $W$ into a normed space $Z$ and a bounded $n$\nobreakdash-linear map $\beta\colon P^n\to Z$.
\item A map $f$ from an open subset $U$ of $P$ to $W$ is \emph{divided-power analytic}, or \emph{DP}, at $x_0\in U$ if
there are $R>0$ and bounded polynomial laws $a_n\colon P\to W$ of degree $n$, $n\ge0$, with $\sup_n\norm{a_n}R^n<\infty$
and $f(x_0+h)=\sum_na_n(h)$ for $\norm h<R$. It is \emph{DP on $U$} if it is DP at every point of $U$.
\end{enumerate}
\end{definition}

Without the bound on $\norm a$, the maps in Definition~\ref{dp:def:laws}(1) are the maps $P\to W$ induced by the
polynomial laws on $(P,W)$ in the sense of Roby~\cite{Ro} that are homogeneous of degree $n$, and each of these maps
determines its law, since $K$ is an infinite field and the linear forms on the vector space $W$ separate its
points~\cite[Theorem~I.1 and Propositions~I.1, I.8 and~IV.6]{Ro}.

\begin{theorem}[divided-power calculus]\label{dp:thm:calculus}
\leavevmode
\begin{enumerate}
\item \textup{(Mixed terms.)} If $b\colon W\to Y$ is a bounded polynomial law of degree $m$, then
\[
 \norm{b^{[\alpha]}(w_1,\dots,w_r)}\le(8/\abs\varpi)^m\norm b\,(\norm{w_1}+\dots+\norm{w_r})^m
 \qquad(w_1,\dots,w_r\in W,\ \abs\alpha=m).
\]
Every bounded polynomial law is continuous and DP on its whole domain.
\item \textup{(Composition.)} Let $f$ map an open subset of $P$ to $W$ and $g$ an open subset of $W$ to $Y$. If $f$ is
DP at $x_0$ and $g$ is DP at $f(x_0)$, then $g\circ f$, which is defined near $x_0$, is DP at $x_0$.
\item \textup{(Comparison.)} A map ambient of degree $n$, for instance the diagonal of a bounded $n$\nobreakdash-linear
map, is a bounded polynomial law of degree $n$, and a map analytic at a point is DP there. If $P$ has finite
coordinates, a map DP at a point is analytic there.
\item \textup{(Alternating maps.)} Let $E$, $D$, $F$ and $F'$ be normed spaces and $k\ge1$. The maps
$A^k_{E,D;F}\colon\cL(E,D)\to\cL(\Alt^k(D;F),\Alt^k(E;F))$, $\cL(E,D)\to\cL(\Lambda^k_\pi E,\Lambda^k_\pi D)$,
$u\mapsto\Lambda^k_\pi u$, and $\cL(E,D)\to\cL(\widehat\Lambda^k_\pi E,\widehat\Lambda^k_\pi D)$,
$u\mapsto\widehat\Lambda^k_\pi u$, are ambient of degree $k$, and the joint action $B$ of
Proposition~\ref{prop:ambient-polynomial}, restricted to alternating maps, is ambient of degree $k+1$ as a map from
$\cL(E,D)\times\cL(F,F')$ to $\cL(\Alt^k(D;F),\Alt^k(E;F'))$. As bounded polynomial laws, all four have norm at most $1$.
\end{enumerate}
\end{theorem}

The proof is in Appendix~\ref{app:deferred}.

\begin{corollary}[divided-power bundles]\label{dp:cor:bundles}
Theorem~\ref{thm:bundle-lifting} holds with $C^n$ replaced by DP on open sets, for manifolds and vector bundles whose
transition maps are DP on open sets. The alternating bifunctor $(E,F)\mapsto\Alt^k(E;F)$ and the functors
$E\mapsto\Lambda^k_\pi E$ and $E\mapsto\widehat\Lambda^k_\pi E$ satisfy its hypothesis. So in every characteristic and for
all normed fibers, the alternating and the projective exterior power bundles of a DP vector bundle are DP vector
bundles, with DP induced morphisms.
\end{corollary}

\begin{proof}
\emph{Step 1 (the calculus).} By Theorem~\ref{dp:thm:calculus}(2), maps that are DP on open sets compose; by~(3) and~(1),
bounded multilinear maps, in particular bounded linear maps and evaluation of operators, and constants are DP on their
whole domains. A map into a product, with the maximum norm, is DP at a point when its components are, since a pair of
bounded polynomial laws of degree $n$ is a bounded polynomial law of degree $n$. Being DP at a point depends only on the
map near that point.

\emph{Step 2 (the theorem).} The proof of Theorem~\ref{thm:bundle-lifting} uses of the class $C^n$ only these
properties: Step~1 there composes the transitions with the maps of the functor on hom spaces and evaluates operators,
Step~2 composes local operator families, Step~3 uses constants and Step~4 restricts to open subsets of $P$.

\emph{Step 3 (the functors).} By Theorem~\ref{dp:thm:calculus}(4) and~(1), the maps of the three functors on hom spaces
are bounded polynomial laws and DP on their whole domains.
\end{proof}

The two notions differ: if $\operatorname{char}K=p>0$ and $k\ge p$, then $A^k_{E,D;F}$ is DP on $\cL(E,D)$ by
Theorem~\ref{dp:thm:calculus}(4) and~(1), and for the spaces of Theorem~\ref{thm:operator-obstructions} it is analytic at
no point.

Bounded polynomial laws need not be ambient. Let $p=\operatorname{char}K>0$ and $c_0=c_0(\N,K)$, and identify $K^m$
with the span of the first $m$ unit vectors, $m\ge1$. Let $\Gamma^p(K^m)\subseteq(K^m)^{\otimes p}$ be the space of
symmetric tensors, put $Q_m=(K^m)^{\otimes p}/\Gamma^p(K^m)$, and let $\bar z\in Q_m$ be the image of
$z\in(K^m)^{\otimes p}$. Let $I_m$ be the set of the $p^{pm}$ tuples $l=(y_l^1,\dots,y_l^p)$ of vectors in $K^m$ with
entries in the prime field $\F_p$, and put $y_l=y_l^1\otimes\dots\otimes y_l^p$. Let $N_m\subseteq\ell^1(I_m,K)$ consist
of the vectors $(\phi(\bar y_l))_{l\in I_m}$, $\phi$ a linear form on $Q_m$, let $G_m=\ell^1(I_m,K)/N_m$ with the
quotient norm, and let $\varepsilon_l\in G_m$ be the image of the unit vector at $l$. Let $G$ be the $\ell^1$-sum of the
spaces $G_m$, $m\ge1$: its elements are the sequences whose $m$-th term lies in $G_m$ and whose norm, the sum of the
norms of the terms, is finite; we read $G_m$ as the $m$-th summand of $G$. Let $H=K\oplus c_0\oplus G$, with the
maximum norm; it is a Banach space if $K$ is complete. Let $V$ be the span of the tensors $u^{\otimes p}\otimes g$, $u\in c_0$, $g\in G$, in $c_0^{\otimes p}\otimes G$,
with the norm
\[
 \nu(t)=\inf\Bigl\{\sum_i\abs{c_i}\norm{u_i}^p\norm{g_i}:t=\sum_ic_i\,u_i^{\otimes p}\otimes g_i\Bigr\},
\]
which dominates the projective norm $\pi$ of $c_0^{\otimes p}\otimes G$. Both are norms, since the coordinate forms
of $c_0$, and the linear forms on the spaces $G_m$ composed with the projections of $G$ onto $G_m$, are bounded and
separate points.

\begin{proposition}[laws that are not ambient]\label{dp:prop:nonambient}
Assume that either (a) $K$ is discretely valued, that is, $\abs{K^\times}=\abs\vartheta^{\Z}$ for some $\vartheta\in K$
with $0<\abs\vartheta<1$, or (b) $p=2$. Put $t_m=\sum_{l\in I_m}y_l\otimes\varepsilon_l$.
\begin{enumerate}
\item $t_m\in V$ and $\pi(t_m)\le p^{pm}$; under~(a),
$\nu(t_m)\ge mp^{pm}/\bigl(2^pM_{\abs\vartheta}(1+\lfloor\log_pm\rfloor)\bigr)$ with
$M_{\abs\vartheta}=\max_{j\ge0}(j+1)\abs\vartheta^{\,j}$, and under~(b), $\nu(t_m)\ge m4^{m-1}$.
\item For $n\ge p+1$, $R_n(\lambda,x,g)=\lambda^{n-p-1}\,x^{\otimes p}\otimes g$, with $\lambda^0=1$, defines a bounded polynomial law
$R_n\colon H\to V$ of degree $n$ and norm at most $1$ that is not ambient of degree $n$, also as a map into the
completion of $V$.
\end{enumerate}
\end{proposition}

The proof is in Appendix~\ref{app:deferred}. Under~(a) it rests on Lemma~\ref{dp:lem:projections}, by which every
projection of $(\kappa^m)^{\otimes p}$ onto the symmetric tensors, $\kappa$ the residue field of $K$, takes, on average
over the tensors $y_l$, which have entries in $\F_p$, values whose $m\times m^{p-1}$ matrices of entries have large rank.
Under~(b) it rests on Proposition~\ref{ocl:prop:sections}, by which every projection $\Psi$ of $(K^m)^{\otimes2}$ onto
the symmetric tensors has some $\Psi(y_l)$ whose determinant has absolute value at least $1$, and on
Lemma~\ref{ocl:lem:pfaffian-bound}(2), which turns this into a lower bound for its projective norm. Condition~(a) holds
under~\textup{(L)}, with the element $\vartheta$ of Section~\ref{obs:ssec:field}, and then $M_{\abs\vartheta}\le M_\varrho$.
In degrees below $p$ every bounded polynomial
law is ambient: a bounded polynomial law $a$ of degree $n<p$ is the diagonal of
$(h_1,\dots,h_n)\mapsto a^{[(1,\dots,1)]}(h_1,\dots,h_n)/n!$. This map is $n$\nobreakdash-linear, as comparing
coefficients in $a(s_1h_1+s_1'h_1'+s_2h_2+\dots+s_nh_n)$ after the substitutions $s_1'=s_1$, $s_1'=0$ and $s_1=0$, and
scaling $s_1$, shows, and likewise in every slot; it has diagonal $a$, since $a((s_1+\dots+s_n)h)=(s_1+\dots+s_n)^na(h)$
gives $a^{[(1,\dots,1)]}(h,\dots,h)=n!\,a(h)$; and it is bounded by Theorem~\ref{dp:thm:calculus}(1). A bounded
polynomial law $a\colon P\to W$ of degree $n$ whose composite $ja$ with a linear isometry $j$ of $W$ into a normed space
is analytic at one point $h_0$ is ambient of degree $n$: if $ja(h_0+h)=\sum_ib_i(h,\dots,h)$ near $h=0$,
Lemma~\ref{lem:unique} on the line $s\mapsto h_0+sh$ gives $b_n(h,\dots,h)=ja(h)$, the coefficient of $s^n$ in the
polynomial $ja(h_0+sh)$. So $R_n$ is DP on $H$ by Theorem~\ref{dp:thm:calculus}(1), while by
Proposition~\ref{dp:prop:nonambient}(2) $jR_n$ is analytic at no point for every linear isometry $j$ of $V$, or of its
completion, into a normed space.

In degree $p$ itself, with $p=\operatorname{char}K>0$ as above, bounded polynomial laws need not be ambient either,
over every such $K$. For a normed space $P$ let $\Gamma^p(P)\subseteq P^{\otimes p}$ be the span of the powers
$h^{\otimes p}$, $h\in P$; for $P=K^m$ it is the space of symmetric tensors, by the paragraph at the beginning of
Section~\ref{dgr:ssec:polarizable}, as $K$ is infinite. Give $\Gamma^p(P)$ the seminorm
\begin{equation}\label{dp:eq:gamma-norm}
 \norm t_\Gamma=\inf\Bigl\{\sum_i\abs{c_i}\norm{h_i}^p:t=\sum_ic_i\,h_i^{\otimes p}\Bigr\},
\end{equation}
which dominates the projective seminorm $\pi$ of $P^{\otimes p}$; both are norms if the bounded linear forms on $P$
separate points. For integers $m\ge1$ and $d\ge0$ let $\N^m_d$ be the set
of $\alpha\in\N^m$ with $\abs\alpha=d$, and let $e_1,\dots,e_m$ be the unit vectors of $K^m$. For a symmetric tensor
$T\in(\F_p^m)^{\otimes p}$ and $\alpha\in\N^m_p$ let $T_\alpha$ be its coefficient at
$e_{i_1}\otimes\dots\otimes e_{i_p}$ for any $(i_1,\dots,i_p)$ in which each $i$ occurs $\alpha_i$ times. For
$\gamma\in\N^m_{p-1}$ let the \emph{slice} $\partial^\gamma T\in\F_p^m$ be the vector whose $i$-th entry is
$T_{\gamma+e_i}$, where $\gamma+e_i\in\N^m_p$ is $\gamma$ with its $i$-th entry raised by $1$. For a real number
$\tau\ge1$ let $\mathcal E_{T,\tau}$ be $K^m$ with the norm
\begin{equation}\label{dp:eq:weighted-norm}
 \norm h_{T,\tau}=\inf\Bigl\{\sum_{i=1}^m\abs{z_i}+\tau\sum_{\gamma\in\N^m_{p-1}}\abs{z_\gamma}:
 h=\sum_{i=1}^mz_ie_i+\sum_{\gamma\in\N^m_{p-1}}z_\gamma\,\partial^\gamma T\Bigr\},
\end{equation}
the infimum over the scalars $z_i,z_\gamma\in K$. It is a norm, and the coordinate forms of $\mathcal E_{T,\tau}$ are bounded:
as the entries of $T$ lie in $\F_p$ and $\tau\ge1$, no entry of $h$ has absolute value above $\norm h_{T,\tau}$. The
tensor $T$ lies in $\Gamma^p(\mathcal E_{T,\tau})$.

\begin{proposition}[laws that are not ambient, over every $K$]\label{dp:prop:degree-p}
\leavevmode
\begin{enumerate}
\item For every $m\ge1$ there are a nonzero symmetric tensor $T\in(\F_p^m)^{\otimes p}$ and a real number $\tau\ge1$
such that, in $\Gamma^p(\mathcal E_{T,\tau})$,
\begin{equation}\label{dp:eq:degree-p-ratio}
 \norm T_\Gamma\ge\frac{(p-1)!}{3p^{3/2}(2p-1)!}\Bigl(\frac m{1+\log_pm}\Bigr)^{1/2}\pi(T).
\end{equation}
\item Choose such $T^{(m)}$ and $\tau^{(m)}$ for every $m\ge1$, and let $\mathcal E$ be the $\ell^1$-sum of the spaces
$\mathcal E_m=\mathcal E_{T^{(m)},\tau^{(m)}}$, $m\ge1$, defined as for $G$; it is a Banach space if $K$ is complete. For $n\ge p$, the
map $(\lambda,h)\mapsto\lambda^{n-p}h^{\otimes p}$, with $\lambda^0=1$, from $K\oplus\mathcal E$, with the maximum norm, to
$\Gamma^p(\mathcal E)$, with the norm $\norm\cdot_\Gamma$ of~\eqref{dp:eq:gamma-norm}, is a bounded polynomial law of degree
$n$ and norm at most $1$ that is not ambient of degree $n$, also as a map into the completion of $\Gamma^p(\mathcal E)$.
\end{enumerate}
\end{proposition}

The proof is in Appendix~\ref{app:deferred}. A representation of $T$ in $\Gamma^p(\mathcal E_{T,\tau})$ gives a symmetric tensor
over $\F_p$ in $m+\binom{m+p-2}{p-1}$ symbols, which stand for the vectors $e_i$ and $\partial^\gamma T$, that maps to
$T$ and whose nonzero coefficients have total weight at most the cost $\sum_i\abs{c_i}\norm{h_i}^p$ of the
representation; a coefficient at two or more slices weighs at least $\tau^2$, the others at least $1$. The coefficients
at exactly one slice act on $T$ through a linear map $\mathsf L^\sharp$ of Lemma~\ref{dp:lem:linear}, an identity for
symmetric tensors in characteristic $p$ by which $\id-\mathsf L^\sharp$ has rank at least $(p-1)!\,m^p/(2p-1)!$. So
for most $T$ every such tensor has many nonzero coefficients, or many at two or more slices. By
Theorem~\ref{dp:thm:calculus}(1) the law of~(2) is DP on $K\oplus\mathcal E$, while, as for $R_n$, its composite with a linear
isometry of $\Gamma^p(\mathcal E)$, or of its completion, into a normed space is analytic at no point. For $n\ge p+1$ this gives
laws with the properties of Proposition~\ref{dp:prop:nonambient}(2) over every $K$, on another space.

For $h_1,\dots,h_p\in P$ the tensor $\sum_{\sigma\in S_p}h_{\sigma(1)}\otimes\dots\otimes h_{\sigma(p)}$ lies in
$\Gamma^p(P)$, by~\eqref{can:eq:polarization}; the elements of the span of these tensors are \emph{hollow}. For
$t\in\Gamma^p(P)$ let
\[
 \norm t_{\mathrm F}=\inf\{\norm{t'}_\Gamma:t'\in\Gamma^p(P),\ t-t'\text{ hollow}\},
\]
the quotient seminorm of the seminorm $\norm\cdot_\Gamma$ of~\eqref{dp:eq:gamma-norm} modulo the hollow tensors. For
$C\ge0$, $P$ satisfies \emph{$(\mathrm N_P)$ with constant $C$} if $\norm t_\Gamma\le C\pi(t)$ for every hollow $t$,
and \emph{$(\mathrm F_P)$ with constant $C$} if $\norm t_{\mathrm F}\le C\pi(t)$ for every $t\in\Gamma^p(P)$. A
\emph{Frobenius law} is a map $a\colon P\to W$ with $a(h+h')=a(h)+a(h')$ and $a(\lambda h)=\lambda^pa(h)$ for
$h,h'\in P$ and $\lambda\in K$, and with $\norm a=\sup_{h\ne0}\norm{a(h)}/\norm h^p$ finite; it is a bounded
polynomial law of degree $p$, as $a(s_1h_1+\dots+s_rh_r)=s_1^pa(h_1)+\dots+s_r^pa(h_r)$.

\begin{theorem}[laws of degree $p$]\label{dp:thm:frobenius}
Let $p=\operatorname{char}K>0$.
\begin{enumerate}
\item $P$ satisfies $(\mathrm F_P)$ with constant $p$, and with constant $\dim P$ if $\dim P<p$.
\item Every Frobenius law $a\colon P\to W$ is ambient of degree $p$, with a $p$\nobreakdash-linear map $\beta$ of norm
at most $p\norm a$ in Definition~\ref{dp:def:laws}(2).
\item If $P$ satisfies $(\mathrm N_P)$ with constant $C$, then $\norm t_\Gamma\le(p+(p+1)C)\pi(t)$ for every
$t\in\Gamma^p(P)$, and every bounded polynomial law $a\colon P\to W$ of degree $p$ is ambient of degree $p$, with a
$p$\nobreakdash-linear map $\beta$ of norm at most $(p+(p+1)C)\norm a$. Conversely, if every bounded polynomial law of
degree $p$ on $P$, with values in any normed space, is ambient of degree $p$, then $P$ satisfies $(\mathrm N_P)$ with
some constant.
\end{enumerate}
\end{theorem}

The proof is in Appendix~\ref{app:deferred}. It rests on Lemma~\ref{dp:lem:bases}: if
$t\in\Gamma^p(P)$ is a sum of elementary tensors with nonzero factors, some of these factors form a basis of the span
of all of them for which $t$ is congruent, modulo the hollow tensors, to a combination of the $p$-th powers of the basis
vectors in which the power of each basis vector $h$, taken as a factor of a term $h_1\otimes\dots\otimes h_p$, has a
coefficient of absolute value at most $\norm{h_1}\cdots\norm{h_p}/\norm h^p$. By~(3), the space $K\oplus\mathcal E$ of
Proposition~\ref{dp:prop:degree-p}(2) does not satisfy $(\mathrm N_{K\oplus\mathcal E})$ with any constant, as the law there
with $n=p$ is not ambient of degree $p$.

For $n\ge1$ let $\Gamma^n(P)\subseteq P^{\otimes n}$ be the span of the powers $h^{\otimes n}$, $h\in P$, with the
seminorm $\norm\cdot_\Gamma$ of~\eqref{dp:eq:gamma-norm} with $n$ in place of $p$, and let $\pi$ be the projective
seminorm of $P^{\otimes n}$; for $n=p$ these are the space and the seminorms above. Let $\mathfrak a_n(P)\in[0,\infty]$
be the infimum of the $C\ge0$ with $\norm t_\Gamma\le C\pi(t)$ for all $t\in\Gamma^n(P)$; if it is finite, it is the
least such $C$. For $m\ge1$ let $[m]=\{1,\dots,m\}$, and let $c_0([m],K)$ be $K^m$ with the maximum norm, read as the span of the first $m$ unit
vectors of $c_0(\N,K)$. Put $\mathfrak c_K=\abs\vartheta$ if $K$ is discretely valued, with $\vartheta$ as in
Proposition~\ref{dp:prop:nonambient}(a), and $\mathfrak c_K=1$ otherwise.

\begin{proposition}[ambience constants]\label{dp:prop:ultrametric}
Let $n\ge1$.
\begin{enumerate}
\item If $\mathfrak a_n(P)<\infty$, then every bounded polynomial law $a\colon P\to W$ of degree $n$ is ambient of degree
$n$, with an $n$\nobreakdash-linear map $\beta$ of norm at most $\mathfrak a_n(P)\norm a$ in
Definition~\ref{dp:def:laws}(2). Conversely, if every bounded polynomial law of degree $n$ on $P$, with values in any
normed space, is ambient of degree $n$, then $\mathfrak a_n(P)<\infty$.
\item $\mathfrak a_n(P)\le(n+1)\,\mathfrak a_{n+1}(K\oplus P)$, with the maximum norm on $K\oplus P$.
\item If $P$ is nonarchimedean, then
$\mathfrak a_n(P)\le\mathfrak c_K^{-n}\sup_{m\ge1}\mathfrak a_n(c_0([m],K))\le\mathfrak c_K^{-n}\,\mathfrak a_n(c_0(\N,K))$.
\item If $P$ is nonarchimedean, then $\mathfrak a_2(P)\le\frac32\,\mathfrak c_K^{-2}$.
\end{enumerate}
\end{proposition}

The proof is in Appendix~\ref{app:deferred}. By~(1) and~(3), if every bounded polynomial law of degree $n$ on
$c_0(\N,K)$ is ambient of degree $n$, then so is every bounded polynomial law of degree $n$ on every nonarchimedean
normed space; by~(1) and~(4), every bounded polynomial law of degree $2$ on a nonarchimedean normed space is ambient of
degree $2$.
%% ---- end sec-10-geometry.tex
%% ---- begin sec-10b-families.tex
\section{Familywise analyticity and bundles}\label{sec:families}

Throughout, $P$, $E$, $D$ are normed spaces and $k\ge1$; $P$ is the model space of a base, since
Theorem~\ref{thm:bundle-lifting} needs only operator families on open subsets of $P$.

\subsection{Families of operators}\label{fam2:ssec:framework}

\begin{definition}\label{fam2:def:good}
The triple $(P;E,D)$ is \emph{good in degree $k$} if for every open $U\subseteq P$ and every analytic
$u\colon U\to\cL(E,\widehat D)$ the map $x\mapsto\Lambda^k_\pi u(x)\in\cL(\Lambda^k_\pi E,\hLambda^k_\pi D)$ is analytic on
$U$. We write $(P;E)$ for $(P;E,E)$. For a Banach space $H$ and an open $U\subseteq P$, an \emph{invertible family} on $U$
is an analytic $g\colon U\to\cL(H)$ whose values are invertible and for which $x\mapsto g(x)^{-1}$ is analytic.
\end{definition}

Here $\Lambda^k_\pi u(x)$ maps $\Lambda^k_\pi E$ into $\Lambda^k_\pi\widehat D\subseteq\hLambda^k_\pi D$ (paragraph before
Theorem~\ref{ob:thm:mechanisms}). By Theorem~\ref{thm:duality}(1) with $X=\Lambda^k_\pi E$, $Y=\hLambda^k_\pi D$, and the
completed form of Proposition~\ref{prop:exterior}, $x\mapsto\Lambda^k_\pi u(x)$ is analytic at a point, or on $U$, if and
only if $x\mapsto A^k_{E,\widehat D;F}(u(x))$ is analytic there for every Banach space $F$. For a linear family
$\theta\colon P\to\cL(E,\widehat D)$ (defined before Theorem~\ref{cmb:thm:decomposition}), $h\mapsto\Lambda^k_\pi\theta(h)$ is analytic at some
point if and only if $\theta$ is liftable in degree $k$, and then it is a continuous polynomial
(Proposition~\ref{prop:homogeneous-lift}, as in the proof of Lemma~\ref{fib:lem:test}).

\begin{proposition}[reductions]\label{fam2:prop:reductions}
\begin{enumerate}
\item Let $H$ be a Banach space. If $x\mapsto\Lambda^k_\pi g(x)$ is analytic for every invertible family $g$ on every open
subset of $P$, then $(P;H)$ is good in degree $k$. More precisely, for an analytic $u\colon U\to\cL(H)$ and $x_0\in U$ there
are distinct $\lambda_0,\dots,\lambda_k\in K$ and a neighborhood $U_0\subseteq U$ of $x_0$ such that every $1+\lambda_iu$,
with $1=\id_H$, is an invertible family on $U_0$ and $\Lambda^k_\pi u$ is a fixed linear combination of the maps
$\Lambda^k_\pi(1+\lambda_iu)$ on $U_0$.
\item If $(P;E,D)$ is good in degree $k$, so is every triple of bounded linear retracts of $P$, $E$ and $D$, respectively.
The triple $(P;E\oplus E',D)$ is good in every degree $j\le k$ if and only if $(P;E,D)$ and $(P;E',D)$ are; the same holds
in the target slot.
\item Let $E$ and $D$ be Banach spaces and $H=D\oplus E$ with the maximum norm. Then $x\mapsto\Lambda^j_\pi g(x)$ is analytic
for every invertible family $g$ into $\cL(H)$ over every open subset of $P$ and every $j\le k$ if and only if
the four triples $(P;E,E)$, $(P;E,D)$, $(P;D,E)$, $(P;D,D)$ are good in every degree $j\le k$. If $\cC$ is a class of Banach spaces closed
under finite direct sums, then, in the single degree $k$, $x\mapsto\Lambda^k_\pi g(x)$ is analytic for every $G\in\cC$ and
every invertible family $g$ into $\cL(G)$ over an open subset of $P$ if and only if $(P;E,D)$ is good in degree $k$ for
all $E,D\in\cC$.
\end{enumerate}
\end{proposition}

The proof is in Appendix~\ref{app:families}.

\begin{corollary}[what bundles need]\label{fam2:cor:bundles}
Let $\cC$ be a class of Banach spaces closed under finite direct sums. The following are equivalent.
\begin{enumerate}
\item For every analytic manifold $M$ modeled on $P$, every analytic vector bundle over $M$ with typical fiber $E\in\cC$ and
every Banach space $F$, the atlas induced on the bundle with typical fiber $\Alt^k(E;F)$ is analytic.
\item The triple $(P;E,D)$ is good in degree $k$ for all $E,D\in\cC$.
\end{enumerate}
If they hold, Theorem~\ref{thm:bundle-lifting} applies to fiberwise alternating maps on analytic bundles with typical
fibers in $\cC$ and analytic bundles with Banach typical fibers, and to their operator-valued analytic morphisms, over every
analytic manifold modeled on $P$.
\end{corollary}

The proof is in Appendix~\ref{app:families}.

\subsection{Homogeneous families}\label{fam2:ssec:reduction}

\begin{definition}\label{fam2:def:lifting}
For $n\ge0$ and $\beta\in\Mult^n(P;\cL(E,\widehat D))$ put $u_\beta(h)=\beta(h,\dots,h)$; for $n=0$, $\beta$ is a constant
operator. For $j\ge1$, a \emph{lift of $u_\beta$ in degree $j$} is a bounded $jn$-linear map
$q\colon P^{jn}\to\cL(\Lambda^j_\pi E,\hLambda^j_\pi D)$ with $q(h,\dots,h)=\Lambda^j_\pi u_\beta(h)$ for $h\in P$, and
$c_j(\beta)\in[0,\infty]$ is the infimum of their norms.
\end{definition}

For $n\ge1$, $c_j(\beta)<\infty$ if and only if $h\mapsto\Lambda^j_\pi u_\beta(h)$ is analytic at some point of $P$
(Lemma~\ref{fmo:lem:homogeneous}). For $j\ge1$ consider the condition
\begin{equation}\label{fam2:eq:G}
\begin{aligned}
 (\mathrm G_j)\colon\ &\text{for every sequence $(\beta_n)_{n\ge1}$ with $\beta_n\in\Mult^n(P;\cL(E,\widehat D))$ and
 $\norm{\beta_n}\le1$,}\\
 &\text{every $c_j(\beta_n)$ is finite and $\limsup_{n\to\infty}c_j(\beta_n)^{1/n}<\infty$.}
\end{aligned}
\end{equation}

\begin{theorem}[reduction to homogeneous families]\label{fam2:thm:reduction}
\begin{enumerate}
\item If $(P;E,D)$ is good in degree $j$, then $(\mathrm G_j)$ holds.
\item If $(\mathrm G_j)$ holds for every $j\le k$, then $(P;E,D)$ is good in degree $k$.
\item Condition $(\mathrm G_j)$ holds whenever $j!\ne0$ in $K$. If $\operatorname{char}K=p>0$, then $(P;E,D)$ is good in
every degree $j\le k$ if and only if $(\mathrm G_{p^i})$ holds for every $i\ge1$ with $p^i\le k$.
\item Let $Y_n$ be the completed projective tensor product of $n$ copies of $P$ and one copy of $E$, and let
$\beta^{\mathrm u}_n\in\Mult^n(P;\cL(E,Y_n))$ be given by
$\beta^{\mathrm u}_n(h_1,\dots,h_n)y=h_1\otimes\dots\otimes h_n\otimes y$. For every normed $D$, $n\ge1$ and
$\beta\in\Mult^n(P;\cL(E,\widehat D))$, $c_j(\beta)\le c_j(\beta^{\mathrm u}_n)\norm\beta^j$. Consequently $(P;E,D)$ is good
in every degree $j\le k$ for every normed $D$ if and only if, for every $j\le k$, every $c_j(\beta^{\mathrm u}_n)$ is finite
and $\limsup_nc_j(\beta^{\mathrm u}_n)^{1/n}<\infty$; if $\operatorname{char}K=p>0$, the degrees $j=p^i\le k$ suffice.
Equivalently, $(P;E,D_{\mathrm u})$ is good in every degree $j\le k$ for the $\ell^1$-sum $D_{\mathrm u}$ of the spaces
$Y_n$.
\end{enumerate}
\end{theorem}

\begin{corollary}[degrees of triples]\label{fam2:cor:degrees}
Let $\operatorname{char}K=p>0$, and let $\fD(P;E,D)$ be the set of $k\ge1$ in which $(P;E,D)$ is good.
\begin{enumerate}
\item $\{1,\dots,p-1\}\subseteq\fD(P;E,D)$, and $a+b\in\fD(P;E,D)$ whenever $a,b\in\fD(P;E,D)$ and the addition of $a$ and
$b$ in base $p$ has no carry. If $\fD(P;E,D)\ne\{1,2,\dots\}$, its least missing element is a power of $p$.
\item If $E\cong E\oplus K$ and $D\cong D\oplus K$ by bounded linear isomorphisms, then $\fD(P;E,D)$ is
$\{1,\dots,p^i-1\}$ for some $i\ge1$ or $\{1,2,\dots\}$.
\item $\fD(P;E,D)\supseteq\fD_B(E,\widehat D)$, with equality for $P=\cL(E,\widehat D)$.
\item The set of $k$ for which $(P;E,D)$ is good in degree $k$ for all normed $E$ and $D$ is $\{1,\dots,p^i-1\}$ for some
$i\ge1$, or $\{1,2,\dots\}$.
\end{enumerate}
\end{corollary}

The proofs are in Appendix~\ref{app:families}.

\subsection{The familywise class}\label{fam2:ssec:class}

The tensor algebra $\mathcal T(P)$ is the Banach space of sequences $(w_n)_{n\ge0}$, $w_n\in\widehat\otimes^n_\pi P$,
$\widehat\otimes^0_\pi P=\Kh$, with norm $\sum_n\norm{w_n}<\infty$. The map $\operatorname{ver}_n\in\Mult^n(P;\mathcal T(P))$
places $h_1\otimes\dots\otimes h_n$ in slot $n$, with $\operatorname{ver}_0$ the unit of slot $0$, and
$\operatorname{ver}(h)=\sum_n\operatorname{ver}_n(h,\dots,h)=(h^{\otimes n})_{n\ge0}$ is analytic for $\norm h<1$
(Lemma~\ref{fmo:lem:series}). Every analytic germ factors through $\operatorname{ver}$: if $u$ is analytic at $x_0$ with
values in $\cL(E,\widehat D)$, then $u(x_0+\lambda h)=\theta(\operatorname{ver}(h))$ for $\norm h<1$, for some
$\lambda\in K^\times$ and a linear family $\theta\colon\mathcal T(P)\to\cL(E,\widehat D)$ (Lemma~\ref{fmo:lem:germs}).

For a Banach space $Y$ and $0<R<1$, $\Mult^\omega_R(P;Y)$ is the Banach space of sequences $(b_N)_{N\ge0}$,
$b_N\in\Mult^N(P;Y)$, with norm $\sup_NR^N\norm{b_N}<\infty$; $\Mult^\omega_{R,0}(P;Y)$ is its closed subspace where every
$b_N$ has zero diagonal, so that $b_0=0$; for $R'\le R$, $\iota_{R'R}$ is the contractive inclusion into
$\Mult^\omega_{R'}(P;Y)$. With $I_1,\dots,I_k$ the consecutive blocks of $\{1,\dots,N\}$ of sizes $n_1,\dots,n_k$ and $h_I$
the tuple of the $h_i$ with $i\in I$, the map $\operatorname{ver}^\sharp\colon\Mult^k(\mathcal T(P);Y)\to\Mult^\omega_R(P;Y)$,
\[
 (\operatorname{ver}^\sharp\Phi)_N(h_1,\dots,h_N)=\sum_{n_1+\dots+n_k=N}
 \Phi\bigl(\operatorname{ver}_{n_1}(h_{I_1}),\dots,\operatorname{ver}_{n_k}(h_{I_k})\bigr),
\]
is bounded and maps $\Mult^k_0(\mathcal T(P);Y)$ into $\Mult^\omega_{R,0}(P;Y)$ (Lemma~\ref{fmo:lem:veronese}). For a linear
family $\theta\colon\mathcal T(P)\to\cL(E,\widehat D)$, with $\ev_\theta$ and $\ob_k(\theta)$ as in
Section~\ref{cmb:ssec:families}, put
\begin{equation}\label{fam2:eq:class}
\begin{split}
 \ob^\omega_k(\theta;R)&=\bigl(\operatorname{ver}^\sharp\ev_\theta|_{\cN^k_E}\bigr)_*[\cS^k_E]
 =(\operatorname{ver}^\sharp)_*\ob_k(\theta)\\
 &\in\Ext\bigl(\hLambda^k_\pi E,\Mult^\omega_{R,0}(P;\hLambda^k_\pi D)\bigr),
\end{split}
\end{equation}
the second equality by Corollary~\ref{ocl:cor:functorial}(b), and
$\ob_{k,N}(\theta)=((\operatorname{pr}_N\operatorname{ver}^\sharp\ev_\theta)|_{\cN^k_E})_*[\cS^k_E]
\in\Ext(\hLambda^k_\pi E,\allowbreak\Mult^N_0(P;\hLambda^k_\pi D))$, with $\operatorname{pr}_N$ the $N$th term. The class
$\ob^\omega_k(\theta)$ \emph{vanishes after shrinking} if $(\iota_{R'R})_*\ob^\omega_k(\theta;R)=0$ for some
$R'\in(0,R]$; this does not depend on $R$.

\begin{theorem}[the familywise class]\label{fam2:thm:class}
\begin{enumerate}
\item For a linear family $\theta\colon\mathcal T(P)\to\cL(E,\widehat D)$ the following are equivalent:
$x\mapsto\Lambda^k_\pi\theta(\operatorname{ver}(x))$ is analytic at $0$; $\ob^\omega_k(\theta)$ vanishes after shrinking;
there are $C,R'>0$ such that for every $N$ the map $(\operatorname{pr}_N\operatorname{ver}^\sharp\ev_\theta)|_{\cN^k_E}$
extends to a bounded linear map $\widehat\otimes^k_\pi E\to\Mult^N_0(P;\hLambda^k_\pi D)$ of norm at most $CR'^{-N}$.
\item $(P;E,D)$ is good in degree $k$ if and only if $\ob^\omega_k(\theta)$ vanishes after shrinking for every linear family
$\theta\colon\mathcal T(P)\to\cL(E,\widehat D)$.
\item Let $\theta_{\mathrm u}\colon\mathcal T(P)\to\cL(E,\mathcal T(P)\widehat\otimes_\pi E)$ be given by
$\theta_{\mathrm u}(w)y=w\otimes y$. For $\theta$ as in~(1), let $\tau_\theta\colon\mathcal T(P)\widehat\otimes_\pi E\to\widehat D$
be the linearization of $(w,y)\mapsto\theta(w)y$, and let $(\hLambda^k_\pi\tau_\theta)_\flat$ apply $\hLambda^k_\pi\tau_\theta$
to every term. Then $\ob^\omega_k(\theta;R)=((\hLambda^k_\pi\tau_\theta)_\flat)_*\ob^\omega_k(\theta_{\mathrm u};R)$. Hence
$(P;E,D)$ is good in degree $k$ for every normed $D$ if and only if $\ob^\omega_k(\theta_{\mathrm u})$ vanishes after
shrinking, and if and only if $x\mapsto\Lambda^k_\pi\mathcal U(x)$ is analytic at $0$ for the universal germ
$\mathcal U(x)y=\operatorname{ver}(x)\otimes y$.
\end{enumerate}
\end{theorem}

The proof is in Appendix~\ref{app:families}.

\subsection{Base models good for all fibers}\label{fam2:ssec:bases}

\begin{definition}\label{fam2:def:reflection}
The space $P$ is \emph{universally good in degree $k$}, written $(\mathrm U^k_P)$, if $(P;E,D)$ is good in degree $k$ for
all normed $E$ and $D$. It has \emph{fixed-degree reflection in degree $k$}, written $(\mathrm{FR}^k_P)$, if for all Banach
spaces $G$ and $Z$, every closed subspace $W\subseteq Z$, every bounded $k$-linear $\Phi\colon G^k\to Z$ with
$\Phi(y,\dots,y)\in W$ for $y\in G$, every open $U\subseteq P$ and every analytic $\gamma\colon U\to G$, the map
$x\mapsto\Phi(\gamma(x),\dots,\gamma(x))$ is analytic as a $W$-valued map on $U$.
\end{definition}

Recall $g_n(P)$ from Definition~\ref{ob:def:sequences}(3).

\begin{theorem}[polarizable base models]\label{fam2:thm:polarizable}
\begin{enumerate}
\item The following are equivalent: $(\mathrm{FR}^k_P)$; $(\mathrm{FR}^k_P)$ restricted to affine maps
$\gamma(x)=y_0+Tx$ on $U=P$, with $T$ bounded linear; $P$ is polarizable in every degree $n\le k$.
\item If $P$ is polarizable in every degree $n\le k$, then $(\mathrm U^k_P)$ holds. Quantitatively, let
$\bar g\ge\max(1,g_1(P),\dots,g_k(P))$, with $\bar g>g_n(P)$ when no projection of norm $g_n(P)$ exists. If
$u\colon U\to\cL(E,\widehat D)$ satisfies $u(x_0+h)=\sum_n\beta_n(h,\dots,h)$ for $\norm h<R$, with
$\norm{\beta_n}\le CR^{-n}$, then $\Lambda^k_\pi u$ has at $x_0$ an expansion of radius at least $R/\bar g$ whose $N$th
coefficient has norm at most $k!\,C^k\binom{N+k-1}{k-1}(\bar g/R)^N$. In particular, Theorem~\ref{thm:bundle-lifting}
applies to fiberwise alternating maps on analytic bundles with arbitrary normed typical fibers and analytic bundles with
Banach typical fibers over every analytic manifold modeled on $P$.
\item If $\operatorname{char}K=0$, or $\operatorname{char}K=p$ and $k<p$, then $(\mathrm U^k_P)$ holds for every $P$. If
$\operatorname{char}K=p$ and $P$ is polarizable in degree $p$, then $(\mathrm U^n_P)$ holds for every $n\ge1$. Every normed
space $P$ with a lower $r$-estimate for some $r<2$ (Section~\ref{spl:ssec:source-target}) satisfies $(\mathrm U^n_P)$ for
every $n\ge1$, over every $K$.
\item Let $K$ be complete and $P$ a Banach space. Then $P$ has universal analytic reflection
(Definition~\ref{fam:def:universal-reflection}) if and only if $\sup_ng_n(P)^{1/n}<\infty$, and $(\mathrm{FR}^k_P)$ holds
for every $k$ if and only if $g_n(P)<\infty$ for every $n$. If $K$ is nonarchimedean and $P=\ell^1(\N,K)$, then
$(\mathrm{FR}^k_P)$ and $(\mathrm U^k_P)$ hold for every $k$, but $P$ does not have universal analytic reflection.
\end{enumerate}
\end{theorem}

\begin{theorem}[base models under (L)]\label{fam2:thm:bases}
Assume~\textup{(L)}.
\begin{enumerate}
\item Let $P$ be a symmetric sequence space in which $c_{00}(\N,K)$ is dense, let $E_0$ be $c_{00}(\N,K)$ with the supremum
norm, and let $\theta_0\colon P\to\cL(E_0,P)$ be given by $\theta_0(a)y=ay$. The following are equivalent:
$(\mathrm U^k_P)$; $a\mapsto\Lambda^k_\pi\theta_0(a)$ is analytic at some point of $P$; $\cP_<$ is bounded on $P\otimes_\pi P$;
$P$ is polarizable in degree $p$; $(\mathrm U^n_P)$ holds for every $n\ge1$. They are equivalent to the conditions of
Theorem~\ref{cls:thm:symmetric}. In particular, for $P=\ell^r(\N,K)$ with $1\le r<\infty$, $(\mathrm U^k_P)$ holds if and
only if $r<2$; and for $P=\ell^r(\N,K)$ with $2\le r<\infty$, for $P=c_0(\N,K)$, and for every symmetric sequence space $P$
in which $c_{00}(\N,K)$ is dense and $\varphi_P(m)^2\le Cm$ for some $C$, $(\mathrm U^n_P)$ holds exactly for $n\le p-1$.
\item Let $P$ be normed, and suppose there are bounded linear maps $i\colon c_{00}(\N,K)\to P$, with the supremum norm on
$c_{00}(\N,K)$, and $w\colon P\to\linf(\N,\Kh)$ with $w\circ i$ the inclusion. Then for the linear family
$x\mapsto\mu_{w(x)}\in\cL(c_{00}(\N,K),\linf(\N,\Kh))$, with $\linf(\N,\Kh)$ regarded over $K$, the map
$x\mapsto\Lambda^k_\pi\mu_{w(x)}$ is analytic at no point, so $(\mathrm U^k_P)$ fails. Such maps exist for $c_0(\Sigma,K)$ and
$\linf(\Sigma,K)$ with $\Sigma$ infinite, and for every $P$ with an equivalent nonarchimedean norm and infinite-dimensional
completion, $\dim_{\Kh}\widehat P=\infty$. For $P$ with an equivalent nonarchimedean norm, $(\mathrm U^k_P)$ holds if and
only if $P$ is polarizable in degree $p$, and if and only if $\dim_{\Kh}\widehat P<\infty$.
\end{enumerate}
\end{theorem}

The proofs are in Appendix~\ref{app:families}.

\subsection{Fibers and mixed mechanisms}\label{fam2:ssec:mixed}

\begin{proposition}[fibers good over every base]\label{fam2:prop:fibers}
\begin{enumerate}
\item If $(E,\widehat D)$ is good for Banach values in degree $k$, with a lift $Q$ as in Lemma~\ref{fib:lem:test}, then
$(P;E,D)$ is good in degree $k$ for every normed $P$, and $c_k(\beta)\le\norm Q\norm\beta^k$ for every $n$ and
$\beta\in\Mult^n(P;\cL(E,\widehat D))$.
\item $(\cL(E,\widehat D);E,D)$ is good in degree $k$ if and only if $(E,\widehat D)$ is good for Banach values in degree
$k$. Hence $(P;E,D)$ is good in degree $k$ for every normed $P$ if and only if $(E,\widehat D)$ is good for Banach values
in degree $k$.
\end{enumerate}
\end{proposition}

The proof is in Appendix~\ref{app:families}.

\begin{definition}\label{fam2:def:gamma-split}
A bounded linear map $T\colon P\to P'$ is \emph{$\Gamma^k$-split} if there is a bounded $k$-linear
$\Pi\colon P^k\to\hGamma^k_\pi P'$ with $\Pi(a,\dots,a)=(Ta)^{\otimes k}$ for $a\in P$.
\end{definition}

By Lemma~\ref{ob:lem:parameter}, $P$ is polarizable in degree $k$ if and only if $\id_P$ is $\Gamma^k$-split, and $T$ is
$\Gamma^k$-split if $P$ or $P'$ is polarizable in degree $k$: compose $T^{\otimes k}$ with a projection onto
$\hGamma^k_\pi P'$, or a projection onto $\hGamma^k_\pi P$ with the extension of $T^{\otimes k}$.

\begin{theorem}[parameter-relative ordering]\label{fam2:thm:gamma-split}
Let $2<r<\infty$.
\begin{enumerate}
\item For every totally ordered set $\Sigma$ and every $k\ge1$, the inclusion $\ell^2(\Sigma,\Kh)\to\ell^r(\Sigma,\Kh)$,
over $K$, is $\Gamma^k$-split.
\item Let $K$ be nonarchimedean, and let $P$ have an $\ell^2$-dominated linearly independent family $(e_t)_{t\in\Sigma}$ with dense span $V$ and
$\sup_t\norm{e_t}<\infty$. Let $G$ be a normed space with an equivalent $r$-norm, $Z$ a Banach space, $W\subseteq Z$ a
closed subspace and $\Phi\colon G^k\to Z$ a bounded $k$-linear map with $\Phi(y,\dots,y)\in W$ for $y\in G$. For every
bounded linear $T\colon P\to G$ there is a bounded $k$-linear $Q\colon P^k\to W$ with $Q(a,\dots,a)=\Phi(Ta,\dots,Ta)$ for
$a\in P$; and for every open $U\subseteq P$ and every analytic $\gamma\colon U\to G$, the map
$x\mapsto\Phi(\gamma(x),\dots,\gamma(x))$ is analytic into $W$.
\item For $K$ and $P$ as in~(2), and for $D$ with an equivalent $r$-norm or an equivalent nonarchimedean norm, $(P;E,D)$ is good in every degree for every normed $E$.
\end{enumerate}
\end{theorem}

Under~\textup{(L)} and for $2<r<\infty$, (3) makes $(\ell^2(\N,K);\ell^r(\N,K))$ good in every degree, although neither
space is polarizable in degree $k$ (Corollary~\ref{cls:cor:symmetric-examples}).

\begin{theorem}[weighted staircase]\label{fam2:thm:staircase}
Assume~\textup{(L)}. Let $E$ be a sequence space with $\norm y_\infty\le\norm y_E$ for $y\in E$, and let $P$ be a normed
space containing $c_{00}(\N,K)$ as a linear subspace. For $G\in\{E,P\}$ let
\[
 \varphi^\pm_G(n)=\sup\bigl\{\norm y_G:y\in c_{00}(\N,K),\ y(t)\in\{0,1,-1\}\text{ for all }t,\
 \abs{\operatorname{supp}y}\le n\bigr\}.
\]
If $\liminf_{m\to\infty}\varphi^\pm_P(2m)^2\varphi^\pm_E(4m)^2/m=0$, then no $k$-linear
$Q\colon c_{00}(\N,K)^k\to\cL(\Lambda^k_\pi E,\hLambda^k_\pi E)$ that is bounded for the norm of $P$ satisfies
$Q(a,\dots,a)=\Lambda^k_\pi\mu_a$ for $a\in c_{00}(\N,K)$. Consequently, if $T\colon P\to\cL(E,\widehat E)$ is bounded
linear with $T(a)=\mu_a$ for $a\in c_{00}(\N,K)$, then $x\mapsto\Lambda^k_\pi T(x)$ is analytic at no point of $P$. For a
symmetric sequence space $G$, $\varphi^\pm_G=\varphi_G$.
\end{theorem}

\begin{corollary}[the table]\label{fam2:cor:table}
Assume~\textup{(L)}. Let $P$ be $\ell^{r_1}(\N,K)$ with $1\le r_1<\infty$, or $c_0(\N,K)$. In~(1) and~(2), let $E$ be
$\ell^{r_2}(\N,K)$ with $1\le r_2<\infty$, or $c_0(\N,K)$, or, if $K$ is complete, an infinite-dimensional Banach space with
an equivalent nonarchimedean norm.
\begin{enumerate}
\item $(P;E)$ is good in every degree if $r_1<2$; if $E=\ell^{r_2}(\N,K)$ with $r_2<2$; or if $r_1=2$ and $E$ is
$\ell^{r_2}(\N,K)$ with $r_2>2$, or $c_0(\N,K)$, or has an equivalent nonarchimedean norm.
\item There is a linear family $P\to\cL(E,\widehat E)$ along which $\Lambda^k_\pi$ is analytic at no point, so $(P;E)$ is
not good in degree $k$, if $P=c_0(\N,K)$ and $E$ is not $\ell^{r_2}(\N,K)$ with $r_2<2$; or if $2<r_1<\infty$ and $E$ is
$c_0(\N,K)$, or has an equivalent nonarchimedean norm, or is $\ell^{r_2}(\N,K)$ with $1/r_1+1/r_2<1/2$.
\item Let $E$ be a symmetric sequence space in which $c_{00}(\N,K)$ is dense. The multiplier family
$c_0(\N,K)\to\cL(E,\widehat E)$, $a\mapsto\mu_a$, is liftable in some degree $k\ge p$ if and only if $(E,E)$ is good for
Banach values in every degree, and if and only if $(c_0(\N,K);E)$ is good in every degree.
\end{enumerate}
\end{corollary}

The proofs are in Appendix~\ref{app:families}.

\subsection{Tangent bundles}\label{fam2:ssec:tangent}

\begin{definition}\label{fam2:def:tangent}
A normed space $P_0$ is \emph{tangent-good in degree $k$} if for every open $U\subseteq P_0$ and every analytic
$f\colon U\to P_0$ whose derivative $Df\colon U\to\cL(P_0,P_0)$ is analytic, the map
$x\mapsto\Lambda^k_\pi Df(x)\in\cL(\Lambda^k_\pi P_0,\hLambda^k_\pi P_0)$ is analytic on $U$.
\end{definition}

Let $K$ be complete and $P_0$ a Banach space, so that analytic maps between open subsets of $P_0$ have analytic
derivatives~\cite{Bou}. For a manifold $M$ modeled on $P_0$ with chart transitions $f_{ij}$, the transitions of
$\hLambda^k_\pi TM$ and $\Alt^k(TM;F)$ are $\Lambda^k_\pi Df_{ij}$ and $A^k_{P_0,P_0;F}(Df_{ji})$; so if $P_0$ is
tangent-good in degree $k$, these bundles are analytic for every Banach space $F$, by the paragraph after
Definition~\ref{fam2:def:good}.

\begin{theorem}[shear charts]\label{fam2:thm:tangent}
\begin{enumerate}
\item Let $K$ be complete, let $P_0$ be a Banach space with a bounded linear isomorphism $J\colon P_1\times D\times E\to P_0$,
let $U\subseteq P_1$ be open and let $u\colon U\to\cL(E,D)$ be analytic. Then $\psi(x,d,e)=(x,d+u(x)e,e)$ is an analytic
diffeomorphism of $U\times D\times E$ onto itself, with inverse $(x,d,e)\mapsto(x,d-u(x)e,e)$, and with $i_E(e)=(0,0,e)$
and $\pi_D(x,d,e)=d$,
\[
 \Lambda^k_\pi u(x)=\hLambda^k_\pi\pi_D\circ\Lambda^k_\pi D\psi(x,d,e)\circ\Lambda^k_\pi i_E.
\]
If $\Lambda^k_\pi u$ is not analytic at $x_0$, then neither $(x,d,e)\mapsto\Lambda^k_\pi D\psi(x,d,e)$ nor
$(x,d,e)\mapsto\Lambda^k_\pi D(\psi^{-1})(x,d,e)$ is analytic at any point $(x_0,d,e)$, and the manifold
$M=J(U\times D\times E)$ with the charts $\id_M$ and $J\circ\psi\circ J^{-1}$ has bundles $\hLambda^k_\pi TM$ and
$\Alt^k(TM;\hLambda^k_\pi D)$ whose induced atlases are not analytic.
\item Let $K$ be complete and $P_0$ a Banach space with $P_0\cong P_0\times P_0\times P_0$. The following are equivalent:
$P_0$ is tangent-good in degree $k$; for every analytic manifold $M$ modeled on $P_0$ the induced atlases of
$\hLambda^k_\pi TM$ and of $\Alt^k(TM;F)$, for every Banach space $F$, are analytic; $(P_0;P_0)$ is good in degree $k$.
\item Assume~\textup{(L)}. For $E=c_0(\N,K)$ and for $E=\ell^r(\N,K)$ with $4<r<\infty$, the polynomial map
$\psi(a,d,e)=(a,d+a\cdot e,e)$ on $E^3$, with $a\cdot e$ the pointwise product, has the polynomial inverse
$(a,d,e)\mapsto(a,d-a\cdot e,e)$, and $x\mapsto\Lambda^k_\pi D\psi(x)$ and $x\mapsto\Lambda^k_\pi D(\psi^{-1})(x)$ are
analytic at no point of $E^3$; since $E\cong E^3$, $E$ is not tangent-good in degree $k$. If $K$ is complete, no
infinite-dimensional Banach space with an equivalent nonarchimedean norm is tangent-good in degree $k$.
\end{enumerate}
\end{theorem}

The proof is in Appendix~\ref{app:families}.

\subsection{Structure groups}\label{fam2:ssec:groups}

\begin{definition}\label{fam2:def:groups}
Let $E$ be a normed space and $\operatorname{GL}(E)$ the group of invertible elements of $\cL(E)$. For a subset
$\mathsf S\subseteq\cL(E)$, an \emph{$\mathsf S$-family} over $P$ is an analytic map $u\colon U\to\cL(E)$, with $U\subseteq P$
open, such that $u(U)\subseteq\mathsf S$; and $\mathsf S$ is \emph{good over $P$ in degree $k$} if $x\mapsto\Lambda^k_\pi u(x)$
is analytic on $U$ for every $\mathsf S$-family $u$ over $P$. For a linear subspace $\mathsf V\subseteq\cL(E)$ and
$\varepsilon>0$ put $\mathsf V(\varepsilon)=\{v\in\mathsf V:\norm v<\varepsilon\}$ and
$1+\mathsf V(\varepsilon)=\{1+v:v\in\mathsf V(\varepsilon)\}$, with $1=\id_E$.
\end{definition}

By Theorem~\ref{thm:bundle-lifting}, a bundle with transitions in a subgroup $\mathsf H\subseteq\operatorname{GL}(E)$ needs
only $\mathsf H$-families.

\begin{theorem}[structure groups]\label{fam2:thm:groups}
\begin{enumerate}
\item Let $\mathsf H\subseteq\operatorname{GL}(E)$ be a subgroup, $\mathsf H_0\subseteq\mathsf H$ a subgroup that is open in
$\mathsf H$ for the operator norm, and $\mathsf V_0\subseteq\cL(E)$ a linear subspace with $\mathsf H_0\subseteq\mathsf V_0$
and $1+\mathsf V_0(\varepsilon)\subseteq\mathsf H_0$ for some $\varepsilon>0$. Then $\mathsf H$ is good over $P$ in degree
$k$ if and only if $\mathsf V_0$ is. Neither $E$ nor $K$ need be complete. If $E$ is a Banach space, $\operatorname{GL}(E)$ is
good over $P$ in degree $k$ if and only if $(P;E)$ is good in degree $k$.
\item If $E$ is a Banach space with a nonarchimedean norm, the group $\operatorname{Isom}(E)$ of surjective
linear isometries of $E$ is good over $P$ in degree $k$ if and only if $(P;E)$ is good in degree $k$.
\item Let $K$ be nonarchimedean and $E=c_0(\N,K)$ or $E=\ell^r(\N,K)$ with $1\le r<\infty$. Let
$\mathsf V_{\mathrm d}=\{\mu_a:a\in\linf(\N,K)\}$, let
$\mathsf H_{\mathrm d}=\{\mu_a:a\in\linf(\N,K),\ \inf_t\abs{a(t)}>0\}$ be the diagonal group and
$\mathsf H_{\mathrm m}=\{\sigma_\sharp\mu_a:\mu_a\in\mathsf H_{\mathrm d},\ \sigma\text{ a bijection of }\N\}$ the monomial
group, where $\sigma_\sharp e_t=e_{\sigma(t)}$. Then $\mathsf H_{\mathrm m}$ is good over $P$ in degree $k$ if and only if
$\mathsf H_{\mathrm d}$ is, and if and only if $\mathsf V_{\mathrm d}$ is.
\item Under~\textup{(L)}, the commutative group $\mathsf H_{\mathrm d}$ of $E=c_0(\N,K)$ is not good over $P=c_0(\N,K)$ in
degree $k$.
\end{enumerate}
\end{theorem}

Corollary~\ref{lau:cor:diagonal-bundle} gives isometric diagonal transitions with the same behavior over Laurent series
fields.

\begin{proposition}[the closed span does not decide]\label{fam2:prop:span}
Assume~\textup{(L)} with $K$ complete, and let $\kappa'$ and $\vartheta$ be as in Lemma~\ref{obs:lem:normalization}, with
$\kappa'$ identified with its image in $K$, so that the nonzero elements of $\kappa'$ have absolute value $1$ and every
element of $K$ is a convergent series $\sum_{n\ge n_0}\xi_n\vartheta^n$ with $\xi_n\in\kappa'$. Let $E=D=c_0(\N,K)$, let
$e^*_i$ be the $i$th coordinate functional, let $\mathsf U\subseteq\cL(E)$ be the additive group of the operators $u$ with
$e^*_i(ue_j)\in\kappa'$ for all $i,j$, and let $\mathsf H=\{g(u):u\in\mathsf U\}\subseteq\operatorname{GL}(D\oplus E)$, with
$g$ as in Proposition~\ref{prop:shear}. Then $\mathsf H$ is a discrete subgroup, every $\mathsf H$-family over every $P$ is
locally constant, and $\mathsf H$ is good over every $P$ in every degree. But the closed linear span of $\mathsf H$
contains $g(\mu_x)$ for every $x\in c_0(\N,K)$, and $x\mapsto\Lambda^k_\pi g(\mu_x)$ is analytic at no point of
$c_0(\N,K)$.
\end{proposition}

The proofs are in Appendix~\ref{app:families}.

\subsection{Characteristic two}\label{fam2:ssec:char-two}

\begin{proposition}[families in characteristic two]\label{fam2:prop:char-two}
Let $K$ be a nontrivially normed field of characteristic~$2$, not necessarily satisfying~\textup{(L)}, and let $k\ge2$.
\begin{enumerate}
\item Let $E$, $P$, $\varphi^\pm_E$ and $\varphi^\pm_P$ be as in Theorem~\ref{fam2:thm:staircase}. If
$\sup_nn/(\varphi^\pm_P(n)^2\varphi^\pm_E(n)^2)=\infty$, the conclusions of Theorem~\ref{fam2:thm:staircase} hold.
\item The conclusion of Corollary~\ref{fam2:cor:table}(2) holds for $P=c_0(\N,K)$ and $E$ equal to $c_0(\N,K)$ or to
$\ell^{r_2}(\N,K)$ with $2<r_2<\infty$, and for $P=\ell^{r_1}(\N,K)$ with $2<r_1<\infty$ and $E$ equal to $c_0(\N,K)$ or
to $\ell^{r_2}(\N,K)$ with $1/r_1+1/r_2<1/2$.
\item For $E=c_0(\N,K)$ and for $E=\ell^r(\N,K)$ with $4<r<\infty$, the maps $x\mapsto\Lambda^k_\pi D\psi(x)$ and
$x\mapsto\Lambda^k_\pi D(\psi^{-1})(x)$ of Theorem~\ref{fam2:thm:tangent}(3) are analytic at no point of $E^3$, and $E$
is not tangent-good in degree $k$.
\item The commutative group $\mathsf H_{\mathrm d}$ of Theorem~\ref{fam2:thm:groups}(3), for $E=c_0(\N,K)$, is not good
over $P=c_0(\N,K)$ in degree $k$.
\end{enumerate}
\end{proposition}

The proof is in Appendix~\ref{app:families}. Item~(2) does not cover the pair $P=c_0(\N,K)$, $E=\ell^2(\N,K)$ of
Corollary~\ref{fam2:cor:table}(2). Items~(2) and~(3) do not cover the spaces with an equivalent nonarchimedean norm in
Corollary~\ref{fam2:cor:table}(2) and in the last sentence of Theorem~\ref{fam2:thm:tangent}(3), whose proofs use
Lemma~\ref{fib:lem:marking}, which assumes~\textup{(L)}.
%% ---- end sec-10b-families.tex
%% ---- begin sec-11-questions.tex
\section{Questions}\label{sec:questions}

The questions below concern $\operatorname{char}K=p>0$ and $k\ge p$. For $1<r<\infty$ and $1\le b\le\infty$, the
symmetric sequence space $\ell^{r,b}$ (Definition~\ref{obs:def:sequence-spaces}) is the closure of $c_{00}$ in the space
of all $x\in\linf(\N,K)$ with $\sup_n\norm{x1_{[n]}}_{r,b}<\infty$, where on $c_{00}$
$\norm x_{r,b}=c_{r,b}\bigl(\sum_{n\ge1}n^{b/r-1}(x^{**}_n)^b\bigr)^{1/b}$ for $b<\infty$,
$\norm x_{r,\infty}=c_{r,\infty}\sup_{n\ge1}n^{1/r}x^{**}_n$, $x^{**}_n$ is the average of the $n$ largest values of
$\abs{x(t)}$, and $c_{r,b}>0$ makes $\norm{e_1}_{r,b}=1$.

\begin{conjecture}[product triangle conjecture]\label{q:conj:pt}
Let $\kappa$ and $K$ be as in \textup{(L)}, and let $E,D$ be symmetric sequence spaces in which $c_{00}$ is dense. If $(E,D)$
satisfies \textup{(PTB)} (Definition~\ref{obs:def:ptb}), then $(E,D)$ is good for Banach values in every degree.
\end{conjecture}

The converse holds (Theorem~\ref{obs:thm:ptb}). The conjecture holds when $D=E$ or $E=c_0$, since indicator functions of
supports in~\eqref{obs:eq:ptb} bound $\cP_<$ on $D\otimes_\pi D$ (Theorem~\ref{cls:thm:symmetric}), and when $D$ has an
equivalent $r$\nobreakdash-norm, $2<r<\infty$, or nonarchimedean norm (Theorem~\ref{cls:thm:rnorm}). The same holds for
spreading sequence spaces (Theorem~\ref{cls:thm:spreading}).

\begin{conjecture}\label{q:conj:g}
Let $\kappa$ be a finite field of characteristic $p$, $K=\kappa((X))$ with the absolute value of~\textup{(L)}, and $1<r<2\le b\le\infty$. Then $\cP_<$ is unbounded on
$\ell^{r,b}\otimes_\pi\ell^{r,b}$.
\end{conjecture}

By Theorem~\ref{cls:thm:symmetric}, the conjecture for given $r,b$ is equivalent to the failure of
$(\mathrm S^p_{\ell^{r,b}})$, of $(\mathrm R^p_{\ell^{r,b}})$, or of $(\Gamma^p_{\ell^{r,b}})$, and to $(\ell^{r,b},\ell^{r,b})$
and $(c_0,\ell^{r,b})$ being good for Banach values in no degree $k\ge p$.

\begin{conjecture}[sequence lattices]\label{q:conj:lattice}
Assume \textup{(L)} with $K$ complete. Let $E$ and $D$ be sequence lattices, that is, each is a Banach space $Y$ with
$c_{00}(\N,K)\subseteq Y\subseteq\linf(\N,K)$ in which $c_{00}$ is dense, with $\norm{\cdot}_\infty\le\norm{\cdot}_Y$,
$\sup_t\norm{e_t}_Y<\infty$, and $\norm x_Y\le\norm y_Y$ whenever $\abs{x(t)}\le\abs{y(t)}$ for every $t$. Then $(E,D)$ is good for Banach values in
degree $k$ if and only if it has a splitting decomposition of degree $k$ (Definition~\ref{cmb:def:decomposition}).
\end{conjecture}

\begin{question}\label{q:quest:generation}
Is every pair that is good for Banach values in all degrees $\le k$ obtained from pairs $(E,D)$ with $(\mathrm S^n_E)$
for all $n\le k$, pairs with $(\mathrm R^n_D)$ for all $n\le k$, relative splittings through $\mathrm S^{\le k}$-split or
$\mathrm R^{\le k}$-split maps, and parameter pieces, by finite decompositions (Theorem~\ref{cmb:thm:decomposition},
Proposition~\ref{cmb:prop:parameter-pieces}), transfer (Proposition~\ref{cmb:prop:transfer}) and block assemblies
(Proposition~\ref{cmb:prop:blocks})?
\end{question}

\begin{question}\label{q:quest:coincide}
Let $\kappa$ and $K$ be as in \textup{(L)}. Do $(\mathrm S^k_Y)$, $(\mathrm R^k_Y)$ and $(\Gamma^k_Y)$ coincide for every
Banach space $Y$ and every $k\ge p$?
\end{question}

They coincide for spreading, in particular symmetric, sequence spaces in which $c_{00}$ is dense
(Theorems~\ref{cls:thm:symmetric} and~\ref{cls:thm:spreading}), and
$(\mathrm S^2_Y)$ is equivalent to $(\Gamma^2_Y)$ for every normed $Y$ (Proposition~\ref{spl:prop:degree-two}).

\begin{question}\label{q:quest:three-space}
Let $W$ be a closed subspace of a Banach space $Z$. If $W$ and $Z/W$ satisfy $(\mathrm S^k)$, does $Z$? The same question
for $(\mathrm R^k)$ and for polarizability in degree $k$.
\end{question}

Under \textup{(L)}, goodness for Banach values with a fixed target has no three-space property in the source slot, even
with an $\ell^1$ kernel (Theorem~\ref{cmb:thm:three-space}).

\begin{question}\label{q:quest:source-split}
Let $k\ge p$. If $(E,D)$ is good for Banach values in degree $k$ for every normed $D$, does $(\mathrm S^k_E)$ hold?
\end{question}

Under \textup{(L)} the answer is yes for spreading, in particular symmetric, sequence spaces $E$ in which $c_{00}$ is
dense (Theorems~\ref{cls:thm:symmetric} and~\ref{cls:thm:spreading}).

\begin{question}\label{q:quest:factorization}
Assume \textup{(L)}. Does every $\mathrm S^k$-split operator factor through a normed space $G$ with $(\mathrm S^k_G)$? Does
the inclusion $\ell^2(\N,K)\to c_0(\N,\Kh)$, which is $\mathrm S^n$-split for every $n$
(Corollary~\ref{spl:cor:relative-ordering-na}), factor through such a space?
\end{question}

\begin{question}\label{q:quest:residue}
Let $K$ be a field of Laurent series in $X$ over an infinite field of characteristic $p$, with an $X$-adic absolute
value, or a densely valued field of characteristic $p$.
\begin{enumerate}
\item If $p$ is odd, is $(c_0(\N,K),c_0(\N,K))$ good for Banach values in degree $p$?
\item If $p=2$, is $(\ell^2(\N,K),\ell^2(\N,K))$ good for Banach values in degree $2$?
\end{enumerate}
\end{question}

For $p=2$, $(c_0(\N,K),c_0(\N,K))$ is good for Banach values in no degree $k\ge2$, over every nontrivially normed field
of characteristic~$2$ (Theorem~\ref{obs:thm:char-two}(4)); for $E=\ell^2(\N,K)$, $\bar\varphi_E(n)^2=n$ and
\eqref{obs:eq:multiplier-bound} gives no contradiction. For odd $p$, a positive answer to~(1) needs a zero-diagonal
correction making the ordering lift bounded (Theorem~\ref{ob:thm:ordering}).

\begin{question}\label{q:quest:height}
If a pair is good in degree $p$, is it good in every degree? The same question for goodness for Banach values. Is every
$\mathrm S^p$-split operator $\mathrm S^n$-split for every $n$, and every $\mathrm R^p$-split operator
$\mathrm R^n$-split for every $n$?
\end{question}

The least degree in which a pair fails to be good is a power of $p$ (Corollary~\ref{fib:cor:degrees}). For the
conditions $(\mathrm S^p_Y)$ and $(\mathrm R^p_Y)$ of one normed space $Y$, Theorem~\ref{cmb:thm:column} answers the third question.

\begin{question}\label{q:quest:pairs}
Let $\operatorname{char}K=p$.
\begin{enumerate}
\item If a normed space $P$ satisfies $(\mathrm U^p_P)$ (Definition~\ref{fam2:def:reflection}), is $P$ polarizable in
degree $p$?
\item Assume \textup{(L)}. Is $(\ell^{r_1}(\N,K);\ell^{r_2}(\N,K))$ good in degree $k$ when $r_1=r_2=2$, when
$2<r_1<\infty$ and $r_2=2$, or when $2<r_1,r_2<\infty$ and $1/r_1+1/r_2\ge1/2$?
\item Do Theorems~\ref{fam2:thm:reduction} and~\ref{fam2:thm:class} hold for one Banach value space $F$, with
$x\mapsto A^k_{E,\widehat D;F}(u(x))$ in place of $x\mapsto\Lambda^k_\pi u(x)$?
\end{enumerate}
\end{question}

Polarizability in the degrees $n\le k$ gives $(\mathrm U^k_P)$ (Theorem~\ref{fam2:thm:polarizable}(2)); under
\textup{(L)}, $(\mathrm U^k_P)$ gives polarizability in degree $p$ when $P$ is a symmetric sequence space in which
$c_{00}(\N,K)$ is dense or has an equivalent nonarchimedean norm (Theorem~\ref{fam2:thm:bases}). The cases in~(2) are those left open by
Corollary~\ref{fam2:cor:table}; for $r_1=r_2=r$ and $K$ complete they ask whether $\ell^r(\N,K)$ is tangent-good in
degree $k$ (Theorem~\ref{fam2:thm:tangent}(2)).

\begin{question}\label{q:quest:linear-families}
Let $P$, $E$ and $D$ be normed spaces. If every linear family $P\to\cL(E,\widehat D)$ is liftable in degree $k$, is
$(P;E,D)$ good in degree $k$?
\end{question}

The answer is yes for $P=\cL(E,\widehat D)$, whose identity is a linear family (Lemma~\ref{fib:lem:test},
Proposition~\ref{fam2:prop:fibers}(2)); for fixed-degree reflection, affine families suffice
(Theorem~\ref{fam2:thm:polarizable}(1)).

\begin{question}\label{q:quest:lower-degrees}
If $(P;E,D)$ is good in degree $k$, is it good in every degree $j<k$?
\end{question}

The answer is yes when $E\cong E\oplus K$ and $D\cong D\oplus K$, and for goodness for all normed $E$ and $D$
(Corollary~\ref{fam2:cor:degrees}(2),~(4)).

\begin{question}\label{q:quest:non-uniform}
Is there a linear family $\theta\colon\mathcal T(P)\to\cL(E,\widehat D)$ for which every class $\ob_{k,N}(\theta)$ vanishes
but $\ob^\omega_k(\theta)$ does not vanish after shrinking?
\end{question}

By Theorem~\ref{fam2:thm:class}(1), $\ob^\omega_k(\theta)$ vanishes after shrinking exactly when the classes
$\ob_{k,N}(\theta)$ vanish with extension costs growing at most geometrically in $N$. For general closed subspaces, Taylor
coefficients can have lifts without a common radius (Corollary~\ref{fam:cor:no-common-radius}).

\begin{question}\label{q:quest:composite}
Let $E$, $D$ and $F$ be normed spaces. If $(E,D)$ and $(D,F)$ are good for Banach values in degree $k$, is the
composite family $\theta_{D,F}\odot\theta_{E,D}$ of Proposition~\ref{cmb:prop:composition} liftable in degree $k$?
\end{question}

The composite family is liftable in degree $k$ when $(E,F)$ is good for Banach values in that degree
(Proposition~\ref{cmb:prop:composition}(2)).

\begin{question}\label{q:quest:garling}
Let $\kappa$ and $K$ be as in \textup{(L)}, $1\le q<2$ and $0<\alpha<1-q/2$. Is $(E_{\alpha,q},E_{\alpha,q})$
(Example~\ref{cls:ex:garling}) good for Banach values in some degree $k\ge p$?
\end{question}

By Theorem~\ref{cls:thm:spreading}, this holds if and only if $\cP_<$ is bounded on $E_{\alpha,q}\otimes_\pi E_{\alpha,q}$,
and then the pair is good for Banach values in every degree. The necessary conditions of Theorem~\ref{cls:thm:symmetric}
hold, and Theorem~\ref{fib:thm:lower} does not apply to the unit vectors in their natural order
(Example~\ref{cls:ex:garling}(2),(4)).

\begin{question}\label{q:quest:polynomial-laws}
Let $K$ be nonarchimedean of characteristic $p>0$. Is every bounded polynomial law on $c_0(\N,K)$ ambient of its degree?
\end{question}

Bounded polynomial laws of degree below $p$ are ambient (Section~\ref{sec:dp}), and for $p=2$ so are those of degree
$2$ on nonarchimedean normed spaces (Proposition~\ref{dp:prop:ultrametric}(1),(4)). By
Proposition~\ref{dp:prop:ultrametric}(1),(3), every bounded polynomial law of degree $n$ on $c_0(\N,K)$ is ambient of
degree $n$ if and only if every bounded polynomial law of degree $n$ on every nonarchimedean normed space over $K$ is.
As $K\oplus c_0(\N,K)$, with the maximum norm, is isometric to $c_0(\N,K)$, Proposition~\ref{dp:prop:ultrametric}(1),(2) shows that the degrees
with a positive answer form an initial segment of the positive integers: a negative answer in some degree gives a
negative answer in every higher degree, and a positive answer in infinitely many degrees gives one in all. Frobenius
laws are ambient, and every bounded polynomial law of degree $p$ on a normed space $P$, with values in any normed space,
is ambient of degree $p$ if and only if $P$ satisfies $(\mathrm N_P)$ with some constant
(Theorem~\ref{dp:thm:frobenius}(2),(3)); so for $p\ge3$, in degree $p$ the question asks whether $c_0(\N,K)$ satisfies
$(\mathrm N_{c_0(\N,K)})$ with some constant. Over every $K$ and in every degree $n\ge p$,
Proposition~\ref{dp:prop:degree-p}(2) gives a law that is not ambient of degree $n$, on a space that is a Banach space if
$K$ is complete; if $K$ is discretely valued or $p=2$, Proposition~\ref{dp:prop:nonambient}(2) gives such laws in every
degree $n\ge p+1$ also on the space $H=K\oplus c_0(\N,K)\oplus G$ of Section~\ref{sec:dp}. These spaces have no
equivalent nonarchimedean norm: they contain $\ell^1$-sums of finite-dimensional quotients of weighted $\ell^1$-spaces,
and both proofs obtain their lower bounds, for $\nu$ and for $\norm\cdot_\Gamma$, by adding up contributions of the
coordinates of these $\ell^1$-spaces.

\paragraph{\textit{Targets, bases and domains.}}
An equivalent nonarchimedean, spherically complete norm makes a Banach space $F$ a universal target
(Section~\ref{obs:ssec:counterexamples}, Theorem~\ref{thm:spherical}), but the given norm of $F$ need not be spherically
complete (Remark~\ref{rem:nonspherical-analytic-target}). Over complete, not spherically complete fields universal targets have
zero dual (Proposition~\ref{prop:universal-zero-dual}); we do not know whether that suffices, nor whether over densely
valued, spherically complete fields every nonarchimedean Banach space is one. We do not know which $P$ satisfy
$(\mathrm U^k_P)$ (Question~\ref{q:quest:pairs}(1)); for Banach $P$ over complete fields universal analytic reflection
suffices (Theorem~\ref{fam:thm:tensor-analytic}) but, over nonarchimedean fields, is not necessary
(Theorem~\ref{fam2:thm:polarizable}(4)). We do not know whether compatibility with all analytic coordinate changes forces
the ambient and intrinsic analytic forms to agree; one product structure does not
(Proposition~\ref{prop:tangent-obstruction}). Over the $t$-adic field $\F_p(t)$ a largest full analytic domain need not
exist (Theorem~\ref{dom:no-largest}); we do not know whether this happens over complete fields. Under \textup{(L)} with
$K$ complete, $\Alt^k(TM;K)$ is analytic on every analytic manifold $M$ (Corollary~\ref{cor:discrete-targets}); we do
not know this over incomplete fields such as $\kappa(X)$, nor for tangent-valued forms on $\ell^r(\N,K)$, $4<r<\infty$
(Theorem~\ref{fam2:thm:tangent}(3)). Nor do we know whether Bourbaki's $\alpha_k$ is a vector functor of class $C^r$
over densely valued, spherically complete fields with $k!=0$.
%% ---- end sec-11-questions.tex
%% ---- begin sec-12-verification.tex
\section{Formal verification}\label{sec:verification}

The Lean proofs are in~\cite{LeanP}. Its file \texttt{challenge.md} is a ledger of 76 claims of this paper. Each claim
has a challenge file, stating it in Lean with \texttt{sorry} as proof, and a solution file proving the same statement.
For a repository identifier \texttt{ID} such as \texttt{Thm4\_4}, \texttt{python3 scripts/check\_claim.py ID} builds
both files and checks that the proof uses only the axioms \texttt{propext}, \texttt{Classical.choice} and
\texttt{Quot.sound} and that the two statements agree when printed with all implicit arguments (\texttt{pp.all}). All 76
claims pass with Lean~\texttt{4.34.0-rc2} and the Mathlib fork \texttt{Deicyde/mathlib4} at \texttt{2b73d98}
(upstream \texttt{0383a80e64} plus the commit of~\cite{PR}). The six statements of the comparator challenge in the
same repository pass the comparator \url{https://github.com/leanprover/comparator} at revision \texttt{19e111e}.

Table~\ref{ver:tab:ids} gives the identifiers of the 24 claims in Sections~\ref{sec:bundles}--\ref{sec:geometry}. The
other 52 claims are the appendix statements \ref{lem:unique}--\ref{lem:findim}, \ref{lem:detinj}--\ref{lem:D},
\ref{rem:FFunit}, \ref{thm:laurent}--\ref{lem:Psi}, \ref{lau:prop:c0-multipliers}--\ref{lau:cor:dichotomy},
\ref{lem:discdist}--\ref{prop:ascent}, \ref{lem:descent}, \ref{prop:nonultra}, \ref{rem:spaces},
\ref{thm:scalar-counter}--\ref{scalar:certificate}, \ref{fam:prop:tensor-homogeneous},
\ref{fam:thm:tensor-analytic}--\ref{fam:cor:no-common-radius}, \ref{dom:split-incoming}--\ref{dom:degree-gap} and
\ref{prop:scalar-countable}--\ref{rem:nonspherical-analytic-target}; the identifier of each is formed from its type and
number, for example \texttt{LemB\_2} for Lemma~\ref{lem:detinj}.

\begin{table}[htbp]
\centering
\small
\caption{Repository identifiers of the ledger claims in Sections~\ref{sec:bundles}--\ref{sec:geometry}}\label{ver:tab:ids}
\begin{tabular}{ll@{\qquad}ll}
\hline
Statement & Identifier & Statement & Identifier\\
\hline
Theorem~\ref{thm:bundle-lifting} & \texttt{Thm2\_1} & Theorem~\ref{thm:summable-parameters}(2) & \texttt{Thm4\_5b}\\
Proposition~\ref{prop:ambient-polynomial} & \texttt{Prop2\_2} & Corollary~\ref{cor:bundle-cases} & \texttt{Cor4\_6}\\
Theorem~\ref{thm:coefficient-descent} & \texttt{Thm3\_1} & Theorem~\ref{thm:operator-obstructions}(1) & \texttt{Thm6\_1a}\\
Lemma~\ref{lem:closed-smooth} & \texttt{Lem3\_2} & Theorem~\ref{thm:operator-obstructions}(2) & \texttt{Thm6\_1b}\\
Proposition~\ref{prop:homogeneous-lift} & \texttt{Prop3\_3} & Corollary~\ref{cor:classification} & \texttt{Cor6\_2}\\
Proposition~\ref{prop:exterior} & \texttt{Prop5\_1} & Theorem~\ref{thm:scalar-classification} & \texttt{Thm6\_3}\\
Theorem~\ref{thm:duality} & \texttt{Thm5\_2} & Proposition~\ref{prop:universal-zero-dual} & \texttt{Sec9\_zeroDual}\\
Proposition~\ref{prop:split-lift} & \texttt{Prop4\_1} & Proposition~\ref{prop:shear} & \texttt{Prop6\_4}\\
Theorem~\ref{thm:spherical} & \texttt{Thm4\_2} & Corollary~\ref{cor:shear-bundle} & \texttt{Prop6\_4\_bundle}\\
Corollary~\ref{cor:discrete-targets} & \texttt{Cor4\_3} & Proposition~\ref{prop:tangent-obstruction} & \texttt{Prop6\_5}\\
Theorem~\ref{thm:finite-parameters} & \texttt{Thm4\_4} & Theorem~\ref{thm:ambient-forms} & \texttt{Thm7\_2}\\
Theorem~\ref{thm:summable-parameters}(1) & \texttt{Thm4\_5a} & Proposition~\ref{prop:ambient-agreement} & \texttt{Prop7\_3}\\
\hline
\end{tabular}
\end{table}

A numbered statement of this paper has a Lean proof of a corresponding statement exactly when it is one of the 24
claims of Table~\ref{ver:tab:ids} or the 52 claims just listed. Excluded are the definitions, the conjectures and
questions of Section~\ref{sec:questions}, Remarks~\ref{rem:FFhyp}, \ref{rem:infkappa} and~\ref{rem:ascent},
Proposition~\ref{prop:tangent-values}, the coefficient bound on retracts in Theorem~\ref{thm:summable-parameters}(1),
and all other theorems, propositions, lemmas, corollaries and remarks of
Sections~\ref{sec:duality}--\ref{sec:questions} and Appendices~\ref{app:ordering}--\ref{app:families}. The excluded
results have written proofs only, except Lemma~\ref{rig:lem:ramsey}, which is cited from~\cite{Ra}, and apart from the
Lean declarations named below.

\texttt{challenge.md} and the docstrings of the challenge files list where a ledger statement differs from the text.
The main differences: Theorem~\ref{thm:ambient-forms} and Proposition~\ref{prop:ambient-agreement} are stated on open
subsets of a normed space, since Mathlib has no differential forms on manifolds; the Lean statement of
Theorem~\ref{thm:bundle-lifting} groups the variables by variance, uses the canonical trivializations and does not state
independence of the trivializing cover; Corollary~\ref{lau:cor:diagonal-bundle} is stated for the explicit transition
and section maps, and Corollary~\ref{cor:bundle-cases} for objects and morphisms without the functor laws.

In~\cite{LeanP}, \texttt{alternatingMapAction} is the joint action $B$ of Section~\ref{sec:bundles}, and
the alternating bundle construction rests on the family-preservation premise: for $n\in\N\cup\{\infty,\omega\}$, a
$C^n$ manifold $M$, normed spaces $E,D,F,F'$ and $B\colon\cL(E,D)\times\cL(F,F')\to\cL(\Alt^k(D;F),\Alt^k(E;F'))$,
\begin{equation}\label{eq:mathlib-family-premise}
 U\subseteq M\text{ open},\quad \gamma\colon M\to\cL(E,D)\times\cL(F,F'),\quad \gamma|_U\text{ is }C^n
 \quad\Longrightarrow\quad (B\circ\gamma)|_U\text{ is }C^n.
\end{equation}
The predicate \texttt{PreservesAlternatingFamilies} states~\eqref{eq:mathlib-family-premise} and typechecks
against~\cite{LeanP}. The declarations
\begin{leancode}
contMDiffOn_continuousAlternatingMapCoordChange_of_family
alternatingVectorPrebundle_isContMDiff_of_family
contMDiffVectorBundle_alternating_of_family
contMDiffVectorBundle_alternating_of_finiteCoordinates
alternatingBundleHom_of_family
\end{leancode}
of~\cite{LeanP} index alternating maps by \texttt{Fin k}. The first three assume~\eqref{eq:mathlib-family-premise}
with $D=E$ and $F'=F$ and show that the alternating coordinate changes, prebundle and bundle are $C^n$. The fourth shows
that over an analytic manifold modeled on a space $P$ with a continuous linear isomorphism $P\to K^d$ the alternating
bundle of two analytic bundles is analytic, with the premise from Theorem~\ref{thm:finite-parameters}. The fifth
assumes~\eqref{eq:mathlib-family-premise} and turns $C^n$ bundle morphisms into a $C^n$ morphism of the induced
alternating bundles when these are $C^n$.
%% ---- end sec-12-verification.tex
\appendix
%% ---- begin app-A-preliminaries.tex
\section{Preliminaries for the obstruction proofs}\label{app:obstruction-preliminaries}
The constructions below use three elementary facts about valued fields and power series. We record their proofs to make explicit the scope of the coefficient argument, including incomplete base fields.

\begin{lemma}[one-variable uniqueness]\label{lem:unique}
Let $K$ be nontrivially normed and $Z$ a normed $K$\nobreakdash-space, not necessarily complete. Let $z_0,z_1,\ldots\in Z$, $\rho>0$
and $M\ge0$ with $\norm{z_n}\rho^n\le M$ for all $n$. If $\sum_nt^nz_n$ converges to $0$ for every $t\in K$ with
$\abs t<\rho$, then $z_n=0$ for all $n$.
\end{lemma}

\begin{proof}
Otherwise let $N$ be least with $z_N\neq0$. For $0<\abs t<\rho/2$ the series $\sum_{n>N}t^{n-N}z_n$ converges, being
$t^{-N}$ times the convergent tail $\sum_{n>N}t^nz_n$, and $t^N\bigl(z_N+\sum_{n>N}t^{n-N}z_n\bigr)=0$. Hence
\[
\norm{z_N}=\Bigl\lVert\sum_{n>N}t^{n-N}z_n\Bigr\rVert\le\sum_{n>N}M\rho^{-n}\abs t^{n-N}
=M\rho^{-N}\frac{\abs t/\rho}{1-\abs t/\rho},
\]
where the norm of a convergent series is bounded by the sum of the norms by continuity of the norm on partial sums.
Since $K$ is nontrivially normed, there are $t\ne0$ with $\abs t$ arbitrarily small, so $\norm{z_N}=0$, a
contradiction.
\end{proof}

\begin{lemma}\label{lem:ultra}
A normed field of characteristic $p>0$ is nonarchimedean.
\end{lemma}

\begin{proof}
Frobenius and the triangle inequality give, for every $m\ge0$,
\[
 \abs{x+y}^{p^m}=\abs{x^{p^m}+y^{p^m}}
 \le2\max(\abs x,\abs y)^{p^m}.
\]
Taking $p^m$\nobreakdash-th roots and letting $m\to\infty$ proves the assertion.
\end{proof}

\begin{lemma}\label{lem:findim}
Let $\mathbb K$ be a complete nontrivially normed field and $U$ a finite-dimensional normed $\mathbb K$\nobreakdash-space. Then every
linear functional on $U$ is bounded, and $U$ is complete.
\end{lemma}

\begin{proof}
We prove both assertions by induction on $n=\dim U$, the case $n=0$ being immediate. Suppose $n\ge1$ and let $g\ne0$
be a linear functional on $U$. Its kernel $H$ has dimension $n-1$, so the induction hypothesis makes $H$ complete
and therefore closed in $U$. Choose $v$ with $g(v)=1$. Since $v\notin H$,
$\delta:=\inf_{h\in H}\norm{v-h}>0$. Write $x=g(x)v+h$ with $h\in H$. If $g(x)\ne0$, then
\[
\norm x=\abs{g(x)}\,\norm{v+g(x)^{-1}h}\ge\abs{g(x)}\delta.
\]
Thus $\abs{g(x)}\le\delta^{-1}\norm x$ for every $x$, proving boundedness.

Choose a basis of $U$. Its coordinate functionals are now bounded, so the coordinate map
$U\to\mathbb K^n$, with the max norm on $\mathbb K^n$, is bounded. The inverse map
$(c_i)\mapsto\sum_ic_ie_i$ is bounded by $\sum_i\norm{e_i}$. Hence $U$ is bi-Lipschitz equivalent to the complete
space $\mathbb K^n$, and is itself complete.
\end{proof}

When the scalar field needs to be displayed, $A^{k,L}$ denotes the degree-$k$ precomposition action formed over $L$. For the precomposition action $A$, a \emph{bounded lift} means a bounded $k$\nobreakdash-linear map on its operator domain whose diagonal is $A$, as in Proposition~\ref{prop:homogeneous-lift}. That proposition applies at every base point; Lemma~\ref{lem:unique} provides the one-variable coefficient comparison used in its proof.
%% ---- end app-A-preliminaries.tex
%% ---- begin app-B-multiplier.tex
\section{The multiplier theorem over a finite field}\label{app:FF}\label{sec:FF}

This appendix proves Theorem~\ref{thm:FF}: over a finite field with $k!=0$, coordinatewise multiplication has no lift
satisfying~\eqref{eq:Pol1} with uniformly bounded support dimension. Ramsey's theorem reduces the coefficients to
cluster values (Lemmas~\ref{lem:H},~\ref{lem:C}), \eqref{eq:Pol1} makes them nonconstant (Lemma~\ref{lem:D}), and the
support bound forbids the resulting staircase (Lemma~\ref{lem:S}).

\subsection{Exterior powers, supports and determinant arrays}\label{ssec:ext}
Let $L$ be a field, $V$ an $L$\nobreakdash-vector space and $k\ge1$. We write $\Lambda^kV=\Lambda^k_LV$ for the algebraic exterior
power, in which $x_1\wedge\dots\wedge x_k=0$ whenever two $x_i$ coincide. We use the standard facts: for a subspace
$W\subseteq V$ the map $\Lambda^kW\to\Lambda^kV$ is injective, and we regard $\Lambda^kW$ as a subspace of $\Lambda^kV$;
if $u_1,\dots,u_d$ is a basis of a finite-dimensional space $U$, the wedges $u_I:=u_{i_1}\wedge\dots\wedge u_{i_k}$,
$I=\{i_1<\dots<i_k\}$, form a basis of $\Lambda^kU$; and every $\omega\in\Lambda^kV$ lies in $\Lambda^kU$ for some
finite-dimensional $U\subseteq V$.

\begin{definition}\label{def:supp}
The \emph{support dimension} of $\omega\in\Lambda^kV$ is
\[
\sdim(\omega):=\min\{\dim W:\ W\subseteq V\text{ finite-dimensional},\ \omega\in\Lambda^kW\}.
\]
\end{definition}

The definition gives $\sdim(0)=0$, $\sdim(c\omega)=\sdim(\omega)$ for $c\neq0$, $\sdim(\omega+\omega')\le\sdim(\omega)+\sdim(\omega')$
(use $W+W'$), and $\sdim(\omega)\le\dim\spn S$ if $\omega$ is a sum of wedges of vectors from a finite set $S$. If
$V\subseteq V'$, then $\sdim$ computed in $V'$ is at most $\sdim$ computed in $V$.

For spaces of functions, exterior vectors can be detected by determinants of coordinate evaluations. These arrays
also turn bounds on support dimension into ordinary matrix-rank bounds.
Let $S$ be a set and $V\subseteq L^S$ a space of functions on $S$. The \emph{determinant array} is the linear map
\[
\Omega\colon\Lambda^kV\to L^{S^k},\qquad \Omega_{y_1\wedge\dots\wedge y_k}(c_1,\dots,c_k):=\det\bigl(y_b(c_a)\bigr)_{a,b=1}^k .
\]
It is well defined because the right side is multilinear and alternating in $(y_1,\dots,y_k)$ (these are the
columns). For each $\omega$, the function $\Omega_\omega$ is alternating in $(c_1,\dots,c_k)$ (the rows); in particular
$\Omega_\omega(c_{\sigma(1)},\dots,c_{\sigma(k)})=\sgn(\sigma)\,\Omega_\omega(c_1,\dots,c_k)$.

\begin{lemma}\label{lem:detinj}
The determinant array $\Omega\colon\Lambda^kV\to L^{S^k}$ is injective.
\end{lemma}

\begin{proof}
Let $\Omega_\omega=0$ and choose a finite-dimensional $U\subseteq V$ with $\omega\in\Lambda^kU$. The evaluations
$\mathrm{ev}_s|_U$, $s\in S$, separate the points of $U$ (a function in $U$ vanishing at every $s$ is $0$), so they span
the dual space $U^*$. Choose $s_1,\dots,s_d\in S$ such that the $\mathrm{ev}_{s_i}|_U$ form a basis of $U^*$, and let
$u_1,\dots,u_d$ be the dual basis of $U$, so that $u_j(s_i)=\delta_{ij}$. For increasing $k$\nobreakdash-subsets $I,J$ of
$\{1,\dots,d\}$ we get $\Omega_{u_I}(s_{j_1},\dots,s_{j_k})=\det(u_{i_b}(s_{j_a}))_{a,b}=\delta_{IJ}$. Writing
$\omega=\sum_Ic_Iu_I$ therefore gives $0=\Omega_\omega(s_{j_1},\dots,s_{j_k})=c_J$ for every $J$, so $\omega=0$.
\end{proof}

\begin{lemma}[support bounds flattening rank]\label{lem:flat}
Let $V\subseteq L^S$, $\omega\in\Lambda^kV$ and $\alpha\in\{1,\dots,k\}$. Let $R\subseteq S$ and
$R'\subseteq S^{\{1,\dots,k\}\setminus\{\alpha\}}$ be finite. The $R\times R'$ matrix $M$ with entries
$M(c_\alpha,c'):=\Omega_\omega(c)$, where $c$ has $c_\alpha$ in slot $\alpha$ and the entries of $c'$ in the other
slots, has rank at most $\sdim(\omega)$.
\end{lemma}

\begin{proof}
Write $\omega=\sum_jy_{j,1}\wedge\dots\wedge y_{j,k}$ with all $y_{j,b}$ in a subspace $W$ of dimension $\sdim(\omega)$.
Expanding each determinant along row $\alpha$,
$\det(y_{j,b}(c_a))_{a,b}=\sum_b\pm y_{j,b}(c_\alpha)\cdot m_{j,b}(c')$, where the minor $m_{j,b}(c')$ does not
involve $c_\alpha$. So for fixed $c'$ the column $c_\alpha\mapsto M(c_\alpha,c')$ is a linear combination of the
restrictions $y_{j,b}|_R$. All columns lie in the restriction of $W$ to $R$, a space of dimension at most
$\sdim(\omega)$.
\end{proof}

A second way to obtain lower bounds on support is to contract all but one exterior factor. The resulting vectors
must lie in every subspace supporting the original exterior vector.

\begin{samepage}
\begin{lemma}[contraction]\label{lem:contr}
Let $k\ge2$ and $\varphi=(\varphi_1,\dots,\varphi_{k-1})\in(V^*)^{k-1}$. There is a linear map $c_\varphi\colon\Lambda^kV\to V$
with
\[
c_\varphi(x_1\wedge\dots\wedge x_k)=\sum_{i=1}^k(-1)^{k+i}\det\bigl(\varphi_a(x_b)\bigr)_{1\le a\le k-1,\ b\ne i}\;x_i ,
\]
and $c_\varphi(\Lambda^kW)\subseteq W$ for every subspace $W$. Consequently
$\sdim(\omega)\ge\dim\spn\{c_\varphi(\omega):\varphi\in(V^*)^{k-1}\}$.
\end{lemma}
\end{samepage}

\begin{proof}
The right side is multilinear in $(x_b)$. For $\psi\in V^*$, applying $\psi$ to it gives, by Laplace expansion along the
last row, the determinant of the $k\times k$ matrix with entries $\varphi_a(x_b)$ in the rows $a<k$ and $\psi(x_b)$ in the last
row; it vanishes when two columns coincide. Since $V^*$ separates the points of $V$, the right side
vanishes whenever $x_a=x_b$ for some $a\ne b$, so it factors through $\Lambda^kV$. Each term is a multiple of some $x_i$, which gives
$c_\varphi(\Lambda^kW)\subseteq W$. If $\omega\in\Lambda^kW$ with $\dim W=\sdim(\omega)$, all $c_\varphi(\omega)$ lie in $W$.
\end{proof}

\begin{proposition}[canonical exterior support]\label{prop:canonical-support}
Let $L$ be any field, $V$ an $L$\nobreakdash-vector space, $k\ge2$, and
$\omega\in\Lambda^kV$. Then
\[
S(\omega):=\spn\{c_\varphi(\omega):\varphi\in(V^*)^{k-1}\}
\]
is the unique smallest subspace supporting $\omega$. In particular,
$\omega\in\Lambda^k S(\omega)$ and
$\sdim(\omega)=\dim S(\omega)$.
\end{proposition}

\begin{proof}
Choose a finite-dimensional supporting subspace $U\subseteq V$.
Lemma~\ref{lem:contr} gives $S(\omega)\subseteq U$, so $S(\omega)$ is
finite-dimensional. Choose a basis $e_1,\dots,e_d$ of $U$ whose first
$r$ vectors form a basis of $S(\omega)$, and extend its coordinate
functionals algebraically to $V$. Write
$\omega=\sum_{|I|=k}a_Ie_I$ in the exterior basis of $U$.
If $j\in I$ and $j>r$, contract against the $k-1$ coordinate functionals
indexed by $I\setminus\{j\}$. The $e_j$\nobreakdash-coordinate of this
contraction is $\pm a_I$: no other basis wedge contributes to that
coordinate. The contraction belongs to $S(\omega)$, so this coordinate
is zero. Hence $a_I=0$ whenever $I$ meets $\{r+1,\dots,d\}$, proving
$\omega\in\Lambda^k S(\omega)$.
Every supporting subspace contains all the contractions by
Lemma~\ref{lem:contr}, and therefore contains $S(\omega)$. This proves
minimality and the dimension formula. Only unit signs occur, so the
argument is valid in every characteristic.
\end{proof}

The same conclusion holds in degree one: under $\Lambda^1V\simeq V$,
the empty-tuple contraction is the identity and $S(\omega)=L\omega$.
In degree zero, every element of $\Lambda^0V=L$ is supported by the
zero subspace, so its support dimension is zero; no contraction to
$V$ is needed.

For matrix calculations, choose any finite-dimensional supporting
subspace $U$ and a basis of it. Form the finite matrix whose columns
are the coordinates of the contractions against $(k-1)$\nobreakdash-tuples
of coordinate functionals, extended to $V$. Arbitrary functionals
restrict to linear combinations of these coordinates on $U$;
multilinearity in the covectors therefore shows that the column span
is exactly $S(\omega)$. The matrix consequently has rank
$\sdim(\omega)$. This uses ordinary matrix rank, and entails no
choice of a rank convention for an infinite matrix.

\subsection{Lifts and three diagonal identities}\label{ssec:diag}
Fix a field $L$, an integer $k\ge1$, $L$\nobreakdash-vector spaces $A$ and $V$, and a linear map $\mu\colon A\to\End_L(V)$,
$a\mapsto \mu_a$. A \emph{lift of $\mu$ in degree $k$} is a $2k$\nobreakdash-linear map $\Psi\colon A^k\times V^k\to\Lambda^kV$ over $L$ that is
alternating in the $V$\nobreakdash-slots; its diagonal identity will be specified separately. For every $n\ge1$ write
$[n]:=\{1,\dots,n\}$, and for $\alpha\in\N^m$ put $\abs\alpha:=\sum_i\alpha_i$. Given $m\ge1$, a map
$f\colon[k]\to[m]$, a family $b=(b_1,\dots,b_m)\in A^m$ and $\lambda\in L^m$, put
\[
b_f:=(b_{f(1)},\dots,b_{f(k)}),\quad \type f:=\bigl(\abs{f^{-1}(1)},\dots,\abs{f^{-1}(m)}\bigr)\in\N^m,\quad
\lambda_f:=\prod_{j=1}^k\lambda_{f(j)}=\lambda^{\type f}.
\]
The \emph{pointwise identity} is
\begin{equation*}\tag{Pw}\label{eq:Pw}
\Psi(a,\dots,a;x_1,\dots,x_k)=\mu_ax_1\wedge\dots\wedge \mu_ax_k\qquad\text{for all }a\in A,\ x\in V^k.
\end{equation*}
The \emph{polarized identity} is: for all $m\ge1$, $b\in A^m$, $\alpha\in\N^m$ with $\abs\alpha=k$, and $x\in V^k$,
\begin{equation*}\tag{Pol}\label{eq:Pol}
\sum_{\type f=\alpha}\Psi(b_f;x)=\sum_{\type f=\alpha}\mu_{b_{f(1)}}x_1\wedge\dots\wedge \mu_{b_{f(k)}}x_k .
\end{equation*}
Its \emph{multilinear instance} is: for all $b_1,\dots,b_k\in A$ and $x\in V^k$,
\begin{equation*}\tag{Pol1}\label{eq:Pol1}
\sum_{\sigma\in S_k}\Psi(b_{\sigma(1)},\dots,b_{\sigma(k)};x)
=\sum_{\sigma\in S_k}\mu_{b_{\sigma(1)}}x_1\wedge\dots\wedge \mu_{b_{\sigma(k)}}x_k .
\end{equation*}
Here $\sum_{\type f=\alpha}$ runs over the maps $f\colon[k]\to[m]$ of type $\alpha$. \eqref{eq:Pol1} is the instance $m=k$,
$\alpha=(1,\dots,1)$ of \eqref{eq:Pol}, since the maps of type $(1,\dots,1)$ are the permutations. For $b\in A^m$ and
$x\in V^k$ put $a_\lambda:=\sum_i\lambda_ib_i$ and
\begin{align*}
\Delta_{b,x}(\lambda)&:=\Psi(a_\lambda,\dots,a_\lambda;x)-\mu_{a_\lambda}x_1\wedge\dots\wedge \mu_{a_\lambda}x_k
=\sum_{\abs\alpha=k}\lambda^\alpha\,\delta_\alpha(b,x),\\
\delta_\alpha(b,x)&:=\sum_{\type f=\alpha}\Bigl[\Psi(b_f;x)-\mu_{b_{f(1)}}x_1\wedge\dots\wedge \mu_{b_{f(k)}}x_k\Bigr].
\end{align*}
Multilinearity of $\Psi$ in the $A$\nobreakdash-slots, linearity of $\mu$ and multilinearity of the wedge give the displayed
expansion. Thus $\Delta_{b,x}$ is a homogeneous polynomial of degree $k$ with coefficients in $\Lambda^kV$.
Identity~\eqref{eq:Pol} says that every coefficient $\delta_\alpha(b,x)$ is zero. Identity~\eqref{eq:Pw} says instead that
$\Delta_{b,x}$ vanishes at every point of $L^m$; conversely, this pointwise vanishing gives~\eqref{eq:Pw} by taking
$m=1$ and $\lambda=1$. The following two identity principles identify the cases in which pointwise vanishing forces
all coefficients to vanish.

\begin{lemma}[identity principle]\label{lem:idprinc}
Let $\mathbb K$ be an infinite field, $Y$ a $\mathbb K$\nobreakdash-vector space and $g(\lambda)=\sum_\alpha\lambda^\alpha y_\alpha$ a
polynomial in $\lambda\in\mathbb K^m$ with finitely many nonzero coefficients $y_\alpha\in Y$. If $g(\lambda)=0$ for all
$\lambda\in\mathbb K^m$, then every $y_\alpha=0$.
\end{lemma}

\begin{proof}
The finitely many coefficients $y_\alpha$ span a finite-dimensional subspace of $Y$. Applying its coordinate
functionals reduces the claim to scalar polynomials. For these, induction on $m$ reduces to the fact that a nonzero
polynomial in one variable has only finitely many roots.
\end{proof}

\begin{lemma}[homogeneous forms over a finite field]\label{lem:forms}
Let $q$ be a prime power, $L'\supseteq\F_q$ a field, $1\le e\le q$, and $g\in L'[\lambda_1,\dots,\lambda_m]$ homogeneous of
degree $e$ with $g(\lambda)=0$ for every $\lambda\in\F_q^m$. Then $g=0$.
\end{lemma}

\begin{proof}
Call a monomial $\lambda^\beta$ \emph{reduced} if every $\beta_i\le q-1$. The indicator function of a point
$a\in\F_q^m$ is $\prod_i\bigl(1-(\lambda_i-a_i)^{q-1}\bigr)$, which expands into reduced monomials. So the $q^m$ reduced
monomials span the $q^m$\nobreakdash-dimensional $L'$\nobreakdash-space of functions $\F_q^m\to L'$, and they are linearly independent as
functions. A monomial of degree $e\le q$ that is not reduced is $\lambda_i^q$ (and then $e=q\ge2$); as a function on
$\F_q^m$ it agrees with the reduced monomial $\lambda_i$, of degree $1<e$. If $e=q$, replace each $\lambda_i^q$ in $g$ by
$\lambda_i$. These degree-one images are distinct, and homogeneity separates them from all the remaining monomials,
which still have degree $q$. The replacement therefore sends the monomials of $g$ injectively to reduced monomials
without changing the associated function. Linear independence of the reduced monomials forces every coefficient to
vanish. If $e<q$, no replacement is needed.
\end{proof}

\begin{proposition}[pointwise versus polarized]\label{prop:diag}
Let $\Psi$ be a lift of $\mu$ in degree $k$ over a field $L$.
\begin{enumerate}
\item \eqref{eq:Pol}$\Rightarrow$\eqref{eq:Pw}$\Rightarrow$\eqref{eq:Pol1}, over every field.
\item If $L$ is infinite, or if $L=\F_q$ and $k\le q$, then \eqref{eq:Pw}$\Rightarrow$\eqref{eq:Pol}.
\item If $L=\F_q$ and $k\ge q+1$, then \eqref{eq:Pw} does not imply \eqref{eq:Pol} in general.
\item If $k!=0$ in $L$, then \eqref{eq:Pol1} does not imply \eqref{eq:Pw} in general.
\end{enumerate}
\end{proposition}

\begin{proof}
\emph{(1) Inclusion--exclusion.} The implication \eqref{eq:Pol}$\Rightarrow$\eqref{eq:Pw} is the case $m=1$ and
$\alpha=(k)$. To prove \eqref{eq:Pw}$\Rightarrow$\eqref{eq:Pol1}, fix $b_1,\dots,b_k\in A$ and $x\in V^k$, and put
$b_S:=\sum_{i\in S}b_i$ for $S\subseteq[k]$. Multilinearity gives
\begin{align*}
\sum_{S\subseteq[k]}(-1)^{k-\abs S}\Psi(b_S,\dots,b_S;x)&=\sum_S(-1)^{k-\abs S}\sum_{f\colon[k]\to S}\Psi(b_f;x)\\
&=\sum_{f\colon[k]\to[k]}\Psi(b_f;x)\sum_{S\supseteq f([k])}(-1)^{k-\abs S}=\sum_{f\text{ bijective}}\Psi(b_f;x),
\end{align*}
because $\sum_{S\supseteq T}(-1)^{k-\abs S}=(1-1)^{k-\abs T}$ is $1$ if $T=[k]$ and $0$ otherwise. The same computation
applies to the multilinear map $(a_1,\dots,a_k)\mapsto \mu_{a_1}x_1\wedge\dots\wedge \mu_{a_k}x_k$, whose diagonal is the right
side of \eqref{eq:Pw}. Evaluating \eqref{eq:Pw} at the $2^k$ subset sums $b_S$ therefore gives \eqref{eq:Pol1}.
The computation is valid in every characteristic and involves no division by $k!$.

\emph{(2) The identity principles.} The coefficients of $\Delta_{b,x}$ span a finite-dimensional subspace of $\Lambda^kV$; in a basis of it each
coordinate of $\Delta_{b,x}$ is a scalar form of degree $k$ vanishing on $L^m$. It is zero by Lemma~\ref{lem:idprinc}
if $L$ is infinite, and by Lemma~\ref{lem:forms} if $L=\F_q$ and $k\le q$.

\emph{(3) Pointwise vanishing without full polarization.} Let $\dim A\ge3$, with $b_1,b_2,b_3\in A$ and functionals $g_1,g_2,g_3$ on $A$ such that
$g_i(b_j)=\delta_{ij}$, and let $y_1,\dots,y_k\in V$ be linearly independent. Put
\[
\theta(a_1,\dots,a_k):=\bigl[g_1(a_1)\cdots g_1(a_q)\,g_2(a_{q+1})-g_1(a_1)\,g_2(a_2)\cdots g_2(a_{q+1})\bigr]
\cdot g_3(a_{q+2})\cdots g_3(a_k)
\]
(for $k=q+1$ the $g_3$\nobreakdash-factor is empty and $\dim A\ge2$ suffices), and $\Theta(a;x):=\theta(a)\,x_1\wedge\dots\wedge x_k$.
Then $\Theta$ is a lift of $\mu=0$. It satisfies \eqref{eq:Pw}, because
$\theta(a,\dots,a)=(g_1(a)^qg_2(a)-g_1(a)g_2(a)^q)\,g_3(a)^{k-q-1}=0$ since $c^q=c$ for $c\in\F_q$. It
violates \eqref{eq:Pol} for $\alpha=(q,1,k-q-1)$ and $b=(b_1,b_2,b_3)$: the first product contributes to
$\sum_{\type f=\alpha}\theta(b_f)$ exactly for $f=(1,\dots,1,2,3,\dots,3)$, with value $1$; the second product would need
$f(2)=\dots=f(q+1)=2$, that is, at least $q\ge2$ values equal to $2$, which is impossible for type $\alpha$. So the
coefficient $\delta_\alpha(b,y)$ of $\Theta$ is $y_1\wedge\dots\wedge y_k\neq0$. When $k=q+1$, the third label and its
functional can be omitted throughout, with $\alpha=(q,1)$.

\emph{(4) The multilinear identity without the pointwise identity.} Take $\bar a\in A$ and a functional $g_1$ with $g_1(\bar a)=1$, and $y_1,\dots,y_k$ linearly independent. The
lift $\Theta'(a;x):=g_1(a_1)\cdots g_1(a_k)\,x_1\wedge\dots\wedge x_k$ of $\mu=0$ satisfies \eqref{eq:Pol1}, since
$\sum_\sigma\Theta'(b_\sigma;x)=k!\,g_1(b_1)\cdots g_1(b_k)\,x_1\wedge\dots\wedge x_k=0$ for all
$b_1,\dots,b_k\in A$ and $x\in V^k$, but $\Theta'(\bar a,\dots,\bar a;y)=y_1\wedge\dots\wedge y_k\neq0$.
\end{proof}

The finite-field obstruction below requires only~\eqref{eq:Pol1}, which follows from either~\eqref{eq:Pw}
or~\eqref{eq:Pol}. The distinction between the latter two remains important: over $\F_q$ they need not agree when
$k>q$. The lift constructed in Appendix~\ref{sec:laurent} nevertheless satisfies the full identity~\eqref{eq:Pol} in every
degree: the identity is polarized over an infinite Laurent-series field before constant coefficients are taken
(Lemma~\ref{lem:Psi}).

\subsection{Statement and proof strategy}
For the proof of the theorem, assume that $L$ is a \emph{finite} field with $k!=0$ in $L$, that is, $\operatorname{char}L=p\le k$; in
particular $k\ge2$. Put $V:=L^\N$ and let $V_{\mathrm{fin}}:=L^{(\N)}\subseteq V$ be the subspace of finitely supported
sequences. For $u,x\in V$ we write $ux$ for the coordinatewise product; so $\mu_ux=ux$ is the multiplier family.

\Needspace{16\baselineskip}
\begin{samepage}
\begin{theorem}[multipliers over a finite field]\label{thm:FF}
Let $L$ be a finite field and $k\ge1$ with $k!=0$ in $L$.
\begin{enumerate}
\item There is no $2k$\nobreakdash-linear map $\Psi\colon V_{\mathrm{fin}}^k\times V_{\mathrm{fin}}^k\to\Lambda^k_LV$ over $L$ such that
\begin{enumerate}
\item $\Psi$ is antisymmetric in the last $k$ slots;
\item $\Psi$ satisfies \eqref{eq:Pol1} for $\mu_ux=ux$: for all $u_1,\dots,u_k,x_1,\dots,x_k\in V_{\mathrm{fin}}$,
\[
\sum_{\sigma\in S_k}\Psi(u_{\sigma(1)},\dots,u_{\sigma(k)};x_1,\dots,x_k)=\sum_{\sigma\in S_k}(u_{\sigma(1)}x_1)\wedge\dots\wedge(u_{\sigma(k)}x_k);
\]
\item for some $d\in\N$, $\sdim(\Psi(u;x))\le d$ for all $u,x\in V_{\mathrm{fin}}^k$ with entries in $\{0,\pm1\}$.
\end{enumerate}
\item Let $E_0=V$ or $E_0=V_{\mathrm{fin}}$. There is no $2k$\nobreakdash-linear map $\Psi\colon E_0^k\times E_0^k\to\Lambda^k_LE_0$ over $L$
that is antisymmetric in the last $k$ slots, satisfies \eqref{eq:Pol1} for all finitely supported $u,x$, and has
$\sdim(\Psi(u;x))\le d$ for all $u,x\in E_0^k$, for some $d$.
\end{enumerate}
\end{theorem}
\end{samepage}

\emph{Reduction to (1).} Given a map as in (2), restrict its inputs to
$V_{\mathrm{fin}}^k\times V_{\mathrm{fin}}^k$ and compose its output with the injective map
$\Lambda^kE_0\to\Lambda^kV$. This preserves antisymmetry and~\eqref{eq:Pol1}. Support dimension can only decrease under
the inclusion of the ambient space (\S\ref{ssec:ext}), so the resulting map satisfies (a)--(c).

We therefore prove (1). Suppose that $\Psi$ satisfies its hypotheses, and fix the bound $d$ in (c). Write
$e_i\in V_{\mathrm{fin}}$ for the unit vectors, $e_a:=(e_{a_1},\dots,e_{a_k})$ for $a\in\N^k$, and $\Omega$ for the
determinant array of \S\ref{ssec:ext} with $S=\N$.

The proof has three steps. Ramsey's theorem first makes the coefficients of $\Omega\Psi$ depend only on the relative
order and equalities among their indices. Special weights on four-point clusters then cancel every contribution
from an occupied cluster containing no output index. This cancellation produces a two-valued staircase matrix from
one exterior vector. Its bounded support forces the two values to agree, and hence makes all cluster order patterns
have a common value. Finally, antisymmetry and~\eqref{eq:Pol1} say that the sum of these $k!$ equal values is $1$,
contradicting $k!=0$ in $L$.

\subsection{Order-homogeneity from Ramsey's theorem}
We use Ramsey's theorem~\cite{Ra} in its infinite form: \emph{for every $n\ge1$ and every coloring of the
$n$\nobreakdash-element subsets of $\N$ with finitely many colors there is an infinite $H\subseteq\N$ all of whose $n$\nobreakdash-element
subsets have the same color.}

For a tuple $z\in\N^N$ with distinct values $y_1<\dots<y_r$, its \emph{pattern} is the map
$i\mapsto j$ determined by $z_i=y_j$. In particular, repeated entries are retained in the pattern. Two tuples $z$
and $z'$ have the same pattern if and only if $z_i<z_j\Leftrightarrow z'_i<z'_j$ and $z_i=z_j\Leftrightarrow z'_i=z'_j$ for
all $i,j$.

\begin{lemma}[pattern homogeneity]\label{lem:H}
For every $N\ge1$ and every function $T\colon\N^N\to L$ there is an infinite $H\subseteq\N$ such that $T(z)=T(z')$
whenever $z,z'\in H^N$ have the same pattern.
\end{lemma}

\begin{proof}
Color an $N$\nobreakdash-element set $s=\{s_1<\dots<s_N\}\subseteq\N$ by the vector
\[
\bigl(T(s_{\beta(1)},\dots,s_{\beta(N)})\bigr)_{\beta\colon[N]\to[N]}\in L^{N^N}.
\]
The coloring records all maps $\beta$, including those that repeat indices. It has finitely many values because
$L$ is finite; this is the only use of the finiteness of $L$ in the obstruction proof. Ramsey's theorem gives an
infinite $H$ on which the color is constant, say equal to $(t_\beta)_\beta$. Let $z\in H^N$ have value set
$\{y_1<\dots<y_r\}$. Since $H$ is unbounded, add $N-r$ elements of $H$ above $y_r$ to obtain an
$N$\nobreakdash-element set $s\subseteq H$. Padding above the largest value leaves the ranks unchanged: $s_j=y_j$ for
$j\le r$, and
$z=(s_{\beta(1)},\dots,s_{\beta(N)})$, where $\beta(i)$ is the rank of $z_i$ in $\{y_1,\dots,y_r\}$. The map $\beta$
depends only on the pattern of $z$, and $T(z)=t_\beta$.
\end{proof}

\emph{Coefficients.} For $(a;y;c)\in\N^k\times\N^k\times\N^k$ put
\[
T(a;y;c):=\Omega_{\Psi(e_a;e_y)}(c_1,\dots,c_k)\in L .
\]
Lemma~\ref{lem:H} with $N=3k$ gives an infinite set $H=\{h_0<h_1<h_2<\cdots\}$ on which $T$ depends only on the pattern.
For $u_j,v_j\in V_{\mathrm{fin}}$ supported in $H$, multilinearity of $\Psi$ and linearity of $\Omega$ give
\begin{equation}\label{eq:expand}
\Omega_{\Psi(u;v)}(c)=\sum_{a,y\in H^k}\ \prod_{j=1}^ku_j(a_j)\prod_{j=1}^kv_j(y_j)\cdot T(a;y;c).
\end{equation}
Although the sum is indexed by $H^k\times H^k$, only finitely many summands are nonzero, because all the input
vectors have finite support. Every output tuple $c$ used below also lies in $H^k$, so pattern homogeneity applies to
all $3k$ indices in~\eqref{eq:expand}.

\subsection{Clusters and the cluster value}
A \emph{cluster} is a set $C=\{h_n,h_{n+1},h_{n+2},h_{n+3}\}$ of four
\emph{consecutive} elements of $H$; we write $C(i):=h_{n+i-1}$ for $i=1,\dots,4$. Since a cluster is an interval of
$H$, two disjoint clusters are \emph{comparable}: every element of one is smaller than every element of the other; we
write $C<C'$. Distinct clusters need not be adjacent in $H$.

We choose two weight vectors, one for each kind of input slot:
\begin{align*}
s&=(s_1,s_2,s_3,s_4):=(1,0,-1,0)&&\text{on the operator slots},\\
w&=(w_1,w_2,w_3,w_4):=(1,-1,1,-1)&&\text{on the vector slots}.
\end{align*}
Over $\Z$, hence over every field,
\begin{equation*}\tag{W}\label{eq:W}
\sum_is_i=\sum_iw_i=\sum_is_iw_i=\sum_{i<i'}s_iw_{i'}=\sum_{i>i'}s_iw_{i'}=0,\qquad s_1w_1=1.
\end{equation*}
Indeed $\sum_is_i=1-1$, $\sum_iw_i=1-1+1-1$, $s_1w_1=1$, and
\begin{align*}
\sum_is_iw_i&=s_1w_1+s_3w_3=1-1=0,\\
\sum_{i<i'}s_iw_{i'}&=s_1(w_2+w_3+w_4)+s_3w_4=-1+1=0,\\
\sum_{i>i'}s_iw_{i'}&=s_3(w_1+w_2)=-(1-1)=0 .
\end{align*}
The first two zero sums will cancel a cluster occupied by only one input index. The other three will cancel a
cluster occupied by an operator index and a vector index, according to their relative order. For a cluster $C$ put
\[
s_C:=\sum_{i=1}^4s_ie_{C(i)},\qquad w_C:=\sum_{i=1}^4w_ie_{C(i)} .
\]
These are vectors of $V_{\mathrm{fin}}$ with entries in $\{0,\pm1\}$, and $s_Cw_C=e_{C(1)}-e_{C(3)}$.

\emph{The cluster value.} For pairwise disjoint clusters $C_1,\dots,C_k$, one for each slot, define
\begin{align*}
\chi(C_1,\dots,C_k)&:=\Omega_{\Psi(s_{C_1},\dots,s_{C_k};\,w_{C_1},\dots,w_{C_k})}\bigl(C_1(1),\dots,C_k(1)\bigr)\\
&=\sum_{i,i'\in[4]^k}\ \prod_{j=1}^ks_{i_j}w_{i'_j}\cdot
T\bigl(C_1(i_1),\dots,C_k(i_k);\,C_1(i'_1),\dots,C_k(i'_k);\,C_1(1),\dots,C_k(1)\bigr).
\end{align*}
Let $\tau\in S_k$ be the \emph{order pattern} of $(C_1,\dots,C_k)$: $\tau(j)$ is the position of $C_j$ among
$C_1,\dots,C_k$. The pattern of every $3k$\nobreakdash-tuple in the sum is determined by $(i,i')$ and $\tau$: entries in different
clusters compare as the clusters do, and the entries $C_j(i_j)$, $C_j(i'_j)$, $C_j(1)$ in the same cluster compare as
their indices do, including equalities. Pattern homogeneity therefore shows that the displayed sum depends only
on $\tau$. Every permutation occurs as the order pattern of some disjoint clusters in $H$; we denote the resulting
scalar by $\chi(\tau)$.

\begin{lemma}[cancellation of output-free clusters]\label{lem:C}
Let $\cC_1,\dots,\cC_k$ be finite families of clusters that are pairwise disjoint as sets of clusters, such that all
clusters in $\bigcup_j\cC_j$ are pairwise disjoint. Put $u_j:=\sum_{C\in\cC_j}s_C$ and $v_j:=\sum_{C\in\cC_j}w_C$. Let
$C_l\in\cC_l$ for $l=1,\dots,k$ and $c:=(C_1(1),\dots,C_k(1))$. Then $\Omega_{\Psi(u;v)}(c)=\chi(C_1,\dots,C_k)$.
\end{lemma}

\begin{proof}
We show that every cluster assignment except the one matching the output clusters has total contribution zero.
Expand by~\eqref{eq:expand}. Since the clusters are disjoint, a nonzero term has, for each slot $j$, an operator index
$a_j=\Cop_j(i_j)$ and a vector index $y_j=\Cvec_j(i'_j)$ with $\Cop_j,\Cvec_j\in\cC_j$, and weight
$\prod_js_{i_j}w_{i'_j}$. Group the terms according to the \emph{cluster assignment} $(\Cop_j,\Cvec_j)_{j}$.

Call an occupied cluster \emph{output-free} if it is none of $C_1,\dots,C_k$. Suppose a group has such a cluster
$\Gamma$. Disjointness of the families assigns $\Gamma$ to exactly one slot $j$, so the only tuple entries it can
contain are $a_j$ and $y_j$. Disjointness of the clusters excludes every output index from $\Gamma$.

Fix the positions of all indices outside $\Gamma$ and sum over the positions inside it. Since $\Gamma$ consists of
consecutive elements of $H$, every outside entry lies either below the whole cluster or above it. Moving an index
within $\Gamma$ cannot change any comparison with an outside entry. By pattern homogeneity, the coefficient $T$
therefore depends on the internal positions only as follows:
\begin{enumerate}
\item[(i)] if only $a_j\in\Gamma$, not at all, and the partial sum is $(\sum_is_i)\cdot T=0$;
\item[(ii)] if only $y_j\in\Gamma$, not at all, and the partial sum is $(\sum_iw_i)\cdot T=0$;
\item[(iii)] if both lie in $\Gamma$, only through whether $i_j<i'_j$, $i_j=i'_j$ or $i_j>i'_j$; with the corresponding
values $T_<,T_=,T_>$ of $T$ the partial sum is
\[
T_<\sum_{i<i'}s_iw_{i'}+T_=\sum_is_iw_i+T_>\sum_{i>i'}s_iw_{i'}=0 .
\]
\end{enumerate}
In each case the weights of the other slots are a common factor, and the displayed partial sum vanishes
by~\eqref{eq:W}. Thus every assignment with an output-free cluster contributes zero. In any surviving assignment,
every occupied cluster must be an output cluster. The only output cluster belonging to $\cC_j$ is $C_j$, which forces
$\Cop_j=\Cvec_j=C_j$ for every $j$. The contribution of this one assignment is exactly
$\chi(C_1,\dots,C_k)$.
\end{proof}

\subsection{The staircase and the diagonal}

We now use the support bound to compare two order patterns that differ by an adjacent exchange. The key is to
realize arbitrarily large staircase matrices as flattenings of a single value of $\Psi$.

\begin{lemma}[staircase]\label{lem:S}
Let $\tau\in S_k$, let $\alpha\neq\beta$ be slots with $\tau(\alpha)=\ell$ and $\tau(\beta)=\ell+1$, and let $\tau'$ be $\tau$
with these two positions exchanged, so $\tau'(\alpha)=\ell+1$, $\tau'(\beta)=\ell$ and $\tau'(\gamma)=\tau(\gamma)$ otherwise.
Then $\chi(\tau)=\chi(\tau')$.
\end{lemma}

\begin{proof}
Let $m\ge1$. Choose pairwise disjoint clusters, in increasing order in $H$: first a cluster $C_\gamma$ for each slot
$\gamma$ with $\tau(\gamma)<\ell$, ordered as $\tau$ prescribes; then $2m$ clusters $A_1<B_1<A_2<B_2<\dots<A_m<B_m$; then a
cluster $C_\gamma$ for each slot $\gamma$ with $\tau(\gamma)>\ell+1$, ordered as $\tau$ prescribes. Put
$\cC_\alpha:=\{A_1,\dots,A_m\}$, $\cC_\beta:=\{B_1,\dots,B_m\}$ and $\cC_\gamma:=\{C_\gamma\}$ for the other slots, and let
$u,v$ be as in Lemma~\ref{lem:C} and $\omega:=\Psi(u;v)$. Disjointness of the clusters ensures that the entries of
$u$ and $v$ still lie in $\{0,\pm1\}$, so $\sdim(\omega)\le d$ by hypothesis~(c) of Theorem~\ref{thm:FF}(1).
For this fixed $m$, keep $u$, $v$ and hence $\omega$ fixed while varying only the output indices.

Consider the $m\times m$ matrix $M$ with rows indexed by $c_\alpha=A_n(1)$ and columns by $c_\beta=B_{n'}(1)$
($n,n'=1,\dots,m$), all other output indices being fixed at $c_\gamma=C_\gamma(1)$, and entries $M_{nn'}:=\Omega_\omega(c)$.
By Lemma~\ref{lem:C}, applied with the output clusters $A_n\in\cC_\alpha$, $B_{n'}\in\cC_\beta$ and $C_\gamma\in\cC_\gamma$,
$M_{nn'}$ is the cluster value of the tuple with $A_n$ in slot $\alpha$, $B_{n'}$ in slot $\beta$, and $C_\gamma$ in
every other slot. The interlacing order gives $A_n<B_{n'}$ exactly when $n\le n'$. The tuple therefore has pattern
$\tau$ on and above the matrix diagonal, and pattern $\tau'$ below it.

Put $\delta:=\chi(\tau)-\chi(\tau')$, let $\mathbf 1\in L^m$ be the all-ones column, and let $U$ be the upper
unitriangular matrix with $U_{nn'}=1$ for $n\le n'$ and $U_{nn'}=0$ otherwise. Then
\[
M=\begin{pmatrix}
\chi(\tau)&\chi(\tau)&\cdots&\chi(\tau)\\
\chi(\tau')&\chi(\tau)&\cdots&\chi(\tau)\\
\vdots&\ddots&\ddots&\vdots\\
\chi(\tau')&\cdots&\chi(\tau')&\chi(\tau)
\end{pmatrix}
=\chi(\tau')\,\mathbf 1\mathbf 1^{\mathsf T}+\delta\,U,
\]
The rows vary only slot $\alpha$; each column fixes all the other slots, including the chosen value in slot $\beta$.
Thus $M$ is a submatrix of the flattening of this one array $\Omega_\omega$ along slot $\alpha$. Lemma~\ref{lem:flat}
gives $\rank M\le\sdim(\omega)\le d$.

If $\delta\ne0$, the matrix $\delta U$ is invertible, whereas the constant matrix
$\chi(\tau')\mathbf 1\mathbf 1^{\mathsf T}$ has rank at most one. Applying rank subadditivity to
$\delta U=M-\chi(\tau')\mathbf 1\mathbf 1^{\mathsf T}$ gives
\[
m=\rank(\delta U)\le\rank M+1\le d+1.
\]
The construction is available for every $m\ge1$; choosing $m=d+2$ is a contradiction. Hence $\delta=0$, as required.
\end{proof}

The support bound has now forced equality between adjacent order patterns. The remaining identity comes from
\eqref{eq:Pol1}; it does not use the support bound.

\begin{lemma}[the diagonal]\label{lem:D}
$\sum_{\tau\in S_k}\chi(\tau)=1$.
\end{lemma}

\begin{proof}
Take clusters $C_1<\dots<C_k$, and put $u_j:=s_{C_j}$, $v_j:=w_{C_j}$ and
$c:=(C_1(1),\dots,C_k(1))$. Apply $\omega\mapsto\Omega_\omega(c)$ to both sides of~\eqref{eq:Pol1}, with the operator
vectors $u_1,\dots,u_k$ and the vector arguments $v_1,\dots,v_k$.

\emph{Right side.} If $\sigma\ne\id$, some slot $j$ has $\sigma(j)\ne j$. Disjoint supports then give
$u_{\sigma(j)}v_j=0$, so the corresponding wedge vanishes. Only $\sigma=\id$ survives. It gives $\Omega_{(u_1v_1)\wedge\dots\wedge(u_kv_k)}(c)=\det\bigl((u_bv_b)(C_a(1))\bigr)_{a,b}$. Since
$u_bv_b=e_{C_b(1)}-e_{C_b(3)}$ and $C_a(1)\ne C_b(3)$ for all $a,b$, this is the determinant of the identity matrix, so
the right side is $1$.

\emph{Left side.} Fix $\sigma\in S_k$ and write $u_\sigma:=(u_{\sigma(1)},\dots,u_{\sigma(k)})$ and
$v_\sigma:=(v_{\sigma(1)},\dots,v_{\sigma(k)})$. To identify the summand with a cluster value, we must permute both the
vector arguments and the output coordinates to match the operator arguments. Each permutation contributes the same
sign. More precisely, antisymmetry of $\Psi$ in the vector slots gives
$\Psi(u_\sigma;v)=\sgn(\sigma)\Psi(u_\sigma;v_\sigma)$. For every $\omega\in\Lambda^kV$ the function $\Omega_\omega$ is
alternating in its arguments, so $\Omega_\omega(c)=\sgn(\sigma)\,\Omega_\omega(c_{\sigma(1)},\dots,c_{\sigma(k)})$; we
apply this with $\omega=\Psi(u_\sigma;v_\sigma)$. Since $\sgn(\sigma)^2=1$,
\begin{align*}
\Omega_{\Psi(u_\sigma;\,v)}(c)&=\sgn(\sigma)\,\Omega_{\Psi(u_\sigma;\,v_\sigma)}(c)
=\sgn(\sigma)^2\,\Omega_{\Psi(u_\sigma;\,v_\sigma)}(c_{\sigma(1)},\dots,c_{\sigma(k)})\\
&=\Omega_{\Psi(u_\sigma;\,v_\sigma)}(c_{\sigma(1)},\dots,c_{\sigma(k)}).
\end{align*}
In the last expression slot $j$ has operator $s_{C_{\sigma(j)}}$, vector $w_{C_{\sigma(j)}}$ and output
$C_{\sigma(j)}(1)$, so it equals $\chi(C_{\sigma(1)},\dots,C_{\sigma(k)})$ by definition. Its order pattern is $\sigma$,
since $C_{\sigma(j)}$ is the $\sigma(j)$\nobreakdash-th cluster in increasing order. Thus the convention for patterns
gives $\sigma$, not $\sigma^{-1}$. No cluster cancellation is needed in this calculation.

As $\sigma$ runs over $S_k$, each order pattern occurs exactly once. Comparing the two sides gives
\[
\sum_{\sigma\in S_k}\Omega_{\Psi(u_\sigma;\,v)}(c)=\sum_{\tau\in S_k}\chi(\tau)=1 .\qedhere
\]
\end{proof}

\begin{proof}[Proof of Theorem~\ref{thm:FF}]
Exchanging adjacent ranks acts on an order pattern by
$\tau\mapsto(\ell\ \ell{+}1)\circ\tau$, $1\le\ell<k$. These transpositions generate $S_k$, and Lemma~\ref{lem:S}
makes $\chi$ invariant under each exchange. Hence $\chi$ is constant; write its value as $\chi_0$.
Lemma~\ref{lem:D} now gives $k!\,\chi_0=1$ in $L$, contradicting $k!=0$. This proves (1), and the preceding reduction
gives (2).
\end{proof}

\begin{remark}[where the hypotheses enter]\label{rem:FFhyp}
Ramsey homogeneity is the only step in the obstruction proof that uses the finiteness of $L$.
Multilinearity then gives the cluster expansion and cancellation. The support bound is used in Lemma~\ref{lem:S}
only at the finitely many test inputs obtained by taking $m=d+2$; all have entries in $\{0,\pm1\}$.
Antisymmetry and~\eqref{eq:Pol1} enter in Lemma~\ref{lem:D}, which does not use that bound. Finally, $k!=0$ is used only
to turn the diagonal identity into a contradiction.

The inclusion--exclusion argument in Proposition~\ref{prop:diag}(1) uses only multilinearity in the
operator slots. The theorem therefore remains true with the pointwise identity \eqref{eq:Pw} at $L$\nobreakdash-points in place of
\eqref{eq:Pol1}, for this family of multipliers. In characteristic $2$ antisymmetry is the same as symmetry, and the two signs in
Lemma~\ref{lem:D} are trivial.
\end{remark}

\begin{remark}[when $k!$ is invertible]\label{rem:FFunit}
The vanishing of $k!$ is essential. If instead $k!\ne0$ in $L$, the \emph{normalized alternatization}
\[
\Psi_{\mathrm{alt}}(u;x):=\frac1{k!}\sum_{\sigma\in S_k}(u_{\sigma(1)}x_1)\wedge\dots\wedge(u_{\sigma(k)}x_k)
\]
is alternating in $x$ and satisfies~\eqref{eq:Pol1}. Every summand is supported in the span of the $k^2$ vectors
$u_ix_j$, so $\sdim(\Psi_{\mathrm{alt}}(u;x))\le k^2$. The preceding homogeneity, cancellation, staircase and diagonal
arguments apply, but now yield $\chi\equiv1/k!$: for pairwise disjoint clusters only $\sigma=\id$
survives in $\Psi_{\mathrm{alt}}(s_{C_1},\dots,s_{C_k};w_{C_1},\dots,w_{C_k})$, and the determinant at
$(C_1(1),\dots,C_k(1))$ is that of the identity. The final equation $k!\cdot(1/k!)=1$ is then consistent. So the
obstruction is exactly the vanishing of $k!$, not positive characteristic as such. For example, at $(p,k)=(5,3)$ every
step is consistent, since $3!=6=1$ in $\F_5$ and $\chi_0=1$.
\end{remark}
%% ---- end app-B-multiplier.tex
%% ---- begin app-C-laurent.tex
\section{The Banach exterior target over a Laurent series field}\label{app:laurent}\label{sec:laurent}

We now realize the finite-field obstruction of Theorem~\ref{thm:FF} in a Banach-space setting.
The space $F_1$ of Theorem~\ref{thm:laurent} is $\widehat\Lambda^k_\pi E_1$, the test target of
Lemma~\ref{fib:lem:test} for the pair $(E_1,E_1)$, so Theorem~\ref{thm:laurent} says that $(E_1,E_1)$ is not good for
Banach values in degree $k$ (Corollary~\ref{obs:cor:linf}).

\begin{theorem}\label{thm:laurent}
Let $\kappa$ be a finite field of characteristic $p$, let $k\ge p$ and $\varrho\in(0,1)$, and let $K_1=\kappa((X))$ with
$\abs X=\varrho$. Let $E_1=\linf(\N,K_1)$, and let $F_1$ be the completion of $\Lambda^k_{K_1}E_1$ under the projective exterior
norm. Then $A^{k,K_1}_{E_1,E_1;F_1}$ admits no bounded $k$\nobreakdash-linear lift over $K_1$. Consequently it is analytic at no point of
$\cL(E_1,E_1)$.
\end{theorem}

\subsection{Setting}\label{ssec:laurentsetting}
Throughout this section $k\ge1$, $\kappa$ is a field, $\varrho\in(0,1)$, and $K_1:=\kappa((X))$ is the Laurent series field with
absolute value $\abs a:=\varrho^{\ord a}$ ($\abs0=0$). This field is complete, nontrivially normed and nonarchimedean, and
$\abs c=1$ for every $c\in\kappa^\times$. Write $\coeff_n(a)\in\kappa$ for the $X^n$\nobreakdash-coefficient of $a\in K_1$.
For a set $S$ and $y\in\linf(S,K_1)$ we take the constant coefficient coordinatewise:
\[
\coeff_0(y):=\bigl(\coeff_0(y(s))\bigr)_{s\in S}\in\kappa^S.
\]
The resulting \emph{constant-coefficient map} $\coeff_0$ is $\kappa$\nobreakdash-linear but not $K_1$\nobreakdash-linear.
A nonzero Laurent series of absolute value less than one has order at least one. Consequently,
\begin{equation}\label{eq:small}
\norm y_\infty<1\ \Longrightarrow\ \coeff_0(y)=0,\qquad\text{hence}\qquad \norm{y-y'}_\infty<1\ \Longrightarrow\ \coeff_0(y)=\coeff_0(y').
\end{equation}
Put $E_1:=\linf(\N,K_1)$ with the sup norm, a nonarchimedean $K_1$\nobreakdash-Banach space, and let $E_0:=\kappa^\N\subseteq E_1$ be
the $\kappa$\nobreakdash-subspace of sequences with entries in $\kappa$. Every nonzero $z\in E_0$ has $\norm z_\infty=1$. For
$x\in E_1$ and $n\in\Z$ let $x_{[n]}:=(\coeff_n(x(i)))_{i\in\N}\in E_0$ be the $n$\nobreakdash-th Laurent component. For $x\ne0$
the set $\{\ord x(i):x(i)\ne0\}$ is bounded below, because $\varrho^{\ord x(i)}\le\norm x_\infty$, so it has a minimum
$\nu(x)$; then $\norm x_\infty=\varrho^{\nu(x)}$ and $x_{[n]}=0$ for $n<\nu(x)$. Finally put
\[
M_\varrho:=\max_{l\in\N}\,(l+1)\varrho^l .
\]
The maximum exists because $(l+1)\varrho^l\to0$. This constant will control the number of Laurent components that can
contribute to a constant coefficient. For $\varrho\le\frac12$ one has $(l+1)\varrho^l\le(l+1)2^{-l}\le1$, so $M_\varrho=1$.

\subsection{The projective exterior norm}\label{ssec:projnorm}
On $\Lambda:=\Lambda^k_{K_1}E_1$ put
\[
\norm\omega_\pi:=\inf\Bigl\{\sum_j\prod_{i=1}^k\norm{x_{j,i}}_\infty\ :\ \omega=\sum_jx_{j,1}\wedge\dots\wedge x_{j,k}\Bigr\},
\]
where the infimum is over finite decompositions and the cost is an ordinary sum. Let $\Omega^{K_1}\colon\Lambda\to\linf(\N^k,K_1)$ be the determinant
array of \S\ref{ssec:ext} with $L=K_1$, $S=\N$ and $V=E_1$, so that
$\Omega^{K_1}(x_1\wedge\dots\wedge x_k)(i_1,\dots,i_k)=\det(x_b(i_a))_{a,b}$; its values are bounded by the next lemma.

\needspace{9\baselineskip}
\begin{lemma}\label{lem:proj}
$\norm\cdot_\pi$ is a $K_1$\nobreakdash-norm on $\Lambda$ with $\norm{x_1\wedge\dots\wedge x_k}_\pi\le\prod_i\norm{x_i}_\infty$, and
$\Omega^{K_1}$ is $K_1$\nobreakdash-linear and injective with $\norm{\Omega^{K_1}(\omega)}_\infty\le\norm\omega_\pi$. Let $F_1$ be the completion of
$(\Lambda,\norm\cdot_\pi)$, a $K_1$\nobreakdash-Banach space, $J\colon F_1\to\linf(\N^k,K_1)$ the continuous extension of $\Omega^{K_1}$ (so
$\norm J\le1$), and $\omega_{E_1}\colon E_1^k\to F_1$, $\omega_{E_1}(x):=x_1\wedge\dots\wedge x_k$. Then $\omega_{E_1}\in\Alt^k_{K_1}(E_1;F_1)$,
$\norm{\omega_{E_1}}\le1$, and $J(\omega_{E_1}(x))=\Omega^{K_1}(x_1\wedge\dots\wedge x_k)$.
\end{lemma}

\begin{proof}
For a pure wedge, each entry of the determinant array is a signed sum of products
$x_1(i_{\sigma(1)})\cdots x_k(i_{\sigma(k)})$. Each product has absolute value at most
$\prod_b\norm{x_b}_\infty$, so the nonarchimedean inequality in $K_1$ gives the same bound for the determinant.
For any finite decomposition of $\omega$, it follows that
\[
\norm{\Omega^{K_1}(\omega)}_\infty
\le\max_j\prod_i\norm{x_{j,i}}_\infty
\le\sum_j\prod_i\norm{x_{j,i}}_\infty.
\]
Thus the array is bounded, and taking the infimum gives
$\norm{\Omega^{K_1}(\omega)}_\infty\le\norm\omega_\pi$.
The map $\Omega^{K_1}$ is linear and injective by Lemma~\ref{lem:detinj}, with $L=K_1$, $S=\N$ and $V=E_1$.
In particular, $\norm\omega_\pi>0$ whenever $\omega\ne0$.

The remaining norm laws follow directly from decompositions. Concatenation gives the triangle inequality.
Since $k\ge1$, a scalar can be absorbed into the first factor of each wedge; for $\lambda\ne0$ this gives
$\norm{\lambda\omega}_\pi\le\abs\lambda\norm\omega_\pi$, and applying the same inequality to $\lambda^{-1}$ gives
equality. The pure-wedge bound follows by using a single summand.

The array space $\linf(\N^k,K_1)$ is complete, so $\Omega^{K_1}$ extends to the contraction $J$ on $F_1$.
The wedge map into $F_1$ is $K_1$\nobreakdash-multilinear and vanishes whenever two arguments coincide. The pure-wedge
bound makes it continuous with norm at most one, and the asserted identity with $J$ holds on its values by
construction.
\end{proof}

For $k\ge2$, $F_1$ is not nonarchimedean (Proposition~\ref{prop:nonultra}).

\subsection{The support estimate}\label{ssec:suppest}
We next show that the constant coefficient of every completed determinant array is represented by a finite
algebraic exterior vector. The quantitative bound comes from counting the Laurent components needed as wedge
\emph{factors}; the number of wedge summands may be much larger.

Let $\Omega^\kappa\colon\Lambda^k_\kappa E_0\to\kappa^{\N^k}$ be the determinant array over $\kappa$; it is injective by
Lemma~\ref{lem:detinj} with $L=\kappa$ and $V=E_0$. For $y_1,\dots,y_k\in E_0$ we have
$\Omega^{K_1}(y_1\wedge\dots\wedge y_k)=\Omega^\kappa(y_1\wedge\dots\wedge y_k)$: the determinants are the same, with entries in
$\kappa\subseteq K_1$. Support dimensions of elements of $\Lambda^k_\kappa E_0$ are computed in $E_0$ over $\kappa$.

\begin{proposition}[support estimate]\label{prop:supp}
For every $b\in F_1$ there is a unique $\eta(b)\in\Lambda^k_\kappa E_0$ with $\Omega^\kappa(\eta(b))=\coeff_0(Jb)$. The map $\eta$ is
$\kappa$\nobreakdash-linear, and
\begin{equation}\label{eq:suppest}
\sdim\bigl(\eta(b)\bigr)\le kM_\varrho\,\norm b_{F_1}\qquad\text{for all }b\in F_1 .
\end{equation}
\end{proposition}

\begin{proof}
\emph{Step 1: one decomposable.} Let $x_1,\dots,x_k\in E_1$. If some $x_b=0$, the wedge is zero and we take $\eta=0$.
Otherwise put $\nu_b:=\nu(x_b)$ and $l:=-\sum_b\nu_b$, so that
$\prod_b\norm{x_b}_\infty=\varrho^{-l}$. The possible exponent tuples contributing to a constant coefficient form the set
\[
\mathcal N:=\Bigl\{n\in\Z^k:\ n_b\ge\nu_b\text{ for all }b,\ \sum_bn_b=0\Bigr\}.
\]
For such a tuple, the nonnegative integers $n_b-\nu_b$ sum to $l$. Hence $\mathcal N$ is empty if $l<0$;
if $l\ge0$, every coordinate satisfies $\nu_b\le n_b\le\nu_b+l$, so $\mathcal N$ is finite.

Fix $(i_1,\dots,i_k)\in\N^k$ and expand each entry as
$x_b(i_a)=\sum_{n\ge\nu_b}\coeff_n(x_b(i_a))X^n$. In the determinant expansion
\[
\det(x_b(i_a))_{a,b}=\sum_\sigma\sgn(\sigma)\prod_bx_b(i_{\sigma(b)}),
\]
the constant coefficient of each product is its finite coefficient convolution over $\mathcal N$.
Exchanging this finite sum with the permutation sum gives
\[
\coeff_0\det\bigl(x_b(i_a)\bigr)_{a,b}
=\sum_{n\in\mathcal N}\det\bigl(x_{b,[n_b]}(i_a)\bigr)_{a,b}.
\]
The same finite set $\mathcal N$ works for every array coordinate. Thus
$\coeff_0\bigl(\Omega^{K_1}(x_1\wedge\dots\wedge x_k)\bigr)=\Omega^\kappa(\eta)$, where
\[
\eta:=\sum_{n\in\mathcal N}x_{1,[n_1]}\wedge\dots\wedge x_{k,[n_k]}\in\Lambda^k_\kappa E_0.
\]
This is a finite sum in the algebraic exterior power. When $l<0$ it is zero. When $l\ge0$, all its wedge factors
belong to the set
$\{x_{b,[n]}:1\le b\le k,\ \nu_b\le n\le\nu_b+l\}$, which has at most $k(l+1)$ elements. Taking their span gives
\[
\sdim(\eta)\le k(l+1)=k\,(l+1)\varrho^l\cdot \varrho^{-l}\le kM_\varrho\prod_b\norm{x_b}_\infty.
\]
The final bound therefore holds in both cases.

\emph{Step 2: an element of $\Lambda$.} Write $\omega=\sum_jw_j$, where
$w_j=x_{j,1}\wedge\dots\wedge x_{j,k}$, and take $\eta_j$ from Step~1. Additivity gives
\[
\coeff_0(\Omega^{K_1}(\omega))
=\sum_j\coeff_0(\Omega^{K_1}(w_j))
=\Omega^\kappa\Bigl(\sum_j\eta_j\Bigr).
\]
Injectivity of $\Omega^\kappa$ shows that $\eta(\omega):=\sum_j\eta_j$ is independent of the chosen decomposition.
Subadditivity of support dimension now yields
\[
\sdim(\eta(\omega))\le\sum_j\sdim(\eta_j)
\le kM_\varrho\sum_j\prod_i\norm{x_{j,i}}_\infty.
\]
The left side is fixed while the decomposition varies. Taking the infimum therefore gives
$\sdim(\eta(\omega))\le kM_\varrho\norm\omega_\pi$, without requiring a minimizing decomposition.

\emph{Step 3: completion.} Choose $\omega_n\in\Lambda$ with $\norm{b-\omega_n}_{F_1}\to0$. Since $J$ is a contraction,
\[
\norm{Jb-\Omega^{K_1}(\omega_n)}_\infty\le\norm{b-\omega_n}_{F_1}<1
\qquad(n\ge n_0).
\]
Equation~\eqref{eq:small} therefore gives the \emph{exact} identities
\[
\coeff_0(Jb)=\coeff_0(\Omega^{K_1}(\omega_n))=\Omega^\kappa(\eta(\omega_n))
\qquad(n\ge n_0).
\]
By injectivity of $\Omega^\kappa$, the algebraic vectors $\eta(\omega_n)$ are eventually equal. Define $\eta(b)$
to be this common value. Its array is $\coeff_0(Jb)$, so the same injectivity also proves independence of the
approximating sequence and uniqueness. In particular, no limit in $\Lambda^k_\kappa E_0$ is involved.

For every $n\ge n_0$ the fixed vector $\eta(b)$ satisfies
$\sdim(\eta(b))\le kM_\varrho\norm{\omega_n}_\pi$. Since $\norm{\omega_n}_\pi\to\norm b_{F_1}$, this gives~\eqref{eq:suppest}.
Finally, the array characterization proves $\kappa$\nobreakdash-linearity: $J$ is $K_1$\nobreakdash-linear, $\coeff_0$ is
$\kappa$\nobreakdash-linear, and $\Omega^\kappa$ is linear and injective.
\end{proof}

\subsection{The coefficient lift}\label{ssec:Psi}
The support estimate converts norm control into an algebraic obstruction. We apply it to a hypothetical lift,
evaluated on coordinatewise multiplication operators and on the universal wedge map $\omega_{E_1}$.

Suppose that $P$ is a bounded $k$\nobreakdash-linear lift of $A^{k,K_1}_{E_1,E_1;F_1}$ over $K_1$. For $a\in E_1$ let
$\mu_a\in\cL(E_1,E_1)$ be coordinatewise multiplication, $\mu_ax=ax$. The map $a\mapsto \mu_a$ is $K_1$\nobreakdash-linear and
$\norm{\mu_a}\le\norm a_\infty$. Define
\begin{align*}
\Phi(a_1,\dots,a_k;x_1,\dots,x_k)&:=P(\mu_{a_1},\dots,\mu_{a_k})(\omega_{E_1})(x_1,\dots,x_k)\in F_1,\\
\Psi&:=\eta\circ\Phi|_{E_0^{2k}}\colon E_0^k\times E_0^k\to\Lambda^k_\kappa E_0 .
\end{align*}
Since $\Omega^\kappa\circ\eta=\coeff_0\circ J$, the map $\Psi$ is $\coeff_0\circ J\circ\Phi|_{E_0^{2k}}$, read in $\Lambda^k_\kappa E_0$
through the injective map $\Omega^\kappa$.

\begin{lemma}\label{lem:Psi}
The map $\Psi$ has the following properties.
\begin{enumerate}
\item[($\Psi$1)] $\Psi$ is $2k$\nobreakdash-linear over $\kappa$.
\item[($\Psi$2)] $\Psi$ is alternating in the last $k$ slots.
\item[($\Psi$3)] $\sdim(\Psi(z))\le d:=\lfloor kM_\varrho\norm P\rfloor$ for \emph{all} $z\in E_0^{2k}$.
\item[($\Psi$4)] $\Psi$ satisfies \eqref{eq:Pol} for the multiplier family $\mu^\kappa\colon E_0\to\End_\kappa(E_0)$,
$\mu^\kappa_ax=ax$: for all $m\ge1$, $b_1,\dots,b_m\in E_0$, $\alpha\in\N^m$ with $\abs\alpha=k$, and $x\in E_0^k$,
\[
\sum_{\type f=\alpha}\Psi(b_f;x)=\sum_{\type f=\alpha}(b_{f(1)}x_1)\wedge\dots\wedge(b_{f(k)}x_k)\qquad\text{in }\Lambda^k_\kappa E_0 .
\]
In particular $\Psi$ satisfies \eqref{eq:Pol1}.
\end{enumerate}
\end{lemma}

\begin{proof}
The map $\Phi$ is $K_1$\nobreakdash-multilinear, so its restriction to $E_0^{2k}$ is $\kappa$\nobreakdash-multilinear.
Composing with the $\kappa$\nobreakdash-linear map $\eta$ proves ($\Psi$1). Since
$P(\mu_{a_1},\dots,\mu_{a_k})(\omega_{E_1})$ belongs to $\Alt^k_{K_1}(E_1;F_1)$, it vanishes when two of the last $k$ arguments
coincide; the same holds after applying $\eta$, proving ($\Psi$2).

For ($\Psi$3), all inputs from $E_0$ have norm at most one. Thus
\[
\norm{\Phi(a;x)}\le\norm P\prod_i\norm{\mu_{a_i}}\,\norm{\omega_{E_1}}\prod_i\norm{x_i}_\infty\le\norm P.
\]
Apply~\eqref{eq:suppest} and use that support dimension is an integer to obtain the stated bound by $d$.

To prove ($\Psi$4), we first compare coefficients over the infinite field $K_1$, while the values still lie in
$F_1$. Fix $b_1,\dots,b_m\in E_0$ and $x\in E_0^k$. For $\lambda\in K_1^m$ put
$a_\lambda:=\sum_i\lambda_ib_i\in E_1$. The diagonal identity for the lift gives
\[
\Phi(a_\lambda,\dots,a_\lambda;x)
=P(\mu_{a_\lambda},\dots,\mu_{a_\lambda})(\omega_{E_1})(x)
=A(\mu_{a_\lambda})(\omega_{E_1})(x)
=\omega_{E_1}(a_\lambda x_1,\dots,a_\lambda x_k).
\]
Expanding both sides by multilinearity gives
$\sum_\alpha\lambda^\alpha(L_\alpha-R_\alpha)=0$ for all $\lambda\in K_1^m$, where
\[
L_\alpha:=\sum_{\type f=\alpha}\Phi(b_f;x)\in F_1,\qquad
R_\alpha:=\sum_{\type f=\alpha}\omega_{E_1}(b_{f(1)}x_1,\dots,b_{f(k)}x_k)\in F_1.
\]
The field $K_1$ is infinite even when $\kappa$ is finite. Lemma~\ref{lem:idprinc}, applied in the
$K_1$\nobreakdash-vector space $F_1$, therefore yields $L_\alpha=R_\alpha$ for every $\alpha$.

We can now apply $\eta$. Additivity gives
$\eta(L_\alpha)=\sum_{\type f=\alpha}\Psi(b_f;x)$. To identify the other side, take $y_1,\dots,y_k\in E_0$.
By Lemma~\ref{lem:proj} and the determinant-array agreement preceding Proposition~\ref{prop:supp},
\[
J(\omega_{E_1}(y))=\Omega^{K_1}(y_1\wedge\dots\wedge y_k)
=\Omega^\kappa(y_1\wedge\dots\wedge y_k).
\]
This array has entries in $\kappa$, so taking its constant coefficient leaves it unchanged. Uniqueness in
Proposition~\ref{prop:supp} then gives
$\eta(\omega_{E_1}(y))=y_1\wedge\dots\wedge y_k\in\Lambda^k_\kappa E_0$.
Using $y_j=b_{f(j)}x_j$ in each summand of $R_\alpha$ proves
\[
\eta(R_\alpha)=\sum_{\type f=\alpha}(b_{f(1)}x_1)\wedge\dots\wedge(b_{f(k)}x_k),
\]
which is the required identity.
\end{proof}

The order of these operations is essential for the full identity ($\Psi$4). Evaluating the diagonal only at
$\kappa$\nobreakdash-points would give \eqref{eq:Pw} there, which for $\kappa=\F_q$ and $k>q$ is strictly weaker than
\eqref{eq:Pol} (Proposition~\ref{prop:diag}(3)). The proof below uses only the consequence~\eqref{eq:Pol1}.

\subsection{No bounded lift over a Laurent series field}

We now combine the support estimate (Proposition~\ref{prop:supp}) and the coefficient lift (Lemma~\ref{lem:Psi}) with the finite-field theorem of Appendix~\ref{app:FF} (Theorem~\ref{thm:FF}) to prove Theorem~\ref{thm:laurent}.

\begin{proof}[Proof of Theorem~\ref{thm:laurent}]
Suppose $P$ exists, and form $\Psi$ as in \S\ref{ssec:Psi}. By Lemma~\ref{lem:Psi},
$\Psi\colon E_0^k\times E_0^k\to\Lambda^k_\kappa E_0$ with $E_0=\kappa^\N$ is $2k$\nobreakdash-linear over $\kappa$, alternating and hence
antisymmetric in the last $k$ slots, satisfies \eqref{eq:Pol1}, and has $\sdim(\Psi(z))\le d$ for all inputs $z$. Moreover
$k!=0$ in $\kappa$, since $k\ge p=\operatorname{char}\kappa$. This contradicts Theorem~\ref{thm:FF}(2) with $L=\kappa$
and $E_0=\kappa^\N$. The last sentence follows from Proposition~\ref{prop:homogeneous-lift}.
\end{proof}

\begin{remark}\label{rem:infkappa}
Lemmas~\ref{lem:proj} and~\ref{lem:Psi} and Proposition~\ref{prop:supp} hold for every field $\kappa$, and nothing in them
uses a normalization of $\varrho$: the constant $M_\varrho$ makes the argument uniform in $\varrho\in(0,1)$. Only the final appeal to
Theorem~\ref{thm:FF} needs $\kappa$ finite. We do not need, and do not discuss, the analogue of
Theorem~\ref{thm:laurent} for infinite $\kappa$ with this particular pair $(E_1,F_1)$. For an infinite $\kappa$ of
characteristic $p$ the field $\kappa((X))$ is covered by Theorem~\ref{thm:operator-obstructions}(1), with the pair obtained by ascent from its
subfield $\F_p((X))$.
\end{remark}

\subsection{Multiplier families and \texorpdfstring{$c_0$}{c0} bases}\label{lau:ssec:multipliers}
The proof of Theorem~\ref{thm:laurent} uses the hypothetical lift only on coordinatewise multiplication operators and
only at the single form $\omega_{E_1}$. Moreover, Theorem~\ref{thm:FF}(1) needs only finitely supported test vectors with
entries in $\{0,\pm1\}$. Both observations give a much smaller bad parameter space. Throughout this subsection
$\kappa$, $k$, $\varrho$, $K_1$, $E_1$, $F_1$ and $\omega_{E_1}$ are as in Theorem~\ref{thm:laurent}, $c_0(\N,K_1)\subseteq E_1$ carries the
supremum norm, and $\mu_a\in\cL(E_1,E_1)$ denotes multiplication by $a$.

\begin{proposition}[multipliers on $c_0$]\label{lau:prop:c0-multipliers}
The homogeneous map
\[
 \sigma\colon c_0(\N,K_1)\to\Alt^k_{K_1}(E_1;F_1),\qquad
 \sigma(a)(x_1,\dots,x_k)=\omega_{E_1}(ax_1,\dots,ax_k)=(ax_1)\wedge\dots\wedge(ax_k),
\]
has no bounded $k$\nobreakdash-linear lift with values in $\Alt^k_{K_1}(E_1;F_1)$. Consequently, for every
$u_0\in\cL(E_1,E_1)$, the map $a\mapsto A^{k,K_1}_{E_1,E_1;F_1}(u_0+\mu_a)(\omega_{E_1})$ is analytic at no point of $c_0(\N,K_1)$. In
particular $a\mapsto A^{k,K_1}_{E_1,E_1;F_1}(u_0+\mu_a)$ is analytic at no point.
\end{proposition}



\begin{proof}[Proof of Proposition~\ref{lau:prop:c0-multipliers}]
Suppose $p\colon c_0(\N,K_1)^k\to\Alt^k_{K_1}(E_1;F_1)$ is a bounded $k$\nobreakdash-linear lift of $\sigma$. Let
$V_{\mathrm{fin}}=\kappa^{(\N)}$ be the finitely supported sequences with entries in $\kappa$; it lies in both $c_0(\N,K_1)$
and $E_0$. Define $\Psi(a;x)=\eta\bigl(p(a_1,\dots,a_k)(x_1,\dots,x_k)\bigr)$ for $a,x\in V_{\mathrm{fin}}^k$. The proof of
Lemma~\ref{lem:Psi} applies verbatim. The map $\Psi$ is $2k$\nobreakdash-linear over $\kappa$ and alternating in the last
$k$ slots. Its inputs have supremum norm at most one, so~\eqref{eq:suppest} gives
$\sdim\Psi(a;x)\le\lfloor kM_\varrho\norm p\rfloor$. The polarization points $a_\lambda=\sum_i\lambda_ib_i$, with
$b_i\in V_{\mathrm{fin}}$, are finitely supported and hence lie in $c_0(\N,K_1)$, so the identity~\eqref{eq:Pol}, and in
particular~\eqref{eq:Pol1}, follows from the diagonal identity as before. This contradicts Theorem~\ref{thm:FF}(1), since
$k!=0$ in $\kappa$.

For the second assertion, suppose that the map is analytic at $a_0$, with expansion $\sum_n p_n$ there. For
$h\in c_0(\N,K_1)$ and $t\in K_1$, multilinearity gives, in $\Mult^k_{K_1}(E_1;F_1)$,
\[
 \omega_{E_1}\circ(u_0+\mu_{a_0}+t\mu_h,\dots,u_0+\mu_{a_0}+t\mu_h)=\sum_{r=0}^kt^rC_r(h),\qquad C_k(h)=\sigma(h).
\]
Comparing this with the given expansion on scalar lines, Lemma~\ref{lem:unique} shows, exactly as in
Proposition~\ref{prop:homogeneous-lift}, that $p_k$ is a bounded lift of $\sigma$. This is impossible. The operator-valued
statement follows by evaluation at $\omega_{E_1}$.
\end{proof}

\begin{corollary}[diagonal transitions over a $c_0$ base]\label{lau:cor:diagonal-bundle}
Let $U$ be the open unit ball of $c_0(\N,K_1)$, an analytic manifold with one chart. Let $\mathcal V\to U$ be the
trivial bundle with fiber $E_1$ and the two global trivializations $\tau_0(a,z)=(a,z)$ and $\tau_1(a,z)=(a,\mu_{1+a}z)$,
and let $\mathcal B=U\times F_1$. Then $\mathcal V$ is an analytic normed vector bundle whose transitions are isometric
diagonal automorphisms of $E_1$. The naturally induced atlas on $\Alt^k(\mathcal V;\mathcal B)$ is not analytic: the section
that is constantly $\omega_{E_1}$ in the chart induced by $\tau_1$ is, in the chart induced by $\tau_0$, the map
$a\mapsto \omega_{E_1}\circ(\mu_{1+a},\dots,\mu_{1+a})$, which is analytic at no point of $U$.
\end{corollary}

\begin{proof}
For $a\in U$ every coordinate satisfies $\abs{1+a_i}=1$, so $\mu_{1+a}$ is an isometry of $E_1$. Its inverse
$\mu_{(1+a)^{-1}}=\sum_{n\ge0}(-\mu_a)^n$ is a power series of radius one in $a$, since $\mu_a^n=\mu_{a^n}$ has norm at most
$\norm a^n$. Hence both transitions are analytic. If $\theta_a\in\Alt^k(\mathcal V_a;F_1)$ has coordinates $m_0(a)$ and
$m_1(a)$ in the two induced charts, then $m_0(a)=m_1(a)\circ(\mu_{1+a},\dots,\mu_{1+a})$. With $m_1\equiv \omega_{E_1}$, this is the
map of Proposition~\ref{lau:prop:c0-multipliers} with $u_0=\id$, which is analytic at no point.
\end{proof}


The obstruction concerns the atlas obtained from the two analytic frames. The underlying pointwise bundle is trivial and has an analytic product structure; the failure is compatibility of that structure with the natural construction from all analytic frames. The complete target $F_1$ admits no equivalent nonarchimedean norm by Proposition~\ref{prop:nonultra}, so this example does not contradict Corollary~\ref{cor:discrete-targets}.

\begin{corollary}[finite dimension is sharp for nonarchimedean bases]\label{lau:cor:dichotomy}
Let $K=\F_q((X))$, where $q$ is a power of $p$, with an $X$\nobreakdash-adic absolute value, and let $k\ge p$. Let $P$ be a $K$\nobreakdash-Banach space admitting an equivalent nonarchimedean norm. The alternating construction preserves analytic bundles and their operator-valued analytic morphisms over every analytic manifold modeled on $P$, with arbitrary normed fibers, if and only if $P$ is finite-dimensional.
\end{corollary}

\begin{proof}
If $P$ is finite-dimensional, Lemma~\ref{lem:findim} supplies continuous coordinates and Corollary~\ref{cor:bundle-cases} applies. Conversely, suppose that $P$ is infinite-dimensional. Round an equivalent nonarchimedean norm upwards to the discrete value group $\abs{K^\times}$; the result is still an equivalent nonarchimedean $K$\nobreakdash-norm, and its nonzero values lie in $\abs{K^\times}$. Serre's orthonormal-basis theorem~\cite[\S1, Proposition~1]{Se} then identifies $P$ topologically with $c_0(I,K)$ for an infinite set $I$. Choosing a countable subset of $I$ gives bounded linear maps
\[
 \iota\colon c_0(\N,K)\longrightarrow P,\qquad
 \pi\colon P\longrightarrow c_0(\N,K),\qquad \pi\iota=\id.
\]
Apply the Laurent construction with $\kappa=\F_q$ and $K_1=K$. On the open set $V=\{x\in P:\norm{\pi x}<1\}$, take the trivial bundle with fiber $E_1$ and two analytic frames whose transition is $\mu_{1+\pi(x)}$, as in Corollary~\ref{lau:cor:diagonal-bundle}, and the trivial target bundle with fiber $F_1$. If the induced alternating transition were analytic at $x_0\in V$, restriction along $b\mapsto x_0+\iota b$ would make
\[
 b\longmapsto A^k_{E_1,E_1;F_1}\bigl(\mu_{1+\pi(x_0)}+\mu_b\bigr)
\]
analytic at zero. This contradicts Proposition~\ref{lau:prop:c0-multipliers} with $u_0=\mu_{1+\pi(x_0)}$.
\end{proof}

The dichotomy permits arbitrary normed target fibers. In the nonarchimedean Banach target category, Corollary~\ref{cor:discrete-targets} gives analyticity over every base for this discretely valued field.
%% ---- end app-C-laurent.tex
%% ---- begin app-D-ascent.tex
\section{Ascent to an arbitrary base field}\label{app:ascent}\label{sec:ascent}

Proposition~\ref{prop:ascent}, applied in Appendix~\ref{app:assembly} with $E_1$ and $F_1$ as in Theorem~\ref{thm:laurent},
carries analyticity at no point from a spherically complete Laurent subfield $K_1$ to the completion $K'$ of the prescribed
field, since a lift after scalar extension descends to $K_1$. Lemma~\ref{lem:descent} then passes from $K'$ to the field itself.

\subsection{The Laurent subfield}
We use spherical completeness in the sense used in Theorem~\ref{thm:spherical}: every nonempty family of pairwise-meeting
closed balls has a common point. For the field case of the following lemma, see
Schneider~\cite[Lemmas~1.3 and~1.6]{Sch}.

\begin{lemma}[discrete distances]\label{lem:discdist}
Let $Y$ be a complete nonarchimedean space whose distances lie in $r^\Z\cup\{0\}$ for some $r\in(0,1)$. Then $Y$ is
spherically complete.
\end{lemma}

\begin{proof}
Write $\cD(a,\rho)=\{y\in Y:\mathrm d(y,a)\le\rho\}$. In a nonarchimedean space, every point of a ball is a center.
Consequently, whenever two closed balls meet, the one with the smaller radius lies in the other.
The discrete distance assumption also gives
\[
\cD(a,\rho)=\cD(a,\rho^\flat),\qquad
\rho^\flat:=\max\{s\in r^\Z\cup\{0\}:s\le\rho\}.
\]
We may therefore replace every radius by a value in $r^\Z\cup\{0\}$.

Let $(\cD(a_i,\rho_i))_{i\in I}$ be a nonempty family of pairwise-meeting closed balls. If its radii have a least
element $\rho_{i_0}$, then $\cD(a_{i_0},\rho_{i_0})$ lies in every ball of the family, so $a_{i_0}$ is a common point.
Otherwise, discreteness implies that all the radii are positive and that their infimum is zero. Choose
$i_1,i_2,\dots$ so that $\rho_{i_n}$ decreases strictly to zero. The corresponding balls are nested, and
\[
\mathrm d(a_{i_{n'}},a_{i_n})\le\rho_{i_n}\qquad(n'\ge n).
\]
Their centers form a Cauchy sequence. By completeness it has a limit $a$, and closedness puts $a$ in every ball
of the sequence. For any $i\in I$, choose $n$ with $\rho_{i_n}\le\rho_i$. The $n$\nobreakdash-th ball then lies in
$\cD(a_i,\rho_i)$, so $a$ belongs to every ball of the original family.
\end{proof}



\begin{lemma}[Laurent subfield]\label{lem:laurentsub}
Let $K'$ be a complete nontrivially normed field of characteristic $p$, and let $t\in K'$ with $0<\abs t=r<1$.
\begin{enumerate}
\item $K'$ is nonarchimedean, and $\abs c=1$ for every $c\in\F_p^\times$.
\item For $a_m,\dots,a_M\in\F_p$, not all zero, $\bigl\lvert\sum_{i=m}^Ma_it^i\bigr\rvert=r^{\min\{i:a_i\ne0\}}$.
\item Give $\F_p((X))$ the absolute value with $\abs X=r$. There is a continuous isometric ring homomorphism
$\mathrm{ev}_t\colon\F_p((X))\to K'$ such that
$\mathrm{ev}_t\bigl(\sum_{i\ge m}a_iX^i\bigr)=\sum_{i\ge m}a_it^i$. Its image
$K_1$ is a closed subfield of $K'$; it is the closure of $\F_p(t)$, and it is isometrically isomorphic to $\F_p((X))$.
\item $K_1$ is spherically complete.
\end{enumerate}
\end{lemma}

\begin{proof}
(1) Lemma~\ref{lem:ultra} makes $K'$ nonarchimedean. For $c\in\F_p^\times$, the identity $c^{p-1}=1$ gives $\abs c=1$.

(2) The nonzero terms $a_it^i$ have pairwise distinct absolute values $r^i$. In a nonarchimedean field,
$\abs{x+y}=\abs x$ whenever $\abs y<\abs x$. Repeatedly applying this identity shows that the sum has the absolute
value of its largest term, namely $r^{\min\{i:a_i\ne0\}}$.

(3) Evaluation on Laurent polynomials defines a ring homomorphism
$\mathrm{ev}_t\colon\F_p[X,X^{-1}]\to K'$, and (2) shows that
$\abs{\mathrm{ev}_t(g)}=r^{\ord g}=\abs g$. Laurent polynomials are dense in $\F_p((X))$, since truncations converge.
Completeness of $K'$ therefore gives a unique continuous extension to $\F_p((X))$. Passing to limits preserves the
isometry and the ring operations. It also gives the stated evaluation formula, because the partial sums of
$\sum_{i\ge m}a_iX^i$ map to the partial sums of $\sum_{i\ge m}a_it^i$.

The image is complete as the isometric image of a complete space, hence closed in $K'$. It is a subfield: for
$g\ne0$, the identity $\mathrm{ev}_t(g)\mathrm{ev}_t(g^{-1})=1$ provides the inverse. This subfield contains
$\F_p(t)$, while the Laurent polynomials $\F_p[t,t^{-1}]\subseteq\F_p(t)$ are dense in it. The image is consequently
exactly the closure of $\F_p(t)$.

(4) By (3), $K_1$ is complete and has value group $\abs{K_1^\times}=r^\Z$. It is nonarchimedean by (1), so its distances
lie in $r^\Z\cup\{0\}$. Lemma~\ref{lem:discdist} gives spherical completeness.
\end{proof}

\subsection{Ingleton's theorem and its consequences}
Throughout the rest of this section we fix fields $K_1\subseteq K'$ with the following properties.
\begin{enumerate}
\item[(H1)] $K_1$, with the restricted absolute value, is nontrivially normed, complete and spherically complete.
\item[(H2)] $K'$ is nonarchimedean.
\end{enumerate}
By multiplication, $K'$ is a nonarchimedean normed $K_1$\nobreakdash-space. In our application $K'=\Kh$ and $K_1$ is the field of
Lemma~\ref{lem:laurentsub}.

\begin{theorem}[Ingleton~\cite{In}; see also~\cite{vR}]\label{thm:ingleton}
Let $K_1$ be a spherically complete nonarchimedean normed field, $V$ a $K_1$\nobreakdash-vector space with a nonarchimedean norm,
$V_0\subseteq V$ a subspace, $C\ge0$, and $T_0\colon V_0\to K_1$ linear with $\abs{T_0v}\le C\norm v$ for $v\in V_0$. Then
$T_0$ extends to a linear map $T\colon V\to K_1$ with $\abs{Tx}\le C\norm x$ for all $x\in V$.
\end{theorem}

Schneider~\cite[Prop.~9.2]{Sch} gives the more general formulation with a dominating nonarchimedean seminorm.

We will use Theorem~\ref{thm:ingleton} on the nonarchimedean $K_1$\nobreakdash-space $K'$ in two ways: to obtain the scalar
retraction below, and to extend finite-dimensional coordinate functionals in the proof of positivity of the
projective tensor norm. Both applications have codomain $K_1$. The spaces $E_1$ and $F_1$ introduced later need not
be nonarchimedean.

\begin{corollary}[Ingleton projection]\label{cor:proj}
There is a $K_1$\nobreakdash-linear map $\varpi\colon K'\to K_1$ with $\varpi|_{K_1}=\id$ and $\abs{\varpi(\lambda)}\le\abs\lambda$ for all
$\lambda\in K'$.
\end{corollary}

\begin{proof}
Theorem~\ref{thm:ingleton} with $V=K'$, $V_0=K_1$, $T_0=\id$ and $C=1$; its hypotheses hold by (H1) and (H2).
\end{proof}

\begin{samepage}
\begin{lemma}[extension]\label{lem:ext}
Let $\mathbb K\in\{K_1,K'\}$ and let $C\ge0$.
\begin{enumerate}
\item Let $V_0$ be a dense $\mathbb K$\nobreakdash-subspace of a normed $\mathbb K$\nobreakdash-space $V$, let $Y$ be a
$\mathbb K$\nobreakdash-Banach space and $T_0\colon V_0\to Y$ a $\mathbb K$\nobreakdash-linear map with $\norm{T_0x}\le C\norm x$. Then $T_0$ extends
uniquely to a continuous map $T\colon V\to Y$, which is $\mathbb K$\nobreakdash-linear with $\norm{Tx}\le C\norm x$.
\item Let $V_1,\dots,V_k$ be normed $K'$\nobreakdash-spaces with dense subspaces $V_i^0$, $Y$ a $K'$\nobreakdash-Banach space, and
$M\colon\prod_iV_i^0\to Y$ a $K'$\nobreakdash-multilinear map with $\norm{M(d)}\le C\prod_i\norm{d_i}$. Then $M$ extends uniquely to a
continuous $K'$\nobreakdash-multilinear map $\overline M\colon\prod_iV_i\to Y$ with $\norm{\overline M}\le C$.
\end{enumerate}
\end{lemma}
\end{samepage}

\begin{proof}
(1) The bound makes $T_0$ Lipschitz, hence uniformly continuous. Density of $V_0$ and completeness of $Y$ give a
unique continuous extension. Linearity and the norm bound pass to limits by continuity of addition, scalar
multiplication and the norm.

(2) The case $k=0$ is immediate. For $k\ge1$, multilinearity gives the telescoping estimate
\[
\norm{M(d)-M(d')}\le C\rho^{k-1}\sum_i\norm{d_i-d'_i}
\]
whenever $\norm{d_i},\norm{d'_i}\le\rho$ for every $i$, with $\rho>0$. For any $v\in\prod_iV_i$, choose tuples
$d^{(n)}\in\prod_iV_i^0$ converging to $v$. Their coordinates are uniformly bounded, so the estimate makes
$M(d^{(n)})$ Cauchy in $Y$. Define $\overline M(v)$ to be its limit. The same estimate, applied to two approximating
sequences, proves that this definition is independent of the choices and is continuous on bounded sets.
Multilinearity and the product norm bound pass to limits. Finally, density of $\prod_iV_i^0$ proves uniqueness.
\end{proof}

\subsection{Projective base change}
The next construction extends scalars while retaining the original norm and a contractive retraction whenever
the original space is complete. The scalar retraction from Corollary~\ref{cor:proj} is the essential ingredient;
the construction applies to arbitrary normed spaces.

For a normed $K_1$\nobreakdash-space $V$ let $V_{K'}:=V\otimes_{K_1}K'$, a $K'$\nobreakdash-vector space through the second factor, and put
\[
\norm u_\pi:=\inf\Bigl\{\sum_j\norm{x_j}\,\abs{\lambda_j}\ :\ u=\sum_jx_j\otimes\lambda_j\Bigr\},
\]
the infimum over finite decompositions, with an ordinary sum. For a bounded $K_1$\nobreakdash-linear $\psi\colon K'\to K_1$ let
$T_\psi\colon V_{K'}\to V$ be the $K_1$\nobreakdash-linear map with $T_\psi(x\otimes\lambda)=\psi(\lambda)x$; it exists because
$(x,\lambda)\mapsto\psi(\lambda)x$ is $K_1$\nobreakdash-bilinear. Let $\iota_V\colon V\to V_{K'}$, $x\mapsto x\otimes1$; it is
$K_1$\nobreakdash-linear, since $cx\otimes1=x\otimes c=c\cdot(x\otimes1)$ for $c\in K_1$.

\begin{lemma}[projective base change]\label{lem:bc}
Let $V$ be a normed $K_1$\nobreakdash-space; no nonarchimedean hypothesis is made on $V$.
\begin{enumerate}
\item $\norm{T_\psi u}\le\norm\psi\,\norm u_\pi$ for every bounded $K_1$\nobreakdash-linear $\psi\colon K'\to K_1$.
\item $\norm\cdot_\pi$ is a $K'$\nobreakdash-norm on $V_{K'}$.
\item $\iota_V$ is isometric.
\item Let $V\widehat\otimes_\pi K'$ be the completion of $(V_{K'},\norm\cdot_\pi)$. If $V$ is complete, then $T_\varpi$, with
$\varpi$ as in Corollary~\ref{cor:proj}, extends to a $K_1$\nobreakdash-linear map $\Pi_V\colon V\widehat\otimes_\pi K'\to V$ with
$\norm{\Pi_V}\le1$ and $\Pi_V\circ\iota_V=\id_V$.
\end{enumerate}
\end{lemma}

\begin{proof}
(1) For every decomposition $u=\sum_jx_j\otimes\lambda_j$,
\[
\norm{T_\psi u}=\Bigl\lVert\sum_j\psi(\lambda_j)x_j\Bigr\rVert
\le\sum_j\norm\psi\abs{\lambda_j}\norm{x_j}.
\]
Taking the infimum over decompositions proves the contraction estimate.

(2) Concatenating decompositions gives the triangle inequality. Multiplying their second factors by
$\mu\in K'$ gives $\norm{\mu u}_\pi\le\abs\mu\norm u_\pi$; for $\mu\ne0$, applying the same inequality with
$\mu^{-1}$ gives equality. The zero-scalar case is immediate.

It remains to prove positivity. Let $u\ne0$. Choose a finite tensor representation and rewrite its second
factors in a basis of their $K_1$\nobreakdash-span. This gives
\[
u=\sum_{j=1}^nx_j\otimes\lambda_j
\]
with $\lambda_1,\dots,\lambda_n$ linearly independent over $K_1$. Some $x_{j_0}$ is nonzero, since otherwise
$u=0$. On the finite-dimensional subspace $U:=\operatorname{span}_{K_1}\{\lambda_1,\dots,\lambda_n\}\subseteq K'$,
take the coordinate functional $\psi_0(\lambda_j)=\delta_{jj_0}$. It is bounded by some $C\ge0$ by
Lemma~\ref{lem:findim}, using completeness of $K_1$ in (H1). Apply Theorem~\ref{thm:ingleton} to the nonarchimedean
$K_1$\nobreakdash-space $K'$ to extend it to $\psi\colon K'\to K_1$ with $\norm\psi\le C$. Then
$T_\psi u=x_{j_0}$, so (1) gives
\[
0<\norm{x_{j_0}}\le C\norm u_\pi.
\]
Thus $\norm u_\pi>0$. This argument detects a nonzero vector coefficient by a functional on the \emph{scalar}
factor $K'$; it requires no separating continuous dual on $V$.

(3) The elementary tensor representation gives $\norm{x\otimes1}_\pi\le\norm x$. For the reverse inequality,
use $\varpi(1)=1$ and $\norm\varpi\le1$ in (1):
\[
\norm x=\norm{T_\varpi(x\otimes1)}\le\norm{x\otimes1}_\pi.
\]

(4) When $V$ is complete, the contraction $T_\varpi$ on the dense tensor subspace extends by
Lemma~\ref{lem:ext}(1), with $\mathbb K=K_1$, to a $K_1$\nobreakdash-linear contraction $\Pi_V$ with values in $V$.
On the embedded copy of $V$ it satisfies
$\Pi_V\iota_Vx=T_\varpi(x\otimes1)=\varpi(1)x=x$.
\end{proof}

We now fix the spaces to which the descent argument will be applied. Let $E_1$ be a normed $K_1$\nobreakdash-space
and $F_1$ a $K_1$\nobreakdash-Banach space, and put
\[
E:=E_1\widehat\otimes_\pi K',\qquad F:=F_1\widehat\otimes_\pi K',
\]
the completions of $((E_1)_{K'},\norm\cdot_\pi)$ and $((F_1)_{K'},\norm\cdot_\pi)$. They are $K'$\nobreakdash-Banach spaces;
$\iota_E\colon E_1\to E$ and $\iota_F\colon F_1\to F$ are $K_1$\nobreakdash-linear isometries, and $\Pi_F\colon F\to F_1$ is $K_1$\nobreakdash-linear
with $\norm{\Pi_F}\le1$ and $\Pi_F\circ\iota_F=\id_{F_1}$ (Lemma~\ref{lem:bc}).

\begin{lemma}[operators]\label{lem:ops}
For $f\in\cL_{K_1}(E_1,E_1)$ the map $f\otimes\id$ on $(E_1)_{K'}$ is $K'$\nobreakdash-linear with
$\norm{(f\otimes\id)u}_\pi\le\norm f\norm u_\pi$. Its extension $f_{K'}\in\cL_{K'}(E,E)$ satisfies $\norm{f_{K'}}\le\norm f$
and $f_{K'}\circ\iota_E=\iota_E\circ f$, and $f\mapsto f_{K'}$ is $K_1$\nobreakdash-linear.
\end{lemma}

\begin{proof}
For any decomposition $u=\sum_jx_j\otimes\lambda_j$, the image under $f\otimes\id$ has the decomposition
$\sum_jf(x_j)\otimes\lambda_j$, whose cost is at most
$\norm f\sum_j\norm{x_j}\abs{\lambda_j}$. Taking the infimum proves the norm bound on the algebraic tensor product.
Since $E$ is complete, Lemma~\ref{lem:ext}(1), with $\mathbb K=K'$, gives the continuous extension $f_{K'}$ with the
same bound. Its action on $x\otimes1$ proves the embedding identity. Additivity and $K_1$\nobreakdash-homogeneity in
$f$ hold on elementary tensors, hence on their span, and then on $E$ by continuity.
\end{proof}

\begin{lemma}[forms]\label{lem:forms-bc}
For $m\in\Alt^k_{K_1}(E_1;F_1)$ there is $m_{K'}\in\Alt^k_{K'}(E;F)$ with $\norm{m_{K'}}\le\norm m$ and
\[
m_{K'}(\iota_Ex_1,\dots,\iota_Ex_k)=\iota_F\bigl(m(x_1,\dots,x_k)\bigr)\qquad(x_i\in E_1),
\]
and $m\mapsto m_{K'}$ is $K_1$\nobreakdash-linear.
\end{lemma}

\begin{proof}
We first define the form on algebraic tensors, then establish a bound independent of their decompositions, and
finally pass to the completions. Fix a $K_1$\nobreakdash-basis $(\lambda_\beta)_{\beta\in\mathfrak B}$ of $K'$.
Every $u\in(E_1)_{K'}$ has a unique finite expansion
$u=\sum_\beta u^\beta\otimes\lambda_\beta$, with $u^\beta\in E_1$, and each coordinate map $u\mapsto u^\beta$ is
$K_1$\nobreakdash-linear. Define
\[
m^{\mathrm{alg}}(u_1,\dots,u_k):=
\sum_{\beta_1,\dots,\beta_k}
m\bigl(u_1^{\beta_1},\dots,u_k^{\beta_k}\bigr)\otimes
\lambda_{\beta_1}\cdots\lambda_{\beta_k}\in(F_1)_{K'}.
\]
The sum is finite, since each input has finitely many nonzero coordinates.

\textup{(i)} Linearity of the coordinate maps and multilinearity of $m$ make $m^{\mathrm{alg}}$
$K_1$\nobreakdash-multilinear.

\textup{(ii)} On elementary tensors, the definition gives
\[
m^{\mathrm{alg}}(x_1\otimes\mu_1,\dots,x_k\otimes\mu_k)
=m(x_1,\dots,x_k)\otimes\mu_1\cdots\mu_k
\qquad(x_i\in E_1,\ \mu_i\in K').
\]
Indeed, writing $\mu_i=\sum_\beta c_{i\beta}\lambda_\beta$ with $c_{i\beta}\in K_1$ gives
$x_i\otimes\mu_i=\sum_\beta(c_{i\beta}x_i)\otimes\lambda_\beta$. Substitution in the defining sum yields
\[
\sum_{\beta_1,\dots,\beta_k}
\Bigl(\prod_ic_{i\beta_i}\Bigr)m(x)\otimes\prod_i\lambda_{\beta_i}
=m(x)\otimes\prod_i\Bigl(\sum_\beta c_{i\beta}\lambda_\beta\Bigr),
\]
which is the asserted formula.

\textup{(iii)} The map $m^{\mathrm{alg}}$ is $K'$\nobreakdash-multilinear. By additivity in each slot, it suffices
to check $K'$\nobreakdash-homogeneity when every argument is an elementary tensor; there it follows from
\textup{(ii)}. The same formula determines the map on all tuples of algebraic tensors, so the initial basis
choice does not affect it.

\textup{(iv)} For arbitrary decompositions $u_i=\sum_jx_{ij}\otimes\mu_{ij}$, multilinearity and \textup{(ii)} give
\[
m^{\mathrm{alg}}(u)=
\sum_{j_1,\dots,j_k}m(x_{1j_1},\dots,x_{kj_k})\otimes\prod_i\mu_{ij_i}.
\]
The cost of this tensor decomposition is at most
$\norm m\prod_i\bigl(\sum_j\norm{x_{ij}}\abs{\mu_{ij}}\bigr)$.
Taking the infimum separately over the decompositions in each slot proves
\[
\norm{m^{\mathrm{alg}}(u_1,\dots,u_k)}_\pi\le\norm m\prod_i\norm{u_i}_\pi.
\]

\textup{(v)} To prove strong alternation, suppose two distinct slots $a,b$ both contain
$u=\sum_jx_j\otimes\mu_j$. Use this \emph{same} decomposition in both slots and expand all arguments into
elementary tensors. Terms using the same index $j$ in slots $a$ and $b$ vanish because $m$ has a repeated
argument. For $j\ne j'$, the term using indices $(j,j')$ in these slots pairs with the term using $(j',j)$.
Their scalar factors agree, while the values of $m$ are negatives of one another, so the pair sums to zero.
Thus $m^{\mathrm{alg}}$ vanishes whenever any two arguments coincide. Pair cancellation remains valid in
characteristic $2$; merely checking $m^{\mathrm{alg}}(u,\dots,u)=0$ would not establish strong alternation for
$k\ge3$.

\textup{(vi)} Regard $m^{\mathrm{alg}}$ as taking values in the completion $F$.
By \textup{(iii)}, \textup{(iv)} and Lemma~\ref{lem:ext}(2), it extends to a continuous
$K'$\nobreakdash-multilinear map $m_{K'}\colon E^k\to F$ with $\norm{m_{K'}}\le\norm m$.
Strong alternation also passes to the limit: if $v_a=v_b$, approximate this common vector by one sequence in
$(E_1)_{K'}$, and approximate the other entries separately. Every approximating tuple has equal entries in
slots $a,b$, so \textup{(v)} and continuity give $m_{K'}(v)=0$.

\textup{(vii)} Taking every scalar factor equal to one in \textup{(ii)} gives
\[
m_{K'}(\iota_Ex_1,\dots,\iota_Ex_k)
=m^{\mathrm{alg}}(x_1\otimes1,\dots,x_k\otimes1)
=m(x)\otimes1=\iota_F(m(x)).
\]
The defining formula is $K_1$\nobreakdash-linear in $m$. These linear identities pass to the continuous
extensions by uniqueness, proving $K_1$\nobreakdash-linearity of $m\mapsto m_{K'}$.
\end{proof}

\begin{proposition}[descent of a lift]\label{prop:ascent}
Assume \textup{(H1)} and \textup{(H2)}. Let $E_1$ be a normed $K_1$\nobreakdash-space, $F_1$ a $K_1$\nobreakdash-Banach space, and $E,F$ as
above. If $A^{k,K'}_{E,E;F}$ admits a bounded $k$\nobreakdash-linear lift $P$ over $K'$, then $A^{k,K_1}_{E_1,E_1;F_1}$ admits a
bounded $k$\nobreakdash-linear lift $P_1$ over $K_1$ with $\norm{P_1}\le\norm P$. Equivalently, by Proposition~\ref{prop:homogeneous-lift} over $K_1$ and over
$K'$: if $A^{k,K_1}_{E_1,E_1;F_1}$ is analytic at no point, then $A^{k,K'}_{E,E;F}$ is analytic at no point either.
\end{proposition}

\begin{proof}
Extend the operator and form inputs to $K'$, apply the given lift, and retract the value to $F_1$. Explicitly,
for $f_1,\dots,f_k\in\cL_{K_1}(E_1,E_1)$, $m\in\Alt^k_{K_1}(E_1;F_1)$ and $x\in E_1^k$, set
\[
P_1(f_1,\dots,f_k)(m)(x_1,\dots,x_k):=
\Pi_F\Bigl(P\bigl((f_1)_{K'},\dots,(f_k)_{K'}\bigr)(m_{K'})(\iota_Ex_1,\dots,\iota_Ex_k)\Bigr).
\]
We check the norm bound and then the diagonal identity.

For fixed $f$ and $m$, the map $P((f_i)_{K'})(m_{K'})$ is a continuous alternating $K'$\nobreakdash-multilinear form
on $E$. Precomposing each slot with the $K_1$\nobreakdash-linear isometry $\iota_E$ and postcomposing with
$\Pi_F$ gives a $K_1$\nobreakdash-multilinear map on $E_1$, still vanishing whenever two arguments coincide.
Lemmas~\ref{lem:bc}--\ref{lem:forms-bc} give
\[
\norm{P_1(f)(m)}
\le\norm{\Pi_F}\norm P\prod_i\norm{(f_i)_{K'}}\norm{m_{K'}}
\le\norm P\prod_i\norm{f_i}\norm m.
\]
In particular, $P_1(f)(m)\in\Alt^k_{K_1}(E_1;F_1)$.
The form extension is $K_1$\nobreakdash-linear in $m$, the operator extension is $K_1$\nobreakdash-linear in each
$f_i$, and $P$ is $K'$\nobreakdash-multilinear with values in $K'$\nobreakdash-linear operators.
Consequently $P_1$ is $K_1$\nobreakdash-linear in $m$ and $K_1$\nobreakdash-multilinear in $(f_1,\dots,f_k)$.
The displayed estimate proves boundedness at both levels and yields $\norm{P_1}\le\norm P$.

On the diagonal, use $P(g,\dots,g)=A^{k,K'}(g)$ with $g=f_{K'}$, followed by the embedding identities in
Lemmas~\ref{lem:ops} and~\ref{lem:forms-bc} and the retraction identity $\Pi_F\iota_F=\id$:
\begin{align*}
P_1(f,\dots,f)(m)(x)
&=\Pi_F\bigl(m_{K'}(f_{K'}\iota_Ex_1,\dots,f_{K'}\iota_Ex_k)\bigr)\\
&=\Pi_F\bigl(m_{K'}(\iota_Efx_1,\dots,\iota_Efx_k)\bigr)\\
&=\Pi_F\bigl(\iota_F(m(fx_1,\dots,fx_k))\bigr)\\
&=m(fx_1,\dots,fx_k)=A^{k,K_1}(f)(m)(x).\qedhere
\end{align*}
\end{proof}

\begin{remark}\label{rem:ascent}
Proposition~\ref{prop:ascent} is characteristic-free and uses no hypothesis on $k$. It transports negative results
upward only, from $K_1$ to $K'$. Ingleton's theorem is used only on $K'$, with no nonarchimedean norm hypothesis on $F_1$. The completeness of $F_1$ is used essentially: it is what makes the retraction $\Pi_F$ land in $F_1$.
\end{remark}

\subsection{Incomplete base fields}
The preceding construction takes place over the completion $\Kh$ of the prescribed field. To return to $K$,
we restrict scalars on the resulting Banach spaces. Density of $K$ in $\Kh$ forces every continuous
$K$\nobreakdash-multilinear map between these spaces to be $\Kh$\nobreakdash-multilinear, so a lift cannot appear under
this restriction.

\begin{lemma}\label{lem:descent}
Let $K$ be a nontrivially normed field with completion $\Kh$, and let $E,E',F$ be $\Kh$\nobreakdash-Banach spaces, regarded as
$K$\nobreakdash-Banach spaces by restriction of scalars (with the same norms). Then $\cL_K(E,E')=\cL_{\Kh}(E,E')$,
$\Alt^k_K(E;F)=\Alt^k_{\Kh}(E;F)$ and $\Alt^k_K(E';F)=\Alt^k_{\Kh}(E';F)$ with the same norms, $A^{k,K}_{E,E';F}=A^{k,\Kh}_{E,E';F}$,
and every bounded $k$\nobreakdash-linear lift of $A^{k,K}_{E,E';F}$ over $K$ is a bounded $k$\nobreakdash-linear lift of $A^{k,\Kh}_{E,E';F}$ over
$\Kh$.
\end{lemma}

\begin{proof}
Let $g$ be a continuous $K$\nobreakdash-linear map between normed spaces that already carry compatible
$\Kh$\nobreakdash-scalar actions. For $\lambda\in\Kh$, choose $\lambda_n\in K$ with $\lambda_n\to\lambda$.
Since $\norm{\lambda_nx-\lambda x}=\abs{\lambda_n-\lambda}\norm x$, continuity gives
\[
g(\lambda x)=\lim_n g(\lambda_nx)=\lim_n\lambda_ng(x)=\lambda g(x).
\]
Thus $g$ is $\Kh$\nobreakdash-linear. Applying the same argument with all but one slot fixed shows that every
continuous $K$\nobreakdash-multilinear map is $\Kh$\nobreakdash-multilinear. Alternation is the same equal-input
vanishing condition over either field.

The spaces of bounded maps therefore coincide as sets. Their operator norms agree as well, since the norms are
defined by the same pointwise bounds on the same vectors. They are $\Kh$\nobreakdash-Banach spaces, and
$A^{k,K}_{E,E';F}$ and $A^{k,\Kh}_{E,E';F}$ agree by their defining precomposition formula.
Finally, apply the slotwise density argument once more to a bounded $k$\nobreakdash-linear lift $P$ between these
map spaces. It becomes $\Kh$\nobreakdash-multilinear with the same norm and the same diagonal identity, and hence
is a lift over $\Kh$.
\end{proof}

\needspace{8\baselineskip}

%% ---- end app-D-ascent.tex
%% ---- begin app-E-assembly.tex
\section{Proof of the operator obstruction theorem}\label{app:assembly}\label{sec:proofs}

We show that the target spaces $F_1$ and $F$ admit no equivalent nonarchimedean norm (Proposition~\ref{prop:nonultra}) and prove Theorem~\ref{thm:operator-obstructions}(1), which with Proposition~\ref{prop:split-lift}(1) gives Corollary~\ref{cor:classification}.

\subsection{The target has no equivalent nonarchimedean norm}

\begin{proposition}\label{prop:nonultra}
Let $\kappa$ be any field, $k\ge2$, $\varrho\in(0,1)$, and let $K_1=\kappa((X))$, $E_1$, $E_0$ and $F_1$ be as in
Appendix~\ref{sec:laurent}. Let $e_i\in E_0$ be the unit vectors and $\omega_j:=e_{kj+1}\wedge\dots\wedge e_{kj+k}\in F_1$
($j\ge0$), wedges on pairwise disjoint blocks of indices. Then $\norm{\omega_j}_{F_1}=1$ and
\[
\frac N{M_\varrho}\ \le\ \Bigl\lVert\sum_{j<N}\omega_j\Bigr\rVert_{F_1}\ \le\ N\qquad(N\ge1).
\]
Hence the norm of $F_1$ is not nonarchimedean, and $F_1$ admits no equivalent nonarchimedean norm. The same holds for
$F=F_1\widehat\otimes_\pi K'$ as in Appendix~\ref{sec:ascent}, for any $K'\supseteq K_1$ satisfying \textup{(H1)} and \textup{(H2)},
regarded as a normed space over $K'$ or over any subfield of $K'$.
\end{proposition}

\begin{proof}
Each $\omega_j$ is a wedge of unit vectors, so $\norm{\omega_j}_{F_1}\le1$ by the defining projective norm. Evaluation
of its determinant array at the block indices gives $\det I=1$, hence
$\norm{\omega_j}_{F_1}\ge\norm{J\omega_j}_\infty\ge1$. The triangle inequality now gives the upper bound $N$.

For the lower bound, write $\xi_N$ for the sum of the same block wedges in $\Lambda^k_\kappa E_0$, and put
$\omega:=\sum_{j<N}\omega_j\in F_1$. Its array $J\omega$ is the constant Laurent array $\Omega^\kappa(\xi_N)$.
Thus $\coeff_0(J\omega)=\Omega^\kappa(\xi_N)$, and uniqueness in Proposition~\ref{prop:supp} gives
$\eta(\omega)=\xi_N$.

We apply Lemma~\ref{lem:contr} to $\xi_N$. Fix one of the first $N$ blocks and an index $i_0$ in it. Take for
$\varphi$ the $k-1$ coordinate functionals corresponding to the other indices of that block. On every other
block all these functionals vanish, so its contraction is zero; here $k\ge2$ ensures that the minors have
positive size. On the chosen block, the coefficient of $e_{i_0}$ is $\pm1$, while every other cofactor contains
the zero column obtained by evaluating the functionals at $e_{i_0}$. Consequently
$c_\varphi(\xi_N)=\pm e_{i_0}$. These contractions span all $kN$ unit vectors in the first $N$ blocks, so
\[
kN\le\sdim(\xi_N)=\sdim(\eta(\omega))\le kM_\varrho\norm\omega_{F_1}
\]
by~\eqref{eq:suppest}. This proves the lower bound.

If $\norm\cdot_u$ were an equivalent nonarchimedean norm, choose positive constants $c,c'$ with
$\norm\cdot_{F_1}\le c\norm\cdot_u\le c'\norm\cdot_{F_1}$. Then
\[
\frac{N}{M_\varrho}\le\Bigl\lVert\sum_{j<N}\omega_j\Bigr\rVert_{F_1}
\le c\max_{j<N}\norm{\omega_j}_u\le c'
\]
for every $N\ge1$, a contradiction. In particular, the original norm is not nonarchimedean. This argument uses
only the nonarchimedean inequality for $\norm\cdot_u$, not its scalar homogeneity.

Finally, $\iota_F\colon F_1\to F$ is a $K_1$\nobreakdash-linear isometry by Lemma~\ref{lem:bc}(3). The images
$\iota_F\omega_j$ satisfy the same unit-norm and sum-growth estimates, so the argument applies to $F$ as well.
Restriction to a subfield leaves these estimates unchanged.
\end{proof}

If $\varrho\le\frac12$, then $M_\varrho=1$ and $\norm{\sum_{j<N}\omega_j}_{F_1}=N$ exactly.

\subsection{Proof of Theorem~\ref{thm:operator-obstructions}(1)}
\begin{proof}
Let $K$ be a nontrivially normed field of characteristic $p>0$, and let $k\ge p$.
\begin{enumerate}
\item \emph{Choose a Laurent subfield.} By Lemma~\ref{lem:ultra}, $K$ is nonarchimedean. Its completion $\Kh$ is
a complete, nontrivially normed, nonarchimedean field of characteristic $p$. Choose $t\in K$ with
$0<\abs t=\varrho<1$, and let $K_1$ be the closure of $\F_p(t)$ in $\Kh$. Lemma~\ref{lem:laurentsub} identifies $K_1$
isometrically with $\F_p((X))$, with $\abs X=\varrho$, by $\mathrm{ev}_t$. In particular, $K_1$ is complete and
spherically complete.
\item \emph{Obtain the obstruction over $K_1$.} Put $E_1:=\linf(\N,K_1)$, and let $F_1$ be the completion of
$\Lambda^k_{K_1}E_1$ under the projective exterior norm. Transport of scalars along $\mathrm{ev}_t$ preserves
the norms, exterior powers and completions in this construction. It also identifies the operator and
alternating-map spaces, their precomposition maps, and bounded lifts. Theorem~\ref{thm:laurent}, with
$\kappa=\F_p$, therefore shows that $A^{k,K_1}_{E_1,E_1;F_1}$ has no bounded $k$\nobreakdash-linear lift.
\item \emph{Ascend to $\Kh$.} The inclusion $K_1\subseteq\Kh$ satisfies (H1) and (H2), and $F_1$ is complete.
Proposition~\ref{prop:ascent}, with $K'=\Kh$, gives Banach spaces
\[
E:=E_1\widehat\otimes_\pi\Kh,\qquad F:=F_1\widehat\otimes_\pi\Kh
\]
for which $A^{k,\Kh}_{E,E;F}$ has no bounded $k$\nobreakdash-linear lift.
\item \emph{Descend to $K$.} Regard the same $E$ and $F$ as $K$\nobreakdash-Banach spaces. By
Lemma~\ref{lem:descent}, a bounded lift of $A^{k,K}_{E,E;F}$ would extend to a bounded lift of
$A^{k,\Kh}_{E,E;F}$, which is impossible. Proposition~\ref{prop:homogeneous-lift} applies over the possibly incomplete
field $K$: $A^{k,K}_{E,E;F}$ is analytic at no point. Reindexing the arguments by a bijection gives the same conclusion, with the same spaces $E,F$, for every finite index set with $k$ elements.
\item \emph{Exclude an equivalent nonarchimedean target norm.} Since $k\ge p\ge2$, Proposition~\ref{prop:nonultra}
applies to $F_1$ and $F$, after the scalar identification in step~(2). Its conclusion persists when $F$ is
regarded over $K$.\qedhere
\end{enumerate}
\end{proof}

\begin{remark}\label{rem:spaces}
If $K_1=\Kh$, as when $K=\F_p((X))$ and $t=X$, the completed tensor products identify isometrically with
$E_1$ and $F_1$: use Lemma~\ref{lem:bc}(3) and $x\otimes\lambda=\lambda x\otimes1$. Thus $E$ is nonarchimedean in
this case. For $K=\F_q((X))$, $q=p^n$, with its $X$\nobreakdash-adic absolute value, one can instead apply
Theorem~\ref{thm:laurent} directly with $\kappa=\F_q$, obtaining $E=\linf(\N,K)$ and $F=F_1$.
The proof over a general base field does not require, or assert, that $E$ is nonarchimedean.
\end{remark}
%% ---- end app-E-assembly.tex
%% ---- begin app-F-scalar.tex
\section{A counterexample with scalar values}\label{app:scalar}

We prove Theorem~\ref{thm:operator-obstructions}(2) in the form below: Lemma~\ref{scalar:fibre} rules out $k$\nobreakdash-linear
lifts of the determinant vanishing on prescribed subspace pairs, and the chain spaces of~\eqref{scalar:spaces} carry
these conditions to Banach spaces. Its spaces $E,E'$ are the pair $(E,D)$ of Theorem~\ref{thm:operator-obstructions}(2).

\begin{theorem}[scalar counterexample]\label{thm:scalar-counter}
Let $K$ be a complete, nontrivially normed field that is not spherically complete, and let $k\ge1$ satisfy
$k! = 0$ in $K$. There are nonarchimedean Banach spaces $E,E'$ for which
\[
 A^k_{E,E';K}\colon\cL(E,E')\longrightarrow
 \cL\bigl(\Alt^k(E';K),\Alt^k(E;K)\bigr)
\]
is analytic at no point. The spaces may be taken to be closed subspaces of
$\linf(\Lambda;K^{k+1})$ and $\linf(\Lambda;K^k)$, respectively, containing the corresponding $c_0$ spaces,
where $\Lambda$ is countable. Both spaces are nonseparable.
\end{theorem}

The factorial hypothesis says precisely that $K$ has characteristic $p>0$ and $k\ge p$. In particular
$k\ge2$, and $K$ is nonarchimedean by Lemma~\ref{lem:ultra}. We first establish the three ingredients
of the construction.

\subsection{The countable multilinear tail property}

For a set $I$, a family $(a_i)_{i\in I}$ is \emph{null} if, for every $\varepsilon>0$, only finitely many
indices satisfy $\abs{a_i}\ge\varepsilon$. Over a complete nonarchimedean field, every null scalar family
is unconditionally summable: its finite partial sums form a Cauchy net, since the norm of a difference is
at most the largest omitted term. Its sum has norm at most $\sup_i\abs{a_i}$, and such sums may be regrouped.

We derive the multilinear tail property from van der Put's description of the dual of countable $\linf$. The induction is independent of the characteristic.

\begin{lemma}[multilinear tails]\label{scalar:tail}
Let $K$ be complete, nonarchimedean and not spherically complete, let $I$ be countably infinite, and let
$d\ge1$. Every bounded $d$\nobreakdash-linear map $\lambda\colon\linf(I,K)^d\to K$ has an expansion
\begin{equation}\label{scalar:array}
 \lambda(x^1,\dots,x^d)
 =\sum_{(i_1,\dots,i_d)\in I^d}a_{i_1,\dots,i_d}x^1_{i_1}\cdots x^d_{i_d},
 \qquad a_{i_1,\dots,i_d}=\lambda(e_{i_1},\dots,e_{i_d}),
\end{equation}
where $a\in c_0(I^d,K)$. Consequently, for every $\varepsilon>0$, there is a finite $S\subseteq I$ such that
$\abs{\lambda(x^1,\dots,x^d)}\le\varepsilon$ whenever all $\norm{x^r}\le1$ and at least one $x^r$ vanishes
on $S$. In particular $\lambda(e_i,\dots,e_i)\to0$ outside finite subsets of $I$.
\end{lemma}

\begin{proof}
After reindexing, take $I=\N$. Van der Put's theorem~\cite[Th\'eor\`eme~4.1]{VdP}, applied to countably
many copies of $K$ (a countable index set satisfies its nonmeasurability hypothesis), gives
\[
 \varphi(x)=\sum_i x_i\varphi(e_i),\qquad (\varphi(e_i))_i\in c_0(\N,K)
 \quad\bigl(\varphi\in\cL(\linf,K)\bigr).
\]
This is the case $d=1$. Suppose the assertion about the array and expansion holds in smaller degrees.
Fixing a coordinate vector in any slot of $\lambda$ gives a form of degree $d-1$. Hence, for every
$\varepsilon>0$, each slice of the set
\[
 W=\{(i_1,\dots,i_d):\abs{a_{i_1,\dots,i_d}}\ge\varepsilon\}
\]
obtained by fixing one coordinate is finite.

If $W$ were infinite, we could choose $w^{(0)},w^{(1)},\ldots\in W$ with the coordinates
$w^{(n)}_r$ distinct as $n$ varies, for each $r$: a finite list of forbidden coordinates excludes only
finitely many elements of $W$. Define an isometric embedding $T_r\colon\linf(\N,K)\to\linf(\N,K)$
by putting $(T_rx)_{w^{(n)}_r}=x_n$ and zero elsewhere. The pulled-back form
$\mu=\lambda\circ(T_1,\dots,T_d)$ has array $b$ with $\abs{b_{n,\dots,n}}\ge\varepsilon$ for all $n$.
Its slices are null by induction. Thus
\[
 F_n=\{v\in\N^d:\abs{b_v}\ge\varepsilon,\ v_r=n\text{ for some }r\}
\]
is finite. Choose $n_1<n_2<\cdots$ so that $n_s$ exceeds every coordinate occurring in
$F_{n_1}\cup\cdots\cup F_{n_{s-1}}$. Every nonconstant tuple in $\{n_s:s\ge1\}^d$ then has coefficient
of norm $<\varepsilon$: its least coordinate puts a large tuple in the corresponding $F_{n_r}$,
whereas its greatest coordinate violates the choice of the subsequence.

Let $y$ be the indicator of $\{n_s:s\ge1\}$ and set $\psi(x)=\mu(x,y,\dots,y)$. The expansion of each
fixed-coordinate slice gives
\[
 \psi(e_{n_s})
 =b_{n_s,\dots,n_s}
   +\sum_{v'\in\{n_r:r\ge1\}^{d-1},\ v'\ne(n_s,\dots,n_s)}b_{n_s,v'}.
\]
The summands in the last sum form a null family and all have norm $<\varepsilon$. Their supremum is
also $<\varepsilon$: outside a finite set they have norm $<\varepsilon/2$, and the remaining finite
maximum is $<\varepsilon$. The nonarchimedean inequality therefore gives
$\abs{\psi(e_{n_s})}\ge\varepsilon$ for every $s$, contradicting the displayed dual formula.
Thus $a$ is null. Applying that formula in the first slot and the inductive expansion in the remaining
slots gives~\eqref{scalar:array}; the null full array permits regrouping the iterated sums.

Finally choose a finite $S$ containing all coordinates occurring in tuples with $\abs{a_{i_1,\dots,i_d}}
>\varepsilon$. If one input vanishes on $S$, all these terms vanish. Every remaining term has norm at
most $\varepsilon$ on the product of unit balls, so the same holds for the sum. Apply this conclusion to
coordinate vectors outside $S$ to obtain the last assertion.
\end{proof}

\subsection{Chains and a finite-dimensional incompatibility}

Let $L$ be a nonempty finite set and $\Lambda=L^{<\omega}$ the set of finite words in $L$, including the
empty word. Write $s\preceq t$ when $t$ extends $s$, and write $sj$ for the word obtained by appending
$j\in L$. A \emph{$j$-chain} is a set $T=\{s_0,s_1,\dots\}$ with $s_rj\preceq s_{r+1}$ for every $r$.
Let $\Gamma_j$ denote all such chains.

\begin{lemma}[chain gap]\label{scalar:gap}
Chains with different labels have at most one common word. If sets $a_j\subseteq\Lambda$, $j\in L$,
contain all but finitely many words of every chain in $\Gamma_j$, then $\bigcap_{j\in L}a_j\ne\varnothing$.
\end{lemma}

\begin{proof}
If distinct words $s,t$ belong to both a $j$-chain and an $h$-chain, take $s$ shorter than $t$.
Then $t$ extends both $sj$ and $sh$, forcing $j=h$.

For each word $r$ and label $j$, there is an extension $s\succeq r$ such that every extension of $sj$
belongs to $a_j$. Otherwise start with $s_0=r$ and choose successively an extension
$s_{m+1}\succeq s_mj$ outside $a_j$. This is a $j$-chain with infinitely many words outside $a_j$,
a contradiction. Enumerate $L=\{j_1,\dots,j_q\}$ and start at the empty word $r_0$. Choose $s_m\succeq
r_{m-1}$ with the asserted property for $j_m$, and put $r_m=s_mj_m$. The word $r_q$ extends every
$s_mj_m$ and hence belongs to every $a_j$.
\end{proof}

The next lemma is purely algebraic; the alternating maps in its statement need no topology.
Over a field $\kappa$, put $V=\kappa^{k+1}$ and $V'=\kappa^k$. Write $e_1,\dots,e_{k+1}$ and
$e'_1,\dots,e'_k$ for their bases, $\varepsilon^a$ for the coordinate forms on $V$, and $\delta'$ for the
determinant on $V'$. Set
\begin{align*}
 N&=\{\varepsilon^a:1\le a\le k+1\}
       \cup\{\varepsilon^a+\varepsilon^b:1\le a<b\le k+1\},\\
 \Sigma_\nu&=\ker\nu,\qquad C'_s=\{z:z_s=0\},\qquad H'_{cd}=\{z:z_c=z_d\},
\end{align*}
and consider the finite family of pairs
\begin{equation}\label{scalar:family}
 \mathcal F=\{(V,H'_{cd}):c<d\}
 \cup\{(\Sigma_\nu,C'_s):\nu\in N,\ 1\le s\le k\}.
\end{equation}
Let $g_0\colon V\to V'$ be the coordinate projection onto the first $k$ coordinates.

\begin{lemma}[finite fiber obstruction]\label{scalar:fibre}
Suppose $k! = 0$ in $\kappa$. There is no $k$\nobreakdash-linear map
\[
 \tau\colon\operatorname{Hom}_\kappa(V,V')^k\longrightarrow\Alt^k_\kappa(V;\kappa)
\]
satisfying both
\begin{enumerate}
\item $\tau(g_0,\dots,g_0)=\delta'\circ(g_0,\dots,g_0)$;
\item for every $(\Sigma,\Sigma')\in\mathcal F$, the restriction
$\tau(g^1,\dots,g^k)|_{\Sigma^k}$ is zero whenever $g^r(\Sigma)\subseteq\Sigma'$ for every $r$.
\end{enumerate}
\end{lemma}

\begin{proof}
Let $\delta$ be the determinant on $V$. Every alternating $k$\nobreakdash-form $\alpha$ has a unique expression
$\alpha=\delta(w_\alpha,\cdot,\dots,\cdot)$: its evaluations on the basis tuples omitting one vector give
\[
 (w_\alpha)_a=(-1)^{a-1}\alpha(e_1,\dots,\widehat e_a,\dots,e_{k+1}).
\]
For a nonzero $\nu$, this description also gives
\begin{equation}\label{scalar:restriction}
 \alpha|_{(\ker\nu)^k}=0\quad\Longleftrightarrow\quad\nu(w_\alpha)=0.
\end{equation}
Indeed choose $\zeta$ with $\nu(\zeta)=1$ and a basis $y_1,\dots,y_k$ of $\ker\nu$. Writing
$w_\alpha=\nu(w_\alpha)\zeta+z$ with $z\in\ker\nu$ gives
$\delta(w_\alpha,y_1,\dots,y_k)=\nu(w_\alpha)\delta(\zeta,y_1,\dots,y_k)$, whose last factor is nonzero.

Suppose $\tau$ exists. For a function $a\colon\{1,\dots,k\}\to\{1,\dots,k\}$ define
\[
 B_a(\varphi_1,\dots,\varphi_k)
 =\tau(\varphi_1\otimes e'_{a(1)},\dots,\varphi_k\otimes e'_{a(k)}),
 \qquad (\varphi\otimes e'_c)(x)=\varphi(x)e'_c.
\]
If $a$ omits $s$, every input map takes $V$ into $C'_s$. Applying condition (2) to
$(\ker\varepsilon^b,C'_s)$ for each $b$ and using~\eqref{scalar:restriction} makes every coordinate of
$w_{B_a(\varphi)}$ zero. Hence $B_a=0$ unless $a$ is a permutation.

For a permutation $a$ and distinct $c,d$, replace the target vector in slots $a^{-1}(c)$ and $a^{-1}(d)$
by $e'_c+e'_d$. All input maps then take $V$ into $H'_{cd}$, so their image under $\tau$ is zero.
Expansion in these two slots gives $B_a+B_{(c\ d)\circ a}$ and two terms with a missing target coordinate.
The latter vanish by the preceding paragraph. Thus, with $B=B_{\id}$,
\begin{equation}\label{scalar:target-sign}
 B_a=\sgn(a)B\qquad(a\in S_k).
\end{equation}

Define the $(k+1)$\nobreakdash-linear form
\[
 \beta(\varphi_1,\dots,\varphi_k;\chi)
   =\chi\bigl(w_{B(\varphi_1,\dots,\varphi_k)}\bigr).
\]
If the $s$th argument is $\nu\in N$, its rank-one input vanishes on $\Sigma_\nu$, and every other
rank-one input takes values in $C'_s$. Conditions (2) and~\eqref{scalar:restriction} give
$\beta(\varphi_1,\dots,\nu,\dots,\varphi_k;\nu)=0$. Use this identity successively for
$\varepsilon^a$, $\varepsilon^b$ and $\varepsilon^a+\varepsilon^b$ to obtain
\[
 \beta(\dots,\varepsilon^a,\dots;\varepsilon^b)
 =-\beta(\dots,\varepsilon^b,\dots;\varepsilon^a).
\]
Consequently
$t(\sigma)=\beta(\varepsilon^{\sigma(1)},\dots,\varepsilon^{\sigma(k)};
\varepsilon^{\sigma(k+1)})$ changes sign under every transposition of a slot with the last slot.
These transpositions generate $S_{k+1}$, so $t(\sigma)=\sgn(\sigma)t(\id)$.

Expand $g_0=\sum_{c=1}^k\varepsilon^c\otimes e'_c$ in condition (1). The missing-coordinate terms vanish,
and~\eqref{scalar:target-sign} gives
\[
 \varepsilon^1\wedge\cdots\wedge\varepsilon^k
 =\sum_{a\in S_k}\sgn(a)B(\varepsilon^{a(1)},\dots,\varepsilon^{a(k)}).
\]
Apply $\alpha\mapsto\varepsilon^{k+1}(w_\alpha)$. Its value on the left is $(-1)^k$; on the right it is
$\sum_{a\in S_k}\sgn(a)^2t(\id)=k!t(\id)=0$. This is a contradiction. All identities remain valid in
characteristic $2$: repeated arguments vanish by alternation, and no step divides by $2$ or by a factorial.
\end{proof}

We will need a quantitative consequence of this finite incompatibility. Give $K^{k+1}$ and $K^k$ their
maximum norms and use the same family~\eqref{scalar:family} over $K$.

\begin{lemma}[finite test certificate]\label{scalar:certificate}
Assume $k!=0$ in $K$, as in Theorem~\ref{thm:scalar-counter}. There are finite sets of tests $\mathcal T_j$, one for each $j=(\Sigma^j,\Sigma'^j)\in\mathcal F$, consisting of
tuples $(g^1,\dots,g^k;\xi_1,\dots,\xi_k)$ with $g^r\Sigma^j\subseteq\Sigma'^j$ and $\xi_r\in\Sigma^j$,
such that every $K$\nobreakdash-multilinear $\tau$ satisfying condition (1) of Lemma~\ref{scalar:fibre} obeys
\begin{equation}\label{scalar:threshold}
 \max_{j\in\mathcal F}\max_{(g;\xi)\in\mathcal T_j}
 \abs{\tau(g^1,\dots,g^k)(\xi_1,\dots,\xi_k)}\ge1.
\end{equation}
\end{lemma}

\begin{proof}
Write $\operatorname{char}K=p$. All spaces in~\eqref{scalar:family} are defined over $\F_p$.
For each $j$, choose an $\F_p$\nobreakdash-basis of the subspace $\Sigma^j$ defined over $\F_p$, and one of
\[
 \{g\in\operatorname{Hom}_{\F_p}(\F_p^{k+1},\F_p^k):g\Sigma^j\subseteq\Sigma'^j\}.
\]
Their scalar extensions are bases of the corresponding $K$\nobreakdash-spaces: row reduction of the
defining prime-field matrices has the same pivot positions over $K$. Take all $k$\nobreakdash-tuples from
each basis to form $\mathcal T_j$. Multilinearity in the maps and the vectors makes these tests equivalent
to condition (2).

Choose coordinates $\theta$ for the finite-dimensional space of possible $\tau$. Condition (1) is a
system $A\theta=b$, and the tests for condition (2) form a system $C\theta=0$, with all matrices and
constants over $\F_p$. Lemma~\ref{scalar:fibre} over $\F_p$ makes the combined system inconsistent.
Row reduction of its augmented matrix therefore gives row vectors $\lambda,\mu$ over $\F_p$ with
\[
 \lambda A+\mu C=0,\qquad \lambda b=1.
\]
For every $K$\nobreakdash-valued solution of $A\theta=b$ it follows that $\mu C\theta=-1$.
Every nonzero coefficient of $\mu$ has absolute value $1$, since $\F_p^\times$ is finite.
The nonarchimedean inequality now gives~\eqref{scalar:threshold}.
\end{proof}

\subsection{The Banach spaces defined by chain limits}

Set $L=\mathcal F$ and $\Lambda=L^{<\omega}$. Write $\Sigma^j\subseteq V=K^{k+1}$ and
$\Sigma'^j\subseteq V'=K^k$ for the pair indexed by $j\in L$. Define
\begin{equation}\label{scalar:spaces}
 \begin{split}
 E&=\{x\in\linf(\Lambda;V):
       \operatorname{dist}(x_i,\Sigma^j)\longrightarrow0
       \text{ along every }T\in\Gamma_j,\ j\in L\},\\
 E'&=\{z\in\linf(\Lambda;V'):
       \operatorname{dist}(z_i,\Sigma'^j)\longrightarrow0
       \text{ along every }T\in\Gamma_j,\ j\in L\}.
 \end{split}
\end{equation}
Limits along a chain are taken outside its finite subsets. We verify the properties of these full
solution spaces that will enter the proof.

Each subspace in the definition admits a projection $R$ with $\norm R\le1$ and $\norm{1-R}\le1$.
For $\ker\varepsilon^a$, delete coordinate $a$; for $\ker(\varepsilon^a+\varepsilon^b)$, replace coordinate
$a$ by the negative of coordinate $b$; for $V$, use the identity. On $V'$, delete coordinate $s$ for
$C'_s$, and replace coordinate $c$ by coordinate $d$ for $H'_{cd}$. These formulas are projections onto
the stated spaces. Their norm bounds follow directly from the maximum norm and the nonarchimedean
inequality. Moreover
\begin{equation}\label{scalar:distance}
 \operatorname{dist}(v,\Sigma)=\norm{(1-R)v}.
\end{equation}
Indeed $Rv\in\Sigma$ gives one inequality, and $(1-R)(v-y)=(1-R)v$ for every $y\in\Sigma$ gives the other.

For fixed $j$ and $T\in\Gamma_j$, the corresponding condition in~\eqref{scalar:spaces} is thus the inverse
image of $c_0(T;V)$ under the bounded linear map $x\mapsto((1-R_j)x_i)_{i\in T}$, and similarly for $V'$.
The space $c_0$ is closed in the sup norm: a uniform approximation within $\varepsilon/2$ by a null family
leaves only finitely many coordinates of norm at least $\varepsilon$. Each condition is therefore closed
and linear. Their intersection is a closed subspace of the Banach space $\linf(\Lambda;V)$, respectively
$\linf(\Lambda;V')$, so $E$ and $E'$ are nonarchimedean Banach spaces.

Both contain their corresponding $c_0$ spaces. They are also modules over $\linf(\Lambda,K)$ under
coordinatewise multiplication, with $\norm{tx}\le\norm t\norm x$: each residual family in
\eqref{scalar:distance} is multiplied by the same bounded scalar family and remains null.
Finally, if $T\in\Gamma_j$, every bounded $\Sigma^j$\nobreakdash-valued family supported on $T$ belongs
to $E$. For chains with label $j$ its residual is zero; for a chain with another label, its support meets
that chain in at most one word by Lemma~\ref{scalar:gap}. The identical assertion holds for $E'$ and
$\Sigma'^j$.

Since $k\ge2$, the label $j=(V,H'_{12})$ occurs and both its spaces contain a vector of norm $1$.
Fix a $j$-chain $T$. Multiplying either such vector by arbitrary bounded scalar families on $T$ gives
an isometric embedding of $\linf(\N,K)$ into $E$, respectively $E'$. The binary-valued sequences form
an uncountable set at pairwise distance $1$, so both spaces are nonseparable.

\subsection{The global contradiction}

\begin{proof}[Proof of Theorem~\ref{thm:scalar-counter}]
Use the spaces~\eqref{scalar:spaces}. Suppose that $A=A^k_{E,E';K}$ has a bounded $k$\nobreakdash-linear
lift $P$. For $i\in\Lambda$, let $\xi_{[i]}\in E$ denote the family equal to $\xi\in V$ at $i$ and zero
elsewhere. Single-coordinate families belong to $c_0\subseteq E$, and likewise in $E'$. For
$g\in\operatorname{Hom}_K(V,V')$ let $g_{[i]}\in\cL(E,E')$ send $x$ to the family equal to $g(x_i)$ at
$i$ and zero elsewhere. Its norm is at most $\norm g$. Set
\[
 \delta'_i(z_1,\dots,z_k)=\delta'(z_{1,i},\dots,z_{k,i})\in\Alt^k(E';K).
\]
The determinant estimate for the maximum norm gives $\norm{\delta'_i}\le1$. Define the fiber map
\begin{equation}\label{scalar:fibre-map}
 \pi_i(g^1,\dots,g^k)(\xi_1,\dots,\xi_k)
 =P(g^1_{[i]},\dots,g^k_{[i]})(\delta'_i)
       (\xi_{1[i]},\dots,\xi_{k[i]}).
\end{equation}
It is $k$\nobreakdash-linear in the maps and alternating in the vectors, and the lift identity gives
\begin{equation}\label{scalar:fibre-diagonal}
 \pi_i(g,\dots,g)=\delta'\circ(g,\dots,g)
 \qquad\text{for every }g.
\end{equation}
In particular every $\pi_i$ satisfies condition (1) of Lemma~\ref{scalar:fibre}.

Fix $j=(\Sigma,\Sigma')\in L$, a chain $T\in\Gamma_j$, maps $g^1,\dots,g^k$ with
$g^r\Sigma\subseteq\Sigma'$, and vectors $\xi_1,\dots,\xi_k\in\Sigma$. We claim that
\begin{equation}\label{scalar:chain-decay}
 \pi_i(g^1,\dots,g^k)(\xi_1,\dots,\xi_k)\longrightarrow0
 \qquad(i\in T).
\end{equation}
For $t\in\linf(\Lambda,K)$ define $u^r_t$ by
\[
 (u^r_t x)_i=\begin{cases}t_i g^r(x_i),&i\in T,\\0,&i\notin T.\end{cases}
\]
This is a bounded map from $E$ to $E'$. To see that its values lie in $E'$, let $R$ be the projection
onto $\Sigma$ above and decompose, on $T$,
\[
 t_i g^r(x_i)=t_i g^r(Rx_i)+t_i g^r((1-R)x_i).
\]
The first term is a bounded $\Sigma'$\nobreakdash-valued family supported on $T$, hence belongs to $E'$.
The second is null on $T$ by~\eqref{scalar:distance}, hence its zero extension belongs to $c_0\subseteq E'$.
The coordinate formula gives $\norm{u^r_t}\le\norm t\norm{g^r}$, so $t\mapsto u^r_t$ is bounded linear
with values in $\cL(E,E')$.

For $c\in\linf(\Lambda,K)$ put
\begin{equation}\label{scalar:chain-form}
 m_c(z_1,\dots,z_k)=\sum_{i\in T}c_i\delta'(z_{1,i},\dots,z_{k,i}).
\end{equation}
This sum converges. In fact, if $R'$ projects onto $\Sigma'$, write
$z_{r,i}=a_{r,i}+b_{r,i}$ with $a_{r,i}=R'z_{r,i}$ and $b_{r,i}=(1-R')z_{r,i}$. Both families are bounded
by $\norm{z_r}$, and $b_{r,i}\to0$ along $T$. Expanding the determinant, the term with only $a$'s is
zero because $\dim\Sigma'=k-1$. Every remaining term contains an error factor, giving
\[
 \abs{\delta'(z_{1,i},\dots,z_{k,i})}
 \le\max_{1\le r\le k}\left(\norm{b_{r,i}}\prod_{q\ne r}\norm{z_q}\right)
 \longrightarrow0.
\]
Multiplication by bounded $c$ preserves this limit. Completeness of $K$ now gives the unconditional
scalar sum in~\eqref{scalar:chain-form}. Linearity in each argument and in $c$ follows by taking limits
of finite partial sums; repeated arguments make every summand zero. Moreover the determinant estimate gives
\[
 \abs{m_c(z_1,\dots,z_k)}\le\norm c\prod_{r=1}^k\norm{z_r}.
\]
Thus $c\mapsto m_c$ is a bounded linear map into $\Alt^k(E';K)$, of norm at most $1$.

Let $x^r$ be the family equal to $\xi_r$ on $T$ and zero elsewhere; it belongs to $E$ by the
chain-generator inclusion already proved. Coordinatewise multiplication defines a bounded linear map
$\rho\mapsto\rho x^r$ from $\linf(\Lambda,K)$ to $E$. Consequently
\[
 \begin{split}
 \Phi(t^1,\dots,t^k,c,r^1,\dots,r^k)
   =P(u^1_{t^1},\dots,u^k_{t^k})(m_c)(r^1x^1,\dots,r^kx^k)
 \end{split}
\]
is a bounded $(2k+1)$\nobreakdash-linear scalar form on $\linf(\Lambda,K)$, with bound
\[
 \norm P\prod_{s=1}^k\bigl(\norm{g^s}\norm{\xi_s}\bigr)
       \prod_{s=1}^k\norm{t^s}\,\norm c\,\prod_{s=1}^k\norm{r^s}.
\]
The set $\Lambda$ is countably infinite, so Lemma~\ref{scalar:tail} gives
$\Phi(e_i,\dots,e_i)\to0$ outside finite sets. For $i\in T$ we have
$u^s_{e_i}=g^s_{[i]}$, $m_{e_i}=\delta'_i$ and $e_i x^s=\xi_{s[i]}$.
Equation~\eqref{scalar:fibre-map} therefore identifies this diagonal value with the expression in
\eqref{scalar:chain-decay}, proving the claim.

Choose the finite tests of Lemma~\ref{scalar:certificate}, and define
\[
 G_j(i)=\max_{(g;\xi)\in\mathcal T_j}
       \abs{\pi_i(g^1,\dots,g^k)(\xi_1,\dots,\xi_k)},
 \qquad a_j=\{i\in\Lambda:G_j(i)<1\}.
\]
For fixed $j$ and $T\in\Gamma_j$, equation~\eqref{scalar:chain-decay} applies to every test in the finite
set $\mathcal T_j$. Hence $G_j(i)\to0$ along $T$, and $T\setminus a_j$ is finite. Lemma~\ref{scalar:gap}
supplies an $i$ belonging to every $a_j$. But~\eqref{scalar:fibre-diagonal} and
Lemma~\ref{scalar:certificate} force $G_j(i)\ge1$ for some $j$, a contradiction.

There is therefore no bounded $k$\nobreakdash-linear lift of $A$. Proposition~\ref{prop:homogeneous-lift} applies in
this degree without dividing by $k!$: analyticity at any point would make the degree-$k$ Taylor coefficient
such a lift. Thus $A$ is analytic at no point, as claimed.
\end{proof}
%% ---- end app-F-scalar.tex
%% ---- begin app-G-reflection.tex
\section{Parameter splittings and universal analytic reflection}\label{app:reflection}\label{fam:sec:tensors}
A normed space $P$ is polarizable in degree $n$ (condition $(\Gamma^n_P)$, Definition~\ref{ob:def:sequences}) if a bounded
projection $R_n$ of $\widehat\otimes^n_\pi P$ onto $\widehat\Gamma^n_\pi P$ exists; for Banach $P$ over complete $K$ this is
reflection in degree $n$ (Proposition~\ref{fam:prop:tensor-homogeneous}). Analytic descent (Section~\ref{sec:descent}) of every
Banach-valued analytic map through every closed inclusion is stronger; we characterize it and separate the two by an example.

Throughout this appendix, $K$ is complete and nontrivially normed, and $P$ is Banach.
For $n\ge1$, let $\widehat\otimes^n_\pi P$ be the completed projective tensor power as in
Section~\ref{ob:ssec:powers}; it represents bounded $n$\nobreakdash-linear maps from $P^n$ to Banach spaces.
Recall from Definition~\ref{ob:def:sequences} the closed subspace
\[
 \widehat\Gamma^n_\pi P=\overline{\spn\{x^{\otimes n}:x\in P\}}\subseteq \widehat\otimes^n_\pi P.
\]

\begin{proposition}[homogeneous reflection]\label{fam:prop:tensor-homogeneous}
For a fixed $n\ge1$, the following conditions are equivalent.
\begin{enumerate}
\item For every Banach space $Z$, closed subspace $W\subseteq Z$, and bounded $n$\nobreakdash-linear map
$B\colon P^n\to Z$ with diagonal values in $W$, there is a bounded $W$\nobreakdash-valued
$n$\nobreakdash-linear map with the same diagonal.
\item The subspace $\widehat\Gamma^n_\pi P$ is the range of a bounded projection in $\widehat\otimes^n_\pi P$.
\end{enumerate}
\end{proposition}

\begin{proof}
Suppose first that a projection $R_n$ exists. Linearize $B$ to a bounded map
$\widetilde B\colon \widehat\otimes^n_\pi P\to Z$. The diagonal hypothesis and closedness of $W$ imply
$\widetilde B(\widehat\Gamma^n_\pi P)\subseteq W$.
Composing $\widetilde B|_{\widehat\Gamma^n_\pi P}$ with $R_n$ and the canonical tensor map gives the required
multilinear representative, of norm at most $\norm B\norm{R_n}$.
Conversely, take $Z=\widehat\otimes^n_\pi P$, $W=\widehat\Gamma^n_\pi P$, and $B$ equal to the canonical tensor map.
Linearizing the representative supplied by the first condition gives a bounded map into $\widehat\Gamma^n_\pi P$
fixing every pure power. It fixes their linear span and, by continuity, their closed span, so it is a projection.
\end{proof}

\begin{definition}\label{fam:def:universal-reflection}
The Banach parameter space $P$ has \emph{universal analytic reflection} if, for every Banach space $Z$, closed
subspace $W\subseteq Z$, every open $U\subseteq P$, and every map $f\colon U\to W$, analyticity of
the composite $U\to W\hookrightarrow Z$ implies analyticity of $f$.
\end{definition}

\begin{theorem}[analytic tensor criterion]\label{fam:thm:tensor-analytic}
A Banach space $P$ has universal analytic reflection if and only if there are projections
$R_n\colon \widehat\otimes^n_\pi P\to\widehat\Gamma^n_\pi P$ for all $n\ge1$ such that
\begin{equation}\label{fam:eq:exponential-projections}
 \limsup_{n\to\infty}\norm{R_n}^{1/n}<\infty.
\end{equation}
\end{theorem}

\begin{proof}
Assume first that the projections exist. Expand an ambient analytic map with values in $W$ at a point.
The scalar-line uniqueness argument in Theorem~\ref{thm:coefficient-descent} puts every homogeneous diagonal in $W$.
The construction in Proposition~\ref{fam:prop:tensor-homogeneous} replaces its coefficient $p_n$ by a
$W$\nobreakdash-valued coefficient of norm at most $\norm{p_n}\norm{R_n}$ and with the same diagonal.
The exponential bound preserves a positive radius of convergence. The degree-zero coefficient is
already in $W$, and the resulting expansion represents the given map.

Conversely, form the Banach spaces with supremum norms
\[
 Z=c_0\bigl(\{\widehat\otimes^n_\pi P\}_{n\ge1}\bigr),\qquad
 W=c_0\bigl(\{\widehat\Gamma^n_\pi P\}_{n\ge1}\bigr)\subseteq Z.
\]
Here $c_0$ denotes families tending to zero in their respective Banach coordinate spaces; $W$ is closed.
On the open unit ball of $P$, define
\[
 f(x)=(x^{\otimes n})_{n\ge1}.
\]
Since $\norm{x^{\otimes n}}\le\norm x^n$, this vector lies in $W$.
As a $Z$\nobreakdash-valued map it is analytic: its degree-$n$ coefficient is the canonical tensor map
placed in coordinate $n$, and has norm at most one. This gives an expansion at zero. To re-expand
at $a$ with $\norm a<1$, choose $r>0$ with $\norm a+r<1$. Expanding each tensor power in $a+h$
gives translated coefficients, and the sum of all term bounds is at most
$\sum_{n\ge1}(\norm a+r)^n<\infty$. Completeness of $Z$ supplies the coefficients and the
rearranged sum, proving analyticity throughout the unit ball. Universal reflection gives a $W$\nobreakdash-valued
expansion $\sum_mq_m(x,\ldots,x)$ near zero with $\norm{q_m}\le C\rho^{-m}$ for some $C,\rho>0$.
Compose this expansion with the projection $\pi_n\colon W\to\widehat\Gamma^n_\pi P$ and compare on scalar lines
with $x\mapsto x^{\otimes n}$. Uniqueness of scalar power-series coefficients gives
\[
 \pi_nq_n(x,\ldots,x)=x^{\otimes n}\qquad(x\in P).
\]
Linearizing $\pi_nq_n$ therefore gives a projection $R_n$ onto $\widehat\Gamma^n_\pi P$ with
$\norm{R_n}\le\norm{q_n}\le C\rho^{-n}$, proving~\eqref{fam:eq:exponential-projections}.
\end{proof}

\subsection{Projection bounds and the special role of fixed degree}
For $P=K^d$ with the maximum norm, the coordinate construction in
Theorem~\ref{thm:finite-parameters}, applied to the canonical tensor map, gives
$\norm{R_n}\le d^n$ for $d\ge1$. A bounded linear change of coordinates introduces only a fixed
distortion factor to the $n$\nobreakdash-th power, so it does not change the exponential-growth condition.

For $P=\ell^1(I,K)$, the proof of Proposition~\ref{spl:prop:gamma-l1}(1) identifies $\widehat\otimes^n_\pi P$
isometrically with $\ell^1(I^n,K)$, with unit vectors $e_w$ for $w\in I^n$, puts the sum of the $e_w$ over each
permutation orbit in $I^n$ into $\widehat\Gamma^n_\pi P$, and gives a projection onto $\widehat\Gamma^n_\pi P$ of norm
at most $n!$. In these coordinates every element of $\widehat\Gamma^n_\pi P$ is constant on each permutation orbit,
since every pure power is.
Over a nonarchimedean field and for infinite $I$, the bound $n!$ is optimal even among projections that
do not preserve the individual orbit blocks.

\begin{proposition}\label{fam:prop:l1-optimal}
If $K$ is nonarchimedean and complete, and $I$ is infinite, then
\[
 \inf\{\norm R:R\colon \widehat\otimes^n_\pi\ell^1(I,K)\to\widehat\Gamma^n_\pi\ell^1(I,K)
                  \text{ is a bounded projection}\}=n!.
\]
Consequently $\ell^1(I,K)$ does not have universal analytic reflection.
\end{proposition}

\begin{proof}
Proposition~\ref{spl:prop:gamma-l1}(1) gives the upper bound. For the lower bound, take an orbit $\mathcal O$
of a tuple with $n$ distinct indices; it has $n!$ elements. For each $w\in\mathcal O$, let $\lambda_w$
be the common coefficient on $\mathcal O$ of $R(e_w)$.
The orbit sum $\sum_{w\in\mathcal O}e_w$ belongs to $\widehat\Gamma^n_\pi P$ and is fixed by $R$, hence
\[
 \sum_{w\in\mathcal O}\lambda_w=1.
\]
The nonarchimedean inequality implies that some $\abs{\lambda_w}\ge1$.
For this $w$, the ordinary sum norm satisfies
$\norm{R(e_w)}_1\ge n!\abs{\lambda_w}\ge n!$.
Thus $\norm R\ge n!$. Finally, $(n!)^{1/n}$ is unbounded, so
Theorem~\ref{fam:thm:tensor-analytic} excludes universal analytic reflection.
\end{proof}

\begin{corollary}[entire polynomial approximation of a nonanalytic smooth map]
\label{fam:cor:no-common-radius}
Let $K$ be a complete nonarchimedean nontrivially normed field and put
$P=\ell^1(\N,K)$, with the ordinary sum norm. There are Banach spaces
$W\subseteq Z$, with $W$ closed, and a map $f$ on the open unit ball of $P$
such that $f$ is analytic into $Z$ and $C^\infty$ into $W$, but is not analytic
into $W$ at zero. Moreover, $f$ is a uniform limit on every smaller ball of
entire $W$\nobreakdash-valued polynomials.
\end{corollary}

\begin{proof}
Take
\[
 Z=c_0\bigl(\{\widehat\otimes^n_\pi P\}_{n\ge1}\bigr),\qquad
 W=c_0\bigl(\{\widehat\Gamma^n_\pi P\}_{n\ge1}\bigr),\qquad
 f(x)=(x^{\otimes n})_{n\ge1}.
\]
The spaces are Banach and $W$ is closed in $Z$. As in the proof of
Theorem~\ref{fam:thm:tensor-analytic}, the ambient expansion has coefficient
norms at most one and defines an analytic map on the open unit ball.
Smooth reflection through the closed isometric inclusion makes $f$
$C^\infty$ as a map into $W$.

Suppose a $W$\nobreakdash-valued expansion at zero has coefficients $q_n$.
Uniqueness on scalar lines gives
\[
 \pi_nq_n(x,\ldots,x)=x^{\otimes n},
\]
where $\pi_n\colon W\to\widehat\Gamma^n_\pi P$ is the coordinate projection.
Linearizing $\pi_nq_n$ gives a projection
$\widehat\otimes^n_\pi P\to\widehat\Gamma^n_\pi P$ of norm at most $\norm{q_n}$.
Proposition~\ref{fam:prop:l1-optimal} therefore gives
$\norm{q_n}\ge n!$. Since $(n!)^{1/n}$ is unbounded, these coefficients
give a formal multilinear series of radius zero, contradicting the assumed analytic expansion.

For each $n$, choose the sorted-orbit projection
$R_n\colon \widehat\otimes^n_\pi P\to\widehat\Gamma^n_\pi P$. Its composition with the canonical
tensor map is a bounded multilinear representative of $x\mapsto x^{\otimes n}$
with values in $\widehat\Gamma^n_\pi P$. Consequently the truncation
\[
 f_N(x)=(x,x^{\otimes2},\ldots,x^{\otimes N},0,\ldots)
\]
is an entire $W$\nobreakdash-valued polynomial. For $0<r<1$ and $\norm x\le r$,
\[
 \norm{f(x)-f_N(x)}_W\le r^{N+1},
\]
which proves the asserted uniform convergence.
\end{proof}


There is no conflict with Theorem~\ref{thm:summable-parameters}(2): the ambient representation there has one fixed
degree $d$, so factoring a general analytic parameter map through a free $\ell^1$ space incurs a single factorial
bound $d!$, rather than a sequence of factorial bounds across Taylor degrees.
Over $\mathbb R$ and $\mathbb C$ every Banach space has universal analytic reflection, by Theorem~\ref{fam:thm:tensor-analytic}: the symmetrization $\frac1{n!}\sum_{\sigma\in S_n}\sigma$ has norm at most one on $\widehat\otimes^n_\pi P$ and, by polarization, is a projection onto $\widehat\Gamma^n_\pi P$.
%% ---- end app-G-reflection.tex
%% ---- begin app-H-domains.tex
\section{Why a largest full analytic domain need not exist}\label{app:no-largest}\label{dom:section}\label{sec:domains}
In this appendix a pair $(E,F)$ consists of a source space and a value space. It is split in degree $k$
(Definition~\ref{dom:split-pair}) when $j_{E,F}$ has a bounded linear retraction, a splitting for the single
value space $F$. Split pairs form an analytic domain (Definition~\ref{dom:definition}), since a splitting at
the destination of a morphism suffices, with no compatibility between pairs (Lemma~\ref{dom:split-incoming}).

There need not be a largest analytic domain: the endomorphism actions at two objects can each be
analytic while their cross-action is not.
As in Proposition~\ref{prop:homogeneous-lift}, a bounded $d$-linear \emph{lift} of a
homogeneous map $f$ is a bounded multilinear map $B$ with $B(x,\ldots,x)=f(x)$.

\begin{definition}\label{dom:definition}
A full subcategory $\mathcal D\subseteq\mathsf{Vec}_K^{\mathrm{op}}\times\mathsf{Vec}_K$ is an
\emph{analytic domain} for $\Alt^k$ if the restriction of $\Alt^k$ to $\mathcal D$ is analytic on every
hom space.
\end{definition}

\Needspace{5\baselineskip}
\begin{definition}\label{dom:split-pair}
A pair $(E,F)$ is \emph{split in degree $k$} if the isometric inclusion
\[
j_{E,F}\colon\Alt^k(E;F)\longrightarrow\Mult^k(E;F)
\]
has a bounded linear retraction.
\end{definition}

\begin{lemma}\label{dom:split-incoming}
Every morphism-space action whose destination is a split pair is analytic. In particular, the full
subcategory of split pairs is an analytic domain.
\end{lemma}

\begin{proof}
Let $(E,F)$ be split, with retraction $R$, and let the other object be $(D,G)$. For $u\colon E\to D$
and $v\colon G\to F$, the action is
\[
m\longmapsto R\bigl(v\circ m\circ(u,\ldots,u)\bigr).
\]
Before applying $R$, replace the $k$ copies of $u$ by independent variables. This gives a bounded
$(k+1)$\nobreakdash-linear lift of the joint action, with values in
$\cL(\Alt^k(D;G),\Mult^k(E;F))$. Postcomposition with $R$ gives a lift with the required alternating
codomain. No compatibility between the retractions chosen for different objects is needed.
\end{proof}

\begin{theorem}[no largest full analytic domain]\label{dom:no-largest}
Let $p$ be prime, let $k\ge p$, and equip $K=\F_p(t)$ with its $t$\nobreakdash-adic absolute value.
There is no largest full analytic domain for
\[
\Alt^k\colon\mathsf{Vec}_K^{\mathrm{op}}\times\mathsf{Vec}_K\longrightarrow\mathsf{Vec}_K.
\]
The assertion remains true if the domains must be closed under isomorphisms, and if their objects are
restricted to pairs of finite-dimensional, nonarchimedean normed $K$\nobreakdash-spaces. In the latter case,
isomorphism closure is understood inside that restricted ambient category.

More precisely, there are objects $X_k=(E_k,G)$ and $Y_k=(D_k,G)$ whose singleton full subcategories
are analytic, such that $Y_k$ is split and the action on morphisms from $Y_k$ to $X_k$ is not analytic.
\end{theorem}

Here finite-dimensional means finite \emph{algebraic} dimension over the incomplete field $K$. These spaces
need not admit continuous coordinates over $K$. The theorem does not give a counterexample over a complete
field, or a counterexample with Banach objects.


We give the entire construction. Put $L=\F_p((t))$, the completion of $K$.
We first work in degree $p$, then extend to every $k\ge p$ by adjoining coordinate summands.

\subsection{The rigid source and the split pair}

We use the following elementary completion principle. Let $K$ be dense in a complete valued extension $L$,
let $(E_i)$ be a finite family of dense $K$\nobreakdash-subspaces of Banach $L$\nobreakdash-spaces $H_i$,
and let $Z$ be Banach over $L$.
A bounded $K$\nobreakdash-multilinear map $\prod_iE_i\to Z$ extends uniquely to a bounded
$L$\nobreakdash-multilinear map $\prod_iH_i\to Z$, with the same norm. Indeed, approximate each input by
vectors in $E_i$. A telescoping expansion and the multilinear bound make the resulting values Cauchy,
independently of the approximations. The bound persists in the limit, and density of $K$ in $L$ gives
$L$\nobreakdash-linearity in each slot. In the applications below the extended codomain is the complete
space $L$ or $A$, while membership of values on the original inputs in the smaller spaces is retained.

Put $L=\F_p((t))$, the completion of $K$, and consider the normed $L$\nobreakdash-algebra
\[
A=L[\epsilon]/(\epsilon^p),\qquad e_i=\epsilon^i\quad(0\le i<p),
\]
with the maximum norm in this basis. Multiplication satisfies $\norm{xy}\le\norm x\,\norm y$.
Choose elements $a_0,\ldots,a_{p-1}\in L$, together with the finitely many auxiliary scalars
$\tau_w$ introduced below, algebraically independent over $K$. Such a choice is possible because $K$ is
countable and $L$ is uncountable: the elements of $L$ algebraic over any finitely generated extension of
$K$ form a countable set. Set
\[
a=\sum_{i=0}^{p-1}a_i e_i,\qquad E=K^p+Ka\subseteq A.
\]
Every subspace used below carries the norm induced by $A$ or $L$. In particular, $E$ has $K$\nobreakdash-dimension
$p+1$ and completion $A$.

\begin{lemma}\label{dom:rigid}
Every bounded endomorphism of $E$ is multiplication by a scalar in $K$, and every bounded linear map
$E\to K$ is zero.
\end{lemma}

\begin{proof}
A bounded $K$\nobreakdash-linear map extends to completions, and the extension is $L$\nobreakdash-linear
by density of $K$ in $L$. If it preserves $E$, its columns, the images of the $e_i$, give a matrix
$M+ac^{\mathsf T}$, with $M\in M_p(K)$ and $c\in K^p$. Its value at $a$ lies in $K^p+Ka$.
Comparison of homogeneous degrees in the algebraically independent $a_i$ shows first that
$a(c^{\mathsf T}a)=0$, hence $c=0$, and then that $M=sI$ for some $s\in K$.

Similarly, the extension of a bounded map $E\to K$ is an $L$\nobreakdash-linear functional on $A$.
Its coefficients belong to $K$, since its values on the $e_i$ do. Its value on $a$ also lies in $K$,
so algebraic independence makes all its coefficients zero.
\end{proof}

Multiplication by $a$ has lower triangular Toeplitz matrix
\[
(M_a)_{ij}=\begin{cases}a_{i-j},&i\ge j,\\0,&i<j.\end{cases}
\]
Define
\[
D_0=\spn_K\{e_i,\epsilon^i a:0\le i<p\}+Ka^2,
\qquad
\mathcal W=\{e_i:0\le i<p\}\cup\{e_i+e_j:0\le i<j<p\},
\]
and put
\[
D=D_0+\sum_{w\in\mathcal W}K\tau_w w,\qquad
G=\spn_K\{\det(d_1,\ldots,d_p):d_i\in D\}\subseteq L.
\]
The determinant is taken in the basis $(e_0,\ldots,e_{p-1})$. These are finite-dimensional
$K$\nobreakdash-spaces, with completions $A$ and $L$: $D$ contains $K^p$, and $G$ contains $1$.
By multilinearity, $G$ is spanned by determinants of the displayed finite generating set for $D$.
We may therefore identify $G$ with a finite-dimensional subspace of
$K[a_0,\ldots,a_{p-1},(\tau_w)_w]$. Every element has degree at most one in each individual $\tau_w$,
since a nonzero determinant cannot repeat the generator $\tau_w w$. Algebraic independence gives
\begin{equation}\label{dom:tau-gap}
\tau_w^2G\cap G=\{0\}\qquad(w\in\mathcal W).
\end{equation}

\begin{lemma}\label{dom:all-alt}
Every bounded $p$\nobreakdash-linear map $D^p\to G$ is alternating. Thus $(D,G)$ is split in degree $p$.
\end{lemma}

\begin{proof}
Compose such a map $m$ with $G\hookrightarrow L$ and extend to a bounded $L$\nobreakdash-multilinear map
on $A^p$. Fix two argument positions. Put the same $w\in\mathcal W$ in those positions and arbitrary
standard basis vectors in the others, and call the value $c\in G$. Replacing both copies of $w$ by
$\tau_w w\in D$ shows that $\tau_w^2c\in G$, so~\eqref{dom:tau-gap} gives $c=0$.
Taking $w=e_i$ gives vanishing on repeated basis vectors. Taking $w=e_i+e_j$ then gives the signed
interchange identity in the two positions. Expanding two equal arbitrary arguments in the basis proves
that the completed map, and hence $m$, is alternating, including when $p=2$. The inclusion of alternating
maps into multilinear maps is therefore an isometric bijection, whose inverse is the desired retraction.
\end{proof}

\subsection{A missing coefficient}

The coefficient space of alternating maps on $E$ is
\[
C=\{\lambda\in G:a_i\lambda\in G\text{ for every }i\}.
\]
Indeed, an alternating map extends to $\lambda\det$ on the $p$\nobreakdash-dimensional completion $A$.
Determinants of tuples from $E$ span $\spn_K\{1,a_0,\ldots,a_{p-1}\}$: a determinant can contain at
most one copy of $a$, and replacing one basis vector by $a$ realizes each $a_i$ up to sign. Consequently
\begin{equation}\label{dom:coefficient-space}
\Alt^p(E;G)=C\det.
\end{equation}
This identification is isometric, since the determinant bound gives one inequality and evaluation on the
standard basis gives the reverse inequality.

Algebraically set every $\tau_w$ to zero. The image of $G$ is
\[
G_0=\spn_K\{\det(d_1,\ldots,d_p):d_i\in D_0\}.
\]
Put $C_0=\{\lambda\in G_0:a_i\lambda\in G_0\text{ for every }i\}$. Specialization sends $C$ into $C_0$;
no continuity of specialization is asserted. The displayed generators of $D_0$ are homogeneous polynomial
vectors of degrees zero, one, and two. Their determinants are homogeneous, so $G_0$ is graded, and the
defining conditions show that $C_0$ is graded as well.

\begin{lemma}\label{dom:degree-gap}
The homogeneous degree-one part of $C_0$ is zero.
\end{lemma}

\begin{proof}
The degree-two part of $G_0$ is spanned by the $2\times2$ minors of $M_a$ and the coordinates of $a^2$.
These correspond respectively to choosing two degree-one columns, or one degree-two column, in a
determinant of generators. On homogeneous quadratic polynomials define the $K$\nobreakdash-linear map
$\Phi\colon K[a_0,\ldots,a_{p-1}]_2\to K[T]$ by
\[
\Phi(a_i a_j)=
\begin{cases}
T^{i+j},&i+j\ge p-1,\\
0,&i+j<p-1.
\end{cases}
\]
Every $2\times2$ minor of $M_a$ is killed by $\Phi$. When both product terms are nonzero, their index
sums coincide. If only one is nonzero, write the rows as $r<s$ and columns as $i<j$. The surviving term
has index sum $r+s-i-j$, and the missing term implies $r<j$. Hence
$r+s-i-j\le s-i-1\le p-2$, so its image is again zero.

The coefficient of $\epsilon^r$ in $a^2$ is $\sum_{i+j=r}a_i a_j$. Its image is zero for $r<p-1$;
for $r=p-1$ there are exactly $p$ summands and its image is $pT^{p-1}=0$. Thus $\Phi$ kills the entire
degree-two part of $G_0$. If $\ell=\sum_jc_j a_j$ lies in the degree-one part of $C_0$, then
$a_{p-1}\ell\in G_0$, whereas
\[
\Phi(a_{p-1}\ell)=\sum_jc_jT^{p-1+j}.
\]
The powers are distinct, so all $c_j$ vanish.
\end{proof}

\subsection{The nonanalytic cross-action}

Set $X=(E,G)$ and $Y=(D,G)$. Lemma~\ref{dom:rigid} identifies $\cL(E,E)$ isometrically with $K$,
and the self-action of $X$ is
\[
(s,v)\longmapsto s^p v_*.
\]
This is a bounded polynomial, so the singleton full subcategory on $X$ is analytic. The singleton on $Y$
is analytic by Lemmas~\ref{dom:all-alt} and~\ref{dom:split-incoming}.

For $x\in E$, multiplication by $x$ in $A$ restricts to a bounded map $M_x\colon E\to D$.
Indeed, writing $x=b+ha$ with $b\in K^p$ and $h\in K$, its values on the generators $e_j,a$ lie in
\[
K^p+\sum_iK\epsilon^ia+Ka^2\subseteq D.
\]
Submultiplicativity gives a bounded linear family $x\mapsto M_x$ from $E$ to $\cL(E,D)$.
Its matrix is lower triangular with diagonal entry $x_0$, the $e_0$\nobreakdash-coordinate of $x$ in $A$, so $\det M_x=x_0^p$. In the cross-action from
$Y$ to $X$, fix the target morphism to $\id_G$, evaluate at $\det\in\Alt^p(D;G)$, and use
\eqref{dom:coefficient-space}. The resulting map is
\begin{equation}\label{dom:bad-coordinate}
f\colon E\longrightarrow C,\qquad f(x)=x_0^p.
\end{equation}
Its values lie in $C$ because it was obtained by pulling back an alternating $G$\nobreakdash-valued map.

Suppose that $f$ were analytic at zero. Since $f(sx)=s^pf(x)$, uniqueness of power-series coefficients
on scalar lines gives a bounded $p$\nobreakdash-linear map $B\colon E^p\to C$ with diagonal $f$.
This is the homogeneous coefficient argument used in Proposition~\ref{prop:homogeneous-lift}; it requires no factorial
division or completeness. Compose $B$ with $C\hookrightarrow L$ and extend to an $L$\nobreakdash-multilinear
map on $A^p$. Put
\[
b_i=B(e_0,\ldots,e_0,e_i)\in C.
\]
There are $p-1$ initial copies of $e_0$, and $b_0=f(e_0)=1$. Evaluation at $a$ in the last position gives
$\sum_i a_i b_i\in C$. Specializing all $\tau_w$ to zero yields elements $d_i\in C_0$ with
\[
d_0=1,\qquad \sum_i a_i d_i\in C_0.
\]
The degree-one part of this last expression is $\sum_i a_i(d_i)_0$, where $(d_i)_0\in K$ is the
constant term. Its $a_0$\nobreakdash-coefficient equals one. This contradicts Lemma~\ref{dom:degree-gap}.
Thus~\eqref{dom:bad-coordinate}, and therefore the cross-action from $Y$ to $X$, is not analytic.

\subsection{Higher degrees and the categorical conclusion}

For $k=p+r$, let
\[
E_k=E\oplus K^r,\qquad D_k=D\oplus K^r
\]
with maximum norms. Their completions have dimension $k$ over $L$. Determinants of tuples from $E_k$
still span $\spn_K\{1,a_0,\ldots,a_{p-1}\}$, so extraction of the top-degree coefficient identifies
\begin{equation}\label{dom:padded-coefficients}
\Alt^k(E_k;G)=C\det_k
\end{equation}
isometrically. In particular, the coefficient takes values in $C$, a fact needed to retain the preceding
obstruction after increasing the degree.

By Lemma~\ref{dom:rigid}, the extension of every bounded endomorphism of $E_k$ to its
$L$\nobreakdash-completion $A\oplus L^r$ has block form
\[
\begin{pmatrix}sI_p&B\\0&T\end{pmatrix},\qquad s\in K,\quad T\in M_r(K).
\]
The scalar $s$ and entries of $T$ are bounded linear functions of the endomorphism, as is seen by
evaluation on the standard basis. Its determinant is $s^p\det T$, so its action on
\eqref{dom:padded-coefficients} is multiplication by this bounded polynomial. Including postcomposition
shows that the singleton full subcategory on $X_k=(E_k,G)$ is analytic.

The pair $Y_k=(D_k,G)$ is split. Given $m\in\Mult^k(D_k;G)$, restrict the first $p$ arguments to
$D$ and fix the last $r$ arguments at the standard auxiliary basis $e'_1,\ldots,e'_r$ of $K^r$.
The resulting $p$-linear map $\alpha_m$ is alternating by Lemma~\ref{dom:all-alt}. Write
$\pi_D\colon D_k\to D$ and $\pi_r\colon D_k\to K^r$ for the two projections. Explicitly, set
\[
 (\mathcal Rm)(v_1,\ldots,v_k)=
 \sum_{\substack{S\subseteq\{1,\ldots,k\}\\|S|=p}}
 \varepsilon(S)\,
 \alpha_m(\pi_Dv_{s_1},\ldots,\pi_Dv_{s_p})
 \det(\pi_rv_{t_1},\ldots,\pi_rv_{t_r}),
\]
where $s_1<\cdots<s_p$ list $S$, $t_1<\cdots<t_r$ list its complement, and $\varepsilon(S)$
is the sign of the permutation $(s_1,\ldots,s_p,t_1,\ldots,t_r)$. The determinant and the sign
are scalars multiplying the $G$-valued factor. In degree $r=0$ the determinant is $1$.
This shuffle expression is alternating by the same repeated-argument cancellation as the wedge
product, involves no factorial division, and gives a bounded linear map
$\mathcal R\colon\Mult^k(D_k;G)\to\Alt^k(D_k;G)$ with norm at most $\binom{k}{p}$.
It fixes every alternating map: after extension to the $k$\nobreakdash-dimensional completion, both maps
are determined by their value on the standard basis, where the shuffle formula has exactly one nonzero
term. This is the required retraction. For $r>0$, we do not need or assert that every $k$\nobreakdash-linear
map on $D_k$ is alternating.

Finally, restrict the cross-action along the bounded affine family
\[
E\longrightarrow\cL(E_k,D_k),\qquad x\longmapsto M_x\oplus I_r.
\]
Evaluation at $\det_k$, which is $G$\nobreakdash-valued on $D_k$ by Laplace expansion along the first $p$ rows, and coefficient extraction through~\eqref{dom:padded-coefficients} give exactly
the nonanalytic map~\eqref{dom:bad-coordinate}. Thus the cross-action from $Y_k$ to $X_k$ is not analytic.

A largest analytic full subcategory would have to contain both individually analytic singleton
subcategories, and fullness would then require this cross-action to be analytic. This is impossible.
If closure under isomorphisms is required, replace each singleton by its full isomorphism closure.
Every action between two copies is transported from the self-action by bounded linear changes of
coordinates, so both closures remain analytic. The same contradiction proves Theorem~\ref{dom:no-largest}.

\paragraph{\textit{Finite-dimensional bases.}}
The same spaces show that Theorem~\ref{thm:finite-parameters} needs continuous coordinates on the base,
not only finite algebraic dimension. Regard $E$ as an analytic manifold with one chart; by Lemma~\ref{dom:rigid} it has no
continuous coordinates. Over it take the trivial bundle with fiber $D_k\oplus E_k$ and the two frames
$\id$ and $S_x(d,e)=\bigl(d+(M_x\oplus I_r)e,e\bigr)$, together with the trivial bundle with fiber $G$;
$S_x$ and its inverse $(d,e)\mapsto\bigl(d-(M_x\oplus I_r)e,e\bigr)$ are affine in $x$, so these bundles
are analytic. By~\eqref{eq:shear-compression}, evaluation at $\det_k$
and~\eqref{dom:padded-coefficients}, analyticity of the induced alternating atlas at $x=0$ would
make~\eqref{dom:bad-coordinate} analytic at zero, which was excluded above.
%% ---- end app-H-domains.tex
%% ---- begin app-I-coordinates.tex
\section{Further coordinate methods}\label{app:coordinate-methods}
The proof of Proposition~\ref{prop:scalar-countable} gives a splitting for scalar values: it writes down an explicit section of unnormalized antisymmetrization.

\subsection{Countable-type scalar sources}
The scalar obstruction of Theorem~\ref{thm:operator-obstructions}(2) requires large source spaces in a stronger sense than infinite dimension.
A normed space is \emph{of countable type} if a countable subset has dense linear span.

\begin{proposition}[countable-type scalar sources]\label{prop:scalar-countable}
Let $K$ be complete and nonarchimedean. If either $E$ or $D$ is of countable type, then
$A^k_{E,D;K}$ is analytic in every degree. Neither completeness nor a nonarchimedean norm
is required of $E,D$. In particular, the two spaces in
Theorem~\ref{thm:operator-obstructions}(2) must both fail to be of countable type.
\end{proposition}

\begin{proof}
The case $k=0$ is constant. For a normed space $V$, set
\[
 q_V(x)=\inf_{x=\sum_i x_i}\max_i\norm{x_i},
\]
where the decompositions use nonempty finite families, allowing a zero term. Scaling and concatenating them show that
$q_V$ is a nonarchimedean seminorm. Let $V^u$ be the completion of its separated quotient;
the canonical map $j_V\colon V\to V^u$ is contractive and has dense range.
Expanding a bounded scalar multilinear form $m$ along decompositions of each argument gives
\[
 \abs{m(x_1,\ldots,x_k)}\le\norm m\prod_a q_V(x_a).
\]
Thus scalar multilinear forms extend uniquely and isometrically to $V^u$, as do
alternating forms. Extension preserves alternation by density of tuples with repeated
entries, and commutes with permutations. Every bounded $u\colon V\to W$ induces
$u^u\colon V^u\to W^u$ with $\norm{u^u}\le\norm u$. Uniqueness on the dense image makes
$u\mapsto u^u$ bounded linear, and extension of forms intertwines precomposition by $u$ and $u^u$.

If $V$ is of countable type, so is $V^u$. By van der Put's basis theorem
\cite[Proposition~1.4]{VdP}, it has a $t$\nobreakdash-orthogonal basis $(b_i)$ for any
fixed $0<t<1$. Writing $f_i$ for its coordinate functionals, define
\[
 S(m)(x_1,\ldots,x_k)
 =\sum_{i_1<\cdots<i_k}m(b_{i_1},\ldots,b_{i_k})
                        \prod_{a=1}^k f_{i_a}(x_a).
\]
The weighted coordinates $\abs{f_i(x)}\norm{b_i}$ form a null family bounded by
$t^{-1}\norm x$. Products of finitely many bounded null families are null, so the
displayed sum converges unconditionally in $K$ and $\norm{S(m)}\le t^{-k}\norm m$.
Let $\mathcal A_V$ be unnormalized antisymmetrization, the signed sum over permutations
of the arguments. It is bounded and alternating-valued in every characteristic.
On finite-coordinate vectors, sorting distinct indices shows
$\mathcal A_V S(m)=m$ for alternating $m$; repeated indices contribute zero.
Density proves the same identity everywhere. Hence $R=\mathcal A_V S$ retracts all
multilinear forms onto alternating forms, while $S|_{\Alt}$ is a bounded right inverse
of $\mathcal A_V$. Transport through $j_V$ gives both operators on the original $V$.

If $E$ is of countable type, a lift of $A$ is
$R_E\bigl(m\circ(u_1,\ldots,u_k)\bigr)$.
If $D$ is of countable type, use instead
$\mathcal A_E\bigl(S_D(m)\circ(u_1,\ldots,u_k)\bigr)$.
Both are bounded multilinear in the $u_i$, linear in $m$, and have diagonal $A(u)(m)$;
the second identity uses that antisymmetrization commutes with simultaneous precomposition.
Proposition~\ref{prop:homogeneous-lift} concludes the proof.
\end{proof}


\subsection{Continuous algebraic polynomials on \texorpdfstring{$c_0$}{c0}}
Theorem~\ref{thm:summable-parameters}(1) starts with bounded multilinear coefficients in an ambient space. For one
homogeneous polynomial, interpolation supplies a lift even when only the diagonal is known to be
continuous. The resulting bound is degreewise; it gives no assertion about a common radius across
an infinite sequence of degrees.

\begin{proposition}[continuous algebraic polynomials on $c_0$]\label{fam:prop:c0-algebraic}
Let $K$ be nontrivially normed and nonarchimedean, let $Z$ be complete and nonarchimedean, and let
$n\ge1$. Suppose $p\colon c_0(I,K)\to Z$ is continuous and is an algebraic homogeneous polynomial:
there exists a $K$\nobreakdash-multilinear map $b\colon c_0(I,K)^n\to Z$, not assumed continuous,
with $p(x)=b(x,\ldots,x)$. Then $p$ has a bounded $n$\nobreakdash-linear lift $q$. More precisely,
there is $C_{K,n}<\infty$, independent of $I,Z,p$, such that
\[
 \norm q\le C_{K,n}\norm p_{\mathrm{diag}},\qquad
 \norm p_{\mathrm{diag}}:=\inf\{D\ge0:\norm{p(x)}\le D\norm x^n\text{ for all }x\}.
\]
No completeness assumption on $K$ is needed.
\end{proposition}
\begin{proof}
Continuity at zero and $p(ax)=a^np(x)$ imply $\norm p_{\mathrm{diag}}<\infty$.
Explicitly, choose $t\in K$ with $0<\abs t<1$ and $\delta>0$ such that $\norm{p(x)}\le1$
when $\norm x<\delta$. For $x\ne0$, choose an integer $m$ with
$\abs t\delta\le\abs{t^m}\norm x<\delta$. Homogeneity gives
$\norm{p(x)}\le(\abs t\delta)^{-n}\norm x^n$.

Fix a total order on $I$. On each finite coordinate subspace, the algebraic formula for $p$ is an
ordinary homogeneous polynomial of degree $n$. Polynomial uniqueness over the infinite field $K$
makes its coefficients compatible as the finite coordinate set varies. Denote the coefficient of
$x^\alpha$ by $c_\alpha$, where $\alpha\colon I\to\N$ is finitely supported and $\sum_i\alpha_i=n$.
Its support has at most $n$ elements.

Choose distinct $a_0,\ldots,a_n$ in the valuation ring of $K$, and let $L_n\ge1$ bound the absolute
values of all entries of the inverse Vandermonde matrix $(a_i^j)_{0\le i,j\le n}$.
If the support of $\alpha$ has $d\le n$ elements, interpolate the restriction of $p$ to those
coordinates on the grid $\{a_0,\ldots,a_n\}^d$. Each grid vector has norm at most one. Applying the
inverse Vandermonde matrix in each coordinate and using the nonarchimedean inequality in $Z$ yields
\[
 \norm{c_\alpha}\le L_n^d\norm p_{\mathrm{diag}}
 \le L_n^n\norm p_{\mathrm{diag}}.
\]
For a nondecreasing word $i_1\le\cdots\le i_n$, write $c_{i_1,\ldots,i_n}$ for the coefficient
with those multiplicities. Define
\[
 q(x_1,\ldots,x_n)=
 \sum_{i_1\le\cdots\le i_n}
 c_{i_1,\ldots,i_n}(x_1)_{i_1}\cdots(x_n)_{i_n}.
\]
The coefficients are uniformly bounded and each input is a null family, so the summands tend to zero
outside finite sets of words. Completeness and the nonarchimedean inequality give convergence,
multilinearity, and the bound $\norm q\le L_n^n\norm p_{\mathrm{diag}}$.
On finite supports its diagonal is $p$; density and continuity extend the identity to all of $c_0(I,K)$.
Thus $C_{K,n}=L_n^n$ suffices.
\end{proof}

This bound is uniform in the number of coordinates, but the interpolation nodes and $L_n$ depend on
the degree. No exponential growth bound for $C_{K,n}$ is asserted. Consequently this proposition does
not justify replacing lifting norms by diagonal norms in an arbitrary infinite analytic expansion.
Theorem~\ref{thm:summable-parameters}(1) obtains its stronger, radius-preserving conclusion from the
already bounded ambient coefficients.


\subsection{A nonspherical target norm with analytic precomposition}
\begin{remark}\label{rem:nonspherical-analytic-target}
Spherical completeness of the given target norm is not necessary for analyticity, even in degrees
at least the characteristic. Let $K=\F_p((t))$ with $\abs t=1/2$ and let $F=c_0(\N,K)$, the space
of sequences tending to zero. Define
\[
 w_j=1+\frac1{j+2},\qquad \norm{x}_w=\sup_{j\ge0}w_j\abs{x_j}.
\]
This is a nonarchimedean norm satisfying
\[
 \norm{x}_\infty\le\norm{x}_w\le\frac32\norm{x}_\infty.
\]
Thus $F$ is Banach for either norm. The usual supremum norm has nonzero values in $2^\Z$,
since every nonzero null sequence attains its maximum absolute value. The discrete-distance
argument in the proof of Corollary~\ref{cor:discrete-targets} therefore makes the supremum norm
spherically complete. Theorem~\ref{thm:spherical} gives analytic precomposition in every finite
degree for that norm, for all normed $E,E'$. Equivalent norms preserve the bounded maps and their
analytic power-series expansions: the norms of coefficients change by a fixed factor and convergence
is unchanged. Thus precomposition remains analytic for $\norm\cdot_w$, also directly by
Corollary~\ref{cor:discrete-targets}.

Nevertheless, the weighted norm is not spherically complete. Write $e_j$ for the coordinate vectors,
set $a_n=\sum_{j<n}e_j$, and consider the closed balls
\[
 B_n=\overline B_w(a_n,w_n).
\]
They are nested because $w_{n+1}<w_n$ and $\norm{a_{n+1}-a_n}_w=w_n$; the nonarchimedean
inequality gives $B_{n+1}\subseteq B_n$. If $x$ belonged to all of them, then, for each $j$ and
any $n>j$,
\[
 \abs{x_j-1}\le\frac{w_n}{w_j}<1.
\]
Hence $\abs{x_j}=1$ for every $j$, contradicting $x_j\to0$. Their intersection is empty.
This example separates spherical completeness of a particular norm from analyticity, which is
invariant under equivalent norms. It does not exclude a characterization using an equivalent
spherically complete norm: this target has exactly such a norm.
\end{remark}
%% ---- end app-I-coordinates.tex
%% ---- begin app-J-ordering.tex
\section{Ordering splittings}\label{app:ordering}

Throughout, $K$ is a nontrivially normed field, $\Sigma$ is a totally ordered set, and $\vartheta\in K$ is the fixed
element with $0<\abs\vartheta<1$.

\subsection{Lower estimates}\label{ord:ssec:lower}

\begin{proof}[Proof of Theorem~\ref{fib:thm:lower}]
(1) \emph{Step 1 (the weight $\nu$).} Let $x,y\in V$ be nonzero and $\Sigma_0\subseteq\Sigma$ a finite set containing
their supports. In this proof an interval is a set of consecutive elements of $\Sigma_0$; disjoint ones lie in disjoint
intervals of $\Sigma$, so~\eqref{fib:eq:lower} applies to them. For an interval $I$ let $\nu_x(I)$ be the supremum of
$\sum_i\norm{\cP_{I_i}x}^r$ over the families of pairwise disjoint intervals $I_i\subseteq I$. Then $\nu_x$ is monotone
and superadditive on disjoint intervals, $\norm{\cP_Ix}^r\le\nu_x(I)$ and $\nu_x(\Sigma_0)\le C^r\norm x^r$. Put
$a_I=\nu_x(I)/\nu_x(\Sigma_0)$, $b_I=\nu_y(I)/\nu_y(\Sigma_0)$ and $\nu(I)=a_I+b_I$, so that $\nu(\Sigma_0)=2$.

\emph{Step 2 (the binary tree).} Build a binary tree of intervals with root $\Sigma_0$ of depth $0$. A node $I$ is a leaf
if it has at most one element or $\nu(I)=0$. Otherwise let $t_I$ be the least $t\in I$ with
$\nu(\{t'\in I:t'\le t\})\ge\nu(I)/2$; the children of $I$, of depth one more, are $I_<=\{t\in I:t<t_I\}$ and
$I_>=\{t\in I:t>t_I\}$. Then $\nu(I_<)<\nu(I)/2$ by minimality of $t_I$ and $\nu(I_>)\le\nu(I)/2$ by superadditivity,
so the nodes of depth $d$ are pairwise disjoint and satisfy $\nu\le2^{1-d}$.

\emph{Step 3 (the sum over the nodes).} For an interval $I$ let $z_I$ be the sum of $x(t)y(t')\,e_t\otimes e_{t'}$ over
$t<t'$ in $I$, so that $\cP_<(x\otimes y)=z_{\Sigma_0}$. At an internal node
\[
 z_I=z_{I_<}+z_{I_>}+\cP_{I_<}x\otimes\cP_{I\setminus I_<}y+\cP_{\{t_I\}}x\otimes\cP_{I_>}y,
\]
and $z_I=0$ at a leaf $I$: it has at most one element, or $\nu(I)=0$ and $x$ vanishes on $I$. So
$\cP_<(x\otimes y)$ is the sum of the last two terms over the internal nodes. At a node $I$ they have projective norm
at most
\[
 2\nu_x(I)^{1/r}\nu_y(I)^{1/r}\le2C^2(a_Ib_I)^{1/r}\norm x\norm y.
\]
At depth $d$, $a_Ib_I\le(\nu(I)/2)^2\le4^{-d}$, and $\sum_I(a_Ib_I)^{1/2}\le1$ by the Cauchy--Schwarz inequality, as
$\sum_Ia_I\le1$ and $\sum_Ib_I\le1$; so $\sum_I(a_Ib_I)^{1/r}\le2^{-d(2/r-1)}$. Summing over $d$ gives
$\norm{\cP_<(x\otimes y)}\le C_\sharp\norm x\norm y$. The bounded bilinear map $(x,y)\mapsto\cP_<(x\otimes y)$ extends
from $V\times V$ to $E\times E$ with values in $E\widehat\otimes_\pi E$; its linearization is the asserted extension.

(2) For $1\le a<k$ let $\cP_<^{(a)}=\id^{\otimes(a-1)}\otimes\cP_<\otimes\id^{\otimes(k-a-1)}$, a map
$\otimes^k_\pi E\to\widehat\otimes^k_\pi E$ of norm at most $C_\sharp$, extended to the completion, and
$\cP^{(k)}=\cP_<^{(k-1)}\circ\dots\circ\cP_<^{(1)}$. Then $\norm{\cP^{(k)}}\le C_\sharp^{k-1}$, and $\cP^{(k)}$ fixes
$e_{t_1}\otimes\dots\otimes e_{t_k}$ if $t_1<\dots<t_k$ and maps it to $0$ otherwise. Put
$s=\cP^{(k)}\circ\operatorname{alt}^k_E$ and $\rho=\widehat\wedge^k_E\circ\cP^{(k)}$. For $x_1,\dots,x_k\in V$,
\[
 s(x_1\wedge\dots\wedge x_k)=\sum_{t_1<\dots<t_k}\det\bigl(x_b(t_a)\bigr)_{a,b}\,e_{t_1}\otimes\dots\otimes e_{t_k},
\]
whose image under $\widehat\wedge^k_E$ is $x_1\wedge\dots\wedge x_k$. Such wedges span a dense subspace of
$\Lambda^k_\pi E$, as $V$ is dense and $\norm{x_1\wedge\dots\wedge x_k}\le\prod_i\norm{x_i}$. Hence
$\widehat\wedge^k_E\circ s=\id$ and $\rho\circ\operatorname{alt}^k_E=\widehat\wedge^k_E\circ s=\id$. If $E=V$, then
$\cP_<$ preserves $V\otimes V$, so $s$ and $\rho$ take values in the uncompleted spaces.

(3) follows from~(2) and Proposition~\ref{fib:prop:split}.
\end{proof}

\subsection{Ordering splittings in every degree}\label{ord:ssec:ordering}

\begin{proof}[Proof of Proposition~\ref{spl:prop:ordering}]
\emph{Step 1 (source and target).} As at the end of part~(1) of the proof of Theorem~\ref{fib:thm:lower}, $\cP_<$
extends to a map $E\otimes_\pi E\to E\widehat\otimes_\pi E$ of norm at most $C$, and part~(2) of that proof, with $C$ in
place of $C_\sharp$ and $n$ in place of $k$, proves~(1) and~(2).

\emph{Step 2 (diagonal operators).} For $n=1$ take $R_1=\id$; let $n\ge2$. Let $(1\,2)$ be the isometry exchanging the
factors of $E\otimes_\pi E$, $\cP_>=(1\,2)\circ\cP_<\circ(1\,2)$ and $\Delta=\id-\cP_<-\cP_>$, so that
$\norm\Delta\le1+2C$. For $T\in\{\cP_<,\Delta\}$ and $1\le a<n$ let
$T^{(a)}=\id^{\otimes(a-1)}\otimes T\otimes\id^{\otimes(n-a-1)}$, extended to $\widehat\otimes^n_\pi E$, of norm at most
$\norm T$. For $\mathbf t\in\Sigma^n$ and $e_{\mathbf t}=e_{t_1}\otimes\dots\otimes e_{t_n}$, the vector
$\cP_<^{(a)}e_{\mathbf t}$, respectively $\Delta^{(a)}e_{\mathbf t}$, is $e_{\mathbf t}$ if $t_a<t_{a+1}$, respectively
$t_a=t_{a+1}$, and $0$ otherwise. So these operators preserve $\otimes^nV$ and commute on it, hence on
$\widehat\otimes^n_\pi E$.

\emph{Step 3 (the projection).} For a composition $c=(c_1,\dots,c_m)$ of $n$ let $J(c)$ be the set of partial sums
$c_1+\dots+c_i$, $i<m$, and $\Pi_c=\prod_{a\in J(c)}\cP_<^{(a)}\prod_{a\notin J(c)}\Delta^{(a)}$, with $a$ ranging over
$\{1,\dots,n-1\}$. Then $\Pi_ce_{\mathbf t}=e_{\mathbf t}$ if $\mathbf t$ is nondecreasing with constant blocks of
lengths $c_1,\dots,c_m$, and $\Pi_ce_{\mathbf t}=0$ otherwise. The stabilizer of such a $\mathbf t$ in $S_n$, acting by
isometric permutations of the factors, is $S_{c_1}\times\dots\times S_{c_m}$; if $\Omega_c$ is a set of representatives
of its left cosets, then $O_{\mathbf t}=\sum_{\sigma\in\Omega_c}\sigma e_{\mathbf t}$. Hence
$R_n=\sum_c\sum_{\sigma\in\Omega_c}\sigma\circ\Pi_c$, the first sum over all compositions of $n$, satisfies
$R_ne_{\mathbf t}=O_{\mathbf t}$ for nondecreasing $\mathbf t$ and $R_ne_{\mathbf t}=0$ otherwise, and
$\norm{R_n}\le\sum_cn!\,(c_1!\cdots c_m!)^{-1}C^{m-1}(1+2C)^{n-m}$.

\emph{Step 4 (the image).} Let $\Gamma_V$ be the span of the $O_{\mathbf t}$ with $\mathbf t$ nondecreasing. Then
$R_n(\otimes^nV)=\Gamma_V$ and $R_n$ fixes $\Gamma_V$, so by continuity $R_n$ is a projection of
$\widehat\otimes^n_\pi E$ onto $\overline{\Gamma_V}$. For $h\in V$,
$h^{\otimes n}=\sum_{\mathbf t\text{ nondecreasing}}h(t_1)\cdots h(t_n)O_{\mathbf t}\in\Gamma_V$; as
$h\mapsto h^{\otimes n}$ is continuous and $V$ is dense, $\hGamma^n_\pi E\subseteq\overline{\Gamma_V}$. Conversely, by
the argument of Step~1 of the proof of Proposition~\ref{spl:prop:gamma-l1}, with the vectors $e_t$ in place of the unit
vectors, $O_{\mathbf t}$ is a linear combination of pure powers. Hence $\overline{\Gamma_V}=\hGamma^n_\pi E$, and
$(\Gamma^n_E)$ holds.
\end{proof}

\subsection{Boxes in chain sets}\label{ord:ssec:boxes}

Throughout this subsection and the next, $2<r<\infty$ is fixed. For finite $T_1,\dots,T_k\subseteq\Sigma$ the
\emph{chain set} is $\Xi(T_1,\dots,T_k)=\{t\in T_1\times\dots\times T_k:t_1<\dots<t_k\}$, with $\Xi()=\{()\}$. A
\emph{box} is a product $J_1\times\dots\times J_k$ with $J_a\subseteq T_a$, of \emph{cost} $\prod_a\abs{J_a}^{1/r}$;
$\chi(T_1,\dots,T_k)$ is the least total cost of a partition of $\Xi(T_1,\dots,T_k)$ into boxes contained in it. So
$\chi=0$ if the chain set is empty, $\chi()=1$, and $\chi(T_1)=\abs{T_1}^{1/r}$, as $y\mapsto y^{1/r}$ is subadditive.

For $y,y'>0$ let $\psi(y,y')=\min(y,y')^{1-1/r}\max(y,y')^{1/r}$. A \emph{matching} of a finite set $S$ is a set
$\mathfrak m$ of pairwise disjoint two-element subsets of $S$, \emph{near-perfect} if it has $\lfloor\abs S/2\rfloor$
pairs. Let $\operatorname{Mat}(S)$ be the set of near-perfect matchings of $S$ and
$\operatorname{Mat}_q=\operatorname{Mat}(\{1,\dots,q\})$. For positive reals $y=(y_a)_{a\in S}$ put
\[
 \langle y\rangle=\max_{\mathfrak m}\prod_{\{a,b\}\in\mathfrak m}\min(y_a,y_b),\qquad
 \Psi(y)=\prod_{a\in S}y_a^{1/r}\,\langle y\rangle^{1-2/r},
\]
the maximum over all matchings of $S$, empty products being $1$. For a matching $\mathfrak m$ let
$\Psi_{\mathfrak m}(y)$ be the product of $\psi(y_a,y_b)$ over the pairs $\{a,b\}\in\mathfrak m$ and of $y_a^{1/r}$ over
the unmatched $a$; as $\psi(y,y')=(yy')^{1/r}\min(y,y')^{1-2/r}$, it is $\Psi(y)$ with $\langle y\rangle$ replaced by the
product of the minima over $\mathfrak m$. Write $y_{S'}=(y_a)_{a\in S'}$ for $S'\subseteq S$. Let
$\lambda_1=1+(1-2^{-1/r})^{-1}$, $\lambda_2=1+2^{1-2/r}(2^{1-2/r}-1)^{-1}$ and $\lambda=\max(\lambda_1,\lambda_2)$. Four
elementary facts are used.
\begin{enumerate}
\item[(F1)] $\psi(y,y')\le y^{1-1/r}y'^{1/r}$ and $\psi(y,y')\le y^{1/r}y'^{1-1/r}$, with equality in the first if
$y\le y'$: one bound equals $\psi(y,y')$, the other is $(\max/\min)^{1-2/r}$ times larger.
\item[(F2)] $\Psi_{\mathfrak m}(y)\le\Psi(y)$ for every matching $\mathfrak m$ of $S$, as $1-2/r>0$. If all $y_a\ge1$, then
$\Psi(y)=\Psi_{\mathfrak m}(y)$ for some $\mathfrak m\in\operatorname{Mat}(S)$, as every matching extends to a
near-perfect one and the added minima are $\ge1$; so $\Psi(y)\le\sum_{\mathfrak m\in\operatorname{Mat}(S)}\Psi_{\mathfrak m}(y)$.
\item[(F3)] $\Psi$ is nondecreasing in each variable.
\item[(F4)] Let $N\ge1$, $y_{b,i}\ge0$ for $b$ in a finite set $S$ and $1\le i\le N$, $\alpha_b\ge0$ and
$\Theta=\sum_b\alpha_b$. Then, with $0^0=1$,
$\sum_{i=1}^N\prod_by_{b,i}^{\alpha_b}\le N^{(1-\Theta)_+}\prod_b(\sum_{i=1}^Ny_{b,i})^{\alpha_b}$. For $\Theta\le1$ this
is H\"older's inequality, with the constant function $1$ if $\Theta<1$; for $\Theta>1$, apply the case $\Theta=1$ to
the exponents $\alpha_b/\Theta$ and use $\sum_iz_i^\Theta\le(\sum_iz_i)^\Theta$ for $z_i=\prod_by_{b,i}^{\alpha_b/\Theta}$.
\end{enumerate}
The \emph{hull} of a finite nonempty $Y\subseteq\Sigma$ is $\operatorname{hull}(Y)=\{t\in\Sigma:\min Y\le t\le\max Y\}$.

\begin{lemma}[pair tree]\label{ord:lem:pair-tree}
Let $T',T''\subseteq\Sigma$ be finite and nonempty, $m=\min(\abs{T'},\abs{T''})$ and $h=\lceil\log_2m\rceil$; call $T'$
the \emph{tree side} if $\abs{T'}\le\abs{T''}$, and $T''$ otherwise. There are a finite set $\mathcal T$ with a partition
$\mathcal T=\mathcal T_0\sqcup\dots\sqcup\mathcal T_{h-1}\sqcup\mathcal T_\ast$ into \emph{levels}, and nonempty sets
$T'_i\subseteq T'$, $T''_i\subseteq T''$ with $\max T'_i<\min T''_i$ for $i\in\mathcal T$, such that
\begin{equation}\label{ord:eq:pair-partition}
 \{(t',t'')\in T'\times T'':t'<t''\}=\bigsqcup_{i\in\mathcal T}T'_i\times T''_i
\end{equation}
and, if $Y_i$ denotes the one of $T'_i,T''_i$ that lies in the tree side:
\textup{(T1)} $\abs{\mathcal T_d}\le2^d$ for $d<h$, and $\abs{\mathcal T_\ast}\le m$;
\textup{(T2)} $\abs{Y_i}\le m2^{-d}$ for $i\in\mathcal T_d$, and $\abs{Y_i}=1$ for $i\in\mathcal T_\ast$;
\textup{(T3)} within each level, the sets $T'_i$ have pairwise disjoint hulls, and so have the sets $T''_i$.
\end{lemma}

\begin{proof}
\emph{Step 1 (the tree).} Assume first $\abs{T'}\le\abs{T''}$. The nodes are pairs $(U,I)$ of a nonempty set $U$ of
consecutive elements of $T'$ and an interval $I$ of $\Sigma$; the root $(T',\Sigma)$ has depth $0$. If $\abs U\ge2$, let
$U_L$ consist of the first $\lceil\abs U/2\rceil$ elements of $U$ and $U_R$ of the others, $t_U=\max U_L$,
$I_L=\{t\in I:t\le t_U\}$ and $I_R=\{t\in I:t>t_U\}$; the node has the children $(U_L,I_L)$ and $(U_R,I_R)$, of depth
one more, and contributes the pair $(U_L,T''\cap I_R)$ to $\mathcal T_d$, $d$ its depth, if $T''\cap I_R\ne\emptyset$. A
node $(\{t_0\},I)$ is a leaf and contributes $(\{t_0\},\{t\in T''\cap I:t>t_0\})$ to $\mathcal T_\ast$ if this set is
nonempty. Each pair satisfies $\max T'_i<\min T''_i$.

\emph{Step 2 (partition).} Put $\Upsilon(U,I)=\{(t',t'')\in U\times(T''\cap I):t'<t''\}$. If $\abs U\ge2$, then
$\Upsilon(U,I)=U_L\times(T''\cap I_R)\sqcup\Upsilon(U_L,I_L)\sqcup\Upsilon(U_R,I_R)$: for $t'\in U_L$ the cases $t''>t_U$
and $t''\le t_U$ give the first two pieces, and every $t'\in U_R$ exceeds $t_U$. At a leaf, $\Upsilon$ is the contributed
pair or empty. Unwinding from $\Upsilon(T',\Sigma)$ gives~\eqref{ord:eq:pair-partition}.

\emph{Step 3 (counting and hulls).} By induction a node of depth $d$ has $\abs U\le\lceil m/2^d\rceil$, as children
have at most $\lceil\abs U/2\rceil$ elements and $\lceil\lceil y\rceil/2\rceil=\lceil y/2\rceil$; there are at most
$2^d$ such nodes. If $\abs U\ge2$, then $m>2^d$, so $d<h$, and $\abs{U_L}\le\lceil m/2^{d+1}\rceil\le m/2^d$, as
$\lceil y\rceil\le2y$ for $y>1/2$. The leaves have distinct singletons $U\subseteq T'$. This proves~(T1) and~(T2).
Children split the block $U$ and the interval $I$ of their parent into disjoint parts, so two nodes neither of which
descends from the other, such as two nodes of one depth or two leaves, have disjoint blocks $U$, one before the
other, and disjoint intervals $I$. A pair contributed by a node has its first member in $U$ and its second in $I$; this
proves~(T3).

\emph{Step 4 (the other case).} If $\abs{T''}<\abs{T'}$, apply Steps~1--3 to $T''$, $T'$ and the reversed order of
$\Sigma$, which changes no hull.
\end{proof}

\begin{lemma}[level summation]\label{ord:lem:level-sum}
Let $m\ge1$, $h=\lceil\log_2m\rceil$, and let $\mathcal T=\mathcal T_0\sqcup\dots\sqcup\mathcal T_{h-1}\sqcup\mathcal T_\ast$
be a finite set with $\abs{\mathcal T_d}\le2^d$ and $\abs{\mathcal T_\ast}\le m$. Let $g\colon\mathcal T\to[0,\infty)$
satisfy $g\le m2^{-d}$ on $\mathcal T_d$ and $g\le1$ on $\mathcal T_\ast$. For $b$ in a finite set $S$ let
$y_b\colon\mathcal T\to[0,\infty)$ have sum at most $Y_b$ over each $\mathcal T_d$ and over $\mathcal T_\ast$. Let
$\alpha_0\ge1/r$, $\alpha_b\ge0$, $\Theta=\sum_{b\in S}\alpha_b$, $Y=\prod_bY_b^{\alpha_b}$ and
$\Phi=\sum_{i\in\mathcal T}g(i)^{\alpha_0}\prod_by_b(i)^{\alpha_b}$.
\textup{(1)} If $\Theta+\alpha_0\ge1+1/r$, then $\Phi\le\lambda_1m^{\alpha_0}Y$.
\textup{(2)} If $\alpha_0=\Theta=1/r$, then $\Phi\le\lambda_2m^{1-1/r}Y$.
\end{lemma}

\begin{proof}
Put $\zeta=(1-\Theta)_+-\alpha_0$. By~(F4), the part of $\Phi$ over $\mathcal T_d$ is at most
$(m2^{-d})^{\alpha_0}2^{d(1-\Theta)_+}Y=m^{\alpha_0}2^{d\zeta}Y$, and the part over $\mathcal T_\ast$ at most
$m^{(1-\Theta)_+}Y$. (1) Here $\zeta\le-1/r$ and $(1-\Theta)_+\le\alpha_0$, as $(1-\Theta)_+=0$ if $\Theta\ge1$ and
$1-\Theta\le\alpha_0-1/r$ if $\Theta<1$; so $\Phi\le m^{\alpha_0}(\sum_{d\ge0}2^{-d/r}+1)Y=\lambda_1m^{\alpha_0}Y$.
(2) Here $(1-\Theta)_+=1-1/r$ and $\zeta=1-2/r>0$. For $m\ge2$, $2^h<2m$ gives
$\sum_{d<h}2^{d\zeta}\le2^\zeta m^\zeta/(2^\zeta-1)$, and $1/r+\zeta=1-1/r$; for $m=1$ there are no levels
$\mathcal T_d$. Adding the two parts gives $\Phi\le\lambda_2m^{1-1/r}Y$.
\end{proof}

\begin{theorem}[weighted boxes]\label{ord:thm:weighted-box}
Let $\eta_0=\eta_1=1$ and $\eta_k=\lambda\sum_{l=2}^k\abs{\operatorname{Mat}_{l-1}}\,\eta_{l-1}\eta_{k-l}$ for $k\ge2$.
For every $k\ge1$ and all finite nonempty $T_1,\dots,T_k\subseteq\Sigma$ with $m_a=\abs{T_a}$ elements,
\[
 \chi(T_1,\dots,T_k)\le\eta_k\Psi(m_1,\dots,m_k)=\eta_k\prod_{a=1}^km_a^{1/r}\,\langle m\rangle^{1-2/r}.
\]
\end{theorem}

\begin{proof}
We induct on $k$; the cases $k\le1$ are the values of $\chi()$ and $\chi(T_1)$ computed after the definition of $\chi$.
Let $k\ge2$ and $m=(m_1,\dots,m_k)$.

\emph{Step 1 (decomposition).} Apply Lemma~\ref{ord:lem:pair-tree} to $T'=T_1$ and $T''=T_2$. For $i\in\mathcal T$ and
$T\subseteq\Sigma$ let $I_i=\operatorname{hull}(T''_i)$ and $T^{>i}=\{t\in T:t>\max T''_i\}$, and for $2\le l\le k$ let
\[
 \Xi_{i,l}=T'_i\times\Xi(T''_i,T_3\cap I_i,\dots,T_l\cap I_i)\times\Xi(T_{l+1}^{>i},\dots,T_k^{>i}).
\]
These sets partition $\Xi(T_1,\dots,T_k)$: for a chain $t$, $(t_1,t_2)$ lies in exactly one $T'_i\times T''_i$
by~\eqref{ord:eq:pair-partition}, and if $l$ is the largest $a$ with $t_a\le\max T''_i$, then $t_a\in T_a\cap I_i$ for
$3\le a\le l$ and $t_a\in T_a^{>i}$ for $a>l$; conversely, elements of $\Xi_{i,l}$ are chains, as
$\max T'_i<\min T''_i$ and $t_l\le\max T''_i<t_{l+1}$. Products of box partitions of the three factors are box
partitions of $\Xi_{i,l}$ with the product cost, so
\begin{equation}\label{ord:eq:box-split}
 \chi(T_1,\dots,T_k)\le\sum_{l=2}^k\sum_{i\in\mathcal T}\abs{T'_i}^{1/r}\,
 \chi(T''_i,T_3\cap I_i,\dots,T_l\cap I_i)\,\chi(T_{l+1}^{>i},\dots,T_k^{>i}).
\end{equation}

\emph{Step 2 (exponents).} Fix $l$, and let $S_l=\{3,\dots,l\}$ and $S_l^+=\{l+1,\dots,k\}$. By induction and~(F3) the
last factor in~\eqref{ord:eq:box-split} is at most $\eta_{k-l}\Psi(m_{S_l^+})$. Put $m_{2,i}=\abs{T''_i}$ and
$m_{a,i}=\abs{T_a\cap I_i}$ for $a\in S_l$. By induction and~(F2), the middle factor is at most
$\eta_{l-1}\sum_{\mathfrak m}\Psi_{\mathfrak m}(m_{2,i},(m_{a,i})_{a\in S_l})$, the sum over
$\mathfrak m\in\operatorname{Mat}(\{2\}\cup S_l)$, or zero if some $m_{a,i}=0$. Fix $\mathfrak m$. For a pair
$\{a,b\}\in\mathfrak m$ with $m_a\le m_b$, ties broken by the index, let $\alpha_a=1-1/r$ and $\alpha_b=1/r$; let
$\alpha_a=1/r$ for an unmatched $a$. By~(F1),
\begin{equation}\label{ord:eq:box-exponents}
 \Psi_{\mathfrak m}\bigl(m_{2,i},(m_{a,i})_{a\in S_l}\bigr)\le\prod_{a\in\{2\}\cup S_l}m_{a,i}^{\alpha_a},\qquad
 \prod_{a\in\{2\}\cup S_l}m_a^{\alpha_a}=\Psi_{\mathfrak m}(m_{\{2\}\cup S_l}).
\end{equation}
The sum $\Theta_{\mathfrak m}$ of the $\alpha_a$ is at least $1$ if $S_l\ne\emptyset$, since $\mathfrak m$ then has a
pair, and is $1/r$ if $S_l=\emptyset$.

\emph{Step 3 (level sums).} Let $\Phi_{\mathfrak m}=\sum_{i\in\mathcal T}\abs{T'_i}^{1/r}\prod_{a\in\{2\}\cup S_l}m_{a,i}^{\alpha_a}$.
By~(T3) the sums of $\abs{T'_i}$, $m_{2,i}$ and $m_{a,i}$ over each level are at most $m_1$, $m_2$ and $m_a$. Apply
Lemma~\ref{ord:lem:level-sum} with $\min(m_1,m_2)$ in the role of $m$, with $g(i)=\abs{Y_i}$ for the set $Y_i$ of
Lemma~\ref{ord:lem:pair-tree} and $\alpha_0$ its exponent in $\Phi_{\mathfrak m}$, which is $1/r$ or $\alpha_2\ge1/r$,
and with the other factors of $\Phi_{\mathfrak m}$ as the functions $y_b$; (T2) gives the bounds on $g$, and the
exponents add up to $1/r+\Theta_{\mathfrak m}$. If $S_l\ne\emptyset$, part~(1) and~\eqref{ord:eq:box-exponents} give
$\Phi_{\mathfrak m}\le\lambda m_1^{1/r}\Psi_{\mathfrak m}(m_{\{2\}\cup S_l})$; if $S_l=\emptyset$, part~(2) gives
$\Phi_{\mathfrak m}\le\lambda\psi(m_1,m_2)$.

\emph{Step 4 (assembly).} By~(F2), $\Psi(m_{S_l^+})=\Psi_{\mathfrak m^+}(m_{S_l^+})$ for some
$\mathfrak m^+\in\operatorname{Mat}(S_l^+)$. Then $m_1^{1/r}\Psi_{\mathfrak m}(m_{\{2\}\cup S_l})\Psi_{\mathfrak m^+}(m_{S_l^+})$
is $\Psi_{\mathfrak m\cup\mathfrak m^+}(m)$, and $\psi(m_1,m_2)\Psi_{\mathfrak m^+}(m_{S_l^+})$ is
$\Psi_{\{\{1,2\}\}\cup\mathfrak m^+}(m)$; both are at most $\Psi(m)$ by~(F2). As $\{2\}\cup S_l$ has $l-1$ elements, the
$l$-th term of~\eqref{ord:eq:box-split} is at most $\lambda\abs{\operatorname{Mat}_{l-1}}\eta_{l-1}\eta_{k-l}\Psi(m)$;
summing over $l$ proves the claim.
\end{proof}

\subsection{Triangular truncation from \texorpdfstring{$\ell^2$}{l\textasciicircum2} into \texorpdfstring{$\ell^r$}{l\textasciicircum r}}\label{ord:ssec:triangle}

Let $G_r$ be $c_{00}(\Sigma,K)$ with the norm $\norm\cdot_r$; its completion is $\ell^r(\Sigma,\Kh)$. Let $\delta_t$ be the
unit vectors, and for $g_1,\dots,g_k\in c_{00}(\Sigma,K)$ let
$\cP^{(k)}(g_1\otimes\dots\otimes g_k)=\sum_{t_1<\dots<t_k}g_1(t_1)\cdots g_k(t_k)\,\delta_{t_1}\otimes\dots\otimes\delta_{t_k}$,
the $k$\nobreakdash-fold triangular truncation for the family $(\delta_t)$.

\begin{lemma}[flat vectors]\label{ord:lem:flat}
If $g_1,\dots,g_k\in c_{00}(\Sigma,K)$ satisfy $\norm{g_a}_\infty\le1$ and have nonempty supports $T_a$ with $m_a$
elements, then
$\norm{\cP^{(k)}(g_1\otimes\dots\otimes g_k)}_{\otimes^k_\pi G_r}\le\chi(T_1,\dots,T_k)\le\eta_k\Psi(m_1,\dots,m_k)$.
\end{lemma}

\begin{proof}
If boxes $J^i_1\times\dots\times J^i_k$ partition $\Xi(T_1,\dots,T_k)$, then
$\cP^{(k)}(g_1\otimes\dots\otimes g_k)=\sum_i(g_11_{J^i_1})\otimes\dots\otimes(g_k1_{J^i_k})$, as both sides have the
coefficient $g_1(t_1)\cdots g_k(t_k)$ at each chain $t$ and no other terms, and $\norm{g_a1_J}_r\le\abs J^{1/r}$.
Theorem~\ref{ord:thm:weighted-box} gives the second inequality.
\end{proof}

\begin{theorem}[triangular truncation]\label{ord:thm:triangle}
Let $k\ge1$, $\tau=\abs\vartheta$, $\gamma=1/2-1/r$ and
\[
 C_{k,r}=\eta_k\abs{\operatorname{Mat}_k}\Bigl(\frac{2\tau^{-2}}{1-\tau^\gamma}\Bigr)^{\lfloor k/2\rfloor}
 \Bigl(\frac{\tau^{-2/r}}{1-\tau^{2\gamma}}\Bigr)^{k-2\lfloor k/2\rfloor}.
\]
For $x_1,\dots,x_k\in c_{00}(\Sigma,K)$,
\begin{equation}\label{ord:eq:triangle}
 \norm{\cP^{(k)}(x_1\otimes\dots\otimes x_k)}_{\otimes^k_\pi G_r}\le C_{k,r}\tau^{-k}\prod_{a=1}^k\norm{x_a}_2 .
\end{equation}
Consequently $\cP^{(k)}$ extends to a bounded linear map from $\otimes^k_\pi\ell^2(\Sigma,K)$ into
$\widehat\otimes^k_\pi\ell^r(\Sigma,\Kh)$ of norm at most $C_{k,r}\tau^{-k}$.
\end{theorem}

\begin{proof}
\emph{Step 1 (bands).} For $x\in c_{00}(\Sigma,K)$ and $i\in\Z$ let $\Sigma_i(x)=\{t:\tau^{i+1}<\abs{x(t)}\le\tau^i\}$,
$m_i(x)=\abs{\Sigma_i(x)}$ and $x^{(i)}=x1_{\Sigma_i(x)}$. Then $x=\sum_ix^{(i)}$ is a finite sum,
$x^{(i)}=\vartheta^ig$ with $\norm g_\infty\le1$, and $\sum_i\tau^{2i}m_i(x)\le\tau^{-2}\norm x_2^2$. Expanding each $x_a$
in bands and applying Lemma~\ref{ord:lem:flat} and~(F2), the left side of~\eqref{ord:eq:triangle} is at most
\[
 \eta_k\sum_{\mathfrak m\in\operatorname{Mat}_k}\,\prod_{\{a,b\}\in\mathfrak m}\Phi_2(x_a,x_b)\prod_{a\text{ unmatched}}\Phi_1(x_a),
\]
where $\Phi_1(x)=\sum_i\tau^im_i(x)^{1/r}$ and $\Phi_2(x,y)=\sum_{i,l}\tau^{i+l}\psi(m_i(x),m_l(y))$, summed over the
indices with $m_i(x),m_l(y)\ge1$; for fixed $\mathfrak m$ the sum over the bands factors, as each slot occurs in one
factor of $\Psi_{\mathfrak m}$.

\emph{Step 2 (single sums).} If some $x_a=0$ there is nothing to prove. Replacing $x_a$ by $\vartheta^cx_a$, $c\in\Z$,
multiplies both sides of~\eqref{ord:eq:triangle} by $\tau^c$, so we may assume $\tau<\norm{x_a}_2\le1$ for every $a$.
If $\norm x_2\le1$, then $m_i(x)\ge1$ forces $\tau^{i+1}<1$, so $i\ge0$, and $\tau^{2i}m_i(x)\le\tau^{-2}$; hence
$\tau^im_i(x)^{1/r}=(\tau^{2i}m_i(x))^{1/r}\tau^{2\gamma i}\le\tau^{-2/r}\tau^{2\gamma i}$ and
$\Phi_1(x)\le\tau^{-2/r}/(1-\tau^{2\gamma})$.

\emph{Step 3 (pair sums).} Let $\norm x_2,\norm y_2\le1$, $\alpha_i=\tau^im_i(x)^{1/2}$ and $\alpha'_l=\tau^lm_l(y)^{1/2}$,
so that $\sum_i\alpha_i^2$ and $\sum_l\alpha_l'^2$ are at most $\tau^{-2}$. For $b',b''>0$,
$\psi(b',b'')=(b'b'')^{1/2}\min(b'/b'',b''/b')^\gamma\le(b'b'')^{1/2}\min(b'/b'',b''/b')^{\gamma/2}$, as the minimum is
at most~$1$. With $m_i(x)/m_l(y)=(\alpha_i/\alpha'_l)^2\tau^{2(l-i)}$, the term of $\Phi_2(x,y)$ at $(i,l)$ is at most
$\alpha_i^{1+\gamma}\alpha_l'^{1-\gamma}\tau^{\gamma(l-i)}$ if $l\ge i$ and at most
$\alpha_i^{1-\gamma}\alpha_l'^{1+\gamma}\tau^{\gamma(i-l)}$ if $l<i$. For fixed $l-i$, H\"older's inequality with
exponents $2/(1\pm\gamma)$ bounds the sum of these terms by $\tau^{-2}\tau^{\gamma\abs{l-i}}$, so
$\Phi_2(x,y)\le\tau^{-2}(1+\tau^\gamma)/(1-\tau^\gamma)\le2\tau^{-2}/(1-\tau^\gamma)$.

\emph{Step 4 (conclusion).} Each $\mathfrak m\in\operatorname{Mat}_k$ has $\lfloor k/2\rfloor$ pairs and
$k-2\lfloor k/2\rfloor$ unmatched slots, so the left side of~\eqref{ord:eq:triangle} is at most
$C_{k,r}<C_{k,r}\tau^{-k}\prod_a\norm{x_a}_2$. As $c_{00}(\Sigma,K)$ is dense in $\ell^2(\Sigma,K)$, the bounded
$k$\nobreakdash-linear map $(x_1,\dots,x_k)\mapsto\cP^{(k)}(x_1\otimes\dots\otimes x_k)$ extends to $\ell^2(\Sigma,K)^k$
with values in $\widehat\otimes^k_\pi G_r$; its linearization, followed by the contraction induced by the isometry
$G_r\to\ell^r(\Sigma,\Kh)$, is the asserted extension.
\end{proof}

\subsection{Relative ordering splittings}\label{ord:ssec:relative}

\begin{proof}[Proof of Proposition~\ref{spl:prop:r-norm}]
Let $\nu$ be an $r$\nobreakdash-norm on $D$ with $a\nu\le\norm\cdot\le b\nu$, $0<a\le b$. As
$\nu(\sum_tz_t)^r\le\sum_t\nu(z_t)^r$ for finite families, the left side of~\eqref{spl:eq:r-norm} is at most
$b(\sum_t\abs{\xi_t}^r\nu(y_t)^r)^{1/r}\le(b/a)(\sum_t\abs{\xi_t}^r)^{1/r}\max_t\norm{y_t}$. Conversely,
assume~\eqref{spl:eq:r-norm} with constant $C$, and let $\nu(y)$ be the infimum of $(\sum_i\norm{y_i}^r)^{1/r}$ over
the finite decompositions $y=\sum_iy_i$. Then $\nu\le\norm\cdot$;
$\nu(\xi y)=\abs\xi\nu(y)$ for $\xi\in K^\times$, as $y_i\mapsto\xi y_i$ matches the decompositions of $y$ and $\xi y$;
and $\nu(y+z)^r\le\nu(y)^r+\nu(z)^r$, by concatenating decompositions. Let $y=\sum_iy_i$ with all $y_i\ne0$, and
$\xi_i=\vartheta^{c_i}$ with $c_i$ the largest integer such that $\abs\vartheta^{c_i}\ge\norm{y_i}$. Then
$\norm{\xi_i^{-1}y_i}\le1$ and $\abs{\xi_i}<\norm{y_i}/\abs\vartheta$, so~\eqref{spl:eq:r-norm} gives
$\norm y\le C(\sum_i\abs{\xi_i}^r)^{1/r}\le(C/\abs\vartheta)(\sum_i\norm{y_i}^r)^{1/r}$, hence
$\norm y\le(C/\abs\vartheta)\nu(y)$. So $\nu$ vanishes only at $0$; it is a norm, since $r\ge1$, and an equivalent
$r$\nobreakdash-norm.
\end{proof}

\begin{proof}[Proof of Theorem~\ref{spl:thm:relative-ordering}]
Let $C_e=\sup_t\norm{e_t}$, let $C_2$ be the constant of $\ell^2$\nobreakdash-domination, fix a degree $k\ge1$ and
$r\in(2,\infty)$, and let $G_r$, $\delta_t$, $\tau$ and $C_{k,r}$ be as in Subsection~\ref{ord:ssec:triangle}. In
Steps~1 and~3 we assume, as in~(2), that $D$ has an $r$\nobreakdash-norm $\nu$ with $c\nu\le\norm\cdot\le c'\nu$.

\emph{Step 1 (factorization).} For $u\in\cL(E,D)$ let $\tilde u\colon G_r\to D$, $\tilde u(\xi)=\sum_t\xi(t)\,ue_t$. Then
$\nu(\tilde u\xi)^r\le\sum_t\abs{\xi(t)}^r\nu(ue_t)^r\le(C_e\norm u/c)^r\norm\xi_r^r$, so
$\norm{\tilde u}\le(c'/c)C_e\norm u$; $\tilde u$ is linear in $u$, and $\tilde u\delta_t=ue_t$. For $x\in V$,
$\tilde u(wx)=\sum_tx(t)\,ue_t=ux$, so the extension $\tilde u\colon\ell^r(\Sigma,\Kh)\to\widehat D$ satisfies
$\tilde u\circ w=u$, as $V$ is dense; with Step~2, $(\ell^r(\Sigma,\Kh),w,u\mapsto\tilde u)$ is a source-side relative
splitting of degree $k$ (Definition~\ref{cmb:def:decomposition}).

\emph{Step 2 (relative section).} For $x\in V$ the coordinate vector $wx=(x(t))_t$ satisfies
$\norm{wx}_r\le\norm{wx}_2\le C_2\norm x$, so $w$ extends to a bounded map $w\colon E\to\ell^r(\Sigma,\Kh)$. By
Theorem~\ref{ord:thm:triangle} the $k$\nobreakdash-linear map $(x_1,\dots,x_k)\mapsto\cP^{(k)}(wx_1\otimes\dots\otimes wx_k)$
from $V^k$ to $\otimes^k_\pi G_r$ has norm at most $C_{k,r}\tau^{-k}C_2^k$, so it extends to $E^k$ with values in
$\widehat\otimes^k_\pi G_r$. Let $L$ be its linearization and $s=L\circ\alt^k_E$; then
$\norm s\le k!\,C_{k,r}\tau^{-k}C_2^k$ and, for $x_1,\dots,x_k\in V$,
\begin{equation}\label{ord:eq:relative-section}
 s(x_1\wedge\dots\wedge x_k)=\sum_{t_1<\dots<t_k}\det\bigl(x_b(t_a)\bigr)_{a,b}\,\delta_{t_1}\otimes\dots\otimes\delta_{t_k}.
\end{equation}
Expanding $wx_b=\sum_tx_b(t)\delta_t$, dropping repeated indices and sorting shows that $\widehat\wedge^k$ maps the
right side to $wx_1\wedge\dots\wedge wx_k$. The wedges $x_1\wedge\dots\wedge x_k$ with $x_a\in V$ span a dense
subspace of $\Lambda^k_\pi E$, so $\widehat\wedge^k\circ s=\Lambda^k_\pi w$, with $s$ read in
$\widehat\otimes^k_\pi\ell^r(\Sigma,\Kh)$. Hence $w$ is $\mathrm S^k$\nobreakdash-split, which proves~(1). If $E=V$, $s$ takes values in
$\otimes^k_\pi G_r$.

\emph{Step 3 (lift).} Put $Q(u_1,\dots,u_k)=\widehat\wedge^k_D\circ(\tilde u_1\otimes\dots\otimes\tilde u_k)\circ s$, with
$\tilde u_1\otimes\dots\otimes\tilde u_k$ extended to $\widehat\otimes^k_\pi G_r$. Then $Q$ is $k$\nobreakdash-linear,
$\norm{Q(u_1,\dots,u_k)}\le\norm s\,((c'/c)C_e)^k\prod_a\norm{u_a}$, and by~\eqref{ord:eq:relative-section}, for
$x_1,\dots,x_k\in V$, $Q$ is given by the \emph{canonical lift formula}
\begin{equation}\label{ord:eq:canonical-lift}
 Q(u_1,\dots,u_k)(x_1\wedge\dots\wedge x_k)=\sum_{t_1<\dots<t_k}\det\bigl(x_b(t_a)\bigr)_{a,b}\,
 u_1e_{t_1}\wedge\dots\wedge u_ke_{t_k}.
\end{equation}
For $u_1=\dots=u_k=u$ the right side of~\eqref{ord:eq:canonical-lift} is $ux_1\wedge\dots\wedge ux_k$, by expanding
$ux_b=\sum_tx_b(t)\,ue_t$; by density, $Q(u,\dots,u)=\Lambda^k_\pi u$ as maps into $\widehat\Lambda^k_\pi D$. So
$(E,D)$ is good for Banach values in degree $k$ by Lemma~\ref{fib:lem:test}. If $E=V$, then $Q$ takes values in
$\cL(\Lambda^k_\pi E,\Lambda^k_\pi D)$ by Step~2, and $(E,D)$ is good in degree $k$ by the same lemma.
\end{proof}
%% ---- end app-J-ordering.tex
%% ---- begin app-K-forcing.tex
\section{Forcing over Laurent series fields}\label{app:forcing}

Throughout this appendix $K$ satisfies~\textup{(L)}; from Section~\ref{frc:ssec:trace} on, $K$ is dense in $\kappa((X))$,
so $\Kh=\kappa((X))$ and $\abs\vartheta=\varrho$. The results of Sections~\ref{sec:sharpness} and~\ref{sec:classification}
proved from Section~\ref{frc:ssec:trace} on, except Lemma~\ref{obs:lem:single-vector}, which assumes this density, depend
only on $K$ as a normed field; Lemma~\ref{obs:lem:normalization} extends them to every $K$ satisfying~\textup{(L)}.
Theorem~\ref{cls:thm:spreading}(3) and the last assertion of Example~\ref{cls:ex:garling}(3) are proved in
Section~\ref{frc:ssec:spreading} over every nontrivially normed field of characteristic~$2$. Lemma~\ref{fib:lem:coeff}
holds, with the same proof, for every field $\kappa$ of characteristic $p$ in place of the finite field
of~\textup{(L)}; the proof of Proposition~\ref{dp:prop:nonambient} uses it in that generality.

\subsection{Coefficients and costs}\label{frc:ssec:costs}
For a normed space $E$ with $c_{00}(\N,K)\subseteq E\subseteq\linf(\N,K)$ and $\norm x_\infty\le\norm x$, the determinant
array $\Omega$ of \S\ref{ssec:ext}, with $S=\N$, satisfies $\abs{\Omega_z(c)}\le\norm z$, as in the proof of
Lemma~\ref{lem:proj}, so it extends to $\widehat\Lambda^k_\pi E$ with values in $\linf(\N^k,\kappa((X)))$. Likewise
$z(t,t')$ denotes the continuous extension to $E\widehat\otimes_\pi E$ of the entry $x(t)y(t')$ of $x\otimes y$.

For a matrix $\mathsf Z$ over $\kappa$ with rows indexed by a finite set $R$ and columns by a finite set $R'$,
$\operatorname{cost}(\mathsf Z)$ is the least value of
$\sum_i\abs{\operatorname{supp}g_i}^{1/2}\abs{\operatorname{supp}h_i}^{1/2}$ over the representations
$\mathsf Z=\sum_ig_ih_i^{\mathsf T}$ with $g_i\in\kappa^R$ and $h_i\in\kappa^{R'}$; it is subadditive.

\begin{lemma}[tournament lemma]\label{fib:lem:tournament}
Let $\mathsf Z$ be an $m\times m$ matrix over $\kappa$ such that every off-diagonal entry of $\mathsf Z-\mathsf Z^{\mathsf T}$ is nonzero. Then
$\operatorname{cost}(\mathsf Z)\ge m\ln m/(1+\ln(2\abs\kappa))$.
\end{lemma}

\begin{proof}
\emph{Step 1 (parts).} Let $\mathsf Z=\sum_ig_ih_i^{\mathsf T}$ with all $g_i,h_i$ nonzero, let $\mathcal V_i$ be the
union of the supports of $g_i$ and $h_i$, and for $t\in\mathcal V_i$ let the part $\mathcal V_i(t)$ be the set of points
of $\mathcal V_i$ at which $(g_i,h_i)$ is a multiple of $(g_i(t),h_i(t))$. The $(t,t')$ entry
$g_i(t)h_i(t')-h_i(t)g_i(t')$ of $g_ih_i^{\mathsf T}-h_ig_i^{\mathsf T}$ is nonzero exactly when $t,t'\in\mathcal V_i$ lie
in different parts; so every pair $t\ne t'$ lies in some $\mathcal V_i$, in different parts.

\emph{Step 2 (deletion).} Following Hansel's deletion argument~\cite{Ha} with weighted choices, choose independently
for each $i$ one part of $\mathcal V_i$, with probability proportional to its size, and delete the points of $\mathcal V_i$
outside it. At most one point survives, since two survivors would lie in different parts of some $\mathcal V_i$. The
point $t$ survives with probability $\exp(-L_t)$, where
$L_t=\sum_{i:t\in\mathcal V_i}\ln(\abs{\mathcal V_i}/\abs{\mathcal V_i(t)})$; so $\sum_t\exp(-L_t)\le1$, and convexity of
the exponential gives $\sum_tL_t\ge m\ln m$. Moreover $\sum_tL_t=\sum_i\operatorname{ent}_i$ with
$\operatorname{ent}_i=\sum_{t\in\mathcal V_i}\ln(\abs{\mathcal V_i}/\abs{\mathcal V_i(t)})$.

\emph{Step 3 (one term).} Exchanging $g_i$ and $h_i$, which does not change the parts, we may assume that the sizes
$N_1$ and $N_2$ of their supports satisfy $N_1\le N_2$, so $N=\abs{\mathcal V_i}\le2N_2$. It remains to show
$\operatorname{ent}_i\le(1+\ln(2\abs\kappa))\sqrt{N_1N_2}$. The points of $\mathcal V_i$ outside the support of $g_i$ form one
part, which contributes $(N-N_1)\ln(N/(N-N_1))\le N_1$. The other parts divide the support of $g_i$ by the value of
$h_i/g_i$ into at most $\abs\kappa$ sets; since a distribution on at most $\abs\kappa$ points has entropy at most
$\ln\abs\kappa$, they contribute at most $N_1\ln(N/N_1)+N_1\ln\abs\kappa$. Hence
\[
 \operatorname{ent}_i\le N_1\Bigl(1+\ln\frac{2\abs\kappa N_2}{N_1}\Bigr)\le(1+\ln(2\abs\kappa))\sqrt{N_1N_2},
\]
the second inequality with $x=N_2/N_1$, since $x^{-1/2}(1+\ln(2\abs\kappa x))$ is decreasing for $x\ge1$, as
$\ln(2\abs\kappa)>1$.
\end{proof}

\begin{lemma}[constant coefficients]\label{fib:lem:coeff}
Let $R,R'$ be finite sets and $r\in\{2,\infty\}$, let $z$ belong to the completion of
$\ell^r(R,K)\otimes_\pi\ell^r(R',K)$, and let $\coeff_0z\in\kappa^{R\times R'}$ be the matrix of constant coefficients of
its entries.
\begin{enumerate}
\item If $r=2$, then $\operatorname{cost}(\coeff_0z)\le\norm z/(1-\varrho^2)$.
\item If $r=\infty$, then $\rank(\coeff_0z)\le M_\varrho\norm z$.
\end{enumerate}
\end{lemma}

\begin{proof}
For $x\in\ell^r(R,K)$ and $n\in\Z$ let $x_{[n]}\in\kappa^R$ have entries $\coeff_n(x(t))$. Then
$\coeff_0(x\otimes y)=\sum_nx_{[n]}y_{[-n]}^{\mathsf T}$, a finite sum.
(1) The vector $x_{[n]}$ is supported in $O_n(x)=\{t:x(t)\ne0,\ \ord x(t)\le n\}$, and
$\sum_n\abs{O_n(x)}\varrho^{2n}=\norm x^2/(1-\varrho^2)$. By the Cauchy--Schwarz inequality,
$\operatorname{cost}\coeff_0(x\otimes y)\le\sum_n\abs{O_n(x)}^{1/2}\abs{O_{-n}(y)}^{1/2}\le\norm x\norm y/(1-\varrho^2)$.
(2) If $\norm x=\varrho^\nu$ and $\norm y=\varrho^{\nu'}$, only the $l+1$ values $\nu\le n\le-\nu'$ contribute, where
$l=-\nu-\nu'$, and none if $l<0$. Since $l+1\le M_\varrho\varrho^{-l}$, this gives
$\rank\coeff_0(x\otimes y)\le M_\varrho\norm x\norm y$.

For general $z$, write $z=\sum_jx_j\otimes y_j$ as an absolutely convergent series with
$\sum_j\norm{x_j}\norm{y_j}$ as close to $\norm z$ as we like, by approximating $z$ quickly with algebraic tensors and decomposing the
first approximation and the successive differences nearly optimally. The terms with $\norm{x_j}\norm{y_j}<1$ have
entries of absolute value at most $\varrho$, since $\abs{x(t)}\le\norm x$ and absolute values lie in $\varrho^\Z$; by
the nonarchimedean inequality so does their sum, which therefore does not change $\coeff_0$. The finitely many
remaining terms give the bounds with any number larger than $\norm z$ in place of $\norm z$, by subadditivity of cost
and rank, and hence with $\norm z$ itself.
\end{proof}

\begin{lemma}[two-slot slice]\label{fib:lem:slice}
Let $E$ be a normed space with $c_{00}(\N,K)\subseteq E\subseteq\linf(\N,K)$ and $\norm x_\infty\le\norm x$, let
$\gamma_1\ne\gamma_2$ be slots in $\{1,\dots,k\}$, and fix $c_\gamma\in\N$ for the remaining slots $\gamma$. There is a
bounded linear $\Upsilon\colon\widehat\Lambda^k_\pi E\to E\widehat\otimes_\pi E$ of norm at most $k(k-1)$ such that
$(\Upsilon z)(t,t')=\Omega_z(c)$, where $c$ has $t$ in slot $\gamma_1$, $t'$ in slot $\gamma_2$ and $c_\gamma$ in each slot
$\gamma$.
\end{lemma}

\begin{proof}
For $x_1,\dots,x_k\in E$, expansion along the rows $\gamma_1,\gamma_2$ gives
\[
 \det(x_b(c_a))_{a,b}=\sum_{b\ne b'}x_b(c_{\gamma_1})x_{b'}(c_{\gamma_2})\det_{bb'}(x),
\]
where $\det_{bb'}(x)$, the sum of
$\sgn(\sigma)\prod_{a\ne\gamma_1,\gamma_2}x_{\sigma(a)}(c_a)$ over the permutations with $\sigma(\gamma_1)=b$ and
$\sigma(\gamma_2)=b'$, does not involve $c_{\gamma_1},c_{\gamma_2}$, and
$\abs{\det_{bb'}(x)}\le\prod_{b''\ne b,b'}\norm{x_{b''}}$ by the nonarchimedean inequality. So the $k$\nobreakdash-linear
map $x\mapsto\sum_{b\ne b'}\det_{bb'}(x)\,x_b\otimes x_{b'}$ into $E\otimes_\pi E$ has norm at most $k(k-1)$, and its
values have the stated determinants as entries. It is alternating, since these entries vanish when two $x_b$ coincide
and a tensor in $E\otimes E$ with all entries zero is zero. Proposition~\ref{prop:exterior} gives $\Upsilon$ on
$\Lambda^k_\pi E$, and continuity extends it and the entry formula to the completion.
\end{proof}

\subsection{The \texorpdfstring{$\ell^r$}{l\textasciicircum r} threshold}\label{frc:ssec:threshold}
\begin{proof}[Proof of Theorem~\ref{fib:thm:ell2}]
\emph{Step 1 (a coefficient form).} For finitely supported $a_1,\dots,a_k,x_1,\dots,x_k$ with entries in $\F_p$ and
$c\in\N^k$ put
\[
 \Xi(a;x)(c)=\coeff_0\bigl(\Omega_{Q(a_1,\dots,a_k)(x_1\wedge\dots\wedge x_k)}(c)\bigr)\in\kappa,
\]
which is $\F_p$\nobreakdash-multilinear in $(a,x)$ and alternating in $x$ and in $c$. For such $b_1,\dots,b_k,x_1,\dots,x_k$
and $\lambda\in K^k$, with $a_\lambda=\sum_i\lambda_ib_i$, the identity
$Q(a_\lambda,\dots,a_\lambda)(x_1\wedge\dots\wedge x_k)=(a_\lambda x_1)\wedge\dots\wedge(a_\lambda x_k)$ is an equality of
polynomials in $\lambda$ with coefficients in $\widehat\Lambda^k_\pi E$. Comparing the coefficients of
$\lambda_1\cdots\lambda_k$ (Lemma~\ref{lem:idprinc}) and applying $\coeff_0\circ\Omega$ gives
\begin{equation}\label{fib:eq:pol1}
 \sum_{\sigma\in S_k}\Xi(b_{\sigma(1)},\dots,b_{\sigma(k)};x)(c)
 =\sum_{\sigma\in S_k}\det\bigl((b_{\sigma(j)}x_j)(c_i)\bigr)_{i,j},
\end{equation}
since these determinants lie in $\F_p$. The proofs of Lemmas~\ref{lem:H}, \ref{lem:C} and~\ref{lem:D} use only
multilinearity, alternation and~\eqref{fib:eq:pol1} of the map $\Psi$ of Appendix~\ref{app:FF}, and the
expansion~\eqref{eq:expand}. For $\Xi$ the coefficient function
$T(d;y;c)=\Xi(e_{d_1},\dots,e_{d_k};e_{y_1},\dots,e_{y_k})(c)$ takes values in the finite field $\kappa$
and~\eqref{eq:expand} follows from $\F_p$\nobreakdash-multilinearity, the inputs there having entries in $\{0,\pm1\}$; so
those three lemmas hold for $\Xi$. Since $k!=0$ in $\kappa$, Lemma~\ref{lem:D} shows that the cluster value $\chi$ is
not constant; as adjacent rank exchanges generate $S_k$, there are $\tau$ and $\ell$ with
$\delta=\chi(\tau)-\chi(\tau')\ne0$, where $\tau'=(\ell\ \ell{+}1)\circ\tau$ as in Lemma~\ref{lem:S}. Put
$\gamma_1=\tau^{-1}(\ell)$ and $\gamma_2=\tau^{-1}(\ell+1)$.

\emph{Step 2 (the quantitative bound).} Fix $m\ge1$, and let $a=(a_1,\dots,a_k)$ and $x=(x_1,\dots,x_k)$ be the inputs in
the proof of Lemma~\ref{lem:S}, with $\gamma_1,\gamma_2$ as the slots $\alpha,\beta$ there. By Lemma~\ref{lem:C}, the
$m\times m$ matrix $\mathsf Z$ of the values $\Xi(a;x)(c)$ with $A_n(1)$ in slot $\gamma_1$, $B_{n'}(1)$ in slot
$\gamma_2$ and $C_\gamma(1)$ in each other slot $\gamma$ is $\mathsf Z=\chi(\tau')\mathbf 1\mathbf 1^{\mathsf T}+\delta U$.
All entries of $a$ and $x$ lie in $\{0,\pm1\}$, and $x_{\gamma_1}$, $x_{\gamma_2}$ have $4m$ nonzero entries and the
other $x_\gamma$ four, so $\zeta=Q(a_1,\dots,a_k)(x_1\wedge\dots\wedge x_k)$ has $\norm\zeta\le\norm Q(4m)^{2/r}4^{(k-2)/r}$.
The slice $\Upsilon$ of Lemma~\ref{fib:lem:slice} for the slots $\gamma_1,\gamma_2$ and $c_\gamma=C_\gamma(1)$, followed by
restriction of both tensor factors to the coordinates $\{A_n(1)\}$ and $\{B_{n'}(1)\}$, gives $z$ with
$\norm z\le k(k-1)\norm\zeta$ and $\coeff_0z=\mathsf Z$. If $r=2$, Lemma~\ref{fib:lem:coeff}(1) gives
$\operatorname{cost}(\mathsf Z)\le k(k-1)2^km\norm Q/(1-\varrho^2)$, while $\mathsf Z-\mathsf Z^{\mathsf T}=\delta(U-U^{\mathsf T})$
and Lemma~\ref{fib:lem:tournament} give $\operatorname{cost}(\mathsf Z)\ge m\ln m/(1+\ln(2\abs\kappa))$. If $r>2$, the
$\ell^r$ norms dominate the supremum norms, and Lemma~\ref{fib:lem:coeff}(2) gives
$m-1\le\rank\mathsf Z\le M_\varrho k(k-1)\norm\zeta$, the first inequality since $\delta U$ is invertible. This
proves~\eqref{fib:eq:quant}, which fails for large $m$.

\emph{Step 3 (lifts and splittings).} A lift $Q_0$ as in Lemma~\ref{fib:lem:test}, for Banach values, would give such a
$Q$, namely $Q(a_1,\dots,a_k)=Q_0(\mu_{a_1},\dots,\mu_{a_k})$, since $\norm{\mu_a}\le\norm a_\infty$.
Lemma~\ref{fib:lem:test}, Theorem~\ref{thm:duality}(1) and the dichotomy of Proposition~\ref{prop:homogeneous-lift} give
the assertion about $A^k_{E,E;F}$, Proposition~\ref{fib:prop:split} the failure of the splittings, and
Corollary~\ref{cor:discrete-targets} the last assertion, since an equivalent nonarchimedean norm on
$\widehat\Lambda^k_\pi E$ would make $A^k_{E,E;\widehat\Lambda^k_\pi E}$ analytic.
\end{proof}

\begin{proof}[Proof of Corollary~\ref{fib:cor:threshold}]
Theorem~\ref{fib:thm:lower}, with the examples following it, and Theorem~\ref{fib:thm:ell2}.
\end{proof}

\begin{proof}[Proof of Corollary~\ref{fib:cor:pairs}]
The positive cases are Theorem~\ref{fib:thm:lower}. In the negative cases, repeat the proof of
Theorem~\ref{fib:thm:ell2} with $Q(a_1,\dots,a_k)=Q_0(\mu_{a_1},\dots,\mu_{a_k})$ for a lift $Q_0$ on $\cL(E,D)$ as in
Lemma~\ref{fib:lem:test}, with $\widehat\Lambda^k_\pi D$ in place of $\widehat\Lambda^k_\pi E$ and the slice formed in
$D\widehat\otimes_\pi D$. If $r_1\le r_2$, then $\norm{\mu_a}\le\norm a_\infty$ and
$\norm\zeta\le\norm{Q_0}(4m)^{2/r_1}4^{(k-2)/r_1}$. If $r_1>r_2$, H\"older's inequality gives
$\norm{\mu_a}\le\norm a_{r_0}$ with $1/r_0=1/r_2-1/r_1$; as $a_{\gamma_1}$, $a_{\gamma_2}$ have $2m$ nonzero entries and
the other $a_\gamma$ two, $\norm\zeta\le\norm{Q_0}(2m)^{2/r_0}2^{(k-2)/r_0}(4m)^{2/r_1}4^{(k-2)/r_1}$. In both cases
$\norm\zeta\le4^k\norm{Q_0}m^{2/\min(r_1,r_2)}$. If $r_2=2\le r_1$, Lemmas~\ref{fib:lem:coeff}(1)
and~\ref{fib:lem:tournament} give $m\ln m/(1+\ln(2\abs\kappa))\le k(k-1)4^km\norm{Q_0}/(1-\varrho^2)$; if $r_1,r_2>2$,
Lemma~\ref{fib:lem:coeff}(2) gives $m-1\le M_\varrho k(k-1)4^k\norm{Q_0}m^{2/\min(r_1,r_2)}$. Both fail for large $m$.
\end{proof}

\subsection{A trace norm over \texorpdfstring{$\kappa((X))$}{a Laurent series field}}\label{frc:ssec:trace}
Let $\mathcal O=\kappa[[X]]$, the closed unit ball of $\kappa((X))$. For finite sets $R,R'$ we identify
$z\in\kappa((X))^{R\times R'}$ with $\sum_{t,t'}z(t,t')\,e_t\otimes e_{t'}$; its projective norm $\norm z_\pi$ over
$\kappa((X))$ for the supremum norms on $\kappa((X))^R$ and $\kappa((X))^{R'}$ is the infimum of
$\sum_a\norm{g_a}_\infty\norm{h_a}_\infty$ over the finite representations $z=\sum_ag_ah_a^{\mathsf T}$. Fix enumerations $t_1,t_2,\dots$ of $R$ and $t'_1,t'_2,\dots$ of $R'$. Since $\mathcal O$ is
a principal ideal domain whose nonzero elements are units times powers of $X$, the elementary divisor theorem, applied to
$X^cz\in\mathcal O^{R\times R'}$ for a large integer $c$, writes every $z$ as
\begin{equation}\label{frc:eq:smith}
 z=U_1\Delta U_2,\qquad U_1\in GL_R(\mathcal O),\quad U_2\in GL_{R'}(\mathcal O),
\end{equation}
where $\Delta$ has the entries $X^{d_1},\dots,X^{d_{\rank z}}$, with integers $d_1\le\dots\le d_{\rank z}$, at the positions
$(t_1,t'_1),\dots,(t_{\rank z},t'_{\rank z})$ and zeros elsewhere. For $l\in\Z$ let $\rank_l(z)$ be the number of indices
$i$ with $d_i\le l$.

\begin{lemma}[a trace norm]\label{frc:lem:trace-norm}
Let $R,R'$ be finite sets and $z\in\kappa((X))^{R\times R'}$.
\begin{enumerate}
\item The numbers $\rank_l(z)$ do not depend on the factorization~\eqref{frc:eq:smith}.
\item \textup{(Level formula.)} $\norm z_\pi=\sum_i\varrho^{d_i}=(1-\varrho)\sum_{l\in\Z}\varrho^l\rank_l(z)$.
\item \textup{(Submatrices.)} If $z_0$ is the submatrix of $z$ with rows in $S\subseteq R$ and columns in
$S'\subseteq R'$, then $\rank_l(z_0)\le\rank_l(z)$ for every $l\in\Z$.
\item \textup{(Rank at a level.)} If $\abs R=\abs{R'}$, $e\in\Z$, and $z=X^ew$ for a matrix $w\in\mathcal O^{R\times R'}$
whose reduction modulo $X$ is invertible over $\kappa$, then $\rank_e(z)=\abs R$.
\end{enumerate}
\end{lemma}

\begin{proof}
For $l\ge0$ let $\mathcal O_l=\mathcal O/X^{l+1}\mathcal O$, a local ring with residue field $\kappa$, and let
$\operatorname{gen}(Y)$ be the least number of generators of a finitely generated $\mathcal O_l$\nobreakdash-module $Y$.
(a) If an $\mathcal O_l$\nobreakdash-module is generated by $N$ elements, so is each submodule: the module is a quotient of
$\mathcal O^N$, and the preimage of the submodule is free of rank at most $N$, since $\mathcal O$ is a principal ideal
domain. (b) $\operatorname{gen}(Y)=\dim_\kappa Y/XY$ by Nakayama's lemma; so
$\bigoplus_{i=1}^N\mathcal O_l/X^{c_i}\mathcal O_l$ needs exactly $N$ generators when every $c_i\ge1$.

\emph{Step 1 (generators).} Let $w\in\mathcal O^{R\times R'}$, $l\ge0$, and let $\operatorname{Col}_l(w)\subseteq
\mathcal O_l^R$ be the submodule generated by the columns of $w$ reduced modulo $X^{l+1}$. If $w=U_1\Delta U_2$ is a
factorization~\eqref{frc:eq:smith}, then $\Delta=U_1^{-1}wU_2^{-1}$ has entries in $\mathcal O$, so all $d_i\ge0$; we
claim $\operatorname{gen}(\operatorname{Col}_l(w))=\rank_l(w)$. The columns of $w$ and of $U_1\Delta$ generate the same
submodule, since $U_2$ is invertible over $\mathcal O$; multiplication by $U_1$ is an automorphism of $\mathcal O_l^R$; and
the columns of $\Delta$ generate a module isomorphic to $\bigoplus_{d_i\le l}\mathcal O_l/X^{l+1-d_i}\mathcal O_l$, which
needs exactly $\rank_l(w)$ generators by~(b).

\emph{Step 2 (item~(1)).} Choose $c\in\Z$ with $X^cz\in\mathcal O^{R\times R'}$. Multiplying~\eqref{frc:eq:smith} by $X^c$
gives a factorization of $X^cz$ with the exponents $d_i+c\ge0$. So $\rank_l(z)$ is $0$ if $l+c<0$ and
$\operatorname{gen}(\operatorname{Col}_{l+c}(X^cz))$ by Step~1 otherwise; neither depends on the factorization.

\emph{Step 3 (item~(2)).} The two right sides agree, since $(1-\varrho)\sum_{l\ge d}\varrho^l=\varrho^d$. Writing $z$ as
the sum over $i$ of $X^{d_i}$ times the product of the column $t_i$ of $U_1$ and the row $t'_i$ of $U_2$, vectors over
$\mathcal O$, gives a representation of cost at most $\sum_i\varrho^{d_i}$. Conversely, let $z\ne0$ and
$z=\sum_{a=1}^Ng_ah_a^{\mathsf T}$ with all $g_a,h_a$ nonzero. Since $\abs{\kappa((X))^\times}=\varrho^\Z$,
$g_a=X^{\nu_a}g'_a$ and $h_a=X^{\nu'_a}h'_a$ with vectors $g'_a,h'_a$ over $\mathcal O$ of supremum norm $1$; with
$m_a=\nu_a+\nu'_a$ the cost is $\sum_a\varrho^{m_a}$. Multiplying $z$ and every $g_a$ by $X^c$ multiplies the cost and
$\sum_i\varrho^{d_i}$ by $\varrho^c$, so we may assume every $m_a\ge0$; then $z\in\mathcal O^{R\times R'}$. For $l\ge0$,
$z\equiv\sum_{m_a\le l}X^{m_a}g'_ah_a'^{\mathsf T}$ modulo $X^{l+1}$, so $\operatorname{Col}_l(z)$ lies in the submodule
generated by the reductions of the $g'_a$ with $m_a\le l$, and (a) with Step~1 gives
$\rank_l(z)\le N_l=\abs{\{a:m_a\le l\}}$; for $l<0$ both are $0$. Hence
$\sum_a\varrho^{m_a}=(1-\varrho)\sum_l\varrho^lN_l\ge(1-\varrho)\sum_l\varrho^l\rank_l(z)$, and the infimum over
representations gives the lower bound.

\emph{Step 4 (item~(3)).} Multiplying $z$ and $z_0$ by the same $X^c$ replaces $l$ by $l+c$ on both sides, so we may
assume $z\in\mathcal O^{R\times R'}$ and $l\ge0$, both sides vanishing for $l<0$. Deleting columns replaces
$\operatorname{Col}_l(z)$ by a submodule, which needs no more generators by~(a); deleting rows replaces it by its image
under the coordinate projection $\mathcal O_l^R\to\mathcal O_l^S$, which is generated by the images of generators. Step~1 gives $\rank_l(z_0)\le\rank_l(z)$.

\emph{Step 5 (item~(4)).} Let $\Pi_0\in\mathcal O^{R\times R'}$ have the entry $1$ at each position $(t_i,t'_i)$ and $0$
elsewhere. The determinant of $\Pi_0^{-1}w$ reduces to a nonzero element of $\kappa$, so it is a unit and
$\Pi_0^{-1}w\in GL_{R'}(\mathcal O)$. Then $z=I\cdot(X^e\Pi_0)\cdot\Pi_0^{-1}w$ is a factorization~\eqref{frc:eq:smith}
with all $\abs R$ exponents equal to $e$, so $\rank_e(z)=\abs R$.
\end{proof}

\subsection{Coefficient functions of multiplier lifts}\label{frc:ssec:coefficients}
Throughout this subsection $G$ and $D$ are sequence spaces (Definition~\ref{obs:def:sequence-spaces}) with
$\norm{e_t}_G\le1$ and $\norm y_\infty\le\norm y_D$, $\Omega$ is the determinant array on $\widehat\Lambda^k_\pi D$ of
Section~\ref{frc:ssec:costs}, and $\widetilde Q\colon c_{00}(\N,K)^k\to\cL(\Lambda^k_\pi G,\widehat\Lambda^k_\pi D)$ is a
\emph{multiplier lift}: a $k$\nobreakdash-linear map with
\begin{equation}\label{frc:eq:multiplier-lift}
 \widetilde Q(b,\dots,b)(x_1\wedge\dots\wedge x_k)=(bx_1)\wedge\dots\wedge(bx_k)\qquad(b,x_1,\dots,x_k\in c_{00}(\N,K))
\end{equation}
and $\norm{\widetilde Q(e_{a_1},\dots,e_{a_k})}\le\varrho^{-l_0}$ for all $a\in\N^k$ and some integer $l_0\ge0$. For
$l\in\Z$ and $a,y,c\in\N^k$ put
\[
 \Omega_{a;y}(c)=\Omega_{\widetilde Q(e_{a_1},\dots,e_{a_k})(e_{y_1}\wedge\dots\wedge e_{y_k})}(c),\qquad
 T_l(a;y;c)=\coeff_l\Omega_{a;y}(c)\in\kappa .
\]
Since $\abs{\Omega_z(c)}\le\norm z$, $\norm{e_{y_1}\wedge\dots\wedge e_{y_k}}\le1$ and the nonzero absolute values of
$\kappa((X))$ lie in $\varrho^\Z$, a nonzero $\Omega_{a;y}(c)$ has order at least $-l_0$, so $T_l=0$ for $l<-l_0$. For
$g_1,\dots,g_k,h_1,\dots,h_k\in c_{00}(\N,K)$ and $c\in\N^k$, multilinearity gives the finite expansion
\begin{equation}\label{frc:eq:expand}
 \Omega_{\widetilde Q(g_1,\dots,g_k)(h_1\wedge\dots\wedge h_k)}(c)=\sum_{a,y\in\N^k}\prod_{\gamma=1}^kg_\gamma(a_\gamma)
 h_\gamma(y_\gamma)\cdot\Omega_{a;y}(c).
\end{equation}

\begin{lemma}[homogeneity]\label{frc:lem:homogeneity}
For every integer $l_1\ge0$ there is an infinite set $I\subseteq\N$ such that for every $l\le l_1$,
$T_l(\mathbf z)=T_l(\mathbf z')$ whenever $\mathbf z,\mathbf z'\in I^{3k}$ have the same pattern in the sense of Appendix~\ref{app:FF}.
\end{lemma}

\begin{proof}
The map $\mathbf z\mapsto(T_l(\mathbf z))_{-l_0\le l\le l_1}$ takes values in $\kappa^{l_0+l_1+1}$, a finite set whose
size is a power of $p$. Identifying it with a finite field of that size, Lemma~\ref{lem:H} with $N=3k$ gives $I$; for
$l<-l_0$, $T_l=0$.
\end{proof}

For the rest of this subsection fix an integer $l_1\ge0$ and $I$ as in Lemma~\ref{frc:lem:homogeneity}. As in
Appendix~\ref{app:FF}, clusters are the sets $C=\{C(1)<C(2)<C(3)<C(4)\}$ of four consecutive elements of $I$, disjoint
clusters are compared by $C<C'$, and the operator and vector weights $o=(1,0,-1,0)$ and $w=(1,-1,1,-1)$
satisfy~\eqref{eq:W}, with $o$ in place of $s$. Put $o_C=\sum_io_ie_{C(i)}$ and $w_C=\sum_iw_ie_{C(i)}$, so that
$o_Cw_C=e_{C(1)}-e_{C(3)}$. For $l\le l_1$ and pairwise disjoint clusters $C_1,\dots,C_k$ put
\begin{align*}
 \chi_l(C_1,\dots,C_k)&=\coeff_l\Omega_{\widetilde Q(o_{C_1},\dots,o_{C_k})(w_{C_1}\wedge\dots\wedge w_{C_k})}
 \bigl(C_1(1),\dots,C_k(1)\bigr)\\
 &=\sum_{i,i'\in[4]^k}\prod_{\gamma=1}^ko_{i_\gamma}w_{i'_\gamma}\cdot
 T_l\bigl(C_1(i_1),\dots,C_k(i_k);C_1(i'_1),\dots,C_k(i'_k);C_1(1),\dots,C_k(1)\bigr),
\end{align*}
by~\eqref{frc:eq:expand}, the weights being constants in $\{0,\pm1\}$. As in Appendix~\ref{app:FF}, the pattern of each
$3k$\nobreakdash-tuple in the sum is determined by $(i,i')$ and the order pattern $\tau\in S_k$ of $(C_1,\dots,C_k)$, so
by Lemma~\ref{frc:lem:homogeneity} the value depends only on $\tau$; we write it $\chi_l(\tau)$. Also $\chi_l=0$ for
$l<-l_0$.

\begin{lemma}[two-level cancellation]\label{frc:lem:clusters}
Let $\cC_1,\dots,\cC_k$ be finite families of clusters that are pairwise disjoint as sets of clusters, such that all
clusters in $\bigcup_\gamma\cC_\gamma$ are pairwise disjoint, and let $\mathrm{lev}^o_\gamma,\mathrm{lev}^w_\gamma\colon
\cC_\gamma\to\Z_{\ge0}$ be functions. Put
\[
 g_\gamma=\sum_{C\in\cC_\gamma}\vartheta^{\mathrm{lev}^o_\gamma(C)}o_C,\qquad
 h_\gamma=\sum_{C\in\cC_\gamma}\vartheta^{\mathrm{lev}^w_\gamma(C)}w_C\qquad(\gamma=1,\dots,k).
\]
Let $C_\gamma\in\cC_\gamma$ for $\gamma=1,\dots,k$, let $c=(C_1(1),\dots,C_k(1))$, and let
$\lambda=\sum_\gamma\bigl(\mathrm{lev}^o_\gamma(C_\gamma)+\mathrm{lev}^w_\gamma(C_\gamma)\bigr)$. Then for every integer
$l\le l_1$,
\[
 \coeff_l\Omega_{\widetilde Q(g_1,\dots,g_k)(h_1\wedge\dots\wedge h_k)}(c)
 =\sum_{l'\ge0}\coeff_{l'}(\vartheta^\lambda)\,\chi_{l-l'}(C_1,\dots,C_k),
\]
a finite sum.
\end{lemma}

\begin{proof}
\emph{Step 1 (expansion).} Expand by~\eqref{frc:eq:expand}. As the clusters are disjoint and $g_\gamma$, $h_\gamma$ are
supported in the union of the clusters of $\cC_\gamma$, a nonzero term has, in each slot $\gamma$, an operator index
$a_\gamma=\Cop_\gamma(i_\gamma)$ and a vector index $y_\gamma=\Cvec_\gamma(i'_\gamma)$ with
$\Cop_\gamma,\Cvec_\gamma\in\cC_\gamma$, and the coefficient $\prod_\gamma o_{i_\gamma}w_{i'_\gamma}\cdot\vartheta^{\lambda'}$,
where $\lambda'=\sum_\gamma\bigl(\mathrm{lev}^o_\gamma(\Cop_\gamma)+\mathrm{lev}^w_\gamma(\Cvec_\gamma)\bigr)\ge0$ depends
only on the \emph{cluster assignment} $(\Cop_\gamma,\Cvec_\gamma)_\gamma$. Since $\vartheta^{\lambda'}$ has nonnegative
order and $T_{l''}=0$ for $l''<-l_0$,
$\coeff_l(\vartheta^{\lambda'}\Omega_{a;y}(c))=\sum_{l'\ge0}\coeff_{l'}(\vartheta^{\lambda'})\,T_{l-l'}(a;y;c)$, a finite
sum.

\emph{Step 2 (output-free clusters).} Group the terms by cluster assignment and by $l'$. Suppose an assignment has an
occupied cluster $C'$ that is none of $C_1,\dots,C_k$. Then $C'$ belongs to exactly one family $\cC_\gamma$, so only
$a_\gamma$ and $y_\gamma$ can lie in $C'$, and no output index does. Fix $l'$ and all positions outside $C'$, and sum over
the positions inside $C'$. The factor $\coeff_{l'}(\vartheta^{\lambda'})$ and the weights of the other slots are common
factors, and $T_{l-l'}$ is pattern-homogeneous on $I$ since $l-l'\le l_1$. As $C'$ consists of consecutive elements of $I$,
cases (i)--(iii) of the proof of Lemma~\ref{lem:C} apply with $T_{l-l'}$ in place of $T$, and~\eqref{eq:W} makes the
partial sum vanish.

\emph{Step 3 (the surviving assignment).} In a surviving assignment every occupied cluster is an output cluster, and the
only output cluster in $\cC_\gamma$ is $C_\gamma$; so $\Cop_\gamma=\Cvec_\gamma=C_\gamma$ for every $\gamma$,
$\lambda'=\lambda$, and the contribution is $\sum_{l'\ge0}\coeff_{l'}(\vartheta^\lambda)\chi_{l-l'}(C_1,\dots,C_k)$.
\end{proof}

\begin{lemma}[the diagonal]\label{frc:lem:diagonal}
For every integer $l\le l_1$, $\sum_{\tau\in S_k}\chi_l(\tau)$ equals $1$ if $l=0$ and $0$ otherwise.
\end{lemma}

\begin{proof}
\emph{Step 1 (polarization).} Let $b_1,\dots,b_k,x_1,\dots,x_k\in c_{00}(\N,K)$ and $a_\alpha=\sum_i\alpha_ib_i$ for
$\alpha\in K^k$. Both sides of~\eqref{frc:eq:multiplier-lift} with $b=a_\alpha$ are polynomials in $\alpha$, and
comparing the coefficients of $\alpha_1\cdots\alpha_k$ by Lemma~\ref{lem:idprinc} gives
\begin{equation}\label{frc:eq:polarized}
 \sum_{\sigma\in S_k}\widetilde Q(b_{\sigma(1)},\dots,b_{\sigma(k)})(x_1\wedge\dots\wedge x_k)
 =\sum_{\sigma\in S_k}(b_{\sigma(1)}x_1)\wedge\dots\wedge(b_{\sigma(k)}x_k).
\end{equation}

\emph{Step 2 (right side).} Take clusters $C_1<\dots<C_k$, $b_\gamma=o_{C_\gamma}$, $x_\gamma=w_{C_\gamma}$ and
$c=(C_1(1),\dots,C_k(1))$, and apply $\Omega_\cdot(c)$ to~\eqref{frc:eq:polarized}. On the right, a permutation
$\sigma\ne\id$ has a slot $\gamma$ with $\sigma(\gamma)\ne\gamma$, and $o_{C_{\sigma(\gamma)}}w_{C_\gamma}=0$ by
disjointness; the identity contributes $\det\bigl((o_{C_b}w_{C_b})(C_a(1))\bigr)_{a,b}=1$, since
$o_Cw_C=e_{C(1)}-e_{C(3)}$.

\emph{Step 3 (left side).} On the left, the sign manipulation in the proof of Lemma~\ref{lem:D}, which uses only that the
wedge and $\Omega_z$ are alternating, identifies the summand of $\sigma$ with
\[
 \Omega_{\widetilde Q(o_{C_{\sigma(1)}},\dots,o_{C_{\sigma(k)}})(w_{C_{\sigma(1)}}\wedge\dots\wedge w_{C_{\sigma(k)}})}
 \bigl(C_{\sigma(1)}(1),\dots,C_{\sigma(k)}(1)\bigr),
\]
whose $l$th coefficient is $\chi_l(\sigma)$, since $(C_{\sigma(1)},\dots,C_{\sigma(k)})$ has order pattern $\sigma$.
Taking $\coeff_l$ gives $\sum_\sigma\chi_l(\sigma)=\coeff_l(1)$.
\end{proof}

\begin{corollary}[an adjacent transposition]\label{frc:cor:transposition}
There are $\tau\in S_k$ and $\ell\in\{1,\dots,k-1\}$ with $\chi_0(\tau)\ne\chi_0(\tau')$, where $\tau'=(\ell\ \ell{+}1)\circ\tau$.
For such $\tau,\ell$ put $\gamma_1=\tau^{-1}(\ell)$, $\gamma_2=\tau^{-1}(\ell+1)$, $\delta_l=\chi_l(\tau)-\chi_l(\tau')$ for
$l\le l_1$, and let $l_*$ be the least $l\le l_1$ with $\delta_l\ne0$. Then $-l_0\le l_*\le0$.
\end{corollary}

\begin{proof}
Otherwise $\chi_0$ is invariant under the adjacent transpositions, which generate $S_k$, hence constant, and
Lemma~\ref{frc:lem:diagonal} with $l=0$ gives $k!\,\chi_0=1$, impossible since $k!=0$ in $\kappa$. As $\delta_0\ne0$ and
$\delta_l=0$ for $l<-l_0$, $-l_0\le l_*\le0$.
\end{proof}

\subsection{The product triangle bound}\label{frc:ssec:ptb}
For a symmetric sequence space $D$ and finite $J\subseteq\N$, $D_J$ is the span of the $e_t$, $t\in J$; we use three facts.
\begin{enumerate}
\item[(P1)] For finite $J,J'\subseteq\N$, a tensor in $D_J\otimes D_{J'}$ has the same projective norm in
$D_J\otimes_\pi D_{J'}$ as in $D\otimes_\pi D$, since multiplication by $1_J$ is a contractive projection of $D$ onto $D_J$
(Lemma~\ref{obs:lem:multiplier}(1)).
\item[(P2)] Relabeling the rows and the columns of finitely supported tensors by injective maps $\N\to\N$, and
multiplying rows and columns by sequences in $K$ of absolute value $1$, are isometries for the norm of $D\otimes_\pi D$:
by symmetry of the norm they map each $D_J$ isometrically onto some $D_{J''}$, and (P1) applies.
\item[(P3)] For finite $J,J'\subseteq\N$, the completion of $D_J\otimes_\pi D_{J'}$ is $\Kh^{J\times J'}$, with a norm
between the maximum and the sum of the absolute values of the entries, since $\norm\cdot_\infty\le\norm\cdot_D$ and
$\norm{e_t}_D=1$; as $K$ is dense in $\Kh$ and this norm is continuous for the maximum norm, $\norm{\lambda z}=\abs\lambda\norm z$
for $\lambda\in\Kh$. Multiplication by $1_J\otimes1_{J'}$ extends to a contraction of $D\widehat\otimes_\pi D$ onto this
completion whose value at $z$ is the restriction of $z$ to $J\times J'$ (Section~\ref{frc:ssec:costs}).
\end{enumerate}

\begin{proof}[Proof of Theorem~\ref{obs:thm:ptb}]
Let $C_Q=\norm Q$, $l_0=\max(0,\lceil\log_{1/\varrho}C_Q\rceil)$ and $C_\star=64\cdot8^{k-2}k(k-1)\,C_Q+\varrho^{-l_0}+1$.
We show \textup{(PTB)} with constant $C_\star$. For general $K$ satisfying~\textup{(L)},
Lemma~\ref{obs:lem:normalization} then gives \textup{(PTB)} with $\abs\vartheta$ in place of $\varrho$ in $C_\star$ and $l_0$.

\emph{Step 0 (a multiplier lift).} For $a_1,\dots,a_k\in c_{00}(\N,K)$ put
$\widetilde Q(a_1,\dots,a_k)=Q(\mu_{a_1},\dots,\mu_{a_k})$, a $k$\nobreakdash-linear map with
\begin{equation}\label{frc:eq:mult-bound}
 \norm{\widetilde Q(a_1,\dots,a_k)}\le C_Q\prod_{\gamma=1}^k\norm{\mu_{a_\gamma}}_{E\to D}.
\end{equation}
Since $Q(u,\dots,u)=\Lambda^k_\pi u$, $\widetilde Q$ satisfies~\eqref{frc:eq:multiplier-lift}; since $\norm{e_t}_E=1$,
$\norm y_\infty\le\norm y_D$ and $\norm{\mu_{e_t}}_{E\to D}=1$ (Lemma~\ref{obs:lem:multiplier}(1),(3)), it is a multiplier
lift with $G=E$, as $\norm{\widetilde Q(e_{a_1},\dots,e_{a_k})}\le C_Q\le\varrho^{-l_0}$. The rest of the proof uses $Q$
only through these properties of $\widetilde Q$ and~\eqref{frc:eq:mult-bound}.

\emph{Step 1 (configuration).} Assume first:
\begin{enumerate}
\item[(A)] $\operatorname{supp}a=\operatorname{supp}x=J_1$ and $\operatorname{supp}b=\operatorname{supp}y=J_2$ with
$J_1\cap J_2=\emptyset$, and every nonzero entry of $a$, $x$, $b$ and $y$ is a nonnegative power of $\vartheta$.
\end{enumerate}
If one of the four vectors is $0$, both sides of~\eqref{obs:eq:ptb} vanish; so let them be nonzero and put
$N_0=\norm{\mu_a}_{E\to D}\norm x_E\norm{\mu_b}_{E\to D}\norm y_E>0$. List $J_1\cup J_2$ as $f_1<\dots<f_m$, and let
$l_1\ge0$ be the least integer with $m^2\varrho^{l_1+1}\le N_0$. Take $I$ as in Lemma~\ref{frc:lem:homogeneity} for this
$l_1$, and $\tau$, $\ell$, $\gamma_1$, $\gamma_2$, $l_*$ as in Corollary~\ref{frc:cor:transposition}. Choose pairwise
disjoint clusters in increasing order: a cluster $C_\gamma$ for each slot $\gamma$ with $\tau(\gamma)<\ell$, ordered as
$\tau$ prescribes; then clusters $Y_1<\dots<Y_m$; then a cluster $C_\gamma$ for each slot $\gamma$ with
$\tau(\gamma)>\ell+1$, ordered as $\tau$ prescribes. Let
\begin{align*}
 \cC_{\gamma_1}&=\{Y_i:f_i\in J_1\},& \mathrm{lev}^o_{\gamma_1}(Y_i)&=\ord a(f_i),& \mathrm{lev}^w_{\gamma_1}(Y_i)&=\ord x(f_i),\\
 \cC_{\gamma_2}&=\{Y_i:f_i\in J_2\},& \mathrm{lev}^o_{\gamma_2}(Y_i)&=\ord b(f_i),& \mathrm{lev}^w_{\gamma_2}(Y_i)&=\ord y(f_i),
\end{align*}
and $\cC_\gamma=\{C_\gamma\}$ with both levels $0$ for $\gamma\ne\gamma_1,\gamma_2$; these families are pairwise disjoint
since $J_1\cap J_2=\emptyset$. Let $g_\gamma$ and $h_\gamma$ be as in Lemma~\ref{frc:lem:clusters}, and
$\zeta=\widetilde Q(g_1,\dots,g_k)(h_1\wedge\dots\wedge h_k)\in\widehat\Lambda^k_\pi D$.

\emph{Step 2 (a bounded slice).} As $\vartheta$ has order $1$, (A) gives $g_{\gamma_1}=g'-g''$ with
$g'=\sum_{f_i\in J_1}a(f_i)e_{Y_i(1)}$ and $g''$ the same sum with $Y_i(3)$ in place of $Y_i(1)$, two relabelings of $a$;
so $\norm{\mu_{g_{\gamma_1}}}_{E\to D}\le2\norm{\mu_a}_{E\to D}$ by Lemma~\ref{obs:lem:multiplier}(3). Likewise
$\norm{\mu_{g_{\gamma_2}}}_{E\to D}\le2\norm{\mu_b}_{E\to D}$ and $\norm{\mu_{g_\gamma}}_{E\to D}\le2$ for the other slots,
and since $h_\gamma$ is a combination, with the coefficients $w_1,\dots,w_4$, of four relabelings of $x$, of $y$ or of a
unit vector,
$\norm{h_{\gamma_1}}_E\le4\norm x_E$, $\norm{h_{\gamma_2}}_E\le4\norm y_E$ and $\norm{h_\gamma}_E\le4$. So
\eqref{frc:eq:mult-bound} gives $\norm\zeta\le64\cdot8^{k-2}C_QN_0$. Let $\Upsilon\colon\widehat\Lambda^k_\pi D\to
D\widehat\otimes_\pi D$ be the slice of Lemma~\ref{fib:lem:slice} for the slots $\gamma_1,\gamma_2$ and
$c_\gamma=C_\gamma(1)$, let $R=\{Y_i(1):f_i\in J_1\}$ and $R'=\{Y_i(1):f_i\in J_2\}$, and let $z\in\Kh^{R\times R'}$ be the
restriction of $\Upsilon\zeta$ to $R\times R'$ (fact (P3)). Then, in the completion of $D_R\otimes_\pi D_{R'}$,
\begin{equation}\label{frc:eq:slice}
 \norm z\le64\cdot8^{k-2}k(k-1)\,C_QN_0 .
\end{equation}

\emph{Step 3 (reading all coefficients).} Let $t=Y_i(1)\in R$ and $t'=Y_{i'}(1)\in R'$. By Lemma~\ref{fib:lem:slice},
$z(t,t')=\Omega_\zeta(c)$, where $c$ has $Y_i(1)$ in slot $\gamma_1$, $Y_{i'}(1)$ in slot $\gamma_2$ and $C_\gamma(1)$ in
every other slot $\gamma$. The clusters $C_\gamma$ lie below or above all of $Y_1,\dots,Y_m$ as $\tau$ prescribes, so the
output clusters have order pattern $\sigma=\tau$ if $f_i<f_{i'}$ and $\sigma=\tau'$ if $f_i>f_{i'}$. The level sum of
Lemma~\ref{frc:lem:clusters} is $\lambda=\ord a(f_i)+\ord x(f_i)+\ord b(f_{i'})+\ord y(f_{i'})$, and
$\vartheta^\lambda=(ax)(f_i)\,(by)(f_{i'})$ by~(A). For $\sigma'\in S_k$ put
$\hat\chi(\sigma')=\sum_{l=-l_0}^{l_1}\chi_l(\sigma')X^l\in\kappa((X))$. For $l\le l_1$,
$\coeff_l(\vartheta^\lambda\hat\chi(\sigma))=\sum_{l'\ge0}\coeff_{l'}(\vartheta^\lambda)\chi_{l-l'}(\sigma)$, since
$\vartheta^\lambda$ has nonnegative order and $\chi_{l''}=0$ for $l''<-l_0$. So by Lemma~\ref{frc:lem:clusters},
$z(t,t')-\vartheta^\lambda\hat\chi(\sigma)$ has no nonzero coefficient of index at most $l_1$, and
\[
 \abs{z(t,t')-(ax)(f_i)\,(by)(f_{i'})\,\hat\chi(\sigma)}\le\varrho^{l_1+1}.
\]
The relabeling $Y_i(1)\mapsto f_i$ of rows and columns is increasing on $R\cup R'$, and by (P2) and continuity it is an
isometry onto the completion of $D_{J_1}\otimes_\pi D_{J_2}$; let $z^\sharp$ be the image of $z$. With $\tilde a=ax$ and
$\tilde b=by$, as matrices indexed by $J_1\times J_2$,
\begin{equation}\label{frc:eq:reading}
 z^\sharp=\hat\chi(\tau')\,(\tilde a\otimes\tilde b)+\hat\delta\cdot\cP_<(\tilde a\otimes\tilde b)+z_1,\qquad
 \hat\delta=\hat\chi(\tau)-\hat\chi(\tau')=\sum_{l=-l_0}^{l_1}\delta_lX^l,
\end{equation}
where $z_1\in\Kh^{J_1\times J_2}$ has entries of absolute value at most $\varrho^{l_1+1}$. The lowest nonzero coefficient
of $\hat\delta$ is $\delta_{l_*}$ with $l_*\le0$, so $\abs{\hat\delta}=\varrho^{l_*}\ge1$, and $\abs{\hat\chi(\tau')}\le\varrho^{-l_0}$.

\emph{Step 4 (conclusion under~(A)).} In the completion of $D_{J_1}\otimes_\pi D_{J_2}$, a $\Kh$\nobreakdash-normed space
by (P3), $\norm{\tilde a\otimes\tilde b}\le\norm{ax}_D\norm{by}_D\le N_0$ by Lemma~\ref{obs:lem:multiplier}(2),
$\norm{z_1}\le m^2\varrho^{l_1+1}\le N_0$ by (P3) and the choice of $l_1$, and $\norm{z^\sharp}=\norm z$ is bounded
by~\eqref{frc:eq:slice}. By~\eqref{frc:eq:reading},
\[
 \norm{\cP_<(\tilde a\otimes\tilde b)}=\abs{\hat\delta}^{-1}\norm{z^\sharp-\hat\chi(\tau')(\tilde a\otimes\tilde b)-z_1}
 \le\bigl(64\cdot8^{k-2}k(k-1)\,C_Q+\varrho^{-l_0}+1\bigr)N_0=C_\star N_0 .
\]
By (P1), the left side is $\norm{\cP_<((ax)\otimes(by))}_{D\otimes_\pi D}$. The constant does not depend on $I$, $\tau$,
$\ell$ or $l_1$, which were chosen after $a,x,b,y$.

\emph{Step 5 (removing~(A)).} Both sides of~\eqref{obs:eq:ptb} are homogeneous of degree one in each of $a,x,b,y$. The
changes (a) and (b) below do not increase $\norm{\mu_a}_{E\to D}$, $\norm x_E$, $\norm{\mu_b}_{E\to D}$ or $\norm y_E$
(Lemma~\ref{obs:lem:multiplier}(1),(3) and symmetry of the norms), and they keep $\cP_<((ax)\otimes(by))$ or change it by an
isometry of (P2).
(a)~\emph{Supports.} Replace $(a,x)$ by $(a1_J,x1_J)$ with $J=\operatorname{supp}(ax)$, and likewise $(b,y)$; then
$\operatorname{supp}a=\operatorname{supp}x=J_1$ and $\operatorname{supp}b=\operatorname{supp}y=J_2$. If
$J_1\cap J_2\ne\emptyset$, list $J_1\cup J_2$ as $f_1<\dots<f_m$ and move the entries of $a$ and $x$ at $f_i\in J_1$ to
the position $2i$ and those of $b$ and $y$ at $f_i\in J_2$ to the position $2i-1$; since $2i<2i'-1$ exactly when
$i<i'$, this relabels the rows and columns of $\cP_<((ax)\otimes(by))$, and the new supports are disjoint.
(b)~\emph{Units.} Since $\abs{K^\times}=\abs\vartheta^\Z$, each nonzero entry of $a$ is
$a(t)=\theta_a(t)\vartheta^{\ord a(t)}$ with $\abs{\theta_a(t)}=1$. Replace $a$ by the vector $a'$ with the entries
$\vartheta^{\ord a(t)}$ on $\operatorname{supp}a$ and $0$ elsewhere, and likewise $x,b,y$; this multiplies the rows of
$\cP_<((ax)\otimes(by))$ by $(\theta_a\theta_x)^{-1}$ and its columns by $(\theta_b\theta_y)^{-1}$.
(c)~\emph{Levels.} By homogeneity we may multiply $a,x,b,y$ by powers of $\vartheta$ that make all exponents
nonnegative.
After (a)--(c), (A) holds, and Steps~1--4 give~\eqref{obs:eq:ptb} with constant $C_\star$.
\end{proof}

\begin{proof}[Proof of Lemma~\ref{obs:lem:single-vector}]
We may assume $x\ne0$. Let $R=\operatorname{supp}x$, and let $z\in K^{R\times R}$ have the entries $z(t,t')=x(t)x(t')$
for $t<t'$ and $0$ otherwise; $\norm\cdot_\pi$ is the norm of Section~\ref{frc:ssec:trace}.

\emph{Step 1 (reduction to matrices).} Restricting to $R$ the vectors of a representation
$\cP_<(x\otimes x)=\sum_ag_a\otimes h_a$ with $g_a,h_a\in D$ gives a representation
$z=\sum_a(g_a|_R)(h_a|_R)^{\mathsf T}$ with $\norm{g_a|_R}_\infty\le\norm{g_a}_D$ and $\norm{h_a|_R}_\infty\le\norm{h_a}_D$.
Hence $\norm z_\pi\le\norm{\cP_<(x\otimes x)}_{D\otimes_\pi D}$.

\emph{Step 2 (normalization).} Write $\norm x_\infty=\varrho^c$ with $c\in\Z$ and put $x'=\vartheta^{-c}x$; the matrix
$z'=\vartheta^{-2c}z$ is formed from $x'$ in the same way, and $\norm z_\pi=\varrho^{2c}\norm{z'}_\pi$. Since
$\abs{K^\times}=\abs\vartheta^\Z$, each nonzero entry is $x'(t)=\theta_t\vartheta^{\lambda_t}$ with an integer $\lambda_t\ge0$
and $\abs{\theta_t}=1$, so
$z'\in\mathcal O^{R\times R}$. For $i\ge0$ let $m_i$ be the number of $t\in R$ with $\lambda_t=i$.

\emph{Step 3 (one level).} Fix $i$ with $m_i\ge2$, and let $t_1<\dots<t_{m_i}$ be the $t\in R$ with $\lambda_t=i$. The
submatrix $z'_i$ of $z'$ with rows $t_1,\dots,t_{m_i-1}$ and columns $t_2,\dots,t_{m_i}$ has the entry
$\theta_{t_a}\theta_{t_{b+1}}\vartheta^{2i}$ in position $(a,b)$ if $a\le b$, and $0$ otherwise. As $\vartheta=X\eta_0$
with a unit $\eta_0$ of $\mathcal O$, $z'_i=X^{2i}\eta_0^{2i}\Theta_1U\Theta_2$, where $\Theta_1$ and $\Theta_2$ are
diagonal with the units $\theta_{t_a}$ and $\theta_{t_{b+1}}$ and $U$ is the upper unitriangular
$(m_i-1)\times(m_i-1)$ matrix of ones; the reduction of $\eta_0^{2i}\Theta_1U\Theta_2$ modulo $X$ is invertible.
Lemma~\ref{frc:lem:trace-norm}(4),(3) gives $\rank_{2i}(z')\ge\rank_{2i}(z'_i)=m_i-1$.

\emph{Step 4 (conclusion).} The indices $2i$ are distinct, so the level formula of Lemma~\ref{frc:lem:trace-norm}(2) gives
$\norm{z'}_\pi\ge(1-\varrho)\sum_{i:m_i\ge1}\varrho^{2i}(m_i-1)$, and
\[
 \norm{x'}_2^2=\sum_{i:m_i\ge1}(m_i-1)\varrho^{2i}+\sum_{i:m_i\ge1}\varrho^{2i}
 \le\frac{\norm{z'}_\pi}{1-\varrho}+\frac1{1-\varrho^2}.
\]
Multiplying by $\varrho^{2c}=\norm x_\infty^2$ and using Step~1 gives the lemma.
\end{proof}

\begin{proof}[Proof of Proposition~\ref{obs:prop:l2-domination}]
For $m\ge1$ and $1\le r'\le\infty$ put
\[
 \psi_{r'}(m)=\sup\{\norm\xi_{r'}/\norm\xi_E:\xi\in c_{00}(\N,K)\setminus\{0\},\ \abs{\operatorname{supp}\xi}\le m\},
\]
which is finite, since $\norm\xi_{r'}\le m^{1/r'}\norm\xi_\infty\le m^{1/r'}\norm\xi_E$ by
Lemma~\ref{obs:lem:multiplier}(1); and $\psi_\infty(m)\le1$.

\emph{Step 1 (multipliers of indicators).} For finite $J\subseteq\N$ and $x\in E$, $\mu_{1_J}x=x1_J$ and
$\norm{x1_J}_D\le C'\norm{x1_J}_r\le C'\psi_r(\abs J)\norm{x1_J}_E\le C'\psi_r(\abs J)\norm x_E$, so
$\norm{\mu_{1_J}}_{E\to D}\le C'\psi_r(\abs J)$.

\emph{Step 2 (one inequality for each support size).} Let $x\in c_{00}(\N,K)\setminus\{0\}$ with $J=\operatorname{supp}x$
of size at most $m$. \textup{(PTB)} with $a=b=1_J$ and $y=x$, and Step~1, give
$\norm{\cP_<(x\otimes x)}_{D\otimes_\pi D}\le CC'^2\psi_r(m)^2\norm x_E^2$; Lemma~\ref{obs:lem:single-vector} and
$\norm x_\infty\le\norm x_E$ give $\norm x_2^2\le(\mathfrak a\psi_r(m)^2+\mathfrak b)\norm x_E^2$. So
\begin{equation}\label{frc:eq:support-size}
 \psi_2(m)^2\le\mathfrak a\,\psi_r(m)^2+\mathfrak b\qquad(m\ge1).
\end{equation}

\emph{Step 3 (interpolation).} If $r=\infty$, then $\psi_\infty\le1$ and~\eqref{frc:eq:support-size} give
$\psi_2(m)^2\le\mathfrak a+\mathfrak b=C_2^2$. Let $r<\infty$. For $\xi$ with $\abs{\operatorname{supp}\xi}\le m$,
$\norm\xi_r^r\le\norm\xi_\infty^{r-2}\norm\xi_2^2\le\psi_2(m)^2\norm\xi_E^r$, so $\psi_r(m)^2\le\psi_2(m)^{4/r}$
and~\eqref{frc:eq:support-size} gives $\psi_2(m)^2\le\mathfrak a\,\psi_2(m)^{4/r}+\mathfrak b$. If
$\psi_2(m)^2>2\mathfrak b$, then $\psi_2(m)^2<2\mathfrak a\,\psi_2(m)^{4/r}$, so $\psi_2(m)^2<(2\mathfrak a)^{r/(r-2)}$.
In both cases $\psi_2(m)^2\le C_2^2$.

\emph{Step 4 (all of $E$).} So $\norm x_2\le C_2\norm x_E$ for $x\in c_{00}(\N,K)$, and for $x\in E$,
$\norm x_2=\sup_n\norm{x1_{[n]}}_2\le C_2\sup_n\norm{x1_{[n]}}_E=C_2\norm x_E$ by
Definition~\ref{obs:def:sequence-spaces}.
\end{proof}

\begin{proof}[Proof of Proposition~\ref{obs:prop:gamma-ptb}]
Write $c_0=c_0(\N,K)$, a symmetric sequence space, and let $g_k(E)$, finite by $(\Gamma^k_E)$, be the infimum of the norms
of projections of $\widehat\otimes^k_\pi E$ onto $\widehat\Gamma^k_\pi E$ (Definition~\ref{ob:def:sequences}).

\emph{Step 1 (a multiplier lift for $(c_0,E)$).} Let $Z=\cL(\widehat\otimes^k_\pi c_0,\widehat\Lambda^k_\pi E)$ and let
$W\subseteq Z$ be the closed subspace of operators vanishing on $\cN^k_{c_0}=\ker\widehat\wedge^k_{c_0}$. Since
$\norm{\mu_a}_{c_0\to E}=\norm a_E$ (Lemma~\ref{obs:lem:multiplier}(4)), $c_{00}(\N,K)$ is dense in $E$ and $Z$ is a Banach space, the map
$\Phi(a_1,\dots,a_k)=\widehat\wedge^k_E\circ(\mu_{a_1}\otimes\dots\otimes\mu_{a_k})$ extends from $c_{00}(\N,K)^k$ to a
$k$\nobreakdash-linear $\Phi\colon E^k\to Z$ of norm at most $1$. Its diagonal lies in $W$, since
$\Phi(a,\dots,a)=\widehat\Lambda^k_\pi\mu_a\circ\widehat\wedge^k_{c_0}$ for $a\in c_{00}(\N,K)$ by naturality of the wedge
and $W$ is closed. Lemma~\ref{ob:lem:parameter}(3), with $P=E$, $n=k$ and $\delta=1$, gives a bounded
$k$\nobreakdash-linear $\Phi'\colon E^k\to W$ with the diagonal of $\Phi$ and $\norm{\Phi'}\le g_k(E)+1$. By
Lemma~\ref{ob:lem:metric-surjection}, $T\mapsto T\circ\widehat\wedge^k_{c_0}$ is a linear isometry of
$\cL(\widehat\Lambda^k_\pi c_0,\widehat\Lambda^k_\pi E)$ onto $W$. For $a_1,\dots,a_k\in c_{00}(\N,K)$ let
$\widetilde Q(a_1,\dots,a_k)$ be the restriction to $\Lambda^k_\pi c_0$ of the preimage of $\Phi'(a_1,\dots,a_k)$. Then
$\widetilde Q$ is $k$\nobreakdash-linear, $\widetilde Q(a,\dots,a)=\Lambda^k_\pi\mu_a$, and
\[
 \norm{\widetilde Q(a_1,\dots,a_k)}\le(g_k(E)+1)\prod_{\gamma=1}^k\norm{\mu_{a_\gamma}}_{c_0\to E}.
\]

\emph{Step 2 (the triangle bound).} By Lemma~\ref{obs:lem:multiplier}(1),(3), $\widetilde Q$ is a multiplier lift with
$G=c_0$, target $E$ and $l_0=\max(0,\lceil\log_{1/\varrho}(g_k(E)+1)\rceil)$, and it satisfies~\eqref{frc:eq:mult-bound}
for $(c_0,E)$ with $C_Q=g_k(E)+1$. So Steps~1--5 of the proof of Theorem~\ref{obs:thm:ptb} give \textup{(PTB)} for
$(c_0,E)$, and the choice for $G=c_0$ in Step~1 of the proof of Theorem~\ref{cls:thm:symmetric} bounds $\cP_<$ on
$E\otimes_\pi E$.
\end{proof}

\subsection{Nonarchimedean fibers}\label{frc:ssec:nonarch}
\begin{proof}[Proof of Lemma~\ref{fib:lem:marking}]
\emph{Step 1 (a normalized basis).} Extend an equivalent nonarchimedean norm of $E$ to $\widehat E$ and round it up to
$\varrho^\Z\cup\{0\}$, as in the proof of Corollary~\ref{cor:discrete-targets}; this gives $\norm\cdot'$. We construct
$f_1,f_2,\dots\in\widehat E$ with $\norm{\sum_n\xi_nf_n}'=\max_n\abs{\xi_n}$. Given $f_1,\dots,f_n$, their span $V_n$ is
closed, being finite-dimensional over the complete field $\Kh$ (Lemma~\ref{lem:findim}). For $v\in\widehat E\setminus V_n$
the distances $\norm{v-w}'$, $w\in V_n$, lie in $\varrho^\Z$ and are at least $\operatorname{dist}(v,V_n)>0$, so some
$w_0$ attains their minimum, and the nonarchimedean inequality gives
$\norm{\lambda(v-w_0)+w}'=\max(\abs\lambda\norm{v-w_0}',\norm w')$ for $\lambda\in\Kh$ and $w\in V_n$. Let $f_{n+1}$ be a
multiple of $v-w_0$ of norm one.

\emph{Step 2 (the projection).} The closed span $V_\infty$ of the $f_n$ is isometric to $c_0(\N,\Kh)$, complete with
distances in $\varrho^\Z\cup\{0\}$, hence spherically complete (Lemma~\ref{lem:discdist}); the extension argument in the
proof of Theorem~\ref{thm:spherical}, over $\Kh$, gives a contractive projection $\Pi$ of $\widehat E$ onto $V_\infty$.
Choose $e_n\in E$ with $\norm{e_n-f_n}'\le\varrho$. Then $\norm{\sum_n\xi_n(e_n-f_n)}'<\max_n\abs{\xi_n}$ for $\xi\ne0$,
which gives the norm formula for the $e_n$. The operator $\Pi'$ on $V_\infty$ with
$\Pi'(\sum_n\xi_nf_n)=\Pi(\sum_n\xi_ne_n)$ satisfies $\norm{\Pi'-\id_{V_\infty}}\le\varrho<1$, so it is invertible with
$\norm{\Pi'^{-1}}\le1$. Let $\theta_E(x)$ be the coordinate sequence of $\Pi'^{-1}\Pi x$ in the basis $(f_n)$; then
$\theta_E(e_n)$ is the $n$th unit vector. If $E$ is complete, then $E=\widehat E$ is a $\Kh$\nobreakdash-Banach space and
the series defining $\widehat\psi_E$ converge; if $K$ is complete, $K=\Kh$.
\end{proof}

\begin{proof}[Proof of Theorem~\ref{fib:thm:marking}]
\emph{Step 1 (the form $\Phi$).} Let $F_1$ and $\eta$ be as in Appendix~\ref{app:laurent} for $K_1=\Kh$, so that
$E_1=\linf_{\Kh}$, and let $\eta'$ be $\eta$ composed with the map $\widehat\Lambda^k_\pi\linf_{\Kh}\to F_1$ induced by
the identity of $\linf_{\Kh}$, a contraction since a decomposition over $K$ is one over $\Kh$ with the same sum
$\sum_j\prod_i\norm{x_{ji}}$. Suppose $Q$ exists, and for $a,x\in c_{00}(\N,K)^k$ put
\[
 \Phi(a;x)=\theta_F\bigl(Q(a_1,\dots,a_k)(\psi x_1,\dots,\psi x_k)\bigr)\in\widehat\Lambda^k_\pi\linf_{\Kh}.
\]
It is $k$\nobreakdash-linear in $a$ and in $x$, alternating in $x$, of norm at most $\norm{\theta_F}\norm Q\norm\psi^k$,
and $\Phi(a,\dots,a;x)=\omega_{\linf_{\Kh}}(ax_1,\dots,ax_k)$ by~\eqref{fib:eq:marking}.

\emph{Step 2 (reduction to $\kappa$).} For $a,x$ with entries in $\F_p$, $\Psi(a;x)=\eta'(\Phi(a;x))\in
\Lambda^k_\kappa(\kappa^\N)$ is $\F_p$\nobreakdash-multilinear and alternating in $x$, and by Proposition~\ref{prop:supp},
$\sdim\Psi(a;x)\le\lfloor kM_\varrho\norm{\theta_F}\norm Q\norm\psi^k\rfloor$ when all entries lie in $\{0,\pm1\}$. As in
Step~1 of the proof of Theorem~\ref{fib:thm:ell2}, polarizing the diagonal identity over the infinite field $K$ and
applying $\eta'$, which sends $\omega_{\linf_{\Kh}}(y_1,\dots,y_k)$ to $y_1\wedge\dots\wedge y_k$ when every $y_j$ has
entries in $\kappa$ (proof of Lemma~\ref{lem:Psi}), gives~\eqref{eq:Pol1} for $\Psi$ on such inputs. Let $\Psi^\kappa$ be
the $\kappa$\nobreakdash-multilinear map on $(\kappa^{(\N)})^{2k}$ that agrees with $\Psi$ on unit vectors. Antisymmetry in
the last $k$ slots and~\eqref{eq:Pol1} are identities between $\kappa$\nobreakdash-multilinear expressions that hold on
unit vectors, so they hold for $\Psi^\kappa$, and $\Psi^\kappa=\Psi$ on inputs with entries in $\{0,\pm1\}$. Since $k!=0$
in $\kappa$, this contradicts Theorem~\ref{thm:FF}(1).

\emph{Step 3 (consequences).} If $A^k_{E,E';F}$ were analytic at a point, Proposition~\ref{prop:homogeneous-lift} would
give a bounded lift $Q_0$, and $Q(a_1,\dots,a_k)=Q_0(u_{a_1},\dots,u_{a_k})(\nu)$ would be excluded. The last assertion
follows from Lemma~\ref{fib:lem:test}.
\end{proof}

\begin{proof}[Proof of Theorem~\ref{fib:thm:local}]
\emph{Step 1 (finite dimension).} If $\dim_{\Kh}\widehat E<\infty$, then $\widehat E$ has a $\Kh$\nobreakdash-basis
$(e_i)$ with continuous coordinates $(\varepsilon_i)$ (Lemma~\ref{lem:findim}). For $u_1,\dots,u_k\in\cL(E,E')$, with
extensions $\widehat u_i$ to the completions, put
\[
 Q(u_1,\dots,u_k)(x_1\wedge\dots\wedge x_k)=\sum_{i_1<\dots<i_k}\det\bigl(\varepsilon_{i_a}(x_b)\bigr)_{a,b}\,
 \widehat\omega_{E'}(\widehat u_1e_{i_1},\dots,\widehat u_ke_{i_k}),
\]
where the continuous extension $\widehat\omega_{E'}$ of $\omega_{E'}$ is $\Kh$\nobreakdash-multilinear by density.
Expanding $x_b=\sum_i\varepsilon_i(x_b)e_i$ shows, as in Proposition~\ref{prop:split-lift}(2), that $Q$ is a bounded lift
of $u\mapsto\Lambda^k_\pi u$ with values in $\cL(\Lambda^k_\pi E,\widehat\Lambda^k_\pi E')$. For $u_i\in\cL(E',E)$ use the
terms $\det(\varepsilon_{i_a}(u_ax_b))_{a,b}\,\widehat\omega_E(e_{i_1},\dots,e_{i_k})$ instead. If $E$ is complete, then
$e_i\in E$, and moving the scalar into the first argument shows that these lifts take values in the uncompleted spaces.
This proves the positive directions; in~(3) with $\dim E'<\infty$, exchange the roles of $E$ and $E'$ in the second
formula.

\emph{Step 2 (infinite dimension).} Let $\dim_{\Kh}\widehat E=\infty$. Take $e_n$, $\theta_E$ and, for complete $E$,
$\widehat\psi_E$ from Lemma~\ref{fib:lem:marking}, and put $\psi_E(\xi)=\sum_n\xi_ne_n$ on $c_{00}(\N,K)$.
Theorem~\ref{fib:thm:marking} applies to the following data, which satisfy~\eqref{fib:eq:marking}:
\begin{enumerate}
\item[(a)] the pair $(E,c_0(\N,\Kh))$, with $\psi=\psi_E$, $\theta$ the inclusion, and $u_a(x)=a\,\theta_E(x)$;
\item[(b)] the pair $(c_0(\N,K),E)$, with $\psi$ the inclusion, $\theta=\theta_E$, and $u_a(x)=\psi_E(ax)$;
\item[(c)] for complete $E$, the pair $(E,E)$, with $\psi=\psi_E$, $\theta=\theta_E$, and
$u_a(x)=\widehat\psi_E(a\,\theta_E(x))$;
\item[(d)] for complete $K$ and infinite-dimensional $E,E'$, the pair $(E,E')$, with $\psi=\psi_E$,
$\theta=\theta_{E'}$, and $u_a(x)=\psi_{E'}(a\,\theta_E(x))$, since $\theta_E$ then takes values in $c_0(\N,K)$.
\end{enumerate}
Over complete $K$, $\dim_{\Kh}\widehat E=\dim E$, since finite-dimensional spaces are complete. An equivalent
nonarchimedean norm on the values would make the action analytic by Corollary~\ref{cor:discrete-targets}, since
$\abs{K^\times}$ is discrete.
\end{proof}

\subsection{Proofs of the classification results}\label{frc:ssec:classification}
As in Section~\ref{sec:classification}, $c_{00}$, $c_0$ and $\ell^r$ stand for $c_{00}(\N,K)$, $c_0(\N,K)$ and
$\ell^r(\N,K)$.

\begin{proof}[Proof of Theorem~\ref{cls:thm:symmetric}]
\emph{Step 1 ((1)$\Rightarrow$(3) and (2)$\Rightarrow$(3)).} Let $G\in\{E,c_0\}$ with $(G,E)$ good for Banach values in
a degree $k\ge p$; Theorem~\ref{obs:thm:ptb} gives \textup{(PTB)} for $(G,E)$ with a constant $C$. For $x,y\in c_{00}$,
\eqref{obs:eq:ptb} with $a=1_{\operatorname{supp}x}$ and $b=1_{\operatorname{supp}y}$ if $G=E$, and with
$(x,1_{\operatorname{supp}x},y,1_{\operatorname{supp}y})$ in place of $(a,x,b,y)$ if $G=c_0$, gives
$\norm{\cP_<(x\otimes y)}_{E\otimes_\pi E}\le C\norm x_E\norm y_E$, since $\norm{\mu_{1_S}}_{E\to E}\le1$, $\norm{1_S}_{c_0}=1$
and $\norm{\mu_x}_{c_0\to E}=\norm x_E$ (Lemma~\ref{obs:lem:multiplier}(1),(4)).

\emph{Step 2 ((3)$\Rightarrow$(4), (5), (6), (7)).} Condition~(3) is the hypothesis of
Proposition~\ref{spl:prop:ordering} for the unit vectors, a linearly independent family indexed by $\N$ with dense span
$c_{00}$. So $(\mathrm S^n_E)$, $(\mathrm R^n_E)$ and $(\Gamma^n_E)$ hold for every $n\ge1$, which gives (4), (5) and (6) with
$n=p$, and Proposition~\ref{fib:prop:split}(1),(2) gives~(7).

\emph{Step 3 ((4)$\Rightarrow$(1), (5)$\Rightarrow$(1), (6)$\Rightarrow$(3), (7)$\Rightarrow$(1), (2)).}
Proposition~\ref{fib:prop:split}(1), respectively~(2), with $D=E$ gives the first two implications, and
Proposition~\ref{obs:prop:gamma-ptb} the third; for the last take $D=E$, respectively $D=c_0$, and the degree $p$.

\emph{Step 4 ($\ell^2$-domination).} Let $C$ be the constant in~(3). Lemma~\ref{obs:lem:single-vector} with $D=E$ and
$\norm x_\infty\le\norm x_E$ (Lemma~\ref{obs:lem:multiplier}(1)) gives $\norm x_2\le C_2\norm x_E$ for $x\in c_{00}$, with
$C_2^2=C/(1-\varrho)+1/(1-\varrho^2)$, hence for $x\in E$ by taking suprema over the $x1_{[n]}$
(Definition~\ref{obs:def:sequence-spaces}).

\emph{Step 5 (the fundamental function).} We may assume $C>0$. Let $m\ge1$ and
$z_m=\cP_<(1_{[m]}\otimes1_{[m]})=\sum_{t<t'}e_t\otimes e_{t'}$, the sum over $t,t'\in[m]$. By~(3) and fact (P1) of
Section~\ref{frc:ssec:ptb}, $z_m$ has norm at most $C\varphi_E(m)^2$ in $E_{[m]}\otimes_\pi E_{[m]}$; the identity of
$E_{[m]}$ onto $\ell^2([m],K)$ has norm at most $C_2$ by Step~4, so $z_m$ has norm at most $C_2^2C\varphi_E(m)^2$ in
$\ell^2([m],K)\otimes_\pi\ell^2([m],K)$. The matrix $\mathsf Z=\coeff_0z_m$ is the strictly upper triangular matrix of
ones, so every off-diagonal entry of $\mathsf Z-\mathsf Z^{\mathsf T}$ is $\pm1$, and Lemmas~\ref{fib:lem:tournament}
and~\ref{fib:lem:coeff}(1) give
\[
 \frac{m\ln m}{1+\ln(2\abs\kappa)}\le\operatorname{cost}(\mathsf Z)\le\frac{C_2^2C\varphi_E(m)^2}{1-\varrho^2},
\]
so $\varphi_E(m)^2\ge c\,m\ln m$ with $c>0$ independent of $m$.
\end{proof}

\begin{proof}[Proof of Theorem~\ref{cls:thm:rnorm}]
(1)$\Rightarrow$(2) is Theorem~\ref{obs:thm:ptb}, and (4)$\Rightarrow$(1) follows with $k=p$.

(2)$\Rightarrow$(3). Let $\nu$ be an equivalent $r$\nobreakdash-norm on $D$, or an equivalent nonarchimedean norm, in
which case put $r=\infty$, scaled so that $\nu(y)\le\norm y_D\le C''\nu(y)$ for $y\in D$. Since $\nu(e_t)\le1$, iterating
$\nu(y+y')^r\le\nu(y)^r+\nu(y')^r$, respectively $\nu(y+y')\le\max(\nu(y),\nu(y'))$, gives
$\norm\xi_D\le C''\nu(\xi)\le C''\norm\xi_r$ for $\xi\in c_{00}$. As $\norm y_\infty\le\norm y_D$
(Lemma~\ref{obs:lem:multiplier}(1)), Proposition~\ref{obs:prop:l2-domination} with $C'=C''$ gives~(3).

(3)$\Rightarrow$(4). The unit vectors of $E$ are a linearly independent family indexed by $\N$, of norm one with dense
span $c_{00}$, and $\ell^2$\nobreakdash-dominated by~(3) (Definition~\ref{obs:def:sequence-spaces}). As $K$ is
nonarchimedean, Theorem~\ref{spl:thm:relative-ordering}(2) for $r<\infty$, and
Corollary~\ref{spl:cor:relative-ordering-na} for a nonarchimedean norm, give~(4).
\end{proof}

\begin{proof}[Proof of Corollary~\ref{cls:cor:lr-table}]
\emph{Positive cases.} If $r_i<2$ for some $i$, then $\ell^{r_i}$ has a lower $r_i$\nobreakdash-estimate with its unit
vectors and $C=1$ (the examples after Theorem~\ref{fib:thm:lower}), and Theorem~\ref{fib:thm:lower}(3) makes the pair good
for Banach values in every degree. If $r_1=2<r_2$, the unit vectors of $\ell^2$ are $\ell^2$\nobreakdash-dominated, of norm one
with dense span; $\ell^{r_2}$ is $r_2$\nobreakdash-normed by its own norm if $r_2<\infty$, as noted
after Corollary~\ref{spl:cor:relative-ordering-na}, and $c_0$ has a nonarchimedean norm, so
Theorem~\ref{spl:thm:relative-ordering}(2), respectively Corollary~\ref{spl:cor:relative-ordering-na}, makes the pair good
for Banach values in every degree.

\emph{Negative cases.} For the remaining pairs, $r_2=2\le r_1$ or $r_1,r_2>2$, Corollary~\ref{fib:cor:pairs}
excludes every degree $k\ge p$ and Corollary~\ref{fib:cor:degrees}(1) gives every degree $k<p$. So
$\fD_B(\ell^{r_1},\ell^{r_2})$ is $\{1,2,\dots\}$ in the positive cases and $\{1,\dots,p-1\}$ otherwise, which gives the
equivalence and the table.
\end{proof}

\subsection{Spreading spaces}\label{frc:ssec:spreading}
For a finite set $S\subseteq\N$, a map $\psi\colon S\to\N$ that is increasing and $x\in c_{00}(\N,K)$ supported in $S$,
put $\psi_*x=\sum_{t\in S}x(t)e_{\psi(t)}$. If $E$ is spreading, then $\abs{\psi_*x}=\sum_{t\in S}\abs{x(t)}e_{\psi(t)}$
and Definition~\ref{cls:def:spreading} give $\norm{\psi_*x}_E=\norm x_E$; in particular $\norm{1_S}_E=\varphi_E(\abs S)$.

\begin{proof}[Proof of Theorem~\ref{cls:thm:spreading}]
\emph{Step 1 (item (1)).} In the proof of Lemma~\ref{obs:lem:multiplier}, invariance under permutations is used only
in~(3), to see that a vector $y$ supported in $\operatorname{supp}a$ and its relabeling $y^\gamma$ have the same norm.
For $\gamma$ increasing on $\operatorname{supp}a$ this is the identity $\norm{\psi_*y}=\norm y$ above, with $\psi$ the
restriction of $\gamma$ to $\operatorname{supp}a$. The rest of the proof uses only monotonicity, $\norm{e_t}=1$ and
$\norm x=\sup_n\norm{x1_{[n]}}$.

\emph{Step 2 (the facts (P1)--(P3)).} They hold for spreading $D$, with (P2) for relabelings of the rows along a map
increasing on the row support and of the columns along a map increasing on the column support: such relabelings map
$D_J$ and $D_{J'}$ isometrically onto $D_{\psi(J)}$ and $D_{\psi'(J')}$, so their tensor product is an isometry for the
projective norms, and (P1) applies.

\emph{Step 3 (Theorem~\ref{obs:thm:ptb}).} Lemmas~\ref{frc:lem:homogeneity}, \ref{frc:lem:clusters}
and~\ref{frc:lem:diagonal} and Corollary~\ref{frc:cor:transposition} concern sequence spaces $G$ and $D$ with
$\norm{e_t}_G\le1$ and $\norm\cdot_\infty\le\norm\cdot_D$, and these properties hold for spreading spaces by Step~1. In
the proof of Theorem~\ref{obs:thm:ptb}, the norms of $E$ and $D$ enter through Lemma~\ref{obs:lem:multiplier}, the facts
(P1)--(P3) and relabelings, and each relabeling is increasing on the support it is applied to: in Step~2, $f_i\mapsto
Y_i(c)$ on $J_1$ and on $J_2$, $c=1,\dots,4$, since $Y_1<\dots<Y_m$; in Step~3, $Y_i(1)\mapsto f_i$ on $R$ and on $R'$; in
Step~5(a), $f_i\mapsto2i$ on $J_1$ and $f_i\mapsto2i-1$ on $J_2$. Steps~5(b) and~(c) use only $\norm x=\norm{\abs x}$ and
homogeneity. So Steps~1 and~2 give the theorem for spreading $E$ and $D$.

\emph{Step 4 (the other results of item (2)).} The proofs of Propositions~\ref{obs:prop:l2-domination}
and~\ref{obs:prop:gamma-ptb} and of Theorems~\ref{cls:thm:symmetric} and~\ref{cls:thm:rnorm} use the symmetry of the
spaces only through Lemma~\ref{obs:lem:multiplier}, Theorem~\ref{obs:thm:ptb} and fact~(P1), together with
$\norm{e_t}_D=1$ in Theorem~\ref{cls:thm:rnorm}; Lemma~\ref{obs:lem:single-vector} concerns every sequence space $D$ with
$\norm\cdot_\infty\le\norm\cdot_D$, and $c_0$ is spreading. Corollaries~\ref{cls:cor:symmetric-examples}
and~\ref{cls:cor:dichotomy}(1),(2) follow from these results as before. Remark~\ref{cls:rem:mechanisms} rests, for these
pairs, on Proposition~\ref{spl:prop:ordering} (Step~2 of the proof of Theorem~\ref{cls:thm:symmetric}) and on the
relative ordering splitting described after Theorem~\ref{cls:thm:rnorm} (Theorem~\ref{spl:thm:relative-ordering}(1) and
Corollary~\ref{spl:cor:relative-ordering-na}), which concern arbitrary normed spaces.

\emph{Step 5 (item (3)).} Here $K$ is any nontrivially normed field of characteristic~$2$; condition~\textup{(L)} is not
used. Steps~7 and~8 of the proof of Theorem~\ref{obs:thm:char-two} use of $E$ only that it is a sequence
space with $\norm{e_t}_E=1$, $\norm x_\infty\le\norm x_E$, $\norm{\mu_a}_{E\to E}\le\norm a_\infty$ and
$\sup_nn/\bar\varphi_E(n)^2=\infty$. For spreading $E$ the first three follow from monotonicity, as there, and
$\bar\varphi_E=\varphi_E$, since $\norm{1_S}_E=\varphi_E(\abs S)$ and $\varphi_E$ is nondecreasing.
\end{proof}

\begin{proof}[Proof of the assertions of Example~\ref{cls:ex:garling}]
Write $\norm\cdot$ for $\norm\cdot_{\alpha,q}$ on $c_{00}(\N,\mathbb R)$, $\varphi$ for $\varphi_{E_{\alpha,q}}$, and put
$\mathfrak h_n=\sum_{j\le n}1/j$.

\emph{Step 1 (spreading).} Minkowski's inequality along each increasing sequence makes $\norm\cdot$ a norm. It depends
only on $\abs f$, is monotone in $\abs f$ and dominates $\norm f_\infty$, and $\norm{e_t}=1$. Let
$f=\sum_{i\le d}f_ie_{t_i}$ with $t_1<\dots<t_d$. An increasing sequence $(u_j)$ contributes only through its terms in
$\{t_1,\dots,t_d\}$, say $u_{j_s}=t_{a_s}$ with $j_1<j_2<\cdots$ and $a_1<a_2<\cdots$; as $j_s\ge s$, its sum is at most
$\sum_ss^{-\alpha}\abs{f_{a_s}}^q$, the sum of the sequence $(t_{a_s})_s$. So $\norm f^q$ is the largest value of
$\sum_ss^{-\alpha}\abs{f_{a_s}}^q$ over the increasing sequences $(a_s)$ in $\{1,\dots,d\}$, which does not depend on
$t_1,\dots,t_d$. The other conditions of Definition~\ref{cls:def:spreading} hold by the definition of $E_{\alpha,q}$.

\emph{Step 2 (decreasing profiles; item (1)).} If $\abs{g(1)}\ge\abs{g(2)}\ge\cdots$, then $\abs{g(u_j)}\le\abs{g(j)}$
for every increasing sequence $(u_j)$, so $\norm g^q=\sum_jj^{-\alpha}\abs{g(j)}^q$. In particular
$\varphi(m)^q=\sum_{j\le m}j^{-\alpha}$, which lies between $m^{1-\alpha}$ and
$1+\int_1^m\tau^{-\alpha}\,d\tau\le m^{1-\alpha}/(1-\alpha)$.

\emph{Step 3 (blocks).} For $n,M\ge1$ let $I_i=\{(i-1)n+1,\dots,in\}$ and
$x=\sum_{i\le M}\sum_{j\le n}j^{-\beta}e_{(i-1)n+j}$, a real vector. By Steps~1 and~2 each $x1_{I_i}$ has norm $\mathfrak
h_n^{1/q}$, since $j^{-\alpha}j^{-\beta q}=j^{-1}$. We show $\norm x^q\le1+\ln(Mn)+(M-1)/\alpha$. It suffices to consider
increasing sequences in the support of $x$, since other terms only raise the indices of the later ones. Fix one, let
$m_i$ be the number of its terms in $I_i$, and $\lambda_i=m_1+\dots+m_{i-1}$. The terms in $I_i$ have indices
$\lambda_i+s$, $s\le m_i$, at positions $(i-1)n+j_s$ with $j_s\ge s$, so they contribute at most
$G_i=\sum_{s\le m_i}(\lambda_i+s)^{-\alpha}s^{\alpha-1}$. If $\lambda_i=0$, then $G_i=\mathfrak h_{m_i}\le1+\ln m_i$. If
$\lambda_i\ge1$, the function $\tau\mapsto(\lambda_i+\tau)^{-\alpha}\tau^{\alpha-1}$ decreases on $(0,\infty)$, so the
substitution $\tau=\lambda_iv$ gives
\[
 G_i\le\int_0^{m_i/\lambda_i}(1+v)^{-\alpha}v^{\alpha-1}\,dv\le\frac1\alpha+\ln\frac{\lambda_i+m_i}{\lambda_i},
\]
using $(1+v)^{-\alpha}\le1$ on $(0,1]$ and $(1+v)^{-\alpha}v^{\alpha-1}\le v^{-1}$ on $[1,\infty)$. Let $i_0$ be the first
$i$ with $m_i>0$. At most $M-1$ later blocks are met, and their logarithms add up to the logarithm of the number of
terms divided by $m_{i_0}$, at most $\ln(Mn/m_{i_0})$. So the sum of the sequence is at most
$1+\ln m_{i_0}+(M-1)/\alpha+\ln(Mn/m_{i_0})$.

\emph{Step 4 (item (2): symmetric norms).} Let $x^*=x\circ\sigma$ be the decreasing rearrangement of $x$, for a
permutation $\sigma$ of $[Mn]$; it takes the value $j^{-\beta}$ at the $M$ positions $(j-1)M+1,\dots,jM$. By Step~2,
\[
 \norm{x^*}^q=\sum_{j\le n}j^{\alpha-1}\sum_{l\le M}((j-1)M+l)^{-\alpha}\ge M^{1-\alpha}\mathfrak h_n,
\]
so by Step~3 the ratio $\norm{x\circ\sigma}^q/\norm x^q$ is at least $M^{1-\alpha}\mathfrak h_n/(1+\ln(Mn)+(M-1)/\alpha)$,
which is unbounded as $n$ and then $M$ grow. For real $f\ge0$ in $c_{00}(\N,\mathbb R)$ let $f'\in c_{00}(\N,K)$ have
$f'(t)=0$ if $f(t)=0$ and $f'(t)=\vartheta^{c(t)}$ otherwise, with $c(t)$ the largest integer such that
$\abs\vartheta^{c(t)}\ge f(t)$. Then $f\le\abs{f'}<\abs\vartheta^{-1}f$, so $\norm f\le\norm{f'}\le\abs\vartheta^{-1}\norm f$,
and $(f\circ\sigma)'=f'\circ\sigma$. Hence $\norm{x'\circ\sigma}/\norm{x'}\ge\abs\vartheta\norm{x\circ\sigma}/\norm x$ is
unbounded, which a norm on $c_{00}$ invariant under permutations and equivalent to that of $E_{\alpha,q}$ would exclude.

\emph{Step 5 (item (2): lower estimates).} Let $r\in(0,2)$ and $C\ge1$, and let $x'$ be as in Step~4. The $I_i$ are
intervals, $\norm{x'1_{I_i}}\ge\norm{x1_{I_i}}=\mathfrak h_n^{1/q}$ and $\norm{x'}\le\abs\vartheta^{-1}\norm x$, so
\eqref{fib:eq:lower} would give $M^{1/r}\mathfrak h_n^{1/q}\le C\abs\vartheta^{-1}(1+\ln(Mn)+(M-1)/\alpha)^{1/q}$. For
$M=\lceil\ln n\rceil$ and large $n$ the right side is at most $C\abs\vartheta^{-1}((2+1/\alpha)\ln n)^{1/q}$ and the left
side at least $(\ln n)^{1/r+1/q}$, so $(\ln n)^{1/r}\le C\abs\vartheta^{-1}(2+1/\alpha)^{1/q}$, which fails for large $n$.

\emph{Step 6 (items (3) and (4)).} If $2\beta\le1$, then $\varphi(m)^2\le(1-\alpha)^{-2/q}m$ by~(1), and
Corollary~\ref{cls:cor:symmetric-examples}, which holds for spreading spaces by Theorem~\ref{cls:thm:spreading}(2), gives
the first assertion of~(3); if $2\beta<1$, then $\sup_mm/\varphi(m)^2=\infty$, and Theorem~\ref{cls:thm:spreading}(3)
gives the second, Steps~1 and~2 holding over every nontrivially normed field. For~(4), let $y\in c_{00}$ and let $y^*_m$
be the $m$th largest value of $\abs y$. The positions of the $m$ largest values, in increasing order, give
$\norm y^q\ge\sum_{s\le m}s^{-\alpha}(y^*_m)^q\ge m^{1-\alpha}(y^*_m)^q$, so $y^*_m\le m^{-\beta}\norm y$ and
$\norm y_2^2\le(\sum_mm^{-2\beta})\norm y^2$. For $y\in E_{\alpha,q}$ take the supremum over the $y1_{[n]}$.
\end{proof}
%% ---- end app-K-forcing.tex
%% ---- begin app-L-degrees.tex
\section{Degrees of splittings}\label{app:degrees}

This appendix proves the results of Section~\ref{sec:combining} on degrees, composition, polarizability and permanence.
Since $K$ is nontrivially normed, it is infinite, so the coefficients of a polynomial map from $K^d$ to a vector space are
linear combinations of its values. Degree $0$ follows the paragraph on split operators in Section~\ref{cmb:ssec:degrees},
with $\Lambda^0_\pi u=\id_K$.

\subsection{Gaps in the degree sets}\label{dgr:ssec:gaps}

\begin{lemma}[retracts and lines]\label{dgr:lem:lines}
Let $E$ and $D$ be normed spaces and $k\ge1$.
\begin{enumerate}
\item Let $E_0$ and $D_0$ be normed spaces, and let $\iota_1\colon E_0\to E$, $\iota'_1\colon E\to E_0$,
$\iota_2\colon D_0\to D$ and $\iota'_2\colon D\to D_0$ be bounded linear maps with $\iota'_1\iota_1=\id_{E_0}$ and
$\iota'_2\iota_2=\id_{D_0}$. If $(E,D)$ is good, or good for Banach values, in degree $k$, then so is $(E_0,D_0)$. In
particular, isomorphic pairs have the same degree sets.
\item If $f\colon E\to K^d$ has as coordinates linearly independent bounded linear functionals $f_1,\dots,f_d$, then $E$
is isomorphic to $\ker f\oplus K^d$ with the maximum norm.
\item Give $E\oplus K$ and $D\oplus K$ the maximum norms. If $(E,D)$ is good, or good for Banach values, in every degree
$n\le k$, then so is $(E\oplus K,D\oplus K)$ in degree $k$.
\end{enumerate}
\end{lemma}

\begin{proof}
(1) For Banach values this is Proposition~\ref{cmb:prop:permanence}(1). For goodness, if $Q$ is a lift for $(E,D)$ as in
Lemma~\ref{fib:lem:test}, then
$Q'(u_1,\dots,u_k)=\Lambda^k_\pi\iota'_2\circ Q(\iota_2u_1\iota'_1,\dots,\iota_2u_k\iota'_1)\circ\Lambda^k_\pi\iota_1$ is a
lift for $(E_0,D_0)$, since $Q'(u,\dots,u)=\Lambda^k_\pi(\iota'_2\iota_2u\iota'_1\iota_1)=\Lambda^k_\pi u$ by
Proposition~\ref{prop:exterior}. Isomorphic pairs are retracts of each other in this sense.

(2) The map $f$ is onto, since a nonzero functional $c$ on $K^d$ vanishing on its image would give $\sum_ac_af_a=0$. For
$x_1,\dots,x_d\in E$ with $f_a(x_b)=\delta_{ab}$, the maps $x\mapsto(x-\sum_af_a(x)x_a,f(x))$ and
$(y,\lambda)\mapsto y+\sum_a\lambda_ax_a$ are mutually inverse bounded linear maps between $E$ and $\ker f\oplus K^d$.

(3) Let $J_E,\Pi_E,J_K,\Pi_K$ be the inclusions and projections of $E\oplus K$, and $J_D,\Pi_D,J'_K,\Pi'_K$ those of
$D\oplus K$. Since $J_E\Pi_E+J_K\Pi_K$ and $J_D\Pi_D+J'_K\Pi'_K$ are identities, every
$u\in\cL(E\oplus K,D\oplus K)$ is $\psi_1(u)+\psi_2(u)+\psi_3(u)$ with $\psi_1(u)=J_D\Pi_DuJ_E\Pi_E$,
$\psi_2(u)=J'_K\Pi'_KuJ_E\Pi_E$ and $\psi_3(u)=uJ_K\Pi_K$. For $n\le k$ and a lift $Q_n$ for $(E,D)$ in degree $n$, the map
$\Lambda^n_\pi J_D\circ Q_n(\Pi_Du_1J_E,\dots,\Pi_Du_nJ_E)\circ\Lambda^n_\pi\Pi_E$, with $\widehat\Lambda^n_\pi J_D$ for
Banach values, lifts $\psi_1$ in degree $n$. The operators $\psi_2(u)$ and $\psi_3(u)$ have rank at most one, so their exterior
powers of degree $n\ge2$ vanish, and $\psi_2$ and $\psi_3$ are liftable in every degree, with the lift $0$ in degrees $n\ge2$. Theorem~\ref{cmb:thm:decomposition}(2)
concludes.
\end{proof}

\begin{theorem}[no gap after the first failure]\label{dgr:thm:gap}
Let $\operatorname{char}K=p>0$, let $E$ and $D$ each have $d$ linearly independent bounded linear functionals, and let
$\fD$ be $\fD(E,D)$ or $\fD_B(E,D)$. If $\fD\ne\{1,2,\dots\}$ and $n$ is its least missing element, then none of
$n,n+1,\dots,n+d$ lies in $\fD$.
\end{theorem}

\begin{proof}
\emph{Step 1 (splitting off lines).} Let $f\colon E\to K^d$ and $g\colon D\to K^d$ have the functionals as coordinates.
For $0\le l\le d$ let $E^{(l)}=\ker f\oplus K^{d-l}$ and $D^{(l)}=\ker g\oplus K^{d-l}$ with the maximum norms, so that
$E^{(l)}=E^{(l+1)}\oplus K$ and $D^{(l)}=D^{(l+1)}\oplus K$. Let $\fD^{(l)}$ be the set of the same kind for
$(E^{(l)},D^{(l)})$, and $n^{(l)}$ its least missing element, or $\infty$ if there is none. By
Lemma~\ref{dgr:lem:lines}(1),(2), $\fD^{(0)}=\fD$ and $n^{(0)}=n$.

\emph{Step 2 (the first failure does not move).} The pair $(E^{(l+1)},D^{(l+1)})$ is a retract of $(E^{(l)},D^{(l)})$,
so $\fD^{(l)}\subseteq\fD^{(l+1)}$ and $n^{(l)}\le n^{(l+1)}$ by Lemma~\ref{dgr:lem:lines}(1). If $i<n^{(l+1)}$, then
$1,\dots,i\in\fD^{(l+1)}$, so $i\in\fD^{(l)}$ by Lemma~\ref{dgr:lem:lines}(3); hence $n^{(l)}\ge n^{(l+1)}$. By induction
$n^{(l)}=n$ for every $l\le d$.

\emph{Step 3 (descent).} Suppose that $k\in\fD$ with $n<k\le n+d$, and put $i=k-n$. By
Corollary~\ref{fib:cor:degrees}(1), $n\ge p\ge2$, so $k,k-1,\dots,n+1\ge2$. Theorem~\ref{fib:thm:hyperplane}, applied with
every normed, respectively Banach, space $F$, turns a degree $k'\ge2$ in $\fD^{(l)}$ into the degree $k'-1$ in
$\fD^{(l+1)}$. Applied $i$ times, it gives $n\in\fD^{(i)}$, contradicting $n^{(i)}=n$.
\end{proof}

\begin{proof}[Proof of Theorem~\ref{cmb:thm:initial-segments}]
Since $E^*$ and $D^*$ are infinite-dimensional, Theorem~\ref{dgr:thm:gap} applies for every $d$. So if $\fD$, which is
$\fD(E,D)$ or $\fD_B(E,D)$, has a least missing element $n$, then no degree $k\ge n$ lies in $\fD$, and
$\fD=\{1,\dots,n-1\}$. By Corollary~\ref{fib:cor:degrees}(2), $n=p^h$ for a unique $h\ge1$.
\end{proof}

\subsection{Column constructions}\label{dgr:ssec:columns}

Throughout this subsection $q,m\ge1$ and $n=qm$. For vectors $x_1,\dots,x_n$ and $C\subseteq\{1,\dots,n\}$, $x_C$ is the
wedge of the $x_c$, $c\in C$, in increasing order. For disjoint sets $C_1,\dots,C_l$, $\epsilon(C_1,\dots,C_l)$ is the sign
of the listing of $C_1$, then $C_2$, and so on, each increasingly, relative to the increasing listing of their union; so
$x_{C_1}\wedge\dots\wedge x_{C_l}=\epsilon(C_1,\dots,C_l)\,x_{C_1\cup\dots\cup C_l}$. A listing
$\lambda=(\lambda_1,\dots,\lambda_n)$ of $\{1,\dots,n\}$ has the sign $\sgn(\lambda)$ as a permutation, and $N_{q,m}$ is as
in Theorem~\ref{cmb:thm:composition}.

For a normed space $G$ let $\wedge^{[q,m]}_G\colon(\Lambda^m_\pi G)^{\otimes_\pi q}\to\Lambda^n_\pi G$ be the iterated
wedge product $\xi_1\otimes\dots\otimes\xi_q\mapsto\xi_1\wedge\dots\wedge\xi_q$, contractive by Lemma~\ref{fib:lem:hopf},
and $\sh^{[q,m]}_G=(\id\otimes\dots\otimes\id\otimes\sh^{m,m}_G)\circ\dots\circ(\id\otimes\sh^{m,(q-2)m}_G)\circ
\sh^{m,(q-1)m}_G$ from $\Lambda^n_\pi G$ to $(\Lambda^m_\pi G)^{\otimes_\pi q}$, so that
\[
 \sh^{[q,m]}_G(x_1\wedge\dots\wedge x_n)=\sum_{(J_1,\dots,J_q)}\epsilon(J_1,\dots,J_q)\,x_{J_1}\otimes\dots\otimes x_{J_q},
\]
the sum over the ordered partitions of $\{1,\dots,n\}$ into $q$ blocks of size $m$. Both maps are natural in bounded
operators, and iterating Lemma~\ref{fib:lem:hopf} gives, on $\otimes^n_\pi G=(\otimes^m_\pi G)^{\otimes_\pi q}$,
\begin{equation}\label{dgr:eq:iterated}
 \wedge^n_G=\wedge^{[q,m]}_G\circ(\wedge^m_G)^{\otimes q},\qquad
 \alt^n_G=(\alt^m_G)^{\otimes q}\circ\sh^{[q,m]}_G .
\end{equation}
Their extensions to completions keep their names, with a hat on $\wedge^{[q,m]}_G$. In
$\otimes^n_\pi G=(\otimes^q_\pi G)^{\otimes_\pi m}$, the $t$th block of $q$ factors is column $t$, and the $i$th factors of
the columns form row $i$; the isometry $\tau_\zeta$ of Section~\ref{sec:canonical}, where $\zeta\in S_n$ sends the position
$(t-1)q+i$ to $(i-1)m+t$, passes to the row-major order, in which row $i$ is the $i$th block of $m$ factors.

\begin{lemma}\label{dgr:lem:unit}
If $q$ and $m$ are powers of a prime $p$, then $p$ does not divide $N_{q,m}$. Hence, if $\operatorname{char}K=p$, the
image of $N_{q,m}$ in $K$ is a nonzero element of the prime field, and its absolute value is $1$.
\end{lemma}

\begin{proof}
By Legendre's formula, the exponent of $p$ in $(p^c)!$ is $p^{c-1}+\dots+1=(p^c-1)/(p-1)$. For $q=p^a$ and $m=p^b$, the
exponent of $p$ in $N_{q,m}=(qm)!/((q!)^m\,m!)$ is therefore
$\bigl((p^{a+b}-1)-p^b(p^a-1)-(p^b-1)\bigr)/(p-1)=0$. Every nonzero element of $\F_p$ is a root of unity, so its
absolute value is $1$.
\end{proof}

\begin{proposition}[source column map]\label{dgr:prop:s-column}
Let $w\colon E\to G$ have a section $s$ in degree $q$. There is a bounded linear map $\operatorname{col}_s$ from
$\Lambda^n_\pi E$ to the completion of $(\Lambda^m_\pi G)^{\otimes_\pi q}$, of norm at most $N_{q,m}\norm s^m$ with
$N_{q,m}$ read as a real number, such that
\[
 \widehat\wedge^{[q,m]}_G\circ\operatorname{col}_s=\eta\,N_{q,m}\,\Lambda^n_\pi w,
\]
where $\eta=\pm1$ depends only on $q$ and $m$. If $s$ takes values in $\otimes^q_\pi G$, then $\operatorname{col}_s$ takes
values in $(\Lambda^m_\pi G)^{\otimes_\pi q}$.
\end{proposition}

\begin{proof}
\emph{Step 1 (regrouping).} Let $\Theta=(\wedge^m_G)^{\otimes q}\circ\tau_\zeta$, a contraction from $\otimes^n_\pi G$
to $(\Lambda^m_\pi G)^{\otimes_\pi q}$. It sends the tensor with columns $g_{1t}\otimes\dots\otimes g_{qt}$ to
$(g_{11}\wedge\dots\wedge g_{1m})\otimes\dots\otimes(g_{q1}\wedge\dots\wedge g_{qm})$. If $\gamma\in S_m$ and
$\tau'_\gamma$ moves column $t$ to column $\gamma(t)$, then $\Theta\circ\tau'_\gamma=\sgn(\gamma)^q\,\Theta$, since $\gamma$
permutes the factors of each of the $q$ row wedges. By~\eqref{dgr:eq:iterated} and Section~\ref{sec:canonical},
$\wedge^{[q,m]}_G\circ\Theta=\wedge^n_G\circ\tau_\zeta=\eta\,\wedge^n_G$ with $\eta=\sgn(\zeta)$.

\emph{Step 2 (the multilinear map).} For $x_1,\dots,x_n\in E$ put
\[
 \Xi_s(x_1,\dots,x_n)=\sum_{\{C_1,\dots,C_m\}}\epsilon(C_1,\dots,C_m)\,
 \Theta\bigl(s(x_{C_1})\otimes\dots\otimes s(x_{C_m})\bigr),
\]
the sum over the partitions of $\{1,\dots,n\}$ into $m$ blocks of size $q$, each with some order of its blocks.
Reordering the blocks by $\gamma\in S_m$ multiplies $\epsilon(C_1,\dots,C_m)$ by $\sgn(\gamma)^{q^2}=\sgn(\gamma)^q$ and,
by Step~1, the $\Theta$-term by $\sgn(\gamma)^q$; so each term depends only on the partition. The map $\Xi_s$ is
$n$-linear of norm at most $N_{q,m}\norm s^m$.

\emph{Step 3 ($\Xi_s$ is alternating).} For a listing $\lambda$ of $\{1,\dots,n\}$, cut into $m$ consecutive groups of $q$
entries, let $\Phi(\lambda)=\sgn(\lambda)\,\Theta\bigl(s(x_{\lambda_1}\wedge\dots\wedge x_{\lambda_q})\otimes
s(x_{\lambda_{q+1}}\wedge\dots\wedge x_{\lambda_{2q}})\otimes\dots\bigr)$. Permuting entries inside a group changes
$\sgn(\lambda)$ and that group's wedge by the same sign, so $\Phi(\lambda)$ is unchanged; and the term of a partition is
$\Phi(\lambda)$ for $\lambda$ listing $C_1$, then $C_2$, and so on. Let $x_a=x_b$ with $a\ne b$. The terms with $a$ and
$b$ in one block vanish. Exchanging $a$ and $b$ is an involution without fixed points on the partitions that separate
them, and it replaces $\lambda$ by the listing $\lambda'$ with the entries $a$ and $b$ exchanged; since
$\sgn(\lambda')=-\sgn(\lambda)$ and $x_a=x_b$, $\Phi(\lambda')=-\Phi(\lambda)$, so these terms cancel in pairs.
Proposition~\ref{prop:exterior} gives $\operatorname{col}_s$ with $\operatorname{col}_s(x_1\wedge\dots\wedge x_n)=\Xi_s(x)$.

\emph{Step 4 (the identity).} By Step~1, \eqref{dgr:eq:iterated} for $m$ blocks of size $q$, and
$\widehat\wedge^q_Gs=\Lambda^q_\pi w$,
\[
\begin{aligned}
 \widehat\wedge^{[q,m]}_G\Theta\bigl(s(x_{C_1})\otimes\dots\otimes s(x_{C_m})\bigr)
 &=\eta\,\Lambda^q_\pi w(x_{C_1})\wedge\dots\wedge\Lambda^q_\pi w(x_{C_m})\\
 &=\eta\,\epsilon(C_1,\dots,C_m)\,\Lambda^n_\pi w(x_1\wedge\dots\wedge x_n),
\end{aligned}
\]
so each of the $N_{q,m}$ terms of $\Xi_s$ contributes $\eta\Lambda^n_\pi w(x_1\wedge\dots\wedge x_n)$. If $s$ takes values
in $\otimes^q_\pi G$, every term lies in $(\Lambda^m_\pi G)^{\otimes_\pi q}$.
\end{proof}

\begin{proposition}[target column map]\label{dgr:prop:r-column}
Let $w\colon G\to D$ have a retraction $\rho$ in degree $q$. There is a bounded linear map $\operatorname{col}^\rho$ from
the completion of $(\widehat\Lambda^m_\pi G)^{\otimes_\pi q}$ to $\widehat\Lambda^n_\pi D$, of norm at most
$(m!)^{q-1}\norm\rho^m$, such that
\[
 \operatorname{col}^\rho\circ\sh^{[q,m]}_G=\eta'\,N_{q,m}\,\Lambda^n_\pi w,
\]
where $\eta'=\pm1$ depends only on $q$ and $m$. If $\rho$ takes values in $\Lambda^q_\pi D$, then
$\operatorname{col}^\rho$ maps $(\Lambda^m_\pi G)^{\otimes_\pi q}$ into $\Lambda^n_\pi D$.
\end{proposition}

\begin{proof}
\emph{Step 1 (the map).} For vectors $y_{i,t}\in G$ in rows $i$ and columns $t$, put
\[
 \Psi_\rho(y)=\sum_{g_2,\dots,g_q\in S_m}\sgn(g_2)\cdots\sgn(g_q)\bigwedge_{t=1}^m
 \rho\bigl(y_{1,t}\otimes y_{2,g_2(t)}\otimes\dots\otimes y_{q,g_q(t)}\bigr)\in\widehat\Lambda^n_\pi D,
\]
the wedge taken in the order $t=1,\dots,m$. It is multilinear of norm at most $(m!)^{q-1}\norm\rho^m$.

\emph{Step 2 (alternation in each row).} Let $y_{i,t}=y_{i,t'}$ with $t\ne t'$. If $i\ge2$, replacing $g_i$ by
$(t\ t')\circ g_i$ keeps every factor and changes the sign. If $i=1$, replacing every $g_l$ by $g_l\circ(t\ t')$ exchanges
the wedge factors in the positions $t$ and $t'$, which gives $(-1)^{q^2}=(-1)^q$, and multiplies the coefficient by
$(-1)^{q-1}$; so the sign changes. Both replacements are involutions without fixed points, so the terms cancel in pairs.
For $q=1$, $\Psi_\rho$ is alternating since $\rho=w$. By Proposition~\ref{prop:exterior} in each row, $\Psi_\rho$ induces
a bounded linear map $\operatorname{col}^\rho$ on $(\Lambda^m_\pi G)^{\otimes_\pi q}$, which extends to the completion.

\emph{Step 3 (arrays).} Let $x_1,\dots,x_n\in G$. For an ordered partition $(J_1,\dots,J_q)$ with
$J_i=\{b_{i,1}<\dots<b_{i,m}\}$, the corresponding term of $\operatorname{col}^\rho\sh^{[q,m]}_G(x_1\wedge\dots\wedge x_n)$
is $\epsilon(J_1,\dots,J_q)\Psi_\rho(y)$ with $y_{i,t}=x_{b_{i,t}}$. An \emph{array} is a bijection
$\alpha\colon\{1,\dots,q\}\times\{1,\dots,m\}\to\{1,\dots,n\}$; $\sgn_{\mathrm{row}}(\alpha)$ and
$\sgn_{\mathrm{col}}(\alpha)$ are the signs of its readings row by row and column by column, and
$\sgn_{\mathrm{row}}=\eta'\sgn_{\mathrm{col}}$ for a sign $\eta'$ depending only on $q$ and $m$. Through
$\alpha(1,t)=b_{1,t}$ and $\alpha(i,t)=b_{i,g_i(t)}$ for $i\ge2$, the pairs $((J_i),(g_i))$ correspond bijectively to the
arrays with increasing first row, and $\epsilon(J_1,\dots,J_q)\prod_i\sgn(g_i)=\sgn_{\mathrm{row}}(\alpha)$. Hence
\[
 \operatorname{col}^\rho\sh^{[q,m]}_G(x_1\wedge\dots\wedge x_n)=\eta'\sum_\alpha\Phi(\alpha),\qquad
 \Phi(\alpha)=\sgn_{\mathrm{col}}(\alpha)\bigwedge_{t=1}^m\rho\bigl(x_{\alpha(1,t)}\otimes\dots\otimes x_{\alpha(q,t)}\bigr),
\]
the sum over the arrays with increasing first row.

\emph{Step 4 (columns).} Permuting the columns of an array by $\gamma\in S_m$ leaves $\Phi$ unchanged, since
$\sgn_{\mathrm{col}}$ and the wedge of $m$ elements of degree $q$ both change by $\sgn(\gamma)^{q^2}$. Group the arrays
by the partition $\{C_1,\dots,C_m\}$ formed by their column sets. For a fixed partition, the arrays with increasing first
row correspond bijectively to the choices of an ordering $o_t$ of each $C_t$, since the first entries then fix the order
of the columns. By the column invariance of $\Phi$, the partition contributes
\[
 \epsilon(C_1,\dots,C_m)\sum_{o_1,\dots,o_m}\prod_t\sgn(o_t)\bigwedge_t\rho(x_{o_t})
 =\epsilon(C_1,\dots,C_m)\bigwedge_t\rho\bigl(\alt^q_G(x_{C_t})\bigr),
\]
where $x_{o_t}$ is the tensor of the $x_c$, $c\in C_t$, in the order $o_t$, and $\sgn(o_t)$ is relative to the increasing
order. Since $\rho\alt^q_G=\Lambda^q_\pi w$, this is
$\epsilon(C_1,\dots,C_m)\Lambda^q_\pi w(x_{C_1})\wedge\dots\wedge\Lambda^q_\pi w(x_{C_m})
=\Lambda^n_\pi w(x_1\wedge\dots\wedge x_n)$. Summing over the $N_{q,m}$ partitions gives the identity on pure wedges,
hence everywhere.
\end{proof}

\begin{proof}[Proof of Theorem~\ref{cmb:thm:composition}]
Put $N=N_{q,m}$, invertible in $K$. (1) Let $s=(\eta N)^{-1}(s_2)^{\otimes q}\circ\operatorname{col}_{s_1}$, with
$\operatorname{col}_{s_1}$ the map of Proposition~\ref{dgr:prop:s-column} for $w_1$ and $(s_2)^{\otimes q}$ from the
completion of $(\Lambda^m_\pi G_1)^{\otimes_\pi q}$ to $\widehat\otimes^n_\pi G_2$, rows in the order $1,\dots,q$.
By~\eqref{dgr:eq:iterated} and naturality of $\wedge^{[q,m]}$,
$\widehat\wedge^n_{G_2}(s_2)^{\otimes q}=\widehat\wedge^{[q,m]}_{G_2}(\Lambda^m_\pi w_2)^{\otimes q}
=\widehat\Lambda^n_\pi w_2\circ\widehat\wedge^{[q,m]}_{G_1}$, so Proposition~\ref{dgr:prop:s-column} gives
$\widehat\wedge^n_{G_2}s=\Lambda^n_\pi(w_2w_1)$, and $\norm s\le N_{q,m}\norm{s_1}^m\norm{s_2}^q/\abs N$.

(2) Let $\rho=(\eta'N)^{-1}\operatorname{col}^{\rho_2}\circ(\rho_1)^{\otimes q}$ on
$\otimes^n_\pi E=(\otimes^m_\pi E)^{\otimes_\pi q}$, with $\operatorname{col}^{\rho_2}$ the map of
Proposition~\ref{dgr:prop:r-column} for $w_2$. By~\eqref{dgr:eq:iterated}, $\rho_1\alt^m_E=\Lambda^m_\pi w_1$ and
naturality of $\sh^{[q,m]}$, the map $\rho$ satisfies $\rho\,\alt^n_E=(\eta'N)^{-1}\operatorname{col}^{\rho_2}\sh^{[q,m]}_{G_1}\Lambda^n_\pi w_1
=\Lambda^n_\pi(w_2w_1)$, and $\norm\rho\le(m!)^{q-1}\norm{\rho_2}^m\norm{\rho_1}^q/\abs N$. Uncompleted data give
uncompleted maps, by the last assertions of Propositions~\ref{dgr:prop:s-column} and~\ref{dgr:prop:r-column}.
\end{proof}

\begin{proof}[Proof of Theorem~\ref{cmb:thm:column}]
Let $s_p$ be a section for $\id_Y$ in degree $p$. If $s_{p^j}$ is one in degree $p^j$, Theorem~\ref{cmb:thm:composition}(1)
with $w_1=w_2=\id_Y$, $q=p$ and $m=p^j$, which applies by Lemma~\ref{dgr:lem:unit}, gives one in degree $p^{j+1}$. So
$(\mathrm S^{p^j}_Y)$ holds for every $j\ge0$. By its base-$p$ expansion, every $n\ge1$ is a carry-free sum of powers of
$p$, so $(\mathrm S^n_Y)$ holds by the Lucas step for split operators in Section~\ref{cmb:ssec:degrees}. For retractions
use Theorem~\ref{cmb:thm:composition}(2) with $w_1=w_2=\id_Y$, $m=p^j$ and $q=p$. Uncompleted data stay uncompleted at each
step.
\end{proof}

\begin{proof}[Proof of Corollary~\ref{cmb:cor:composition}]
(1) Let $w=w_h\cdots w_1$ with each $w_l$ $\mathrm S^p$-split, and put $\gamma_l=w_l\cdots w_1$. Then $\gamma_l$ is
$\mathrm S^{p^l}$-split for $l\le h$: for $l=1$ by hypothesis, and for $l\ge2$ by Theorem~\ref{cmb:thm:composition}(1)
applied to $\gamma_{l-1}$ and $w_l$ with $q=p^{l-1}$ and $m=p$, by Lemma~\ref{dgr:lem:unit}. So $w=(w_h\cdots w_{l+1})\gamma_l$ is
$\mathrm S^{p^l}$-split for $l\le h$, by Proposition~\ref{ob:prop:split-operators}(2). Every $n<p^{h+1}$ is a carry-free
sum of powers $p^l$ with $l\le h$, so $w$ is $\mathrm S^n$-split by the Lucas step for split operators in
Section~\ref{cmb:ssec:degrees}. For retractions use Theorem~\ref{cmb:thm:composition}(2) with the inner map
$\gamma_{l-1}$, $m=p^{l-1}$, and the outer map $w_l$, $q=p$.
(2) If $w=v_2v_1$ with $v_1\colon E\to Y$, $v_2\colon Y\to G$ and $(\mathrm S^p_Y)$, then $\id_Y$ is $\mathrm S^n$-split for
every $n$ by Theorem~\ref{cmb:thm:column}, hence so is $w=v_2\,\id_Y\,v_1$ by Proposition~\ref{ob:prop:split-operators}(2).
The same argument applies to $(\mathrm R^p_Y)$.
\end{proof}

\subsection{Polarizability in every degree}\label{dgr:ssec:polarizable}

A bounded linear map $v\colon P\to P'$ of normed spaces is \emph{polarizable in degree $n$} if there is a bounded
$n$-linear map $\Pi\colon P^n\to\hGamma^n_\pi P'$ with $\Pi(h,\dots,h)=(vh)^{\otimes n}$ for $h\in P$; by
Lemma~\ref{ob:lem:parameter}, $\id_P$ is polarizable in degree $n$ if and only if $P$ is. Each isometry $\tau_\sigma$ of
Section~\ref{sec:canonical} fixes every pure power, hence every point of $\hGamma^n_\pi P'$. Every tensor $y$ in the
algebraic tensor power $\otimes^n_KP'$ with $\tau_\sigma y=y$ for all $\sigma\in S_n$ is a linear combination of pure
powers: if $e_1,\dots,e_d$ is a basis of a finite-dimensional subspace containing the factors of $y$, then $y$ is a
combination of the sums of the tensors $e_{i_1}\otimes\dots\otimes e_{i_n}$ over the rearrangements of a fixed index tuple,
and each such sum is a coefficient of the polynomial $\lambda\mapsto(\sum_i\lambda_ie_i)^{\otimes n}$ on $K^d$. For
$a,b\ge0$, $y\in\widehat\otimes^a_\pi P'$ and $z\in\widehat\otimes^b_\pi P'$, the \emph{shuffle product}
$y*z\in\widehat\otimes^{a+b}_\pi P'$ is the sum of the $\binom{a+b}a$ tensors obtained from $y\otimes z$ by distributing the
factors of $y$ and of $z$ over the $a+b$ slots in their orders; it is bilinear with
$\norm{y*z}\le\binom{a+b}a\norm y\norm z$.

\begin{lemma}[polarizable maps]\label{dgr:lem:polarizable}
Let $v\colon P\to P'$ be a bounded linear map of normed spaces.
\begin{enumerate}
\item The shuffle product maps $\hGamma^a_\pi P'\times\hGamma^b_\pi P'$ into $\hGamma^{a+b}_\pi P'$. If $v$ is polarizable
in degrees $a,b\ge1$, with $\Pi_a$ and $\Pi_b$, and $\binom{a+b}a\ne0$ in $K$, then $v$ is polarizable in degree $a+b$, with
$(h_1,\dots,h_{a+b})\mapsto\binom{a+b}a^{-1}\Pi_a(h_1,\dots,h_a)*\Pi_b(h_{a+1},\dots,h_{a+b})$.
\item If $n!\ne0$ in $K$, then $v$ is polarizable in degree $n$. Hence, if $\operatorname{char}K=p>0$, the degrees in which
$v$ is polarizable contain $1,\dots,p-1$ and are closed under carry-free sums.
\end{enumerate}
\end{lemma}

\begin{proof}
(1) For $h,h'\in P'$, $h^{\otimes a}*h'^{\otimes b}$ is the coefficient of $\lambda^a$ in
$(\lambda h+h')^{\otimes(a+b)}$, hence a combination of pure powers; by bilinearity and continuity the shuffle product maps
$\hGamma^a_\pi P'\times\hGamma^b_\pi P'$ into $\hGamma^{a+b}_\pi P'$. The stated map has diagonal
$\binom{a+b}a^{-1}(vh)^{\otimes a}*(vh)^{\otimes b}=(vh)^{\otimes(a+b)}$. (2) By Proposition~\ref{can:prop:natural}(2),
$R_n=\frac1{n!}\sum_\sigma\tau_\sigma$ is a projection onto $\hGamma^n_\pi P'$, and
$\Pi(h_1,\dots,h_n)=R_n(vh_1\otimes\dots\otimes vh_n)$ has diagonal $(vh)^{\otimes n}$. By Lucas's theorem,
$\binom{a+b}a\ne0$ in $K$ when $a$ and $b$ add without carry in base $p$, and~(1) gives the last assertion.
\end{proof}

\begin{lemma}[composition of polarizable maps]\label{dgr:lem:polarizable-composition}
Let $q,m\ge1$ and $n=qm$, and suppose that the number $N_{m,q}$ of partitions of $\{1,\dots,n\}$ into $q$ blocks of size
$m$ is nonzero in $K$. If $v_1\colon P\to P'$ is polarizable in degree $q$, with $\Pi_1$, and $v_2\colon P'\to P''$ is
polarizable in degree $m$, with $\Pi_2$, then $v_2v_1$ is polarizable in degree $n$, with a map of norm at most
$N_{m,q}\norm{\Pi_1}^m\norm{\Pi_2}^q/\abs{N_{m,q}}$.
\end{lemma}

\begin{proof}
\emph{Step 1 (the map).} Use the rows, columns and $\tau_\zeta$ of Subsection~\ref{dgr:ssec:columns}, with $P'$ and
$P''$ in place of $G$. The stabilizer $S_{\mathrm{row}}\subseteq S_n$ of the partition into rows consists of the
permutations inside rows composed with permutations of whole rows, and it has $N_{m,q}$ left cosets. For representatives $\sigma_1,\dots,\sigma_{N_{m,q}}$ put
$T=\sum_l\tau_{\sigma_l}$ on $\widehat\otimes^n_\pi P''$, of norm at most $N_{m,q}$. If $\tau_\sigma Y=Y$ for all
$\sigma\in S_{\mathrm{row}}$, then $TY$ does not depend on the representatives, and $\tau_\sigma TY=TY$ for all
$\sigma\in S_n$. For an array $(h_{i,t})$ in $P$ let $c_t=\Pi_1(h_{1,t},\dots,h_{q,t})\in\hGamma^q_\pi P'$ and
\[
 Y(h)=(\Pi_2^\otimes)^{\otimes q}\,\tau_\zeta(c_1\otimes\dots\otimes c_m),\qquad \Pi(h)=N_{m,q}^{-1}\,T\,Y(h),
\]
where the linearization $\Pi_2^\otimes\colon\widehat\otimes^m_\pi P'\to\hGamma^m_\pi P''$ of $\Pi_2$ acts on each row.
The map $\Pi$ is $n$-linear with the stated bound.

\emph{Step 2 (invariance).} A permutation of whole rows permutes the $q$ factors of every column block in the same way,
which fixes each $c_t$, so it fixes $\tau_\zeta(c_1\otimes\dots\otimes c_m)$. The operator $(\Pi_2^\otimes)^{\otimes q}$
commutes with these permutations, and its values are fixed by permutations inside rows, since $\Pi_2^\otimes$ takes values
in $\hGamma^m_\pi P''$. Hence $\tau_\sigma Y(h)=Y(h)$ for $\sigma\in S_{\mathrm{row}}$.

\emph{Step 3 (the diagonal).} If every $h_{i,t}$ is $h$, then $c_t=(v_1h)^{\otimes q}$,
$\tau_\zeta(c_1\otimes\dots\otimes c_m)=(v_1h)^{\otimes n}$, and $\Pi_2^\otimes$ sends each row $(v_1h)^{\otimes m}$ to
$(v_2v_1h)^{\otimes m}$. So $Y(h)=(v_2v_1h)^{\otimes n}$, $TY(h)=N_{m,q}(v_2v_1h)^{\otimes n}$ and
$\Pi(h,\dots,h)=(v_2v_1h)^{\otimes n}$.

\emph{Step 4 (values).} Fix the array $h$, and choose $c^{(j)}_t\to c_t$ in the linear span of the pure powers in
$\otimes^q_\pi P'$. Expanding, $\tau_\zeta(c^{(j)}_1\otimes\dots\otimes c^{(j)}_m)$ is a linear combination of tensors
$\tau_\zeta(h'^{\otimes q}_1\otimes\dots\otimes h'^{\otimes q}_m)=y^{\otimes q}$ with $y=h'_1\otimes\dots\otimes h'_m$, and
$(\Pi_2^\otimes)^{\otimes q}(y^{\otimes q})=z^{\otimes q}$ with $z=\Pi_2^\otimes(y)\in\hGamma^m_\pi P''$. For $z_l\to z$ in
the linear span of the pure powers, $z_l^{\otimes q}$ lies in the algebraic tensor power and is fixed by $\tau_\sigma$ for
$\sigma\in S_{\mathrm{row}}$, so $Tz_l^{\otimes q}$ is fixed by every $\tau_\sigma$ and is a linear combination of pure
powers; hence $Tz^{\otimes q}=\lim_lTz_l^{\otimes q}$ lies in $\hGamma^n_\pi P''$. By linearity and continuity, so does
$TY(h)$.
\end{proof}

\begin{proof}[Proof of Theorem~\ref{cmb:thm:polarizable-height}]
If $P$ is polarizable in degrees $p$ and $p^i$, then Lemma~\ref{dgr:lem:polarizable-composition} with $v_1=v_2=\id_P$,
$q=p$ and $m=p^i$, which applies by Lemma~\ref{dgr:lem:unit}, makes $P$ polarizable in degree $p^{i+1}$. By induction $P$
is polarizable in every degree $p^i$, and Lemma~\ref{dgr:lem:polarizable} gives every degree, since every $n\ge1$ is a
carry-free sum of powers of $p$.
\end{proof}

\begin{proof}[Proof of Corollary~\ref{cmb:cor:parameter-degrees}]
(1) The parameter space of a parameter piece of degree $p$ is polarizable in degree $p$, hence in every degree by
Theorem~\ref{cmb:thm:polarizable-height}; so the piece, and every piece of a generalized decomposition of degree $p$ into
parameter pieces, is a parameter piece of every degree. (2) By Theorem~\ref{cmb:thm:polarizable-height}, $\cL(E,D)$ is
polarizable in every degree, and Corollary~\ref{ob:cor:mechanisms} applies in each degree.
\end{proof}

\subsection{Decompositions of degree \texorpdfstring{$p$}{p}}\label{dgr:ssec:decompositions}

\begin{proof}[Proof of Theorem~\ref{cmb:thm:mixed-degrees}]
Let $w$ be $\mathrm S^{\le p}$-split, as every $\mathrm S^p$-split map is, since every map is $\mathrm S^n$-split for $n<p$.
Every $n<p^2$ is $c_1p+c_0$ with base-$p$ digits $c_0,c_1<p$, a carry-free sum of copies of $p$ and of $1$, so $w$ is
$\mathrm S^n$-split by the Lucas step for split operators in Section~\ref{cmb:ssec:degrees}. Thus $w$ is
$\mathrm S^{\le p^2-1}$-split, and the same holds for $\mathrm R$. Hence every piece of the decomposition, including the
single piece of a relative splitting, is of the same type in every degree $n<p^2$. Theorem~\ref{cmb:thm:decomposition}(3)
gives the goodness; the Lucas step keeps uncompleted data uncompleted.
\end{proof}

\begin{proof}[Proof of Proposition~\ref{cmb:prop:target-pieces}]
Choose $f\in E^*$ and $x_0\in E$ with $f(x_0)=1$, and let the pieces be $(G_i,G'_i,w_i,w'_i,\varepsilon_i)$. For $y\in D$,
the operator $u_y=f(\cdot)\,y\in\cL(E,D)$ has $\norm{u_y}\le\norm f\norm y$, and~\eqref{cmb:eq:decomposition} gives
$y=u_yx_0=\sum_iw'_i\chi_i(y)$ with the bounded linear maps $\chi_i\colon D\to G'_i$, $\chi_i(y)=\varepsilon_i(u_y)w_ix_0$.
So the inclusion $\iota_D$ of $D$ in $\widehat D$ is $\sum_iw'_i\chi_i$, which is $\mathrm R^{\le k}$-split, since every
$w'_i$ is, by Proposition~\ref{ob:prop:split-operators}(2) and the paragraph on split operators in
Section~\ref{cmb:ssec:degrees}. A
retraction for $\iota_D$ in degree $k$ is a bounded linear
$\rho\colon\otimes^k_\pi D\to\widehat\Lambda^k_\pi\widehat D=\widehat\Lambda^k_\pi D$ such that $\rho\alt^k_D$ is the
inclusion of $\Lambda^k_\pi D$ in its completion; this is $(\mathrm R^k_D)$. Without completion, the maps $w'_i$ take values
in $D$ and their retractions in the uncompleted exterior powers, and so does $\rho$. For $k=p$,
Corollary~\ref{cmb:cor:column} gives the last assertion.
\end{proof}

\subsection{Permanence and its failure}\label{dgr:ssec:permanence}

\begin{proof}[Proof of Theorem~\ref{cmb:thm:cantor}]
\emph{Step 0 (notation).} Order the words of each length $m$ lexicographically. For a normed space $U$, a step function of
level $m$ on $\Omega^k$ is a function $z=\sum_\xi1_{C_{\xi_1}\times\dots\times C_{\xi_k}}u_\xi$, the sum over $k$-tuples
$\xi$ of words of length $m$, with $u_\xi\in U$; put $\norm z_1=\sum_\xi2^{-mk}\norm{u_\xi}$ and
$\norm z_\infty=\max_\xi\norm{u_\xi}$, and let $L^1(\Omega^k;U)$ be the completion of the step functions, in which the
$U_0$-valued ones are dense if $U_0$ is dense in $U$. Let $\Delta_m\subseteq\Omega^k$ be the union of the cells
$C_{\xi_1}\times\dots\times C_{\xi_k}$ of level $m$ with $\xi_1<\dots<\xi_k$, and $T_m$ the union of those with two equal
words; $T_m$ has measure at most $\binom k22^{-m}$, so $\norm{1_{T_m}z}_1\le\binom k22^{-m}\norm z_\infty$ for step
functions $z$. Every extension of a word $\xi$ precedes every extension of a later word of the same length, so
$\Delta_m\subseteq\Delta_{m'}$ and $\Delta_{m'}\setminus\Delta_m\subseteq T_m$ for $m\le m'$.

\emph{Step 1 (tensor powers).} For a finite set $I$ with weights $c_t>0$, the map $\sum_te_t\otimes u_t\mapsto(u_t)_t$ is
an isometry of $\ell^1(I,c)\otimes_\pi U$ onto $U^I$ with the norm $\sum_tc_t\norm{u_t}$, since this decomposition costs
$\sum_tc_t\norm{u_t}$ and $(\eta,u)\mapsto(\eta_tu)_t$ is contractive. With associativity and commutativity of
$\otimes_\pi$, the contractive $k$-linear map $(f_1,\dots,f_k)\mapsto[\zeta\mapsto f_1(\zeta_1)\otimes\dots\otimes
f_k(\zeta_k)]$ therefore linearizes to an isometry of the $k$th projective tensor power of the step functions of level $m$
onto the $\otimes^k_\pi G$-valued step functions of level $m$ on $\Omega^k$; since the tensor power at level $m$ maps
contractively into that of all step functions, it is an isometry of the latter onto all $\otimes^k_\pi G$-valued step
functions. Passing to completions, $\widehat\otimes^k_\pi Y=L^1(\Omega^k;\widehat\otimes^k_\pi G)$; then
$\widehat\wedge^k_Y$ sends $1_{C_1\times\dots\times C_k}(g_1\otimes\dots\otimes g_k)$ to
$(1_{C_1}g_1)\wedge\dots\wedge(1_{C_k}g_k)$, and $\alt^k_Y$ sends $f_1\wedge\dots\wedge f_k$ to
$\zeta\mapsto\sum_\sigma\sgn(\sigma)f_{\sigma(1)}(\zeta_1)\otimes\dots\otimes f_{\sigma(k)}(\zeta_k)$.

\emph{Step 2 (ordered antisymmetrization).} For step functions $f_a=\sum_\xi1_{C_\xi}g_{a,\xi}$ of level at most $m$,
let $\Phi_m(f)$ be the step function equal to $\sum_\sigma\sgn(\sigma)f_{\sigma(1)}(\zeta_1)\otimes\dots\otimes
f_{\sigma(k)}(\zeta_k)$ on $\Delta_m$ and to $0$ elsewhere. Since $(\xi,\sigma)\mapsto\upsilon$ with
$\upsilon_{\sigma(a)}=\xi_a$ is a bijection from the increasing $k$-tuples of words times $S_k$ onto the $k$-tuples of
distinct words, $\norm{\Phi_m(f)}_1\le\sum_\upsilon2^{-mk}\prod_b\norm{g_{b,\upsilon_b}}\le\prod_b\norm{f_b}_1$. For
$m'\ge m$, $\Phi_{m'}(f)$ and $\Phi_m(f)$ agree outside $T_m$ by Step~0, so
$\norm{\Phi_{m'}(f)-\Phi_m(f)}_1\le k!\binom k22^{-m}\prod_a\norm{f_a}_\infty$. Hence $\Phi(f)=\lim_m\Phi_m(f)$ exists; it
is $k$-linear, alternating, since the terms $\sigma$ and $\sigma\circ(a\ b)$ cancel when $f_a=f_b$, and bounded by
$\prod_a\norm{f_a}_1$, so it extends to a contractive alternating map on $Y^k$.

\emph{Step 3 (sections).} For the $f_a$ of Step~2, $f_1\wedge\dots\wedge f_k$ is the sum over $k$-tuples of words of
$\bigwedge_a1_{C_{\xi_a}}g_{a,\xi_a}$. A tuple of distinct words is $(\xi_{\sigma^{-1}(1)},\dots,\xi_{\sigma^{-1}(k)})$
for a unique increasing tuple $\xi$ and a unique $\sigma\in S_k$, and its wedge is
$\sgn(\sigma)\bigwedge_b1_{C_{\xi_b}}g_{\sigma(b),\xi_b}$; by Step~1 these terms add up to $\widehat\wedge^k_Y\Phi_m(f)$.
Since $\norm{1_{C_\xi}g}_1=2^{-m}\norm g$, the tuples with $\xi_a=\xi_b$ for some $a<b$ contribute at most
$\sum_{a<b}2^{-m}\norm{f_a}_\infty\norm{f_b}_1\prod_{c\ne a,b}\norm{f_c}_1$ in norm. Letting $m\to\infty$ gives
$f_1\wedge\dots\wedge f_k=\widehat\wedge^k_Y\Phi(f)$. By Proposition~\ref{prop:exterior}, $\Phi$ induces
$s\colon\Lambda^k_\pi Y\to\widehat\otimes^k_\pi Y$ with $\norm s\le1$, and $\widehat\wedge^k_Ys=\id$ on wedges of step
functions, hence everywhere, since these span a dense subspace (the paragraph before Theorem~\ref{ob:thm:mechanisms}).

\emph{Step 4 (retractions).} For a step function $z$ of level at most $m$ and $m'\ge m$,
$\norm{1_{\Delta_{m'}}z-1_{\Delta_m}z}_1\le\binom k22^{-m}\norm z_\infty$ by Step~0. So $\cP_\Delta z=\lim_m1_{\Delta_m}z$
exists, with $\norm{\cP_\Delta z}_1\le\norm z_1$, and $\cP_\Delta$ extends to a contraction of
$L^1(\Omega^k;\widehat\otimes^k_\pi G)=\widehat\otimes^k_\pi Y$. Let $\rho=\widehat\wedge^k_Y\circ\cP_\Delta$ on
$\otimes^k_\pi Y$, so $\norm\rho\le1$. For step functions $f_a$, Step~1 gives
$1_{\Delta_m}\alt^k_Y(f_1\wedge\dots\wedge f_k)=\Phi_m(f)$, so
$\rho\alt^k_Y(f_1\wedge\dots\wedge f_k)=\widehat\wedge^k_Y\Phi(f)=f_1\wedge\dots\wedge f_k$ by Step~3; by density
$\rho\alt^k_Y=\id$.

\emph{Step 5 (polarization).} On step functions $\tau_\sigma z(\zeta)=\tau_\sigma\bigl(z(\zeta_{\sigma(1)},\dots,
\zeta_{\sigma(k)})\bigr)$, so $\tau_\sigma(1_Sz)=1_{\sigma S}\tau_\sigma z$ with
$\sigma S=\{\zeta:(\zeta_{\sigma(1)},\dots,\zeta_{\sigma(k)})\in S\}$; the sets $\sigma\Delta_m$, $\sigma\in S_k$, are
disjoint with union $\Omega^k\setminus T_m$. Put $R_k=\sum_\sigma\tau_\sigma\cP_\Delta$, of norm at most $k!$. For a step
function $h$, $h^{\otimes k}$ is fixed by every $\tau_\sigma$, so
$R_kh^{\otimes k}=\lim_m1_{\Omega^k\setminus T_m}h^{\otimes k}=h^{\otimes k}$; by continuity $R_k$ fixes every point of
$\hGamma^k_\pi Y$. For cells $C_1<\dots<C_k$ of level $m$ and $g_a\in G$,
$\sum_\sigma\tau_\sigma(1_{C_1}g_1\otimes\dots\otimes1_{C_k}g_k)$ is the coefficient of $\lambda_1\cdots\lambda_k$ in
$(\sum_a\lambda_a1_{C_a}g_a)^{\otimes k}$, hence a linear combination of pure powers. For a $\otimes^k_\pi G$-valued step
function $z$ of level at most $m$, $1_{\Delta_m}z$ is a finite sum of such tensors, so
$R_kz=\lim_m\sum_\sigma\tau_\sigma1_{\Delta_m}z$ lies in $\hGamma^k_\pi Y$; by continuity $R_k$ maps
$\widehat\otimes^k_\pi Y$ onto $\hGamma^k_\pi Y$.
\end{proof}

\begin{proof}[Proof of Corollary~\ref{cmb:cor:no-hereditary}]
(1) The isometry $g\mapsto1_\Omega g$ has complete, hence closed, image. By Theorem~\ref{cmb:thm:cantor},
$L^1(\Omega;E)$ and $L^1(\Omega;D)$ satisfy $(\mathrm S^n)$ and $(\mathrm R^n)$ for every $n$, and
Proposition~\ref{fib:prop:split}(1),(2) gives the goodness statements.

(2) The map $q_E\colon\ell^1(\operatorname{Ball}(E),\Kh)\to E$, $e_x\mapsto x$, is contractive. Fix $c\in K$ with
$0<\abs c<1$. For $x\ne0$ in $E$ choose $j\in\Z$ with $\abs c^{j+1}<\norm x\le\abs c^j$; then $x/c^j\in\operatorname{Ball}(E)$,
$q_E(c^je_{x/c^j})=x$ and $\norm{c^je_{x/c^j}}_1=\abs c^j<\abs c^{-1}\norm x$. So $q_E$ maps each open ball of radius
$\delta$ onto a set containing the open ball of radius $\abs c\delta$; it is surjective and open, and $E$ is isomorphic to
$\ell^1(\operatorname{Ball}(E),\Kh)/\ker q_E$. The unit vectors of $\ell^1(\operatorname{Ball}(E),\Kh)$ are linearly
independent over $K$ with dense $K$-span, on which $\norm{\cP_<(x\otimes y)}\le\norm x_1\norm y_1$ for every total order
of $\operatorname{Ball}(E)$. So Proposition~\ref{spl:prop:ordering} gives $(\mathrm S^n)$, $(\mathrm R^n)$ and
polarizability in degree $n$ for every $n$, and Proposition~\ref{fib:prop:split}(1),(2) gives the goodness statements.
The same holds for $D$.

(3) By Theorem~\ref{fib:thm:ell2} with $r=\infty$ and Corollary~\ref{cls:cor:c0-not-polarizable},
$(c_0(\N,K),c_0(\N,K))$ is not good for Banach values in degree $k$, and $c_0(\N,K)$ satisfies neither $(\mathrm S^k)$ nor
$(\mathrm R^k)$ and is not polarizable in degree $k$. Proposition~\ref{cmb:prop:permanence}(3) transfers the last three
statements to the completion $Y_0=c_0(\N,\Kh)$. If $(Y_0,Y_0)$ were good for Banach values in degree $k$, then
Proposition~\ref{cmb:prop:permanence}(4), applied to the source with the complete target $Y_0$ and then to the target,
would make $(c_0(\N,K),c_0(\N,K))$ good for Banach values in degree $k$. By~(1) and~(2) with $E=D=Y_0$, $Y_0$ is a closed
subspace of $L^1(\Omega;Y_0)$ and isomorphic to a quotient of $\ell^1(\operatorname{Ball}(Y_0),\Kh)$; these two spaces
satisfy the three conditions in degree $k$, and replacing either entry of $(Y_0,Y_0)$ by one of them gives a pair that is
good for Banach values in degree $k$. Since the three conditions and goodness for Banach values are invariant under
isomorphisms (Proposition~\ref{cmb:prop:permanence}(1)), the last assertion follows.
\end{proof}
%% ---- end app-L-degrees.tex
%% ---- begin app-M-rigid.tex
\section{Rigid Banach spaces}\label{app:rigid}

This appendix proves Theorems~\ref{bey:thm:rigid}, \ref{bey:thm:rigid-good} and~\ref{bey:thm:rigid-sum}. It transfers
the rigid subspace of $L_0$ of Kalton and Roberts~\cite[\S3]{KR}: their quasi-norms become $r$\nobreakdash-norms, their
exponents increasing to $1$ become exponents $r_i\uparrow\infty$, and forcing goes through $L^r$ readouts.

\subsection{Blocks and the Nakano sum}\label{rig:ssec:nakano}

Throughout Subsections~\ref{rig:ssec:nakano}--\ref{rig:ssec:spaces}, $K$ is a complete nonarchimedean field with
nontrivial absolute value that is separable as a metric space, and $\vartheta\in K$ is fixed with $0<\abs\vartheta<1$.

For a finite-dimensional normed space $Y$ and $1\le r<\infty$, $L^r((0,1);Y)$ is the completion of the step functions
$f\colon(0,1)\to Y$ modulo null functions for $\norm f_r=(\int_0^1\norm{f(t)}^r\,dt)^{1/r}$, and $1\otimes y$ is the
constant function with value $y$. By the Riesz--Fischer argument, $Y$ being complete (Lemma~\ref{lem:findim}), each
element is a measurable function, unique up to null sets, with the same formula for the norm. If the norm of $Y$ is
nonarchimedean, $\norm\cdot_r$ is an $r$\nobreakdash-norm (Section~\ref{sec:positive}), since
$\norm{f(t)+g(t)}^r\le\max(\norm{f(t)},\norm{g(t)})^r\le\norm{f(t)}^r+\norm{g(t)}^r$. An $r$\nobreakdash-norm is an
$r'$\nobreakdash-norm for $1\le r'\le r$, since $(c^r+c'^r)^{1/r}\le(c^{r'}+c'^{r'})^{1/r'}$ for reals $c,c'\ge0$. The
space $K^{1+N}$ carries the maximum norm, with unit vectors $e_1,\dots,e_{1+N}$.

Let $1\le r_0\le r_1\le\cdots$ be real numbers. For a sequence $a=(a_i)_{i\ge0}$ of reals $a_i\ge0$ put
\[
 \nu(a)=\inf\Bigl\{\varsigma>0:\sum_{i\ge0}(a_i/\varsigma)^{r_i}\le1\Bigr\}\in[0,\infty].
\]
Convexity of $\varsigma\mapsto\varsigma^{r_i}$, with the weights $\varsigma/(\varsigma+\varsigma')$ and
$\varsigma'/(\varsigma+\varsigma')$, gives $\nu(a+a')\le\nu(a)+\nu(a')$. Moreover $\nu$ is positively homogeneous and monotone, $\nu(a)=0$ only for $a=0$,
$\nu(a)=a_i$ if $a_i$ is the only nonzero entry, and $\nu(a)=\sup_h\nu(a^{\le h})$, where $a^{\le h}$ replaces the
entries with index above $h$ by zero. For finitely supported $a$ the sum is continuous and decreasing in $\varsigma$, so
\begin{equation}\label{rig:eq:modular}
 \nu(a)\le\varsigma\iff\textstyle\sum_i(a_i/\varsigma)^{r_i}\le1\qquad(\varsigma>0).
\end{equation}

Given Banach spaces $L_i$, $i\ge0$, the \emph{Nakano sum} $\mathcal Z$ is the normed space of the sequences
$f=(f_i)_{i\ge0}$, $f_i\in L_i$, with $\norm f:=\nu((\norm{f_i})_{i\ge0})<\infty$. It is complete: a Cauchy sequence
converges in each $L_i$, since $\norm{f_i}\le\norm f$, and to its coordinatewise limit in $\mathcal Z$, since
$\nu(a)=\sup_h\nu(a^{\le h})$. We identify $L_i$ isometrically with the $i$\nobreakdash-th block of $\mathcal Z$. The
projections $\Pi_i$, $\Pi_{<i}$, $\Pi_{\le i}$ and $\Pi_{>i}$ keep the coordinates with index equal to, less than, at
most and greater than $i$ and replace the others by zero; they are contractive, since $\nu$ is monotone, and
$\id=\Pi_{<i}+\Pi_i+\Pi_{>i}$. Let $\mathcal Z_{<i}$, $\mathcal Z_{\le i}$ and $\mathcal Z_{>i}$ be the ranges of the
last three.

\begin{lemma}\label{rig:lem:nakano}
If the norm of every $L_m$ is an $r_m$\nobreakdash-norm and $i\ge0$, then
$\norm{f^{(1)}+\dots+f^{(d)}}^{r_{i+1}}\le\norm{f^{(1)}}^{r_{i+1}}+\dots+\norm{f^{(d)}}^{r_{i+1}}$ for
$f^{(1)},\dots,f^{(d)}\in\mathcal Z_{>i}$.
\end{lemma}

\begin{proof}
By induction it suffices to take $d=2$, $f=f^{(1)}$ and $g=f^{(2)}$. Put $r=r_{i+1}$ and, for $m>i$, $a_m=\norm{f_m}$,
$a'_m=\norm{g_m}$ and $a''_m=\norm{f_m+g_m}$; then $a''^r_m\le a^r_m+a'^r_m$, since the norm of $L_m$ is an
$r_m$\nobreakdash-norm and $r_m\ge r$. Let $\nu'$ be defined like $\nu$ on sequences indexed by $m>i$, with the
exponents $r_m/r\ge1$; like $\nu$, it is monotone and subadditive. With $a^r=(a^r_m)_{m>i}$, the substitution
$\varsigma=\varsigma'^r$ gives $\nu'(a^r)=\inf\{\varsigma'^r:\sum_{m>i}(a_m/\varsigma')^{r_m}\le1\}=\nu(a)^r$. Hence
\[
 \norm{f+g}^r=\nu'(a''^r)\le\nu'(a^r+a'^r)\le\nu'(a^r)+\nu'(a'^r)=\norm f^r+\norm g^r.
\]
\end{proof}

\subsection{Operators between blocks}\label{rig:ssec:operators}

\begin{lemma}[compare~{\cite[Lemma~2.2]{KR}}]\label{rig:lem:rolewicz}
Let $1\le r\le r'$, let $Y$ be a finite-dimensional normed space, $N'\ge0$ an integer and $h\in L^r((0,1);Y)$ a step
function. There is a bounded linear map $\Phi\colon L^{r'}((0,1);K^{1+N'})\to L^r((0,1);Y)$ with
$\Phi(1\otimes e_1)=h$ and $\norm{\Phi g}_r\le\norm h_r\norm g_{r'}$ for every $g$.
\end{lemma}

\begin{proof}
If $h=0$, take $\Phi=0$. Otherwise let $e^*_1\colon K^{1+N'}\to K$ be the first coordinate, so $e^*_1(e_1)=1$ and
$\abs{e^*_1(y)}\le\norm y$. The function $\xi(t)=\norm h_r^{-r}\int_0^t\norm{h(t')}^r\,dt'$ maps $[0,1]$ onto $[0,1]$
and is continuous, nondecreasing and piecewise affine, with positive slope exactly where $h\ne0$. For a step function
$g$ the step function $(\Phi g)(t)=e^*_1(g(\xi(t)))\,h(t)$ depends only on the class of $g$, since $\Phi g=0$ where
$h=0$ and preimages of null sets under $\xi$ are null where $h\ne0$. On the intervals where $h\ne0$, which $\xi$ maps
onto $(0,1)$ up to finitely many points, the substitution $t''=\xi(t)$ has $\norm{h(t)}^r\,dt=\norm h_r^r\,dt''$. Hence
\[
 \norm{\Phi g}_r^r=\norm h_r^r\int_0^1\abs{e^*_1(g(t''))}^r\,dt''\le\norm h_r^r\norm g_r^r\le\norm h_r^r\norm g_{r'}^r,
\]
the last step by Jensen's inequality on the probability space $(0,1)$, since $r\le r'$. Extend $\Phi$ by continuity;
then $\Phi(1\otimes e_1)=e^*_1(e_1)\,h=h$.
\end{proof}

From here to the end of Subsection~\ref{rig:ssec:spaces}, $N_i\ge i$ are integers, $L_i=L^{r_i}((0,1);K^{1+N_i})$ and
$\mathbf e_i=1\otimes e_1\in L_i$ for $i\ge0$, and $\mathcal Z$ is the Nakano sum of the $L_i$.

\begin{lemma}[compare~{\cite[Lemma~3.2]{KR}}]\label{rig:lem:onto-earlier}
Let $i\ge1$, and let $f\in\mathcal Z_{<i}$ have step functions as coordinates and $\norm f\le1$. There is a bounded
linear map $\Phi\colon L_i\to\mathcal Z_{<i}$ with $\Phi\mathbf e_i=f$ and $\norm\Phi\le1$.
\end{lemma}

\begin{proof}
By~\eqref{rig:eq:modular}, $\sum_{m<i}\norm{\Pi_mf}^{r_m}\le1$. For $m<i$, since $r_m\le r_i$,
Lemma~\ref{rig:lem:rolewicz} gives $\Phi_m\colon L_i\to L_m$ with $\Phi_m\mathbf e_i=\Pi_mf$ and
$\norm{\Phi_mg}\le\norm{\Pi_mf}\norm g$. Put $\Phi g=\sum_{m<i}\Phi_mg$, with $\Phi_mg$ in block $m$. Then
$\Phi\mathbf e_i=f$, and for $g\ne0$, $\sum_{m<i}(\norm{\Phi_mg}/\norm g)^{r_m}\le\sum_{m<i}\norm{\Pi_mf}^{r_m}\le1$, so
$\norm{\Phi g}\le\norm g$ by~\eqref{rig:eq:modular}.
\end{proof}

\subsection{The inductive choices}\label{rig:ssec:choices}

Fix an integer $h_0\ge0$ with $\sum_{h\ge0}\abs\vartheta^{h+h_0}\le1/4$ and put $\zeta_h=\vartheta^{h+h_0}$ for $h\ge0$.
Given reals $\eta_i>0$ and integers $d_i\ge1$, let $\mathbf e_{i,l}\in L_i$, $1\le l\le d_i$,
equal $e_1$ on $[(l-1)/d_i,l/d_i)$ and $0$ elsewhere; then $\sum_l\mathbf e_{i,l}=\mathbf e_i$ and
$\norm{\mathbf e_{i,l}}=d_i^{-1/r_i}$. Let $Y_i=\spn(e_2,\dots,e_{1+N_i})\subseteq K^{1+N_i}$ and
\[
 \mathcal U_i=\spn\{\mathbf e_{i,1},\dots,\mathbf e_{i,d_i}\}+\{1\otimes y:y\in Y_i\}\subseteq L_i,\qquad
 \mathcal U_{<i}=\mathcal U_0+\dots+\mathcal U_{i-1}\subseteq\mathcal Z_{<i},
\]
finite-dimensional spaces of step functions. For $i\ge1$ let $C_i\ge1$ satisfy $\norm{\sum_lx_l}\le C_i\max_l\norm{x_l}$
for every finite family $(x_l)$ in $\mathcal U_{<i}$; it exists because the norm of $\mathcal U_{<i}$ is equivalent to
the maximum of the coordinates in a basis (Lemma~\ref{lem:findim}).

\begin{lemma}[compare~{\cite[Lemma~3.1]{KR}}]\label{rig:lem:choices}
Let real numbers $\underline r_i\ge1$, $i\ge0$, be given. One can choose $r_i$, $\eta_i$ and $d_i$ inductively so that
for every $i\ge0$:
\begin{enumerate}
\item $r_i\ge\max(r_{i-1},\underline r_i,i)$, where $r_{-1}=1$;
\item $d_i^{-1/r_i}<\eta_i$, that is, $\norm{\mathbf e_{i,l}}<\eta_i$ for every $l$;
\item $(\sum_{l=1}^{d_i}\norm{\mathbf e_{i,l}}^{r_{i+1}})^{1/r_{i+1}}<\eta_i$;
\item $iC_i\eta_i<\abs{\zeta_i}$ if $i\ge1$.
\end{enumerate}
\end{lemma}

\begin{proof}
At stage $i$ the data with indices below $i$ are fixed. If $i\ge1$, then $(\sum_l\norm{\mathbf e_{i-1,l}}^r)^{1/r}$, a
finite sum, tends to $\max_l\norm{\mathbf e_{i-1,l}}<\eta_{i-1}$ as $r\to\infty$, by~(2) for $i-1$; choose $r_i$
satisfying~(1) and so large that~(3) holds for $i-1$. This determines $\mathcal U_{<i}$ and its norm, which involves
only $r_0,\dots,r_{i-1}$, hence $C_i$.
Choose $\eta_i>0$ satisfying~(4), and then $d_i$ satisfying~(2).
\end{proof}

\subsection{The spaces \texorpdfstring{$\fX_\alpha$}{X\_alpha}}\label{rig:ssec:spaces}

Throughout this subsection the data of Subsection~\ref{rig:ssec:choices} are chosen by Lemma~\ref{rig:lem:choices} for
some lower bounds $\underline r_i$. The norm of each $L_m$ is an $r_m$\nobreakdash-norm, since $K^{1+N_m}$ is
nonarchimedean, so Lemma~\ref{rig:lem:nakano} applies.

\emph{Index sets.} Partition $\{0,1,2,\dots\}=\bigsqcup_{\alpha\ge1}\mathcal I_\alpha$ into infinite sets, enumerate
$\mathcal I_\alpha$ as $i_\alpha(0)<i_\alpha(1)<\cdots$, and put $\mathcal J_{i_\alpha(a)}=\{i_\alpha(2^a(2a'+1)):a'\ge0\}$
for $a\ge0$. Since every positive integer is $2^a(2a'+1)$ for exactly one pair $(a,a')$ and $2^a>a$, the sets
$\mathcal J_i$, $i\in\mathcal I_\alpha$, are pairwise disjoint infinite subsets of $\mathcal I_\alpha$ with union
$\mathcal I_\alpha\setminus\{\min\mathcal I_\alpha\}$, and $\min\mathcal J_i>i$.

\emph{Dense sequences.} For $i\in\mathcal I_\alpha$, the set $\mathcal F_i$ of the
$f\in\sum_{m\in\mathcal I_\alpha,\,m\le i}\mathcal U_m$ with $\abs\vartheta<\norm f\le1$ contains $\mathbf e_i$ and lies
in a finite-dimensional space over the separable field $K$; let $\mathcal F^0_i\ni\mathbf e_i$ be a countable dense
subset. Through a bijection of $\mathcal J_i$ onto $\mathcal F^0_i\times\N$, choose $\gamma_h\in\mathcal F^0_i$,
$h\in\mathcal J_i$, such that every element of $\mathcal F^0_i$ equals $\gamma_h$ for infinitely many $h$. Since
$\gamma_h\in\mathcal Z_{\le i}\subseteq\mathcal Z_{<h}$ has step functions as coordinates and norm at most $1$,
Lemma~\ref{rig:lem:onto-earlier} gives $\Phi_h\colon L_h\to\mathcal Z_{<h}$ with $\Phi_h\mathbf e_h=\gamma_h$ and
$\norm{\Phi_h}\le1$. For $h=\min\mathcal I_\alpha$ put $\Phi_h=0$.

\emph{The construction.} For $\alpha\ge1$ put $\Delta_\alpha=\sum_{h\in\mathcal I_\alpha}\zeta_h\Phi_h\Pi_h$; let
$\mathcal E_\alpha$ and $\mathcal Y_\alpha$ be the closed spans of $\{\mathbf e_i:i\in\mathcal I_\alpha\}$ and of
$\bigcup_{i\in\mathcal I_\alpha}\mathcal U_i$; put $\mathcal Q_\alpha=(\id-\Delta_\alpha)(\mathcal E_\alpha)$ and
$\fX_\alpha=\mathcal Y_\alpha/\mathcal Q_\alpha$; and let $[f]_\alpha$ be the class of $f\in\mathcal Z$ in
$\mathcal Z/\mathcal Q_\alpha$. The series $\Delta_\alpha$ converges absolutely, with norm at most
$\sum_h\abs{\zeta_h}\le1/4$, so the Neumann series inverts $\id-\Delta_\alpha$, with
$\norm{(\id-\Delta_\alpha)^{-1}}\le4/3$. Since $\Phi_h$ maps into $\mathcal Z_{<h}$, $\Delta_\alpha$ maps
$\mathcal Z_{\le i}$ into $\mathcal Z_{<i}$, so $\Pi_i\Delta_\alpha\Pi_{\le i}=\Pi_{>i}\Delta_\alpha\Pi_{\le i}=0$, for
every $i$. For $h\in\mathcal I_\alpha$ the vector $(\id-\Delta_\alpha)\mathbf e_h$ is $\mathbf e_h-\zeta_h\gamma_h$, or
$\mathbf e_h$ if $h=\min\mathcal I_\alpha$, and lies in $\mathcal Y_\alpha$. So $\mathcal Q_\alpha$, closed as the image
of a closed subspace under an isomorphism, lies in $\mathcal Y_\alpha$, and the Banach space $\fX_\alpha$ is the closed
subspace of $\mathcal Z/\mathcal Q_\alpha$ of the classes $[f]_\alpha$, $f\in\mathcal Y_\alpha$. For $i\in\mathcal I_\alpha$ and $h\in\mathcal J_i$,
\begin{equation}\label{rig:eq:quotient}
 [\mathbf e_h]_\alpha=\zeta_h[\gamma_h]_\alpha ,
\end{equation}
and $\Pi_i\mathcal E_\alpha$ lies in the line $K\mathbf e_i$, which is closed by Lemma~\ref{lem:findim}.

\begin{lemma}[compare~{\cite[Lemma~3.3]{KR}}]\label{rig:lem:projection}
If $i\in\mathcal I_\alpha$ and $f\in\mathcal Z_{\le i}$, then
$\inf_{\lambda\in K}\norm{\Pi_if-\lambda\mathbf e_i}\le\frac43\norm{[f]_\alpha}$.
\end{lemma}

\begin{proof}
Let $t>\norm{[f]_\alpha}$, choose $g\in\mathcal E_\alpha$ with $\norm{f-(\id-\Delta_\alpha)g}<t$, and let
$\Pi_ig=\lambda\mathbf e_i$. Since $\Pi_i\Delta_\alpha\Pi_{\le i}=0$, the vector
$\Pi_if-\lambda\mathbf e_i=\Pi_i(f-(\id-\Delta_\alpha)g)-\Pi_i\Delta_\alpha\Pi_{>i}g$ has norm less than
$t+\frac14\norm{\Pi_{>i}g}$. Since $\Pi_{>i}f=0$ and $\Pi_{>i}\Delta_\alpha\Pi_{\le i}=0$, the vector
$\Pi_{>i}((\id-\Delta_\alpha)g-f)=\Pi_{>i}g-\Pi_{>i}\Delta_\alpha\Pi_{>i}g$ has norm less than $t$, so
$\norm{\Pi_{>i}g}<t+\frac14\norm{\Pi_{>i}g}$, that is, $\norm{\Pi_{>i}g}<\frac43t$. Hence
$\norm{\Pi_if-\lambda\mathbf e_i}<\frac43t$; let $t$ decrease to $\norm{[f]_\alpha}$.
\end{proof}

\begin{lemma}[compare~{\cite[Lemma~3.5]{KR}}]\label{rig:lem:dense}
The set $\bigcup_{i\in\mathcal I_\alpha}K[\mathbf e_i]_\alpha$ is dense in $\fX_\alpha$.
\end{lemma}

\begin{proof}
Let $\mathcal C$ be the closure of this set; it is closed under scalar multiplication. Let $i\in\mathcal I_\alpha$. Each
$g\in\mathcal F^0_i$ equals $\gamma_h$ for some $h\in\mathcal J_i$, so $[g]_\alpha=\zeta_h^{-1}[\mathbf e_h]_\alpha\in\mathcal C$
by~\eqref{rig:eq:quotient}; by density, $[f]_\alpha\in\mathcal C$ for every $f\in\mathcal F_i$. A nonzero $f$ in the span
of $\bigcup_{m\in\mathcal I_\alpha}\mathcal U_m$, which is dense in $\mathcal Y_\alpha$, lies in
$\sum_{m\in\mathcal I_\alpha,\,m\le i}\mathcal U_m$ for some $i\in\mathcal I_\alpha$, and a power of $\vartheta$ times $f$
lies in $\mathcal F_i$; so $[f]_\alpha\in\mathcal C$.
\end{proof}

\begin{lemma}\label{rig:lem:readouts}
Let $i\in\mathcal I_\alpha$, and let $\iota_i\colon Y_i\to L^{r_i}((0,1);Y_i)$, $y\mapsto1\otimes y$, be the constant
embedding, which is isometric. There are bounded linear maps $J_i\colon Y_i\to\fX_\alpha$ and
$\Xi_i\colon\fX_\alpha\to L^{r_i}((0,1);Y_i)$ with $\norm{J_i}\le1$, $\norm{\Xi_i}\le4/3$ and $\Xi_iJ_i=\iota_i$.
\end{lemma}

\begin{proof}
Put $J_iy=[1\otimes y]_\alpha$, where $1\otimes y\in\mathcal U_i\subseteq\mathcal Y_\alpha$ has norm $\norm y$. The
pointwise coordinate projection $\operatorname{pr}_{Y_i}\colon L_i\to L^{r_i}((0,1);Y_i)$ of $K^{1+N_i}$ onto $Y_i$ along
$Ke_1$ is contractive and kills $\mathbf e_i$. So the map $z\mapsto\operatorname{pr}_{Y_i}(\Pi_i(\id-\Delta_\alpha)^{-1}z)$
on $\mathcal Z$, of norm at most $4/3$, vanishes on $\mathcal Q_\alpha$, since
$(\id-\Delta_\alpha)^{-1}\mathcal Q_\alpha=\mathcal E_\alpha$ and $\Pi_i\mathcal E_\alpha\subseteq K\mathbf e_i$; let
$\Xi_i$ be the restriction to $\fX_\alpha$ of the induced map on $\mathcal Z/\mathcal Q_\alpha$. For $y\in Y_i$,
$(\id-\Delta_\alpha)^{-1}(1\otimes y)-1\otimes y=\sum_{c\ge1}\Delta_\alpha^c(1\otimes y)$ lies in $\mathcal Z_{<i}$,
since $\Delta_\alpha$ maps $\mathcal Z_{\le i}$ into $\mathcal Z_{<i}$. Hence
$\Xi_iJ_iy=\operatorname{pr}_{Y_i}(1\otimes y)=1\otimes y$.
\end{proof}

\begin{lemma}\label{rig:lem:nonzero}
Each $\fX_\alpha$ is nonzero and infinite-dimensional.
\end{lemma}

\begin{proof}
For $i\in\mathcal I_\alpha$ and $y\in Y_i$, Lemma~\ref{rig:lem:readouts} gives
$\norm y=\norm{\Xi_iJ_iy}\le\frac43\norm{J_iy}$. So $\fX_\alpha$ contains the subspaces $J_i(Y_i)$ of dimension
$N_i\ge i$, for the infinitely many $i\in\mathcal I_\alpha$.
\end{proof}

\begin{proof}[Proof of Theorem~\textup{\ref{bey:thm:rigid}}]
The argument follows~\cite[Theorem~3.6]{KR}. Make the construction with any integers $N_i\ge i$ and any lower bounds
$\underline r_i$; the Banach spaces $\fX_\alpha$ are infinite-dimensional by Lemma~\ref{rig:lem:nonzero}. Let
$u\in\cL(\fX_\alpha,\fX_\beta)$; multiplying $u$ by a power of $\vartheta$, we may assume $\norm u<1$.

\emph{Step 1 (approximation).} Fix $i\in\mathcal I_\alpha$; the set $\mathcal J'_i=\{h\in\mathcal J_i:\gamma_h=\mathbf e_i\}$
is infinite. Let $h\in\mathcal J'_i$; then $h>i\ge0$ and $[\mathbf e_h]_\alpha=\zeta_h[\mathbf e_i]_\alpha$
by~\eqref{rig:eq:quotient}. Since $\norm{u[\mathbf e_{h,l}]_\alpha}\le\norm u\norm{\mathbf e_{h,l}}<\norm{\mathbf e_{h,l}}$, there are
$g_{h,l}\in\mathcal Y_\beta$, $1\le l\le d_h$, with $[g_{h,l}]_\beta=u[\mathbf e_{h,l}]_\alpha$ and
$\norm{g_{h,l}}<\norm{\mathbf e_{h,l}}<\eta_h$ (Lemma~\ref{rig:lem:choices}(2)); put $g_h=\sum_lg_{h,l}$, so that
$[g_h]_\beta=u[\mathbf e_h]_\alpha$. Since $\Pi_m$ maps $\mathcal Y_\beta$ into the closed space $\mathcal U_m$ if
$m\in\mathcal I_\beta$ and to zero otherwise, $\Pi_{<h}g_{h,l}\in\mathcal U_{<h}$ and
$\norm{\Pi_{<h}g_h}\le C_h\max_l\norm{\Pi_{<h}g_{h,l}}<C_h\eta_h$. Lemmas~\ref{rig:lem:nakano}
and~\ref{rig:lem:choices}(3) give $\norm{\Pi_{>h}g_h}\le(\sum_l\norm{\mathbf e_{h,l}}^{r_{h+1}})^{1/r_{h+1}}<\eta_h$. So
$\norm{g_h-\Pi_hg_h}<(C_h+1)\eta_h\le2C_h\eta_h<2\abs{\zeta_h}/h$ by Lemma~\ref{rig:lem:choices}(4), and
\begin{equation}\label{rig:eq:approx}
 \bigl\lVert u[\mathbf e_i]_\alpha-\zeta_h^{-1}[\Pi_hg_h]_\beta\bigr\rVert
 =\abs{\zeta_h}^{-1}\norm{[g_h-\Pi_hg_h]_\beta}<2/h\qquad(h\in\mathcal J'_i).
\end{equation}

\emph{Step 2 (the case $\alpha\ne\beta$).} Then $h\notin\mathcal I_\beta$, so $\Pi_hg_h=0$ and
$\norm{u[\mathbf e_i]_\alpha}<2/h$ for infinitely many $h$. Hence $u[\mathbf e_i]_\alpha=0$ for every
$i\in\mathcal I_\alpha$, and $u=0$ by Lemma~\ref{rig:lem:dense}.

\emph{Step 3 (the case $\alpha=\beta$: eigenvectors).} Let $h'<h$ be elements of $\mathcal J'_i$. The vector
$f=\zeta_h^{-1}\Pi_hg_h-\zeta_{h'}^{-1}\Pi_{h'}g_{h'}\in\mathcal Z_{\le h}$ has $\Pi_hf=\zeta_h^{-1}\Pi_hg_h$ and, by~\eqref{rig:eq:approx}
for $h$ and $h'$, $\norm{[f]_\alpha}<2/h+2/h'$. Lemma~\ref{rig:lem:projection} gives $\xi\in K$ with
$\norm{\zeta_h^{-1}\Pi_hg_h-\xi\mathbf e_h}<\frac43(2/h+2/h')$. Passing to classes and using
$[\mathbf e_h]_\alpha=\zeta_h[\mathbf e_i]_\alpha$ and~\eqref{rig:eq:approx}, we get
$\operatorname{dist}(u[\mathbf e_i]_\alpha,K[\mathbf e_i]_\alpha)<2/h+\frac43(2/h+2/h')$. Letting $h\to\infty$ and then
$h'\to\infty$ gives $u[\mathbf e_i]_\alpha=\lambda_i[\mathbf e_i]_\alpha$ for some $\lambda_i\in K$, since
$K[\mathbf e_i]_\alpha$ is closed (Lemma~\ref{lem:findim}). If $[\mathbf e_i]_\alpha\ne0$, then $\lambda_i$ is unique
and $\abs{\lambda_i}\le\norm u$.

\emph{Step 4 (one eigenvalue).} Let $[\mathbf e_i]_\alpha$ and $[\mathbf e_{i'}]_\alpha$ be linearly independent and
$x=[\mathbf e_i]_\alpha+[\mathbf e_{i'}]_\alpha$. By Lemma~\ref{rig:lem:dense}, nonzero
$y_\ell\in K[\mathbf e_{i_\ell}]_\alpha$, $i_\ell\in\mathcal I_\alpha$, converge to $x$. Since
$uy_\ell=\lambda_{i_\ell}y_\ell$ and $\abs{\lambda_{i_\ell}}\le\norm u$,
$\norm{\lambda_{i_\ell}x-ux}\le\abs{\lambda_{i_\ell}}\norm{x-y_\ell}+\norm u\norm{y_\ell-x}\to0$. So $ux$ lies in the
closed line $Kx$, say $ux=\lambda x$, and comparing coefficients gives $\lambda_i=\lambda=\lambda_{i'}$. Nonzero
dependent $[\mathbf e_i]_\alpha$, $[\mathbf e_{i'}]_\alpha$ have equal eigenvalues directly. So $u$ is a scalar
multiple of the identity on $\bigcup_iK[\mathbf e_i]_\alpha$, hence on $\fX_\alpha$ by Lemma~\ref{rig:lem:dense}.
\end{proof}

\subsection{Forcing through an \texorpdfstring{$L^r$}{L\textasciicircum r} readout}\label{rig:ssec:forcing}

Throughout this subsection condition~\textup{(L)} holds with $K=\kappa((X))$; so $K$ is complete, and separable since
the Laurent polynomials over $\kappa$ are dense. Let $c_0(N)=K^N$, $N\ge1$, carry the maximum norm, with unit vectors
$e_1,\dots,e_N$, and let $\mu_a x=ax$ be the coordinatewise product of $a\in K^N$ and $x\in K^N$; then
$\norm{\mu_a}=\norm a_\infty$. For $r\ge k$ put $G_{N,r}=L^r((0,1);c_0(N))$, with the constant embedding
$\iota_{N,r}\colon c_0(N)\to G_{N,r}$, $x\mapsto1\otimes x$. A \emph{readout lift} for $\iota_{N,r}$ is a $k$\nobreakdash-linear
map $Q\colon(K^N)^k\to\cL(\Lambda^k_\pi c_0(N),\widehat\Lambda^k_\pi G_{N,r})$ with
$Q(a,\dots,a)=\Lambda^k_\pi(\iota_{N,r}\mu_a)$ for $a\in K^N$; its norm is
$\norm Q=\sup\{\norm{Q(a_1,\dots,a_k)}:\norm{a_j}_\infty\le1\}$.

\begin{lemma}[Ramsey~{\cite[Theorem~B]{Ra}}]\label{rig:lem:ramsey}
For integers $g,h,C\ge1$ there is an integer $\operatorname{Ram}(g,h,C)$ such that for every coloring of the
$g$\nobreakdash-element subsets of $\{1,\dots,\operatorname{Ram}(g,h,C)\}$ with $C$ colors there is an
$h$\nobreakdash-element set all of whose $g$\nobreakdash-element subsets have the same color.
\end{lemma}

\begin{lemma}\label{rig:lem:wedge-readout}
Let $Y$ be a finite-dimensional normed space and $r\ge k$. The map sending $(g_1,\dots,g_k)$ to the function
$t\mapsto g_1(t)\wedge\dots\wedge g_k(t)$ is a $k$\nobreakdash-linear alternating map from $L^r((0,1);Y)^k$ to
$L^{r/k}((0,1);\Lambda^k_\pi Y)$ of norm at most $1$. It induces a contraction
\[
 \operatorname{pw}\colon\widehat\Lambda^k_\pi L^r((0,1);Y)\to L^{r/k}((0,1);\Lambda^k_\pi Y)
\]
with $\operatorname{pw}(\Lambda^k_\pi\iota(\xi))=1\otimes\xi$ for $\xi\in\Lambda^k_\pi Y$, where $\iota$ is the constant
embedding of $Y$ into $L^r((0,1);Y)$.
\end{lemma}

\begin{proof}
For step functions on a common partition the value is a step function with
$\norm{g_1(t)\wedge\dots\wedge g_k(t)}\le\prod_j\norm{g_j(t)}$, so its $L^{r/k}$\nobreakdash-norm is at most
$\prod_j\norm{g_j}_r$, by H\"older's inequality with $k$ exponents equal to $k$; it vanishes if two of the $g_j$
coincide. Since step functions are dense and $r/k\ge1$, the map extends to a $k$\nobreakdash-linear alternating map of
norm at most $1$ into the Banach space $L^{r/k}((0,1);\Lambda^k_\pi Y)$, and Proposition~\ref{prop:exterior}, in its
form for complete targets, gives $\operatorname{pw}$. On constant functions $g_j=1\otimes x_j$ the value is
$1\otimes(x_1\wedge\dots\wedge x_k)$.
\end{proof}

\begin{lemma}[test lemma]\label{rig:lem:test}
For every $m\ge1$ there are an integer $N(m)\ge m$ and a real number $r^*(m)\ge k$, depending only on $m$, $k$ and
$\abs\kappa$, such that for all $N\ge N(m)$ and $r\ge r^*(m)$ every readout lift $Q$ for $\iota_{N,r}$ satisfies
$\norm Q\ge m/(2M_\varrho\,k(k-1))$.
\end{lemma}

\begin{proof}
\emph{Step 0 (parameters and reduction).} Put $h=4(2m+2+k)+3k$ and
$N(m)=\operatorname{Ram}(3k,h,\abs\kappa^{(3k)^{3k}})\ge h\ge m$, and choose $r^*(m)\ge k$ with
$\binom{N(m)}h\,k!\,(k-1)\,2^{-r^*(m)/k}<1$. Let $N\ge N(m)$, $r\ge r^*(m)$, and $Q$ a readout lift for $\iota_{N,r}$.
The extension by zero $\operatorname{ext}\colon c_0(N(m))\to c_0(N)$ and the pointwise restriction
$\operatorname{res}\colon G_{N,r}\to G_{N(m),r}$ to the first $N(m)$ coordinates are contractive, and $\operatorname{res}\,\iota_{N,r}\,\mu_{\operatorname{ext}a}\operatorname{ext}=\iota_{N(m),r}\,\mu_a$. So
$(a_1,\dots,a_k)\mapsto\widehat\Lambda^k_\pi\operatorname{res}\circ Q(\operatorname{ext}a_1,\dots,\operatorname{ext}a_k)
\circ\Lambda^k_\pi\operatorname{ext}$ is a readout lift for $\iota_{N(m),r}$ of norm at most $\norm Q$, and we may
assume $N=N(m)$.

\emph{Step 1 (readout).} For $a\in(K^N)^k$ and $\xi\in\Lambda^k_\pi c_0(N)$, Lemma~\ref{rig:lem:wedge-readout} gives an element
$\widetilde Q(a;\xi)=\operatorname{pw}(Q(a)(\xi))$ of $L^{r/k}((0,1);\Lambda^k_\pi c_0(N))$, $k$\nobreakdash-linear in $a$,
linear in $\xi$, with $\norm{\widetilde Q(a;\xi)}_{r/k}\le\norm Q\prod_j\norm{a_j}_\infty\norm\xi$. The
$k$\nobreakdash-linear maps $a\mapsto Q(a)(x_1\wedge\dots\wedge x_k)$ and
$a\mapsto\iota_{N,r}(a_1x_1)\wedge\dots\wedge\iota_{N,r}(a_kx_k)$ agree on the diagonal, so the inclusion--exclusion
computation in the proof of Proposition~\ref{prop:diag}(1), followed by $\operatorname{pw}$, gives for all
$a_1,\dots,a_k,x_1,\dots,x_k\in K^N$
\begin{equation}\label{rig:eq:pol}
 \sum_{\sigma\in S_k}\widetilde Q(a_{\sigma(1)},\dots,a_{\sigma(k)};x_1\wedge\dots\wedge x_k)
 =1\otimes\sum_{\sigma\in S_k}(a_{\sigma(1)}x_1)\wedge\dots\wedge(a_{\sigma(k)}x_k).
\end{equation}

\emph{Step 2 (representatives).} For $i\in\{1,\dots,N\}^k$ write $e_i=(e_{i_1},\dots,e_{i_k})$. For each $i$ and each
strictly increasing $y\in\{1,\dots,N\}^k$ fix an everywhere defined representative
$f_{i,y}\colon(0,1)\to\Lambda^k_\pi c_0(N)$ of $\widetilde Q(e_i;e_{y_1}\wedge\dots\wedge e_{y_k})$
(Subsection~\ref{rig:ssec:nakano}); put $f_{i,y}=\sgn(\sigma)f_{i,y\circ\sigma}$ if $y$ has distinct entries and
$y\circ\sigma$ is increasing, and $f_{i,y}=0$ if $y$ has a repeated entry. Then $f_{i,y}$ represents
$\widetilde Q(e_i;e_{y_1}\wedge\dots\wedge e_{y_k})$ for every $y$. For $t\in(0,1)$ and $a,x\in(\F_p^N)^k$ put
\[
 \Psi_t(a;x)=\sum_{i,y\in\{1,\dots,N\}^k}\ \prod_ja_j(i_j)\prod_jx_j(y_j)\,f_{i,y}(t)\in\Lambda^k_\pi c_0(N).
\]
Then $t\mapsto\Psi_t(a;x)$ represents $\widetilde Q(a;x_1\wedge\dots\wedge x_k)$, and each $\Psi_t$ is
$\F_p$\nobreakdash-multilinear and alternating in $x$: if $x_\gamma=x_{\gamma'}$ with $\gamma\ne\gamma'$, the terms of $y$
and of $y\circ(\gamma\ \gamma')$ have equal weights and $f_{i,y\circ(\gamma\,\gamma')}=-f_{i,y}$, and the terms with
$y_\gamma=y_{\gamma'}$ vanish.

\emph{Step 3 (test inputs and the point $t_0$).} For an $h$\nobreakdash-element set $\mathcal H\subseteq\{1,\dots,N\}$
let $\mathcal H'$ consist of its $h-3k$ smallest elements; clusters of $\mathcal H'$ are sets of four consecutive
elements of $\mathcal H'$, as in Appendix~\ref{app:FF}. Let $(a^{\mathcal H};x^{\mathcal H})$ be the inputs of the proof
of Lemma~\ref{lem:D} for the first $k$ clusters of $\mathcal H'$; by~\eqref{rig:eq:pol} and Step~2, the set
$O_{\mathcal H}$ of the $t$ with
$\sum_\sigma\Psi_t(a^{\mathcal H}_{\sigma(1)},\dots,a^{\mathcal H}_{\sigma(k)};x^{\mathcal H})
\ne\sum_\sigma(a^{\mathcal H}_{\sigma(1)}x^{\mathcal H}_1)\wedge\dots\wedge(a^{\mathcal H}_{\sigma(k)}x^{\mathcal H}_k)$ is
null. For $\tau\in S_k$ and $1\le\ell<k$, let $(a^{\mathcal H,\tau,\ell};x^{\mathcal H,\tau,\ell})$, with entries in
$\{0,\pm1\}$, be the inputs of the proof of Lemma~\ref{lem:S} for $\tau$ and $\ell$, with $m+1$ pairs of clusters, all
clusters taken consecutively from the smallest element of $\mathcal H'$; there are enough, since $4(2m+k)\le h-3k$. For
these inputs $(a;x)$, the set $O_{\mathcal H,\tau,\ell}$ of the $t$ with
$\norm{\Psi_t(a;x)}>2\norm{\widetilde Q(a;x_1\wedge\dots\wedge x_k)}_{r/k}$ has measure at most $2^{-r/k}$, by
Chebyshev's inequality for $t\mapsto\norm{\Psi_t(a;x)}^{r/k}$, and is null if
$\widetilde Q(a;x_1\wedge\dots\wedge x_k)=0$. There are at most $\binom Nh\,k!\,(k-1)$ sets $O_{\mathcal H,\tau,\ell}$, so
the union of all of them and of the $O_{\mathcal H}$ has measure at most $\binom Nh\,k!\,(k-1)\,2^{-r/k}<1$. Fix
$t_0\in(0,1)$ outside this union, and write $\Psi=\Psi_{t_0}$.

\emph{Step 4 (Ramsey at $t_0$).} Let $\Omega$ be the determinant array of~\S\ref{ssec:ext} on $\Lambda^k_\pi c_0(N)$,
with $S=\{1,\dots,N\}$, and put $T(i;y;c)=\coeff_0\Omega_{\Psi(e_i;e_y)}(c)\in\kappa$ for $i,y,c\in\{1,\dots,N\}^k$.
Color each $3k$\nobreakdash-element set $\{d_1<\dots<d_{3k}\}\subseteq\{1,\dots,N\}$ by the family
$(T(d_{\eta(1)},\dots,d_{\eta(3k)}))_\eta$ over all maps $\eta\colon\{1,\dots,3k\}\to\{1,\dots,3k\}$, with at most
$\abs\kappa^{(3k)^{3k}}$ colors. Lemma~\ref{rig:lem:ramsey} gives a homogeneous $h$\nobreakdash-element set
$\mathcal H$. Padding with the $3k$ elements of $\mathcal H\setminus\mathcal H'$, which lie above $\mathcal H'$, as in
the proof of Lemma~\ref{lem:H}, shows that $T(i;y;c)$ depends only on the pattern of $(i,y,c)$ when all its entries lie
in $\mathcal H'$.

\emph{Step 5 (the argument of Appendix~\ref{app:FF} at $t_0$).} For inputs supported in $\mathcal H'$ with entries in
$\{0,\pm1\}$,
$\coeff_0\Omega_{\Psi(a;x)}(c)=\sum_{i,y\in\mathcal H'^k}\prod_ja_j(i_j)\prod_jx_j(y_j)\,T(i;y;c)$, as
in~\eqref{eq:expand}, since $\Psi$ is $\F_p$\nobreakdash-multilinear, $\Omega$ is linear and $\coeff_0$ is
$\F_p$\nobreakdash-linear. Define the cluster values $\chi(\tau)\in\kappa$ for clusters of $\mathcal H'$ as in
Appendix~\ref{app:FF}, with $\coeff_0\Omega_{\Psi(\cdot)}$ in place of $\Omega_{\Psi(\cdot)}$. The proofs of
Lemmas~\ref{lem:C} and~\ref{lem:D} use only this expansion, pattern homogeneity of $T$, that clusters are intervals of
the set on which $T$ is homogeneous, \eqref{eq:W}, alternation of $\Psi$ in its vector slots and of $\Omega_z$ in its
arguments, and~\eqref{eq:Pol1} at $(a^{\mathcal H};x^{\mathcal H})$; the last holds since $t_0\notin O_{\mathcal H}$ and
the determinants on the right of~\eqref{eq:Pol1} lie in $\F_p$. So both lemmas hold, and $\sum_\tau\chi(\tau)=1$ in
$\kappa$. Since $k!=0$ in $\kappa$, $\chi$ is not constant; since exchanges of adjacent ranks generate $S_k$, some $\tau$
and $\ell$ have $\delta=\chi(\tau)-\chi(\tau')\ne0$, where $\tau'=(\ell\ \ell{+}1)\circ\tau$.

Let $(a;x)=(a^{\mathcal H,\tau,\ell};x^{\mathcal H,\tau,\ell})$, let $\gamma_1=\tau^{-1}(\ell)$ and
$\gamma_2=\tau^{-1}(\ell+1)$ be the exchanged slots, let $I_1$ and $I_2$ be the sets of first elements of the $m+1$
clusters given to $\gamma_1$ and $\gamma_2$ in the proof of Lemma~\ref{lem:S}, and let $c_\gamma$ be the first element
of the cluster given to any other slot $\gamma$. By Lemma~\ref{lem:C}, as in that proof, the $(m+1)\times(m+1)$ matrix
$\mathrm{St}$ of the values $\coeff_0\Omega_{\Psi(a;x)}(c)$, with an element of $I_1$ in slot $\gamma_1$, an element of
$I_2$ in slot $\gamma_2$ and $c_\gamma$ in each other slot $\gamma$ of $c$, is
$\chi(\tau')\mathbf 1\mathbf 1^{\mathsf T}+\delta U$ with $U$ upper unitriangular, so $\rank\mathrm{St}\ge m$. The
extension by zero $\operatorname{ext}_\infty\colon c_0(N)\to c_0(\N,K)$ makes $\Lambda^k_\pi\operatorname{ext}_\infty$
a contraction that keeps the determinant array on $\{1,\dots,N\}^k$. Let $\Upsilon$ be the slice of
Lemma~\ref{fib:lem:slice} for $c_0(\N,K)$, the slots $\gamma_1,\gamma_2$ and the entries $c_\gamma$, and obtain $z$
from $\Upsilon\Lambda^k_\pi\operatorname{ext}_\infty(\Psi(a;x))$ by restricting its two tensor factors to the
coordinates in $I_1$ and in $I_2$. Then $\coeff_0z=\mathrm{St}$ and $\norm z\le k(k-1)\norm{\Psi(a;x)}$, so
Lemma~\ref{fib:lem:coeff}(2) gives $m\le\rank\mathrm{St}\le M_\varrho k(k-1)\norm{\Psi(a;x)}$. Since
$t_0\notin O_{\mathcal H,\tau,\ell}$ and the entries of $a$ and $x$ lie in $\{0,\pm1\}$,
$\norm{\Psi(a;x)}\le2\norm{\widetilde Q(a;x_1\wedge\dots\wedge x_k)}_{r/k}\le2\norm Q$, and $m\le2M_\varrho k(k-1)\norm Q$.
\end{proof}

\begin{corollary}\label{rig:cor:readout}
Let $m\ge1$, $N\ge N(m)$ and $r\ge r^*(m)$. Let $E$ be a normed space, and $J\colon c_0(N)\to E$ and
$\Xi\colon E\to G_{N,r}$ bounded linear maps with $\Xi J=\iota_{N,r}$. Then every $s$ as in $(\mathrm S^k_E)$ satisfies
$\norm s\ge m/(2M_\varrho k(k-1)(\norm\Xi\norm J)^k)$, and every $\rho$ as in $(\mathrm R^k_E)$ satisfies
$\norm\rho\ge m/(2M_\varrho k(k-1)k!\,(\norm\Xi\norm J)^k)$.
\end{corollary}

\begin{proof}
Let $v_a$ be pointwise multiplication by $a\in K^N$ on $G_{N,r}$; then $\norm{v_a}\le\norm a_\infty$ and
$v_a\iota_{N,r}=\iota_{N,r}\mu_a$. Given $s$, the $k$\nobreakdash-linear map
$Q(a_1,\dots,a_k)=\widehat\wedge^k_{G_{N,r}}\circ((v_{a_1}\Xi)\otimes\dots\otimes(v_{a_k}\Xi))\circ s\circ\Lambda^k_\pi J$
has norm at most $(\norm\Xi\norm J)^k\norm s$ and, by naturality of $\wedge^k$, diagonal
$\widehat\Lambda^k_\pi(v_a\Xi)\circ\widehat\wedge^k_E\circ s\circ\Lambda^k_\pi J=\Lambda^k_\pi(v_a\Xi J)=\Lambda^k_\pi(\iota_{N,r}\mu_a)$.
Given $\rho$, the $k$\nobreakdash-linear map
$Q(a_1,\dots,a_k)=\widehat\Lambda^k_\pi\Xi\circ\rho\circ((J\mu_{a_1})\otimes\dots\otimes(J\mu_{a_k}))\circ\alt^k_{c_0(N)}$
has norm at most $k!\,(\norm\Xi\norm J)^k\norm\rho$, since $\norm{\alt^k_{c_0(N)}}\le k!$, and, by naturality of
$\alt^k$, diagonal
$\widehat\Lambda^k_\pi\Xi\circ\rho\circ\alt^k_E\circ\Lambda^k_\pi(J\mu_a)=\Lambda^k_\pi(\Xi J\mu_a)=\Lambda^k_\pi(\iota_{N,r}\mu_a)$.
Lemma~\ref{rig:lem:test} applies to both.
\end{proof}

\subsection{Proofs of the goodness theorems}\label{rig:ssec:proofs}

Throughout this subsection $\kappa$, $K$ and $\varrho$ are as in Subsection~\ref{rig:ssec:forcing}. For $k'\ge p$ and
$m\ge1$ let $N_{k'}(m)$ and $r^*_{k'}(m)$ be the numbers of Lemma~\ref{rig:lem:test} in degree $k'$, and make the
construction of Subsections~\ref{rig:ssec:choices} and~\ref{rig:ssec:spaces} with
\begin{equation}\label{rig:eq:choices}
 N_i=\max_{p\le k'\le p+i}N_{k'}(i+1),\qquad\underline r_i=\max_{p\le k'\le p+i}r^*_{k'}(i+1)\qquad(i\ge0);
\end{equation}
then $N_i\ge i+1$, as Subsection~\ref{rig:ssec:operators} requires.

\begin{proof}[Proof of Theorem~\textup{\ref{bey:thm:rigid-good}}]
\emph{Step 1 (no splitting).} Let $\alpha\ge1$, $k\ge p$, and $i\in\mathcal I_\alpha$ with $i\ge k-p$. Identifying $Y_i$
with $c_0(N_i)$ through the basis $e_2,\dots,e_{1+N_i}$ turns the map $\iota_i$ of Lemma~\ref{rig:lem:readouts} into
$\iota_{N_i,r_i}$. By~\eqref{rig:eq:choices} and Lemma~\ref{rig:lem:choices}(1), $N_i\ge N_k(i+1)$ and
$r_i\ge r^*_k(i+1)$. Corollary~\ref{rig:cor:readout}, with $m=i+1$, $J=J_i$ and $\Xi=\Xi_i$, shows that every $s$ as in
$(\mathrm S^k_{\fX_\alpha})$ and every $\rho$ as in $(\mathrm R^k_{\fX_\alpha})$ have norm at least
$(3/4)^k(i+1)/(2M_\varrho k(k-1)k!)$. Since $\mathcal I_\alpha$ is infinite, neither $(\mathrm S^k_{\fX_\alpha})$ nor
$(\mathrm R^k_{\fX_\alpha})$ holds.

\emph{Step 2 (the four items).} Let $\fX=\fX_1$. Item~(1) is Theorem~\ref{bey:thm:rigid}. Since $\id_\fX\ne0$
(Lemma~\ref{rig:lem:nonzero}), Proposition~\ref{bey:prop:one-dim}(1) with $E=D=\fX$ and $u_0=\id_\fX$ makes $(\fX,\fX)$
good in every degree, hence good for Banach values by Lemma~\ref{fib:lem:test}; this is~(2). Item~(3) is Step~1 for
$\alpha=1$. Since $\fX$ is complete, by Proposition~\ref{bey:prop:one-dim}(2) a splitting decomposition of degree $k$
would make $\id_\fX$ $\mathrm S^k$-split or $\mathrm R^k$-split, that is, give $(\mathrm S^k_\fX)$ or
$(\mathrm R^k_\fX)$, contrary to~(3). This proves~(4).
\end{proof}

\begin{proof}[Proof of Theorem~\textup{\ref{bey:thm:rigid-sum}}]
The inclusion $\mathrm{in}_\alpha\colon\fX_\alpha\to E$ of the $\alpha$\nobreakdash-th summand and the coordinate
projection $\mathrm{pr}_\alpha\colon E\to\fX_\alpha$ are contractive, and $x=\sum_\alpha\mathrm{in}_\alpha\mathrm{pr}_\alpha x$
for $x\in E$.

\emph{Step 1 (item~(1)).} Let $u\in\cL(E,E)$. By Theorem~\ref{bey:thm:rigid}, $\mathrm{pr}_\beta u\,\mathrm{in}_\alpha=0$
for $\beta\ne\alpha$ and $\mathrm{pr}_\alpha u\,\mathrm{in}_\alpha=\lambda_\alpha\id_{\fX_\alpha}$ for some
$\lambda_\alpha\in K$, with $\abs{\lambda_\alpha}\le\norm u$ since $\fX_\alpha\ne0$ (Lemma~\ref{rig:lem:nonzero}). Hence
$u\,\mathrm{in}_\alpha=\sum_\beta\mathrm{in}_\beta\mathrm{pr}_\beta u\,\mathrm{in}_\alpha=\lambda_\alpha\mathrm{in}_\alpha$
and $ux=\sum_\alpha\lambda_\alpha\mathrm{in}_\alpha\mathrm{pr}_\alpha x=u_\lambda x$. Conversely $\norm{u_\lambda x}=\sum_\alpha\abs{\lambda_\alpha}\norm{\mathrm{pr}_\alpha x}\le\sup_\alpha\abs{\lambda_\alpha}\norm x$,
with $\norm{u_\lambda x}=\abs{\lambda_\alpha}\norm x$ on the summand $\fX_\alpha\ne0$. So $\lambda\mapsto u_\lambda$ is an
isometric isomorphism onto $\cL(E,E)$ from the space $\linf$ of bounded sequences $(\lambda_\alpha)_{\alpha\ge1}$ in~$K$.

\emph{Step 2 (item~(2)).} Let $k\ge1$. Proposition~\ref{cmb:prop:blocks} applies with $I=\{1,2,\dots\}$ in its natural
order, $E_\alpha=D_\alpha=\fX_\alpha$, $D=E$, inclusions $\mathrm{in}_\alpha$, block projections $\mathrm{pr}_\alpha$,
$J'_\alpha=\mathrm{in}_\alpha$, $\mathcal U=\cL(E,E)$, $u_\alpha=\mathrm{pr}_\alpha u\,\mathrm{in}_\alpha$ and $c=C=1$:
$\norm{u_\alpha}\le\norm u$ and $ux=\sum_\alpha\mathrm{in}_\alpha u_\alpha\mathrm{pr}_\alpha x$ by Step~1, and for
$n\le k$, since $\cL(\fX_\alpha,\fX_\alpha)=K\id_{\fX_\alpha}$ with $\id_{\fX_\alpha}\ne0$,
Proposition~\ref{bey:prop:one-dim}(1) with $u_0=\id_{\fX_\alpha}$ gives the lift
$Q^\alpha_n(\theta_1\id_{\fX_\alpha},\dots,\theta_n\id_{\fX_\alpha})=\theta_1\cdots\theta_n\Lambda^n_\pi\id_{\fX_\alpha}$,
$\theta_j\in K$, which satisfies $Q^\alpha_n(v,\dots,v)=\Lambda^n_\pi v$ and $\norm{Q^\alpha_n}\le1$. So $(E,E)$ is good
for Banach values in degree $k$.

\emph{Step 3 (item~(3)).} Suppose that $(E,E)$ has a generalized decomposition of degree $k\ge p$
(Definition~\ref{cmb:def:parameter-piece}) with pieces $(G_l,G'_l,w_l,w'_l,\varepsilon_l)$, so that
$u=\sum_lw'_l\varepsilon_l(u)w_l$ for $u\in\cL(E,E)$; the maps $w'_l$ take values in the complete space $E$. Fix $\alpha$
and $u$; the compression $\mathrm{pr}_\alpha w'_l\varepsilon_l(u)w_l\mathrm{in}_\alpha$ lies in
$\cL(\fX_\alpha,\fX_\alpha)=K\id_{\fX_\alpha}$. For a piece of source or target type, and for the single piece of a
relative splitting, $w_l$ is $\mathrm S^k$-split or $w'_l$ is $\mathrm R^k$-split, so the compression is
$\mathrm S^k$-split or $\mathrm R^k$-split by the ideal property (Proposition~\ref{ob:prop:split-operators}(2)). If it
were a nonzero multiple of $\id_{\fX_\alpha}$, the ideal property would make $\id_{\fX_\alpha}$ $\mathrm S^k$-split or
$\mathrm R^k$-split, contrary to Step~1 of the proof of Theorem~\ref{bey:thm:rigid-good}; so it is zero.

For a parameter piece, $\varepsilon_l=\varepsilon'_l\circ\psi_l$ with bounded linear $\psi_l\colon\cL(E,E)\to P_l$ and
$\varepsilon'_l\colon P_l\to\cL(G_l,G'_l)$, where $P_l$ is polarizable in degrees $\le k$. The linear maps
$g_{l,\alpha}\colon P_l\to K$ given by
$\mathrm{pr}_\alpha w'_l\varepsilon'_l(y)w_l\mathrm{in}_\alpha=g_{l,\alpha}(y)\id_{\fX_\alpha}$ satisfy
$\abs{g_{l,\alpha}(y)}\le\norm{w'_l}\norm{\varepsilon'_l}\norm{w_l}\norm y$ for every $\alpha$, since $\fX_\alpha\ne0$.
Compressing $u_\lambda=\sum_lw'_l\varepsilon_l(u_\lambda)w_l$ to the summand $\alpha$ gives
$\lambda_\alpha=\sum_lg_{l,\alpha}(\psi_l(u_\lambda))$ for every $\lambda\in\linf$ and every $\alpha$, the sum over the
parameter pieces. Let $P$ be the direct sum of their spaces $P_l$ with the maximum norm. By Step~1 and the uniform
bound, $\mathcal A\colon\linf\to P$, $\mathcal A\lambda=(\psi_l(u_\lambda))_l$, and $\mathcal A'\colon P\to\linf$,
$\mathcal A'(y_l)_l=(\sum_lg_{l,\alpha}(y_l))_{\alpha\ge1}$, are bounded linear, and $\mathcal A'\mathcal A=\id_{\linf}$.

The subspace $c_0\subseteq\linf$ of null sequences, a copy of $c_0(\N,K)$, is complete with distances in
$\varrho^\Z\cup\{0\}$, hence spherically complete by Lemma~\ref{lem:discdist}. The extension argument behind Ingleton's
Theorem~\ref{thm:ingleton}, in the form for spherically complete values given in the proof of
Theorem~\ref{thm:spherical}, extends $\id_{c_0}$ to a contractive linear map $\mathcal B\colon\linf\to c_0$ on the
nonarchimedean space $\linf$. Then $\mathcal B\mathcal A'\circ\mathcal A|_{c_0}=\id_{c_0}$, so $c_0$ is a retract of
$P$, which is polarizable in degrees $\le k$ by Proposition~\ref{cmb:prop:parameter-pieces}(2). By
Proposition~\ref{cmb:prop:parameter-pieces}(3), $c_0$ is polarizable in degree $k$, contrary to
Corollary~\ref{cls:cor:c0-not-polarizable}, since $K$ satisfies~\textup{(L)} and $k\ge p$.
\end{proof}
%% ---- end app-M-rigid.tex
%% ---- begin app-N-obstruction.tex
\section{Properties of the obstruction class}\label{app:obstruction-class}

This appendix contains the proofs to which Sections~\ref{sec:canonical}, \ref{sec:combining} and~\ref{sec:sharpness}
refer. Throughout, $K$ is a nontrivially normed field, hence infinite.

\subsection{Pullbacks and morphisms}\label{ocl:ssec:sequences}

Exact sequences, equivalences, classes and pushouts are as in Section~\ref{ob:ssec:sequences}, morphisms and pullbacks
as in Section~\ref{cmb:ssec:families}. Throughout this subsection $\mathcal Y\colon0\to Y_1\xrightarrow iY_2\xrightarrow
qY_3\to0$ and $\mathcal Y'\colon0\to Y'_1\xrightarrow{i'}Y'_2\xrightarrow{q'}Y'_3\to0$ are exact, with a constant $C$
for $\mathcal Y$ as in Section~\ref{ob:ssec:sequences}.

\begin{lemma}\label{ocl:lem:pullback}
For a bounded linear map $\gamma\colon Y'\to Y_3$ from a Banach space $Y'$, $\gamma^*\mathcal Y$ is exact, its class
depends only on $[\mathcal Y]$, and $\gamma^*[\mathcal Y]=0$ if and only if $\gamma=q\circ L$ for a bounded linear
$L\colon Y'\to Y_2$.
\end{lemma}

\begin{proof}
The space $Y_\gamma$ is the kernel of $(y,y')\mapsto qy-\gamma y'$, hence complete. The first map has image
$\{(y,0):qy=0\}$, the kernel of the second, and $\norm{(iy_1,0)}\ge\norm{y_1}/C$. For $y'\in Y'$ some $y$ with
$qy=\gamma y'$ has $\norm{(y,y')}\le(1+C\norm\gamma)\norm{y'}$. An equivalence $T$ of $\mathcal Y$ with an exact
sequence $\mathcal Y''$ induces the equivalence $(y,y')\mapsto(Ty,y')$ of the pullbacks. The splittings of
$\gamma^*\mathcal Y$, the bounded linear right inverses of $(y,y')\mapsto y'$, are the maps $y'\mapsto(Ly',y')$ with $L$
bounded linear and $qL=\gamma$.
\end{proof}

\begin{lemma}\label{ocl:lem:morphism}
For a morphism $(\alpha_1,\alpha_2,\alpha_3)\colon\mathcal Y\to\mathcal Y'$, $\alpha_{1*}[\mathcal
Y]=\alpha_3^*[\mathcal Y']$ in $\Ext(Y_3,Y'_1)$.
\end{lemma}

\begin{proof}
The map $T[(y'_1,y)]=(i'y'_1+\alpha_2y,qy)$ from the middle space of $\alpha_{1*}\mathcal Y$, the quotient of
$Y'_1\oplus Y_2$ by the vectors $(\alpha_1y_1,-iy_1)$, to that of $\alpha_3^*\mathcal Y'$ is well defined, since
$i'\alpha_1=\alpha_2i$ and $qi=0$, bounded, and takes values in the pullback, since $q'(i'y'_1+\alpha_2y)=\alpha_3qy$. It
sends $[(y'_1,0)]$ to $(i'y'_1,0)$ and has second coordinate $qy$, so it is an equivalence
(Lemmas~\ref{ob:lem:split-criteria}(3) and~\ref{ocl:lem:pullback}).
\end{proof}

\begin{corollary}\label{ocl:cor:functorial}
\begin{enumerate}
\item[(a)] If $(\alpha_1,\alpha_2,\id)\colon\mathcal Y\to\mathcal Y'$ is a morphism, then $[\mathcal
Y']=\alpha_{1*}[\mathcal Y]$. If $\mathcal Y''\colon0\to Y_1\xrightarrow{i''}Y''_2\to Y''_3\to0$ is exact and
$(\id,\alpha_2,\gamma)\colon\mathcal Y''\to\mathcal Y$ is a morphism, then $[\mathcal Y'']=\gamma^*[\mathcal Y]$.
\item[(b)] For bounded linear maps $\alpha,\alpha',\gamma,\gamma'$ of Banach spaces for which the expressions are defined,
$(\alpha'\alpha)_*=\alpha'_*\alpha_*$, $(\gamma\gamma')^*=\gamma'^*\gamma^*$, $\id_*=\id^*=\id$ and
$\alpha_*\gamma^*=\gamma^*\alpha_*$.
\item[(c)] For a bounded linear $\alpha\colon Y_1\to Y'$ into a Banach space, $\alpha_*0=0$, and $(c\alpha)_*[\mathcal
Y]=0$ if and only if $\alpha_*[\mathcal Y]=0$, for $c\in K$, $c\ne0$.
\end{enumerate}
\end{corollary}

\begin{proof}
(a) Lemma~\ref{ocl:lem:morphism} gives $\alpha_{1*}[\mathcal Y]=\id^*[\mathcal Y']$ and $\id_*[\mathcal
Y'']=\gamma^*[\mathcal Y]$, and $(y,y_3)\mapsto y$ and $[(y_1,y)]\mapsto i''y_1+y$ are equivalences of $\id^*\mathcal
Y'$ with $\mathcal Y'$ and of $\id_*\mathcal Y''$ with $\mathcal Y''$. (b) The triples
$(\alpha,y\mapsto[(0,y)],\id)\colon\mathcal Y\to\alpha_*\mathcal Y$ and $(\id,(y,y')\mapsto
y,\gamma)\colon\gamma^*\mathcal Y\to\mathcal Y$ are morphisms, since $[(0,iy_1)]=[(\alpha y_1,0)]$. Composites of such
morphisms, the identity morphism and~(a) give the first three identities; the composite $\gamma^*\mathcal Y\to\mathcal
Y\to\alpha_*\mathcal Y$ and Lemma~\ref{ocl:lem:morphism} give the last. (c) By Lemma~\ref{ob:lem:split-criteria}(3),
$\alpha_*[\mathcal Y]=0$ if and only if $\alpha=\tilde\alpha i$ for a bounded linear $\tilde\alpha$, that is,
$c\alpha=(c\tilde\alpha)i$. If $[\mathcal Y]=0$, then $i$ has a bounded linear left inverse $\lambda$
(Lemma~\ref{ob:lem:split-criteria}(2)), and $\alpha=(\alpha\lambda)i$ gives $\alpha_*0=0$.
\end{proof}

\subsection{Permutation groups}\label{ocl:ssec:permutations}

The notation is that of Section~\ref{sec:canonical}, before Theorem~\ref{can:thm:permutation} and in
Proposition~\ref{can:prop:tate}.

\begin{proof}[Proof of Theorem~\ref{can:thm:permutation}]
\emph{Step 1 (notation).} For a \emph{tuple} $t\colon[k]\to I$ put $e_t=e_{t(1)}\otimes\dots\otimes e_{t(k)}$; these
form a basis of $\otimes^k_KV$, and $h^{\otimes k}e_t=e_{h\circ t}$. The \emph{type} of $t$ is the $k$-multiset
$i\mapsto\abs{t^{-1}(i)}$; if $t$ has type $\mathbf m$, then $h\circ t$ has type $h\mathbf m$. For a $k$-subset $T$ let
$\operatorname{Bij}(T)$ be the set of bijections $[k]\to T$. Then $\wedge^ke_f=\sgn(f_T^{-1}f)\eta_T$ for
$f\in\operatorname{Bij}(T)$, $\wedge^ke_t=0$ for $t$ not injective, the $\eta_T$ form a basis of $\Lambda^k_KV$, and
$\chi_T(h)=\sgn(f_T^{-1}hf_T)$ depends only on the permutation that $h$ induces on $T$, so $\chi_T$ is also a
homomorphism on $\bar{\mathsf H}_T$. Since $f_T^{-1}hf=(f_T^{-1}hf_T)(f_T^{-1}f)$ in $S_k$,
\begin{equation}\label{ocl:eq:sign}
 \sgn(f_T^{-1}hf)=\chi_T(h)\sgn(f_T^{-1}f)\qquad(h\in\mathsf H_T,\ f\in\operatorname{Bij}(T)).
\end{equation}
The groups $\bar{\mathsf H}_T$ and $\bar{\mathsf H}_{\mathbf m}$ act freely, by $f\mapsto h\circ f$, on
$\operatorname{Bij}(T)$ and on the tuples of type $\mathbf m$. This gives the \emph{orbit count}: if $c_f\in\Kh$ satisfy
$c_{h\circ f}=\chi_T(h)c_f$ for $h\in\mathsf H_T$, then by~\eqref{ocl:eq:sign} each orbit contributes $\abs{\bar{\mathsf
H}_T}c_f\sgn(f_T^{-1}f)$ to $\sum_{f\in\operatorname{Bij}(T)}c_f\sgn(f_T^{-1}f)$, which therefore lies in
$\abs{\bar{\mathsf H}_T}\Kh$; if $c_t\in\Kh$ satisfy $c_{h\circ t}=c_t$ for the tuples $t$ of type $\mathbf m$ and
$h\in\mathsf H_{\mathbf m}$, then $\sum_tc_t\in\abs{\bar{\mathsf H}_{\mathbf m}}\Kh$. Let $O_{\mathbf m}$ be the sum of
the $e_t$ over the tuples of type $\mathbf m$, as in Proposition~\ref{spl:prop:ordering}. Then $h^{\otimes k}O_{\mathbf
m}=O_{h\mathbf m}$, and the $O_{\mathbf m}$ form a basis of $\Gamma^k_KV$, the space of tensors fixed by every
$\tau_\sigma$, since $\tau_\sigma e_t=e_{t\circ\sigma^{-1}}$.

\emph{Step 2 (necessity).} Assume the hypotheses of~(4). For a tuple $t$ let $e^*_t$ be the bounded functional on
$\widehat\otimes^k_\pi E$ linearizing $(x_j)\mapsto\prod_je^*_{t(j)}(x_j)$; then $e^*_t(e_{t'})=\delta_{tt'}$, and
$e^*_t\circ h^{\otimes k}=e^*_{h^{-1}\circ t}$ by~\eqref{can:eq:compatible}. Let $\ell_T$ be the bounded functional on
$\widehat\Lambda^k_\pi E$ linearizing $(x_{j'})\mapsto\det(e^*_{f_T(j)}(x_{j'}))_{j,j'}$
(Proposition~\ref{prop:exterior}). Then $\ell_T(\eta_T)=1$; $\ell_T\circ\widehat\Lambda^k_\pi h=\chi_T(h)\ell_T$ for
$h\in\mathsf H_T$, since~\eqref{can:eq:compatible} permutes the rows of the determinant by $f_T^{-1}h^{-1}f_T$; and
$\ell_T\circ\widehat\wedge^k_E=\sum_{f\in\operatorname{Bij}(T)}\sgn(f_T^{-1}f)e^*_f$ by the Leibniz formula.

Let $\abs{\bar{\mathsf H}_T}=0$ in $K$. For $s$ as in~(4), $z=s(\eta_T)$ satisfies $z=\chi_T(h)h^{\otimes k}z$ for
$h\in\mathsf H_T$, so $c_f=e^*_f(z)$ satisfies $c_{h\circ f}=\chi_T(h)c_f$, and
$1=\ell_T(\widehat\wedge^k_Ez)=\sum_fc_f\sgn(f_T^{-1}f)=0$ by the orbit count. For $\rho$ as in~(4),
$a_f=\ell_T(\rho(e_f))$ satisfies $a_{h\circ f}=\chi_T(h)a_f$, and $\alt^k_E\eta_T=\sum_f\sgn(f_T^{-1}f)e_f$ gives
$1=\ell_T(\rho\alt^k_E\eta_T)=\sum_f\sgn(f_T^{-1}f)a_f=0$.

Let $\abs{\bar{\mathsf H}_{\mathbf m}}=0$ in $K$, $R$ be as in~(4), and $t_{\mathbf m}$ a tuple of type $\mathbf m$. As
$e^*_t(x^{\otimes k})=\prod_je^*_{t(j)}(x)$ depends only on the type of $t$, every $e^*_t$ with $t$ of type $\mathbf m$
agrees on $\widehat\Gamma^k_\pi E$ with $\varphi_{\mathbf m}=e^*_{t_{\mathbf m}}$. So for $h\in\mathsf H_{\mathbf m}$,
$\varphi_{\mathbf m}\circ h^{\otimes k}=\varphi_{\mathbf m}$ on $\widehat\Gamma^k_\pi E$, and $c_t=\varphi_{\mathbf
m}(Re_t)$ satisfies $c_{h\circ t}=c_t$. By Step~1, $O_{\mathbf m}$ is a combination of pure powers
of vectors in the span of the $e_i$, so $O_{\mathbf m}\in\widehat\Gamma^k_\pi E$, and $1=\varphi_{\mathbf m}(RO_{\mathbf
m})=\sum_tc_t=0$ by the orbit count. With the algebraic powers of $V$, its coordinate functionals and the determinant
functional on $\Lambda^k_KV$, the same argument gives the necessity in (1)--(3).

\emph{Step 3 (sufficiency in (1)).} Let $\abs{\bar{\mathsf H}_T}\ne0$ in $K$ for all $T$, and let $\mathcal T$ be a
set of representatives of the $\mathsf H$-orbits of $k$-subsets. For $T\in\mathcal T$ put $z_T=\abs{\bar{\mathsf
H}_T}^{-1}\sum_{g\in\bar{\mathsf H}_T}\chi_T(g)e_{g\circ f_T}\in U_T$; then $\wedge^kz_T=\eta_T$ by~\eqref{ocl:eq:sign},
and $h'^{\otimes k}z_T=\chi_T(h')z_T$ for $h'\in\mathsf H_T$ by reindexing. The vectors
$\Lambda^k_Kg\,\eta_T=\pm\eta_{gT}$, one for each $k$-subset $gT$, form a basis of $\Lambda^k_KV$; put
$s(\Lambda^k_Kg\,\eta_T)=g^{\otimes k}z_T$. If $gT=g'T$, then $h'=g^{-1}g'\in\mathsf H_T$, and $\Lambda^k_Kg'\eta_T$ and
$g'^{\otimes k}z_T$ are $\chi_T(h')$ times $\Lambda^k_Kg\,\eta_T$ and $g^{\otimes k}z_T$; so $s$ is well defined. It is
$\mathsf H$-equivariant by construction and a section, since $\wedge^kg^{\otimes k}z_T=\Lambda^k_Kg\,\wedge^kz_T$. For
$\lambda\in(K^\times)^I$ let $d_\lambda e_i=\lambda_ie_i$. Then $d_\lambda^{\otimes k}$ and $\Lambda^k_Kd_\lambda$ act
on $U_{T'}$ and on $K\eta_{T'}$ by $\prod_{i\in T'}\lambda_i$, and $s(K\eta_{T'})\subseteq U_{T'}$, so $s$ commutes with
$d_\lambda$.

\emph{Step 4 (sufficiency in (2) and (3)).} For $T\in\mathcal T$ let $\varphi_T$ be the functional on $U_T$ with
$\varphi_T(e_{g\circ f_T})=\abs{\bar{\mathsf H}_T}^{-1}\chi_T(g)$ for $g\in\bar{\mathsf H}_T$ that vanishes on the other
$e_f$. Then $\varphi_T\circ h'^{\otimes k}=\chi_T(h')\varphi_T$ on $U_T$ for $h'\in\mathsf H_T$, and
$\varphi_T(\alt^k\eta_T)=1$ by~\eqref{ocl:eq:sign}. Put $\rho(e_t)=0$ for $t$ not injective and $\rho(e_{g\circ
f})=\varphi_T(e_f)\Lambda^k_Kg\,\eta_T$ for $g\in\mathsf H$, $T\in\mathcal T$ and $f\in\operatorname{Bij}(T)$; this is
well defined, since $g\circ f=g'\circ f'$ implies $h'=g^{-1}g'\in\mathsf H_T$ and $f=h'\circ f'$. The map $\rho$ is
equivariant, $\rho(\alt^k\Lambda^k_Kg\,\eta_T)=\rho(g^{\otimes k}\alt^k\eta_T)=\Lambda^k_Kg\,\eta_T$, and
$\rho(U_{T'})\subseteq K\eta_{T'}$. For~(3), choose representatives $\mathbf m$ of the $\mathsf H$-orbits of
$k$-multisets and tuples $t_{\mathbf m}$ of type $\mathbf m$. Let $\psi_{\mathbf m}(e_t)$ be $\abs{\bar{\mathsf
H}_{\mathbf m}}^{-1}$ for $t\in\bar{\mathsf H}_{\mathbf m}\circ t_{\mathbf m}$ and $0$ for the other $t$ of type
$\mathbf m$; then $\psi_{\mathbf m}$ is $\mathsf H_{\mathbf m}$-invariant and $\psi_{\mathbf m}(O_{\mathbf m})=1$. Put
$R(e_{g\circ t})=\psi_{\mathbf m}(e_t)O_{g\mathbf m}$ for $g\in\mathsf H$ and $t$ of type $\mathbf m$. As for $\rho$,
$R$ is well defined and equivariant; it maps into $\Gamma^k_KV$, and $R(O_{g\mathbf m})=\psi_{\mathbf m}(O_{\mathbf
m})O_{g\mathbf m}=O_{g\mathbf m}$, so it is a projection onto $\Gamma^k_KV$. It maps $e_t$ into $KO_{\mathbf m'}$,
$\mathbf m'$ the type of $t$, on which $d_\lambda^{\otimes k}$ acts by the same factor as on $e_t$.
\end{proof}

\begin{proof}[Proof of Proposition~\ref{can:prop:tate}]
\emph{Step 1 (augmentation).} Let $\operatorname{aug}_T\colon U_T\to K$ be linear with
$\operatorname{aug}_T(e_f)=\sgn(f_T^{-1}f)$. Then $\wedge^kz=\operatorname{aug}_T(z)\eta_T$ for $z\in U_T$, so
$\cN_T=\ker\operatorname{aug}_T$; and $\operatorname{aug}_T(h\cdot z)=\operatorname{aug}_T(z)$ by~\eqref{ocl:eq:sign}, so
$h\cdot e_{f_T}-e_{f_T}\in\cN_T$. As $h\mapsto h\cdot y-y$ satisfies the cocycle identity for every $y\in U_T$,
$\delta_T$ is a well-defined linear map. Call $z\in U_T$ \emph{invariant} if $h\cdot z=z$ for all $h\in\mathsf H_T$,
that is, $h^{\otimes k}z=\chi_T(h)z$.

\emph{Step 2 (item (1)).} For $a\in K$, $\delta_T(a)=0$ means that $z=ae_{f_T}-n$ is invariant for some $n\in\cN_T$, and
then $\operatorname{aug}_T(z)=a$; conversely, for an invariant $z$ with $\operatorname{aug}_T(z)=a$, the vector
$n=ae_{f_T}-z$ lies in $\cN_T$ and gives $\delta_T(a)=0$. For $a=1$ this is~(1).

\emph{Step 3 (items (2) and (3)).} By Step~2, $\ker\delta_T$ is the image under $\operatorname{aug}_T$ of the invariant
vectors. A vector $\sum_fc_fe_f$ is invariant if and only if $c_{h\circ f}=\chi_T(h)c_f$, so by the orbit count this
image lies in $\abs{\bar{\mathsf H}_T}K$; it contains $\abs{\bar{\mathsf H}_T}$, the image of $\sum_{g\in\bar{\mathsf
H}_T}\chi_T(g)e_{g\circ f_T}$. So $\ker\delta_T=\abs{\bar{\mathsf H}_T}K$, which gives~(2), since $\delta_T$ is linear.
By~(2), $\delta_T(1)=0$ if and only if $\abs{\bar{\mathsf H}_T}\ne0$ in $K$, and Theorem~\ref{can:thm:permutation}(1)
gives~(3).
\end{proof}

\subsection{Operator families}\label{ocl:ssec:families}

The notation is that of Section~\ref{cmb:ssec:families}.

\begin{proof}[Proof of Theorem~\ref{cmb:thm:family}(2)]
The squares commute, since $\ev_\theta(\cN^k_E)\subseteq\Mult^k_0(P;\widehat\Lambda^k_\pi D)$ and
$\overline\ev_\theta\circ\widehat\wedge^k_E=q\circ\ev_\theta$ by the definition of $\overline\ev_\theta$ before
Theorem~\ref{cmb:thm:family}. Lemma~\ref{ocl:lem:morphism} gives the formula.
\end{proof}

\begin{proof}[Proof of Proposition~\ref{cmb:prop:naturality}]
\emph{Step 1 (items (a)--(c)).} On pure tensors, hence everywhere, $\ev_{\theta\circ\psi}=\psi^\sharp\circ\ev_\theta$,
$\ev_{\theta_w}=\ev_\theta\circ w^{\otimes k}$ and $\ev_{w'\theta}=(\widehat\Lambda^k_\pi w')_\flat\circ\ev_\theta$.
Corollary~\ref{ocl:cor:functorial}(b) gives~(a) and~(c). For~(b), $(w^{\otimes k}|_{\cN^k_{E'}},w^{\otimes
k},\widehat\Lambda^k_\pi w)$ is a morphism $\cS^k_{E'}\to\cS^k_E$ (proof of Proposition~\ref{ob:prop:split-operators}),
so Corollary~\ref{ocl:cor:functorial}(b) and Lemma~\ref{ocl:lem:morphism} give
\begin{align*}
 \ob_k(\theta_w)&=(\ev_\theta|_{\cN^k_E})_*(w^{\otimes k}|_{\cN^k_{E'}})_*[\cS^k_{E'}]
 =(\ev_\theta|_{\cN^k_E})_*\operatorname{sp}_k(w)\\
 &=(\ev_\theta|_{\cN^k_E})_*(\widehat\Lambda^k_\pi w)^*[\cS^k_E]=(\widehat\Lambda^k_\pi w)^*\ob_k(\theta).
\end{align*}

\emph{Step 2 (item (d)).} The map $(\alt^k_D)_\flat$ preserves zero diagonals, so it induces a bounded linear
$\overline\alt\colon\Poly^k(P;\widehat\Lambda^k_\pi D)\to\Poly^k(P;\widehat\otimes^k_\pi D)$ and a morphism
$\cM^k(P;\widehat\Lambda^k_\pi D)\to\cM^k(P;\widehat\otimes^k_\pi D)$. By Theorem~\ref{cmb:thm:family}(2),
Lemma~\ref{ocl:lem:morphism} and Corollary~\ref{ocl:cor:functorial}(b), the class
$((\alt^k_D)_\flat)_*\ob_k(\theta)$ equals $(\overline\alt\circ\overline\ev_\theta)^*[\cM^k(P;\widehat\otimes^k_\pi D)]$. The
bounded linear map $L\colon\widehat\Lambda^k_\pi E\to\Mult^k(P;\widehat\otimes^k_\pi D)$,
$L(y)(h_1,\dots,h_k)=(\theta(h_1)\otimes\dots\otimes\theta(h_k))\alt^k_Ey$, has the diagonal
$\theta(h)^{\otimes k}\alt^k_Ey=\alt^k_D\widehat\Lambda^k_\pi\theta(h)y$, by naturality of $\alt^k$; so its composite
with the quotient map of $\cM^k(P;\widehat\otimes^k_\pi D)$ is $\overline\alt\circ\overline\ev_\theta$, and
Lemma~\ref{ocl:lem:pullback} gives the first assertion. For the second, extend $w'$ to $\widehat G$, regard $\theta_0$ as
a family $P\to\cL(E,\widehat G)$, and let $\rho\colon\otimes^k_\pi G\to\widehat\Lambda^k_\pi D$ satisfy
$\rho\circ\alt^k_G=\Lambda^k_\pi w'$, so that $\widehat\Lambda^k_\pi w'=\rho\circ\alt^k_G$ on the completions. By~(c),
Corollary~\ref{ocl:cor:functorial}(b),(c) and the first assertion for $\theta_0$,
$\ob_k(w'\theta_0)=(\rho_\flat)_*((\alt^k_G)_\flat)_*\ob_k(\theta_0)=0$, and Theorem~\ref{cmb:thm:family}(1) concludes.
\end{proof}

\begin{lemma}\label{ocl:lem:idempotents}
Let $P_1$ and $P_2$ be normed spaces, $P=P_1\oplus P_2$ with the maximum norm, $Y$ a Banach space, and $\Phi_S$,
$\Pi^{(a)}$ as in Theorem~\ref{cmb:thm:cartan}. For $\Phi\in\Mult^k(P;Y)$, $\Phi=\sum_{S\subseteq[k]}\Phi_S$,
$(\Phi_S)_{S'}$ is $\Phi_S$ if $S'=S$ and $0$ otherwise, and $\norm{\Phi_S}\le\norm\Phi$. Hence the $\Pi^{(a)}$ are
bounded commuting idempotents with sum $\id$. Each maps $\Mult^k_0(P;Y)$ into itself.
\end{lemma}

\begin{proof}
Expanding $h^j=(h^j_1,0)+(0,h^j_2)$ multilinearly gives $\Phi=\sum_S\Phi_S$. The map $(\Phi_S)_{S'}$ evaluates $\Phi$ at
a tuple whose $j$th entry is $0$ for $j$ in exactly one of $S$ and $S'$, and equals $\Phi_S$ if $S'=S$. Since
$(h^j_1,0)$ and $(0,h^j_2)$ have norm at most $\norm{h^j}$, $\norm{\Phi_S}\le\norm\Phi$. For $\Phi\in\Mult^k_0(P;Y)$,
$h=(h_1,h_2)\in P$ and $\lambda\in K$,
\[
 0=\Phi((\lambda h_1,h_2),\dots,(\lambda h_1,h_2))=\sum_S\lambda^{\abs S}\Phi_S(h,\dots,h)
 =\sum_{a=0}^k\lambda^a(\Pi^{(a)}\Phi)(h,\dots,h).
\]
Since $K$ is infinite, the Vandermonde determinant for $k+1$ distinct values of $\lambda$ shows that
$(\Pi^{(a)}\Phi)(h,\dots,h)=0$ for every $a$.
\end{proof}

\begin{proof}[Proof of Theorem~\ref{cmb:thm:cartan}]
Put $\Mult^k_0=\Mult^k_0(P;Y)$, $\Mult^{(a)}=\Pi^{(a)}\Mult^k(P;Y)$ and
$\ev=\ev_{\theta_1\boxplus\theta_2}$. By
Lemma~\ref{ocl:lem:idempotents}, $\Mult^{(a)}$ and $\Mult^{(a)}_0$ are closed and
$\Mult^{(a)}\cap\Mult^k_0=\Mult^{(a)}_0$; by Corollary~\ref{ocl:cor:functorial}(b),
$\ob^{(a)}=(\Pi^{(a)}\circ\ev|_{\cN^k_E})_*[\cS^k_E]$.

\emph{Step 1 (item (1)).} By Lemma~\ref{ob:lem:split-criteria}(3), $\ob_k(\theta_1\boxplus\theta_2)=0$ if and only if
$\ev|_{\cN^k_E}$ extends to a bounded linear $\xi\colon\widehat\otimes^k_\pi E\to\Mult^k_0$, and $\ob^{(a)}=0$ if and
only if $\Pi^{(a)}\ev|_{\cN^k_E}$ extends to a bounded linear $\xi_a\colon\widehat\otimes^k_\pi E\to\Mult^{(a)}_0$. An
extension $\xi$ gives the extensions $\Pi^{(a)}\xi$, and extensions $\xi_a$ give the extension $\sum_a\xi_a$. For
$S\ne[k]$, $\Phi_S\circ\iota_1^{\times k}$ evaluates $\Phi$ at tuples with a zero entry, so
$\iota_1^\sharp=\iota_1^\sharp\Pi^{(k)}$. With the projections $\operatorname{pr}_l\colon P\to P_l$, $l=1,2$,
$\Phi_{[k]}=(\iota_1^\sharp\Phi)\circ\operatorname{pr}_1^{\times k}$, so $\iota_1^\sharp$ is an isometry of $\Mult^{(k)}$
onto $\Mult^k(P_1;Y)$ with inverse $\operatorname{pr}_1^\sharp$, and both preserve zero diagonals. By
Corollary~\ref{ocl:cor:functorial}(b) and Proposition~\ref{cmb:prop:naturality}(a),
$(\iota_1^\sharp)_*\ob^{(k)}=(\iota_1^\sharp)_*\ob_k(\theta_1\boxplus\theta_2)=\ob_k((\theta_1\boxplus\theta_2)\circ\iota_1)=\ob_k(\theta_1)$.
The case of $\iota_2$, with $S=\emptyset$ and $\operatorname{pr}_2$, is the same.

\emph{Step 2 (item (2)).} Let $\Theta_a(h_1,h_2)$ be the right side of the identity in~(2). By
Lemma~\ref{ob:lem:metric-surjection}, the relation $\Pi^{(a)}\ev-\xi_a=\Psi\circ\widehat\wedge^k_E$ matches the
extensions $\xi_a$ of Step~1 with the bounded linear $\Psi\colon\widehat\Lambda^k_\pi E\to\Mult^{(a)}$ such that
$\Psi(\widehat\wedge^k_Ez)-\Pi^{(a)}\ev(z)\in\Mult^k_0$ for all $z$. For $z=x_1\otimes\dots\otimes x_k$ and
$h=(h_1,h_2)$, $\ev(z)_S(h,\dots,h)$ is the wedge, in the order of $j$, of the vectors $\theta_1(h_1)x_j$, $j\in S$, and
$\theta_2(h_2)x_j$, $j\notin S$; by the expansion before~\eqref{cmb:eq:binomial}, the sum over $\abs S=a$ is
$\Theta_a(h_1,h_2)(x_1\wedge\dots\wedge x_k)$. By continuity, the condition on $\Psi$ says
$\Psi(y)(h,\dots,h)=\Theta_a(h_1,h_2)y$. Given $Q^{(a)}$, the map
$\Psi(y)(h^1,\dots,h^k)=Q^{(a)}(h^1_1,\dots,h^a_1;h^{a+1}_2,\dots,h^k_2)y$ satisfies it and lies in $\Mult^{(a)}$, since
$\Psi(y)=\Psi(y)_{[a]}$ (Lemma~\ref{ocl:lem:idempotents}). Conversely, given $\Psi$, for $g\in P_1^a$, $g'\in
P_2^{k-a}$ and $\abs S=a$ let $(g;g')_S\in P^k$ have the entries $(g_1,0),\dots,(g_a,0)$ in the slots of $S$ and
$(0,g'_1),\dots,(0,g'_{k-a})$ in the others, in increasing order, and put $Q^{(a)}(g;g')y=\sum_{\abs
S=a}\Psi(y)((g;g')_S)$; its diagonal is $(\Pi^{(a)}\Psi(y))(h,\dots,h)=\Theta_a(h_1,h_2)y$.

\emph{Step 3 (items (3)--(5)).} (3) For lifts $Q_a$ of $\theta_1$ and $Q'_{k-a}$ of $\theta_2$, the map
$\wedge^{a,k-a}_{\widehat D}\circ(Q_a(g)\otimes Q'_{k-a}(g'))\circ\operatorname{sh}^{a,k-a}_E$, extended to completions,
is bounded and $k$-linear by Lemma~\ref{fib:lem:hopf}, with diagonal $\Theta_a(h_1,h_2)$; apply Step~2. (4) If all
$\ob^{(a)}$ vanish, then $\ob_k(\theta_1\boxplus\theta_2)=0$ by~(1), so
$\ob_k(\theta_1+\theta_2)=(\Delta_{P_1}^\sharp)_*\ob_k(\theta_1\boxplus\theta_2)=0$ by
Proposition~\ref{cmb:prop:naturality}(a) and Corollary~\ref{ocl:cor:functorial}(c), and Theorem~\ref{cmb:thm:family}(1)
applies. If $\theta_1$ and $\theta_2$ are liftable in every degree at most $k$, then
$\ob_k(\theta_1)=\ob_k(\theta_2)=0$ by Theorem~\ref{cmb:thm:family}(1); applying $(\operatorname{pr}_1^\sharp)_*$
and $(\operatorname{pr}_2^\sharp)_*$ to the identities of~(1) gives $\ob^{(k)}=\ob^{(0)}=0$ by
Corollary~\ref{ocl:cor:functorial}(b),(c), and the other $\ob^{(a)}$ vanish by~(3). (5) For $\abs S=a$,
$(\ev_{\theta\boxplus\theta}(z))_S\circ\Delta_{P_1}^{\times k}=\ev_\theta(z)$, on pure tensors since
$(\theta\boxplus\theta)(h,0)=(\theta\boxplus\theta)(0,h)=\theta(h)$, hence everywhere; summing over the $\binom ka$ sets
$S$ gives the identity. By Corollary~\ref{ocl:cor:functorial}(b),
$(\Delta_{P_1}^\sharp|_{\Mult^{(a)}_0})_*\ob^{(a)}=(\binom ka\ev_\theta|_{\cN^k_E})_*[\cS^k_E]$, which by
Corollary~\ref{ocl:cor:functorial}(c) vanishes if and only if $\ob_k(\theta)=0$. If $\theta$ is liftable in degrees $a$
and $k-a$, then $\ob^{(a)}=0$ by~(3), so $\ob_k(\theta)=0$, and Theorem~\ref{cmb:thm:family}(1) applies.
\end{proof}

\begin{proof}[Proof of Proposition~\ref{cmb:prop:composition}]
(1) As $(\theta'\odot\theta)(\psi_hh')=\theta'(h')\circ\theta(h)$, the family $(\theta'\odot\theta)\circ\psi_h$ is the
family $\theta'_w$ of Proposition~\ref{cmb:prop:naturality}(b) for $\theta'$ on $\widehat D$ and
$w=\theta(h)\in\cL(E,\widehat D)$, with $\widehat\Lambda^k_\pi w$ mapping into $\widehat\Lambda^k_\pi\widehat
D=\widehat\Lambda^k_\pi D$; items~(a) and~(b) give the first identity. Likewise $(\theta'\odot\theta)\circ\psi'_{h'}$ is
the family $w'\theta$ of Proposition~\ref{cmb:prop:naturality}(c) for $w'=\theta'(h')\colon\widehat D\to\widehat F$,
and~(a) and~(c) give the second. (2) On pure tensors $(\theta_{D,F}\odot\theta_{E,D})(v\otimes u)=vu$, so
$\theta_{D,F}\odot\theta_{E,D}=\theta_{E,F}\circ\operatorname{comp}$, and Proposition~\ref{cmb:prop:naturality}(a)
gives the identity. If $(E,F)$ is good for Banach values in degree $k$, then $\ob_k(E,F)=0$ by
Theorem~\ref{ob:thm:obstruction}, so the class of $\theta_{D,F}\odot\theta_{E,D}$ vanishes by
Corollary~\ref{ocl:cor:functorial}(c), and Theorem~\ref{cmb:thm:family}(1) applies.
\end{proof}

\subsection{A three-space counterexample}\label{ocl:ssec:three-space}

The Cantor set $\Omega$, its dyadic cells and the spaces $L^1(\Omega;U)$ are as in Theorem~\ref{cmb:thm:cantor}.

\begin{lemma}\label{ocl:lem:no-operators}
Let $K$ be nonarchimedean, $U$ a Banach space, $U_0$ a closed subspace of $L^1(\Omega;U)$, and $Y$ a normed space with
an equivalent nonarchimedean norm. Then every bounded linear map $L^1(\Omega;U)/U_0\to Y$ is zero.
\end{lemma}

\begin{proof}
Let $\psi$ be such a map, $\psi_0$ its composite with the quotient map, and $\nu$ a nonarchimedean norm on $Y$ with
$\nu\le c_\nu\norm\cdot$. For a dyadic cell $C_\xi$ of level $m$, $u\in U$ and $j\ge0$, $1_{C_\xi}u$ is the sum of the
$2^j$ functions $1_{C_{\xi'}}u$, $C_{\xi'}\subseteq C_\xi$ of level $m+j$, each of norm $2^{-m-j}\norm u$. So
$\nu(\psi_0(1_{C_\xi}u))\le\max_{\xi'}\nu(\psi_0(1_{C_{\xi'}}u))\le c_\nu\norm{\psi_0}2^{-m-j}\norm u$, and $j\to\infty$
gives $\psi_0(1_{C_\xi}u)=0$. Step functions are dense, so $\psi_0=0$, and $\psi=0$.
\end{proof}

\begin{lemma}\label{ocl:lem:dilution}
Let $Y$, $Y'$, $Y''$ be Banach spaces, $T\colon Y\to Y'$ bounded linear such that every $y'\in Y'$ equals $Ty$ for some
$y$ with $\norm y\le C_T\norm{y'}$, $R=\ker T$, and $\eta\colon Y\to Y''$ bounded linear and bounded below. Give
$Y'\oplus Y''$ the maximum norm. Then $y\mapsto(-Ty,\eta y)$ and $(y',y'')\mapsto[y''+\eta(y_{y'})]$, where
$Ty_{y'}=y'$, form an exact sequence $0\to Y\to Y'\oplus Y''\to Y''/\eta(R)\to0$.
\end{lemma}

\begin{proof}
The subspace $\eta(R)$ is complete, hence closed. Two choices of $y_{y'}$ differ by an element of $R$, so the second
map is well defined and linear; it is bounded for the choice $\norm{y_{y'}}\le C_T\norm{y'}$, and it maps $(0,y'')$ to
$[y'']$, where $y''$ can be chosen with $\norm{y''}\le2\norm{[y'']}$. If $y''+\eta y_{y'}=\eta r$ with $r\in R$, then
$(y',y'')=(-Tw,\eta w)$ for $w=r-y_{y'}$; conversely $(-Tw,\eta w)$ maps to $[\eta w-\eta w]=0$, with $y_{-Tw}=-w$.
The first map is bounded below, as $\eta$ is.
\end{proof}

\begin{proof}[Proof of Theorem~\ref{cmb:thm:three-space}]
\emph{Step 1 (exactness).} As $\varpi_D$ is open (Corollary~\ref{cmb:cor:no-hereditary}(2)), for some $\delta>0$ the
image of the open unit ball contains the open ball of radius $\delta$. For $d\ne0$ and $j\in\Z$ with
$\abs\vartheta^{j+1}\delta\le\norm d<\abs\vartheta^j\delta$, $\vartheta$ as in Section~\ref{obs:ssec:field}, some $y$ with
$\norm y<1$ has $\varpi_Dy=\vartheta^{-j}d$, and $y_d=\vartheta^jy$ has $\norm{y_d}<\norm d/(\abs\vartheta\delta)$. As
$\eta$ is isometric, Lemma~\ref{ocl:lem:dilution} applies with $T=\varpi_D$.

\emph{Step 2 (items (1) and (2)).} By Proposition~\ref{spl:prop:projective}(3), $\ell^1(\operatorname{Ball}(D),K)$
satisfies $(\mathrm S^n)$ for every $n$, and it is dense in $Y_1$; so Proposition~\ref{cmb:prop:permanence}(3) gives
$(\mathrm S^n_{Y_1})$, and Proposition~\ref{fib:prop:split}(1) gives~(1). Under~(L), $K$ is nonarchimedean, so
Lemma~\ref{ocl:lem:no-operators} with $U=Y_1$, the closed subspace $U_0=\eta(\ker\varpi_D)$ and $Y=D$ gives
$\cL(Y_3,D)=0$; the zero map on $\cL(Y_3,D)^n$ is then a lift as in Lemma~\ref{fib:lem:test} for every $n\ge1$.

\emph{Step 3 (item (3)).} Let $J\colon D\to Y_2$ be $d\mapsto(d,0)$ and $\Pi'\colon Y_2\to D$ the first coordinate. If
$(Y_2,D)$ were good for Banach values in degree $k$, Lemma~\ref{cmb:lem:subpair}(2), with $J$, $\Pi'$ and
$J'=\Pi=\id_D$, would make $(D,D)$ good for Banach values in degree $k$, contrary to Theorem~\ref{fib:thm:local}(2),
since $D$ is complete, has an equivalent nonarchimedean norm and $\dim_{\Kh}D=\infty$.
\end{proof}

\subsection{Characteristic two}\label{ocl:ssec:char-two}

Throughout this subsection the notation is that of Section~\ref{obs:ssec:char-two}, $\Bbbk$ is a field of
characteristic~$2$, and $K$ is a nontrivially normed field of characteristic~$2$. By Lemma~\ref{lem:ultra}, $K$ is
nonarchimedean, and so is its completion $\widehat K$; $K$ is infinite, as a finite field has only the trivial absolute
value. For $z$ in $\Lambda^kK^{n}$ or in $K^{n}\otimes K^{n}$, with coordinates $z_c$,
\begin{equation}\label{ocl:eq:coordinates}
\max_c\abs{z_c}\le\norm z_\pi\le\sum_c\abs{z_c},
\end{equation}
since each coordinate of $x_1\wedge\dots\wedge x_k$ is a $k\times k$ minor, of absolute value at most $\prod_a\norm{x_a}$
by the ultrametric inequality. So the completions of $\Lambda^k_\pi c_0(n)$ and $c_0(n)\otimes_\pi c_0(n)$ are coordinate
spaces over $\widehat K$ with norms equivalent to the coordinate supremum norm, and their elements are identified with
their coordinate vectors.

\begin{definition}\label{ocl:def:polarization}
Let $J\subseteq[n]$, let $x|_J$ be the restriction of $x\in\Bbbk^n$ to $J$, and let $\operatorname{Sym}_J(\Bbbk)$ be the
space of symmetric $J\times J$ matrices. A \emph{$J$-polarization} is a $\Bbbk$-bilinear map $\mathsf
N\colon\Bbbk^n\times\Bbbk^n\to\operatorname{Sym}_J(\Bbbk)$ with $\mathsf N(x,x)=x|_J(x|_J)^{\mathsf T}$. Regarding the
entries of $x$ and $y$ also as indeterminates, $\mathsf N(x,y)=\sum_{t,t'}x_ty_{t'}\mathsf N(e_t,e_{t'})$ is a matrix
over $\Bbbk[x,y]$. For $L,L'\subseteq[n]$ put $x_Ly_{L'}=\prod_{t\in L}x_t\prod_{t'\in L'}y_{t'}$, and for
$g\in\Bbbk[x,y]$ let $c_g(L,L')$ be the coefficient of $x_Ly_{L'}$ in $g$.
\end{definition}

\begin{lemma}\label{ocl:lem:update}
Let $J$ be finite, $\mathsf M$ a symmetric $J\times J$ matrix over a commutative ring of characteristic $2$, $\lambda$ a
ring element, and $\mathsf M[J']$ the principal submatrix on $J'\subseteq J$, with $\det\mathsf M[\emptyset]=1$. For
$a\ne a'$ in $J$,
\begin{gather*}
 \det(\mathsf M+\lambda(e_ae_{a'}^{\mathsf T}+e_{a'}e_a^{\mathsf T}))=\det\mathsf M+\lambda^2\det\mathsf
M[J\setminus\{a,a'\}],\\
 \det(\mathsf M+\lambda e_ae_a^{\mathsf T})=\det\mathsf M+\lambda\det\mathsf M[J\setminus\{a\}].
\end{gather*}
\end{lemma}

\begin{proof}
The second identity is the expansion along row $a$. For the first, expand multilinearly in the rows $a$ and $a'$, which
change by $\lambda e_{a'}^{\mathsf T}$ and $\lambda e_a^{\mathsf T}$. The two first-order terms are $\lambda$ times the
cofactors of $\mathsf M$ at $(a,a')$ and $(a',a)$, which are equal as $\mathsf M$ is symmetric, so they cancel; the
second-order term is $\pm\lambda^2\det\mathsf M[J\setminus\{a,a'\}]$, and signs do not matter in characteristic $2$.
\end{proof}

\begin{lemma}\label{ocl:lem:polarizations}
Let $J\subseteq[n]$, and let $\mathsf N$ and $\mathsf N'$ be $J$-polarizations. Then $c_{\det\mathsf
N}(L,L')=c_{\det\mathsf N}(L',L)$ and $c_{\det\mathsf N}(L,L)=c_{\det\mathsf N'}(L,L)$ for all $L,L'\subseteq[n]$.
\end{lemma}

\begin{proof}
\emph{Step 1 (swap symmetry).} Expanding $\mathsf N(x+y,x+y)$ at unit vectors gives $\mathsf N(y,x)=\mathsf
N(x,y)+x|_J(y|_J)^{\mathsf T}+y|_J(x|_J)^{\mathsf T}$, and $x|_J(y|_J)^{\mathsf T}+y|_J(x|_J)^{\mathsf T}$, whose
diagonal vanishes, is
the sum of the matrices $(x_ay_{a'}+x_{a'}y_a)(e_ae_{a'}^{\mathsf T}+e_{a'}e_a^{\mathsf T})$ over $a<a'$ in $J$. Adding
them one at a time, Lemma~\ref{ocl:lem:update} changes the determinant by terms $(x_a^2y_{a'}^2+x_{a'}^2y_a^2)g$, whose
monomials contain a squared variable. So the coefficient of $x_Ly_{L'}$ in $\det\mathsf N(x,y)$ equals that in
$\det\mathsf N(y,x)$, which is $c_{\det\mathsf N}(L',L)$.

\emph{Step 2 (invariance).} The map $\mathsf A=\mathsf N'-\mathsf N$ is bilinear with symmetric values and $\mathsf
A(x,x)=0$, so $\mathsf A(e_t,e_t)=0$, $\mathsf A(e_{t'},e_t)=\mathsf A(e_t,e_{t'})$ and $\mathsf
A(x,y)=\sum_{t<t'}(x_ty_{t'}+x_{t'}y_t)\mathsf A(e_t,e_{t'})$. Writing each $\mathsf A(e_t,e_{t'})$ as a combination of
the matrices $e_ae_a^{\mathsf T}$ and $e_ae_{a'}^{\mathsf T}+e_{a'}e_a^{\mathsf T}$, we pass from $\mathsf N$ to
$\mathsf N'$ through $J$-polarizations, each obtained from the previous one, $\mathsf N_1$, by adding
$\lambda(x_ty_{t'}+x_{t'}y_t)\mathsf E$ with $\lambda\in\Bbbk$ and $\mathsf E$ one of these matrices. For $\mathsf
E=e_ae_{a'}^{\mathsf T}+e_{a'}e_a^{\mathsf T}$, Lemma~\ref{ocl:lem:update} changes the determinant by
$\lambda^2(x_t^2y_{t'}^2+x_{t'}^2y_t^2)g$, which has no term $x_Ly_L$. For $\mathsf E=e_ae_a^{\mathsf T}$, the change is
$\lambda(x_ty_{t'}+x_{t'}y_t)g$ with $g=\det\mathsf N_1(x,y)[J\setminus\{a\}]$; its coefficient of $x_Ly_L$ is $0$
unless $t,t'\in L$, when it is $\lambda(c_g(L\setminus\{t\},L\setminus\{t'\})+c_g(L\setminus\{t'\},L\setminus\{t\}))=0$
by Step~1 for the $(J\setminus\{a\})$-polarization $\mathsf N_1[J\setminus\{a\}]$.
\end{proof}

\begin{lemma}\label{ocl:lem:reference}
The map with $\mathsf N_<(x,y)_{tt'}=x_{\max(t,t')}y_{\min(t,t')}$ is an $[n]$-polarization,
\[
 \det\mathsf N_<(x,y)=y_1x_n\prod_{t=1}^{n-1}(x_ty_{t+1}+x_{t+1}y_t),
\]
and $c_{\det\mathsf N_<}([n],[n])=1$.
\end{lemma}

\begin{proof}
The map is bilinear with symmetric values, and $\mathsf N_<(x,x)_{tt'}=x_tx_{t'}$. For $n=1$ the determinant is
$x_1y_1$. For $n\ge2$, over the field of rational functions, subtracting $x_n/x_{n-1}$ times row $n-1$ from row $n$
leaves the last row $(0,\dots,0,(x_n/x_{n-1})(x_{n-1}y_n+x_ny_{n-1}))$ below the matrix $\mathsf N_<$ for $n-1$. By
induction $\det\mathsf N_<=x_1y_1\prod_{t=2}^n(x_t/x_{t-1})(x_{t-1}y_t+x_ty_{t-1})$, which telescopes
to the formula. A term of the product equal to $x_{[n]}y_{[n]}$ takes $x_{n-1}y_n$ from the factor of index $n-1$, since
$x_n$ is used, then $x_{n-2}y_{n-1}$ from the factor of index $n-2$, and so on; so the coefficient is $1$.
\end{proof}

\begin{lemma}\label{ocl:lem:parity}
Let $g$ be a polynomial in $x$ and $y$ over a commutative ring of characteristic $2$, of degree at most $n$ in $x$ and
at most $n$ in $y$. Then $\sum_{S,T\subseteq[n]}g(1_S,1_T)$ is the coefficient of $x_{[n]}y_{[n]}$ in $g$. Likewise, if
$g$ is a polynomial in $x$ of degree at most $n$, then $\sum_{S\subseteq[n]}g(1_S)$ is the coefficient of $x_1\cdots
x_n$.
\end{lemma}

\begin{proof}
A monomial $x^\alpha y^\beta$ takes the value $1$ at $(1_S,1_T)$ if $\operatorname{supp}\alpha\subseteq S$ and
$\operatorname{supp}\beta\subseteq T$, and $0$ otherwise, so its sum over $S,T$ is
$2^{2n-\abs{\operatorname{supp}\alpha}-\abs{\operatorname{supp}\beta}}$. This vanishes unless both supports are $[n]$,
which under the degree bounds happens only for $x_{[n]}y_{[n]}$. The second statement is proved in the same way.
\end{proof}

\begin{proposition}\label{ocl:prop:sections}
Let $n\ge2$ and $\sigma\colon\Lambda^2\Bbbk^n\to\Bbbk^{n\times n}$ be $\Bbbk$-linear with $\wedge\circ\sigma=\id$. Then
$\mathsf N(x,y)=\sigma(x\wedge y)+xy^{\mathsf T}$ is an $[n]$-polarization, $\sum_{S,T\subseteq[n]}\det\mathsf
N(1_S,1_T)=1$, and $\rank\sigma(1_S\wedge1_T)\ge n-1$ for some $S,T\subseteq[n]$.
\end{proposition}

\begin{proof}
The map $\Phi(x,y)=\sigma(x\wedge y)$ is bilinear, $\Phi(x,x)=0$, and $\Phi(x,y)+\Phi(x,y)^{\mathsf T}=\Omega_{x\wedge
y}=xy^{\mathsf T}+yx^{\mathsf T}$; so $\mathsf N$ is bilinear, $\mathsf N(x,y)^{\mathsf T}=\mathsf N(x,y)$ and
$\mathsf N(x,x)=xx^{\mathsf T}$. By Lemmas~\ref{ocl:lem:polarizations}
and~\ref{ocl:lem:reference}, $c_{\det\mathsf N}([n],[n])=c_{\det\mathsf N_<}([n],[n])=1$, and $\det\mathsf N(x,y)$ has
degree $n$ in $x$ and in $y$, so Lemma~\ref{ocl:lem:parity} gives the identity. Hence some $\mathsf N(1_S,1_T)$ is
invertible, and $\sigma(1_S\wedge1_T)=\mathsf N(1_S,1_T)+1_S1_T^{\mathsf T}$ has rank at least $n-1$.
\end{proof}

\begin{proposition}\label{ocl:prop:retractions}
Let $n\ge2$ and $\tau\colon\Bbbk^{n\times n}\to\Lambda^2\Bbbk^n$ be $\Bbbk$-linear with $\tau(\Omega_w)=w$ for
$w\in\Lambda^2\Bbbk^n$. If $n$ is even, $\sum_{S\subseteq[n]}\operatorname{Pf}\Omega_{\tau(1_S1_S^{\mathsf T})}=1$, and
some $\Omega_{\tau(1_S1_S^{\mathsf T})}$ is invertible. For every $n$, some $S\subseteq[n]$ has
$\rank\Omega_{\tau(1_S1_S^{\mathsf T})}\ge2\lfloor n/2\rfloor$.
\end{proposition}

\begin{proof}
\emph{Step 1 (even $n$).} Let $n$ be even and $w_t=\tau(e_te_t^{\mathsf T})$. Since $\tau(e_te_{t'}^{\mathsf
T})+\tau(e_{t'}e_t^{\mathsf T})=\tau(\Omega_{e_t\wedge e_{t'}})=e_t\wedge e_{t'}$, $q(x)=\tau(xx^{\mathsf T})$ equals
$\sum_tx_t^2w_t+\sum_{t<t'}x_tx_{t'}e_t\wedge e_{t'}$, and $\Omega_{q(x)}$ has the off-diagonal entries
$x_tx_{t'}+\sum_{t''}x_{t''}^2(\Omega_{w_{t''}})_{tt'}$. In a product over a partition into pairs, a factor $x_{t''}^2$
gives a squared variable, so the coefficient of $x_1\cdots x_n$ in $\operatorname{Pf}\Omega_{q(x)}$, a polynomial of
degree $n$, is the number $(n-1)(n-3)\cdots1$ of these partitions, which is odd. Lemma~\ref{ocl:lem:parity} gives the
identity, so some $\operatorname{Pf}\Omega_{q(1_S)}$ is nonzero, and $\Omega_{q(1_S)}$ is invertible.

\emph{Step 2 (odd $n$).} For odd $n$ let $\iota'\colon\Bbbk^{n-1}\to\Bbbk^n$ be the inclusion of the first $n-1$
coordinates and $\tau'(\mathsf Z)=\Lambda^2\iota'^{\mathsf T}\tau(\iota'\mathsf Z\iota'^{\mathsf T})$. As
$\iota'\Omega_w\iota'^{\mathsf T}=\Omega_{\Lambda^2\iota'w}$ and $\Omega_{\Lambda^2\iota'^{\mathsf T}w}=\iota'^{\mathsf
T}\Omega_w\iota'$, the map $\tau'$ satisfies the hypothesis for $n-1$, and $\rank\Omega_{\tau'(\mathsf
Z)}\le\rank\Omega_{\tau(\iota'\mathsf Z\iota'^{\mathsf T})}$; the even case for $\tau'$ and $\iota'1_S1_S^{\mathsf
T}\iota'^{\mathsf T}=1_S1_S^{\mathsf T}$ give the last assertion.
\end{proof}

\begin{proposition}[Pfaffian identity]\label{ocl:prop:pfaffian-identity}
Let $n\ge2$ be even and $\Psi$ a multiplier lift on $\Bbbk^n$. Then \eqref{obs:eq:pfaffian-identity} holds, and
\begin{equation}\label{ocl:eq:pfaffian-sum}
\sum_{L}\ \sum_{a,x}\ \sum_{b,y}\operatorname{Pf}\mathsf M_\Psi(a,b,x,y)=1,
\end{equation}
the sums over the $L\subseteq[n]$ with $\abs L=n/2$, over the $a,x\in\{0,1\}^n$ supported in $L$, and over the
$b,y\in\{0,1\}^n$ supported in $L^c$.
\end{proposition}

\begin{proof}
\emph{Step 1 (decomposition).} For $t\ne t'$ put $\Psi^{\mathrm{cor}}_{tt'}=\Psi_{tt'}-\Psi^<_{tt'}$. By
\eqref{obs:eq:multiplier-lift}, $\Psi^{\mathrm{cor}}_{tt'}+\Psi^{\mathrm{cor}}_{t't}=\pi_{tt'}-\pi_{tt'}=0$, so
$\Psi^{\mathrm{cor}}_{tt'}=\Psi^{\mathrm{cor}}_{t't}$; write $\Psi^{\mathrm{cor}}_{\{t,t'\}}$ for it. For a $2$-subset
$\varepsilon=\{t,t'\}$ of $[n]$ with $t<t'$ put $e_\varepsilon=e_t\wedge e_{t'}$, and write
$\Psi^{\mathrm{cor}}_{\{t,t'\}}[\varepsilon,\varepsilon']$ for the entry of $\Psi^{\mathrm{cor}}_{\{t,t'\}}$ in the row of
$e_\varepsilon$ and the column of $e_{\varepsilon'}$. Fix $L$ with $\abs L=n/2$. In the monomial $a_Lb_{L^c}x_Ly_{L^c}$
every $a$-variable has its index in $L$ and every $b$-variable in $L^c$, so only the $\Psi_{tt'}$ with $t\in L$ and
$t'\in L^c$, in particular $t\ne t'$, contribute. As the coordinate of $x\wedge y$ at $e_{\{u,u'\}}$ is
$x_uy_{u'}+x_{u'}y_u$, for $t\ne t'$ and $u\ne u'$ the coefficient of $a_tb_{t'}x_uy_{u'}$ in the entry of
$\mathsf M_\Psi$ at a $2$-subset $\varepsilon$ is
\begin{equation}\label{ocl:eq:coefficient}
O_\varepsilon(t,t',u,u')+\Psi^{\mathrm{cor}}_{\{t,t'\}}[\varepsilon,\{u,u'\}],
\end{equation}
where $O_\varepsilon(t,t',u,u')=1$ if $t<t'$ and $\{t,t'\}=\{u,u'\}=\varepsilon$, and $O_\varepsilon(t,t',u,u')=0$
otherwise.

\emph{Step 2 (configurations).} Expanding the Pfaffian over the perfect matchings $\mathfrak m$ of $[n]$, whose elements
we call edges, and each entry by \eqref{ocl:eq:coefficient}, the coefficient of $a_Lb_{L^c}x_Ly_{L^c}$ in
$\operatorname{Pf}\mathsf M_\Psi$ is the sum of the values of the \emph{configurations} for $L$. A configuration
consists of a perfect matching $\mathfrak m$, an assignment $\varepsilon\mapsto(t_\varepsilon,t'_\varepsilon,
u_\varepsilon,u'_\varepsilon)$ on its edges such that each element of $L$ occurs once among the $t_\varepsilon$ and once
among the $u_\varepsilon$, and each element of $L^c$ once among the $t'_\varepsilon$ and once among the
$u'_\varepsilon$, and a choice at each edge of one of the two summands of \eqref{ocl:eq:coefficient}, where the first
summand may be chosen only if it equals $1$. Its value is the product of the chosen summands. An edge is an
\emph{ordering edge} or a \emph{correction edge} according to the chosen summand.

\emph{Step 3 (ordering edges).} An ordering edge $\varepsilon$ has
$\{t_\varepsilon,t'_\varepsilon\}=\{u_\varepsilon,u'_\varepsilon\}=\varepsilon$ and $t_\varepsilon<t'_\varepsilon$. As
$t_\varepsilon,u_\varepsilon\in L$ and $t'_\varepsilon,u'_\varepsilon\notin L$, this forces
$t_\varepsilon=u_\varepsilon=\min\varepsilon\in L$ and $t'_\varepsilon=u'_\varepsilon=\max\varepsilon\notin L$, and its
summand is $1$.

\emph{Step 4 (correction data).} Let $\varepsilon_1,\dots,\varepsilon_m$ be the correction edges of a configuration, and
put $J_i=\{t_{\varepsilon_i},t'_{\varepsilon_i}\}$ and $J'_i=\{u_{\varepsilon_i},u'_{\varepsilon_i}\}$. The value of the
configuration is
\begin{equation}\label{ocl:eq:value}
\prod_{i=1}^m\Psi^{\mathrm{cor}}_{J_i}[\varepsilon_i,J'_i],
\end{equation}
and the set $\Delta=\{(\varepsilon_i,J_i,J'_i):i\le m\}$ is its \emph{correction data}. Each element of $[n]$ occurs
exactly once among the $t_\varepsilon,t'_\varepsilon$, once among the $u_\varepsilon,u'_\varepsilon$ and once among the
edges of $\mathfrak m$, and the ordering edges use the same two elements in all three roles (Step~3). So
$J_1,\dots,J_m$, the sets $J'_1,\dots,J'_m$ and the edges $\varepsilon_1,\dots,\varepsilon_m$ are three partitions of one
set $U\subseteq[n]$, and the ordering edges form a perfect matching $\mathfrak m'$ of $[n]\setminus U$. Moreover
$L\cap U$ meets each $J_i$ and each $J'_i$ in exactly one element.

\emph{Step 5 (counting).} Fix a set $\Delta=\{(\varepsilon_i,J_i,J'_i):i\le m\}$, $m\ge0$, of triples of $2$-subsets in
which the $\varepsilon_i$, the $J_i$ and the $J'_i$ partition the same set $U$, and call a subset of $U$ meeting each
$J_i$ and each $J'_i$ in exactly one element a \emph{transversal} of $\Delta$. The configurations, for all $L$
together, with correction data $\Delta$ correspond bijectively to the pairs of a transversal $L_U$ and a perfect
matching $\mathfrak m'$ of $[n]\setminus U$: given the pair, put $L=L_U\cup\{\min\varepsilon:\varepsilon\in\mathfrak m'\}$
and $\mathfrak m=\{\varepsilon_1,\dots,\varepsilon_m\}\cup\mathfrak m'$, let $t_{\varepsilon_i}$ and $u_{\varepsilon_i}$
be the elements of $J_i$ and $J'_i$ in $L_U$ and $t'_{\varepsilon_i}$, $u'_{\varepsilon_i}$ the other ones, and use
Step~3 on the edges of $\mathfrak m'$; by Steps 3 and 4 every configuration arises from exactly one pair. All of them
have the value \eqref{ocl:eq:value}, and $[n]\setminus U$ has $(n-2m-1)(n-2m-3)\cdots1$ perfect matchings (one if
$U=[n]$), an odd number. The multigraph on $U$ with the edges $J_1,\dots,J_m,J'_1,\dots,J'_m$ has degree $2$ at every
vertex, with one edge of each kind, so it is a disjoint union of $\gamma$ cycles of even length alternating between the
kinds ($J_i=J'_j$ gives a cycle of length $2$). A transversal contains exactly one endpoint of each of these edges, so
on each cycle it is one of the two colour classes of its proper $2$-colouring; there are $2^\gamma$ transversals, and
$\gamma\ge1$ if $m\ge1$.

\emph{Step 6 (summation).} Summing over $L$ sums the values of all configurations. Grouped by their correction data,
those with $m\ge1$ contribute $2^\gamma$ times an odd number times \eqref{ocl:eq:value}, which is $0$, and those with
$m=0$ are the pairs $(\emptyset,\mathfrak m')$, of value $1$, contributing $(n-1)(n-3)\cdots1=1$. This is
\eqref{obs:eq:pfaffian-identity}.

\emph{Step 7 ($0/1$ inputs).} Fix $L$ and restrict $\operatorname{Pf}\mathsf M_\Psi$ to $a,x$ supported in $L$ and $b,y$
supported in $L^c$. The result is a polynomial of degree at most $2n$ in the $2n$ variables $a_t,x_t$ ($t\in L$) and
$b_{t'},y_{t'}$ ($t'\in L^c$). As in the proof of Lemma~\ref{ocl:lem:parity}, the sum of the values of such a polynomial over
$\{0,1\}^{2n}$ is the coefficient of the product of all its variables, here of $a_Lb_{L^c}x_Ly_{L^c}$. Summing over $L$,
\eqref{obs:eq:pfaffian-identity} gives \eqref{ocl:eq:pfaffian-sum}.
\end{proof}

\begin{lemma}[Pfaffian and determinant bounds]\label{ocl:lem:pfaffian-bound}
\leavevmode
\begin{enumerate}
\item Let $n\ge2$ be even. For $z$ in $\Lambda^2_\pi c_0(n)$ or in its completion, $\norm z_\pi\ge\frac n2\,
\abs{\operatorname{Pf}\Omega_z}^{2/n}$.
\item Let $n\ge1$. For $\mathsf Z$ in $c_0(n)\otimes_\pi c_0(n)$ or in its completion, $\norm{\mathsf Z}_\pi\ge n\,
\abs{\det\mathsf Z}^{1/n}$.
\end{enumerate}
\end{lemma}

\begin{proof}
\emph{Step 1 (expansion).} Let $z=\sum_{i\in I_0}x_i\wedge y_i$ with $x_i,y_i\in K^n$ and $I_0$ a finite ordered set.
Then
\begin{equation}\label{ocl:eq:pfaffian-expansion}
\operatorname{Pf}\Omega_z=\sum_{I\subseteq I_0,\ \abs I=n/2}\det(x_{i_1},y_{i_1},\dots,x_{i_{n/2}},y_{i_{n/2}}),
\qquad I=\{i_1<\dots<i_{n/2}\}.
\end{equation}
Indeed, read with the Pfaffian of the alternating matrix of $z$ over $\mathbb Z$, \eqref{ocl:eq:pfaffian-expansion} is
a polynomial identity with integer coefficients in the entries of the $x_i,y_i$. It holds over $\mathbb Q$: there
$z^{\wedge n/2}=(n/2)!\operatorname{Pf}(\Omega_z)\,e_1\wedge\dots\wedge e_n$, the $2$-forms $x_i\wedge y_i$ commute and
have zero square, so $z^{\wedge n/2}=(n/2)!\sum_{\abs I=n/2}\bigwedge_{i\in I}(x_i\wedge y_i)$, and
$\bigwedge_{i\in I}(x_i\wedge y_i)=\det(x_{i_1},y_{i_1},\dots)\,e_1\wedge\dots\wedge e_n$. Hence it holds over
$\mathbb Z$, and after reduction mod $2$ over every commutative ring of characteristic $2$.

\emph{Step 2 (bound).} Put $\beta_i=\norm{x_i}\,\norm{y_i}$. A determinant is a sum of products containing one entry from each
column, so by the ultrametric inequality each determinant in \eqref{ocl:eq:pfaffian-expansion} has absolute value at
most $\prod_{i\in I}\beta_i$. By the inequality of arithmetic and geometric means,
$\prod_{i\in I}\beta_i\le\bigl(\frac2n\sum_{i\in I}\beta_i\bigr)^{n/2}\le\bigl(\frac2n\sum_{i\in I_0}\beta_i\bigr)^{n/2}$.
By the ultrametric inequality again, $\sum_{i\in I_0}\beta_i\ge\frac n2\abs{\operatorname{Pf}\Omega_z}^{2/n}$. Taking the
infimum over the representations of $z$ gives (1) on $\Lambda^2_\pi c_0(n)$. By \eqref{ocl:eq:coordinates} both
sides are continuous for the coordinate topology, in which $\Lambda^2K^n$ is dense in the completion, so (1) holds on
the completion.

\emph{Step 3 (item (2)).} For $\mathsf Z=\sum_{i\in I_0}x_iy_i^{\mathsf T}$, the Cauchy--Binet formula gives
\begin{equation}\label{ocl:eq:cauchy-binet}
\det\mathsf Z=\sum_{I\subseteq I_0,\ \abs I=n}\det(x_i:i\in I)\,\det(y_i:i\in I),
\end{equation}
and Step~2 applies to \eqref{ocl:eq:cauchy-binet} with $n/2$ replaced by $n$; the parity of $n$ plays no role.
\end{proof}

\begin{lemma}[contractions]\label{ocl:lem:contraction}
Let $n'\ge k\ge1$, and let $l\colon K^{n'}\to K$ be linear with $\abs{l(x)}\le\norm x_\infty$ for all $x$. The contraction
$l\mathbin{\lrcorner}\colon\Lambda^kK^{n'}\to\Lambda^{k-1}K^{n'}$,
$l\mathbin{\lrcorner}(x_1\wedge\dots\wedge x_k)=\sum_j(-1)^{j-1}l(x_j)\,x_1\wedge\dots\wedge\widehat{x_j}\wedge\dots\wedge x_k$,
satisfies $\norm{l\mathbin{\lrcorner} z}_\pi\le\norm z_\pi$ for $z\in\Lambda^k_\pi c_0(n')$, and so does its extension to the
completions.
\end{lemma}

\begin{proof}
It suffices to show $\norm{l\mathbin{\lrcorner}\omega}_\pi\le\prod_j\norm{x_j}$ for $\omega=x_1\wedge\dots\wedge x_k$ with all $x_j\ne0$.
If $l(x_j)=0$ for all $j$, then $l\mathbin{\lrcorner}\omega=0$. Otherwise choose $j_0$ with $\abs{l(x_{j_0})}/\norm{x_{j_0}}$ maximal, so
that $l(x_{j_0})\ne0$, and put $x'_j=x_j-\bigl(l(x_j)/l(x_{j_0})\bigr)x_{j_0}$ for $j\ne j_0$. Then $l(x'_j)=0$, and by the
ultrametric inequality and the choice of $j_0$,
$\norm{x'_j}\le\max\bigl(\norm{x_j},\,\abs{l(x_j)}\,\norm{x_{j_0}}/\abs{l(x_{j_0})}\bigr)=\norm{x_j}$. Replacing $x_j$ by $x'_j$ for $j\ne j_0$
does not change $\omega$, so $l\mathbin{\lrcorner}\omega=\pm l(x_{j_0})\bigwedge_{j\ne j_0}x'_j$, whose projective norm is at most
$\abs{l(x_{j_0})}\prod_{j\ne j_0}\norm{x_j}\le\prod_j\norm{x_j}$.
\end{proof}

\begin{proof}[Proof of Theorem~\ref{obs:thm:char-two}]
\emph{Step 1 (item (1)).} The first assertion is Proposition~\ref{ocl:prop:pfaffian-identity}, with $\det=\operatorname{Pf}^2$.
The assertions on $\sigma$ and $\tau$ are Propositions~\ref{ocl:prop:sections} and~\ref{ocl:prop:retractions}. The map $(x,y)\mapsto\cP_<(xy^{\mathsf T}+
yx^{\mathsf T})$ is bilinear and vanishes for $x=y$, so it induces $\sigma^<$. The matrix $\mathsf Z=\cP_<
\Omega_{x\wedge y}$ is strictly upper triangular, of rank at most $n-1$, and $\Omega_{\wedge\mathsf Z}=\mathsf Z+
\mathsf Z^{\mathsf T}=\Omega_{x\wedge y}$, as $\Omega_{x\wedge y}$ has zero diagonal.

\emph{Step 2 (item (2): a multiplier lift on $K^n$).} Let $\iota_n\colon K^n\to E$ be the inclusion of the first $n$
coordinates, for $n'\ge1$ let $\operatorname{res}_{n'}\colon E\to c_0(n')$ be the restriction to the first $n'$
coordinates, and for $n\le n'$ let $\operatorname{res}^{n'}_n\colon c_0(n')\to c_0(n)$ be the restriction. The
restrictions are contractive since $\norm x_\infty\le\norm x_E$, so their exterior powers (Proposition~\ref{prop:exterior})
and tensor powers are contractive; they extend to the completions, and the extensions are written with the same
symbols.
The $0/1$ vectors $1_S$, $S\subseteq[n]$, have $\norm{1_S}_E\le\bar\varphi_E(n)$ and $\norm{1_S}_\infty\le1$.

Let $n'=n+k-2$, let $1_{>n}$ be the indicator of $\{n+1,\dots,n'\}$ and $e_{>n}=e_{n+1}\wedge\dots\wedge e_{n'}$, the
empty product for $k=2$. Let $\eta_{>n}\colon\Lambda^kK^{n'}\to\Lambda^2K^{n'}$ be the composite of the contractions by
the coordinate functionals of $n+1,\dots,n'$, in this order, and the identity for $k=2$; then
$\eta_{>n}(w\wedge e_{>n})=w$ for $w\in\Lambda^2K^n$. Put $\theta(w)=(\Lambda^2_\pi\iota_nw)\wedge
e_{>n}\in\Lambda^k_\pi E$ for $w\in\Lambda^2K^n$. As $\Lambda^k_\pi\operatorname{res}_{n'}(\theta w)=w\wedge e_{>n}$ in
$\Lambda^kK^{n'}$,
\begin{equation}\label{ocl:eq:reduction-identity}
\Lambda^2_\pi\operatorname{res}^{n'}_n\circ\eta_{>n}\circ\Lambda^k_\pi\operatorname{res}_{n'}\circ\theta=\id
\qquad\text{on }\Lambda^2K^n.
\end{equation}
For $a,b\in K^n$, regarded in $c_{00}(\N,K)$, let $Q_2(a,b)$ be the sum, over the pairs of slots $j<j'$, of $Q$
evaluated with $a$ in slot $j$, $b$ in slot $j'$ and $1_{>n}$ in the other slots, and put
\begin{equation}\label{ocl:eq:reduced-lift}
\Psi(a,b)=\Lambda^2_\pi\operatorname{res}^{n'}_n\circ\eta_{>n}\circ\Lambda^k_\pi\operatorname{res}_{n'}\circ
Q_2(a,b)\circ\theta .
\end{equation}
This is a $K$-linear map $\Lambda^2K^n\to\Lambda^2\widehat K^n$, $K$-bilinear in $(a,b)$; extend it $\widehat K$-linearly
in $w$ and $\widehat K$-bilinearly in $(a,b)$. For $\lambda\in K$, multilinearity gives
$Q(a+\lambda1_{>n},\dots,a+\lambda1_{>n})=\sum_{j=0}^k\lambda^{k-j}q_j(a)$, where $q_j(a)$ is the sum, over the
$j$-element sets of slots, of $Q$ with $a$ in these slots and $1_{>n}$ elsewhere; so $q_2(a)=Q_2(a,a)$. As
$\mu_{a+\lambda1_{>n}}$ multiplies $e_t$ by $a_t$ for $t\le n$ and by $\lambda$ for $n<t\le n'$,
$\Lambda^k_\pi\mu_{a+\lambda1_{>n}}(\theta w)=\lambda^{k-2}\theta(\Lambda^2\mu_aw)$. As $K$ is infinite, comparing the
coefficients of $\lambda^{k-2}$ gives $Q_2(a,a)(\theta w)=\theta(\Lambda^2\mu_aw)$, and
\eqref{ocl:eq:reduction-identity} gives $\Psi(a,a)=\Lambda^2\mu_a$ for $a\in K^n$. Taking $a=e_t$, $a=e_{t'}$ and
$a=e_t+e_{t'}$, $t\ne t'$, \eqref{obs:eq:lambda-mu} gives
$\Psi(e_t,e_{t'})+\Psi(e_{t'},e_t)=\Lambda^2\mu_{e_t+e_{t'}}-\Lambda^2\mu_{e_t}-\Lambda^2\mu_{e_{t'}}=\pi_{tt'}$, so
$\Psi$ is a multiplier lift on $\widehat K^n$.

\emph{Step 3 (item (2): the bound).} Let $a,b,x,y\in\{0,1\}^n$. Each of the $k(k-1)/2$ terms of $Q_2(a,b)$ has norm at
most $\norm Q$; $\Lambda^2_\pi\operatorname{res}^{n'}_n$ and $\Lambda^k_\pi\operatorname{res}_{n'}$ are contractive;
$\eta_{>n}$ is contractive by Lemma~\ref{ocl:lem:contraction}; and $\theta(x\wedge y)$ is a single decomposable vector,
so $\norm{\theta(x\wedge y)}_\pi\le\norm x_E\norm y_E\prod_{t=n+1}^{n'}\norm{e_t}_E\le\bar\varphi_E(n)^2$. Hence
\begin{equation}\label{ocl:eq:upper}
\norm{\Psi(a,b)(x\wedge y)}_\pi\le\tfrac{k(k-1)}2\,\norm Q\,\bar\varphi_E(n)^2 ,
\end{equation}
the norm being that of the completion of $\Lambda^2_\pi c_0(n)$. By \eqref{ocl:eq:pfaffian-sum} over $\widehat K$, a
finite sum of the values $\operatorname{Pf}\mathsf M_\Psi(a,b,x,y)$ at $0/1$ inputs is $1$. As $\widehat K$ is
nonarchimedean, one of them has absolute value at least $1$, and
Lemma~\ref{ocl:lem:pfaffian-bound}(1) gives $\norm{\Psi(a,b)(x\wedge y)}_\pi\ge n/2$ for that input. With
\eqref{ocl:eq:upper} this is \eqref{obs:eq:multiplier-bound}.

\emph{Step 4 (item (3): sections).} Let $\sigma=(\operatorname{res}_n\otimes\operatorname{res}_n)\circ s\circ
\Lambda^2_\pi\iota_n$, extended $\widehat K$-linearly to $\Lambda^2\widehat K^n\to\widehat K^{n\times n}$. By naturality
of $\wedge$, $\wedge\circ\sigma=\Lambda^2_\pi\operatorname{res}_n\circ\widehat\wedge{}^2_E\circ s\circ
\Lambda^2_\pi\iota_n=\id$ on $\Lambda^2K^n$, hence on $\Lambda^2\widehat K^n$ by linearity.
Proposition~\ref{ocl:prop:sections} over $\widehat K$ gives $\sum_{S,T\subseteq[n]}\det\mathsf N(1_S,1_T)=1$ for
$\mathsf N(x,y)=\sigma(x\wedge y)+xy^{\mathsf T}$, so some $\abs{\det\mathsf N(1_S,1_T)}\ge1$, and
Lemma~\ref{ocl:lem:pfaffian-bound}(2) gives $n\le\norm{\mathsf N(1_S,1_T)}_\pi\le\norm s\,\bar\varphi_E(n)^2+1$, so
$\norm s\ge(n-1)/\bar\varphi_E(n)^2$. For $c_0(n)$ itself take $\iota_n$ and $\operatorname{res}_n$ to be the identity
here and in Steps 5 and 6.

\emph{Step 5 (item (3): retractions).} Let $n$ be even; for odd $n$ apply the even case to $n-1$, as $\bar\varphi_E$ is
nondecreasing. Let $\tau=\Lambda^2_\pi\operatorname{res}_n\circ\rho\circ(\iota_n\otimes\iota_n)$, extended
$\widehat K$-linearly to $\widehat K^{n\times n}\to\Lambda^2\widehat K^n$. For $w\in\Lambda^2K^n$,
$\alt^2w=\Omega_w$, so $\tau(\Omega_w)=\Lambda^2_\pi\operatorname{res}_n\circ\rho\circ
\alt^2_E(\Lambda^2_\pi\iota_nw)=w$, and by linearity $\tau(\Omega_w)=w$ for $w\in\Lambda^2\widehat K^n$.
Proposition~\ref{ocl:prop:retractions} over $\widehat K$ gives
$\sum_{S\subseteq[n]}\operatorname{Pf}\Omega_{\tau(1_S1_S^{\mathsf T})}=1$, so some
$\abs{\operatorname{Pf}\Omega_{\tau(1_S1_S^{\mathsf T})}}\ge1$, and Lemma~\ref{ocl:lem:pfaffian-bound}(1) gives
$n/2\le\norm{\tau(1_S1_S^{\mathsf T})}_\pi\le\norm\rho\,\bar\varphi_E(n)^2$.

\emph{Step 6 (item (3): projections).} Let $\mathsf N(x,y)=(\operatorname{res}_n\otimes\operatorname{res}_n)
R_2(\iota_nx\otimes\iota_ny)$, extended $\widehat K$-bilinearly. Its values are symmetric matrices, since
$\widehat\Gamma{}^2_\pi E$ is the closed span of the pure squares and $\operatorname{res}_n\otimes\operatorname{res}_n$
maps pure squares to symmetric matrices, which form a closed subspace. As $R_2$ fixes pure squares,
$\mathsf N(e_t,e_t)=e_te_t^{\mathsf T}$ and, polarizing $(e_t+e_{t'})^{\otimes2}$,
$\mathsf N(e_t,e_{t'})+\mathsf N(e_{t'},e_t)=e_te_{t'}^{\mathsf T}+e_{t'}e_t^{\mathsf T}$; so
$\mathsf N(x,x)=xx^{\mathsf T}$ for $x\in\widehat K^n$ by bilinear expansion. So $\mathsf N$ is an $[n]$-polarization
over $\widehat K$ (Definition~\ref{ocl:def:polarization}), and Lemmas~\ref{ocl:lem:polarizations}, \ref{ocl:lem:reference} and~\ref{ocl:lem:parity} give $\sum_{S,T}\det\mathsf N(1_S,1_T)=1$. As in Step~4,
$n\le\norm{\mathsf N(1_S,1_T)}_\pi\le\norm{R_2}\,\bar\varphi_E(n)^2$.

\emph{Step 7 (item (4): goodness).} Each listed $E$ is a sequence space with $\norm{e_t}_E=1$ and
$\norm{\mu_a}_{E\to E}\le\norm a_\infty$ for finitely supported $a$, and $\sup_nn/\bar\varphi_E(n)^2=\infty$. For a
symmetric sequence space, Definition~\ref{obs:def:sequence-spaces} gives a norm that is monotone in $\abs x$ and equal
to $1$ at the unit vectors; so $\bar\varphi_E=\varphi_E$, $\norm x_\infty\le\norm x_E$, and $\abs{ax}\le\abs{\lambda x}$
for a coordinate $\lambda$ of $a$ with $\abs\lambda=\norm a_\infty$ gives $\norm{\mu_a}\le\norm a_\infty$. As
$\bar\varphi_E$ is nondecreasing, $(n-1)/\bar\varphi_E(n-1)^2\ge(n-1)/\bar\varphi_E(n)^2$, so $n/\bar\varphi_E(n)^2$ is
also unbounded along even $n$. If $(E,E)$ were good for Banach values in degree $k$, Lemma~\ref{fib:lem:test} would give
a bounded $k$-linear $\widetilde Q\colon\cL(E,E)^k\to\cL(\Lambda^k_\pi E,\widehat\Lambda^k_\pi E)$ with
$\widetilde Q(u,\dots,u)=\Lambda^k_\pi u$, and $Q(a_1,\dots,a_k)=\widetilde Q(\mu_{a_1},\dots,\mu_{a_k})$ would satisfy
the hypotheses of item (2) with $\norm Q\le\norm{\widetilde Q}$, contradicting \eqref{obs:eq:multiplier-bound} for large
even $n$. Lemma~\ref{fib:lem:test}, Theorem~\ref{thm:duality}(1) and the dichotomy of Proposition~\ref{prop:homogeneous-lift} then give the assertion on $A^k_{E,E;F}$, and
Proposition~\ref{fib:prop:split}(1),(2) the failure of $(\mathrm S^k_E)$ and $(\mathrm R^k_E)$. Item (3) gives the failure of
$(\Gamma^2_E)$.

\emph{Step 8 (item (4): nonarchimedean norms).} Let $h\ge1$, $n=2h$, $n'=n+k-2$, and
$z_h=\sum_{i=1}^he_{2i-1}\wedge e_{2i}\wedge e_{>n}\in\Lambda^k_\pi E$, a sum of $h$ decomposable vectors of norm at
most $1$.
The contractive maps $\Lambda^k_\pi\operatorname{res}_{n'}$, $\eta_{>n}$ and $\Lambda^2_\pi\operatorname{res}^{n'}_n$ send
it to $\sum_{i\le h}e_{2i-1}\wedge e_{2i}$, whose Pfaffian is $1$, so $\norm{z_h}_\pi\ge h$ by
Lemma~\ref{ocl:lem:pfaffian-bound}(1). If $\norm\cdot'$ were a nonarchimedean norm on $\widehat\Lambda^k_\pi E$ with
$c\norm\cdot_\pi\le\norm\cdot'\le C\norm\cdot_\pi$, then $ch\le\norm{z_h}'\le C$ for all $h$, which is impossible.

\emph{Step 9 (item (4): pairs).} Let $2<r_1,r_2\le\infty$, write $\ell^{r_i}$ for $\ell^{r_i}(\N,K)$, or for
$c_0(\N,K)$ if $r_i=\infty$, and put $1/\infty=0$. If $(\ell^{r_1},\ell^{r_2})$ were good for Banach values in degree
$k$, Lemma~\ref{fib:lem:test} would give a bounded $k$-linear
$\widetilde Q\colon\cL(\ell^{r_1},\ell^{r_2})^k\to\cL(\Lambda^k_\pi\ell^{r_1},\widehat\Lambda^k_\pi\ell^{r_2})$ with
$\widetilde Q(u,\dots,u)=\Lambda^k_\pi u$. Run Steps 2 and 3 with
$Q(a_1,\dots,a_k)=\widetilde Q(\mu_{a_1},\dots,\mu_{a_k})$, where $\mu_a$ is now a map $\ell^{r_1}\to\ell^{r_2}$, with
$\theta_1(w)=(\Lambda^2_\pi\iota_nw)\wedge e_{>n}$ in $\Lambda^k_\pi\ell^{r_1}$ in place of $\theta$ on the right of
\eqref{ocl:eq:reduced-lift} and the same expression $\theta_2(w)$ in $\Lambda^k_\pi\ell^{r_2}$ in the identity
$\Lambda^k_\pi\mu_{a+\lambda1_{>n}}\circ\theta_1=\lambda^{k-2}\theta_2\circ\Lambda^2\mu_a$, and with
$\operatorname{res}_{n'}$ defined on $\ell^{r_2}$. Put $\delta=\max(0,1/r_2-1/r_1)$. H\"older's inequality gives
$\norm{\mu_a}\le\abs{\operatorname{supp}a}^\delta$ for $0/1$ vectors $a$, and $\norm{1_S}_{r_1}\le n^{1/r_1}$, so
$n/2\le\frac{k(k-1)}2\norm{\widetilde Q}(k-2)^{(k-2)\delta}n^{2\delta+2/r_1}$, where $(k-2)^{(k-2)\delta}=1$ for $k=2$.
As $2\delta+2/r_1=2/\min(r_1,r_2)<1$, this fails for large even $n$.
\end{proof}

%% ---- end app-N-obstruction.tex
%% ---- begin app-O-families.tex
\section{Proofs for familywise analyticity}\label{app:families}

This appendix proves the results of Section~\ref{sec:families} not proved there, in its notation. We write
$h^{\times n}=(h,\dots,h)\in P^n$, $h_I$ for the tuple of the $h_i$ with $i\in I$, and $\Mult^0(P;Y)=Y$.

\subsection{Power series and mixed exterior powers}\label{fmo:ssec:tools}

\begin{lemma}\label{fmo:lem:series}
\begin{enumerate}
\item If $Y$ is a Banach space and $b_n\in\Mult^n(P;Y)$, $\norm{b_n}\le CR^{-n}$, then $x\mapsto\sum_nb_n(x^{\times n})$ is
analytic on the open ball of radius $R$ about $0$.
\item Let $\Phi\colon Y_1\times\dots\times Y_m\to Y$ be bounded $m$-linear, the spaces normed, and $f_a\colon U\to Y_a$ with
$f_a(x_0+h)=\sum_nb^a_n(h^{\times n})$ for $\norm h<R$, $\norm{b^a_n}\le CR^{-n}$. Then
$\Phi(f_1(x_0+h),\allowbreak\dots,f_m(x_0+h))=
\sum_N\Phi_N(h^{\times N})$ for $\norm h<R$, where
$\Phi_N(h_1,\dots,h_N)=\sum_{n_1+\dots+n_m=N}\Phi(b^1_{n_1}(h_{I_1}),\allowbreak\dots,
b^m_{n_m}(h_{I_m}))$, $I_a$ consecutive of sizes $n_a$, has norm at most $\norm\Phi\binom{N+m-1}{m-1}C^mR^{-N}$. So the
composite is analytic at $x_0$, and on $U$ if the $f_a$ are.
\end{enumerate}
\end{lemma}

\begin{proof}
(1) For $\norm{x_1}<R$, the sum $b'_j$ over $n\ge j$ and $j$-element $S\subseteq\{1,\dots,n\}$ of $b_n$ with
$h_1,\dots,h_j$ in the slots of $S$ and $x_1$ elsewhere converges in $\Mult^j(P;Y)$ with a geometric bound, as
$\sum_{n\ge j}\binom njR^{-n}\norm{x_1}^{n-j}=R(R-\norm{x_1})^{-j-1}$; for $\norm h<R-\norm{x_1}$ the majorant
$\sum_n\norm{b_n}(\norm{x_1}+\norm h)^n$ of the multilinear expansion gives $\sum_nb_n((x_1+h)^{\times n})=
\sum_jb'_j(h^{\times j})$.
(2) The bound counts compositions of $N$. For $\norm h<R$ the terms $\Phi(b^1_{n_1}(h^{\times n_1}),\dots,b^m_{n_m}(h^{\times
n_m}))$ have the majorant $\norm\Phi\prod_a\sum_nC(\norm h/R)^n$; their sums over $\max_an_a\le L$ are values of $\Phi$ at
partial sums of the $f_a$, and differ from their sums over $n_1+\dots+n_m\le L$ by terms with $\max_an_a>L/m$, whose
majorant tends to $0$.
\end{proof}

\begin{lemma}\label{fmo:lem:multishuffle}
Let $G$ be a Banach space, $v_1,\dots,v_m\in\cL(E,G)$, and $\mathbf i\in\N_0^m$ a \emph{type}, $i_1+\dots+i_m=k\ge1$; a map
$\sigma\colon\{1,\dots,k\}\to\{1,\dots,m\}$ has type $\mathbf i$ if $\abs{\sigma^{-1}(l)}=i_l$ for all $l$, and there are
$\binom k{\mathbf i}=k!/(i_1!\cdots i_m!)$ of them. Entries $i_l=0$ use $\Lambda^0_\pi v=\id_K$ and the conventions stated
with~\eqref{cmb:eq:binomial}.
\begin{enumerate}
\item There is a unique $\Lambda^{\mathbf i}_\pi(v_1,\dots,v_m)\in\cL(\Lambda^k_\pi E,\hLambda^k_\pi G)$ sending
$x_1\wedge\dots\wedge x_k$ to $\sum_\sigma v_{\sigma(1)}x_1\wedge\dots\wedge v_{\sigma(k)}x_k$, over the $\sigma$ of type
$\mathbf i$; its norm is at most $\binom k{\mathbf i}\prod_l\norm{v_l}^{i_l}$.
\item The iterated shuffle $\sh^{\mathbf i}_E=(\id\otimes\sh^{(i_2,\dots,i_m)}_E)\circ\sh^{i_1,k-i_1}_E$, $\sh^{(k)}_E=\id$, into
$\Lambda^{i_1}_\pi E\otimes_\pi\dots\otimes_\pi\Lambda^{i_m}_\pi E$ has norm at most $\binom k{\mathbf i}$, the iterated wedge
product $\wedge^{\mathbf i}_G$ from $\hLambda^{i_1}_\pi G\hot_\pi\dots\hot_\pi\hLambda^{i_m}_\pi G$ to $\hLambda^k_\pi G$ is
contractive, and $\Lambda^{\mathbf i}_\pi(v_1,\dots,v_m)=\wedge^{\mathbf i}_G\circ(\Lambda^{i_1}_\pi v_1\otimes\dots\otimes
\Lambda^{i_m}_\pi v_m)\circ\sh^{\mathbf i}_E$.
\end{enumerate}
\end{lemma}

\begin{proof}
Lemma~\ref{fib:lem:hopf}, iterated and extended to completions, and $\norm{\id\otimes T}\le\norm T$ give the norm bounds
of~(2). By induction on $m$ and~\eqref{fib:eq:shuffle}, $\sh^{\mathbf i}_E(x_1\wedge\dots\wedge x_k)$ is the sum of the
$\epsilon(S_1,\dots,S_m)\,x_{S_1}\otimes\dots\otimes x_{S_m}$ over the ordered partitions with $\abs{S_l}=i_l$, where
$\epsilon(S_1,\dots,S_m)$, the sign of the listing of $S_1,\dots,S_m$ in turn, each increasingly, is $\epsilon(S_1)$ times
the sign of the listing of $S_2,\dots,S_m$ relative to the increasing listing of the complement of $S_1$. With
$\sigma(c)=l$ for $c\in S_l$ and $y_c=v_{\sigma(c)}x_c$, the composite in~(2) sends $x_1\wedge\dots\wedge x_k$ to the sum of
the $\epsilon(S_1,\dots,S_m)\,y_{S_1}\wedge\dots\wedge y_{S_m}=y_1\wedge\dots\wedge y_k$, and ordered partitions correspond
bijectively to maps of type $\mathbf i$. So the composite is the operator of~(1), unique as pure wedges span, and
$\norm{\Lambda^{i_l}_\pi v_l}\le\norm{v_l}^{i_l}$ gives its bound.
\end{proof}

\begin{lemma}\label{fmo:lem:vandermonde}\emergencystretch=3em
Let $G$ be a Banach space and $a,b\in\cL(E,G)$. Then $\Lambda^k_\pi(a+\lambda b)=\sum_{j=0}^k\lambda^j\Lambda^{(k-j,j)}_\pi
(a,b)$ for $\lambda\in K$, where $\Lambda^{(k-j,j)}_\pi(a,b)=\wedge^{k-j,j}_G\circ(\Lambda^{k-j}_\pi a\otimes\Lambda^j_\pi b)
\circ\sh^{k-j,j}_E$ has norm at most $\binom kj\norm a^{k-j}\norm b^j$, and $\Lambda^{(0,k)}_\pi(a,b)=\Lambda^k_\pi b$. If
$\lambda_0,\dots,\lambda_k\in K$ are distinct and $(\chi_{ji})$ is the inverse of $(\lambda_i^j)_{0\le i,j\le k}$, then
$\Lambda^k_\pi b=\sum_i\chi_{ki}\Lambda^k_\pi(a+\lambda_ib)$.
\end{lemma}

\begin{proof}
Equation~\eqref{cmb:eq:binomial} with $w_1=a$, $w_2=\lambda b$ and $\Lambda^j_\pi(\lambda b)=\lambda^j\Lambda^j_\pi b$ gives
the expansion, and Lemma~\ref{fib:lem:hopf} the bound. The Vandermonde matrix being invertible, the identities
$\Lambda^k_\pi(a+\lambda_ib)=\sum_j\lambda_i^j\Lambda^{(k-j,j)}_\pi(a,b)$ give $\Lambda^{(k-j,j)}_\pi(a,b)=
\sum_i\chi_{ji}\Lambda^k_\pi(a+\lambda_ib)$; take $j=k$.
\end{proof}

\begin{lemma}\label{fmo:lem:neumann}
Let $H$ be a Banach space, $U'\subseteq P$ open and $w\colon U'\to\cL(H)$ analytic with $\sup_{U'}\norm w<1$. Then $1+w$ is
an invertible family on $U'$.
\end{lemma}

\begin{proof}
Fix $x_1\in U'$ and let $g_1=1+w(x_1)$, so $\norm{g_1^{-1}}\le(1-\sup_{U'}\norm w)^{-1}$ by the Neumann series. Write
$w(x_1+h)-w(x_1)=\delta(h)=\sum_{n\ge1}b_n(h^{\times n})$, $\norm{b_n}\le CR^{-n}$. With $C'=\norm{g_1^{-1}}C$ and
$t=\norm h/R<1/(1+C')$, $\norm{g_1^{-1}\delta(h)}\le C't/(1-t)<1$, so $(1+w(x_1+h))^{-1}=\sum_l(-g_1^{-1}\delta(h))^l
g_1^{-1}$, whose multilinear expansion has the majorant $\norm{g_1^{-1}}\sum_l(C't/(1-t))^l$; its degree-$N$
coefficient, the sum over the compositions $(n_1,\dots,n_l)$ of $N$ of
$(-1)^lg_1^{-1}b_{n_1}(h_{I_1})g_1^{-1}\cdots b_{n_l}(h_{I_l})g_1^{-1}$, has norm at most $\norm{g_1^{-1}}(1+C')^NR^{-N}$.
\end{proof}

\subsection{Reductions and what bundles need}\label{fmo:ssec:reductions}

\begin{proof}[Proof of Proposition~\ref{fam2:prop:reductions}]
\emph{Step 1 (item~(1)).} Let $u\colon U\to\cL(H)$ be analytic and $x_0\in U$. Choose an open $U_0\subseteq U$
containing $x_0$ with $\norm u\le C=\norm{u(x_0)}+1$ on $U_0$, an integer $m\ge0$ with $\abs\vartheta^mC<1$, and the
distinct $\lambda_i=\vartheta^{m+i}$, $0\le i\le k$. By Lemma~\ref{fmo:lem:neumann} every $1+\lambda_iu$ is an invertible
family on $U_0$, and Lemma~\ref{fmo:lem:vandermonde} with $G=H$, $a=1$, $b=u(x)$ gives
$\Lambda^k_\pi u=\sum_i\chi_{ki}\Lambda^k_\pi(1+\lambda_iu)$ on $U_0$.

\emph{Step 2 (item~(2)).} Let $\alpha_1\colon E_0\to E$, $\alpha_2\colon D_0\to D$, $\psi\colon P_0\to P$ have bounded linear
left inverses $\alpha'_1$, $\alpha'_2$, $\psi'$. For analytic $u_0\colon U'\to\cL(E_0,\widehat D_0)$, $U'\subseteq P_0$
open, $u=\hat\alpha_2u_0(\psi'\cdot)\alpha'_1$ is analytic on $\psi'^{-1}(U')$, hats denoting extensions to completions;
hence so is $\Lambda^k_\pi u$, and $\Lambda^k_\pi u_0(y)=\hLambda^k_\pi\hat\alpha'_2\circ\Lambda^k_\pi u(\psi y)\circ
\Lambda^k_\pi\alpha_1$ (Proposition~\ref{prop:exterior}). For $E\oplus E'$, ``only if'' is the retract case; conversely,
for analytic $u\colon U\to\cL(E\oplus E',\widehat D)$, the injections $\iota_l$ and projections $\pi_l$ of $E\oplus E'$,
$u_l=u\iota_l$ and $k'\le k$, $u=u_1\pi_1+u_2\pi_2$, \eqref{cmb:eq:binomial} and Proposition~\ref{prop:exterior} give
\[\Lambda^{k'}_\pi u=\sum_{j=0}^{k'}\wedge^{j,k'-j}_{\widehat D}\circ\bigl((\Lambda^j_\pi u_1\circ\Lambda^j_\pi\pi_1)\otimes
(\Lambda^{k'-j}_\pi u_2\circ\Lambda^{k'-j}_\pi\pi_2)\bigr)\circ\sh^{j,k'-j}_{E\oplus E'},\]
bounded bilinear (linear for $j\in\{0,k'\}$) in analytic maps; apply Lemma~\ref{fmo:lem:series}(2). The target slot is the
same, with $u=\sum_l\iota'_l\pi'_lu$ for the injections $\iota'_l$ and projections $\pi'_l$ of
$\widehat D\oplus\widehat{D'}$ and $\Lambda^j_\pi(\iota'_l\pi'_lu)=\hLambda^j_\pi\iota'_l\circ\Lambda^j_\pi(\pi'_lu)$.

\emph{Step 3 (item~(3)).} Let $g$ be as in Proposition~\ref{prop:shear}, with $H=D\oplus E$, $\pi_D$ the projection onto
$D$ and $i_E$ the injection of $E$. For analytic $u\colon U\to\cL(E,D)$, $x\mapsto g(u(x))$ is an invertible family, with
inverse $g(-u(x))$, and $\pi_Dg(u(x))i_E=u(x)$ and Proposition~\ref{prop:exterior} give
\begin{equation}\label{fmo:eq:shear}
 \Lambda^k_\pi u(x)=\hLambda^k_\pi\pi_D\circ\Lambda^k_\pi g(u(x))\circ\Lambda^k_\pi i_E .
\end{equation}
In degrees $j\le k$, analyticity along invertible families into $\cL(H)$ gives goodness of $(P;H)$ by~(1), hence of the
four triples by~(2); conversely, goodness of the four triples gives that of $(P;H)$ by~(2) in both slots, and invertible
families are analytic. For $\cC$, goodness of $(P;G,G)$ for $G\in\cC$ covers invertible families into $\cL(G)$;
conversely, for $E,D\in\cC$, $D\oplus E\in\cC$ and~\eqref{fmo:eq:shear} applies.
\end{proof}

\begin{proof}[Proof of Corollary~\ref{fam2:cor:bundles}]
\emph{Step 1 ((2) implies the last assertion and (1)).} Let $E,D\in\cC$, let $F,F'$ be Banach spaces, and let
$u\colon U\to\cL(E,D)$ and $v\colon U\to\cL(F,F')$ be analytic on an open $U\subseteq P$. By~(2) and the paragraph after
Definition~\ref{fam2:def:good}, $x\mapsto A^k_{E,D;F}(u(x))$ is analytic. The joint action $B(u(x),v(x))$, which is
$A^k_{E,D;F}(u(x))$ followed by postcomposition with $v(x)$, is bounded bilinear in the two families, hence analytic
(Lemma~\ref{fmo:lem:series}(2)). The last assertion of Theorem~\ref{thm:bundle-lifting}, for restricted categories of
fibers as in the paragraph after it, gives the last assertion; (1) is its case of a trivial bundle with fiber $F$.

\emph{Step 2 ((1) implies (2)).} Let $E,D\in\cC$, let $u\colon U\to\cL(E,D)$ be analytic on an open $U\subseteq P$, and let
$H=D\oplus E\in\cC$ and $g$ be as in Proposition~\ref{prop:shear}. The trivial bundle $U\times H$ with the trivializations
$\id$ and $(x,z)\mapsto(x,g(u(x))z)$ is analytic, since its transitions $g(\pm u(x))$ are affine in $u(x)$. By~(1), the
induced transition $x\mapsto g(u(x))^*$ on $\Alt^k(H;F)$ is analytic for every Banach space $F$. The
compression~\eqref{eq:shear-compression} makes $x\mapsto A^k_{E,D;F}(u(x))$ analytic, so $\Lambda^k_\pi u$ is analytic on
$U$ by Theorem~\ref{thm:duality}(1) with $Y=\hLambda^k_\pi D$; as $D$ is complete, this is~(2).
\end{proof}

\subsection{Homogeneous families}\label{fmo:ssec:reduction}

A \emph{content} is a finitely supported $\mathbf m\colon\N_0\to\N_0$ with $\sum_n\mathbf m(n)=k$; its degree is
$N(\mathbf m)=\sum_nn\mathbf m(n)$, and its type is $(\mathbf m(n_1),\dots,\mathbf m(n_m))$ for its support $n_1<\dots<n_m$.
The content of $(n'_1,\dots,n'_k)\in\N_0^k$ counts repeated values, so at most $\binom{N+k-1}{k-1}$ contents have degree $N$.

\begin{lemma}\label{fmo:lem:homogeneous}
Let $n\ge0$, $j\ge1$ and $\beta\in\Mult^n(P;\cL(E,\widehat D))$.
\begin{enumerate}
\item $c_j(\beta)$ depends only on $u_\beta$, and $c_j(\lambda\beta)=\abs\lambda^jc_j(\beta)$ for $\lambda\in K^\times$.
\item $c_1(\beta)\le\norm\beta$, and $c_j(\beta)\le(j!/\abs{j!})\norm\beta^j$ if $j!\ne0$ in $K$, the integer $j!$ in the
numerator.
\item For $n\ge1$, $h\mapsto\Lambda^j_\pi u_\beta(h)$ is analytic at some point of $P$ if and only if $c_j(\beta)<\infty$,
and it is then a continuous polynomial.
\item $c_j(b\beta a)\le\norm b^j\norm a^jc_j(\beta)$ for bounded linear $a\colon E'\to E$, $b\colon\widehat D\to\widehat{D'}$,
with $(b\beta a)(h_1,\allowbreak\dots,h_n)=b\,\beta(h_1,\allowbreak\dots,h_n)\,a$.
\end{enumerate}
\end{lemma}

\begin{proof}
(1) The definition involves only $\Lambda^j_\pi u_\beta$, and $\lambda^jq$ lifts $u_{\lambda\beta}$ if $q$ lifts $u_\beta$;
apply this also to $\lambda^{-1}$. (2) $\beta$ lifts $u_\beta$ in degree $1$. If $j!\ne0$, the section $s=\alt^j_E/j!$
of $\wedge^j_E$ has norm at most
$j!/\abs{j!}$ (Proposition~\ref{can:prop:natural}(1)), and $q(h_1,\dots,h_{jn})=\widehat\wedge^j_D\circ(\beta(h_{I_1})\otimes
\dots\otimes\beta(h_{I_j}))\circ s$, blocks of size $n$, has diagonal $\Lambda^j_\pi u_\beta(h)\wedge^j_Es=
\Lambda^j_\pi u_\beta(h)$ by naturality of the wedge, and norm at most $(j!/\abs{j!})\norm\beta^j$. (3) $T\mapsto
T\circ\wedge^j_E$ is an isometry of
$\cL(\Lambda^j_\pi E,\hLambda^j_\pi D)$ onto the operators in $\cL(\otimes^j_\pi E,\hLambda^j_\pi D)$ vanishing on
$\ker\wedge^j_E$ (Lemma~\ref{ob:lem:metric-surjection}), a closed subspace, and $(h_1,\dots,h_{jn})\mapsto
\widehat\wedge^j_D\circ(\beta(h_{I_1})\otimes\dots\otimes\beta(h_{I_j}))$ is an ambient representative of degree $jn\ge1$;
apply Proposition~\ref{prop:homogeneous-lift}. (4) If $q$ lifts $u_\beta$, then $\hLambda^j_\pi b\circ q\circ\Lambda^j_\pi a$
lifts $u_{b\beta a}$ (Proposition~\ref{prop:exterior}).
\end{proof}

\begin{lemma}\label{fmo:lem:assembly}
Let $u(x_0+h)=\sum_{n\ge0}\beta_n(h^{\times n})$ for $\norm h<R$, $\beta_n\in\Mult^n(P;\cL(E,\widehat D))$,
$\norm{\beta_n}\le CR^{-n}$, $a=\beta_0$, and let $C_0,C_1\ge1$ and lifts $q_{n,j}$ of $u_{\beta_n}$ in degree $j$ with
$\norm{q_{n,j}}\le C_0C_1^n$ be given for $j\le k$, $n\ge1$. Then $\Lambda^k_\pi u$ is analytic at $x_0$, with $N$th
coefficient of norm at most $k!\binom{N+k-1}{k-1}\max(1,\norm a)^kC_0^kC_1^N$.
\end{lemma}

\begin{proof}
\emph{Step 1 (one content).} For a content $\mathbf m$ of support $n_1<\dots<n_m$ and type $\mathbf i$ let
$Q_l=q_{n_l,i_l}$ if $n_l\ge1$, $Q_l=\Lambda^{i_l}_\pi a$ if $n_l=0$, and, with consecutive blocks $I_l$ of sizes $n_li_l$,
$q_{\mathbf m}(h_1,\dots,h_{N(\mathbf m)})=\wedge^{\mathbf i}_{\widehat D}\circ(Q_1(h_{I_1})\otimes\dots\otimes Q_m(h_{I_m}))
\circ\sh^{\mathbf i}_E$. By Lemma~\ref{fmo:lem:multishuffle}(2), $q_{\mathbf m}(h^{\times N(\mathbf m)})=\Lambda^{\mathbf i}_\pi
(u_{\beta_{n_1}}(h),\dots,u_{\beta_{n_m}}(h))$ and $\norm{q_{\mathbf m}}\le\binom k{\mathbf i}\prod_l\norm{Q_l}\le
k!\max(1,\norm a)^kC_0^kC_1^{N(\mathbf m)}$, as $m\le k$ and $C_0,C_1\ge1$. The sum $q_N$ of the $q_{\mathbf m}$ of degree
$N$ obeys the stated bound, and $\sum_N\norm{q_N}\norm h^N<\infty$ for $\norm h<1/(2C_1)$.

\emph{Step 2 (the expansion).} For $\norm h<R$ and $y=x_1\wedge\dots\wedge x_k$, as in Lemma~\ref{fmo:lem:series}(2),
$\Lambda^k_\pi u(x_0+h)y=\omega_{\widehat D}(\sum_nu_{\beta_n}(h)x_1,\dots,\sum_nu_{\beta_n}(h)x_k)$ is the absolutely
convergent sum of the $u_{\beta_{\sigma(1)}}(h)x_1\wedge\dots\wedge u_{\beta_{\sigma(k)}}(h)x_k$, $\sigma\colon\{1,\dots,k\}
\to\N_0$. Grouped by the content of $(\sigma(1),\dots,\sigma(k))$ and then by degree, they give $\sum_Nq_N(h^{\times N})y$
(Lemma~\ref{fmo:lem:multishuffle}(1)); as pure wedges span, $\Lambda^k_\pi u(x_0+h)=\sum_Nq_N(h^{\times N})$ for
$\norm h<\min(R,1/(2C_1))$.
\end{proof}

\begin{lemma}\label{fmo:lem:steps}
\begin{enumerate}
\item If $(P;E,D)$ is good in degrees $a,b\ge1$ and $\binom{a+b}a\ne0$ in $K$, it is good in degree $a+b$.
\item Let $k\ge2$, with maximum norms on $E\oplus K$, $D\oplus K$. If $(P;E\oplus K,D\oplus K)$ is good in degree $k$, then
$(P;E,D)$ is good in degree $k-1$.
\end{enumerate}
\end{lemma}

\begin{proof}
(1) Let $n=a+b$. For $v\in\cL(E,\widehat D)$, naturality of the wedge product and~\eqref{fib:eq:hopf}, extended to
completions, give $\Lambda^n_\pi v=\binom na^{-1}\wedge^{a,b}_{\widehat D}\circ(\Lambda^a_\pi v\otimes\Lambda^b_\pi v)\circ
\sh^{a,b}_E$, bounded bilinear in $(\Lambda^a_\pi v,\Lambda^b_\pi v)$; apply Lemma~\ref{fmo:lem:series}(2).
(2) Let $u\colon U\to\cL(E,\widehat D)$ be analytic and $F$ a Banach space. As $u\oplus1$ is analytic into
$\cL(E\oplus K,\widehat D\oplus\Kh)$ and $\widehat D\oplus\Kh$ completes $D\oplus K$, goodness and the prose after
Definition~\ref{fam2:def:good} make $x\mapsto A^k_{E\oplus K,\widehat D\oplus\Kh;F}(u(x)\oplus1)$ analytic, hence also
$x\mapsto A^k_{E\oplus K,\widehat D\oplus K;F}(u(x)\oplus1)$, restriction $\Alt^k(\widehat D\oplus\Kh;F)\to
\Alt^k(\widehat D\oplus K;F)$ being an isometric isomorphism by density. The proof of Theorem~\ref{fib:thm:hyperplane},
with $\widehat D$ for $D$, gives
fixed bounded linear $R_h$, $J_h$ with $A^{k-1}_{E,\widehat D;F}(v)=R_h\circ A^k_{E\oplus K,\widehat D\oplus K;F}(v\oplus1)
\circ J_h$; so $x\mapsto A^{k-1}_{E,\widehat D;F}(u(x))$ is analytic for every Banach $F$, and the paragraph after
Definition~\ref{fam2:def:good} applies.
\end{proof}

\begin{proof}[Proof of Theorem~\ref{fam2:thm:reduction}]
\emph{Step 1 (item~(1)).} Let $(\beta_n)$ be as in~\eqref{fam2:eq:G}. Each $u_{\beta_n}$ is analytic
(Proposition~\ref{prop:homogeneous-lift}), so $c_j(\beta_n)<\infty$ by goodness and Lemma~\ref{fmo:lem:homogeneous}(3). If
$\limsup_nc_j(\beta_n)^{1/n}=\infty$, choose $n_1<n_2<\dots$ with $n_{i+1}>jn_i$ and $c_j(\beta_{n_i})^{1/n_i}\ge i$. By
Lemma~\ref{fmo:lem:series}(1) and goodness, $u(h)=\sum_i\beta_{n_i}(h^{\times n_i})$ is analytic on the open unit ball and
$\Lambda^j_\pi u(h)=\sum_Nq_N(h^{\times N})$ for $\norm h<R'$, $\norm{q_N}\le C'R'^{-N}$. As in Step~2 of the proof of
Lemma~\ref{fmo:lem:assembly}, $\Lambda^j_\pi u(h)(x_1\wedge\dots\wedge x_j)=\sum_Ns_N(h)$, with $s_N(h)$ the sum of the
$u_{\beta_{\sigma(1)}}(h)x_1\wedge\dots\wedge u_{\beta_{\sigma(j)}}(h)x_j$ over the $\sigma\colon\{1,\dots,j\}\to
\{n_1,n_2,\dots\}$ with $\sum_a\sigma(a)=N$, and Lemma~\ref{lem:unique} on the lines $t\mapsto th$ gives
$s_N(h)=q_N(h^{\times N})(x_1\wedge\dots\wedge x_j)$. For $N=jn_i$ only $\sigma\equiv n_i$ qualifies, as $n_{i+1}>jn_i$
and $jn_{i-1}<n_i$; so $q_{jn_i}$ lifts $u_{\beta_{n_i}}$, and $c_j(\beta_{n_i})^{1/n_i}\le C'^{1/n_i}R'^{-j}$ stays
bounded, a contradiction.

\emph{Step 2 (item~(2)).} Let $u\colon U\to\cL(E,\widehat D)$ be analytic, $u(x_0+h)=\sum_n\beta_n(h^{\times n})$,
$\norm{\beta_n}\le CR^{-n}$. For integers $m,m'\ge0$ with $\abs\vartheta^{m'}<R$ and $\abs\vartheta^{-m}\ge\max(C,1)$ and for
$\lambda=\vartheta^{m'}$, $\norm{\vartheta^m\lambda^n\beta_n}\le1$, so~\eqref{fam2:eq:G} gives $C'_j,C''_j\ge1$ with
$c_j(\vartheta^m\lambda^n\beta_n)<C'_jC''^n_j$ for $n\ge1$, $j\le k$, and Lemma~\ref{fmo:lem:homogeneous}(1) gives lifts of
$u_{\beta_n}$ in degree $j$ of norm at most $\abs\vartheta^{-mj}C'_j(C''_j\abs\lambda^{-j})^n$; apply
Lemma~\ref{fmo:lem:assembly} with $C_0=\max_j\abs\vartheta^{-mj}C'_j$ and $C_1=\max_jC''_j\abs\lambda^{-j}$.

\emph{Step 3 (item~(3)).} The first assertion is Lemma~\ref{fmo:lem:homogeneous}(2). In the second, ``only if'' is~(1);
conversely, we show goodness in degree $n\le k$ by induction. If $n<p$, apply~(2) and the first assertion. If $n>p$ is not
a power of $p$ and $p^l$ is the largest power of $p$ below $n$, then $n=p^l+(n-p^l)$ is carry-free, so $\binom n{p^l}\ne0$
in $K$ by Lucas's theorem (proof of Corollary~\ref{fib:cor:degrees}); apply Lemma~\ref{fmo:lem:steps}(1). If $n=p^i$,
$(\mathrm G_j)$ holds for $j<n$ by induction and~(1), and for $j=n$ by hypothesis; apply~(2).

\emph{Step 4 (item~(4)).} The linearization $\tau_\beta\colon Y_n\to\widehat D$ of $(h_1,\dots,h_n,y)\mapsto
\beta(h_1,\dots,h_n)y$ has norm $\norm\beta$ and $\beta=\tau_\beta\circ\beta^{\mathrm u}_n$; apply
Lemma~\ref{fmo:lem:homogeneous}(4). If $(P;E,D_{\mathrm u})$ is good in degrees $j\le k$, then, with $J_n$ the $n$th
injection of $D_{\mathrm u}$, $(\mathrm G_j)$ for $\beta_n=J_n\circ\beta^{\mathrm u}_n$, of norm at most $1$, and $c_j(\beta^{\mathrm u}_n)\le
c_j(\beta_n)$, by Lemma~\ref{fmo:lem:homogeneous}(4) with the $n$th projection, make every $c_j(\beta^{\mathrm u}_n)$ finite
with geometric growth; with the inequality, this gives $(\mathrm G_j)$ for every normed $D$, hence goodness in degrees
$j\le k$ by~(2). In characteristic $p$ the degrees $p^i\le k$ give $(\mathrm G_{p^i})$ for every $D$, and~(3) applies.
\end{proof}

\begin{proof}[Proof of Corollary~\ref{fam2:cor:degrees}]
(1) Degrees below $p$ are good by Theorem~\ref{fam2:thm:reduction}(2),(3), carry-free sums by Lemma~\ref{fmo:lem:steps}(1)
and Lucas's theorem; a least missing $N$ that is not a power of $p$ would be the carry-free sum $p^l+(N-p^l)$ of smaller
elements, $p^l$ the largest power of $p$ below $N$. (2) By Proposition~\ref{fam2:prop:reductions}(2), $(P;E,D)$ and
$(P;E\oplus K,D\oplus K)$ are good in the same degrees, so Lemma~\ref{fmo:lem:steps}(2) makes $\fD(P;E,D)$ closed under
$k\mapsto k-1$, $k\ge2$; apply~(1). (3) is Proposition~\ref{fam2:prop:fibers}. (4) Lemma~\ref{fmo:lem:steps}(2), for all
normed $E$, $D$ at once, gives closure under $k\mapsto k-1$; by~(1) the set contains $1,\dots,p-1$ and is closed under
carry-free sums, so its least missing element is a power of $p$.
\end{proof}

\subsection{The familywise class}\label{fmo:ssec:class}

\begin{lemma}\label{fmo:lem:germs}
Let $u(x_0+h)=\sum_n\beta_n(h^{\times n})$ for $\norm h<R$, $\beta_n\in\Mult^n(P;\cL(E,\widehat D))$,
$\norm{\beta_n}\le CR^{-n}$, and $\lambda\in K$, $0<\abs\lambda<R$. Let $\beta^\otimes_n\colon\hot^n_\pi P\to\cL(E,\widehat D)$
be the linearization of $\beta_n$ for $n\ge1$, and $\beta^\otimes_0(c)=c\beta_0$ for $c$ in slot $0$, $\Kh$, by the
$\Kh$-vector space structure of $\cL(E,\widehat D)$. Then $\theta((w_n)_n)=\sum_n\lambda^n\beta^\otimes_n(w_n)$ is a linear
family on $\mathcal T(P)$, $\norm\theta\le C$, and $u(x_0+\lambda h)=\theta(\operatorname{ver}(h))$ for $\norm h<1$.
Conversely, $\theta\circ\operatorname{ver}$ is analytic on the open unit ball for every linear family $\theta$ on
$\mathcal T(P)$.
\end{lemma}

\begin{proof}
$\norm{\lambda^n\beta^\otimes_n(w_n)}\le C(\abs\lambda/R)^n\norm{w_n}$, and $\theta(\operatorname{ver}(h))=\sum_n\beta_n
((\lambda h)^{\times n})$. The converse is Lemma~\ref{fmo:lem:series}(1), as $\norm{\theta\circ\operatorname{ver}_n}\le
\norm\theta$.
\end{proof}

\begin{lemma}\label{fmo:lem:veronese}\emergencystretch=3em
Let $Y$ be a Banach space, $0<R'\le R<1$, and $C_{k,R}=\sup_N\binom{N+k-1}{k-1}R^N<\infty$. For
$\Phi\in\Mult^k(\mathcal T(P);Y)$, $\norm{\operatorname{ver}^\sharp\Phi}_R\le C_{k,R}\norm\Phi$ and
$\Phi(\operatorname{ver}(h),\dots,\operatorname{ver}(h))=\sum_N(\operatorname{ver}^\sharp\Phi)_N(h^{\times N})$ for
$\norm h<1$; if $\Phi$ has zero diagonal, $\operatorname{ver}^\sharp\Phi\in\Mult^\omega_{R,0}(P;Y)$; and
$\iota_{R'R}\circ\operatorname{ver}^\sharp$ is $\operatorname{ver}^\sharp$ with values in $\Mult^\omega_{R'}(P;Y)$. Hence
vanishing after shrinking does not depend on $R$.
\end{lemma}

\begin{proof}
There are at most $\binom{N+k-1}{k-1}$ compositions of $N$ into $k$ parts, and $\norm{\operatorname{ver}_n}\le1$; the identity
is Lemma~\ref{fmo:lem:series}(2) with radius $1$. For zero diagonal, $\sum_Nt^N(\operatorname{ver}^\sharp\Phi)_N(h^{\times N})=0$
for $\abs t<1/\norm h$; apply Lemma~\ref{lem:unique}. The inclusion keeps components, so $(\iota_{R'R})_*\ob^\omega_k(\theta;R)=
\ob^\omega_k(\theta;R')$ by~\eqref{fam2:eq:class} and Corollary~\ref{ocl:cor:functorial}(b), and pushouts of $0$ vanish.
\end{proof}

\begin{proof}[Proof of Theorem~\ref{fam2:thm:class}]
\emph{Step 1 (item~(1)).} The second and third conditions are equivalent: bounded linear maps into
$\Mult^\omega_{R',0}(P;\hLambda^k_\pi D)$ are sequences of maps into the $\Mult^N_0(P;\hLambda^k_\pi D)$ with
$\sup_NR'^N\norm\cdot<\infty$, $\iota_{R'R}$ keeps components, and a pushout vanishes if and only if the map extends
(Lemma~\ref{ob:lem:split-criteria}(3)). Each $q_N\in\Mult^N(P;\cL(\Lambda^k_\pi E,\hLambda^k_\pi D))$ gives
$q_N^\sharp\colon\hLambda^k_\pi E\to\Mult^N(P;\hLambda^k_\pi D)$, $q_N^\sharp(y)(h_1,\dots,h_N)=q_N(h_1,\dots,h_N)y$, of the
same norm, and conversely. If $\Lambda^k_\pi\theta(\operatorname{ver}(h))=\sum_Nq_N(h^{\times N})$ for $\norm h<R_0$,
$\norm{q_N}\le C'R_0^{-N}$, then for $z\in\hot^k_\pi E$ the $N$th terms of $\operatorname{ver}^\sharp\ev_\theta(z)$ and
$q_N^\sharp\widehat\wedge^k_Ez$ have the same diagonal, both giving the $t^N$-coefficient of
$\ev_\theta(z)(\operatorname{ver}(th),\dots,\operatorname{ver}(th))=\hLambda^k_\pi\theta(\operatorname{ver}(th))
(\widehat\wedge^k_Ez)$ (Lemmas~\ref{fmo:lem:veronese} and~\ref{lem:unique}). So
$\xi_N=\operatorname{pr}_N\operatorname{ver}^\sharp\ev_\theta-q_N^\sharp\widehat\wedge^k_E$ takes values in
$\Mult^N_0$, extends $(\operatorname{pr}_N\operatorname{ver}^\sharp
\ev_\theta)|_{\cN^k_E}$ and has norm at most $C_{k,R}\norm\theta^kR^{-N}+C'R_0^{-N}$. Conversely, extensions $\xi_N$ of norm
at most $CR'^{-N}$ make $\operatorname{pr}_N\operatorname{ver}^\sharp\ev_\theta-\xi_N$ vanish on $\cN^k_E$, so it is
$q_N^\sharp\widehat\wedge^k_E$ with $\norm{q_N^\sharp}\le(C_{k,R}\norm\theta^k+C)\min(R,R')^{-N}$
(Lemma~\ref{ob:lem:metric-surjection}), and $\sum_Nq_N(h^{\times N})=\Lambda^k_\pi\theta(\operatorname{ver}(h))$ for
$\norm h<\min(R,R')$ by Lemma~\ref{fmo:lem:veronese}.

\emph{Step 2 (item~(2)).} For analytic $u$ and $x_0$, Lemma~\ref{fmo:lem:germs} gives $\lambda$, $\theta$ with
$u(x_0+\lambda h)=\theta(\operatorname{ver}(h))$, and~(1), with $h=\lambda^{-1}(x-x_0)$, gives analyticity of
$\Lambda^k_\pi u$ at $x_0$. Conversely, $\theta\circ\operatorname{ver}$ is analytic on the open unit ball
(Lemma~\ref{fmo:lem:germs}); apply goodness and~(1).

\emph{Step 3 (item~(3)).} As $\theta(w)=\tau_\theta\circ\theta_{\mathrm u}(w)$, $\ev_\theta=(\hLambda^k_\pi\tau_\theta)_\flat
\ev_{\theta_{\mathrm u}}$ on pure tensors. The map $(\hLambda^k_\pi\tau_\theta)_\flat$ commutes with
$\operatorname{ver}^\sharp$ and $\iota_{R'R}$, preserves zero diagonals and has norm at most $\norm{\tau_\theta}^k$; so
Corollary~\ref{ocl:cor:functorial}(b) gives the identity, vanishing after shrinking passes from $\theta_{\mathrm u}$ to every
$\theta$, and~(2) gives goodness for every $D$. Conversely, (2) for the Banach space $D=\mathcal T(P)\hot_\pi E$ gives
vanishing for $\theta_{\mathrm u}$. Finally $\theta_{\mathrm u}\circ\operatorname{ver}=\mathcal U$; apply~(1).
\end{proof}

\subsection{Polarizable base models}\label{fmo:ssec:polarizable}

\begin{lemma}\label{fmo:lem:regroup}
Let $n,i\ge1$ and $R_i$ a bounded projection of $\hot^i_\pi P$ onto $\hGamma^i_\pi P$. Identify $\hot^{ni}_\pi P$ with
$(\hot^i_\pi P)^{\hot n}$ by grouping the positions $(a,1),\dots,(a,i)$ into the $a$th factor, and let $\operatorname{reg}$ be
the isometry of $\hot^{ni}_\pi P$ moving position $(a,b)$ to $(b,a)$, with target $(\hot^n_\pi P)^{\hot i}$. Then
$q_{n,i}(h_1,\dots,h_{ni})=\operatorname{reg}\circ R_i^{\otimes n}(h_1\otimes\dots\otimes h_{ni})$ is bounded $ni$-linear,
$\norm{q_{n,i}}\le\norm{R_i}^n$, with values in $\hGamma^i_\pi(\hot^n_\pi P)$, and $q_{n,i}(h^{\times ni})=
(h^{\otimes n})^{\otimes i}$.
\end{lemma}

\begin{proof}
$R_i^{\otimes n}$ has norm at most $\norm{R_i}^n$, takes values in the closed span of the $y_1^{\otimes i}\otimes\dots\otimes
y_n^{\otimes i}$, which $\operatorname{reg}$ sends to $(y_1\otimes\dots\otimes y_n)^{\otimes i}$, and fixes
$(h^{\otimes i})^{\otimes n}$, which $\operatorname{reg}$ sends to $(h^{\otimes n})^{\otimes i}$.
\end{proof}

\begin{lemma}\label{fmo:lem:per}
Let $G$ be a Banach space and $\mathbf m$ a content of type $\mathbf i=(i_1,\dots,i_m)$, let $\Omega_{\mathbf m}$ represent
the left cosets in $S_k$ of the subgroup preserving the consecutive blocks of sizes $i_1,\dots,i_m$, acting on
$\hot^k_\pi G$ by permuting factors, and $\operatorname{per}_{\mathbf m}=\sum_{\sigma\in\Omega_{\mathbf m}}\sigma$. Then
$\norm{\operatorname{per}_{\mathbf m}}\le\binom k{\mathbf i}$; on tensors invariant under that subgroup,
$\operatorname{per}_{\mathbf m}$ does not depend on $\Omega_{\mathbf m}$; for $y_1,\dots,y_m\in G$,
$\operatorname{per}_{\mathbf m}(y_1^{\otimes i_1}\otimes\dots\otimes y_m^{\otimes i_m})$ is the coefficient of
$t_1^{i_1}\cdots t_m^{i_m}$ in $(\sum_lt_ly_l)^{\otimes k}$ and lies in the span of the powers $y^{\otimes k}$, $y\in G$;
and $\operatorname{per}_{\mathbf m}$ maps the closed span of such tensors into $\hGamma^k_\pi G$.
\end{lemma}

\begin{proof}
The bound is the triangle inequality; on invariant tensors $\sigma$ and $\sigma\sigma'$, $\sigma'$ in the subgroup, agree.
The cosets correspond bijectively to the words in which $l$ occurs $i_l$ times, whose sum is the coefficient. A linear
functional on the span of the coefficients of $(\sum_lt_ly_l)^{\otimes k}$ vanishing on these powers vanishes on every
coefficient, by Lemma~\ref{lem:idprinc}, as in Step~1 of the proof of Proposition~\ref{spl:prop:gamma-l1}. Continuity gives
the last claim.
\end{proof}

\begin{proof}[Proof of Theorem~\ref{fam2:thm:polarizable}]
\emph{Step 1 (affine maps suffice).} Affine maps are analytic. Conversely, fix $n\le k$, let $G=\Kh\oplus\widehat P$ with
the maximum norm, $\gamma(x)=(1,x)$, $Z=\hot^n_\pi P=\hot^n_\pi\widehat P$ and $W=\hGamma^n_\pi P=\hGamma^n_\pi\widehat P$
(density; paragraph before Theorem~\ref{ob:thm:mechanisms}). Scalar multiplication on $Z$ extends from $K$ to $\Kh$
by continuity and preserves the closed subspace $W$. The $k$-linear map $\Phi((\lambda_1,x_1),\dots,(\lambda_k,x_k))=
\lambda_{n+1}\cdots\lambda_k\,x_1\otimes\dots\otimes x_n$ has $\norm\Phi\le1$ and $\Phi((\lambda,x),\dots,(\lambda,x))=
\lambda^{k-n}x^{\otimes n}\in W$, so
$x\mapsto x^{\otimes n}=\Phi(\gamma(x),\dots,\gamma(x))$ is analytic at $0$ into $W$; Proposition~\ref{prop:homogeneous-lift},
with the canonical map $P^n\to Z$, and Lemma~\ref{ob:lem:parameter} give polarizability in degree $n$.

\emph{Step 2 (the ambient sum).} Let $G,Z,W,\Phi,U,\gamma$ be as in Definition~\ref{fam2:def:reflection}, $x_0\in U$ and
$\gamma(x_0+h)=\sum_n\gamma_n(h^{\times n})$ for $\norm h<R$, $\norm{\gamma_n}\le CR^{-n}$. Let $\gamma^\otimes_n\colon
\hot^n_\pi P\to G$ and $\Phi^\otimes\colon\hot^k_\pi G\to Z$ be the linearizations, with $\hot^0_\pi P=\Kh$ and
$\gamma^\otimes_0(c)=c\gamma_0$; as $W$ is closed, $\Phi^\otimes(\hGamma^k_\pi G)\subseteq W$. By
Lemma~\ref{fmo:lem:series}(2), for $\norm h<R$, $\Phi(\gamma(x_0+h),\dots,\gamma(x_0+h))$ is the absolutely convergent sum of
the $\Phi(\gamma_{n'_1}(h^{\times n'_1}),\dots,\gamma_{n'_k}(h^{\times n'_k}))$, $(n'_1,\dots,n'_k)\in\N_0^k$, and by
Lemma~\ref{fmo:lem:per} the terms of content $\mathbf m$, of support $n_1<\dots<n_m$ and type $\mathbf i$, sum to
$\Phi^\otimes\operatorname{per}_{\mathbf m}(\gamma_{n_1}(h^{\times n_1})^{\otimes i_1}\otimes\dots\otimes
\gamma_{n_m}(h^{\times n_m})^{\otimes i_m})$.

\emph{Step 3 (lifting each content).} Fix $\bar g$ as in~(2) and projections $R_i$ onto $\hGamma^i_\pi P$ with
$\norm{R_i}\le\bar g$, $i\le k$. Let $q_l=q_{n_l,i_l}$ (Lemma~\ref{fmo:lem:regroup}) if $n_l\ge1$, $q_l=1\in\Kh$ if $n_l=0$,
and, with consecutive blocks $I_l$ of sizes $n_li_l$, $Q_{\mathbf m}(h_1,\dots,h_{N(\mathbf m)})=\Phi^\otimes
\operatorname{per}_{\mathbf m}(\bigotimes_l(\gamma^\otimes_{n_l})^{\otimes i_l}q_l(h_{I_l}))$. Then $\norm{Q_{\mathbf m}}\le
\norm\Phi\binom k{\mathbf i}\prod_l\norm{\gamma_{n_l}}^{i_l}\bar g^{n_l}\le\norm\Phi k!C^k(\bar g/R)^{N(\mathbf m)}$, as
$\bar g\ge1$. Its values lie in $W$, as $q_l$ takes values in $\hGamma^{i_l}_\pi(\hot^{n_l}_\pi P)$, which
$(\gamma^\otimes_{n_l})^{\otimes i_l}$ maps into $\hGamma^{i_l}_\pi G$, and Lemma~\ref{fmo:lem:per} applies. By
Lemma~\ref{fmo:lem:regroup}, $Q_{\mathbf m}(h^{\times N(\mathbf m)})$ is the content-$\mathbf m$ sum of Step~2.

\emph{Step 4 (item~(1)).} The sum $q_N$ of the $Q_{\mathbf m}$ of degree $N$ is bounded $N$-linear into $W$, with
\begin{equation}\label{fmo:eq:reflection}
 \norm{q_N}\le\norm\Phi\,k!\,C^k\tbinom{N+k-1}{k-1}(\bar g/R)^N ,
\end{equation}
and for $\norm h<R/\bar g$, $\sum_Nq_N(h^{\times N})$ converges in the Banach space $W$ to $\Phi(\gamma(x_0+h),\dots,
\gamma(x_0+h))$, by Step~2.

\emph{Step 5 (items~(2)--(4)).} (2) For analytic $u\colon U\to\cL(E,\widehat D)$ take $G=\cL(E,\widehat D)$,
$Z=\Mult^k(E;\hLambda^k_\pi D)$, $W$ the image of $\cL(\Lambda^k_\pi E,\hLambda^k_\pi D)$ under the closed isometry
$T\mapsto T\circ\omega_E$ (Proposition~\ref{prop:exterior}), and $\Phi(u_1,\dots,u_k)(x_1,\dots,x_k)=u_1x_1\wedge\dots\wedge
u_kx_k$, of norm at most $1$, with $\Phi(u,\dots,u)=\Lambda^k_\pi u\circ\omega_E$; apply~(1) and~\eqref{fmo:eq:reflection}.
For bundles, $x\mapsto A^k_{E,\widehat D;F}(u(x))$ is analytic for Banach $F$ (prose after Definition~\ref{fam2:def:good}),
and the last assertion of Theorem~\ref{thm:bundle-lifting} applies as in the proof of Corollary~\ref{fam2:cor:bundles}.
(3) If $n!\ne0$ in $K$, $P$
is polarizable in degree $n$ (Proposition~\ref{can:prop:natural}(2)); so~(2) applies if $\operatorname{char}K=0$ or
$k<p$. Polarizability in
degree $p$ gives it in every degree (Theorem~\ref{cmb:thm:polarizable-height}). A lower $r$-estimate, $r<2$, bounds
$\cP_<$ (Theorem~\ref{fib:thm:lower}(1)), and Proposition~\ref{spl:prop:ordering}(3) gives polarizability in every degree.
(4) Projections with $\norm{R_n}\le2g_n(P)$ and Theorem~\ref{fam:thm:tensor-analytic} give the first equivalence, and~(1)
the second. For nonarchimedean $K$, $g_n(\ell^1(\N,K))=n!$ (Propositions~\ref{spl:prop:gamma-l1}(1)
and~\ref{fam:prop:l1-optimal}) is finite, while $(n!)^{1/n}$ is unbounded; apply~(1) and~(2).
\end{proof}

\begin{proof}[Proof of Theorem~\ref{fam2:thm:bases}]
\emph{Step 1 (item~(1)).} As the statement depends only on $K$ as a normed field, Lemma~\ref{obs:lem:normalization} lets us
assume $K$ dense in $\kappa((X))$, as in the proof of Theorem~\ref{obs:thm:ptb}. By Lemma~\ref{obs:lem:multiplier}(1) for
$P$, $\norm{ay}_P\le\norm y_\infty\norm a_P$ for $a\in P$ and $y\in E_0$, so $\theta_0$ is a linear family, and
$\norm{\theta_0(a)}=\norm{\mu_a}_{E_0\to P}=\norm a_P$ for $a\in c_{00}(\N,K)$ (Lemma~\ref{obs:lem:multiplier}(2),(4)). The
fifth condition gives the first, the first the second. If $\Lambda^k_\pi\theta_0$ is analytic at a point, $\theta_0$ has a
lift $Q$ in degree $k$ (prose after Definition~\ref{fam2:def:good}). Its restriction to $c_{00}(\N,K)^k$ is a multiplier lift for
$(E_0,P)$ (Section~\ref{frc:ssec:coefficients}) satisfying~\eqref{frc:eq:mult-bound} with $C_Q=\norm Q$, and $E_0$, $P$ are
symmetric sequence spaces; so Steps~1--5 of the proof of Theorem~\ref{obs:thm:ptb} apply to $(E_0,P)$, as in Step~2 of the
proof of Proposition~\ref{obs:prop:gamma-ptb}, and $x=1_{\operatorname{supp}a}$, $y=1_{\operatorname{supp}b}$
in~\eqref{obs:eq:ptb} bound $\cP_<$ on $P\otimes_\pi P$. This bound, item~(3) of Theorem~\ref{cls:thm:symmetric}, gives
polarizability in degree $p$ (Proposition~\ref{spl:prop:ordering}(3)), hence the fifth condition
(Theorem~\ref{fam2:thm:polarizable}(3)). The examples follow from Corollary~\ref{cls:cor:symmetric-examples} and
Theorem~\ref{fam2:thm:polarizable}(3), every $k\ge p$ being admissible in~(L).

\emph{Step 2 (item~(2)).} Apply Theorem~\ref{fib:thm:marking} with $E=c_{00}(\N,K)$, $E'=\linf(\N,\Kh)$, $\psi=\id$,
$\theta=\id$, $u_a=\mu_a=\mu_{w(i(a))}$, $F=\hLambda^k_\pi E'$, $\nu=\omega_{E'}$, $\theta_F=\id$: no bounded $k$-linear $Q$ on
$c_{00}(\N,K)^k$ has $Q(a^{\times k})=\omega_{E'}\circ(\mu_a,\dots,\mu_a)$. Analyticity of $x\mapsto\Lambda^k_\pi\mu_{w(x)}$
at a point would give a lift $Q_P$ of the linear family $x\mapsto\mu_{w(x)}$ (prose after Definition~\ref{fam2:def:good}),
and $(a_1,\dots,a_k)\mapsto Q_P(ia_1,\dots,ia_k)\circ\omega_E$ would be such a $Q$. For $c_0(\Sigma,K)$ and
$\linf(\Sigma,K)$ take $\N\subseteq\Sigma$, $i$ extension by zero and $w$ restriction; for $P$
with an equivalent nonarchimedean norm and $\dim_{\Kh}\widehat P=\infty$, take $i(\xi)=\sum_n\xi_ne_n$ and $w=\theta_P|_P$
with $e_n$, $\theta_P$ from Lemma~\ref{fib:lem:marking}. Then $(\mathrm U^k_P)$ fails, and with it polarizability in degree
$p$ (Theorem~\ref{fam2:thm:polarizable}(3)). If $\dim_{\Kh}\widehat P=d<\infty$, a bounded $\Kh$-linear isomorphism
$T\colon\widehat P\to\ell^1(\{1,\dots,d\},\Kh)$ exists as $\Kh$ is complete, and $(T^{-1})^{\otimes n}R_nT^{\otimes n}$
transfers projections (Step~3 of the proof of Proposition~\ref{spl:prop:gamma-l1}); as passing to a dense subspace changes
neither $\hot^n_\pi$ nor $\hGamma^n_\pi$, Proposition~\ref{spl:prop:gamma-l1}(1) for $\ell^1(\{1,\dots,d\},K)$ makes $P$
polarizable in every degree; apply Theorem~\ref{fam2:thm:polarizable}(2).
\end{proof}

\subsection{Fibers, parameter-relative ordering and the staircase}\label{fmo:ssec:mixed}

\begin{proof}[Proof of Proposition~\ref{fam2:prop:fibers}]
(1) For an analytic $u\colon U\to\cL(E,\widehat D)$, $\Lambda^k_\pi u(x)=Q(u(x),\dots,u(x))$ is analytic on $U$ by
Lemma~\ref{fmo:lem:series}(2). For $\beta$ as stated, $q(h_1,\dots,h_{kn})=Q(\beta(h_{I_1}),\dots,\beta(h_{I_k}))$, with
$I_1,\dots,I_k$ the consecutive blocks of size $n$, is a lift of $u_\beta$ in degree $k$ of norm at most
$\norm Q\norm\beta^k$. (2) The ``if'' direction is~(1). Conversely, the identity of the open set $\cL(E,\widehat D)$ is an
analytic family, so $u\mapsto\Lambda^k_\pi u$ is analytic at some point of $\cL(E,\widehat D)$; since
$\hLambda^k_\pi\widehat D=\hLambda^k_\pi D$, Lemma~\ref{fib:lem:test} applies to $(E,\widehat D)$. The second assertion
follows from~(1) and the first.
\end{proof}

\begin{proof}[Proof of Theorem~\ref{fam2:thm:gamma-split}]
\emph{Step 1 (the map $\Pi$).} Let $G_r$ and $\delta_t$ be as in Section~\ref{ord:ssec:triangle}; by density
$\hot^k_\pi G_r=\hot^k_\pi\ell^r(\Sigma,\Kh)$ and $\hGamma^k_\pi G_r=\hGamma^k_\pi\ell^r(\Sigma,\Kh)$. For a composition
$\mathbf l=(l_1,\dots,l_m)$ of $k$ let $I_1,\dots,I_m$ be the consecutive blocks of sizes $l_i$ and $\Omega_{\mathbf l}$
represent the left cosets of $S_{l_1}\times\dots\times S_{l_m}$ in $S_k$, acting by permuting factors. For
$a_1,\dots,a_k\in c_{00}(\Sigma,K)$ put
\[\Pi_{\mathbf l}(a_1,\dots,a_k)=\sum_{t_1<\dots<t_m}\prod_{i=1}^m\prod_{j\in I_i}a_j(t_i)\,\delta_{t_1}^{\otimes l_1}
\otimes\dots\otimes\delta_{t_m}^{\otimes l_m},\qquad\Pi=\sum_{\mathbf l}\sum_{\sigma\in\Omega_{\mathbf l}}\sigma\circ
\Pi_{\mathbf l}.\]
As $\sum_{\sigma\in\Omega_{\mathbf l}}\sigma\delta_{\mathbf t}=O_{\mathbf t}$ for nondecreasing $\mathbf t$ with constant
blocks of sizes $\mathbf l$, $\Pi$ is $R_k$ of Proposition~\ref{spl:prop:ordering}(3), for $(\delta_t)$, applied to
$a_1\otimes\dots\otimes a_k$. So $\Pi(a,\dots,a)=a^{\otimes k}$, and $\Pi$ takes values in the span of the $O_{\mathbf t}$,
which lie in the span of the powers by Step~4 of the proof of that proposition in Appendix~\ref{app:ordering}.

\emph{Step 2 (one composition).} Let $\Delta=\{i:l_i\ge2\}$, $j_i$ the element of $I_i$ for $i\notin\Delta$, and
$b_i=\prod_{j\in I_i}a_j$ for $i\in\Delta$, so $\norm{b_i}_1\le\prod_{j\in I_i}\norm{a_j}_2$ by H\"older's inequality and
$\norm\cdot_{l_i}\le\norm\cdot_2$. For fixed strictly increasing $(t_i)_{i\in\Delta}$, the $i\notin\Delta$ form maximal runs
of consecutive indices, and for a run the admissible $t_i$ are the strictly increasing tuples in the open interval $J$ of
$\Sigma$ between the neighboring fixed values. So $\Pi_{\mathbf l}(a_1,\dots,a_k)$ is the sum over $(t_i)_{i\in\Delta}$ of
$\prod_{i\in\Delta}b_i(t_i)$ times the tensor product, in block order, of the $\delta_{t_i}^{\otimes l_i}$, $i\in\Delta$, of
norm $1$, and, for each run of length $n$, of $\cP^{(n)}(\bigotimes_i
a_{j_i}1_J)$, of norm at most $C'_n\prod_i\norm{a_{j_i}}_2$ with $C'_n=\max(1,C_{n,r}\abs\vartheta^{-n})$
(Theorem~\ref{ord:thm:triangle}). Hence $\norm{\Pi_{\mathbf l}(a_1,\dots,a_k)}\le C_{\mathbf l}\prod_j\norm{a_j}_2$, with
$C_{\mathbf l}$ the product of the $C'_n$ of the runs.

\emph{Step 3 (item~(1)).} So $\Pi$ extends from the dense subspace $c_{00}(\Sigma,K)$ of $\ell^2(\Sigma,\Kh)$, with values
in the closed subspace $\hGamma^k_\pi G_r$, and $\Pi(a,\dots,a)=a^{\otimes k}$ persists by continuity.

\emph{Step 4 (item~(2), linear maps).} Let $T\colon P\to G$ be bounded linear, and extend $\Phi$ to $\widehat G^k$; its
diagonal stays in $W$. The coordinate map $V\to\ell^2(\Sigma,\Kh)$ extends to $w\colon P\to\ell^2(\Sigma,\Kh)$,
$\norm w\le C_2$; write $h(t)=(wh)(t)$ and $\iota\colon\ell^2\to\ell^r$ for the inclusion. As in Step~1 of the proof of
Theorem~\ref{spl:thm:relative-ordering}, with $P$, $G$, $T$ for $E$, $D$, $u$, the map $\xi\mapsto\sum_t\xi(t)T(e_t)$
extends from $c_{00}(\Sigma,K)$ to $\widetilde T\colon\ell^r(\Sigma,\Kh)\to\widehat G$ with $T=\widetilde T\circ\iota\circ w$.
The linearization $\Phi_r^\otimes$ of $(\xi_1,\dots,\xi_k)\mapsto\Phi(\widetilde T\xi_1,\dots,\widetilde T\xi_k)$, extended
to $\hot^k_\pi G_r$, maps $\hGamma^k_\pi G_r$ into $W$, and $Q(a_1,\dots,a_k)=\Phi_r^\otimes\Pi(wa_1,\dots,wa_k)$ satisfies
$Q(a,\dots,a)=\Phi_r^\otimes((\iota wa)^{\otimes k})=\Phi(Ta,\dots,Ta)$, by continuity from $c_{00}(\Sigma,K)$.

\emph{Step 5 (item~(2), analytic maps).} Let $\gamma(x_0+h)=\sum_n\gamma_n(h^{\times n})$ for $\norm h<R$,
$\norm{\gamma_n}\le CR^{-n}$, let $\Sigma_*$ be the disjoint union of the $\Sigma^n$, $n\ge0$, totally ordered in any way,
and $\lambda\in K$ with $0<\abs\lambda\le R/C_e$, $C_e=\sup_t\norm{e_t}$. For $\norm h<\abs\lambda/C_2$ the vector
$\xi(h)\in\ell^2(\Sigma_*,\Kh)$
with entries $\xi(h)_{(n,\mathbf t)}=\lambda^{-n}h(t_1)\cdots h(t_n)$ has $n$th block of norm
$(\norm{wh}_2/\abs\lambda)^n$, and $\xi$ is analytic at $0$ with $n$th coefficient of norm at most $(C_2/\abs\lambda)^n$. As
$\norm{\lambda^n\gamma_n(e_{t_1},\dots,e_{t_n})}\le C$, \eqref{spl:eq:r-norm} for $\widehat G$
(Proposition~\ref{spl:prop:r-norm}) extends $\delta_{(n,\mathbf t)}\mapsto
\lambda^n\gamma_n(e_{t_1},\dots,e_{t_n})$ to $T_*\colon\ell^r(\Sigma_*,\Kh)\to\widehat G$, and $T_*\iota\xi(h)=
\sum_n\gamma_n(h^{\times n})=\gamma(x_0+h)$ for $\norm h<\min(R,\abs\lambda/C_2)$: first for $h\in V$, where the blocks of
$\xi(h)$ are finitely supported, then by continuity. Step~4 for
$\ell^2(\Sigma_*,\Kh)$ with its unit vectors, $\widehat G$ and $T_*\circ\iota$ gives a bounded $k$-linear $Q_{T_*}$ into $W$
with $Q_{T_*}(\eta,\dots,\eta)=\Phi(T_*\iota\eta,\dots,T_*\iota\eta)$; so $\Phi(\gamma(x_0+h),\dots,\gamma(x_0+h))=
Q_{T_*}(\xi(h),\dots,\xi(h))$ near $0$, analytic by Lemma~\ref{fmo:lem:series}(2).

\emph{Step 6 (item~(3)).} The space $\widehat D$ satisfies~\eqref{spl:eq:r-norm}: by Proposition~\ref{spl:prop:r-norm} and
continuity if $D$ has an equivalent $r$-norm, and, if $D$ has an equivalent nonarchimedean norm, as
$\max_t\abs{\xi_t}\le(\sum_t\abs{\xi_t}^r)^{1/r}$, as in the proof of Corollary~\ref{spl:cor:relative-ordering-na}. Hence
so does $G=\cL(E,\widehat D)$, as $\norm{\sum_t\xi_tu_t}=\sup_{\norm x\le1}\norm{\sum_t\xi_tu_tx}$, and $G$ has an equivalent
$r$-norm (Proposition~\ref{spl:prop:r-norm}). Let $Z=\cL(\hot^k_\pi E,\hLambda^k_\pi D)$, $W$ its operators vanishing on
$\cN^k_E$, and $\Phi(u_1,\dots,u_k)=\widehat\wedge^k_D\circ(u_1\otimes\dots\otimes u_k)$, so $\Phi(u,\dots,u)=
\hLambda^k_\pi u\circ\widehat\wedge^k_E$. For analytic $u\colon U\to G$, (2) makes $x\mapsto\hLambda^k_\pi u(x)\circ
\widehat\wedge^k_E$ analytic into $W$; the inverse of the isometry $S\mapsto S\circ\widehat\wedge^k_E$ onto $W$
(Lemma~\ref{ob:lem:metric-surjection}) and restriction to $\Lambda^k_\pi E$ give analyticity of $\Lambda^k_\pi u$.
\end{proof}

\begin{proof}[Proof of Theorem~\ref{fam2:thm:staircase}]
\emph{Step 1 (the input from Theorem~\ref{fib:thm:ell2}).} Suppose such a $Q$ exists. Step~1 of the proof of
Theorem~\ref{fib:thm:ell2} in Appendix~\ref{app:forcing} uses only that $Q$ is $k$-linear on $c_{00}(\N,K)^k$ with
$Q(a,\dots,a)=\Lambda^k_\pi\mu_a$, and that $\norm y_\infty\le\norm y_E$, so that the determinant array of
Section~\ref{frc:ssec:costs} is defined on $\hLambda^k_\pi E$; it gives $\tau$, slots $\gamma_1\ne\gamma_2$ and $\delta\ne0$,
independent of $m$.

\emph{Step 2 (the norm bound for $\zeta$).} Fix $m\ge1$, and let $a_1,\dots,a_k$ and $x_1,\dots,x_k$ be the inputs of
Step~2 of that proof: their entries lie in $\{0,\pm1\}$, $a_{\gamma_1}$ and $a_{\gamma_2}$ have $2m$ nonzero entries and the
other $a_\gamma$ two, and $x_{\gamma_1}$, $x_{\gamma_2}$ have $4m$ and the other $x_\gamma$ four. So
$\zeta=Q(a_1,\dots,a_k)(x_1\wedge\dots\wedge x_k)$ has
$\norm\zeta\le\norm Q\varphi^\pm_P(2)^{k-2}\varphi^\pm_E(4)^{k-2}\varphi^\pm_P(2m)^2\varphi^\pm_E(4m)^2$.

\emph{Step 3 (the rank bound).} The slice of Lemma~\ref{fib:lem:slice} for the slots $\gamma_1,\gamma_2$ and
$c_\gamma=C_\gamma(1)$, followed by restriction to the coordinates $A_n(1)$ and $B_{n'}(1)$, contractive as
$\norm y_\infty\le\norm y_E$, gives $z$ with $\coeff_0z=\mathsf Z$ and $\norm z\le k(k-1)\norm\zeta$; so
Lemma~\ref{fib:lem:coeff}(2) and $\rank\mathsf Z\ge m-1$ give $m-1\le M_\varrho k(k-1)\norm\zeta$ for every $m$, contradicting
the hypothesis by Step~2.

\emph{Step 4 (the linear family and symmetric spaces).} If $\Lambda^k_\pi T$ were analytic at a point,
$T$ would have a lift (prose after Definition~\ref{fam2:def:good}), whose restriction to $c_{00}(\N,K)^k$ Steps~1--3
exclude. For a symmetric sequence space $G$ and $y\in c_{00}(\N,K)$ with entries
in $\{0,\pm1\}$, $\norm y_G=\norm{1_{\operatorname{supp}y}}_G=\varphi_G(\abs{\operatorname{supp}y})$ by permutation
invariance, and $\varphi_G$ is nondecreasing, so $\varphi^\pm_G=\varphi_G$.
\end{proof}

\begin{proof}[Proof of Corollary~\ref{fam2:cor:table}]
\emph{Step 1 (item~(1)).} For $r_1<2$, $P$ has a lower $r_1$-estimate (example after Theorem~\ref{fib:thm:lower}); apply
Theorem~\ref{fam2:thm:polarizable}(3). For $E=\ell^{r_2}(\N,K)$, $r_2<2$, $(E,\widehat E)$ is good for Banach values
(Theorem~\ref{fib:thm:lower}(3)); apply Proposition~\ref{fam2:prop:fibers}. For $r_1=2$ apply
Theorem~\ref{fam2:thm:gamma-split}(3) with the unit vectors of $\ell^2(\N,K)$, and with $r=r_2$ and the $r_2$-norm
$\norm\cdot_{r_2}$ of $\ell^{r_2}(\N,K)$ over the nonarchimedean $K$ (paragraph after
Corollary~\ref{spl:cor:relative-ordering-na}), or with any $r\in(2,\infty)$ and an equivalent nonarchimedean norm of $E$.

\emph{Step 2 (item~(2)).} For $E=\ell^{r_2}(\N,K)$ or $c_0(\N,K)$, $\norm{\mu_a}_{E\to E}\le\norm a_\infty\le\norm a_P$,
so $a\mapsto\mu_a$ extends from $c_{00}(\N,K)$ to a linear family $T\colon P\to\cL(E,\widehat E)$, and analyticity of
$\Lambda^k_\pi T$ at a point would give a lift (prose after Definition~\ref{fam2:def:good}). For $P=c_0(\N,K)$ and $E$
equal to $\ell^{r_2}(\N,K)$, $r_2\ge2$, or $c_0(\N,K)$, its restriction to $c_{00}(\N,K)^k$ is bounded for the supremum
norm, which Theorem~\ref{fib:thm:ell2} excludes. For $2<r_1<\infty$ and $E$ equal to $c_0(\N,K)$ or $\ell^{r_2}(\N,K)$ with
$1/r_1+1/r_2<1/2$, Theorem~\ref{fam2:thm:staircase} applies with $\varphi_{\ell^{r_1}}(n)=n^{1/r_1}$ and $\varphi_E(n)=1$
or $n^{1/r_2}$, as $(2m)^{2/r_1}(4m)^{2/r_2}/m\to0$. If $K$ is complete and $E$ has an equivalent nonarchimedean norm,
Lemma~\ref{fib:lem:marking} gives bounded linear $\widehat\psi_E\colon c_0(\N,K)\to E$ and $\theta_E\colon E\to c_0(\N,K)$
with $\theta_E\widehat\psi_E=\id$, and $T_E(a)=\widehat\psi_E\mu_a\theta_E$ is a linear family on $P$; a lift $Q_{T_E}$ of it
gives the lift $\hLambda^k_\pi\theta_E\circ Q_{T_E}(a_1,\dots,a_k)\circ\Lambda^k_\pi\widehat\psi_E$ of $a\mapsto
\Lambda^k_\pi\mu_a$ on $c_0(\N,K)$, which the case $E=c_0(\N,K)$ excludes.

\emph{Step 3 (item~(3)).} As $\norm{\mu_a}_{E\to E}=\norm a_\infty$ for $a\in c_{00}(\N,K)$
(Lemma~\ref{obs:lem:multiplier}(1),(3)), the argument of Step~1 of the proof of Theorem~\ref{fam2:thm:bases}, with $(E,E)$
for $(E_0,P)$ and with $a=1_{\operatorname{supp}x}$, $b=1_{\operatorname{supp}y}$ in~\eqref{obs:eq:ptb}, shows that a lift in
a degree $k\ge p$ bounds $\cP_<$ on $E\otimes_\pi E$; Theorem~\ref{cls:thm:symmetric}, (3)$\Rightarrow$(7), then makes
$(E,E)$ and $(E,\widehat E)$ good for Banach values in every degree. Conversely, if $(E,E)$ is
good for Banach values in a degree $k\ge p$, (1)$\Rightarrow$(7) of that theorem gives the same; then
Proposition~\ref{fam2:prop:fibers}(1) makes $(c_0(\N,K);E)$ good in every degree, and the prose after
Definition~\ref{fam2:def:good} makes the multiplier family liftable in every degree.
\end{proof}

\subsection{Tangent bundles and structure groups}\label{fmo:ssec:tangent}

\begin{proof}[Proof of Theorem~\ref{fam2:thm:tangent}]
\emph{Step 1 (item~(1), the shear chart).} The factors $P_1$, $D$, $E$ of $P_0$ are Banach spaces. As $(x,e)\mapsto u(x)e$
composes $(x,e)\mapsto(u(x),e)$ with the bounded bilinear evaluation, $\psi$ and $(x,d,e)\mapsto(x,d-u(x)e,e)$ are
analytic (Lemma~\ref{fmo:lem:series}(2)); they map $U\times D\times E$ to itself, are mutually inverse, and have analytic
derivatives (paragraph after Definition~\ref{fam2:def:tangent}). By the product rule $D\psi(x,d,e)(h,b,c)=(h,b+(Du(x)h)e+
u(x)c,c)$, so $\pi_DD\psi(x,d,e)i_E=u(x)$, and Proposition~\ref{prop:exterior} gives the identity. Analyticity of
$\Lambda^k_\pi D\psi$ at $(x_0,d_0,e_0)$, composed with $x\mapsto(x,d_0,e_0)$ and $T\mapsto\hLambda^k_\pi\pi_D\circ T\circ
\Lambda^k_\pi i_E$, would make $\Lambda^k_\pi u$ analytic at $x_0$; for $\psi^{-1}$ use $\pi_DD(\psi^{-1})i_E=-u$.

\emph{Step 2 (item~(1), the manifold $M$).} The chart change from $\id_M$ to $J\circ\psi\circ J^{-1}$ is
$J\circ\psi\circ J^{-1}$, with derivative $J\circ D\psi(J^{-1}y)\circ J^{-1}$ at $y\in M$; so the transitions of
$\hLambda^k_\pi TM$ are $y\mapsto\hLambda^k_\pi J\circ\Lambda^k_\pi D\psi(J^{-1}y)\circ\Lambda^k_\pi J^{-1}$ and the same
expression with $\psi^{-1}$ (Proposition~\ref{prop:exterior}). They are analytic exactly where $\Lambda^k_\pi D\psi$ and
$\Lambda^k_\pi D(\psi^{-1})$ are, so not at the points $J(x_0,d,e)$, by Step~1. For $\Alt^k(TM;\hLambda^k_\pi D)$ use
$A^k_{E,D;F}(u)=A^k(i_E)\circ A^k(D\psi)\circ A^k(\pi_D)$, $A^k(J\circ D\psi\circ J^{-1})=A^k(J^{-1})\circ A^k(D\psi)\circ
A^k(J)$ and the single target $F=\hLambda^k_\pi D$ of Theorem~\ref{thm:duality}(1).

\emph{Step 3 (item~(2)).} If $(P_0;P_0)$ is good, then for $f$ as in Definition~\ref{fam2:def:tangent}, $Df$ is an
analytic $\cL(P_0)$-valued family, so $P_0$ is tangent-good; tangent goodness gives the bundle statement (sentence after
Definition~\ref{fam2:def:tangent}); if the bundles are analytic, (1) with $P_1=D=E=P_0$ and $P_0\cong P_0^3$ makes
$\Lambda^k_\pi u$ analytic for every analytic $u$ into $\cL(P_0)$.

\emph{Step 4 (item~(3)).} For $E=c_0(\N,K)$ and $E=\ell^r(\N,K)$, $4<r<\infty$, $\norm{a\cdot e}_E\le\norm a_\infty\norm e_E$,
so $(a,e)\mapsto a\cdot e$ is bounded bilinear, $D\psi(a,d,e)(h,b,c)=(h,b+h\cdot e+a\cdot c,c)$ and, with $i_3(c)=(0,0,c)$
and $\pi_2(h,b,c)=b$, $\pi_2D(\psi^{\pm1})(a,d,e)i_3=\pm\mu_a$. As in Step~1, analyticity at a point would make
$a\mapsto\Lambda^k_\pi\mu_a$ analytic at a point of $E$, which Theorem~\ref{fam2:thm:staircase} with $P=E$ excludes, as
$\varphi_E(n)$ is $1$ or $n^{1/r}$ and $(2m)^{2/r}(4m)^{2/r}/m\to0$; interleaving coordinates gives $E\cong E^3$. If $K$ is
complete and $P_0$ is an infinite-dimensional Banach space with an equivalent nonarchimedean norm, Lemma~\ref{fib:lem:marking}
gives $\widehat\psi_{P_0}$, $\theta_{P_0}$ with $\theta_{P_0}\widehat\psi_{P_0}=\id$, so $P_0\cong c_0\oplus P_2$ with
$c_0=c_0(\N,K)$ and $P_2=\ker\theta_{P_0}$, and $P_0\cong c_0\times c_0\times(c_0\oplus P_2)$; the chart
$(a,d,e,z)\mapsto(a,d+a\cdot e,e,z)$ again has $\pi_2D\psi i_3=\mu_a$.
\end{proof}

\begin{proof}[Proof of Theorem~\ref{fam2:thm:groups}]
\emph{Step 1 (item~(1)).} Let $g$ be an $\mathsf H$-family and $x_0\in U$. Left multiplication by $g(x_0)$ is a
homeomorphism of $\mathsf H$, so $g(U')\subseteq g(x_0)\mathsf H_0$ for an open $U'\ni x_0$; then
$x\mapsto g(x_0)^{-1}g(x)$ is a $\mathsf V_0$-family on $U'$, and $\Lambda^k_\pi g(x)=\hLambda^k_\pi g(x_0)\circ
\Lambda^k_\pi(g(x_0)^{-1}g(x))$ (Proposition~\ref{prop:exterior}). Conversely, for a $\mathsf V_0$-family $u$ and $x_0\in U$
choose an open $U'\ni x_0$ and $N\ge0$ with $\abs\vartheta^N\norm u<\varepsilon$ on $U'$, and $\lambda_i=\vartheta^{N+i}$,
$0\le i\le k$. Each $1+\lambda_iu$ is an $\mathsf H$-family on $U'$, with values in $1+\mathsf V_0(\varepsilon)$, and
Lemma~\ref{fmo:lem:vandermonde} with $G=\widehat E$, $a$ the inclusion of $E$ and $b=u(x)$ gives
$\Lambda^k_\pi u=\sum_i\chi_{ki}\Lambda^k_\pi(1+\lambda_iu)$. For Banach $E$ take $\mathsf H=\mathsf H_0=\operatorname{GL}(E)$,
$\mathsf V_0=\cL(E)$ and $\varepsilon=1$ (Neumann series).

\emph{Step 2 (item~(2)).} For $\norm w<1$, $\norm{(1+w)x}=\norm x$ by the nonarchimedean inequality and $1+w$ is
invertible; apply~(1) with $\mathsf H=\mathsf H_0=\operatorname{Isom}(E)$, $\mathsf V_0=\cL(E)$, $\varepsilon=1$.

\emph{Step 3 (item~(3)).} As $\sigma_\sharp\mu_a\sigma'_\sharp\mu_b=(\sigma\sigma')_\sharp\mu_{(a\circ\sigma')b}$ and
$(\sigma_\sharp\mu_a)^{-1}=\mu_{1/a}\sigma^{-1}_\sharp$, $\mathsf H_{\mathrm m}$ is a group. If $\sigma(t)\ne t$, the $t$th
coordinate of $(\sigma_\sharp\mu_a-1)e_t$ is $-1$, so $\norm{\sigma_\sharp\mu_a-1}\ge1$; hence $\{h\in\mathsf H_{\mathrm m}:
\norm{h-1}<1\}\subseteq\mathsf H_{\mathrm d}$, and its translates show that $\mathsf H_{\mathrm d}$ is open in
$\mathsf H_{\mathrm m}$. As $\norm{\mu_a}=\norm a_\infty$, the nonarchimedean inequality gives $1+\mathsf V_{\mathrm d}(1)
\subseteq\mathsf H_{\mathrm d}$; apply~(1) with $\mathsf H_0=\mathsf H_{\mathrm d}$, $\mathsf V_0=\mathsf V_{\mathrm d}$,
$\varepsilon=1$, and $\mathsf H$ equal to $\mathsf H_{\mathrm m}$ or $\mathsf H_{\mathrm d}$.

\emph{Step 4 (item~(4)).} By~(3) it suffices that $\Lambda^k_\pi\mu_x$ is analytic at no point for the
$\mathsf V_{\mathrm d}$-family $x\mapsto\mu_x$ on $c_0(\N,K)$, by Proposition~\ref{prop:homogeneous-lift} and
Theorem~\ref{fib:thm:ell2} with $r=\infty$.
\end{proof}

\begin{proof}[Proof of Proposition~\ref{fam2:prop:span}]
\emph{Step 1 ($\mathsf H$ is a discrete group).} As $g(u)g(u')=g(u+u')$, $\mathsf H$ is a group. For $u\ne u'$ in
$\mathsf U$ some $e_i^*((u-u')e_j)$ is a nonzero element of $\kappa'$, so
$\norm{g(u)-g(u')}\ge\norm{u-u'}\ge\abs{e_i^*((u-u')e_j)}=1$; hence $\mathsf H$ is discrete, and $\mathsf H$-families and
their exterior powers are locally constant.

\emph{Step 2 ($g(\mu_x)$ lies in the closed span).} For $x\in c_0(\N,K)$, the expansions of the coordinates give
$x=\sum_{n\ge n_0}\vartheta^nx_{[n]}$, $\abs\vartheta^{n_0}\ge\norm x_\infty$, with $x_{[n]}\in\kappa'^{\N}$ finitely supported
as $x(t)\to0$; so $\mu_{x_{[n]}}\in\mathsf U$, $\norm{\vartheta^n\mu_{x_{[n]}}}\le\abs\vartheta^n$, and $g(\mu_x)-1=
\sum_n\vartheta^n(g(\mu_{x_{[n]}})-1)$ lies in the closed span of $\mathsf H$.

\emph{Step 3 (no point of analyticity).} By~\eqref{fmo:eq:shear},
$\Lambda^k_\pi\mu_x=\hLambda^k_\pi\pi_D\circ\Lambda^k_\pi g(\mu_x)\circ\Lambda^k_\pi i_E$, so analyticity of
$x\mapsto\Lambda^k_\pi g(\mu_x)$ at a point is excluded by Step~4 of the proof of Theorem~\ref{fam2:thm:groups}.
\end{proof}

\begin{proof}[Proof of Proposition~\ref{fam2:prop:char-two}]
\emph{Step 1 (item~(1)).} Suppose that $Q$ is as in Theorem~\ref{fam2:thm:staircase}, with
$\norm{Q(a_1,\dots,a_k)}\le\norm Q\prod_i\norm{a_i}_P$. Steps~2 and~3 of the proof of Theorem~\ref{obs:thm:char-two}
apply to $Q$ with $\norm{a_i}_P$ in place of $\norm{a_i}_\infty$. For $0/1$ vectors $a,b,x,y$ supported in
$\{1,\dots,n\}$, $\norm a_P,\norm b_P\le\varphi^\pm_P(n)$, $\norm{1_{>n}}_P\le\varphi^\pm_P(k-2)$,
$\norm x_E,\norm y_E\le\varphi^\pm_E(n)$ and $\norm{e_t}_E\le\varphi^\pm_E(1)$. So for every even $n\ge2$
\[
 n\le k(k-1)\,c_k\,\norm Q\,\varphi^\pm_P(n)^2\varphi^\pm_E(n)^2,\qquad
 c_k=\varphi^\pm_P(k-2)^{k-2}\varphi^\pm_E(1)^{k-2},
\]
with $c_2=1$. As $\varphi^\pm_P$ and $\varphi^\pm_E$ are nondecreasing, $n/(\varphi^\pm_P(n)^2\varphi^\pm_E(n)^2)$ is
unbounded along even $n$, which is a contradiction. The assertion on $T$ follows as in Step~4 of the proof of
Theorem~\ref{fam2:thm:staircase}.

\emph{Step 2 (items~(2)--(4)).} For $n\ge1$, $\varphi^\pm_{c_0(\N,K)}(n)=1$ and $\varphi^\pm_{\ell^r(\N,K)}(n)=n^{1/r}$.
Item~(2) follows from item~(1) as in Step~2 of the proof of Corollary~\ref{fam2:cor:table}, since
$n^{1-2/r_1-2/r_2}\to\infty$ when $1/r_1+1/r_2<1/2$, with $1/r_1=0$ for $P=c_0(\N,K)$ and $1/r_2=0$ for $E=c_0(\N,K)$.
Item~(3) follows as in Step~4 of the proof of Theorem~\ref{fam2:thm:tangent}, with item~(1) for $P=E$ in place of
Theorem~\ref{fam2:thm:staircase}, since $n^{1-4/r}\to\infty$ for $r>4$. Item~(4) follows as in Step~4 of the proof of
Theorem~\ref{fam2:thm:groups}, with item~(1) for $P=E=c_0(\N,K)$ in place of Theorem~\ref{fib:thm:ell2}.
\end{proof}
%% ---- end app-O-families.tex
%% ---- begin app-Q-head.tex
\section{Deferred proofs}\label{app:deferred}

This appendix contains the proofs to which Sections~\ref{sec:duality}--\ref{sec:geometry} refer, in the order of
their statements; the notation is that of the corresponding section.
%% ---- end app-Q-head.tex
%% ---- begin app-Q-sec03.tex
\subsection{Proofs for Section~\ref{sec:duality}}\label{dfr:ssec:duality}

The notation is that of Section~\ref{sec:duality}.

\begin{proof}[Proof of Proposition~\ref{fib:prop:split}]
The maps $Q$ are bounded and $k$\nobreakdash-linear, and naturality gives
$Q(u,\dots,u)=\Lambda^k_\pi u\circ\widehat\wedge^k_E\circ s=\Lambda^k_\pi u$ in~(1) and
$Q(u,\dots,u)=\rho\circ\operatorname{alt}^k_D\circ\Lambda^k_\pi u=\Lambda^k_\pi u$ in~(2); Lemma~\ref{fib:lem:test}
proves (1)--(3). For~(4), given $s$, send $m\in\Mult^k(E;F)$ to the alternating map $m^\otimes\circ s\circ\omega_E$. If
$m=m'\circ\omega_E$ with $m'\in\cL(\Lambda^k_\pi E,F)$, then $m^\otimes=m'\circ\wedge^k_E$ and $m^\otimes\circ s=m'$; so
this is a retraction, and it commutes with postcomposition since $(g\circ m)^\otimes=g\circ m^\otimes$. Conversely, the
retraction for $F=\widehat\otimes^k_\pi E$ sends the canonical $k$\nobreakdash-linear map to $s\circ\omega_E$ for some
$s\in\cL(\Lambda^k_\pi E,\widehat\otimes^k_\pi E)$. Postcomposition by $\widehat\wedge^k_E$ turns the canonical map into
the alternating map $\omega_E$, so $\widehat\wedge^k_E\circ s\circ\omega_E=\omega_E$ and $\widehat\wedge^k_E\circ s=\id$. Without completion use
$F=\otimes^k_\pi E$.
\end{proof}

\begin{proof}[Proof of Lemma~\ref{ob:lem:split-criteria}]
(1) Given $\tilde\alpha_1$, $\alpha_2-i'\tilde\alpha_1$ vanishes on $\ker q$, so
$\tilde\alpha_3(qy)=(\alpha_2-i'\tilde\alpha_1)y$ defines a linear map, bounded because $q$ is open, with
$q'\tilde\alpha_3q=q'\alpha_2=\alpha_3q$. Given $\tilde\alpha_3$, $q'(\alpha_2-\tilde\alpha_3q)=0$, so
$\alpha_2-\tilde\alpha_3q=i'\tilde\alpha_1$ for a linear $\tilde\alpha_1$, bounded because $i'$ is bounded below, and
$i'\tilde\alpha_1i=\alpha_2i=i'\alpha_1$. (2) Apply (1) to $\mathcal Y'=\mathcal Y$ and the identities. (3) The space
$Y_\alpha$ is a Banach space, and $q_\alpha$ is well defined, surjective and open, with kernel
$\{[(y',iy)]\}=\{[(y'+\alpha y,0)]\}=i_\alpha(Y')$; and $\norm{y'}\le\norm{y'+\alpha y}+C\norm\alpha\norm{iy}$ for all $y$
shows that $i_\alpha$ is injective and bounded below. An extension $\tilde\alpha$ gives the left inverse
$[(y',y)]\mapsto y'+\tilde\alpha y$ of $i_\alpha$, and a left inverse $\lambda$ gives the extension $y\mapsto\lambda[(0,y)]$, since $[(0,iy)]=[(\alpha y,0)]$. Now apply (2).
\end{proof}

\begin{proof}[Proof of Lemma~\ref{ob:lem:metric-surjection}]
Since $\wedge^k_E=\omega_E^\otimes$ is contractive, it suffices to find for $y$ and $\delta>0$ a preimage $z$ with
$\norm z<\norm y+3\delta$. For $y\in\Lambda^k_\pi E$ take $z=\sum_jx_{j1}\otimes\dots\otimes x_{jk}$ for a representation
$y=\sum_jx_{j1}\wedge\dots\wedge x_{jk}$ with $\sum_j\prod_a\norm{x_{ja}}<\norm y+\delta$. For
$y\in\widehat\Lambda^k_\pi E$ write $y=\sum_{l\ge1}y_l$ with $y_l\in\Lambda^k_\pi E$, $\norm{y_1}<\norm y+\delta$ and
$\sum_{l\ge2}\norm{y_l}<\delta$, and choose preimages $z_l$ with $\norm{z_l}<\norm{y_l}+2^{-l}\delta$; the series
$z=\sum_lz_l$ converges, $\widehat\wedge^k_Ez=y$ and $\norm z<\norm y+3\delta$.
\end{proof}

\begin{proof}[Proof of Lemma~\ref{ob:lem:target}]
(1)$\Rightarrow$(2): a retraction $\rho$, extended to $\widehat\otimes^k_\pi D$, satisfies $\rho\alt^k_D=\id$ on
$\widehat\Lambda^k_\pi D$; so $\norm y\le\norm\rho\norm{\alt^k_Dy}$, $\rho$ is a left inverse of the first map of
$\cR^k_D$, and Lemma~\ref{ob:lem:split-criteria}(2) applies. (2)$\Rightarrow$(1): restrict a bounded left inverse to
$\otimes^k_\pi D$. (1)$\Leftrightarrow$(3): $\alt^k_D\omega_D$ is the antisymmetrization of
$(y_1,\dots,y_k)\mapsto y_1\otimes\dots\otimes y_k$. So $\eta(y_1,\dots,y_k)=\rho(y_1\otimes\dots\otimes y_k)$ gives
$\mathcal A_D\eta=\rho\alt^k_D\omega_D=\omega_D$; conversely $\rho=\eta^\otimes$ satisfies
$\rho\alt^k_D\omega_D=\mathcal A_D\eta=\omega_D$, and the wedges span $\Lambda^k_\pi D$.
\end{proof}

\begin{proof}[Proof of Lemma~\ref{ob:lem:parameter}]
(1)$\Rightarrow$(3): choose a projection $R_n$ onto $\widehat\Gamma^n_\pi P$ with $\norm{R_n}\le g_n(P)+\delta$. The
linearization $\Phi^\otimes\colon\widehat\otimes^n_\pi P\to Z$ of $\Phi$, of norm $\norm\Phi$, maps each $h^{\otimes n}$,
hence $\widehat\Gamma^n_\pi P$, into $W$. So $Q(h_1,\dots,h_n)=\Phi^\otimes R_n(h_1\otimes\dots\otimes h_n)$ takes values in
$W$ and has diagonal $\Phi^\otimes(h^{\otimes n})$. (3)$\Rightarrow$(2): take $W=\widehat\Gamma^n_\pi P$, $Z=\widehat\otimes^n_\pi P$ and
$\Phi(h_1,\dots,h_n)=h_1\otimes\dots\otimes h_n$. (2)$\Rightarrow$(1): the linearization of the $n$-linear map maps
$\widehat\otimes^n_\pi P$ into $\widehat\Gamma^n_\pi P$ and fixes every $h^{\otimes n}$, hence every point of
$\widehat\Gamma^n_\pi P$. For the last assertion, for
$\Phi\in\Mult^n(P;Y)$ the $n$-linear map $(h_1,\dots,h_n)\mapsto\Phi^\otimes R_n(h_1\otimes\dots\otimes h_n)$ has the
diagonal of $\Phi$, and it is $0$ when $\Phi$ has zero diagonal, since $\Phi^\otimes$ then vanishes on
$\widehat\Gamma^n_\pi P$. So it defines a right inverse of the quotient map, of norm at most $\norm{R_n}$, that commutes
with postcomposition.
\end{proof}

\begin{proof}[Proof of Proposition~\ref{ob:prop:split-operators}]
Naturality of $\wedge^k$, extended to completions, gives
$\widehat\wedge^k_Gw^{\otimes k}=\widehat\Lambda^k_\pi w\,\widehat\wedge^k_E$. So $w^{\otimes k}(\cN^k_E)\subseteq\cN^k_G$,
and $w^{\otimes k}|_{\cN^k_E}$, $w^{\otimes k}$, $\widehat\Lambda^k_\pi w$ satisfy the hypotheses of
Lemma~\ref{ob:lem:split-criteria}(1) for $\cS^k_E$ and $\cS^k_G$. Now (a)$\Leftrightarrow$(b) because sections extend by
continuity and restrict, (b)$\Leftrightarrow$(c) is Lemma~\ref{ob:lem:split-criteria}(1), and (c)$\Leftrightarrow$(d)
is Lemma~\ref{ob:lem:split-criteria}(3). In (2), naturality of $\wedge^k$ and $\alt^k$ gives
$\widehat\wedge^k_{G_1}v^{\otimes k}s\,\Lambda^k_\pi v'=\widehat\Lambda^k_\pi v\,\widehat\wedge^k_Gs\,\Lambda^k_\pi v'$ and
$\widehat\Lambda^k_\pi v\,\rho\,v'^{\otimes k}\alt^k_{G_0}=\widehat\Lambda^k_\pi v\,\rho\alt^k_E\Lambda^k_\pi v'$; in
either case the result is $\widehat\Lambda^k_\pi v\,\Lambda^k_\pi w\,\Lambda^k_\pi v'=\Lambda^k_\pi(vwv')$.
\end{proof}

From here on, $q$ and $\overline\ev$ are the maps defined after Lemma~\ref{ob:lem:evaluation}.

\begin{proof}[Proof of Theorem~\ref{ob:thm:mechanisms}]
In each case we give an extension as in Theorem~\ref{ob:thm:obstruction}(2), or a lift.

\emph{Step 1 (source).} By Proposition~\ref{ob:prop:split-operators}(1), $w^{\otimes k}|_{\cN^k_E}$ has a bounded linear
extension $\xi_w\colon\widehat\otimes^k_\pi E\to\cN^k_G$, and $\tau\xi_w$ extends $\ev|_{\cN^k_E}$. For $(\mathrm S^k_E)$
take $w=\id_E$ and $\tau=\ev|_{\cN^k_E}$. For the data $(G,w,\varepsilon)$, the formula for $\tau$ in part~(1) defines a
bounded linear map on $\widehat\otimes^k_\pi G$ with values in $\widehat\Lambda^k_\pi\widehat D=\widehat\Lambda^k_\pi D$ (the
paragraph before Theorem~\ref{ob:thm:mechanisms}) and $\norm\tau\le\norm\varepsilon^k$. Naturality of $\wedge^k$ gives
$\tau(\zeta)(u,\dots,u)=\widehat\Lambda^k_\pi\varepsilon(u)\,\widehat\wedge^k_G\zeta$, which vanishes for
$\zeta\in\cN^k_G$. On pure tensors,
$\tau(w^{\otimes k}(x_1\otimes\dots\otimes x_k))(u_1,\dots,u_k)=\varepsilon(u_1)wx_1\wedge\dots\wedge\varepsilon(u_k)wx_k
=u_1x_1\wedge\dots\wedge u_kx_k$, so $\tau w^{\otimes k}=\ev$ by \eqref{ob:eq:evaluation} and continuity.

\emph{Step 2 (parameter).} Let $\gamma$ be a bounded linear right inverse of $q$. Then
$\xi=\ev-\gamma\overline{\ev}\,\widehat\wedge^k_E$ equals $\ev$ on $\cN^k_E$, and $q\xi=q\ev-\overline\ev\,\widehat\wedge^k_E=0$.
If $(\Gamma^k_{\cL(E,D)})$ holds, $\cM^k(\cL(E,D);\widehat\Lambda^k_\pi D)$ splits by Lemma~\ref{ob:lem:parameter}. In the
general form, the operators vanishing on $\cN^k_E$ form a closed subspace $W$ of
$Z=\cL(\widehat\otimes^k_\pi E,\widehat\Lambda^k_\pi D)$, and $T\mapsto T\widehat\wedge^k_E$ is an isometry of
$\cL(\widehat\Lambda^k_\pi E,\widehat\Lambda^k_\pi D)$ onto $W$, by Lemma~\ref{ob:lem:metric-surjection}.
Lemma~\ref{ob:lem:parameter}(3) gives a bounded $k$-linear
$Q'\colon P^k\to\cL(\widehat\Lambda^k_\pi E,\widehat\Lambda^k_\pi D)$ with $Q'(h,\dots,h)\widehat\wedge^k_E=\Phi(h,\dots,h)$.
Then $Q(u_1,\dots,u_k)=Q'(\psi u_1,\dots,\psi u_k)|_{\Lambda^k_\pi E}$ is bounded and $k$-linear with
$Q(u,\dots,u)=\Lambda^k_\pi u$, and Lemma~\ref{fib:lem:test} applies.

\emph{Step 3 (target).} The formula
$\xi'(z)(u_1,\dots,u_k)=\alt^k_D\bigl(\ev(z)(u_1,\dots,u_k)\bigr)-(u_1\otimes\dots\otimes u_k)\alt^k_E(\widehat\wedge^k_Ez)$
defines a bounded linear $\xi'\colon\widehat\otimes^k_\pi E\to\Mult^k(\cL(E,D);\widehat\otimes^k_\pi D)$ that equals
$\alt_*\ev$ on $\cN^k_E$. Its values have zero diagonal, since
$\xi'(z)(u,\dots,u)=\alt^k_D\widehat\Lambda^k_\pi u\,\widehat\wedge^k_Ez-u^{\otimes k}\alt^k_E\widehat\wedge^k_Ez=0$ by
naturality of $\alt^k$. Lemma~\ref{ob:lem:split-criteria}(3) gives the first assertion. If $\rho$ is a retraction of
$\alt^k_D$, extended to $\widehat\otimes^k_\pi D$, then $\rho\circ\xi'$ extends $\rho\alt^k_D\ev|_{\cN^k_E}=\ev|_{\cN^k_E}$.

\emph{Step 4 (relative target).} Let $\rho'\colon\otimes^k_\pi G\to\widehat\Lambda^k_\pi\widehat D=\widehat\Lambda^k_\pi D$
satisfy $\rho'\alt^k_G=\Lambda^k_\pi w'$, and put
\[
 \xi(z)(u_1,\dots,u_k)=\ev(z)(u_1,\dots,u_k)-\rho'\bigl(\varepsilon(u_1)\otimes\dots\otimes\varepsilon(u_k)\bigr)
 \alt^k_E(\widehat\wedge^k_Ez).
\]
This is a bounded linear map, of norm at most $1+k!\norm{\rho'}\norm\varepsilon^k$, that equals $\ev$ on $\cN^k_E$. By
naturality of $\alt^k$ for $\varepsilon(u)\colon E\to G$,
$\rho'\varepsilon(u)^{\otimes k}\alt^k_E=\rho'\alt^k_G\widehat\Lambda^k_\pi\varepsilon(u)
=\widehat\Lambda^k_\pi w'\,\widehat\Lambda^k_\pi\varepsilon(u)=\widehat\Lambda^k_\pi u$, so $\xi(z)(u,\dots,u)=0$.
\end{proof}

\begin{proof}[Proof of Theorem~\ref{ob:thm:ordering}]
\emph{Step 1 (part~(1)).} By~\eqref{ob:eq:evaluation} and~\eqref{ob:eq:ordering-lift},
$\ev(j_\otimes e^\otimes_t)=Q^\sharp_<(e^\wedge_t)$ for $t\in\Sigma^k_<$, which is the first identity on a basis. Hence
$q\circ Q^\sharp_<=\overline\ev\circ\widehat\wedge^k_E\circ j_\otimes\circ s_<=\overline\ev\circ j_\Lambda\circ\wedge_V\circ s_<
=\overline\ev\circ j_\Lambda$.

\emph{Step 2 (part~(2)).} Since $Q$ is a lift, $q\circ Q^\sharp=\overline\ev$, so
$q\circ L=\overline\ev\circ j_\Lambda-q\circ Q^\sharp_<=0$ by part~(1). Thus $L$ takes values in $\ker q=\Mult^k_0$, and
$\Phi$ is $k$-linear with zero diagonal.

\emph{Step 3 (part~(3)).} Put $T=Q^\sharp_<+L$. Since $\norm{Ty}\le\norm T\norm{j_\Lambda y}$, $j_\Lambda(\Lambda^k_KV)$ is
dense and $\Mult^k$ is complete, $T=\widehat T\circ j_\Lambda$ for a bounded linear $\widehat T\colon\hLambda^k_\pi E\to\Mult^k$
with $\norm{\widehat T}=\norm T$. By part~(1) and $q\circ L=0$, $q\circ\widehat T=\overline\ev$ on a dense subspace, hence
everywhere. So $\widehat T=Q^\sharp$ for a lift $Q$, with $\norm Q=\norm T$, and $Q$ is unique by density. Conversely, a
lift $Q$ gives the correction $L$ of part~(2), with $\norm{Q^\sharp_<+L}=\norm{Q^\sharp\circ j_\Lambda}=\norm Q$ by density.
This proves the formula for $c_k(E,D)$ and the last assertion.

\emph{Step 4 (part~(4)).} Since $q\circ\ev=\overline\ev\circ\widehat\wedge^k_E$, part~(1) gives
$q(\ev(x_1\otimes\dots\otimes x_k))=\overline\ev(j_\Lambda(x_1\wedge\dots\wedge x_k))=q(Q^\sharp_<(x_1\wedge\dots\wedge x_k))$,
so $\Omega_<$ takes values in $\Mult^k_0$. For $x\in V^k$,
\[
 \Omega_<(x)-L(x_1\wedge\dots\wedge x_k)=\ev(x_1\otimes\dots\otimes x_k)-(Q^\sharp_<+L)(x_1\wedge\dots\wedge x_k);
\]
since $\norm\ev\le1$ and $\norm{x_1\wedge\dots\wedge x_k}_\pi\le\prod_a\norm{x_a}$, this gives the second inequality. For the
first, let $\delta>0$ and write $y\in\Lambda^k_KV$ as $\sum_ix_{i1}\wedge\dots\wedge x_{ik}$ with $x_{ia}\in V$ and
$\sum_i\prod_a\norm{x_{ia}}<\norm y_\pi+\delta$. Summing the identity for $(x_{i1},\dots,x_{ik})$ over $i$ gives
$\norm{(Q^\sharp_<+L)y}\le(1+\norm{\Omega_<-L\circ\wedge})(\norm y_\pi+\delta)$; let $\delta\to0$. The alternating $k$-linear
maps $V^k\to\Mult^k_0$ are the maps $L\circ\wedge$, by the universal property of $\Lambda^k_KV$. So part~(3) and the two
inequalities give the remaining assertions.
\end{proof}

\begin{proof}[Proof of Lemma~\ref{ob:lem:repeated}]
\emph{Step 1 (algebra).} Let $I_E\subseteq\otimes^k_KE$ be the span of the pure tensors with $x_a=x_{a+1}$ for some $a<k$.
Modulo $I_E$, exchanging two adjacent entries changes the sign of a pure tensor, since
$x\otimes y+y\otimes x=(x+y)\otimes(x+y)-x\otimes x-y\otimes y$. So $I_E$ contains every pure tensor with two equal entries,
and it is the kernel of the quotient map $\otimes^k_KE\to\Lambda^k_KE$, by the universal property of $\Lambda^k_KE$. The
same holds for $V$.

\emph{Step 2 (closure).} Let $I$ be the image of $I_E$ in $\otimes^k_\pi E$, let $z\in\otimes^k_\pi E$ be the class of
$\tilde z\in\otimes^k_KE$, and let $\delta>0$. Write the image of $\tilde z$ in $\Lambda^k_KE$ as
$\sum_ix_{i1}\wedge\dots\wedge x_{ik}$ with $\sum_i\prod_a\norm{x_{ia}}<\norm{\wedge^k_Ez}+\delta$; by Step~1,
$\tilde z-\sum_ix_{i1}\otimes\dots\otimes x_{ik}\in I_E$, so the distance from $z$ to $I$ is at most $\norm{\wedge^k_Ez}$. If
$z\in\cN^k_E$ and $z_l\in\otimes^k_\pi E$ converge to $z$, then $\norm{\wedge^k_Ez_l}=\norm{\widehat\wedge^k_E(z_l-z)}\to0$, so
$z$ lies in the closure of $I$. As $I\subseteq\cN^k_E$ and $\cN^k_E$ is closed, $\cN^k_E$ is the closure of $I$; with
Step~1 this proves the first assertion.

\emph{Step 3 (entries in $V$).} A pure tensor $x_1\otimes\dots\otimes x_k$ with $x_a=x_{a+1}$ is a limit of pure tensors with
entries in $V$ and equal entries in the slots $a$ and $a+1$, since $V$ is dense and
$(x_1,\dots,x_k)\mapsto x_1\otimes\dots\otimes x_k$ is continuous. With Step~2 this proves the last assertion.
\end{proof}

\begin{proof}[Proof of Corollary~\ref{ob:cor:adjacent}]
By linearization and continuity, the bounded $k$-linear maps $\Psi\colon E^k\to\Mult^k_0$ correspond to the bounded linear
maps $\xi\colon\widehat\otimes^k_\pi E\to\Mult^k_0$, by $\xi(x_1\otimes\dots\otimes x_k)=\Psi(x_1,\dots,x_k)$, with
$\norm\xi=\norm\Psi$. By~\eqref{ob:eq:evaluation} the displayed identity says
$\xi(x_1\otimes\dots\otimes x_k)=\ev(x_1\otimes\dots\otimes x_k)$. If it holds whenever $x_a=x_{a+1}$ for some $a$, then
$\xi=\ev$ on $\cN^k_E$ by Lemma~\ref{ob:lem:repeated} and continuity, so it holds whenever two entries coincide;
conversely, an extension of $\ev|_{\cN^k_E}$ gives the identity on these tuples, whose pure tensors lie in $\cN^k_E$. So the
maps $\Psi$ correspond to the extensions $\xi$ of Theorem~\ref{ob:thm:obstruction}(2), and that theorem with its constants
gives the corollary.
\end{proof}

\begin{proof}[Proof of Corollary~\ref{ob:cor:z-linear}]
By Theorem~\ref{ob:thm:ordering}(1),
$q\circ\ev\circ\nu\circ j_\Lambda=\overline\ev\circ\widehat\wedge^k_E\circ\nu\circ j_\Lambda=\overline\ev\circ j_\Lambda
=q\circ Q^\sharp_<$, so $\Omega^\nu_<$ takes values in $\Mult^k_0$; it is homogeneous because $\nu$ is homogeneous and $\ev$,
$j_\Lambda$, $Q^\sharp_<$ are linear. Since $Q^\sharp_<$ is linear,
$\Omega^\nu_<(\sum_iy_i)-\sum_i\Omega^\nu_<(y_i)=\ev\bigl(\nu(\sum_ij_\Lambda y_i)-\sum_i\nu(j_\Lambda y_i)\bigr)$; the argument
of $\ev$ lies in $\cN^k_E$ and has norm at most $2(1+\delta)\sum_i\norm{y_i}_\pi$ by the bound on $\nu$, and $\norm\ev\le1$.
For a linear $L\colon\Lambda^k_KV\to\Mult^k_0$, $\Omega^\nu_<-L=\ev\circ\nu\circ j_\Lambda-(Q^\sharp_<+L)$ with
$\norm{\ev(\nu(j_\Lambda y))}\le(1+\delta)\norm y_\pi$. So the least constant for $\Omega^\nu_<-L$ and $\norm{Q^\sharp_<+L}$
each exceed the other by at most $1+\delta$, and Theorems~\ref{ob:thm:ordering}(3) and~\ref{ob:thm:obstruction} give the
remaining assertions.
\end{proof}

\begin{proof}[Proof of Theorem~\ref{ob:thm:ultrametric}]
Put $C_\alpha=2(c_+/c_-)\norm\alpha(C_cC_e)^k$. For
$z\in\otimes^k_KV$, $j_\otimes(z-s_<(\wedge_Vz))$ lies in $\cN^k_E$, since $\widehat\wedge^k_E\circ j_\otimes=j_\Lambda\circ\wedge_V$
and $\wedge_V\circ s_<=\id$; let $\alpha_0(z)$ be its image under $\alpha$.

\emph{Step 1 (expansion).} For $t\in\Sigma^k$ with a repeated entry put $\zeta_t=e^\otimes_t$. For $t\in\Sigma^k$ with
distinct entries let $\sigma\in S_k$ be the permutation with $t\circ\sigma=(t_{\sigma(1)},\dots,t_{\sigma(k)})\in\Sigma^k_<$,
and put $\zeta_t=e^\otimes_t-\sgn(\sigma)\,e^\otimes_{t\circ\sigma}$. Since $e_{t_1}\wedge\dots\wedge e_{t_k}$ is
$\sgn(\sigma)\,e^\wedge_{t\circ\sigma}$ for distinct entries and $0$ otherwise, $\wedge_V\zeta_t=0$; and
$\norm{j_\otimes\zeta_t}\le2C_e^k$. For $x_1,\dots,x_k\in V$, expanding $x_l=\sum_tx_l(t)e_t$ gives the finite sum
\[
 x_1\otimes\dots\otimes x_k-s_<(x_1\wedge\dots\wedge x_k)=\sum_{t\in\Sigma^k}\prod_{l=1}^kx_l(t_l)\,\zeta_t.
\]

\emph{Step 2 (bound).} Since $\lvert\!\lvert\!\lvert\cdot\rvert\!\rvert\!\rvert_Y$ is a nonarchimedean norm, Step~1 gives
\[
 \lvert\!\lvert\!\lvert\alpha_0(x_1\otimes\dots\otimes x_k)\rvert\!\rvert\!\rvert_Y\le\max_t\prod_{l=1}^k\abs{x_l(t_l)}\,\lvert\!\lvert\!\lvert\alpha(j_\otimes\zeta_t)\rvert\!\rvert\!\rvert_Y
 \le C_c^k\prod_{l=1}^k\norm{x_l}\cdot\frac{2C_e^k\norm\alpha}{c_-},
\]
the maximum over the finitely many $t$ with $\prod_lx_l(t_l)\ne0$. As $\norm\cdot\le c_+\lvert\!\lvert\!\lvert\cdot\rvert\!\rvert\!\rvert_Y$, this gives
$\norm{\alpha_0(x_1\otimes\dots\otimes x_k)}\le C_\alpha\prod_l\norm{x_l}$, and so
$\norm{\alpha_0(z)}\le C_\alpha\sum_i\prod_l\norm{x_{il}}$ for every representation $z=\sum_ix_{i1}\otimes\dots\otimes x_{ik}$ of
$z\in\otimes^k_KV$ with entries in $V$. By the paragraph before Theorem~\ref{ob:thm:mechanisms},
$\norm{\alpha_0(z)}\le C_\alpha\norm{j_\otimes z}$. Since $j_\otimes(\otimes^k_KV)$ is dense and $Y$ is complete, there is a
bounded linear $\tilde\alpha\colon\widehat\otimes^k_\pi E\to Y$ with $\tilde\alpha\circ j_\otimes=\alpha_0$ and
$\norm{\tilde\alpha}\le C_\alpha$.

\emph{Step 3 (agreement).} If $z\in\otimes^k_KV$ and $\wedge_Vz=0$, then $\tilde\alpha(j_\otimes z)=\alpha_0(z)=\alpha(j_\otimes z)$.
The pure tensors with entries in $V$ and two equal adjacent entries lie in $\ker\wedge_V$, and their images span a dense
subspace of $\cN^k_E$ by Lemma~\ref{ob:lem:repeated}. So $\tilde\alpha=\alpha$ on $\cN^k_E$ by continuity, and
$\alpha_*[\cS^k_E]=0$ by Lemma~\ref{ob:lem:split-criteria}(3).

\emph{Step 4 (the class).} Let $\lvert\!\lvert\!\lvert\cdot\rvert\!\rvert\!\rvert_D$ be a nonarchimedean norm on $\hLambda^k_\pi D$ with
$c_-\lvert\!\lvert\!\lvert\cdot\rvert\!\rvert\!\rvert_D\le\norm\cdot\le c_+\lvert\!\lvert\!\lvert\cdot\rvert\!\rvert\!\rvert_D$. Take $Y=\Mult^k_0$ and, for $\Phi\in Y$, let
$\lvert\!\lvert\!\lvert\Phi\rvert\!\rvert\!\rvert_Y$ be the least $C$ with $\lvert\!\lvert\!\lvert\Phi(u_1,\dots,u_k)\rvert\!\rvert\!\rvert_D\le C\prod_l\norm{u_l}$ for all $u_1,\dots,u_k\in\cL(E,D)$.
Then $\lvert\!\lvert\!\lvert\cdot\rvert\!\rvert\!\rvert_Y$ is a nonarchimedean norm on the Banach space $Y$, since $\lvert\!\lvert\!\lvert\cdot\rvert\!\rvert\!\rvert_D$ is one, and $c_-\lvert\!\lvert\!\lvert\cdot\rvert\!\rvert\!\rvert_Y\le\norm\cdot\le c_+\lvert\!\lvert\!\lvert\cdot\rvert\!\rvert\!\rvert_Y$,
since $\norm\Phi$ is the least $C$ with $\norm{\Phi(u_1,\dots,u_k)}\le C\prod_l\norm{u_l}$. The first part, with this $Y$ and
$\alpha=\ev|_{\cN^k_E}$, gives $\ob_k(E,D)=0$.
\end{proof}
%% ---- end app-Q-sec03.tex
%% ---- begin app-Q-sec04.tex
\subsection{Proofs for Section~\ref{sec:canonical}}\label{dfr:ssec:canonical}

The notation is that of Section~\ref{sec:canonical}.

\begin{proof}[Proof of Theorem~\ref{can:thm:monomial}]
Since $k!=0$ in $K$, $p=\operatorname{char}K$ satisfies $2\le p\le k$. Let $\nu\in S_k$ be a cycle of length $p$ and $g$
the permutation matrix with $ge_i=e_{\nu(i)}$, which is monomial. Corollary~\ref{can:cor:p-cycle}(2) excludes (1)--(3)
already for this $g$. For the algebraic powers, let $\mathsf H$ be the group generated by $\nu$, acting on $I=[k]$, so
that $V=K^k$. As in the proof of Corollary~\ref{can:cor:p-cycle}(2), a map equivariant for $g$ is
$\mathsf H$\nobreakdash-equivariant, and $\bar{\mathsf H}_{[k]}=\bar{\mathsf H}_{\mathbf m}$ has order $p$ for the
$k$\nobreakdash-multiset $\mathbf m$ with value $1$ on $[k]$; so Theorem~\ref{can:thm:permutation}(1)--(3) excludes such
maps.
\end{proof}

\begin{proof}[Proof of Corollary~\ref{can:cor:p-cycle}]
(1) We treat (1) and (2) of Theorem~\ref{can:thm:permutation}; for (3) replace $T$ by $\operatorname{supp}\mathbf m$ and
the stabilizers of $T$ by those of $\mathbf m$. The first statement implies the second by restriction. If $h\in\mathsf H$
maps a $k$\nobreakdash-subset $T$ onto itself and induces on it a permutation of order $p=\operatorname{char}K$, then $p$
divides $\abs{\bar C_T}$ for the group $C$ generated by $h$, and Theorem~\ref{can:thm:permutation} excludes a
$C$\nobreakdash-equivariant map; so the second statement implies the third. If the first fails,
Theorem~\ref{can:thm:permutation} gives a $k$\nobreakdash-subset $T$ such that $p=\operatorname{char}K>0$ divides
$\abs{\bar{\mathsf H}_T}$, and by Cauchy's theorem some element of $\mathsf H_T$ induces on $T$ a permutation of order
$p$; so the third fails. If $\operatorname{char}K=0$, all $\abs{\bar{\mathsf H}_T}$ and $\abs{\bar{\mathsf H}_{\mathbf m}}$
are nonzero in $K$.

(2) Let $\mathsf H$ be the group generated by $\nu$, acting on $I=[k]$ and on $E_0$ by $he_i=e_{h(i)}$; these operators
are bounded since the $e^*_i$ are, and satisfy~\eqref{can:eq:compatible} since
$e^*_i(he_j)=\delta_{i,h(j)}=e^*_{h^{-1}(i)}(e_j)$. For $T=I$ and the $k$\nobreakdash-multiset $\mathbf m$ with value $1$
on $I$, $\mathsf H_T=\mathsf H_{\mathbf m}=\mathsf H$, and $\bar{\mathsf H}_T=\bar{\mathsf H}_{\mathbf m}$ has order $p$,
which is $0$ in $K$. An equivariance identity for $g$ holds for its powers, since $(g^j)^{\otimes k}=(g^{\otimes k})^j$
and $\Lambda^k_\pi(g^j)=(\Lambda^k_\pi g)^j$, hence for all of $\mathsf H$, which Theorem~\ref{can:thm:permutation}(4)
excludes.
\end{proof}

\begin{proof}[Proof of Proposition~\ref{can:prop:asymmetry}]
(1) This is Proposition~\ref{prop:diag}(1) applied to $Q$. With $u_S:=\sum_{i\in S}u_i$, every $k$\nobreakdash-linear
$\Phi$ satisfies $\sum_{S\subseteq[k]}(-1)^{k-\abs S}\Phi(u_S,\dots,u_S)=\sum_\sigma\Phi(u_{\sigma(1)},\dots,u_{\sigma(k)})$,
as in~\eqref{can:eq:polarization}. Apply this to $Q(\cdot)\circ\omega_E$ and to
$(u_1,\dots,u_k)\mapsto\bigl((x_b)\mapsto u_1x_1\wedge\dots\wedge u_kx_k\bigr)$, which agree on the diagonal since
$Q(u,\dots,u)=\Lambda^k_\pi u$: the right sides give
$\sum_\sigma Q(u_{\sigma(1)},\dots,u_{\sigma(k)})\circ\omega_E=\Lsym(u_1,\dots,u_k)\circ\omega_E$, and pure wedges span
$\Lambda^k_\pi E$.
(2) A symmetric $Q$ would give $\Lsym=k!\,Q=0$ by~(1). Let $f_1,\dots,f_k$ be the functionals, choose
$x_1,\dots,x_k\in E$ with $f_i(x_j)=\delta_{ij}$ and, since $\Lambda^k_\pi D$ is dense in $\widehat\Lambda^k_\pi D\ne0$,
$d_1,\dots,d_k\in D$ with $d_1\wedge\dots\wedge d_k\ne0$. For $u_i=f_i(\cdot)\,d_i$ only $\sigma=\id$ contributes
to~\eqref{can:eq:sym}, so $\Lsym(u_1,\dots,u_k)(x_1\wedge\dots\wedge x_k)=d_1\wedge\dots\wedge d_k\ne0$.
(3) By~\eqref{can:eq:signs},
$\wedge^k_D=\rho\circ\operatorname{alt}^k_D\circ\wedge^k_D=\sum_\sigma\sgn(\sigma)\rho\circ\tau_\sigma=k!\,\rho=0$, so
$\Lambda^k_\pi D=\wedge^k_D(\otimes^k_\pi D)=0$.
(4) For $w\in\Lambda^k_\pi E$, \eqref{can:eq:signs} gives
$\operatorname{alt}^k_E(w)=\operatorname{alt}^k_E\widehat\wedge^k_Es(w)=\sum_\sigma\sgn(\sigma)\tau_\sigma s(w)=k!\,s(w)=0$.
With $f_i$ and $x_j$ as in~(2), the bounded functional on $\widehat\otimes^k_\pi E$ linearizing
$(y_a)\mapsto\prod_af_a(y_a)$ takes the value $\sum_\sigma\sgn(\sigma)\prod_af_a(x_{\sigma(a)})=1$ at
$\operatorname{alt}^k_E(x_1\wedge\dots\wedge x_k)$, a contradiction.
\end{proof}
%% ---- end app-Q-sec04.tex
%% ---- begin app-Q-sec05.tex
\subsection{Proofs for Section~\ref{sec:positive}}\label{dfr:ssec:positive}

\begin{proof}[Proof of Theorem~\ref{thm:spherical}]
\emph{Step 1 (simultaneous extension).} Let $Y$ be a nonarchimedean seminormed space, $G$ a nonarchimedean
spherically complete space, $T_\alpha\colon Y\to G$ a nonempty family of bounded operators and $r_\alpha\ge0$ radii with
$\norm{T_\alpha-T_\beta}\le\max(r_\alpha,r_\beta)$, and $T_0$ a linear map on a subspace $Y_0$ with
$\norm{T_0d-T_\alpha d}\le r_\alpha\norm d$ for all $d\in Y_0$ and all $\alpha$. We extend $T_0$ to $Y$ with all
these bounds. For $x\notin Y_0$, the balls
\[
 \overline B\bigl(T_\alpha(d+x)-T_0d,\ r_\alpha\norm{d+x}\bigr),
 \qquad \alpha,\ d\in Y_0,
\]
meet pairwise: for $r_\alpha\le r_\beta$, the centers for $(\alpha,d)$ and $(\beta,e)$ differ by at most
$\max(r_\alpha\norm{d+x},r_\beta\norm{e+x})$, by the hypotheses and the nonarchimedean inequality. A common point, given
by spherical completeness, is a value at $x$ extending $T_0$ to $Y_0+Kx$ with all bounds, and Zorn's lemma extends $T_0$
to $Y$. With the single operator $0$ and radius one, contractions extend to contractions: this is the extension
mechanism of Ingleton's theorem~\cite{In}.

\emph{Step 2 ($\Alt^k(E;F)$ is spherically complete).} Let $Y$ be the free vector space on the symbols
$[x_1,\ldots,x_k]$, $x\in E^k$, with the seminorm $\norm{\sum_x c_x[x]}=\max_x\abs{c_x}\prod_i\norm{x_i}$, and let $Y_0$
be the span of the multilinearity and repeated-input relations. Alternating maps $E^k\to F$ correspond isometrically to
bounded linear maps $Y\to F$ vanishing on $Y_0$: one direction uses the nonarchimedean inequality in $F$, the other
evaluates on generators. Pairwise-meeting balls $\overline B(m_\alpha,r_\alpha)$ in $\Alt^k(E;F)$ satisfy
$\norm{m_\alpha-m_\beta}\le\max(r_\alpha,r_\beta)$, since this operator norm is nonarchimedean. So the maps $T_\alpha$
corresponding to the $m_\alpha$, and $T_0=0$ on $Y_0$, satisfy the hypotheses of Step~1 with $G=F$; the extension of
$T_0$ vanishes on $Y_0$, and its alternating map lies in every ball.

\emph{Step 3 (the retraction).} The operator norm of $\Mult^k(E;F)$ is nonarchimedean. By Step~1 with
$G=\Alt^k(E;F)$, spherically complete by Step~2, the inverse of the isometric inclusion $j_{E,F}$ extends to a contraction
of $\Mult^k(E;F)$ onto $\Alt^k(E;F)$, which is the required retraction. Proposition~\ref{prop:split-lift} finishes the
proof. In degree zero the alternating space is $F$.
\end{proof}

\begin{proof}[Proof of Theorem~\ref{fib:thm:ultra}]
\emph{Step 1 (ultrametric envelopes).} As in the proof of Proposition~\ref{prop:scalar-countable}, with the
nonarchimedean inequality of $F$ in place of that of $K$, every bounded $k$\nobreakdash-linear map $m$ into $F$ satisfies
$\norm{m(x_1,\dots,x_k)}\le\norm m\prod_aq_Y(x_a)$. Hence $m\mapsto m\circ j_Y^{\times k}$ is an isometry of
$\Alt^k(Y^u;F)$ onto $\Alt^k(Y;F)$, every $u\in\cL(E,D)$ induces $u^u\in\cL(E^u,D^u)$ with $\norm{u^u}\le\norm u$,
linearly in $u$, and $A^k_{E,D;F}(u)$ corresponds to $A^k_{E^u,D^u;F}(u^u)$. A lift for $(E^u,D^u)$ therefore gives one
for $(E,D)$ of no larger norm. Since $(Y^u)^u=Y^u$, we may assume $E=E^u$ and $D=D^u$.

\emph{Step 2 (the expansion on $c_0(I,K)$).} Order $I$, and let $(e_i)$ and $(e^*_i)$ be the unit vectors and
coordinate functionals of $c_0(I,K)$. For $m\in\Alt^k(c_0(I,K);F)$ and $x_1,\dots,x_k\in c_0(I,K)$, the family
$\det(e^*_{i_a}(x_b))_{a,b}\,m(e_{i_1},\dots,e_{i_k})$, $i_1<\dots<i_k$, is null: if a determinant has absolute value at
least $\delta>0$, every $i_a$ lies in the finite set where some $\abs{e^*_i(x_b)}\ge\delta/\max_b\norm{x_b}^{k-1}$. Its
sum is
\begin{equation}\label{fib:eq:c0-expansion}
 m(x_1,\dots,x_k)=\sum_{i_1<\dots<i_k}\det\bigl(e^*_{i_a}(x_b)\bigr)_{a,b}\,m(e_{i_1},\dots,e_{i_k}),
\end{equation}
by sorting the multilinear expansion for finitely supported $x_b$, and by density in general, both sides being bounded
and multilinear.

\emph{Step 3 (lifts).} The maps
\[
 Q(u_1,\dots,u_k)(m)(x)=\sum_{i_1<\dots<i_k}\det\bigl(e^*_{i_a}(x_b)\bigr)_{a,b}\,m(u_1e_{i_1},\dots,u_ke_{i_k}),
\]
for $u_a\in\cL(c_0(I,K),D)$, and
\[
 Q'(u_1,\dots,u_k)(m)(x)=\sum_{i_1<\dots<i_k}\det\bigl(e^*_{i_a}(u_ax_b)\bigr)_{a,b}\,m(e_{i_1},\dots,e_{i_k}),
\]
for $u_a\in\cL(E,c_0(I,K))$, are well defined by Step~2, with $u_ax_b$ in place of $x_b$ in row $a$ for $Q'$. They are
alternating in $x$, of norm at most one, and have diagonal $A^k(u)(m)$ by~\eqref{fib:eq:c0-expansion}. If $Q_0$ lifts
$A^k_{c_0(I,K),D;F}$ and $\iota'\iota=\id_E$, then $Q_0(u_1\iota',\dots,u_k\iota')(m)\circ\iota^{\times k}$ lifts
$A^k_{E,D;F}$; if $Q_0$ lifts $A^k_{E,c_0(I,K);F}$ and $\iota'\iota=\id_D$, then
$Q_0(\iota u_1,\dots,\iota u_k)(m\circ\iota'^{\times k})$ does. Both have norm at most
$(\norm\iota\norm{\iota'})^k\norm{Q_0}$.

\emph{Step 4 (countable type).} If $E$ is of countable type, so is $E^u$, which then has a $t$\nobreakdash-orthogonal
basis for every $0<t<1$ \cite[Proposition~1.4]{VdP}: nonzero vectors $b_i$ with dense span and
$\norm{\sum_i\lambda_ib_i}\ge t\max_i\abs{\lambda_i}\norm{b_i}$ for finitely many scalars $\lambda_i$. The formulas for
$Q$ and~\eqref{fib:eq:c0-expansion}, with this basis and its coordinate functionals in place of $(e_i)$ and $(e^*_i)$,
give a lift of norm at most $t^{-k}$, as in the proof of Proposition~\ref{prop:scalar-countable}; similarly for $D$
with $Q'$.
\end{proof}

\begin{proof}[Proof of Proposition~\ref{spl:prop:projective}]
(1) The map $\Phi\colon\delta_i\mapsto e_{i_1}\wedge\dots\wedge e_{i_k}$ is contractive on finitely supported vectors.
The alternating map $\Omega(x_1,\dots,x_k)=(\det(x_b(i_a))_{a,b})_{i\in J^k_<}$ satisfies
$\norm{\Omega(x_1,\dots,x_k)}_1\le\sum_{i\in J^k_<}\sum_{\sigma\in S_k}\prod_a\abs{x_{\sigma(a)}(i_a)}\le\prod_b\norm{x_b}_1$,
since $(i,\sigma)\mapsto i\circ\sigma^{-1}$ is injective into $J^k$, tuples being maps $\{1,\dots,k\}\to J$. By
Proposition~\ref{prop:exterior}, $\Omega$ induces a contraction
$\widetilde\Omega\colon\Lambda^k_\pi\ell^1(J,K)\to\ell^1(J^k_<,\widehat K)$ with $\widetilde\Omega\Phi\delta_i=\delta_i$,
because $(e_{i_b}(i'_a))_{a,b}$ is the identity for $i'=i$ and has a zero row for every other $i'\in J^k_<$. So $\Phi$ is
isometric. Its range is dense, since wedges of finitely supported vectors span a dense subspace and expand into the
$e_{i_1}\wedge\dots\wedge e_{i_k}$, $i\in J^k_<$. Completion gives~(1).

(2) By Lemma~\ref{ob:lem:metric-surjection}, $\widehat\wedge^k_E\colon\widehat\otimes^k_\pi E\to\widehat\Lambda^k_\pi E$
is a bounded linear surjection of Banach spaces. Projectivity gives a bounded linear $s'$ with
$\widehat\wedge^k_E\circ s'=\id$, whose restriction to $\Lambda^k_\pi E$ satisfies $(\mathrm S^k_E)$.

(3) Order $J$, and let $\iota\colon E\to\ell^1(J,K)$ and $\iota'\colon\ell^1(J,K)\to E$ be bounded linear with
$\iota'\iota=\id_E$. The inverse of the isomorphism of~(1), followed by $\delta_i\mapsto e_{i_1}\otimes\dots\otimes
e_{i_n}$, restricts to a contraction $s_0\colon\Lambda^n_\pi\ell^1(J,K)\to\widehat\otimes^n_\pi\ell^1(J,K)$ with
$\widehat\wedge^n\circ s_0=\id$, by the density shown in~(1). By naturality of $\wedge^n$, the map
$s=\iota'^{\otimes n}\circ s_0\circ\Lambda^n_\pi\iota$ satisfies $\widehat\wedge^n_E\circ s=\widehat\Lambda^n_\pi\iota'
\circ\Lambda^n_\pi\iota=\id$, and $\norm s\le(\norm\iota\norm{\iota'})^n$.
\end{proof}

\begin{proof}[Proof of Corollary~\ref{spl:cor:relative-ordering-na}]
Let $\nu$ be a nonarchimedean norm on $D$ with $a\nu(y)\le\norm y\le b\nu(y)$. Then
$\norm{\sum_t\xi_ty_t}\le b\max_t\abs{\xi_t}\nu(y_t)\le(b/a)(\sum_t\abs{\xi_t}^3)^{1/3}\max_t\norm{y_t}$, so
\eqref{spl:eq:r-norm} holds with $r=3$, $D$ has an equivalent $3$-norm by Proposition~\ref{spl:prop:r-norm}, and
Theorem~\ref{spl:thm:relative-ordering}(2) gives the first assertion. The extension of the coordinate map is the map $w$ of
Theorem~\ref{spl:thm:relative-ordering}(1) for $r=3$ followed by the contractive inclusion
$\ell^3(\Sigma,\widehat K)\to c_0(\Sigma,\widehat K)$, hence $\mathrm S^n$-split by the ideal property in
Proposition~\ref{ob:prop:split-operators}.
\end{proof}

\begin{proof}[Proof of Theorem~\ref{thm:summable-parameters}]
\emph{Step 1 (the space $c_0$).} Order $I$ and group the coordinate coefficients of each ambient $b_n$ into
nondecreasing words. Finite-coordinate uniqueness, as in the proof of Theorem~\ref{thm:finite-parameters}, puts the
grouped coefficients $c_{i_1,\ldots,i_n}$ in $W$, and the nonarchimedean inequality bounds their norms by $\norm{b_n}$.
The formula
\[
 q_n(x_1,\ldots,x_n)
 =\sum_{i_1\le\cdots\le i_n}
 c_{i_1,\ldots,i_n}(x_1)_{i_1}\cdots(x_n)_{i_n}
\]
converges, since outside finite sets of words some coordinate factor is arbitrarily small while the others are bounded;
completeness of $W$ supplies the sum, of norm at most $\norm{b_n}\prod_a\norm{x_a}$. The diagonals of $q_n$ and $b_n$
agree on finite supports, hence by density. Theorem~\ref{thm:coefficient-descent} gives~(1).

\emph{Step 2 (sorting in $\ell^1$).} For every set $J$, the sorting of Step~1 in degree $d$ gives each bounded
$d$-linear $B'\colon\ell^1(J,K)^d\to Z$ whose diagonal lies in $W$ a lift into $W$ with the same diagonal and of norm at
most $d!\norm{B'}$: $W$ is complete, being closed in $Z$, each grouped coefficient has at most $d!$ terms, and the
resulting series is absolutely summable, bounded by $d!\norm{B'}\prod_a\norm{x_a}_1$.

\emph{Step 3 (factorization of $\gamma$).} Fix $x_0\in U$, write $\gamma(x_0+x)=\sum_na_n(x,\ldots,x)$, and choose
$\lambda\in K$ with $0<\abs\lambda$ within the radius of this expansion. Set $J=\coprod_{n\ge0}I^n$ and
\[
 T(e_{n,i_1,\ldots,i_n})=\lambda^na_n(e_{i_1},\ldots,e_{i_n}),
 \qquad
 g(x)_{n,i_1,\ldots,i_n}=\lambda^{-n}x_{i_1}\cdots x_{i_n}.
\]
The values of $T$ on unit vectors are uniformly bounded, so $T$ extends by absolute summation to a bounded linear map
$\ell^1(J,K)\to\widehat H$. The degree-$n$ multilinear coefficient of $g$ has entries
$\lambda^{-n}(x_1)_{i_1}\cdots(x_n)_{i_n}$ in its $n$th block and zero elsewhere, so its norm is at most
$\abs\lambda^{-n}$; its diagonal block has sum norm $(\norm{x}_1/\abs\lambda)^n$. So $g$ has a convergent multilinear
expansion for $\norm{x}_1<\abs\lambda$, without completeness of $K$, and $T(g(x))=\gamma(x_0+x)$ there.

\emph{Step 4 (conclusion of~(2)).} Extend $B$ to $\widehat H^d$ by completeness of $Z$; its diagonal still lies in $W$,
by closedness. By Step~2, the composite of this extension with $T^{\times d}$ has a bounded $d$-linear lift $q$ into $W$
with the same diagonal. Then $ja(\gamma(x_0+x))=jq(g(x),\ldots,g(x))$ for $\norm{x}_1<\abs\lambda$, so near $x_0$ the
map $a\gamma$ is the diagonal of $q$, a bounded polynomial, composed with the analytic map $g$. Hence $a\gamma$ is
analytic at every point $x_0$.

\emph{Step 5 (retracts).} If $ri=\id_P$ for bounded linear $i\colon P\to V$ and $r\colon V\to P$, apply~(1) or~(2) to
the family composed with $r$ on $r^{-1}(U)$, and restrict along $i$; in~(1) this multiplies the norms of the lifted
coefficients by at most $(\norm r\norm i)^n$.
\end{proof}

\begin{proof}[Proof of Proposition~\ref{spl:prop:gamma-l1}]
\emph{Step 1 (orbit sums).} Let $L=\ell^1(J,K)$ with its unit vectors $e_i$, and order $J$. For nondecreasing
$t\in J^n$, let $O_t$ be the sum of the tensors $e_{\tilde t_1}\otimes\dots\otimes e_{\tilde t_n}$ over the distinct
rearrangements $\tilde t$ of $t$, let $b_1,\dots,b_m$ be the distinct entries of $t$, and for $\alpha\in\N^m$ with
$\alpha_1+\dots+\alpha_m=n$ let $t(\alpha)$ be the nondecreasing tuple with $\alpha_i$ entries $b_i$. For $\lambda\in K^m$,
expanding the tensor power gives $(\sum_i\lambda_ie_{b_i})^{\otimes n}=\sum_\alpha\lambda^\alpha O_{t(\alpha)}$. A linear
functional $\varphi$ on the span of the $O_{t(\alpha)}$ that vanishes on these pure powers makes the polynomial function
$\lambda\mapsto\sum_\alpha\lambda^\alpha\varphi(O_{t(\alpha)})$ vanish on $K^m$, so $\varphi(O_{t(\alpha)})=0$ for all
$\alpha$, since $K$ is infinite. Hence $O_t$ lies in the span of the pure powers, and in $\widehat\Gamma^n_\pi L$. This
argument is algebraic: it holds for arbitrary vectors $e_{b_1},\dots,e_{b_m}$ of a vector space over $K$ in place of
the unit vectors.

\emph{Step 2 (the projection).} As in the proof of Proposition~\ref{spl:prop:projective}(1), with the coordinate map
$(x_1,\dots,x_n)\mapsto(x_1(t_1)\cdots x_n(t_n))_{t\in J^n}$ in place of $\Omega$, the map
$e_{t_1}\otimes\dots\otimes e_{t_n}\mapsto\delta_t$ onto the unit vectors $\delta_t$ of $\ell^1(J^n,\widehat K)$ extends
to an isometric isomorphism $\widehat\otimes^n_\pi L\cong\ell^1(J^n,\widehat K)$. Let $R_n\delta_t=O_t$ for
nondecreasing $t$ and $R_n\delta_t=0$ otherwise. Then $\norm{R_n}\le n!$, since $O_t$ is a sum of at most $n!$ unit
vectors, and $R_n$ maps into the closed subspace $\widehat\Gamma^n_\pi L$ by Step~1. For finitely supported $h\in L$, the
coefficients of $h^{\otimes n}=\sum_th(t_1)\cdots h(t_n)\delta_t$ are constant on rearrangements of $t$, so
$R_nh^{\otimes n}=h^{\otimes n}$. By continuity $R_n$ fixes every pure power, hence $\widehat\Gamma^n_\pi L$. This
proves~(1).

\emph{Step 3 (finite coordinates).} A normed space $P$ with finite continuous coordinates has a bounded linear
isomorphism $T$ onto $\ell^1(\{1,\dots,d\},K)$. The extension of $T^{\otimes n}$ to the completions is an isomorphism
mapping $\widehat\Gamma^n_\pi P$ onto $\widehat\Gamma^n_\pi\ell^1(\{1,\dots,d\},K)$, since
$T^{\otimes n}h^{\otimes n}=(Th)^{\otimes n}$. So $(T^{-1})^{\otimes n}\circ R_n\circ T^{\otimes n}$, with $R_n$ from
Step~2 for $J=\{1,\dots,d\}$, is a projection onto $\widehat\Gamma^n_\pi P$.
\end{proof}

\begin{proof}[Proof of Proposition~\ref{spl:prop:degree-two}]
Alternation is vanishing on repeated arguments, so $\Lambda^2_KP$ is the quotient of $P\otimes_KP$ by the span of the
tensors $h\otimes h$. Hence $\widehat\wedge^2_P(h\otimes h)=0$, and $\widehat\Gamma^2_\pi P\subseteq\cN^2_P$ because
$\cN^2_P$ is closed. Conversely, let $z\in\otimes^2_\pi P$ with $\wedge^2_Pz=0$, represented by $\tilde z\in P\otimes_KP$.
The image of $\tilde z$ in $\Lambda^2_KP$ has seminorm zero, so for $\epsilon>0$ it equals $\sum_jx_j\wedge y_j$ with
$\sum_j\norm{x_j}\norm{y_j}<\epsilon$; then $\tilde z-\sum_jx_j\otimes y_j$ lies in the span of the tensors $h\otimes h$,
so $z$ lies within $\epsilon$ of that span, and $z\in\widehat\Gamma^2_\pi P$. For $z\in\cN^2_P$ choose $z_m\in\otimes^2_\pi P$ with $z_m\to z$. Then
$\wedge^2_Pz_m\to0$, Lemma~\ref{ob:lem:metric-surjection} gives $w_m\in\otimes^2_\pi P$ with $\wedge^2_Pw_m=\wedge^2_Pz_m$
and $\norm{w_m}\le2\norm{\wedge^2_Pz_m}$, and $z_m-w_m\in\widehat\Gamma^2_\pi P$ converges to $z$. The norm bounds are
those of Lemma~\ref{ob:lem:source}.
\end{proof}
%% ---- end app-Q-sec05.tex
%% ---- begin app-Q-sec06.tex
\subsection{Proofs for Section~\ref{sec:combining}}\label{dfr:ssec:combining}

The notation and the standing hypotheses are those of Section~\ref{sec:combining}.

\begin{proof}[Proof of Lemma~\ref{fib:lem:hopf}]
Concatenating representations gives $\norm{z\wedge w}\le\norm z\norm w$. The right side of~\eqref{fib:eq:shuffle} is
$n$\nobreakdash-linear in $(x_1,\dots,x_n)$, of norm at most $\binom na$. It is alternating: if $x_i=x_j$ with $i\ne j$,
the terms with $i$ and $j$ in one block vanish, and the others cancel in pairs by exchanging $i$ and $j$ between the
blocks. Proposition~\ref{prop:exterior} gives $\operatorname{sh}^{a,b}_E$; naturality is checked on pure wedges. Each of
the $\binom na$ terms of $\wedge^{a,b}_E\operatorname{sh}^{a,b}_E(x_1\wedge\dots\wedge x_n)$ is
$\epsilon(S)x_S\wedge x_{S^c}=x_1\wedge\dots\wedge x_n$. The factorization of $\wedge^n_E$ holds on pure tensors; that of
$\operatorname{alt}^n_E$ holds because every permutation of $\{1,\dots,n\}$ is uniquely a shuffle of two blocks followed
by permutations within the blocks, and its sign is the product of their signs.
\end{proof}

\begin{proof}[Proof of Theorem~\ref{fib:thm:lucas}]
Identify $\Alt^j(Y;F')$ with $\cL(\Lambda^j_\pi Y,F')$ by Proposition~\ref{prop:exterior}. For $T\in\Alt^n(D;F)$
define $T^\flat\in\Alt^a(D;G)$ by $T^\flat(z)(w)=T(z\wedge w)$; then $\norm{T^\flat}\le\norm T$. Let $Q_a$ and $Q_b$ be
bounded lifts of $A^a_{E,D;G}$ and $A^b_{E,D;F}$ (Proposition~\ref{prop:homogeneous-lift}). For
$u_1,\dots,u_n\in\cL(E,D)$ the bilinear map on $\Lambda^a_\pi E\times\Lambda^b_\pi E$ given by
\[
 (z,w)\longmapsto\bigl[Q_b(u_{a+1},\dots,u_n)\bigl(Q_a(u_1,\dots,u_a)(T^\flat)(z)\bigr)\bigr](w)
\]
is bounded. Composing its linearization with
$\binom na^{-1}\operatorname{sh}^{a,b}_E$ gives $Q_n(u_1,\dots,u_n)(T)\in\Alt^n(E;F)$, bounded and $n$\nobreakdash-linear
in $(u_1,\dots,u_n)$, with values in $\cL(\Alt^n(D;F),\Alt^n(E;F))$. On the diagonal the bilinear map is
$(z,w)\mapsto T(\Lambda^a_\pi u\,z\wedge\Lambda^b_\pi u\,w)$, so naturality and~\eqref{fib:eq:hopf} give
$Q_n(u,\dots,u)(T)=T\circ\Lambda^n_\pi u=A^n_{E,D;F}(u)(T)$. For the last assertion take $a=1$:
$A^1_{E,D;G}(u)(T)=T\circ u$ is linear in $u$.
\end{proof}

\begin{proof}[Proof of Corollary~\ref{fib:cor:degrees}]
By Lucas's theorem, $\binom{a+b}a\ne0$ in $K$ exactly when the base-$p$ addition of $a$ and $b$ has no carry: over
$\F_p$, $(1+X)^{a+b}=\prod_i(1+X^{p^i})^{c_i}$, where the $c_i$ are the base-$p$ digits of $a+b$.

(1) Degrees below $p$ belong to $\fD$ by Proposition~\ref{prop:split-lift}(1). For a carry-free sum apply
Theorem~\ref{fib:thm:lucas}; if $F$ is normed or Banach, so is $G=\Alt^b(D;F)$. For the scalar set,
Theorem~\ref{thm:duality}(2) shows that $A^j_{E,D;K}$ is analytic exactly when $A^j_{E,D;Y^*}$ is analytic for every
normed $Y$; apply Theorem~\ref{fib:thm:lucas} with $F=Y^*$, for which $G=\Alt^b(D;Y^*)$ is the dual of
$\Lambda^b_\pi D\otimes_\pi Y$.

(2) Let $N$ be the least missing element; $N\ge p$ by~(1). If $N$ is not a power of $p$ and $p^j$ is the largest power
of $p$ below $N$, then $N=p^j+(N-p^j)$ is a carry-free sum with $0<N-p^j<N$, so $N\in\fD$, a contradiction. The second
assertion follows.

(3) The pair $(E,D)$ is not good in degree $n_F$; by~(2), the least degree in which it is not good is $p^i$ for some
$i\ge1$, and $n_F\ge p^i$. If $n_F$ had a nonzero digit in a position $j<i$, then $(E,D)$ would be good in degree $p^j$,
$A^{n_F-p^j}_{E,D;F}$ would be analytic, and $n_F=p^j+(n_F-p^j)$ would be carry-free, so
Theorem~\ref{fib:thm:lucas} would make $A^{n_F}_{E,D;F}$ analytic.
\end{proof}

\begin{proof}[Proof of Theorem~\ref{fib:thm:hyperplane}]
Let $y=(0,1)\in E\oplus K$, and let $\operatorname{pr}_D$ and $\operatorname{pr}_K$ be the coordinate projections of
$D\oplus K$ onto $D$ and $K$. Define bounded linear maps $R\colon\Alt^k(E\oplus K;F)\to\Alt^{k-1}(E;F)$, of norm at most
one, and $J\colon\Alt^{k-1}(D;F)\to\Alt^k(D\oplus K;F)$, of norm at most $k$, by
$R(m)(x_1,\dots,x_{k-1})=m((x_1,0),\dots,(x_{k-1},0),y)$ and
\[
 J(m)(z_1,\dots,z_k)=\sum_{i=1}^k(-1)^{k-i}\operatorname{pr}_K(z_i)\,
 m(\operatorname{pr}_Dz_1,\dots,\widehat{\operatorname{pr}_Dz_i},\dots,\operatorname{pr}_Dz_k).
\]
The form $J(m)$ is alternating in every characteristic: if $z_a=z_b$ with $a<b$, only the terms $i=a$ and $i=b$ can be
nonzero, their argument lists differ by moving one entry across $b-a-1$ others, and
$(-1)^{k-a}+(-1)^{k-b}(-1)^{b-a-1}=0$. For $u\in\cL(E,D)$, $(u\oplus1)(x_i,0)=(ux_i,0)$ and $(u\oplus1)y=(0,1)$, so only
the term $i=k$ survives and
\[
 A^{k-1}_{E,D;F}(u)=R\circ A^k_{E\oplus K,D\oplus K;F}(u\oplus1)\circ J.
\]
The right side is analytic in $u$, since $u\mapsto u\oplus1$ is affine and bounded.
\end{proof}

\begin{proof}[Proof of Theorem~\ref{cmb:thm:decomposition}]
\emph{Step 1 (pieces).} Let $n\le k$. For a piece of source type and a section $s_n$ for $w$ in degree $n$, the map
$Q_n(u_1,\dots,u_n)=\widehat\wedge^n_{\widehat D}\circ(w'\varepsilon(u_1)\otimes\dots\otimes w'\varepsilon(u_n))\circ s_n$
has diagonal $\Lambda^n_\pi(w'\varepsilon(u))\circ\Lambda^n_\pi w$, as in the proof of
Proposition~\ref{cmb:prop:relative}(1). For a piece of target type and a retraction $\rho_n$ for $w'$ in degree $n$, the
map $Q_n(u_1,\dots,u_n)=\rho_n\circ(\varepsilon(u_1)w\otimes\dots\otimes\varepsilon(u_n)w)\circ\alt^n_E$ has diagonal
$\Lambda^n_\pi w'\circ\Lambda^n_\pi(\varepsilon(u)w)$, as in the proof of Proposition~\ref{cmb:prop:relative}(2). Both
diagonals equal $\Lambda^n_\pi(w'\varepsilon(u)w)$, which proves~(1).

\emph{Step 2 (two maps).} Let $\theta,\tilde\theta\colon\cL(E,D)\to\cL(E,\widehat D)$ be liftable in every degree
$a\le k$, with lifts $Q_a$ and $\tilde Q_a$, and put $Q_0=\tilde Q_0=\id_K$. For $n\le k$ the map
\[
 Q^+_n(u_1,\dots,u_n)=\sum_{a=0}^n\wedge^{a,n-a}_{\widehat D}\circ\bigl(Q_a(u_1,\dots,u_a)\otimes
 \tilde Q_{n-a}(u_{a+1},\dots,u_n)\bigr)\circ\operatorname{sh}^{a,n-a}_E ,
\]
with $\wedge^{a,n-a}_{\widehat D}$ extended to the completions, is bounded and $n$-linear with diagonal
$\Lambda^n_\pi(\theta(u)+\tilde\theta(u))$, by~\eqref{cmb:eq:binomial} for $G=\widehat D$.

\emph{Step 3 (conclusion).} By Step~2 and induction on $m$, the map $u\mapsto\sum_i\theta_i(u)=u$ is liftable in degree
$k$, and Lemma~\ref{fib:lem:test} gives~(2); under the hypotheses of its last sentence, the formula of Step~2 keeps values
in $\Lambda^n_\pi D$. Item~(3) follows from~(1) and~(2) for decompositions into pieces of source and target type, with
Step~1 for decompositions without completion, and from Proposition~\ref{cmb:prop:relative} for relative splittings.
\end{proof}

\begin{proof}[Proof of Proposition~\ref{cmb:prop:normal-form}]
\emph{Step 1 (the two bilinear maps).} Given a decomposition~\eqref{cmb:eq:decomposition} into pieces of source and
target type, let $b_S$ and $b_R$ be the sums of the bounded bilinear maps $b_i(u,x)=w'_i\varepsilon_i(u)w_ix$ over the
pieces of source type and of target type; then $b_S(u,x)+b_R(u,x)=ux$. For a piece of source type, the map
$\xi_i\colon G_i\to\cL(\cL(E,D),\widehat D)$, $\xi_i(g)(u)=w'_i\varepsilon_i(u)g$, is bounded linear and
$w_{b_i}=\xi_iw_i$. For a piece of target type, the linear map $\lambda_i\colon\cL(E,D)\otimes_\pi E\to G'_i$ with
$\lambda_i(u\otimes x)=\varepsilon_i(u)w_ix$ is bounded and $w'_{b_i}=w'_i\lambda_i$. By the paragraph on split operators
in Section~\ref{cmb:ssec:degrees}, $w_{b_S}=\sum w_{b_i}$ is $\mathrm S^{\le k}$-split and $w'_{b_R}=\sum w'_{b_i}$ is
$\mathrm R^{\le k}$-split.

\emph{Step 2 (the two pieces).} Conversely, given $b_S$ and $b_R$, the piece
$(\cL(\cL(E,D),\widehat D),\widehat D,w_{b_S},\id_{\widehat D},\varepsilon_S)$ with $\varepsilon_S(u)(T)=T(u)$ is of
source type, the piece $(E,\cL(E,D)\otimes_\pi E,\id_E,w'_{b_R},\varepsilon_R)$ with $\varepsilon_R(u)(x)=u\otimes x$ is
of target type, and $\varepsilon_S(u)w_{b_S}x+w'_{b_R}\varepsilon_R(u)x=b_S(u,x)+b_R(u,x)=ux$.
\end{proof}

\begin{proof}[Proof of Lemma~\ref{cmb:lem:subpair}]
(1) By~\eqref{cmb:eq:decomposition}, $v=\Pi\Phi(v)J=\sum_i\Pi w'_i\,\varepsilon_i(\Phi v)\,w_iJ$, and by the ideal property
(Proposition~\ref{ob:prop:split-operators}(2)) the pieces $(G_i,G'_i,w_iJ,\Pi w'_i,\varepsilon_i\circ\Phi)$ have the types
of the given pieces. A source-side relative splitting $(G,w,\varepsilon)$ gives $(G,wJ,v\mapsto\Pi\varepsilon(\Phi v))$,
and a target-side one $(G,w',\varepsilon)$ gives $(G,\Pi w',v\mapsto\varepsilon(\Phi v)J)$. (2) Apply
Proposition~\ref{cmb:prop:transfer} with $G=E$, $G'=D$, $w=J$, $w'=\Pi|_D$ and $\varepsilon=\Phi$.
\end{proof}

\begin{proof}[Proof of Proposition~\ref{cmb:prop:parameter-pieces}]
\emph{Step 1 (items (1) and (3)).} Item~(1) is Proposition~\ref{spl:prop:gamma-l1}. For~(3), let $J\colon P_0\to P$ and
$\Pi'\colon P\to P_0$ be bounded linear with $\Pi'J=\id_{P_0}$, and let $R_n$ be a projection of $\widehat\otimes^n_\pi P$
onto $\hGamma^n_\pi P$. Then ${\Pi'}^{\otimes n}R_nJ^{\otimes n}$ takes values in $\hGamma^n_\pi P_0$, since
${\Pi'}^{\otimes n}$ maps pure powers to pure powers, and it fixes $h^{\otimes n}={\Pi'}^{\otimes n}(Jh)^{\otimes n}$ for
$h\in P_0$.

\emph{Step 2 (item (2)).} Let $n\le k$, let $J_1,J_2$ and $\Pi_1,\Pi_2$ be the inclusions and projections of
$P=P_1\oplus P_2$, and let $R^{(1)}_a$, $R^{(2)}_a$ be projections onto $\hGamma^a_\pi P_1$, $\hGamma^a_\pi P_2$ for
$a\le n$, with $R^{(1)}_0=R^{(2)}_0=\id_K$. For $0\le a\le n$ let
$\operatorname{sh}_a\colon\widehat\otimes^a_\pi P_1\hot_\pi\widehat\otimes^{n-a}_\pi P_2\to\widehat\otimes^n_\pi P$ send
$y\otimes z$ to the sum of the $\binom na$ tensors obtained by distributing the factors of $J_1^{\otimes a}y$ and of
$J_2^{\otimes(n-a)}z$ over the $n$ slots, keeping their order, and put
\[
 R=\sum_{a=0}^n\operatorname{sh}_a\circ\bigl(R^{(1)}_a\otimes R^{(2)}_{n-a}\bigr)\circ
 \bigl(\Pi_1^{\otimes a}\otimes\Pi_2^{\otimes(n-a)}\bigr).
\]
For $h=J_1h_1+J_2h_2$, $(\Pi_1^{\otimes a}\otimes\Pi_2^{\otimes(n-a)})(h^{\otimes n})=h_1^{\otimes a}\otimes
h_2^{\otimes(n-a)}$ and $h^{\otimes n}=\sum_a\operatorname{sh}_a(h_1^{\otimes a}\otimes h_2^{\otimes(n-a)})$, so $R$ fixes
$h^{\otimes n}$, hence $\hGamma^n_\pi P$. The tensor $\operatorname{sh}_a(y^{\otimes a}\otimes z^{\otimes(n-a)})$ is the
coefficient of $\lambda^a$ in the polynomial $\lambda\mapsto(\lambda J_1y+J_2z)^{\otimes n}$, hence a linear combination of
its values, $K$ being infinite. Since $R^{(1)}_a\otimes R^{(2)}_{n-a}$ takes values in the closed span of the tensors
$y^{\otimes a}\otimes z^{\otimes(n-a)}$, the map $R$ takes values in $\hGamma^n_\pi P$.

\emph{Step 3 (item (4)).} Let $n\le k$ and let $R_n$ be a projection onto $\hGamma^n_\pi P$. The operators that vanish on
$\cN^n_E$ form a closed subspace $W$ of $Z=\cL(\widehat\otimes^n_\pi E,\widehat\Lambda^n_\pi D)$, and
$T\mapsto T\circ\widehat\wedge^n_E$ is an isometry of $\cL(\widehat\Lambda^n_\pi E,\widehat\Lambda^n_\pi D)$ onto $W$
(Lemma~\ref{ob:lem:metric-surjection}). Let $\Theta\colon\widehat\otimes^n_\pi P\to Z$ be the linearization of the bounded
$n$-linear map
$(y_1,\dots,y_n)\mapsto\widehat\wedge^n_{\widehat D}\circ\bigl(w'\varepsilon'(y_1)w\otimes\dots\otimes w'\varepsilon'(y_n)w\bigr)$.
By naturality of $\wedge^n$, $\Theta(y^{\otimes n})=\widehat\Lambda^n_\pi(w'\varepsilon'(y)w)\circ\widehat\wedge^n_E$ lies in
$W$, so $\Theta$ maps $\hGamma^n_\pi P$ into $W$. Hence
$Q_n(u_1,\dots,u_n)=\Theta\bigl(R_n(\psi u_1\otimes\dots\otimes\psi u_n)\bigr)$, read in
$\cL(\widehat\Lambda^n_\pi E,\widehat\Lambda^n_\pi D)$ and restricted to $\Lambda^n_\pi E$, is bounded and $n$-linear, and
its diagonal is $\Lambda^n_\pi(w'\varepsilon(u)w)$, since $R_n$ fixes $(\psi u)^{\otimes n}$. The last assertion follows
from Theorem~\ref{cmb:thm:decomposition}(1)--(3).
\end{proof}

\begin{proof}[Proof of Corollary~\ref{cmb:cor:transport}]
(1) Extend $w'$ to $\widehat G'$. Then $\theta_{E,D}(u)=w'\circ\theta_{G,G'}(\varepsilon(u))\circ w$ for $u\in\cL(E,D)$.
Proposition~\ref{cmb:prop:naturality}(b), (c) and~(a), applied in this order to $w$, $w'$ and $\psi=\varepsilon$, give
$\ob_k(E,D)=(\widehat\Lambda^k_\pi w)^*((\widehat\Lambda^k_\pi w')_\flat)_*(\varepsilon^\sharp)_*\ob_k(G,G')$, and pullbacks
commute with pushouts (Corollary~\ref{ocl:cor:functorial}). (2) The family $\varepsilon$ on $\cL(E,D)$ satisfies
$\theta_{E,D}=\varepsilon_w$, so Proposition~\ref{cmb:prop:naturality}(b) gives the formula. If $w$ is $\mathrm S^k$-split,
then $\operatorname{sp}_k(w)=0$ by Proposition~\ref{ob:prop:split-operators}(1), and the pushout of a split sequence splits by
Lemma~\ref{ob:lem:split-criteria}(2) and~(3).
\end{proof}

\begin{proof}[Proof of Proposition~\ref{cmb:prop:blocks}]
\emph{Step 1 (the alternating map).} Fix $u^1,\dots,u^k\in\mathcal U$, and let $E_{\mathrm{fin}}\subseteq E$ be the
subspace of vectors with finitely many nonzero blocks. For $\nu=(\nu_1,\dots,\nu_k)\in I^k$ let $i_1<\dots<i_m$ be the
distinct entries of $\nu$, let $S_l=\{a:\nu_a=i_l\}$ and $n_l=\abs{S_l}$, and let $\epsilon(S_1,\dots,S_m)$ be the sign of
the permutation of $\{1,\dots,k\}$ that lists $S_1$, then $S_2$, and so on, each in increasing order. For
$x_1,\dots,x_k\in E_{\mathrm{fin}}$ put
\[
 \Phi(x_1,\dots,x_k)=\sum_{\nu\in I^k}\epsilon(S_1,\dots,S_m)\bigwedge_{l=1}^m\widehat\Lambda^{n_l}_\pi J'_{i_l}\,
 Q^{i_l}_{n_l}\bigl(u^{(l)}_{i_l}\bigr)\Bigl(\bigwedge_{a\in S_l}\Pi_{i_l}x_a\Bigr),
\]
a finite sum in $\widehat\Lambda^k_\pi D$. Here $u^{(l)}_{i_l}$ is the tuple of the $u^j_{i_l}$ with
$n_1+\dots+n_{l-1}<j\le n_1+\dots+n_l$, in increasing order of $j$; the wedges over $S_l$ are taken in increasing order;
and the outer wedge is the iterated wedge product, extended to completions. Each $u^j$ enters one slot of one map
$Q^{i_l}_{n_l}$, determined by the sizes $n_l$, so $\Phi$ is $k$-linear in $(x_1,\dots,x_k)$ and in $(u^1,\dots,u^k)$.
It is alternating: let $x_a=x_b$ with $a\ne b$. The terms with $\nu_a=\nu_b$ vanish, since one block wedge contains
$\Pi_{\nu_a}x_a$ twice. The term of $\nu\circ(a\ b)$ is minus the term of $\nu$ with $x_a$ and $x_b$ exchanged: the two
tuples have the same entries and block sizes, the permutations defining their signs differ by the transposition $(a\ b)$
followed by permutations inside the blocks, and each block factor is alternating in the vectors of its block; so the
remaining terms cancel in pairs.

\emph{Step 2 (the bound).} The wedge product is contractive and $\norm{\widehat\Lambda^n_\pi J'_i}\le1$, so the term of
$\nu$ has norm at most $C^m\prod_j(c\norm{u^j})\prod_a\norm{\Pi_{\nu_a}x_a}\le C^kc^k\prod_j\norm{u^j}\prod_a
\norm{\Pi_{\nu_a}x_a}$. Summing over $\nu$ gives $\norm{\Phi(x_1,\dots,x_k)}\le C^kc^k\prod_j\norm{u^j}\prod_a\norm{x_a}$,
since $\sum_\nu\prod_a\norm{\Pi_{\nu_a}x_a}=\prod_a\norm{x_a}$. Hence $\Phi$ extends by continuity to a bounded
$k$-linear map on $E^k$, alternating because a tuple in $E^k$ with two equal entries is a limit of such tuples in
$E_{\mathrm{fin}}^k$, and Proposition~\ref{prop:exterior} gives $Q(u^1,\dots,u^k)\in\cL(\Lambda^k_\pi E,\widehat\Lambda^k_\pi D)$,
$k$-linear in $(u^1,\dots,u^k)$, with $\norm Q\le C^kc^k$.

\emph{Step 3 (the diagonal).} Let $u^j=u$ for every $j$. The $l$-th block factor of the term of $\nu$ is then
$\bigwedge_{a\in S_l}J'_{i_l}u_{i_l}\Pi_{i_l}x_a$, so by the choice of the sign the term is
$\bigwedge_{a=1}^kJ'_{\nu_a}u_{\nu_a}\Pi_{\nu_a}x_a$. Summing over $\nu$ gives
$\bigwedge_a\bigl(\sum_iJ'_iu_i\Pi_ix_a\bigr)=ux_1\wedge\dots\wedge ux_k$ for $x_1,\dots,x_k\in E_{\mathrm{fin}}$. Wedges
of vectors of $E_{\mathrm{fin}}$ span a dense subspace of $\Lambda^k_\pi E$, so $Q(u,\dots,u)=\Lambda^k_\pi u$. The last
assertion is Lemma~\ref{fib:lem:test}.
\end{proof}

\begin{proof}[Proof of Proposition~\ref{cmb:prop:permanence}]
\emph{Step 1 (retracts and sums).} Let $J\colon Y_0\to Y$ and $\Pi'\colon Y\to Y_0$ be bounded linear with
$\Pi'J=\id_{Y_0}$. Since $\id_{Y_0}=\Pi'\,\id_Y\,J$, the conditions $(\mathrm S^k)$ and $(\mathrm R^k)$ pass to $Y_0$ by
Proposition~\ref{ob:prop:split-operators}(2), and polarizability by Proposition~\ref{cmb:prop:parameter-pieces}(3). For
pairs, apply Lemma~\ref{cmb:lem:subpair}(2) in the situation of its second sentence. In~(2), the conditions pass from the
sum to $Y_1$ and $Y_2$ by~(1). Conversely, $\id_{Y_1\oplus Y_2}$ is the sum of the inclusion-projection composites of the
two summands, each $\mathrm S^{\le k}$-split, respectively $\mathrm R^{\le k}$-split, by the ideal property, so their sum
is split in the same degrees by the paragraph on split operators in Section~\ref{cmb:ssec:degrees}; polarizability is
Proposition~\ref{cmb:prop:parameter-pieces}(2). For the pairs, let $J_a,\Pi_a$ be the inclusion and projection of $E_a$
in $E_1\oplus E_2$, and $J'_b,\Pi'_b$ those of $D_b$ in $D_1\oplus D_2$. Every $u\in\cL(E_1\oplus E_2,D_1\oplus D_2)$ is
$\sum_{a,b}J'_b(\Pi'_buJ_a)\Pi_a$. For each $(a,b)$ the map $u\mapsto J'_b(\Pi'_buJ_a)\Pi_a$ is liftable in every degree
$n\le k$, by the formula of Proposition~\ref{cmb:prop:transfer} for the pair $(E_a,D_b)$, with $w=\Pi_a$, $w'=J'_b$ and
$\varepsilon(u)=\Pi'_buJ_a$. Theorem~\ref{cmb:thm:decomposition}(2), applied in each degree $n\le k$, concludes.

\emph{Step 2 (completion of a space).} By the paragraph before Theorem~\ref{ob:thm:mechanisms}, applied to the dense
subspace $Y$ of $\widehat Y$, the inclusion induces isometries $\otimes^k_\pi Y\to\otimes^k_\pi\widehat Y$ and
$\Lambda^k_\pi Y\to\Lambda^k_\pi\widehat Y$ with dense images. So the completed powers of $Y$ and $\widehat Y$ agree, as do
the maps $\widehat\wedge^k$ and $\alt^k$ on them, and $\hGamma^k_\pi Y=\hGamma^k_\pi\widehat Y$, since
$h\mapsto h^{\otimes k}$ is continuous. Sections and retractions for $\id_Y$ extend by continuity to sections and
retractions for $\id_{\widehat Y}$, since both sides of the defining identities are continuous, and those for
$\id_{\widehat Y}$ restrict to ones for $\id_Y$. The projections onto $\hGamma^k_\pi Y=\hGamma^k_\pi\widehat Y$ are the
same.

\emph{Step 3 (completion of a pair).} Let $D$ be complete. Restriction is an isometry of $\cL(\widehat E,D)$ onto
$\cL(E,D)$; write $\widehat u$ for the extension of $u\in\cL(E,D)$. By Step~2, $\Lambda^k_\pi E$ is a dense subspace of
$\Lambda^k_\pi\widehat E$ with the induced norm. A lift $\widehat Q$ for $(\widehat E,D)$ gives the lift
$Q(u_1,\dots,u_k)=\widehat Q(\widehat u_1,\dots,\widehat u_k)|_{\Lambda^k_\pi E}$ for $(E,D)$. A lift $Q$ for $(E,D)$ gives the lift for $(\widehat E,D)$ whose value at $(v_1,\dots,v_k)$ is the extension
by continuity of $Q(v_1|_E,\dots,v_k|_E)$; its diagonal is $\Lambda^k_\pi v$, since $\Lambda^k_\pi v$ extends
$\Lambda^k_\pi(v|_E)$. Finally $\cL(E,D)\subseteq\cL(E,\widehat D)$ isometrically for any normed $D$, and
$\widehat\Lambda^k_\pi\widehat D=\widehat\Lambda^k_\pi D$, so a lift for $(E,\widehat D)$ restricts to a lift for
$(E,D)$. Lemma~\ref{fib:lem:test} concludes.
\end{proof}
%% ---- end app-Q-sec06.tex
%% ---- begin app-Q-sec07.tex
\subsection{Proofs for Section~\ref{sec:sharpness}}\label{dfr:ssec:sharpness}

The notation is that of Section~\ref{sec:sharpness}.

\begin{proof}[Proof of Lemma~\ref{obs:lem:normalization}]
\emph{Step 1 (value group and residue field).} The field $\Kh$ is closed in the complete field $\kappa((X))$, hence
complete. Put $\mathcal O=\{a\in\Kh:\abs a\le1\}$ and $\mathfrak m=\{a\in\Kh:\abs a<1\}$. For $a\in\Kh^\times$ there is
$a'\in K$ with $\abs{a-a'}<\abs a$, hence $\abs{a'}=\abs a$; so $\abs{\Kh^\times}=\abs{K^\times}=\abs\vartheta^\Z$, and
$\abs a\le\abs\vartheta$ for every $a\in\mathfrak m$. Since $\mathcal O=\Kh\cap\kappa[[X]]$ and
$\mathfrak m=\Kh\cap X\kappa[[X]]$, the map $\coeff_0$ restricts to a ring homomorphism $\mathcal O\to\kappa$ with kernel
$\mathfrak m$. So $\mathcal O/\mathfrak m$ is isomorphic to a subfield $\kappa_1$ of $\kappa$, and $\abs{\kappa_1}$ is a
power of $p$.

\emph{Step 2 (Teichm\"uller representatives).} Since $\Kh$ has characteristic $p$,
$\operatorname{Fr}(a)=a^{\abs{\kappa_1}}$ is a continuous ring endomorphism of $\mathcal O$; let $\kappa'$ be its subring
of fixed points. A nonzero $b\in\kappa'$ satisfies $b^{\abs{\kappa_1}-1}=1$, since $\mathcal O$ is a domain; so $\kappa'$
is a field, $\abs b=1$, and the reduction map $\kappa'\to\mathcal O/\mathfrak m$ is injective. For surjectivity let
$a\in\mathcal O$ and $a_n=\operatorname{Fr}^n(a)$. Since $y^{\abs{\kappa_1}}=y$ on $\kappa_1$, $w=\operatorname{Fr}(a)-a$
lies in $\mathfrak m$, and $\abs{a_{n+1}-a_n}=\abs{w^{\abs{\kappa_1}^n}}\le\abs\vartheta^{\abs{\kappa_1}^n}$. Hence
$a_n\to b\in\mathcal O$ with $\operatorname{Fr}(b)=b$ and $\abs{b-a}<1$. So $\kappa'\cong\kappa_1$ is a finite field of
characteristic $p$.

\emph{Step 3 (expansions).} Let $a\in\mathcal O$. By Step~2 there is $b_0\in\kappa'$ with $a-b_0\in\mathfrak m$, and
$(a-b_0)/\vartheta\in\mathcal O$ by Step~1. Repeating this gives $b_0,b_1,\ldots\in\kappa'$ with
$a-\sum_{i<n}b_i\vartheta^i\in\vartheta^n\mathcal O$ for every $n$, so $a=\sum_{i\ge0}b_i\vartheta^i$. If all $b_i$ lie
in $\kappa'$ and $b_n\ne0$, then $\abs{\sum_{i\ge n}b_i\vartheta^i}=\abs\vartheta^n$, since $\abs{b_n\vartheta^n}=
\abs\vartheta^n$ and the later terms have absolute value at most $\abs\vartheta^{n+1}$. Applied to the difference of two
expansions, this shows that the expansion is unique.

\emph{Step 4 (the isomorphism).} For $g=\sum_{i\ge m}b_iX^i$ in $\kappa'((X))$ put $\sigma(g)=\sum_{i\ge m}b_i\vartheta^i$,
which converges since $\abs{b_i\vartheta^i}\le\abs\vartheta^i$. The map $\sigma$ is additive, and it is multiplicative
since in the complete nonarchimedean field $\Kh$ the product of two convergent series is the sum of their Cauchy
product, whose coefficients $\sum_{i+j=n}b_ib'_j$ lie in $\kappa'$. If $g\ne0$ has order $m$, Step~3 applied to
$\vartheta^{-m}\sigma(g)$ gives $\abs{\sigma(g)}=\abs\vartheta^m$; so $\sigma$ is isometric, hence injective. For
$a\in\Kh^\times$ choose $n$ with $\abs a=\abs\vartheta^n$ (Step~1); then $\vartheta^{-n}a=\sum_{i\ge0}b_i\vartheta^i$ by
Step~3, and $a=\sigma(\sum_ib_iX^{i+n})$. So $\sigma$ is an isometric field isomorphism with $\sigma(X)=\vartheta$. It
restricts to an isometric isomorphism of the dense subfield $\sigma^{-1}(K)$ of $\kappa'((X))$ onto $K$, and every
statement that depends only on $K$ as a normed field transfers along this isomorphism.
\end{proof}

\begin{proof}[Proof of Lemma~\ref{obs:lem:multiplier}]
(1) For $x\in c_{00}(\N,K)$ and $t\in\N$, $\abs{x(t)}e_t\le\abs x$ pointwise, so $\abs{x(t)}\le\norm x_E$ by monotonicity
and $\norm{e_t}_E=1$. For $x\in E$ and $t\in[n]$ this gives $\abs{x(t)}\le\norm{x1_{[n]}}_E\le\norm x_E$. If $b$ is
finitely supported with $\abs b\le1$ and $[n]$ contains its support, then $bx=b(x1_{[n]})$ and $\abs{bx}\le\abs{x1_{[n]}}$,
so $\norm{bx}_E\le\norm{x1_{[n]}}_E\le\norm x_E$.

(2) For $x\in E$, $ax$ is finitely supported and $\norm{ax}_D\le\sum_{t\in\operatorname{supp}a}\abs{a(t)}\norm{e_t}_D
\norm x_E$ by~(1); so $\mu_a$ is bounded. Moreover $ax=a(x1_{\operatorname{supp}a})$, where $x1_{\operatorname{supp}a}$ is
supported in $\operatorname{supp}a$ and has norm at most $\norm x_E$ by~(1). This gives the formula.

(3) The map $a\mapsto\mu_a$ is linear, and $\norm{\mu_ae_t}_D=\abs{a(t)}$ gives $\norm{\mu_a}_{E\to D}\ge\norm a_\infty$;
so $a\mapsto\norm{\mu_a}_{E\to D}$ is a norm. Since $\mu_{e_t}x=x(t)e_t$, (1) gives $\norm{\mu_{e_t}}_{E\to D}=1$. For an
injective map $\gamma\colon\N\to\N$ and finitely supported $y$, let $y^\gamma(\gamma(t))=y(t)$ and $y^\gamma=0$ off
$\gamma(\N)$. Then $\norm{y^\gamma}=\norm y$ in $E$ and in $D$, since $\gamma$ restricted to $\operatorname{supp}y$
extends to a permutation of $\N$. The vectors supported in $\operatorname{supp}a^\gamma$ are the $x^\gamma$ with
$\operatorname{supp}x\subseteq\operatorname{supp}a$, and $a^\gamma x^\gamma=(ax)^\gamma$; so~(2) gives
$\norm{\mu_{a^\gamma}}_{E\to D}=\norm{\mu_a}_{E\to D}$. Multiplication of $a$ by a sequence of absolute value $1$ does not
change $\abs{ax}$. If $\abs a\le\abs b$, let $c=a/b$ on $\operatorname{supp}b$ and $c=0$ elsewhere; then $\abs c\le1$,
$ax=b(cx)$ and $\norm{cx}_E\le\norm x_E$ by~(1).

(4) For $x\in c_{00}(\N,K)$, $\abs{ax}\le\norm a_\infty\abs x$ pointwise, so
$\norm{ax}_D\le C\norm{ax}_E\le C\norm a_\infty\norm x_E$. If $D$ is symmetric and $\norm x_\infty\le1$, then
$\abs{ax}\le\abs a$, so $\norm{ax}_D\le\norm a_D$, with equality at $x=1_{\operatorname{supp}a}$. Finally
$\norm{ax}_\infty\le\norm a_\infty\norm x_\infty\le\norm a_\infty\norm x_E$ by~(1), and~(3) gives the reverse inequality.
\end{proof}
%% ---- end app-Q-sec07.tex
%% ---- begin app-Q-sec09.tex
\subsection{Proofs for Section~\ref{sec:beyond}}\label{dfr:ssec:beyond}

\begin{proof}[Proof of Proposition~\ref{bey:prop:one-dim}]
\emph{Step 1 (the lift).} Each $u\in\cL(E,D)$ is $\lambda u_0$ for a unique $\lambda\in K$, and
$\abs\lambda=\norm u/\norm{u_0}$. Hence $Q$ is a well-defined $k$\nobreakdash-linear map
$\cL(E,D)^k\to\cL(\Lambda^k_\pi E,\Lambda^k_\pi D)$ of norm at most $\norm{\Lambda^k_\pi u_0}/\norm{u_0}^k\le1$. Its
diagonal is $\lambda^k\Lambda^k_\pi u_0=\Lambda^k_\pi(\lambda u_0)$, so Lemma~\ref{fib:lem:test} gives~(1).

\emph{Step 2 (sufficiency in~(2)).} Put $\varepsilon(\lambda u_0)=\lambda\,\id_D$ and
$\varepsilon'(\lambda u_0)=\lambda\,\id_E$. Then $\varepsilon(u)\circ u_0=u=u_0\circ\varepsilon'(u)$ for every
$u\in\cL(E,D)$. If $u_0$ is $\mathrm S^k$-split, $(D,u_0,\varepsilon)$ is a source-side relative splitting; if $u_0$ is
$\mathrm R^k$-split, $(E,u_0,\varepsilon')$ is a target-side relative splitting. Either is a splitting decomposition of
degree $k$ (Definition~\ref{cmb:def:decomposition}).

\emph{Step 3 (necessity in~(2), one piece).} A splitting decomposition with one piece writes $u_0$ as a composite of
bounded operators one of which is $\mathrm S^k$-split or $\mathrm R^k$-split (Definition~\ref{cmb:def:decomposition}).
By Proposition~\ref{ob:prop:split-operators}(2), $u_0$ is $\mathrm S^k$-split or $\mathrm R^k$-split.

\emph{Step 4 (necessity in~(2), several pieces).} For a decomposition into at least two pieces of source and target
type, Proposition~\ref{cmb:prop:normal-form} gives, in its notation, bounded bilinear maps $b_S,b_R$ with
$b_S(u,x)+b_R(u,x)=ux$, $w_{b_S}$ $\mathrm S^{\le k}$-split and $w'_{b_R}$ $\mathrm R^{\le k}$-split; their values lie in
$D$, since $D$ is complete. Put $u_S=b_S(u_0,\cdot)$ and $u_R=b_R(u_0,\cdot)$. Since $u_S$ is $w_{b_S}$ followed by
evaluation at $u_0$, and $u_R$ is $x\mapsto u_0\otimes x$ followed by $w'_{b_R}$,
Proposition~\ref{ob:prop:split-operators}(2) makes $u_S$ $\mathrm S^k$-split and $u_R$ $\mathrm R^k$-split. Both
lie in $K\,u_0$ and $u_S+u_R=u_0$, so there is $a\in K$ with $u_S=a\,u_0$ and $u_R=(1-a)\,u_0$. If $a\ne0$, then
$u_0=a^{-1}u_S$ is $\mathrm S^k$-split; otherwise $u_0=u_R$ is $\mathrm R^k$-split.

\emph{Step 5 (the dual).} If $f$ is a nonzero bounded linear functional on $E$, the operators $x\mapsto f(x)\,d$ with
$d\in D$ form a subspace of $\cL(E,D)$ of dimension $\dim D\ge2$, contrary to $\cL(E,D)=K\,u_0$.
\end{proof}
%% ---- end app-Q-sec09.tex
%% ---- begin app-Q-sec10.tex
\subsection{Proofs for Section~\ref{sec:geometry}}\label{dfr:ssec:geometry}

\begin{proof}[Proof of Proposition~\ref{prop:tangent-obstruction}]
Choose $E,D$ from Theorem~\ref{thm:operator-obstructions}(2), so that $A=A^k_{E,D;K}$ is analytic at no point.

\emph{Step 1 (assertion~(1)).} Put $P=\cL(E,D)\times D\times E$ and use the charts $\id_P$ and $\psi(u,d,e)=(u,d+ue,e)$. Operator evaluation is bounded bilinear, so $\psi$ and $\psi^{-1}(u,d,e)=(u,d-ue,e)$ are polynomial analytic, with $D\psi(u,d,e)(v,a,b)=(v,a+ve+ub,b)$. Let $J$ extend forms from $D$ along the projection $\pi_D\colon P\to D$, and let $R$ restrict forms on $P$ along $i_E(b)=(0,0,b)$. These bounded linear maps give
\[
 R\circ A^k_{P,P;K}(D\psi(u,d,e))\circ J=A(u).
\]
If the transition were analytic at $(u_0,d_0,e_0)$, restriction to the affine slice $u\mapsto(u,d_0,e_0)$ would make $A$ analytic at $u_0$. The reverse transition is treated in the same way with $-u$.

\emph{Step 2 (assertion~(2)).} Let $U=\Alt^k(D;K)$, $Y=U\times D$ and $N=\cL(E,D)\times U\times E$. Define the bounded linear map $\omega_0(m,d)=m\circ\pi_D^{\times k}$, with $\pi_D\colon Y\to D$ the projection, and the polynomial map $h(u,m,e)=(m,ue)$. Since $Dh(u,m,e)(v,a,b)=(a,ve+ub)$, restricting the values of $h^*\omega_0$ to the $E$\nobreakdash-directions gives $f(u,m)=A(u)(m)$. If $h^*\omega_0$ were analytic as an $\Alt^k(N;K)$\nobreakdash-valued map at $(u_0,m_0,e_0)$, fixing $e=e_0$ would make $f$ analytic near $(u_0,m_0)$; since Banach-valued power series differentiate termwise on a smaller ball, $\partial_mf(u,m)=A(u)$ would be analytic near $u_0$, a contradiction. In $\Mult^k(N;K)$ the pullback is a bounded multilinear expression in $Dh$ and $\omega_0\circ h$, hence analytic.
\end{proof}

\begin{proof}[Proof of Proposition~\ref{prop:tangent-values}]
Throughout, for Banach spaces $E$ and $D'$ put $P=\cL(E,D')\times D'\times E$, with the maximum norm, and use the charts $\id_P$ and $\psi(u,d,e)=(u,d+ue,e)$ of Step~1 of the proof of Proposition~\ref{prop:tangent-obstruction}. Both charts are polynomial analytic, and $g=D\psi(u,d,e)$ is $g(v,a,b)=(v,a+ve+ub,b)$, with $g^{-1}(v,a,b)=(v,a-ve-ub,b)$.

\emph{Step 1 (values in $F$).} For~(1), take $E,F$ from Theorem~\ref{thm:operator-obstructions}(1) and $D'=E$. The maps $J,R$ of Step~1 of the proof of Proposition~\ref{prop:tangent-obstruction}, with values in $F$, give
\[
 R\circ A^k_{P,P;F}(D\psi(w,d,e))\circ J=A^k_{E,E;F}(w)\qquad(w\in\cL(E,E)),
\]
which is analytic at no point; restriction to affine slices, as in that step, shows that neither transition is analytic at any point. Only bounded linear compression and affine slices are used, so $K$ need not be complete, and $F$ has no equivalent nonarchimedean norm by Theorem~\ref{thm:operator-obstructions}(1). For $K=\F_q((X))$ with its $X$\nobreakdash-adic absolute value, Remark~\ref{rem:spaces} allows $E=\linf(\N,K)$ and $F=F_1$, which has no equivalent nonarchimedean norm by Proposition~\ref{prop:nonultra}. Then $P$ is nonarchimedean, since operator norms into a nonarchimedean space and maximum norms of products are nonarchimedean.

\emph{Step 2 (tangent values).} Suppose that $A^k_{E,D;F}$ is analytic at no point, and put $D'=D\times F$. As functions of the point $(u,d,e)$, the induced transitions on tangent-valued forms are $m\mapsto g\circ m\circ(g^{-1},\ldots,g^{-1})$ and $m\mapsto g^{-1}\circ m\circ(g,\ldots,g)$. Let $\pi_D\colon P\to D$ and $\pi_F\colon P\to F$ be the components of the middle factor, and put $i_F(y)=(0,(0,y),0)$ and $i_E(b)=(0,0,b)$. The bounded linear map
\[
 \Phi(T)(m)=\pi_F\circ T\bigl(i_F\circ m\circ\pi_D^{\times k}\bigr)\circ i_E^{\times k}
\]
sends operators $T$ on $\Alt^k(P;P)$ to operators $\Alt^k(D;F)\to\Alt^k(E;F)$, and it sends $m\mapsto h\circ m\circ(f,\ldots,f)$, for $f,h\in\cL(P,P)$, to $m\mapsto(\pi_Fhi_F)\circ m\circ(\pi_Dfi_E,\ldots,\pi_Dfi_E)$. Since $\pi_Fg^{\pm1}i_F=\id_F$ and $\pi_Dg^{\pm1}i_E=\pm u_D$, where $u_D\in\cL(E,D)$ is the $D$\nobreakdash-component of $u$, $\Phi$ sends the two transitions to $A^k_{E,D;F}(\mp u_D)=(\mp1)^kA^k_{E,D;F}(u_D)$. If either transition were analytic at $(u_0,d_0,e_0)$, restriction to the affine slice $w\mapsto(u_0+(w,0),d_0,e_0)$, $w\in\cL(E,D)$, would make $A^k_{E,D;F}$ analytic at the $D$\nobreakdash-component of $u_0$.

\emph{Step 3 (the models).} For~(2), apply Step~2 with $E,F$ from Theorem~\ref{thm:operator-obstructions}(1) and $D=E$; then $P$ contains the isometric copy $i_F(F)$ of $F$, so it admits no equivalent nonarchimedean norm. For~(3), apply Step~2 with $E,D$ from Theorem~\ref{thm:operator-obstructions}(2) and $F=K$; then $P$ is a nonarchimedean Banach space, as in Step~1.
\end{proof}

\begin{proof}[Proof of Theorem~\ref{thm:ambient-forms}]
\emph{Step 1 (sheaf and pullback).} Locality of analytic coefficient maps gives the sheaf property. Addition and multiplication by an analytic function are bounded multilinear operations on the coefficient spaces. For an analytic map $h$, the formula
\[
 h^*\eta=(\eta\circ h)\circ(Dh,\ldots,Dh)
\]
is analytic in the ambient space and has alternating values. The chain rule gives identity and composition laws for pullback.

\emph{Step 2 (the product).} The shuffle formula~\eqref{eq:form-shuffle} has $\binom{k+l}{k}$ terms, each of norm at most $\norm\eta\norm\zeta$, so it is a bounded bilinear operation on multilinear coefficient spaces, of norm at most $\binom{k+l}{k}$. For alternating inputs it is alternating: terms with two repeated arguments in one block vanish, and the remaining terms cancel in pairs by exchanging those arguments between blocks; this argument includes characteristic two. Formula~\eqref{eq:form-shuffle} commutes with simultaneous precomposition, hence glues in charts and commutes with pullback. Complementing $S$ proves graded commutativity. Expanding either bracketing of a triple product yields the same signed sum over ordered three-block partitions, proving associativity.

\emph{Step 3 ($d\eta$ is a form).} On an open subset of a model space $P$, the operation~\eqref{eq:exterior-derivative} applied to $D\eta$ is bounded linear with norm at most $k+1$, so $d\eta$ is ambient analytic. Each $D\eta(y)(v_i)$ is alternating: difference quotients have alternating values, and that subspace is closed. If $v_a=v_b$ with $a<b$, the summands other than $i=a,b$ vanish; alternation moves the remaining repeated argument across $b-a-1$ positions, making those two terms cancel. Hence $d\eta$ is a form.

\emph{Step 4 ($d^2=0$).} The second derivative of an analytic map is symmetric in every characteristic: if $c_2$ is the bilinear coefficient of degree two of a power series, the second derivative at its center is $(v,w)\mapsto c_2(v,w)+c_2(w,v)$. In the expansion of $d(d\eta)$, pair the terms differentiated in directions $(v_a,v_b)$ and $(v_b,v_a)$ for $a<b$. Their values agree and their signs are $(-1)^{a+b-1}$ and $(-1)^{a+b}$, so $d^2\eta=0$.

\emph{Step 5 (naturality and gluing of $d$).} Differentiate the pullback formula. Terms differentiating $\eta$ give $h^*(d\eta)$ after the signed sum~\eqref{eq:exterior-derivative}. The other terms insert $D^2h(v_a,v_b)$ into one argument of $\eta$, and symmetry of $D^2h$ and alternation of $\eta$ cancel them in pairs. Therefore $d(h^*\eta)=h^*(d\eta)$; applied to chart changes, this identity shows that the local derivatives glue to $d\colon\Omega^k_{\mathrm{amb}}(M)\to\Omega^{k+1}_{\mathrm{amb}}(M)$.

\emph{Step 6 (Leibniz rule).} Differentiating the shuffle formula~\eqref{eq:form-shuffle} gives the Leibniz rule: placing the derivative index in the second block moves it past the first $k$ arguments and contributes $(-1)^k$. Every identity in Steps~1--6 uses finite sums and cancellation, with no division by a factorial.
\end{proof}

\begin{proof}[Proof of Proposition~\ref{fib:prop:derham}]
Clearly (2) implies (1). If (1) holds, put $T=D\eta(y)$ and $\beta(v_0,\dots,v_{p-1})=T(v_0)(v_1,\dots,v_{p-1})$. Grouping the permutations by the image of $0$ and using alternation of each $T(v_i)$,
\[
 \mathcal A_Y\beta(v_0,\dots,v_{p-1})=(p-1)!\sum_{i=0}^{p-1}(-1)^iT(v_i)(v_0,\dots,\widehat{v_i},\dots,v_{p-1})
 =(p-1)!\,\omega(v_0,\dots,v_{p-1}),
\]
and $(p-1)!=-1$ in $K$ by Wilson's theorem, so $\omega=\mathcal A_Y(-\beta)$. If (3) holds, the bounded linear map $\eta(y)(v_1,\dots,v_{p-1})=\sum_\sigma\sgn(\sigma)\,\beta(y,v_{\sigma(1)},\dots,v_{\sigma(p-1)})$, summed over the permutations $\sigma$ of $\{1,\dots,p-1\}$, satisfies $D\eta(y)=\eta$, and the same grouping gives $d\eta=\mathcal A_Y\beta=\omega$. For $F=\widehat\Lambda^p_\pi Y$ the equivalence with $(\mathrm R^p_Y)$ is Lemma~\ref{ob:lem:target}, whose proof also gives the uncompleted case, and Proposition~\ref{fib:prop:split} gives goodness.
\end{proof}

\begin{proof}[Proof of Theorem~\ref{dp:thm:calculus}]
\emph{Step 1 (nodes).} For an integer $a\ge0$ with binary digits $a_l\in\{0,1\}$, so that $a=\sum_la_l2^l$, put
$\xi_a=\sum_la_l\varpi^l\in K$. Then $\abs{\xi_a}\le1$, and for $a\ne b$ the least $l$ with $a_l\ne b_l$ is the exponent
$v(a-b)$ of $2$ in $a-b$, so $\abs{\xi_a-\xi_b}=\abs\varpi^{v(a-b)}$ by the ultrametric inequality. For $0\le a\le i$ the
numbers $\abs{a-b}$, $b\le i$, $b\ne a$, are $1,\dots,a$ and $1,\dots,i-a$, and the exponent of $2$ in a factorial is at
most its argument; so
\[
 \prod_{b\le i,\ b\ne a}\abs{\xi_a-\xi_b}^{-1}=\abs\varpi^{-v(a!)-v((i-a)!)}\le\abs\varpi^{-i}.
\]
Hence the divided difference $\sum_{a\le i}\psi(\xi_a)\prod_{b\le i,\,b\ne a}(\xi_a-\xi_b)^{-1}$ of order $i$ of a map
$\psi$ from $K$ to a normed space is a combination of the $i+1$ values $\psi(\xi_a)$, $a\le i$, with coefficients of
absolute value at most $\abs\varpi^{-i}$.

\emph{Step 2 (item (1): the bound).} If $m=0$, $b$ is constant and the bound is clear. Arguments $w_c$ with $\alpha_c=0$ do
not change $b^{[\alpha]}(w)$, so we may assume $\alpha_c\ge1$ for all $c$; then $r\le m$. If $r=1$, then
$b^{[\alpha]}(w_1)=b(w_1)$. Let $r\ge2$ and $N=r-1$. The map $\Phi(s)=b(s_1w_1+\dots+s_Nw_N+w_r)$ on $K^N$ is a polynomial
of total degree at most $m$ with coefficients in $Y$, and $b^{[\alpha]}(w)$ is its coefficient of
$s_1^{\alpha_1}\cdots s_N^{\alpha_N}$. Newton's interpolation formula on the nodes $\xi_0,\dots,\xi_m$ in each variable gives
\[
 \Phi(s)=\sum_{i\in\{0,\dots,m\}^N}\delta_i\prod_{c=1}^Nq_{i_c}(s_c),\qquad q_i(s)=\prod_{b<i}(s-\xi_b),
\]
where $\delta_i\in Y$ is the divided difference of $\Phi$ of order $i_c$ in the variable $s_c$, on the nodes
$\xi_0,\dots,\xi_{i_c}$, for $c=1,\dots,N$. A divided difference of order $i$ of a polynomial of degree less than $i$
vanishes, and the monomials of $\Phi$ have total degree at most $m$, so $\delta_i=0$ unless $\abs i\le m$. By Step~1,
$\delta_i$ is a combination of at most $\prod_c(i_c+1)\le2^{\abs i}$ values of $\Phi$ at points with coordinates $\xi_a$,
with coefficients of absolute value at most $\abs\varpi^{-\abs i}$, and each of these values has norm at most
$\norm b(\norm{w_1}+\dots+\norm{w_r})^m$, since $\abs{\xi_a}\le1$. The polynomials $q_i$ have coefficients of absolute
value at most $1$. So the coefficient of $s_1^{\alpha_1}\cdots s_N^{\alpha_N}$ is a sum of at most
$\binom{m+N}N\le4^m$ terms, one for each $i$ with $\abs i\le m$, each of norm at most
$(2/\abs\varpi)^m\norm b(\norm{w_1}+\dots+\norm{w_r})^m$.

\emph{Step 3 (item (1): expansions).} Let $x,y\in W$ and $0\le i\le m$. If $x,y\ne0$, choose $\lambda,\mu\in K$ with
$\abs\varpi<\abs\lambda\norm x\le1$ and $\abs\varpi<\abs\mu\norm y\le1$. As
$b^{[(m-i,i)]}(\lambda x,\mu y)=\lambda^{m-i}\mu^ib^{[(m-i,i)]}(x,y)$, Step~2 gives
\begin{equation}\label{dp:eq:binomial}
 \norm{b^{[(m-i,i)]}(x,y)}\le(16/\abs\varpi^2)^m\norm b\,\norm x^{m-i}\norm y^i,
\end{equation}
which holds trivially if $x=0$ or $y=0$. Since $b(x+y)=\sum_{i\le m}b^{[(m-i,i)]}(x,y)$, $b$ is continuous. For fixed
$x$, the map $y\mapsto b^{[(m-i,i)]}(x,y)$ is a bounded polynomial law of degree $i$: on the span of finitely many
vectors it is a coefficient of a polynomial with coefficients in $Y$, and~\eqref{dp:eq:binomial} bounds it. So
$b(x+y)=\sum_{i\le m}b^{[(m-i,i)]}(x,y)$ is a finite expansion at $x$, and $b$ is DP on $W$.

\emph{Step 4 (item (2)).} Let $f(x_0+h)=\sum_na_n(h)$ for $\norm h<R$ and $g(f(x_0)+w)=\sum_mb_m(w)$ for $\norm w<R'$, with
$\norm{a_n}\le MR^{-n}$ and $\norm{b_m}\le M'R'^{-m}$, where $M,M'>0$. Put
$\sigma=\min(R/2,\abs\varpi RR'/(32M))$ and $\Delta(h)=\sum_{n\ge1}a_n(h)=f(x_0+h)-f(x_0)$. For $\norm h\le\sigma$,
$\sum_{n\ge1}\norm{a_n(h)}\le2M\sigma/R\le\abs\varpi R'/16$. For a finitely supported sequence $\alpha=(\alpha_n)_{n\ge1}$
of nonnegative integers put $\abs\alpha=\sum_n\alpha_n$, $\deg\alpha=\sum_nn\alpha_n$ and
$\phi_\alpha(h)=b_{\abs\alpha}^{[\alpha]}(a_1(h),a_2(h),\dots)$, computed with $a_1(h),\dots,a_L(h)$ for any $L$ at least
the largest index in the support of $\alpha$. On the span of finitely many vectors the $a_n$ are polynomials whose
coefficients span a finite-dimensional subspace of $W$, on which $b_{\abs\alpha}^{[\alpha]}$ is a polynomial in the
coordinates of its arguments with coefficients in $Y$; and $\phi_\alpha(\lambda h)=\lambda^{\deg\alpha}\phi_\alpha(h)$. With~(1),
$\phi_\alpha$ is a bounded polynomial law of degree $\deg\alpha$, and for $\norm h\le\sigma$
\[
 \norm{\phi_\alpha(h)}\le(8/\abs\varpi)^{\abs\alpha}M'R'^{-\abs\alpha}(2M\sigma/R)^{\abs\alpha}\le M'2^{-\abs\alpha}.
\]
By homogeneity, $\norm{\phi_\alpha(h)}\le M'(\norm h/(\abs\varpi\sigma))^{\deg\alpha}$ for all $h$. At most $2^{d-1}$ sequences
$\alpha$ have $\deg\alpha=d\ge1$, one for each partition of $d$. So for $\norm h\le\tau=\abs\varpi\sigma/4$ the family
$(\phi_\alpha(h))_\alpha$ is absolutely summable in the completion $\widehat Y$. For each $L$,
$b_m(\sum_{n\le L}a_n(h))$ is the sum of the $\phi_\alpha(h)$ with $\abs\alpha=m$ and support in $\{1,\dots,L\}$; as
$L\to\infty$ the left side tends to $b_m(\Delta(h))$, by continuity of $b_m$. Summing over $m$, the family has sum
$\sum_mb_m(\Delta(h))=g(f(x_0+h))$, since $\norm{\Delta(h)}<R'$. Grouped by degree instead, it has the finite sums
$\theta_d(h)=\sum_{\deg\alpha=d}\phi_\alpha(h)$, bounded polynomial laws of degree $d$ with values in $Y$ and
$\norm{\theta_d}\le2^{d-1}M'(\abs\varpi\sigma)^{-d}$ for $d\ge1$. So $\sum_d\theta_d(h)$ converges in $\widehat Y$ to
$g(f(x_0+h))\in Y$ for $\norm h\le\tau$; its partial sums lie in $Y$, so it converges in $Y$, and $g\circ f$ is DP at
$x_0$.

\emph{Step 5 (item (3)).} Let $ja(h)=\beta(h,\dots,h)$ as in Definition~\ref{dp:def:laws}(2). For $h_1,\dots,h_r\in P$,
the map $s\mapsto\beta(\sum_cs_ch_c,\dots,\sum_cs_ch_c)$ is a homogeneous polynomial of degree $n$ with coefficients in
$Z$. As $K$ is infinite, its coefficients are linear combinations of its values, by interpolation in each variable, and
these values lie in $j(W)$. So $a$ is a bounded polynomial law of degree $n$, with $\norm a\le\norm\beta$; in particular
analytic maps are DP. Now let $T\colon K^r\to P$ be a bounded linear isomorphism, $K^r$ with the maximum norm and unit
vectors $e_1,\dots,e_r$, and let $a\colon P\to W$ be a bounded polynomial law of degree $n$. Then
$a(Tz)=\sum_{\abs\gamma=n}v_\gamma z^\gamma$ with $v_\gamma=(a\circ T)^{[\gamma]}(e_1,\dots,e_r)$, and
$\norm{v_\gamma}\le(8r/\abs\varpi)^n\norm T^n\norm a$ by~(1). The $n$\nobreakdash-linear map
\[
 \beta_a(h^1,\dots,h^n)=\sum_{t_1\le\dots\le t_n}v_{\gamma(t)}\prod_{c=1}^n(T^{-1}h^c)_{t_c},
\]
summed over the nondecreasing words $t$ in $\{1,\dots,r\}$, where $\gamma(t)$ counts the letters of $t$, has diagonal
$a$, since each monomial $z^\gamma$ comes from exactly one such word, and
$\norm{\beta_a}\le\binom{n+r-1}n(8r\norm T\norm{T^{-1}}/\abs\varpi)^n\norm a\le2^{r-1}(16r\norm T\norm{T^{-1}}/\abs\varpi)^n\norm a$.
Applied to the coefficients of a DP expansion, this bound, geometric in $n$, gives an analytic expansion with a
smaller positive radius.

\emph{Step 6 (item (4)).} By Proposition~\ref{prop:ambient-polynomial} and its proof, the restriction of $B$ to
alternating maps is the diagonal of a bounded $(k+1)$\nobreakdash-linear map of norm at most $1$ with values in
$\cL(\Alt^k(D;F),\Mult^k(E;F'))$, and its values lie in $\cL(\Alt^k(D;F),\Alt^k(E;F'))$, which postcomposition by
$j_{E,F'}$ embeds isometrically; the same holds for $A^k_{E,D;F}(u)=B(u,\id_F)$ in degree $k$. By
Lemma~\ref{ob:lem:metric-surjection}, $\Theta\mapsto\Theta\circ\wedge^k_E$ is a linear isometry of
$\cL(\Lambda^k_\pi E,\widehat\Lambda^k_\pi D)$ into $\cL(\otimes^k_\pi E,\widehat\Lambda^k_\pi D)$, and restriction to the dense
subspace $\Lambda^k_\pi E$ is an isometric isomorphism of $\cL(\widehat\Lambda^k_\pi E,\widehat\Lambda^k_\pi D)$ onto
$\cL(\Lambda^k_\pi E,\widehat\Lambda^k_\pi D)$, as $\widehat\Lambda^k_\pi D$ is complete. These isometries map
$\Lambda^k_\pi u$ and $\widehat\Lambda^k_\pi u$ to $\wedge^k_D\circ(u\otimes\dots\otimes u)$, the diagonal of the
$k$\nobreakdash-linear map $(u_1,\dots,u_k)\mapsto\wedge^k_D\circ(u_1\otimes\dots\otimes u_k)$ of norm at most $1$. Item~(3)
bounds the four norms by $1$.
\end{proof}

For the proof of Proposition~\ref{dp:prop:nonambient} we use the following notation. For a field $\Bbbk$ and integers
$m,n\ge1$, read a tensor $t\in(\Bbbk^m)^{\otimes n}$ as the function
$t(i_1,\dots,i_n)$ on $[m]^n$, where $[m]=\{1,\dots,m\}$, and for linear forms $\phi_1,\dots,\phi_n$ on $\Bbbk^m$ let $t(\phi_1,\dots,\phi_n)$ be
the value of $\phi_1\otimes\dots\otimes\phi_n$ at $t$. Let $\Gamma^n(\Bbbk^m)$ be the space of symmetric tensors, let
$\operatorname{diag}(t)\in\Bbbk^m$ have the entries $t(i,\dots,i)$, and let the \emph{flattening rank}
$\operatorname{frk}(t)$ be the rank of the $m\times m^{n-1}$ matrix of the entries $t(i_1,\dots,i_n)$, with rows indexed by
$i_1$; it is subadditive and does not change under extension of scalars. A \emph{projection onto the symmetric tensors}
is a linear map $\Psi$ of $(\Bbbk^m)^{\otimes n}$ into $\Gamma^n(\Bbbk^m)$ with $\Psi(t)=t$ for $t\in\Gamma^n(\Bbbk^m)$.
Symmetric tensors with entries in a commutative $\Bbbk$\nobreakdash-algebra are combinations of symmetric tensors over
$\Bbbk$, so the linear extension of $\Psi$ to such entries maps tensors with these entries to symmetric ones and fixes
the symmetric ones.

\begin{lemma}[projections onto symmetric tensors]\label{dp:lem:projections}
Let $m\ge1$.
\begin{enumerate}
\item Let $\Bbbk$ be a field, $a\ge1$, $x=(x_1,\dots,x_m)^{\mathsf T}$ a column of indeterminates, and $\mathsf M(x)$
an $m\times m$ matrix over $\Bbbk[x_1,\dots,x_m]$ with $\mathsf M(x)x=(x_1^a,\dots,x_m^a)^{\mathsf T}$. Then $\det\mathsf M(x)$ has the coefficient $1$ at
$x_1^{a-1}\cdots x_m^{a-1}$.
\item Let $\Bbbk$ have characteristic $p>0$, let $C\ge0$ be an integer with $\abs\Bbbk\ge2C+2$, and let
$\mathcal X\subseteq\Gamma^p(\Bbbk^m)$ be a linear space whose elements have flattening rank at most $C$. Then
$\dim\operatorname{diag}(\mathcal X)\le C$.
\item Let $f$ be a function from an abelian group into $[0,\infty)$ with $f(z+z')\le f(z)+f(z')$, and let $\Phi$ be a
map from $(\F_p^m)^p$ into this group that is additive in each argument. Then the average of $f\circ\Phi$ over
$(\F_p^m)^p$ is at least $2^{-p}$ times its value at every point.
\item Let $\Bbbk$ be a field of characteristic $p>0$, and let $\Psi$ be a projection of $(\Bbbk^m)^{\otimes p}$ onto
$\Gamma^p(\Bbbk^m)$. Then the average of $\operatorname{frk}\Psi(y^1\otimes\dots\otimes y^p)$ over
$(y^1,\dots,y^p)\in(\F_p^m)^p$ is at least $m/\bigl(2^p(1+\lfloor\log_pm\rfloor)\bigr)$.
\end{enumerate}
\end{lemma}

\begin{proof}
\emph{Step 1 (item (1)).} Let $\mathsf D$ be the diagonal matrix with the entries $x_1^{a-1},\dots,x_m^{a-1}$, and put
$\mathsf N=\mathsf M-\mathsf D$, so $\mathsf Nx=0$. As the determinant is linear in each row,
$\det\mathsf M=\sum_{S\subseteq[m]}\det\mathsf M_S$, where $\mathsf M_S$ has the rows of $\mathsf N$ in the positions in
$S$ and those of $\mathsf D$ elsewhere, and $\det\mathsf M_\emptyset=x_1^{a-1}\cdots x_m^{a-1}$. Let $S\ne\emptyset$, and
let $J_S$ be the ideal generated by the $x_i^a$, $i\notin S$. The vector $\mathsf M_Sx$ has the entries $0$ in the
positions in $S$ and $x_i^a$ elsewhere, so multiplication by the adjugate of $\mathsf M_S$ gives
$\det(\mathsf M_S)x_j\in J_S$ for every $j$. For $j\in S$, $x_j$ is a nonzerodivisor modulo $J_S$, which is generated
by monomials in the other variables; so $\det\mathsf M_S\in J_S$. Every monomial of an element of $J_S$ is divisible by
some $x_i^a$, so $x_1^{a-1}\cdots x_m^{a-1}$ does not occur in $\det\mathsf M_S$.

\emph{Step 2 (item (2): a common kernel).} For $t\in\Gamma^p(\Bbbk^m)$ and a linear form $\phi$ on $\Bbbk^m$ let
$t^\flat(\phi)\in(\Bbbk^m)^{\otimes(p-1)}$ have the entries $\sum_i\phi(e_i)\,t(i,i_2,\dots,i_p)$, $e_i$ the unit
vectors; the linear map $t^\flat$ has rank $\operatorname{frk}(t)$, and $t^\flat(\phi)(\phi,\dots,\phi)=t(\phi,\dots,\phi)$.
Let $t_0\in\mathcal X$ have the largest flattening rank $r_0\le C$, let $\mathcal U$ be the kernel of $t_0^\flat$, of
codimension $r_0$, with a complement $\mathcal U'$, and let $\mathcal I$ be the image of $t_0^\flat$, with a complement
$\mathcal I'$. For $t\in\mathcal X$ write $t^\flat$ in blocks $\mathsf B_{11}(t)\colon\mathcal U'\to\mathcal I$,
$\mathsf B_{12}(t)\colon\mathcal U\to\mathcal I$, $\mathsf B_{21}(t)\colon\mathcal U'\to\mathcal I'$ and
$\mathsf B_{22}(t)\colon\mathcal U\to\mathcal I'$; then $\mathsf B_{11}(t_0)$ is invertible and the other blocks of
$t_0^\flat$ vanish. Put $\mathsf A_\mu=\mathsf B_{11}(t_0)+\mu\mathsf B_{11}(t)$ for $\mu\in\Bbbk$. If $\mu\ne0$ and
$\det\mathsf A_\mu\ne0$, the Schur complement of $\mathsf A_\mu$ in $(t_0+\mu t)^\flat$ vanishes, since
$\operatorname{frk}(t_0+\mu t)\le r_0$; that is,
\begin{equation}\label{dp:eq:schur}
 \det(\mathsf A_\mu)\,\mathsf B_{22}(t)=\mu\,\mathsf B_{21}(t)\operatorname{adj}(\mathsf A_\mu)\mathsf B_{12}(t).
\end{equation}
Both sides are polynomials in $\mu$ of degree at most $r_0$, and $\det\mathsf A_\mu$ has at most $r_0$ roots, since it is
nonzero at $\mu=0$. So~\eqref{dp:eq:schur} holds for all $\mu$ outside a set of at most $r_0+1$ elements, hence, as
$\abs\Bbbk\ge2r_0+2$, for at least $r_0+1$ values of $\mu$, and so identically. At $\mu=0$, \eqref{dp:eq:schur} gives
$\mathsf B_{22}(t)=0$: $t^\flat(\phi)\in\mathcal I$ for $\phi\in\mathcal U$.

\emph{Step 3 (item (2): diagonal forms).} The elements of $\mathcal I$ vanish when one of their arguments lies in
$\mathcal U$, as $t_0$ is symmetric. So $\delta_t(\phi)=t(\phi,\dots,\phi)=t^\flat(\phi)(\phi,\dots,\phi)$
vanishes for $\phi\in\mathcal U$ and $t\in\mathcal X$. As $t$ is symmetric and the binomial coefficients
$\binom p1,\dots,\binom p{p-1}$ vanish in $\Bbbk$, $\delta_t$ is additive, and $\delta_t(c\phi)=c^p\delta_t(\phi)$ for
$c\in\Bbbk$. Let $\phi_1,\dots,\phi_{r_0}$ be a basis of $\mathcal U'$, and write the coordinate forms as
$e_i^*=\psi_i+\sum_jc_{ij}\phi_j$ with $\psi_i\in\mathcal U$. Then $\operatorname{diag}(t)$ has the entries
$\delta_t(e_i^*)=\sum_jc_{ij}^p\delta_t(\phi_j)$, so it is the image of the vector $(\delta_t(\phi_j))_j\in\Bbbk^{r_0}$,
which depends linearly on $t$, under a matrix that does not depend on $t$. Hence
$\dim\operatorname{diag}(\mathcal X)\le r_0\le C$.

\emph{Step 4 (item (3)).} For $z$ in a finite subgroup $B$ of the group, the average of $f$ over $B$ is at least
$f(z)/2$, since $f(b)+f(z-b)\ge f(z)$ and $b\mapsto z-b$ permutes $B$. For fixed other arguments,
$y^j\mapsto\Phi(y^1,\dots,y^p)$ is additive on $\F_p^m$, so as $y^j$ runs through $\F_p^m$ its value runs equally often
through each element of its image, a finite subgroup. Hence the average of $f\circ\Phi$ over $y^j$ is at least half its
value at any fixed $y^j$. Applying this in the slots $j=1,\dots,p$ in turn gives~(3).

\emph{Step 5 (item (4): an invertible diagonal).} Let $x=(x_1,\dots,x_m)$ be a vector of indeterminates, extend $\Psi$
linearly to entries in $\Bbbk[x_1,\dots,x_m]$, and let $\mathsf M(x)$ be the matrix over $\Bbbk[x_1,\dots,x_m]$ of the
linear map $v\mapsto\operatorname{diag}\Psi(x^{\otimes(p-1)}\otimes v)$, whose columns are
$\operatorname{diag}\Psi(x^{\otimes(p-1)}\otimes e_i)$. Its value at a point $\xi\in\Bbbk^m$ is the matrix of the same
map with $\xi$ in place of $x$. As $\Psi$ fixes the symmetric tensor $x^{\otimes p}$,
$\mathsf M(x)x=\operatorname{diag}(x^{\otimes p})=(x_1^p,\dots,x_m^p)^{\mathsf T}$. By~(1), the form
$\det\mathsf M(x)$ of degree $m(p-1)$ has the coefficient $1$ at $x_1^{p-1}\cdots x_m^{p-1}$. The sum
$\sum_{c\in\F_p}c^b$ is $-1$ if $b>0$ and $p-1$ divides $b$, and $0$ otherwise, and $x_1^{p-1}\cdots x_m^{p-1}$ is the
only monomial of degree $m(p-1)$ whose exponents are positive multiples of $p-1$; so
$\sum_{\xi\in\F_p^m}\det\mathsf M(\xi)=(-1)^m$, and $\mathsf M(\xi_0)$ is invertible for some $\xi_0\in\F_p^m$.

\emph{Step 6 (item (4): a value of large rank).} Let $\Bbbk''\supseteq\Bbbk$ be a field with $\abs{\Bbbk''}\ge2m$,
with $\Bbbk''=\Bbbk$ if $\Bbbk$ is infinite, and let $\Psi''$ be the linear extension of $\Psi$ to entries in $\Bbbk''$.
The tensors $\Psi''(\xi_0^{\otimes(p-1)}\otimes u)$, $u\in\Bbbk''^m$, form a linear space whose diagonals form
$\mathsf M(\xi_0)\Bbbk''^m=\Bbbk''^m$, so by~(2) with $C=m-1$ one of them has flattening rank $m$. Hence some $m\times m$
minor $\mathsf P(u)$ of the matrix of entries of $\Psi''(\xi_0^{\otimes(p-1)}\otimes u)$, a polynomial of degree at most
$m$ in each coordinate of $u$, is not the zero polynomial. Let $d=1+\lfloor\log_pm\rfloor$, the least integer with
$p^d>m$. As $\abs{\Bbbk''}>m$, and $\abs{\Bbbk''}$ is a power of $p$ if $\Bbbk''$ is finite, $\Bbbk''$ contains an
$\F_p$\nobreakdash-subspace $\mathcal W$ of dimension $d$, so $\abs{\mathcal W}=p^d>m$. A nonzero polynomial in $r$
variables of degree at most $m$ in each variable does not vanish on all of $\mathcal W^r$: for $r=1$ it has at most
$m$ roots, and for $r>1$ write it as $\sum_{e=0}^m\mathsf P_eu_r^e$ with polynomials $\mathsf P_e$ in the other variables;
some $\mathsf P_e$ is nonzero, by induction it does not vanish at some point of $\mathcal W^{r-1}$, and there the
polynomial is a nonzero polynomial in $u_r$ of degree at most $m$. With $r=m$ and the minor $\mathsf P$, there is
$v\in\mathcal W^m$ with $\operatorname{frk}\Psi''(\xi_0^{\otimes(p-1)}\otimes v)=m$. Write
$v=\sum_{i=1}^d\omega_iv_i$ with an $\F_p$\nobreakdash-basis $\omega_1,\dots,\omega_d$ of $\mathcal W$ and $v_i\in\F_p^m$.
Then $\Psi''(\xi_0^{\otimes(p-1)}\otimes v)=\sum_i\omega_i\Psi(\xi_0^{\otimes(p-1)}\otimes v_i)$, and subadditivity gives
$\operatorname{frk}\Psi(\xi_0^{\otimes(p-1)}\otimes v_{i_0})\ge m/d$ for some $i_0$.

\emph{Step 7 (item (4): averaging).} The flattening rank is subadditive, and
$(y^1,\dots,y^p)\mapsto\Psi(y^1\otimes\dots\otimes y^p)$ is additive in each argument, so by~(3) the average in~(4) is
at least $2^{-p}$ times the flattening rank of the value at $(\xi_0,\dots,\xi_0,v_{i_0})$, hence at least $m/(2^pd)$.
\end{proof}

\begin{proof}[Proof of Proposition~\ref{dp:prop:nonambient}]
\emph{Step 0 (the field under (a)).} Assume~(a). The completion $\Kh$ has characteristic $p$, so its valuation ring is a
complete discrete valuation ring containing the prime field $\F_p$, with $\vartheta$ as a uniformizer. By Cohen's
structure theorem for complete regular local rings containing a field~\cite{Co}, this ring is isomorphic to the power
series ring $\kappa[[X]]$ over its residue field $\kappa$, by an isomorphism taking $\vartheta$ to $X$; it extends to an
isomorphism of $\Kh$ onto $\kappa((X))$, which is isometric for the absolute value $\abs a=\abs\vartheta^{\ord a}$ on
$\kappa((X))$. All objects in Proposition~\ref{dp:prop:nonambient} are defined from $K$ as a normed field, so under~(a)
we may assume that $K$ is a subfield of $\kappa((X))$ with this absolute value.

\emph{Step 1 (item (1): $t_m$).} Let $(f_k)$ be a basis of $Q_m$ with dual basis $(f_k^*)$, and let $e_l$ be the unit
vector of $\ell^1(I_m,K)$ at $l$, so that $\varepsilon_l$ is its image in $G_m$. Then $\sum_l\bar y_l\otimes e_l=\sum_kf_k\otimes(f_k^*(\bar y_l))_l$ lies in
$Q_m\otimes N_m$, so $\sum_l\bar y_l\otimes\varepsilon_l=0$ in $Q_m\otimes G_m$. Hence $t_m$ lies in the kernel of the
map $z\otimes g\mapsto\bar z\otimes g$ of $(K^m)^{\otimes p}\otimes G_m$ onto $Q_m\otimes G_m$, which is
$\Gamma^p(K^m)\otimes G_m$. As $K$ is infinite, $\Gamma^p(K^m)$ is spanned by the powers $u^{\otimes p}$, so $t_m\in V$.
The vectors $y_l^1,\dots,y_l^p$ have entries in $\F_p$, hence norm at most $1$, and $\norm{\varepsilon_l}\le1$, so
$\pi(t_m)\le\abs{I_m}=p^{pm}$.

\emph{Step 2 (a representation gives a projection).} Let $t_m=\sum_ic_iu_i^{\otimes p}\otimes g_i$ with $c_i\in K$,
$u_i\in c_0$, $g_i\in G$, and let $\delta>0$. The coordinate projection of $c_0$ onto $K^m$ and the projection of $G$
onto $G_m$ are contractive; the $p$-th tensor power of the first, tensored with the second, fixes $t_m$ and does not
increase $\sum_i\abs{c_i}\norm{u_i}^p\norm{g_i}$, so we may assume $u_i\in K^m$ and $g_i\in G_m$. Choose
$\hat g_i\in\ell^1(I_m,K)$ with image $g_i$ and $\norm{\hat g_i}\le(1+\delta)\norm{g_i}$, and put
$\sigma_l=\sum_ic_i\hat g_i(l)\,u_i^{\otimes p}\in\Gamma^p(K^m)$ for $l\in I_m$. With $\pi_m$ the projective norm of the
$p$\nobreakdash-fold tensor power of $K^m$ with the maximum norm,
\begin{equation}\label{dp:eq:norm-sum}
 \sum_{l\in I_m}\pi_m(\sigma_l)\le\sum_i\abs{c_i}\norm{u_i}^p\norm{\hat g_i}\le(1+\delta)\sum_i\abs{c_i}\norm{u_i}^p\norm{g_i}.
\end{equation}
The tensors $\sum_l\sigma_l\otimes e_l=\sum_ic_iu_i^{\otimes p}\otimes\hat g_i$ and $\sum_ly_l\otimes e_l$ both map to
$t_m$ in $(K^m)^{\otimes p}\otimes G_m$, so their difference lies in $(K^m)^{\otimes p}\otimes N_m$. The vectors
$(f_k^*(\bar y_l))_l$ span $N_m$, so the difference is $\sum_k\tau_k\otimes(f_k^*(\bar y_l))_l$ for some
$\tau_k\in(K^m)^{\otimes p}$. So $\sigma_l=y_l+\sum_kf_k^*(\bar y_l)\tau_k$ for $l\in I_m$. Put
$\Psi_K(z)=z+\sum_kf_k^*(\bar z)\tau_k$ for $z\in(K^m)^{\otimes p}$; then $\Psi_K(y_l)=\sigma_l$. The linear map
$\Theta(q)=q+\sum_kf_k^*(q)\bar\tau_k$ of $Q_m$
vanishes at every $\bar y_l$, since $\bar\sigma_l=0$, and these span $Q_m$, as $I_m$ contains the tuples of unit
vectors; so $\Theta=0$, and the image of $\Psi_K(z)$ in $Q_m$ is $\Theta(\bar z)=0$. As $\Psi_K(z)=z$ when $\bar z=0$,
$\Psi_K$ is a projection of $(K^m)^{\otimes p}$ onto $\Gamma^p(K^m)$.

\emph{Step 3 (item (1) under (a): constant coefficients).} Let $K\subseteq\kappa((X))$ be as in Step~0. By the
paragraph before Lemma~\ref{dp:lem:projections}, with $\Bbbk=K$, the
linear extension $\widetilde\Psi$ of $\Psi_K$ to entries in $\kappa((X))$ is a projection of
$(\kappa((X))^m)^{\otimes p}$ onto $\Gamma^p(\kappa((X))^m)$. For $z\in(\kappa^m)^{\otimes p}$ put
$\Psi(z)=\coeff_0\widetilde\Psi(z)$. Then $\Psi$ is $\kappa$\nobreakdash-linear with symmetric values and fixes the symmetric
tensors over $\kappa$, so it is a projection of $(\kappa^m)^{\otimes p}$ onto $\Gamma^p(\kappa^m)$. As the $y_l$ have entries in $\F_p\subseteq K\cap\kappa$, $\widetilde\Psi(y_l)=\Psi_K(y_l)=\sigma_l$, and
$\Psi(y_l)=\coeff_0\sigma_l$. The matrix of the entries of
$\sigma_l$, with rows indexed by the first index, read in the completion of
$\linf([m],K)\otimes_\pi\linf([m]^{p-1},K)$, has norm at most $\pi_m(\sigma_l)$, since an elementary tensor of $p-1$ vectors
of $K^m$, read as a function on $[m]^{p-1}$, has as maximum norm the product of their maximum norms. So
Lemma~\ref{fib:lem:coeff}(2), which holds for every field $\kappa$ of characteristic $p$ (see the beginning of
Appendix~\ref{app:forcing}), with $\abs\vartheta$ in place of $\varrho$ and the index sets $[m]$ and $[m]^{p-1}$, gives
$\operatorname{frk}\Psi(y_l)\le M_{\abs\vartheta}\pi_m(\sigma_l)$, and by~\eqref{dp:eq:norm-sum}
\begin{equation}\label{dp:eq:rank-sum}
 \sum_{l\in I_m}\operatorname{frk}\Psi(y_l)\le(1+\delta)M_{\abs\vartheta}\sum_i\abs{c_i}\norm{u_i}^p\norm{g_i}.
\end{equation}

\emph{Step 4 (item (1) under (a)).} The tensors $y_l$, $l\in I_m$, are the $y^1\otimes\dots\otimes y^p$ with
$(y^1,\dots,y^p)\in(\F_p^m)^p$, so by Lemma~\ref{dp:lem:projections}(4) the left side of~\eqref{dp:eq:rank-sum} is at
least $mp^{pm}/\bigl(2^p(1+\lfloor\log_pm\rfloor)\bigr)$. As the representation of $t_m$ and $\delta>0$ are arbitrary,
$\nu(t_m)\ge mp^{pm}/\bigl(2^pM_{\abs\vartheta}(1+\lfloor\log_pm\rfloor)\bigr)$. With Step~1 this is~(1) under~(a).

\emph{Step 5 (item (1) under (b)).} Let $p=2$. For $S\subseteq[m]$ let $1_S\in\F_2^m$ have the entries $1$ at the
elements of $S$ and $0$ elsewhere, so that $I_m$ consists of the pairs $l=(1_S,1_T)$, $S,T\subseteq[m]$. Identify
$(K^m)^{\otimes2}$ with $K^{m\times m}$ by $a\otimes b\mapsto ab^{\mathsf T}$; then $\Gamma^2(K^m)$ is the space of
symmetric matrices and $y_l=1_S1_T^{\mathsf T}$ for $l=(1_S,1_T)$. First, $\abs{\det\sigma_{l_0}}\ge1$ for some
$l_0\in I_m$. If $m=1$, then $\Psi_K=\id$ and $\sigma_{l_0}=1$ for $l_0=(1_{\{1\}},1_{\{1\}})$. Let $m\ge2$, and let
$\wedge\colon K^{m\times m}\to\Lambda^2K^m$ be the linear map of Theorem~\ref{obs:thm:char-two}(1), with $m$ in place of
its $n$; by its formula there, it is onto, its kernel is the space of symmetric matrices, and
$\wedge(ab^{\mathsf T})=a\wedge b$. As $\Psi_K-\id$ vanishes on the symmetric matrices, it is $\chi\circ\wedge$ for a
linear map $\chi\colon\Lambda^2K^m\to K^{m\times m}$, and $\wedge\circ\chi=\id$, since $\Psi_K$ takes symmetric values
and $-1=1$ in $K$. So $\Psi_K(ab^{\mathsf T})=\chi(a\wedge b)+ab^{\mathsf T}$, and Proposition~\ref{ocl:prop:sections},
applied over the field $K$ to $\chi$, with $m$ in place of its $n$, gives
$\sum_{S,T\subseteq[m]}\det\Psi_K(1_S1_T^{\mathsf T})=1$. As $K$ is nonarchimedean, some term has absolute value at
least $1$.

By Lemma~\ref{ocl:lem:pfaffian-bound}(2), $\pi_m(\sigma_{l_0})\ge m$. By Lemma~\ref{dp:lem:projections}(3), applied to the
norm $\pi_m$ and to the map $(y^1,y^2)\mapsto\Psi_K(y^1\otimes y^2)$, the average of $\pi_m(\sigma_l)$ over the $4^m$
elements $l\in I_m$ is at least $\pi_m(\sigma_{l_0})/4$, so $\sum_{l\in I_m}\pi_m(\sigma_l)\ge m4^{m-1}$. As the representation of $t_m$ and $\delta>0$ are arbitrary,
\eqref{dp:eq:norm-sum} gives $\nu(t_m)\ge m4^{m-1}$. With Step~1 this is~(1) under~(b).

\emph{Step 6 (item (2)).} On the span of finitely many vectors of $H$, $R_n$ is a homogeneous polynomial of degree $n$
with coefficients in $V$: the coefficients of $x^{\otimes p}$, as a polynomial in the coordinates of $x$ in a
finite-dimensional subspace of $c_0$, are symmetric tensors, and symmetric tensors in $c_0^{\otimes p}$ are combinations
of powers, as $K$ is infinite. Moreover
$\nu(R_n(\lambda,x,g))\le\abs\lambda^{n-p-1}\norm x^p\norm g\le\norm{(\lambda,x,g)}^n$. So $R_n$ is a bounded polynomial
law of degree $n$ with $\norm{R_n}\le1$. Suppose $jR_n(h)=\beta(h,\dots,h)$ for a linear isometry $j$ of $V$, or of its
completion, into a normed space $Z$ and a bounded $n$\nobreakdash-linear $\beta\colon H^n\to Z$, and let
$\beta_\otimes\colon\otimes^n_\pi H\to Z$ be the linear map with $\beta_\otimes(h_1\otimes\dots\otimes
h_n)=\beta(h_1,\dots,h_n)$, of norm $\norm\beta$. Put $\mathbf 1=(1,0,0)\in H$ and $d_n=n!/(p!\,(n-p-1)!)$. For
$t\in(K^m)^{\otimes p}\otimes G$, read in $\otimes^{p+1}H$, let $\iota(t)\in\otimes^nH$ be the sum, over the $d_n$ words
$w\in\{1,2,3\}^n$ with $p$ letters $1$ and one letter $2$, of the tensor obtained from
$t\otimes\mathbf 1^{\otimes(n-p-1)}$ by moving the first $p$ factors, in order, to the positions of the letters $1$ of
$w$, the $(p+1)$-st factor to the position of the letter $2$, and the factors $\mathbf 1$ to the positions of the
letters $3$. Since the inclusions of $c_0$ and $G$ into $H$ are isometric, permutations of the factors are isometries of
$\otimes^n_\pi H$ and $\norm{\mathbf 1}=1$, $\iota(t)$ has norm at most $d_n\pi(t)$ in $\otimes^n_\pi H$, with $\pi(t)$ the
projective norm of $t$ in $c_0^{\otimes p}\otimes G$. For $x\in K^m$, $g\in G$ and $s=(s_1,s_2,s_3)\in K^3$ put
$h_s=(s_3,s_1x,s_2g)$. Then $\beta_\otimes(h_s^{\otimes n})=\beta(h_s,\dots,h_s)=jR_n(h_s)=s_1^ps_2s_3^{n-p-1}\,
j(x^{\otimes p}\otimes g)$, and $\iota(x^{\otimes p}\otimes g)$ is the coefficient of $s_1^ps_2s_3^{n-p-1}$ in the
polynomial $h_s^{\otimes n}$. As $K$ is infinite, comparing coefficients gives
$\beta_\otimes(\iota(x^{\otimes p}\otimes g))=j(x^{\otimes p}\otimes g)$. So $\beta_\otimes(\iota(t))=j(t)$ for $t$ in the span of the
$x^{\otimes p}\otimes g$, $x\in K^m$, $g\in G$, which contains $t_m$ by Step~1. Hence
$\nu(t_m)=\norm{j(t_m)}\le d_n\norm\beta\pi(t_m)\le d_n\norm\beta p^{pm}$, and~(1) gives
$m\le2^pd_nM_{\abs\vartheta}\norm\beta(1+\lfloor\log_pm\rfloor)$ under~(a) and $m\le4d_n\norm\beta$ under~(b), for every $m$,
which is impossible.
\end{proof}

For the proof of Proposition~\ref{dp:prop:degree-p} we extend the notation before Lemma~\ref{dp:lem:projections} from
$\Bbbk^m$ to $\Bbbk^{\mathcal J}$, for a field $\Bbbk$ and a finite set $\mathcal J$. A tensor
$t\in(\Bbbk^{\mathcal J})^{\otimes d}$ is read as a function on the words $w=(w_1,\dots,w_d)\in\mathcal J^d$, and
$\Gamma^d(\Bbbk^{\mathcal J})$ is the space of symmetric tensors, with $\Gamma^0(\Bbbk^{\mathcal J})=\Bbbk$. Let
$\N^{\mathcal J}_d$ be the set of $\alpha\in\N^{\mathcal J}$ with $\abs\alpha=d$; for $\mathcal J=[m]$ it is the set
$\N^m_d$ of Section~\ref{sec:dp}. The \emph{content} of a word $w\in\mathcal J^d$ is the $\alpha\in\N^{\mathcal J}_d$
whose entry at each element of $\mathcal J$ is the number of letters of $w$ equal to it. Let
$O_\alpha\in\Gamma^d(\Bbbk^{\mathcal J})$ be the sum, over the words $w$ of content $\alpha$, of the tensor product of
the unit vectors of $\Bbbk^{\mathcal J}$ at the letters $w_1,\dots,w_d$. The $O_\alpha$, $\alpha\in\N^{\mathcal J}_d$,
form a basis of $\Gamma^d(\Bbbk^{\mathcal J})$, and $t=\sum_\alpha t_\alpha O_\alpha$, where $t_\alpha$ is the value of
$t$ at any word of content $\alpha$.

The \emph{product} of $t\in\Gamma^d(\Bbbk^{\mathcal J})$ and $t'\in\Gamma^{d'}(\Bbbk^{\mathcal J})$ is the sum
$t\cdot t'\in\Gamma^{d+d'}(\Bbbk^{\mathcal J})$ of the $\binom{d+d'}d$ tensors obtained from $t\otimes t'$ by the
permutations of the $d+d'$ factors that keep the order of the first $d$ factors and the order of the last $d'$; a
scalar in $\Gamma^0(\Bbbk^{\mathcal J})$ multiplies. For $\Bbbk=K$ this is the shuffle product of
Section~\ref{dgr:ssec:polarizable}. As tensor powers of linear maps commute with permutations of the
factors, $q^{\otimes(d+d')}(t\cdot t')=q^{\otimes d}(t)\cdot q^{\otimes d'}(t')$ for every linear map $q$ of
$\Bbbk^{\mathcal J}$ into some $\Bbbk^{\mathcal J'}$. Let $\mathcal J=[m]$, $\alpha\in\N^m_d$ and $\alpha'\in\N^m_{d'}$.
The coefficient of $O_\alpha\cdot O_{\alpha'}$ at a word of content $\alpha+\alpha'$ is the number of ways to choose
the $d$ of its letters that come from the first factor, which must contain $\alpha_i$ of the $\alpha_i+\alpha'_i$
letters equal to $i$, for each $i$; so
\begin{equation}\label{dp:eq:orbit-product}
 O_\alpha\cdot O_{\alpha'}=\prod_{i=1}^m\binom{\alpha_i+\alpha'_i}{\alpha_i}\,O_{\alpha+\alpha'},
\end{equation}
and the product is commutative and associative on the spaces $\Gamma^d(\Bbbk^m)$, $d\ge0$. For
$t\in\Gamma^{d+d'}(\Bbbk^m)$ and $\gamma\in\N^m_d$ let $\partial^\gamma t\in\Gamma^{d'}(\Bbbk^m)$ have the coefficients
$(\partial^\gamma t)_\alpha=t_{\alpha+\gamma}$, $\alpha\in\N^m_{d'}$; for $d=p-1$ and $d'=1$ this is the slice of
Section~\ref{sec:dp}. Then $\partial^\gamma\partial^{\gamma'}=\partial^{\gamma+\gamma'}$, and $\partial^\gamma O_\alpha$
is $O_{\alpha-\gamma}$ if $\gamma\le\alpha$ entrywise and $0$ otherwise. With~\eqref{dp:eq:orbit-product} and
Vandermonde's identity $\sum_\gamma\prod_i\binom{\alpha_i}{\gamma_i}=\binom{d+d'}d$ for $\alpha\in\N^m_{d+d'}$, the sum
over the $\gamma\in\N^m_d$ with $\gamma\le\alpha$, this gives
\begin{equation}\label{dp:eq:contraction-sum}
 \sum_{\gamma\in\N^m_d}(\partial^\gamma t)\cdot O_\gamma=\binom{d+d'}d\,t\qquad(t\in\Gamma^{d+d'}(\Bbbk^m)).
\end{equation}

\begin{lemma}[an identity for symmetric tensors]\label{dp:lem:linear}
Let $\Bbbk$ be a field of characteristic $p>0$ and $m\ge1$, let $\mathsf L$ be a matrix over $\Bbbk$ with rows and
columns indexed by $\N^m_{p-1}$, and let $\mathsf L^\sharp$ be the linear map of $\Gamma^p(\Bbbk^m)$ with
\[
 \mathsf L^\sharp(t)=\sum_{\gamma,\gamma'\in\N^m_{p-1}}\mathsf L_{\gamma\gamma'}\,(\partial^\gamma t)\cdot O_{\gamma'}.
\]
\begin{enumerate}
\item $t=\sum_{\gamma\in\N^m_{p-1}}\bigl((\id-\mathsf L^\sharp)(\partial^\gamma t)\bigr)\cdot O_\gamma$ for every
$t\in\Gamma^{2p-1}(\Bbbk^m)$.
\item $\rank(\id-\mathsf L^\sharp)\ge\binom{m+2p-2}{2p-1}\big/\binom{m+p-2}{p-1}\ge(p-1)!\,m^p/(2p-1)!$.
\end{enumerate}
\end{lemma}

\begin{proof}
\emph{Step 1 (the identity part).} By~\eqref{dp:eq:contraction-sum} with $d=p-1$ and $d'=p$,
$\sum_\gamma(\partial^\gamma t)\cdot O_\gamma=\binom{2p-1}{p-1}t$, the sum over $\N^m_{p-1}$. The integer
$\binom{2p-1}{p-1}(p-1)!$ is the product $(p+1)(p+2)\cdots(2p-1)$, which is congruent to $(p-1)!$ modulo $p$, and
$(p-1)!$ is prime to $p$; so $\binom{2p-1}{p-1}=1$ in $\Bbbk$, and $\sum_\gamma(\partial^\gamma t)\cdot O_\gamma=t$.

\emph{Step 2 (the part of $\mathsf L^\sharp$).} For $\gamma\in\N^m_{p-1}$,
$\mathsf L^\sharp(\partial^\gamma t)=\sum_{\gamma',\gamma''}\mathsf L_{\gamma'\gamma''}\,
(\partial^{\gamma'}\partial^\gamma t)\cdot O_{\gamma''}$. As the product is commutative and associative and
$\partial^{\gamma'}\partial^\gamma=\partial^\gamma\partial^{\gamma'}$,
\[
 \sum_\gamma\mathsf L^\sharp(\partial^\gamma t)\cdot O_\gamma
 =\sum_{\gamma',\gamma''}\mathsf L_{\gamma'\gamma''}\,O_{\gamma''}\cdot
 \sum_\gamma\bigl(\partial^\gamma(\partial^{\gamma'}t)\bigr)\cdot O_\gamma,
\]
all sums over $\N^m_{p-1}$. By~\eqref{dp:eq:contraction-sum} with $d=p-1$ and $d'=1$, applied to
$\partial^{\gamma'}t\in\Gamma^p(\Bbbk^m)$, the inner sum is $p\,\partial^{\gamma'}t=0$. With Step~1 this is~(1).

\emph{Step 3 (item (2)).} By~(1), the linear map
$t\mapsto\bigl((\id-\mathsf L^\sharp)(\partial^\gamma t)\bigr)_{\gamma\in\N^m_{p-1}}$ from $\Gamma^{2p-1}(\Bbbk^m)$ into
the product of $\abs{\N^m_{p-1}}$ copies of the image of $\id-\mathsf L^\sharp$ is injective. As
$\dim\Gamma^d(\Bbbk^m)=\abs{\N^m_d}=\binom{m+d-1}d$, this gives
$\binom{m+2p-2}{2p-1}\le\binom{m+p-2}{p-1}\rank(\id-\mathsf L^\sharp)$. The quotient of the two binomial coefficients is
$(p-1)!/(2p-1)!$ times the product of the $p$ numbers $m+p-1,\dots,m+2p-2$, hence at least $(p-1)!\,m^p/(2p-1)!$.
\end{proof}

\begin{proof}[Proof of Proposition~\ref{dp:prop:degree-p}]
Let $m\ge1$, and let $\mathcal J_m$ be the disjoint union of $[m]$ and $\N^m_{p-1}$. For a field $\Bbbk$ containing
$\F_p$ write $e_i$, $i\in[m]$, and $e_\gamma$, $\gamma\in\N^m_{p-1}$, for the unit vectors of $\Bbbk^{\mathcal J_m}$,
and read $\Bbbk^m$ as the span of the $e_i$. For $T\in\Gamma^p(\F_p^m)$ let
$q_T\colon\Bbbk^{\mathcal J_m}\to\Bbbk^m$ be the linear map with $q_T(e_i)=e_i$ and $q_T(e_\gamma)=\partial^\gamma T$;
its matrix has entries in $\F_p$. For a real number $\tau\ge1$ and $z=\sum_iz_ie_i+\sum_\gamma z_\gamma e_\gamma$ in
$K^{\mathcal J_m}$ put $\norm z_\tau=\sum_i\abs{z_i}+\tau\sum_\gamma\abs{z_\gamma}$; by~\eqref{dp:eq:weighted-norm},
$\norm h_{T,\tau}=\inf\{\norm z_\tau:q_T(z)=h\}$. Call the elements $\mathbf a$ of $\N^{\mathcal J_m}_p$
\emph{positions}, let $\abs{\mathbf a}_\partial$ be the sum of the entries of $\mathbf a$ at the elements of
$\N^m_{p-1}$, and call $\mathbf a$ \emph{nonlinear} if $\abs{\mathbf a}_\partial\ge2$. For a word $w\in\mathcal J_m^p$
let $\abs w_\partial$ be the number of its letters in $\N^m_{p-1}$, and for $S\in\Gamma^p(\Bbbk^{\mathcal J_m})$ let
$\operatorname{supp}S$ be the set of positions $\mathbf a$ with $S_{\mathbf a}\ne0$. Put
$\mathfrak b_m=(p-1)!\,m^p/(2p-1)!$.

\emph{Step 1 (norms).} If $h=q_T(z)$, the $i$-th entry of $h$ is $z_i+\sum_\gamma z_\gamma T_{\gamma+e_i}$, of absolute
value at most $\norm z_\tau$, as $\abs{T_\alpha}\le1$ and $\tau\ge1$. So no entry of $h\in \mathcal E_{T,\tau}$ has absolute
value above $\norm h_{T,\tau}$: the coordinate forms of $\mathcal E_{T,\tau}$ are bounded and separate points, so $\pi$ is a norm
on $\mathcal E_{T,\tau}^{\otimes p}$, and so is $\norm\cdot_\Gamma$ on $\Gamma^p(\mathcal E_{T,\tau})$. A representation
$t=\sum_lc_lh_l^{\otimes p}$ in $\Gamma^p(P)$, for a normed space $P$, gives $\pi(t)\le\sum_l\abs{c_l}\norm{h_l}^p$;
so $\norm t_\Gamma\ge\pi(t)$, by~\eqref{dp:eq:gamma-norm}.

\emph{Step 2 (the projective norm).} In $\mathcal E_{T,\tau}$, $\norm{e_i}\le1$ and $\norm{\partial^\gamma T}\le\tau$,
by~\eqref{dp:eq:weighted-norm}. For a word $(i_1,\dots,i_{p-1})\in[m]^{p-1}$ let $\gamma(i_1,\dots,i_{p-1})$ be its
content. The value of $T$ at $(i_1,\dots,i_{p-1},i)$ is the $i$-th entry of $\partial^{\gamma(i_1,\dots,i_{p-1})}T$, so
\[
 T=\sum_{(i_1,\dots,i_{p-1})\in[m]^{p-1}}e_{i_1}\otimes\dots\otimes e_{i_{p-1}}\otimes\partial^{\gamma(i_1,\dots,i_{p-1})}T,
\]
a sum of $m^{p-1}$ elementary tensors whose factors have norms with product at most $\tau$. Hence
$\pi(T)\le\tau m^{p-1}$ in $\mathcal E_{T,\tau}^{\otimes p}$.

\emph{Step 3 (a representation gives a symmetric tensor over $\F_p$).} The elements of $K$ of absolute value below $1$
form an $\F_p$\nobreakdash-linear subspace that does not contain $1$, so there is an $\F_p$\nobreakdash-linear map
$\rho\colon K\to\F_p$ with $\rho(1)=1$ that vanishes on it; then $\rho$ fixes $\F_p$, and $\abs x\ge1$ for every
$x\in K$ with $\rho(x)\ne0$. Let $\delta>0$ and $T=\sum_lc_lh_l^{\otimes p}$ with $c_l\in K$, $h_l\in \mathcal E_{T,\tau}$ and
$\sum_l\abs{c_l}\norm{h_l}_{T,\tau}^p\le\norm T_\Gamma+\delta$. Choose $\hat h_l\in K^{\mathcal J_m}$ with
$q_T(\hat h_l)=h_l$ and $\norm{\hat h_l}_\tau\le(1+\delta)\norm{h_l}_{T,\tau}$. The tensor
$\widetilde S=\sum_lc_l\hat h_l^{\otimes p}$ is symmetric, and $q_T^{\otimes p}(\widetilde S)=T$. For
$z\in K^{\mathcal J_m}$ the sum of the numbers $\tau^{\abs w_\partial}\abs{z^{\otimes p}(w)}$, $w\in\mathcal J_m^p$, is
$\norm z_\tau^p$; so, by the triangle inequality,
\begin{equation}\label{dp:eq:weighted-lift}
 \sum_{w\in\mathcal J_m^p}\tau^{\abs w_\partial}\abs{\widetilde S(w)}\le\sum_l\abs{c_l}\norm{\hat h_l}_\tau^p
 \le(1+\delta)^p(\norm T_\Gamma+\delta).
\end{equation}
Let $S\in\Gamma^p(\F_p^{\mathcal J_m})$ have the values $S(w)=\rho(\widetilde S(w))$. As the matrix of $q_T$ has
entries in $\F_p$, each value of $q_T^{\otimes p}(S)$ is the image under $\rho$ of the corresponding value of
$q_T^{\otimes p}(\widetilde S)=T$, which lies in $\F_p$; so $q_T^{\otimes p}(S)=T$. For $\mathbf a\in\operatorname{supp}S$
and a word $w$ of content $\mathbf a$, $\rho(\widetilde S(w))=S_{\mathbf a}\ne0$, so $\abs{\widetilde S(w)}\ge1$; and
$\abs w_\partial=\abs{\mathbf a}_\partial$. Taking one such word for each $\mathbf a\in\operatorname{supp}S$,
\eqref{dp:eq:weighted-lift} gives
$\sum_{\mathbf a\in\operatorname{supp}S}\tau^{\abs{\mathbf a}_\partial}\le(1+\delta)^p(\norm T_\Gamma+\delta)$. As
$\Gamma^p(\F_p^{\mathcal J_m})$ is finite, one $S$ occurs for arbitrarily small $\delta$; so
\begin{equation}\label{dp:eq:support-bound}
 \norm T_\Gamma\ge\sum_{\mathbf a\in\operatorname{supp}S}\tau^{\abs{\mathbf a}_\partial}\quad\text{for some }
 S\in\Gamma^p(\F_p^{\mathcal J_m})\text{ with }q_T^{\otimes p}(S)=T.
\end{equation}

\emph{Step 4 (supports with few nonlinear positions).} Let $S\in\Gamma^p(\F_p^{\mathcal J_m})$, and let
$\mathcal K_S\subseteq\N^m_p$ be the set of the $\gamma+e_i$, $i\in[m]$, for the $\gamma\in\N^m_{p-1}$ at which some
nonlinear $\mathbf a\in\operatorname{supp}S$ has a nonzero entry. A position has nonzero entries at no more than $p$
elements of $\N^m_{p-1}$, so $\abs{\mathcal K_S}$ is at most $pm$ times the number of nonlinear positions in
$\operatorname{supp}S$. We show that at most $p^{\abs{\mathcal K_S}+\binom{m+p-1}p-\mathfrak b_m}$ tensors
$T'\in\Gamma^p(\F_p^m)$ satisfy $q_{T'}^{\otimes p}(S)=T'$. Write $S=S^{(0)}+S^{(1)}+S^{(2)}$, where $S^{(0)}$,
$S^{(1)}$ and $S^{(2)}$ are supported in the positions $\mathbf a$ with $\abs{\mathbf a}_\partial=0$,
$\abs{\mathbf a}_\partial=1$ and $\abs{\mathbf a}_\partial\ge2$. The positions with $\abs{\mathbf a}_\partial=0$ are the
$\alpha\in\N^m_p$, and $q_{T'}$ fixes the $e_i$, so $q_{T'}^{\otimes p}(S^{(0)})=\sum_{\alpha\in\N^m_p}S_\alpha O_\alpha$
does not depend on $T'$. For $\gamma,\gamma'\in\N^m_{p-1}$ let $\mathbf a(\gamma,\gamma')$ be the position with the
entry $1$ at $\gamma$, the entries of $\gamma'$ on $[m]$ and $0$ elsewhere. These are the positions with
$\abs{\mathbf a}_\partial=1$, and $O_{\mathbf a(\gamma,\gamma')}=e_\gamma\cdot O_{\gamma'}$, as each word of content
$\mathbf a(\gamma,\gamma')$ arises from exactly one of the $p$ permutations that define this product, applied to
exactly one tensor $e_\gamma\otimes e_{i_1}\otimes\dots\otimes e_{i_{p-1}}$ with $(i_1,\dots,i_{p-1})$ of content
$\gamma'$. As the product commutes with $q_{T'}$,
\[
 q_{T'}^{\otimes p}(S^{(1)})=\sum_{\gamma,\gamma'\in\N^m_{p-1}}S_{\mathbf a(\gamma,\gamma')}\,
 (\partial^\gamma T')\cdot O_{\gamma'}=\mathsf L^\sharp(T'),
\]
with $\mathsf L_{\gamma\gamma'}=S_{\mathbf a(\gamma,\gamma')}$ in Lemma~\ref{dp:lem:linear}. Finally,
$q_{T'}^{\otimes p}(S^{(2)})$ depends only on the vectors $q_{T'}(e_\gamma)=\partial^\gamma T'$ for the $\gamma$ at which
some $\mathbf a\in\operatorname{supp}S^{(2)}$ has a nonzero entry, that is, on the coefficients $T'_\alpha$,
$\alpha\in\mathcal K_S$. So if $q_{T'}^{\otimes p}(S)=T'$, then $(\id-\mathsf L^\sharp)(T')$ equals
$q_{T'}^{\otimes p}(S^{(0)})+q_{T'}^{\otimes p}(S^{(2)})$, which depends only on $(T'_\alpha)_{\alpha\in\mathcal K_S}$.
For each of the $p^{\abs{\mathcal K_S}}$ choices of these coefficients the solutions $T'$ lie in one coset of the kernel
of $\id-\mathsf L^\sharp$, which has at most $p^{\binom{m+p-1}p-\mathfrak b_m}$ elements by
Lemma~\ref{dp:lem:linear}(2), as $\dim\Gamma^p(\F_p^m)=\binom{m+p-1}p$.

\emph{Step 5 (item (1)).} Put $\mathfrak n_1=\mathfrak b_m/(3p^2(1+\log_pm))$, $\mathfrak n_2=\mathfrak b_m/(3pm)$ and
$\tau=\max\{1,(m/(p(1+\log_pm)))^{1/2}\}$. There are $\binom{\abs{\mathcal J_m}+p-1}p\le\abs{\mathcal J_m}^p$
positions, and $\abs{\mathcal J_m}=m+\binom{m+p-2}{p-1}\le m+m^{p-1}\le2m^{p-1}$; so $p$ times the number of positions
is at most $p\,2^pm^{p(p-1)}\le p^{p^2}m^{p^2}=p^{p^2(1+\log_pm)}$, as $p\,2^p\le p^{p^2}$. Every
$S\in\Gamma^p(\F_p^{\mathcal J_m})$ that is nonzero at fewer than $\mathfrak n_1$ positions is the sum of
$\lceil\mathfrak n_1\rceil-1$ tensors $O_{\mathbf a}$, $\mathbf a$ a position, each multiplied by an element of $\F_p$;
so there are at most $p^{\mathfrak n_1p^2(1+\log_pm)}=p^{\mathfrak b_m/3}$ of them. If moreover fewer than
$\mathfrak n_2$ of its positions are nonlinear, then $\abs{\mathcal K_S}\le pm\mathfrak n_2=\mathfrak b_m/3$, and by
Step~4 at most $p^{\binom{m+p-1}p-2\mathfrak b_m/3}$ tensors $T'$ satisfy $q_{T'}^{\otimes p}(S)=T'$. So at most
$p^{\binom{m+p-1}p-\mathfrak b_m/3}$ of the $p^{\binom{m+p-1}p}$ elements $T'$ of $\Gamma^p(\F_p^m)$ satisfy
$q_{T'}^{\otimes p}(S)=T'$ for such an $S$. Choose $T\in\Gamma^p(\F_p^m)$ outside them: every $S$ with
$q_T^{\otimes p}(S)=T$ is nonzero at no fewer than $\mathfrak n_1$ positions, or at no fewer than $\mathfrak n_2$
nonlinear positions. The weights $\tau^{\abs{\mathbf a}_\partial}$ in~\eqref{dp:eq:support-bound} are at least $1$, and
at least $\tau^2$ for nonlinear $\mathbf a$, and $\tau^2\mathfrak n_2\ge\mathfrak n_1$, as
$\tau^2\ge m/(p(1+\log_pm))$; so $\norm T_\Gamma\ge\mathfrak n_1>0$ in $\Gamma^p(\mathcal E_{T,\tau})$, and $T\ne0$. If $\tau>1$,
Step~2 gives
\[
 \frac{\norm T_\Gamma}{\pi(T)}\ge\frac{\mathfrak n_1}{\tau m^{p-1}}
 =\frac{(p-1)!}{3p^{3/2}(2p-1)!}\Bigl(\frac m{1+\log_pm}\Bigr)^{1/2}.
\]
If $\tau=1$, then $m\le p(1+\log_pm)$, so the right side of~\eqref{dp:eq:degree-p-ratio} is at most
$(p-1)!\,\pi(T)/(3p(2p-1)!)<\pi(T)\le\norm T_\Gamma$, by Step~1.

\emph{Step 6 (item (2): the law).} If $K$ is complete, the finite-dimensional spaces $\mathcal E_m$ are complete, and so is
their $\ell^1$-sum $\mathcal E$. By Step~1 the coordinate forms of the spaces $\mathcal E_m$, composed with the projections of $\mathcal E$ onto
them, are bounded; they separate points, so $\pi$ is a norm on $\mathcal E^{\otimes p}$ and $\norm\cdot_\Gamma\ge\pi$ is a norm
on $\Gamma^p(\mathcal E)$. For finitely many vectors $(\lambda_l,h_l)$ of $K\oplus \mathcal E$ and scalars $s_l$, the value of the map at
$\sum_ls_l(\lambda_l,h_l)$ is $(\sum_ls_l\lambda_l)^{n-p}(\sum_ls_lh_l)^{\otimes p}$, a homogeneous polynomial of degree
$n$ in the $s_l$ whose coefficients are combinations of symmetric tensors of the span of the $h_l$, which lie in
$\Gamma^p(\mathcal E)$ as $K$ is infinite. Moreover
$\norm{\lambda^{n-p}h^{\otimes p}}_\Gamma\le\abs\lambda^{n-p}\norm h^p\le\norm{(\lambda,h)}^n$. So the map is a bounded
polynomial law of degree $n$ and norm at most $1$.

\emph{Step 7 (item (2): the ratio on $\Gamma^p(\mathcal E)$).} The inclusion of $\mathcal E_m$ in $\mathcal E$ is isometric, and the projection of
$\mathcal E$ onto $\mathcal E_m$ is a contractive left inverse. So, for $t\in\Gamma^p(\mathcal E_m)$, $\pi(t)$ computed in $\mathcal E^{\otimes p}$ is at
most $\pi(t)$ computed in $\mathcal E_m^{\otimes p}$, and $\norm t_\Gamma$ computed in $\Gamma^p(\mathcal E)$ is at least
$\norm t_\Gamma$ computed in $\Gamma^p(\mathcal E_m)$, as the $p$-th tensor power of the projection fixes $t$ and does not
increase the cost $\sum_l\abs{c_l}\norm{h_l}^p$ of a representation. By~(1), applied to $t=T^{(m)}\ne0$,
$\norm\cdot_\Gamma/\pi$ is unbounded on the nonzero elements of $\Gamma^p(\mathcal E)$.

\emph{Step 8 (item (2): comparing coefficients).} Suppose that $j(\lambda^{n-p}h^{\otimes p})=\beta((\lambda,h),\dots,(\lambda,h))$ for all
$(\lambda,h)\in K\oplus \mathcal E$, for a linear isometry $j$ of $\Gamma^p(\mathcal E)$, or of its completion, into a normed space $Z$
and a bounded $n$\nobreakdash-linear map $\beta\colon(K\oplus \mathcal E)^n\to Z$. Put $\mathbf 1=(1,0)$ and $h^+=(0,h)$ for
$h\in \mathcal E$. For $h_1,\dots,h_p\in \mathcal E$ let $\beta'(h_1,\dots,h_p)$ be the sum, over the $\binom np$ sets of $p$ of the $n$
slots, of the value of $\beta$ at the $n$-tuple with $h_1^+,\dots,h_p^+$ in these slots, in this order, and
$\mathbf 1$ in the others. Then $\beta'$ is $p$\nobreakdash-linear, and $\norm{\beta'}\le\binom np\norm\beta$, as
$\norm{\mathbf 1}=1$ and $\norm{h^+}=\norm h$. For $h\in \mathcal E$ and $s=(s_1,s_2)\in K^2$, the value $\beta((s_2,s_1h),\dots,(s_2,s_1h))$, with
$(s_2,s_1h)=s_2\mathbf 1+s_1h^+$, is a homogeneous polynomial of degree $n$ in $s$ whose
coefficient of $s_1^ps_2^{n-p}$ is $\beta'(h,\dots,h)$, while $j(s_2^{n-p}(s_1h)^{\otimes p})=s_1^ps_2^{n-p}j(h^{\otimes p})$.
As $K$ is infinite, comparing coefficients, as in Step~6 of the proof of Proposition~\ref{dp:prop:nonambient}, gives
$j(h^{\otimes p})=\beta'(h,\dots,h)$. Let $\beta'_\otimes\colon\otimes^p_\pi \mathcal E\to Z$ be the linear map with
$\beta'_\otimes(h_1\otimes\dots\otimes h_p)=\beta'(h_1,\dots,h_p)$, of norm $\norm{\beta'}$. As $\Gamma^p(\mathcal E)$ is spanned
by powers, $j(t)=\beta'_\otimes(t)$ for $t\in\Gamma^p(\mathcal E)$, so
$\norm t_\Gamma=\norm{j(t)}\le\binom np\norm\beta\,\pi(t)$ for $t\in\Gamma^p(\mathcal E)$, against Step~7.
\end{proof}

For the proof of Theorem~\ref{dp:thm:frobenius} let $p=\operatorname{char}K>0$, with $\Gamma^p(P)$,
$\norm\cdot_\Gamma$, $\norm\cdot_{\mathrm F}$, $\pi$, the hollow tensors and the sets $\N^r_p$ as in
Section~\ref{sec:dp}. For a subspace $P_0$ of $P$ we read $P_0^{\otimes p}$ as a subspace of $P^{\otimes p}$, as over
every field. For vectors $h_1,\dots,h_r\in P$ and $\alpha\in\N^r_p$ let $O_\alpha(h_1,\dots,h_r)$ be the sum of the
tensors $h_{w_1}\otimes\dots\otimes h_{w_p}$ over the words $w\in[r]^p$ whose content, defined before
Lemma~\ref{dp:lem:linear}, is $\alpha$; the same sum with the unit vectors of $K^r$ in place of $h_1,\dots,h_r$ is the
tensor $O_\alpha$ there, for $\Bbbk=K$ and $\mathcal J=[r]$. Expanding the power gives
\begin{equation}\label{dp:eq:orbit}
 (s_1h_1+\dots+s_rh_r)^{\otimes p}=\sum_{\alpha\in\N^r_p}s_1^{\alpha_1}\cdots s_r^{\alpha_r}\,O_\alpha(h_1,\dots,h_r)
 \qquad(s_1,\dots,s_r\in K),
\end{equation}
and $O_\alpha(h_1,\dots,h_r)$ lies in $\Gamma^p(P)$, since the coefficients of this polynomial are linear combinations
of its values, by interpolation in each variable, as in Step~5 of the proof of Theorem~\ref{dp:thm:calculus}.

\begin{lemma}[linearization]\label{dp:lem:linearization}
Let $a\colon P\to W$ be a bounded polynomial law of degree $p$.
\begin{enumerate}
\item There is a unique linear map $L_a\colon\Gamma^p(P)\to W$ with $L_a(h^{\otimes p})=a(h)$ for $h\in P$, and
$\norm{L_a(t)}\le\norm a\norm t_\Gamma$ for $t\in\Gamma^p(P)$. If $a$ is a Frobenius law, $L_a$ vanishes on the
hollow tensors, and $\norm{L_a(t)}\le\norm a\norm t_{\mathrm F}$.
\item If $M\ge0$ and $\norm{L_a(t)}\le M\pi(t)$ for $t\in\Gamma^p(P)$, then $a$ is ambient of degree $p$, with a
$p$\nobreakdash-linear map $\beta$ of norm at most $M$ in Definition~\ref{dp:def:laws}(2).
\item If $P$ satisfies $(\mathrm F_P)$ with constant $C_{\mathrm F}$ and $(\mathrm N_P)$ with constant $C_{\mathrm N}$,
then $\norm t_\Gamma\le\bigl(C_{\mathrm F}+(1+C_{\mathrm F})C_{\mathrm N}\bigr)\pi(t)$ for $t\in\Gamma^p(P)$.
\end{enumerate}
\end{lemma}

\begin{proof}
\emph{Step 1 (item (1): the map $L_a$).} Let $\sum_i\lambda_ih_i^{\otimes p}=0$, a finite sum with $\lambda_i\in K$
and $h_i\in P$. Let $b_1,\dots,b_m$ be a basis of the span of the $h_i$, write $h_i=\sum_ls_{i,l}b_l$, and put
$s_i^\alpha=s_{i,1}^{\alpha_1}\cdots s_{i,m}^{\alpha_m}$ for $\alpha\in\N^m_p$. By~\eqref{dp:eq:orbit} and
Definition~\ref{dp:def:laws}(1), $h_i^{\otimes p}=\sum_\alpha s_i^\alpha O_\alpha(b_1,\dots,b_m)$ and
$a(h_i)=\sum_\alpha s_i^\alpha a^{[\alpha]}(b_1,\dots,b_m)$, the sums over $\N^m_p$. The tensors
$O_\alpha(b_1,\dots,b_m)$ are sums over disjoint sets of words of the linearly independent tensors
$b_{w_1}\otimes\dots\otimes b_{w_p}$, hence linearly independent, so $\sum_i\lambda_is_i^\alpha=0$ for every $\alpha$,
and $\sum_i\lambda_ia(h_i)=0$. As the powers span $\Gamma^p(P)$,
$L_a(\sum_i\lambda_ih_i^{\otimes p})=\sum_i\lambda_ia(h_i)$ defines the unique such linear map, and
$\norm{L_a(t)}\le\sum_i\abs{\lambda_i}\norm{a(h_i)}\le\norm a\sum_i\abs{\lambda_i}\norm{h_i}^p$ for every
representation $t=\sum_i\lambda_ih_i^{\otimes p}$; the infimum over these is the bound, by~\eqref{dp:eq:gamma-norm}.

\emph{Step 2 (item (1): Frobenius laws).} Let $a$ be a Frobenius law and $h_1,\dots,h_p\in P$.
By~\eqref{dp:eq:orbit}, for $s_1,\dots,s_p\in K$,
\[
 \sum_{\alpha\in\N^p_p}s_1^{\alpha_1}\cdots s_p^{\alpha_p}L_a\bigl(O_\alpha(h_1,\dots,h_p)\bigr)
 =a(s_1h_1+\dots+s_ph_p)=s_1^pa(h_1)+\dots+s_p^pa(h_p).
\]
As $K$ is infinite and $p\ge2$, comparing the coefficients of $s_1\cdots s_p$ gives
$L_a(O_{(1,\dots,1)}(h_1,\dots,h_p))=0$. The words of content $(1,\dots,1)$ are the orderings of $[p]$, so
$O_{(1,\dots,1)}(h_1,\dots,h_p)$ is one of the tensors that span the hollow tensors, and $L_a$ vanishes on the hollow
tensors. If $t-t'$ is hollow, then $\norm{L_a(t)}=\norm{L_a(t')}\le\norm a\norm{t'}_\Gamma$; the infimum over $t'$ is
the bound.

\emph{Step 3 (item (2)).} For $v\in W$ and $y\in P^{\otimes p}$ let $\norm{(v,y)}_Z$ be the infimum of
$\norm{v+L_a(t)}+M\pi(y-t)$ over $t\in\Gamma^p(P)$, the quotient seminorm of $\norm v+M\pi(y)$ on
$W\oplus P^{\otimes p}$ modulo the subspace of the pairs $(L_a(t),-t)$. Let $Z$ be the quotient of
$W\oplus P^{\otimes p}$ by the kernel of $\norm\cdot_Z$, with the norm $\norm\cdot_Z$, and $[v,y]$ the class of
$(v,y)$. The linear map $j(v)=[v,0]$ is an isometry: $t=0$ gives $\norm{[v,0]}_Z\le\norm v$, and
$\norm{v+L_a(t)}+M\pi(-t)\ge\norm v-\norm{L_a(t)}+M\pi(t)\ge\norm v$ for every $t$, by the hypothesis on $L_a$. The
map $\beta(h_1,\dots,h_p)=[0,h_1\otimes\dots\otimes h_p]$ is $p$\nobreakdash-linear with
$\norm{\beta(h_1,\dots,h_p)}_Z\le M\pi(h_1\otimes\dots\otimes h_p)\le M\norm{h_1}\cdots\norm{h_p}$, and
$\beta(h,\dots,h)-ja(h)=[-a(h),h^{\otimes p}]=0$, as $t=h^{\otimes p}$ shows.

\emph{Step 4 (item (3)).} Let $t\in\Gamma^p(P)$ and $\delta>0$. By $(\mathrm F_P)$ there is $t'\in\Gamma^p(P)$ with
$t-t'$ hollow and $\norm{t'}_\Gamma\le C_{\mathrm F}\pi(t)+\delta$. Then
$\pi(t-t')\le\pi(t)+\norm{t'}_\Gamma\le(1+C_{\mathrm F})\pi(t)+\delta$, as $\pi\le\norm\cdot_\Gamma$, and by
$(\mathrm N_P)$, $\norm t_\Gamma\le\norm{t'}_\Gamma+C_{\mathrm N}\pi(t-t')\le C_{\mathrm F}\pi(t)+\delta
+C_{\mathrm N}\bigl((1+C_{\mathrm F})\pi(t)+\delta\bigr)$. Let $\delta\to0$.
\end{proof}

\begin{lemma}[bases of representations]\label{dp:lem:bases}
Let $P_0$ be a finite-dimensional subspace of $P$ and $t\in\Gamma^p(P)\cap P_0^{\otimes p}$, and for linear forms
$\phi_1,\dots,\phi_p$ on $P_0$ let $t(\phi_1,\dots,\phi_p)$ be the value of $\phi_1\otimes\dots\otimes\phi_p$ at $t$,
as before Lemma~\ref{dp:lem:projections}.
\begin{enumerate}
\item Let $b_1,\dots,b_m$ be a basis of $P_0$, with dual basis $b_1^*,\dots,b_m^*$. Then
$t-\sum_{i=1}^mt(b_i^*,\dots,b_i^*)\,b_i^{\otimes p}$ is hollow.
\item Let $t=\sum_{k=1}^rx_{k,1}\otimes\dots\otimes x_{k,p}$ with nonzero vectors $x_{k,c}$ that span $P_0$. There
is a set $\mathcal B\subseteq[r]\times[p]$ such that the vectors $x_{k,c}$, $(k,c)\in\mathcal B$, form a basis of
$P_0$ whose dual basis $(x_{k,c}^*)_{(k,c)\in\mathcal B}$ satisfies
$\abs{t(x_{k,c}^*,\dots,x_{k,c}^*)}\,\norm{x_{k,c}}^p\le\norm{x_{k,1}}\cdots\norm{x_{k,p}}$ for every
$(k,c)\in\mathcal B$.
\end{enumerate}
\end{lemma}

\begin{proof}
\emph{Step 1 (item (1)).} Let $q$ be a linear projection of $P$ onto $P_0$. As $q^{\otimes p}$ fixes
$P_0^{\otimes p}$ and maps powers to powers, $t=q^{\otimes p}(t)$ is a combination of powers $h^{\otimes p}$,
$h\in P_0$. For a word $w\in[m]^p$ put $b_w=b_{w_1}\otimes\dots\otimes b_{w_p}$; these tensors form a basis of
$P_0^{\otimes p}$, and $t=\sum_wt(b_{w_1}^*,\dots,b_{w_p}^*)\,b_w$. For a power $h^{\otimes p}$, $h\in P_0$, the
coefficient $b_{w_1}^*(h)\cdots b_{w_p}^*(h)$ of $b_w$ depends only on the content $\alpha\in\N^m_p$ of $w$; hence so
does that of $t$, and $t$ is a combination of the tensors $O_\alpha(b_1,\dots,b_m)$. If every $\alpha_i<p$, then
$\alpha_1!\cdots\alpha_m!$ is invertible in $K$, and for a word $w$ of content $\alpha$ the tensor
$O_{(1,\dots,1)}(b_{w_1},\dots,b_{w_p})$, hollow as in Step~2 of the proof of Lemma~\ref{dp:lem:linearization}, is
$\alpha_1!\cdots\alpha_m!\,O_\alpha(b_1,\dots,b_m)$, as the reorderings of the letters of $w$ give each word of content
$\alpha$ exactly $\alpha_1!\cdots\alpha_m!$ times; so $O_\alpha(b_1,\dots,b_m)$ is hollow. Otherwise $\alpha_i=p$ for
one $i$, the tensor $O_\alpha(b_1,\dots,b_m)$ is $b_i^{\otimes p}$, and its coefficient is $t(b_i^*,\dots,b_i^*)$.

\emph{Step 2 (item (2): the choice of $\mathcal B$).} Call $\mathcal B\subseteq[r]\times[p]$ a \emph{basis set} if the
vectors $x_{k,c}$, $(k,c)\in\mathcal B$, form a basis of $P_0$; basis sets exist, as these vectors span $P_0$. Fix a
basis of $P_0$, and for a basis set $\mathcal B$ let $\det_{\mathcal B}$ be the determinant of the coordinates of the
vectors $x_{k,c}$, $(k,c)\in\mathcal B$, in a fixed order of $[r]\times[p]$. Put
$\upsilon_{k,c}=(\norm{x_{k,1}}\cdots\norm{x_{k,p}})^{1/p}/\norm{x_{k,c}}>0$, so that
$\upsilon_{k,1}\cdots\upsilon_{k,p}=1$ for each $k$, and choose, among the finitely many basis sets, a basis set
$\mathcal B$ with the largest value of $\abs{\det_{\mathcal B}}\prod_{(k,c)\in\mathcal B}\upsilon_{k,c}$. Then
\begin{equation}\label{dp:eq:exchange}
 \abs{x_{k,c}^*(x_{k',c'})}\,\upsilon_{k',c'}\le\upsilon_{k,c}\qquad
 \bigl((k,c)\in\mathcal B,\ (k',c')\in[r]\times[p]\bigr).
\end{equation}
This is clear if $(k',c')\in\mathcal B$ or $x_{k,c}^*(x_{k',c'})=0$. Otherwise
$x_{k',c'}=\sum_{(k'',c'')\in\mathcal B}x_{k'',c''}^*(x_{k',c'})\,x_{k'',c''}$, whose coefficient at $(k,c)$ is
$x_{k,c}^*(x_{k',c'})\ne0$, so replacing $(k,c)$ by $(k',c')$ gives a basis set $\mathcal B'$, and
$\det_{\mathcal B'}=\pm x_{k,c}^*(x_{k',c'})\det_{\mathcal B}$ by multilinearity and alternation of the determinant,
the sign coming from the order of $[r]\times[p]$. The choice of $\mathcal B$ gives
$\abs{x_{k,c}^*(x_{k',c'})}\,\upsilon_{k',c'}/\upsilon_{k,c}\le1$.

\emph{Step 3 (item (2): the bound).} Let $(k,c)\in\mathcal B$. Then
$t(x_{k,c}^*,\dots,x_{k,c}^*)=\sum_{k'=1}^r\prod_{c'=1}^px_{k,c}^*(x_{k',c'})$, so by the ultrametric inequality
and~\eqref{dp:eq:exchange}
\[
 \abs{t(x_{k,c}^*,\dots,x_{k,c}^*)}\le\max_{k'}\prod_{c'=1}^p\frac{\upsilon_{k,c}}{\upsilon_{k',c'}}=\upsilon_{k,c}^p
 =\frac{\norm{x_{k,1}}\cdots\norm{x_{k,p}}}{\norm{x_{k,c}}^p}.
\]
\end{proof}

\begin{proof}[Proof of Theorem~\ref{dp:thm:frobenius}]
\emph{Step 1 (item (1)).} Let $t\in\Gamma^p(P)$ and $\delta>0$; we may assume $t\ne0$. Choose a representation
$t=\sum_{k=1}^rx_{k,1}\otimes\dots\otimes x_{k,p}$ with $\sum_k\norm{x_{k,1}}\cdots\norm{x_{k,p}}\le\pi(t)+\delta$;
dropping the terms with a zero factor, all $x_{k,c}$ are nonzero. Let $P_0$ be their span, so $\dim P_0\le\dim P$ and
$t\in P_0^{\otimes p}$, and take $\mathcal B$ from Lemma~\ref{dp:lem:bases}(2). By Lemma~\ref{dp:lem:bases}(1),
$t-t'$ is hollow for $t'=\sum_{(k,c)\in\mathcal B}t(x_{k,c}^*,\dots,x_{k,c}^*)\,x_{k,c}^{\otimes p}$, so
\begin{align*}
 \norm t_{\mathrm F}\le\norm{t'}_\Gamma
 &\le\sum_{(k,c)\in\mathcal B}\abs{t(x_{k,c}^*,\dots,x_{k,c}^*)}\,\norm{x_{k,c}}^p\\
 &\le\sum_{k=1}^r\abs{\mathcal B\cap(\{k\}\times[p])}\,\norm{x_{k,1}}\cdots\norm{x_{k,p}}
 \le\min(p,\dim P_0)\,(\pi(t)+\delta),
\end{align*}
since $\mathcal B$ has $\dim P_0$ elements. Let $\delta\to0$.

\emph{Step 2 (item (2)).} By Lemma~\ref{dp:lem:linearization}(1) and item~(1),
$\norm{L_a(t)}\le\norm a\norm t_{\mathrm F}\le p\norm a\,\pi(t)$ for $t\in\Gamma^p(P)$; apply
Lemma~\ref{dp:lem:linearization}(2) with $M=p\norm a$.

\emph{Step 3 (item (3): the bound).} By item~(1) and Lemma~\ref{dp:lem:linearization}(3) with $C_{\mathrm F}=p$ and
$C_{\mathrm N}=C$, $\norm t_\Gamma\le(p+(p+1)C)\pi(t)$ for $t\in\Gamma^p(P)$. So
$\norm{L_a(t)}\le(p+(p+1)C)\norm a\,\pi(t)$ by Lemma~\ref{dp:lem:linearization}(1); apply
Lemma~\ref{dp:lem:linearization}(2) with $M=(p+(p+1)C)\norm a$.

\emph{Step 4 (item (3): the converse).} Let $W_0$ be the quotient of $\Gamma^p(P)$ by the kernel of
$\norm\cdot_\Gamma$, normed by $\norm\cdot_\Gamma$, and $a_\Gamma(h)\in W_0$ the class of $h^{\otimes p}$.
By~\eqref{dp:eq:orbit} and~\eqref{dp:eq:gamma-norm}, $a_\Gamma$ is a bounded polynomial law of degree $p$ with
$\norm{a_\Gamma}\le1$; by hypothesis $ja_\Gamma(h)=\beta(h,\dots,h)$ for a linear isometry $j$ of $W_0$ into a normed
space $Z$ and a bounded $p$\nobreakdash-linear $\beta\colon P^p\to Z$. Let $\beta_\otimes\colon P^{\otimes p}\to Z$ be
the linear map with $\beta_\otimes(h_1\otimes\dots\otimes h_p)=\beta(h_1,\dots,h_p)$; then
$\norm{\beta_\otimes(y)}\le\norm\beta\pi(y)$ for $y\in P^{\otimes p}$. As $\beta_\otimes(h^{\otimes p})=ja_\Gamma(h)$
and the powers span $\Gamma^p(P)$, $\beta_\otimes$ agrees on $\Gamma^p(P)$ with $j$ composed with the quotient map, so
$\norm t_\Gamma\le\norm\beta\pi(t)$ for $t\in\Gamma^p(P)$, in particular for hollow $t$.
\end{proof}

For the proof of Proposition~\ref{dp:prop:ultrametric} let $n\ge1$, with $\Gamma^n(P)$, $\norm\cdot_\Gamma$, $\pi$,
$\mathfrak a_n(P)$, $[m]$, $c_0([m],K)$ and $\mathfrak c_K$ as in Section~\ref{sec:dp}. The statements of the
paragraph before Lemma~\ref{dp:lem:linearization}, the first assertion of Lemma~\ref{dp:lem:linearization}(1) and
Lemma~\ref{dp:lem:linearization}(2) hold with $n$ in place of $p$, with the same proofs, which do not use that $p$ is the
characteristic of $K$. In
particular~\eqref{dp:eq:orbit} holds with $n$ in place of $p$, with $O_\alpha(h_1,\dots,h_r)\in\Gamma^n(P)$ for
$\alpha\in\N^r_n$. As in Step~4 of the proof of Theorem~\ref{dp:thm:frobenius}, let $W_0$ be the quotient of
$\Gamma^n(P)$ by the kernel of $\norm\cdot_\Gamma$, normed by $\norm\cdot_\Gamma$, and let $a_\Gamma(h)\in W_0$ be the
class of $h^{\otimes n}$. As after~\eqref{dp:eq:gamma-norm}, $\pi\le\norm\cdot_\Gamma$ on $\Gamma^n(P)$.
When two spaces or two fields occur, we say in which space and over which field $\pi$, $\norm\cdot_\Gamma$ and $\mathfrak a_n$ are computed. The completion $\widehat P$ of $P$ is a normed space
over $\Kh$, the scalar multiplication extending by continuity, and it is nonarchimedean if $P$ is.

\begin{lemma}[properties of the ambience constants]\label{dp:lem:constants}
Let $P'$ be a normed space.
\begin{enumerate}
\item $a_\Gamma$ is a bounded polynomial law of degree $n$ with $\norm{a_\Gamma}\le1$. If $ja_\Gamma(h)=\beta(h,\dots,h)$
for all $h\in P$, for a linear isometry $j$ of $W_0$ into a normed space and a bounded $n$\nobreakdash-linear map
$\beta$ on $P^n$, then $\mathfrak a_n(P)\le\norm\beta$.
\item If $\Xi$ is a bounded linear bijection of $P$ onto $P'$ with bounded inverse, then
$\mathfrak a_n(P')\le(\norm\Xi\norm{\Xi^{-1}})^n\mathfrak a_n(P)$.
\item If $P'$ is a subspace of $P$, with the restricted norm, and some linear projection of $P$ onto $P'$ has norm at
most $1$, then $\mathfrak a_n(P')\le\mathfrak a_n(P)$.
\item $\mathfrak a_n(P)$ is at most the supremum of $\mathfrak a_n(P_0)$ over the finite-dimensional subspaces $P_0$ of
$P$, with the restricted norm.
\end{enumerate}
\end{lemma}

\begin{proof}
\emph{Step 1 (item (1)).} This is Step~4 of the proof of Theorem~\ref{dp:thm:frobenius}, with $n$ in place of $p$; it
does not use that $p$ is the characteristic of $K$.

\emph{Step 2 (item (2)).} We may assume that $\mathfrak a_n(P)$ is finite. Let $t'\in\Gamma^n(P')$ and $t=(\Xi^{-1})^{\otimes n}(t')$. The maps $\Xi^{\otimes n}$ and
$(\Xi^{-1})^{\otimes n}$ carry powers to powers and elementary tensors to elementary tensors, so $t\in\Gamma^n(P)$, and
they multiply the cost of a representation by at most $\norm\Xi^n$ and $\norm{\Xi^{-1}}^n$, respectively. Hence
$\norm{t'}_\Gamma\le\norm\Xi^n\norm t_\Gamma\le\norm\Xi^n\mathfrak a_n(P)\pi(t)
\le(\norm\Xi\norm{\Xi^{-1}})^n\mathfrak a_n(P)\pi(t')$.

\emph{Step 3 (item (3)).} We may assume that $\mathfrak a_n(P)$ is finite. Let $\Pi$ be such a projection and $t\in\Gamma^n(P')$. Then $\Pi^{\otimes n}$ fixes $t$ and
carries a representation of $t$ in $\Gamma^n(P)$ to one in $\Gamma^n(P')$ of no larger cost, and representations of $t$
in $P'^{\otimes n}$ are representations in $P^{\otimes n}$. So $\norm t_\Gamma$ computed in $\Gamma^n(P')$ is at most
$\norm t_\Gamma$ computed in $\Gamma^n(P)$, which is at most $\mathfrak a_n(P)\pi(t)$ with $\pi(t)$ computed in
$P^{\otimes n}$, hence at most $\mathfrak a_n(P)\pi(t)$ with $\pi(t)$ computed in $P'^{\otimes n}$.

\emph{Step 4 (item (4)).} Let $t\in\Gamma^n(P)$ and $\delta>0$, and let $t=\sum_kx_{k,1}\otimes\dots\otimes x_{k,n}$ with
$\sum_k\norm{x_{k,1}}\cdots\norm{x_{k,n}}\le\pi(t)+\delta$. Let $P_0$ be the span of the $x_{k,c}$. As in Step~1 of
the proof of Lemma~\ref{dp:lem:bases}, $t=q^{\otimes n}(t)\in\Gamma^n(P_0)$ for a linear projection $q$ of $P$ onto
$P_0$. Then $\norm t_\Gamma$ computed in $\Gamma^n(P)$ is at most $\norm t_\Gamma$ computed in $\Gamma^n(P_0)$, hence
at most $\mathfrak a_n(P_0)(\pi(t)+\delta)$. Let $\delta\to0$.
\end{proof}

\begin{lemma}[orthogonal bases]\label{dp:lem:orthogonal}
Let $K$ be complete, and let $P$ be nonarchimedean, of finite dimension $m\ge1$. Then
$\mathfrak a_n(P)\le\mathfrak c_K^{-n}\mathfrak a_n(c_0([m],K))$.
\end{lemma}

\begin{proof}
\emph{Step 1 (distances).} Let $P_1$ be a subspace of $P$, $x\in P\setminus P_1$ and
$\operatorname{dist}(x,P_1)=\inf_{y\in P_1}\norm{x-y}$. By Lemma~\ref{lem:findim}, $P_1$ is complete, hence closed in $P$,
so $\operatorname{dist}(x,P_1)>0$. Let $K$ be discretely valued. If nonzero vectors $u_1,\dots,u_k\in P$ have norms in
distinct cosets of the subgroup $\abs\vartheta^{\Z}$ of the multiplicative group $(0,\infty)$, then the nonzero numbers
among $\abs{\lambda_i}\norm{u_i}$, for $\lambda_1,\dots,\lambda_k\in K$, are distinct, so the nonarchimedean norm gives
$\norm{\lambda_1u_1+\dots+\lambda_ku_k}=\max_i\abs{\lambda_i}\norm{u_i}$; hence the $u_i$ are linearly independent, and $k\le m$. So the
nonzero values of the norm of $P$ lie in at most $m$ cosets of $\abs\vartheta^{\Z}$, and only finitely many of them lie
in the interval $[\operatorname{dist}(x,P_1),\norm x]$. Hence $\norm{x-y}=\operatorname{dist}(x,P_1)$ for some $y\in P_1$.

\emph{Step 2 (orthogonal bases).} Let $0<\mathfrak o\le1$, with $\mathfrak o<1$ unless $K$ is discretely valued. We show
by induction on $m$ that $P$ has a basis $b_1,\dots,b_m$ with
\begin{equation}\label{dp:eq:orthogonal}
 \norm{\lambda_1b_1+\dots+\lambda_mb_m}\ge\mathfrak o\max_i\abs{\lambda_i}\norm{b_i}\qquad(\lambda_1,\dots,\lambda_m\in K).
\end{equation}
For $m=1$ every nonzero vector will do. Let $m\ge2$, let $P_1$ be a hyperplane of $P$, and let $b_1,\dots,b_{m-1}$ be a
basis of $P_1$ with~\eqref{dp:eq:orthogonal} for $\mathfrak o^{1/2}$ in place of $\mathfrak o$. Let $x\in P\setminus P_1$.
By Step~1 there is $y\in P_1$ with $\norm{x-y}\le\mathfrak o^{-1/2}\operatorname{dist}(x,P_1)$: if $\mathfrak o<1$ since
$\operatorname{dist}(x,P_1)>0$, and if $\mathfrak o=1$ by its last sentence. Put $b_m=x-y$, so that
$\operatorname{dist}(b_m,P_1)=\operatorname{dist}(x,P_1)\ge\mathfrak o^{1/2}\norm{b_m}$. Let $v=v'+\lambda_mb_m$ with
$v'=\sum_{i<m}\lambda_ib_i$. As $-v'\in P_1$,
$\norm v\ge\abs{\lambda_m}\operatorname{dist}(b_m,P_1)\ge\mathfrak o^{1/2}\abs{\lambda_m}\norm{b_m}$, and so
$\norm{v'}\le\max(\norm v,\abs{\lambda_m}\norm{b_m})\le\mathfrak o^{-1/2}\norm v$. Hence
$\mathfrak o\max_{i<m}\abs{\lambda_i}\norm{b_i}\le\mathfrak o^{1/2}\norm{v'}\le\norm v$ and
$\mathfrak o\abs{\lambda_m}\norm{b_m}\le\norm v$, which is~\eqref{dp:eq:orthogonal}.

\emph{Step 3 (the bound).} Multiplying the $b_i$ by nonzero scalars keeps~\eqref{dp:eq:orthogonal}. If $K$ is discretely
valued, take $\mathfrak o=1$ in Step~2 and scale the $b_i$ by powers of $\vartheta$ so that $\abs\vartheta<\norm{b_i}\le1$,
and put $\mathfrak o'=\abs\vartheta$. Otherwise $\abs{K^\times}$ is a subgroup of $(0,\infty)$ that contains
$\abs\varpi$ and is not of the form $\abs\vartheta^{\Z}$, hence dense, since a subgroup of $(0,\infty)$ that contains a
number other than $1$ and is not dense is generated by its largest element below $1$; for $0<\mathfrak o<1$ take the
$b_i$ of Step~2 and scale them so that $\mathfrak o<\norm{b_i}\le1$, and put $\mathfrak o'=\mathfrak o$. The linear bijection
$\Xi(\lambda)=\sum_i\lambda_ib_i$, $\lambda=(\lambda_1,\dots,\lambda_m)$, of $c_0([m],K)$ onto $P$ satisfies
$\norm{\Xi(\lambda)}\le\max_i\abs{\lambda_i}\norm{b_i}\le\norm\lambda$, as $P$ is nonarchimedean, and
$\norm{\Xi(\lambda)}\ge\mathfrak o\max_i\abs{\lambda_i}\norm{b_i}\ge\mathfrak o\mathfrak o'\norm\lambda$,
by~\eqref{dp:eq:orthogonal}. Lemma~\ref{dp:lem:constants}(2) gives
$\mathfrak a_n(P)\le(\mathfrak o\mathfrak o')^{-n}\mathfrak a_n(c_0([m],K))$. This is the claim if $K$ is discretely
valued; otherwise let $\mathfrak o\to1$.
\end{proof}

\begin{lemma}[completion]\label{dp:lem:completion}
\leavevmode
\begin{enumerate}
\item $\mathfrak a_n(P)$ is at most $\mathfrak a_n(\widehat P)$ computed over $\Kh$.
\item For $m\ge1$, $\mathfrak a_n(c_0([m],\Kh))$ computed over $\Kh$ is at most $\mathfrak a_n(c_0([m],K))$.
\end{enumerate}
\end{lemma}

\begin{proof}
\emph{Step 1 (continuity).} Let $a\colon P\to W$ be a bounded polynomial law of degree $n$, and let $R>0$. For
$x,y\in P$ with $\norm x,\norm y\le R$, $a(x+y)-a(x)=\sum_{i=1}^na^{[(n-i,i)]}(x,y)$, and~\eqref{dp:eq:binomial} gives
\begin{equation}\label{dp:eq:lipschitz}
 \norm{a(x+y)-a(x)}\le n(16/\abs\varpi^2)^n\norm a\,R^{n-1}\norm y.
\end{equation}
If $\norm x\le R/2$ and $\norm{x+y}\le R/2$, then $\norm y\le R$; so, by~\eqref{dp:eq:lipschitz}, $a$ is Lipschitz on
the set of vectors of $P$ of norm at most $R/2$, for every $R>0$. Hence $a$ is uniformly continuous on every bounded
subset of $P$, and it extends uniquely to a continuous map $\hat a\colon\widehat P\to\widehat W$, which satisfies
$\norm{\hat a(x)}\le\norm a\norm x^n$ by continuity.

\emph{Step 2 (laws extend).} Let $a$ and $\hat a$ be as in Step~1, $r\ge1$ and $\alpha\in\N^r_n$. By interpolation in each variable on $n+1$ distinct elements of $K$, as in Step~5 of the proof of
Theorem~\ref{dp:thm:calculus}, there are finitely many scalars $\lambda_l$ and $\xi_{l,1},\dots,\xi_{l,r}$ in $K$,
depending only on $r$ and $\alpha$, with $a^{[\alpha]}(h_1,\dots,h_r)=\sum_l\lambda_la(\xi_{l,1}h_1+\dots+\xi_{l,r}h_r)$
for all $h_1,\dots,h_r\in P$. Define
$\hat a^{[\alpha]}(h_1,\dots,h_r)=\sum_l\lambda_l\hat a(\xi_{l,1}h_1+\dots+\xi_{l,r}h_r)$ for $h_1,\dots,h_r\in\widehat P$.
The identity
\[
 \hat a(s_1h_1+\dots+s_rh_r)=\sum_{\alpha\in\N^r_n}s_1^{\alpha_1}\cdots s_r^{\alpha_r}\,\hat a^{[\alpha]}(h_1,\dots,h_r)
\]
holds for $s\in K^r$ and $h_1,\dots,h_r\in P$, by Definition~\ref{dp:def:laws}(1) and the choice of the $\lambda_l$
and $\xi_{l,i}$; both sides are continuous in
$(s,h_1,\dots,h_r)\in\Kh^r\times\widehat P^r$, so it holds there. Hence $\hat a\colon\widehat P\to\widehat W$ is a bounded
polynomial law of degree $n$ over $\Kh$, with $\norm{\hat a}\le\norm a$.

\emph{Step 3 (item (1)).} Let $C=\mathfrak a_n(\widehat P)$, computed over $\Kh$, be finite. By Step~2, $a_\Gamma$
extends to a bounded polynomial law $\hat a_\Gamma\colon\widehat P\to\widehat W_0$ of degree $n$ over $\Kh$, of norm at
most $1$. As $\norm{\hat a_\Gamma}\le1$ and $\norm t_\Gamma\le C\pi(t)$ on $\Gamma^n(\widehat P)$, the first assertion of
Lemma~\ref{dp:lem:linearization}(1) and Lemma~\ref{dp:lem:linearization}(2), with $n$ in place of $p$, over $\Kh$ and
with $M=C$, give a $\Kh$\nobreakdash-linear isometry $\hat\jmath$ of
$\widehat W_0$ into a normed space $Z$ over $\Kh$ and a bounded $n$\nobreakdash-linear map $\hat\beta$ on $\widehat P^n$
over $\Kh$, of norm at most $C$, with $\hat\jmath\hat a_\Gamma(h)=\hat\beta(h,\dots,h)$ for $h\in\widehat P$. The
restrictions of $\hat\jmath$ to $W_0$ and of $\hat\beta$ to $P^n$ are a linear isometry of $W_0$ into $Z$, read as a
normed space over $K$, and a bounded $n$\nobreakdash-linear map of norm at most $C$, and
$\hat\jmath(a_\Gamma(h))=\hat\beta(h,\dots,h)$ for $h\in P$. Lemma~\ref{dp:lem:constants}(1) gives $\mathfrak a_n(P)\le C$.

\emph{Step 4 (projective norms in $c_0([m],K)$).} Read $K^m\subseteq\Kh^m$, so that the tensors
$e_{w_1}\otimes\dots\otimes e_{w_n}$, $w\in[m]^n$, are a basis of $(K^m)^{\otimes n}$ over $K$ and of
$(\Kh^m)^{\otimes n}$ over $\Kh$; call the coefficients of a tensor in this basis its coefficients. For
$z\in(K^m)^{\otimes n}$ we show that $\pi(z)$ is the same over $K$ and over $\Kh$, the spaces being $c_0([m],K)$ and
$c_0([m],\Kh)$. Over $\Kh$ it is at most its value over $K$. Over $K$, $\pi(z)$ is at most the sum of the absolute
values of the coefficients of $z$. Let $z=\sum_kx_{k,1}\otimes\dots\otimes x_{k,n}$ with $x_{k,c}\in\Kh^m$, let
$\mathfrak r=\sum_k\norm{x_{k,1}}\cdots\norm{x_{k,n}}$, and let $0<\delta\le1$. Choose $x'_{k,c}\in K^m$ with
$\norm{x'_{k,c}-x_{k,c}}\le\delta\norm{x_{k,c}}$; then $\norm{x'_{k,c}}\le\norm{x_{k,c}}$. Every coefficient of
\[
 x_{k,1}\otimes\dots\otimes x_{k,n}-x'_{k,1}\otimes\dots\otimes x'_{k,n}
 =\sum_{c=1}^nx'_{k,1}\otimes\dots\otimes x'_{k,c-1}\otimes(x_{k,c}-x'_{k,c})\otimes x_{k,c+1}\otimes\dots\otimes x_{k,n}
\]
has absolute value at most $\delta\norm{x_{k,1}}\cdots\norm{x_{k,n}}$. So
$z'=\sum_kx'_{k,1}\otimes\dots\otimes x'_{k,n}$ lies in $(K^m)^{\otimes n}$, $\pi(z')\le\mathfrak r$ over $K$, and
$z-z'$ has coefficients of absolute value at most $\delta\mathfrak r$. Over $K$ this gives
$\pi(z)\le\mathfrak r+m^n\delta\mathfrak r$. Let $\delta\to0$ and
take the infimum over the representations of $z$ over $\Kh$.

\emph{Step 5 (item (2)).} Let $C=\mathfrak a_n(c_0([m],K))$ be finite, let $t\in\Gamma^n(c_0([m],\Kh))$, computed over
$\Kh$, and let $\delta>0$. Write $O_\alpha=O_\alpha(e_1,\dots,e_m)$ for $\alpha\in\N^m_n$. By the paragraph before
Lemma~\ref{dp:lem:linear}, with $\Bbbk=\Kh$ and $\mathcal J=[m]$, $t=\sum_\alpha t_\alpha O_\alpha$ with
$t_\alpha\in\Kh$, as $t$ is symmetric. As noted before Lemma~\ref{dp:lem:constants}, the $O_\alpha$ lie in $\Gamma^n(c_0([m],K))$, computed over $K$; let $M_0$ be the
sum of their seminorms $\norm{O_\alpha}_\Gamma$ there. Choose $t'_\alpha\in K$ with $\abs{t_\alpha-t'_\alpha}\le\delta$
and put $t'=\sum_\alpha t'_\alpha O_\alpha$. Representations over $K$ are representations over $\Kh$, so, over $\Kh$,
$\norm{t-t'}_\Gamma\le\delta M_0$, hence $\pi(t-t')\le\delta M_0$ and $\pi(t')\le\pi(t)+\delta M_0$. With Step~4,
\[
 \norm t_\Gamma\le\norm{t'}_\Gamma+\delta M_0\le C\pi(t')+\delta M_0\le C(\pi(t)+\delta M_0)+\delta M_0,
\]
where $\norm t_\Gamma$ and $\pi(t)$ are computed over $\Kh$, $\norm{t'}_\Gamma$ over $K$, which is at least its value
over $\Kh$, and $\pi(t')$ over either field. Let $\delta\to0$.
\end{proof}

\begin{lemma}[degree two]\label{dp:lem:degree-two}
For every $m\ge1$, $\mathfrak a_2(c_0([m],K))\le3/2$.
\end{lemma}

\begin{proof}
Identify $(K^m)^{\otimes2}$ with $K^{m\times m}$ by $x\otimes x'\mapsto xx'^{\mathsf T}$, as in Step~5 of the proof of
Proposition~\ref{dp:prop:nonambient}. Then the elements of $\Gamma^2(c_0([m],K))$ are symmetric matrices,
$\norm{\mathsf S}_\Gamma$ is the infimum of $\sum_l\abs{c_l}\norm{u_l}^2$ over the representations
$\mathsf S=\sum_lc_lu_lu_l^{\mathsf T}$, and $\pi(\mathsf S)$ is the infimum of $\sum_l\norm{x_l}\norm{x'_l}$ over the
representations $\mathsf S=\sum_lx_lx_l'^{\mathsf T}$. A matrix is \emph{integral} if its entries have absolute value at
most $1$. For a matrix $\mathsf X$ and $k\ge1$ let $\mathfrak d_k(\mathsf X)$ be the largest absolute value of a
$k\times k$ minor of $\mathsf X$, or $0$ if there is none, and $\mathfrak d_0(\mathsf X)=1$. A determinant is a sum of
products with one entry from each column, so by the ultrametric inequality the minors of an integral matrix have
absolute value at most $1$, and by the Cauchy--Binet formula $\mathfrak d_k(\mathsf X)$ does not increase when
$\mathsf X$ is multiplied on either side by an integral matrix.

\emph{Step 1 (a pivot).} Let $\mathsf S\in K^{m\times m}$ be symmetric and nonzero, with entries $\mathsf S_{ij}$, and
let $\mathfrak f$ be the largest absolute value of an entry. If $\abs{\mathsf S_{ii}}=\mathfrak f$ for some $i$, put
$\mathcal V=\{i\}$; otherwise put $\mathcal V=\{i,j\}$ for some $i\ne j$ with $\abs{\mathsf S_{ij}}=\mathfrak f$. Let
$\mathsf S_{\mathcal V}$ be the submatrix of $\mathsf S$ with rows and columns in $\mathcal V$, and $\mathsf S_{[m]\mathcal
V}$ the one with columns in $\mathcal V$. Then $\abs{\det\mathsf S_{\mathcal V}}=\mathfrak f^{\abs{\mathcal V}}$; in the
second case $\det\mathsf S_{\mathcal V}=\mathsf S_{ii}\mathsf S_{jj}-\mathsf S_{ij}^2$ with
$\abs{\mathsf S_{ii}\mathsf S_{jj}}<\mathfrak f^2$. So $\mathsf S_{\mathcal V}^{-1}$, its adjugate divided by its
determinant, has entries of absolute value at most $\mathfrak f^{-1}$, and
$\mathsf G=\mathsf S_{[m]\mathcal V}\mathsf S_{\mathcal V}^{-1}$ is integral, with the identity matrix as its rows in
$\mathcal V$. Put $\mathsf S'=\mathsf S-\mathsf G\mathsf S_{\mathcal V}\mathsf G^{\mathsf T}$. It is symmetric with
entries of absolute value at most $\mathfrak f$, and its rows and columns in $\mathcal V$ vanish, as the rows in
$\mathcal V$ of $\mathsf G\mathsf S_{\mathcal V}\mathsf G^{\mathsf T}$ form $\mathsf S_{\mathcal V}\mathsf G^{\mathsf T}=
\mathsf S_{[m]\mathcal V}^{\mathsf T}$. If $\mathcal V=\{i\}$, then
$\mathsf G\mathsf S_{\mathcal V}\mathsf G^{\mathsf T}=\mathsf S_{ii}uu^{\mathsf T}$ for the column $u$ of $\mathsf G$. If
$\mathcal V=\{i,j\}$ and $u_1$, $u_2$ are the columns of $\mathsf G$ at $i$ and $j$, then
\[
 \mathsf G\mathsf S_{\mathcal V}\mathsf G^{\mathsf T}=(\mathsf S_{ii}-\mathsf S_{ij})u_1u_1^{\mathsf T}
 +(\mathsf S_{jj}-\mathsf S_{ij})u_2u_2^{\mathsf T}+\mathsf S_{ij}(u_1+u_2)(u_1+u_2)^{\mathsf T},
\]
three terms with coefficients of absolute value at most $\mathfrak f$ and vectors of norm at most $1$. In both cases
\begin{equation}\label{dp:eq:pivot-cost}
 \norm{\mathsf G\mathsf S_{\mathcal V}\mathsf G^{\mathsf T}}_\Gamma\le\tfrac32\abs{\mathcal V}\mathfrak f.
\end{equation}
Let $\mathsf U\in K^{m\times m}$ have the columns of $\mathsf G$ at the indices in $\mathcal V$ and the columns of the
identity matrix $\id$ elsewhere, and let $\mathsf E\in K^{m\times m}$ agree with $\mathsf S$ at the positions in
$\mathcal V\times\mathcal V$ and vanish elsewhere. The nonzero columns of $\mathsf U-\id$ lie at indices in
$\mathcal V$, and its rows in $\mathcal V$ vanish; so $(\mathsf U-\id)^2=0$, $\mathsf U^{-1}=2\id-\mathsf U$ is
integral, and $\mathsf U\mathsf S'=\mathsf S'=\mathsf S'\mathsf U^{\mathsf T}$, since the rows and columns of
$\mathsf S'$ in $\mathcal V$ vanish. As $\mathsf U\mathsf E\mathsf U^{\mathsf T}=\mathsf G\mathsf S_{\mathcal
V}\mathsf G^{\mathsf T}$, this gives $\mathsf S=\mathsf U(\mathsf E+\mathsf S')\mathsf U^{\mathsf T}$, so
$\mathfrak d_k(\mathsf S)\ge\mathfrak d_k(\mathsf E+\mathsf S')$. For $0\le k_1\le\min(k,\abs{\mathcal V})$, a
$k_1\times k_1$ minor of $\mathsf S_{\mathcal V}$ and a $(k-k_1)\times(k-k_1)$ minor of $\mathsf S'$ with rows and
columns outside $\mathcal V$ together form a $k\times k$ minor of $\mathsf E+\mathsf S'$ whose matrix is, after
reordering its rows and columns, block diagonal. As the nonzero minors of $\mathsf S'$ have rows and columns outside
$\mathcal V$, $\mathfrak d_k(\mathsf S)\ge\mathfrak d_{k_1}(\mathsf S_{\mathcal V})\,\mathfrak d_{k-k_1}(\mathsf S')$;
here $\mathfrak d_1(\mathsf S_{\mathcal V})=\mathfrak f$, and $\mathfrak d_2(\mathsf S_{\mathcal V})=\mathfrak f^2$ if
$\abs{\mathcal V}=2$.

\emph{Step 2 (the pivot sequence).} Apply Step~1 to $\mathsf S$, then to $\mathsf S'$ if it is nonzero, and so on. A row
that vanishes stays zero, as the row of $\mathsf G$ at its index, and so that of
$\mathsf G\mathsf S_{\mathcal V}\mathsf G^{\mathsf T}$, vanishes; and the rows in $\mathcal V$, which contain an entry
of absolute value $\mathfrak f$, vanish after the step. So after at most $m$ steps the matrix is zero, and $\mathsf S$
is the sum of the matrices $\mathsf G\mathsf S_{\mathcal V}\mathsf G^{\mathsf T}$ of the steps. List the numbers
$\mathfrak f$ of the steps, each once if $\abs{\mathcal V}=1$ and twice if $\abs{\mathcal V}=2$, as
$\mathfrak f_1,\dots,\mathfrak f_r$. As the entries of $\mathsf S'$ have absolute value at most $\mathfrak f$,
$\mathfrak f_1\ge\dots\ge\mathfrak f_r>0$. By~\eqref{dp:eq:pivot-cost},
$\norm{\mathsf S}_\Gamma\le\frac32(\mathfrak f_1+\dots+\mathfrak f_r)$. By the last sentence of Step~1 and induction on
the number of steps, applied to $\mathsf S'$, whose sequence is $\mathfrak f_{\abs{\mathcal V}+1},\dots,\mathfrak f_r$,
\begin{equation}\label{dp:eq:minor-lower}
 \mathfrak d_k(\mathsf S)\ge\mathfrak f_1\cdots\mathfrak f_k\qquad(1\le k\le r):
\end{equation}
for $k\le\abs{\mathcal V}$ take $k_1=k$, and for $k>\abs{\mathcal V}$ take $k_1=\abs{\mathcal V}$.

\emph{Step 3 (products and sums).} Let $\mathfrak s_1,\dots,\mathfrak s_r>0$ with
$\mathfrak f_1\cdots\mathfrak f_k\le\mathfrak s_1\cdots\mathfrak s_k$ for $1\le k\le r$. Then
$\mathfrak f_1+\dots+\mathfrak f_r\le\mathfrak s_1+\dots+\mathfrak s_r$. Indeed,
$\log(\mathfrak s_k/\mathfrak f_k)\le\mathfrak s_k/\mathfrak f_k-1$, so
$\mathfrak s_k-\mathfrak f_k\ge\mathfrak f_k\log(\mathfrak s_k/\mathfrak f_k)$. Put
$\mathfrak l_k=\sum_{i\le k}\log(\mathfrak s_i/\mathfrak f_i)\ge0$, $\mathfrak l_0=0$ and $\mathfrak f_{r+1}=0$;
summation by parts gives
\[
 \sum_{k=1}^r\mathfrak f_k\log(\mathfrak s_k/\mathfrak f_k)=\sum_{k=1}^r\mathfrak f_k(\mathfrak l_k-\mathfrak l_{k-1})
 =\sum_{k=1}^r(\mathfrak f_k-\mathfrak f_{k+1})\mathfrak l_k\ge0.
\]

\emph{Step 4 (the projective norm).} Let $\mathsf S=\sum_{l=1}^{r'}x_lx_l'^{\mathsf T}$, and let
$\mathfrak s_1\ge\mathfrak s_2\ge\dots\ge\mathfrak s_{r'}$ be the numbers $\norm{x_l}\norm{x'_l}$ in nonincreasing
order. By the Cauchy--Binet formula, a $k\times k$ minor of $\mathsf S$ is the sum, over the sets of $k$ indices $l$, of
the products of the corresponding $k\times k$ minors of the matrices with the columns $x_l$ and $x'_l$, $l$ in the set;
by the ultrametric inequality such a minor of the first matrix has absolute value at most the product of the
$\norm{x_l}$, and likewise for the second. So $\mathfrak d_k(\mathsf S)\le\mathfrak s_1\cdots\mathfrak s_k$ if
$k\le r'$, and $\mathfrak d_k(\mathsf S)=0$ if $k>r'$. By~\eqref{dp:eq:minor-lower}, $r'\ge r$ and
$\mathfrak f_1\cdots\mathfrak f_k\le\mathfrak s_1\cdots\mathfrak s_k$ for $k\le r$, so
$\mathfrak s_1,\dots,\mathfrak s_r>0$, and Step~3 gives
$\mathfrak f_1+\dots+\mathfrak f_r\le\mathfrak s_1+\dots+\mathfrak s_r\le\sum_l\norm{x_l}\norm{x'_l}$. Taking the
infimum over the representations, $\mathfrak f_1+\dots+\mathfrak f_r\le\pi(\mathsf S)$, and Step~2 gives
$\norm{\mathsf S}_\Gamma\le\frac32\pi(\mathsf S)$, also for $\mathsf S=0$.
\end{proof}

\begin{proof}[Proof of Proposition~\ref{dp:prop:ultrametric}]
\emph{Step 1 (item (1)).} Let $\mathfrak a_n(P)<\infty$ and $a\colon P\to W$ be a bounded polynomial law of degree $n$.
By the first assertion of Lemma~\ref{dp:lem:linearization}(1), with $n$ in place of $p$,
$\norm{L_a(t)}\le\norm a\norm t_\Gamma\le\mathfrak a_n(P)\norm a\,\pi(t)$ for $t\in\Gamma^n(P)$; apply
Lemma~\ref{dp:lem:linearization}(2), with $n$ in place of $p$ and $M=\mathfrak a_n(P)\norm a$. Conversely, if every
bounded polynomial law of degree $n$ on $P$ is ambient of degree $n$, so is $a_\Gamma$, and
Lemma~\ref{dp:lem:constants}(1) gives $\mathfrak a_n(P)<\infty$.

\emph{Step 2 (item (2)).} We may assume that $C=\mathfrak a_{n+1}(K\oplus P)$ is finite. The map $a_+(\lambda,h)=\lambda
a_\Gamma(h)$ from $K\oplus P$ to $W_0$ is a bounded polynomial law of degree $n+1$ with $\norm{a_+}\le1$: for finitely
many $(\lambda_l,h_l)\in K\oplus P$ and scalars $s_l$,
$a_+(\sum_ls_l(\lambda_l,h_l))=(\sum_ls_l\lambda_l)\,a_\Gamma(\sum_ls_lh_l)$, and $\norm{a_+(\lambda,h)}\le\abs\lambda\norm
h^n\le\norm{(\lambda,h)}^{n+1}$. By~(1), $ja_+(\lambda,h)=\beta((\lambda,h),\dots,(\lambda,h))$ for a linear isometry $j$
of $W_0$ into a normed space $Z$ and a bounded $(n+1)$\nobreakdash-linear map $\beta\colon(K\oplus P)^{n+1}\to Z$ with
$\norm\beta\le C$. Put $\mathbf 1=(1,0)$ and $h^+=(0,h)$ for $h\in P$. For $h_1,\dots,h_n\in P$ let
$\beta'(h_1,\dots,h_n)$ be the sum, over the $n+1$ slots, of the value of $\beta$ at the
$(n+1)$\nobreakdash-tuple with $\mathbf 1$ in this slot and $h_1^+,\dots,h_n^+$ in the others, in this order. Then
$\beta'$ is $n$\nobreakdash-linear, and $\norm{\beta'}\le(n+1)\norm\beta$, as $\norm{\mathbf 1}=1$ and
$\norm{h^+}=\norm h$. For $h\in P$ and $s\in K$, $\beta((1,sh),\dots,(1,sh))$, with $(1,sh)=\mathbf 1+sh^+$, is a
polynomial in $s$ whose coefficient of $s^n$ is $\beta'(h,\dots,h)$, while $ja_+(1,sh)=s^nja_\Gamma(h)$. As $K$ is infinite,
comparing coefficients gives $ja_\Gamma(h)=\beta'(h,\dots,h)$, and Lemma~\ref{dp:lem:constants}(1) gives
$\mathfrak a_n(P)\le\norm{\beta'}\le(n+1)C$.

\emph{Step 3 (item (3)).} The second inequality follows from Lemma~\ref{dp:lem:constants}(3), applied to the coordinate
projection of $c_0(\N,K)$ onto $c_0([m],K)$. For the first, by Lemma~\ref{dp:lem:constants}(4), and as
$\mathfrak a_n(\{0\})=0$, it suffices to show
$\mathfrak a_n(P_0)\le\mathfrak c_K^{-n}\sup_{m\ge1}\mathfrak a_n(c_0([m],K))$ for every finite-dimensional subspace
$P_0\ne\{0\}$ of $P$. Its completion $\widehat P_0$ is a nonarchimedean normed space over $\Kh$. The $\Kh$\nobreakdash-span of
a basis of $P_0$ in $\widehat P_0$ is complete by Lemma~\ref{lem:findim}, hence closed, and it contains the dense subspace $P_0$;
so it is $\widehat P_0$, whose dimension $m'$ over $\Kh$ satisfies $1\le m'<\infty$. If $\lambda_k\to\lambda$ in $\Kh$ with
$\lambda_k\in K$ and $\lambda\ne0$, then $\abs{\lambda_k-\lambda}<\abs\lambda$ and so $\abs{\lambda_k}=\abs\lambda$ for large $k$; hence $\abs{\Kh^\times}=\abs{K^\times}$,
and $\mathfrak c_{\Kh}=\mathfrak c_K$. Lemma~\ref{dp:lem:completion}(1), Lemma~\ref{dp:lem:orthogonal} over $\Kh$ and
Lemma~\ref{dp:lem:completion}(2) give
\[
 \mathfrak a_n(P_0)\le\mathfrak a_n(\widehat P_0)\le\mathfrak c_K^{-n}\,\mathfrak a_n(c_0([m'],\Kh))
 \le\mathfrak c_K^{-n}\,\mathfrak a_n(c_0([m'],K)),
\]
where the second and third constants are computed over $\Kh$.

\emph{Step 4 (item (4)).} Apply item~(3) with $n=2$ and Lemma~\ref{dp:lem:degree-two}.
\end{proof}
%% ---- end app-Q-sec10.tex
%% ---- begin bib.tex
\begin{thebibliography}{99}
\bibitem{AAW} F.~Albiac, J.~L.~Ansorena and B.~Wallis, \emph{Garling sequence spaces}, J.\ London Math.\ Soc.\ (2018),
doi:10.1112/jlms.12129; arXiv:1703.07772.

\bibitem{BGN} W.~Bertram, H.~Gl\"ockner and K.-H.~Neeb,
\emph{Differential calculus, manifolds and Lie groups over arbitrary infinite fields},
arXiv:math/0303300 (2003); published as \emph{Differential calculus over general base fields and rings}, Expo.\ Math.\ \textbf{22} (2004), 213--282,\\
{\small\url{https://arxiv.org/abs/math/0303300}}. We cite the numbering of the arXiv version.

\bibitem{Bou} N.~Bourbaki, \emph{Vari\'et\'es diff\'erentielles et analytiques. Fascicule de r\'esultats}, Hermann, Paris,
1967 (\S\S1--7) and 1971 (\S\S8--15); reprinted by Springer, 2007. We cite the reprint.

\bibitem{Br} K.~S.~Brown, \emph{Cohomology of Groups}, Graduate Texts in Mathematics 87, Springer, New York, 1982.

\bibitem{CSC} F.~Cabello S\'anchez and J.~M.~F.~Castillo, \emph{Homological Methods in Banach Space Theory},
Cambridge Studies in Advanced Mathematics 203, Cambridge University Press, Cambridge, 2023.

\bibitem{Co} I.~S.~Cohen, \emph{On the structure and ideal theory of complete local rings}, Trans.\ Amer.\ Math.\ Soc.\
\textbf{59} (1946), 54--106.

\bibitem{Ga} D.~J.~H.~Garling, \emph{Symmetric bases of locally convex spaces}, Studia Math.\ \textbf{30} (1968),
163--181.

\bibitem{G1} S.~Gou\"ezel, \emph{In characteristic $0$, precomposition by a linear map is analytic on the space of
alternating maps}, pull request \#43338 to Mathlib, merged 30 September 2026,\\ {\small\url{https://github.com/leanprover-community/mathlib4/pull/43338}}.

\bibitem{G2} S.~Gou\"ezel, \emph{The bundle of continuous alternating maps is $C^n$}, pull request \#43234 to Mathlib,\\
{\small\url{https://github.com/leanprover-community/mathlib4/pull/43234}}.

\bibitem{Gr} L.~Gruson, \emph{Th\'eorie de Fredholm $p$-adique}, Bull.\ Soc.\ Math.\ France \textbf{94} (1966), 67--95.

\bibitem{Ha} G.~Hansel, \emph{Nombre minimal de contacts de fermeture n\'ecessaires pour r\'ealiser une fonction
bool\'eenne sym\'etrique de $n$ variables}, C.~R.\ Acad.\ Sci.\ Paris \textbf{258} (1964), 6037--6040.

\bibitem{In} A.~W.~Ingleton, \emph{The Hahn--Banach theorem for non-Archimedean-valued fields}, Proc.\ Cambridge Philos.\
Soc.\ \textbf{48} (1952), 41--45.

\bibitem{Ka} N.~J.~Kalton, \emph{The three space problem for locally bounded F-spaces}, Compositio Math.\ \textbf{37}
(1978), 243--276.

\bibitem{KP} N.~J.~Kalton and N.~T.~Peck, \emph{Twisted sums of sequence spaces and the three space problem}, Trans.\
Amer.\ Math.\ Soc.\ \textbf{255} (1979), 1--30.

\bibitem{KR} N.~J.~Kalton and J.~W.~Roberts, \emph{A rigid subspace of $L_0$}, Trans.\ Amer.\ Math.\ Soc.\ \textbf{266}
(1981), 645--654.

\bibitem{La} S.~Lang, \emph{Fundamentals of Differential Geometry}, Graduate Texts in Mathematics 191, Springer, 1999.

\bibitem{docs} The mathlib Community, \emph{Mathlib documentation}, \texttt{docs\#ContDiffWithinAt},\\
{\footnotesize\url{https://leanprover-community.github.io/mathlib4_docs/find/?pattern=ContDiffWithinAt#doc}}.

\bibitem{LeanP} J.~McCarthy, \emph{Lean formalization of alternating precomposition in positive characteristic},
commit \texttt{863fa6b},\\
{\small\url{https://github.com/Deicyde/alternating-analytic-characteristic-p}}.

\bibitem{PR} J.~McCarthy, \emph{Smoothness reflects along a linear isometry with closed range},\\ pull request \#43548
to Mathlib,\\ {\small\url{https://github.com/leanprover-community/mathlib4/pull/43548}}.

\bibitem{VdP} M.~van der Put, \emph{Espaces de Banach non archim\'ediens}, Bull.\ Soc.\ Math.\ France
\textbf{97} (1969), 309--320,\\
{\small\url{https://www.numdam.org/article/BSMF_1969__97__309_0.pdf}}.

\bibitem{Ra} F.~P.~Ramsey, \emph{On a problem of formal logic}, Proc.\ London Math.\ Soc.\ (2) \textbf{30} (1930), 264--286.

\bibitem{Ro} N.~Roby, \emph{Lois polynomes et lois formelles en th\'eorie des modules}, Ann.\ Sci.\ \'Ecole Norm.\
Sup.\ (3) \textbf{80} (1963), 213--348.

\bibitem{vR} A.~C.~M.~van Rooij, \emph{Non-Archimedean Functional Analysis}, Marcel Dekker, 1978.

\bibitem{Schi} W.~H.~Schikhof, \emph{Banach spaces over nonarchimedean valued fields}, Topology Proc.\
\textbf{24} (1999), 547--581,\\
{\small\url{https://topology.nipissingu.ca/tp/reprints/v24/tp24224.pdf}}.

\bibitem{Sch} P.~Schneider, \emph{Nonarchimedean Functional Analysis}, Springer Monographs in Mathematics,
Springer, Berlin, 2002, \href{https://doi.org/10.1007/978-3-662-04728-6}{doi:10.1007/978-3-662-04728-6}.
Online version of 25 October 2005,\\
{\small\url{https://ivv5hpp.uni-muenster.de/u/pschnei/publ/lectnotes/nfa.pdf}}.

\bibitem{Se} J.-P.~Serre, \emph{Endomorphismes compl\`etement continus des espaces de Banach $p$\nobreakdash-adiques}, Publ.\ Math.\
IH\'ES \textbf{12} (1962), 69--85.

\bibitem{Yuan} Y.~Yuan, \emph{Formalization of non-Archimedean functional analysis 1: spherically complete spaces},\\
arXiv:2601.21734v3 (2026), accepted for publication in the Journal of Symbolic Logic,\\
{\small\url{https://arxiv.org/abs/2601.21734v3}}.

\end{thebibliography}
%% ---- end bib.tex
\end{document}
